\section*{Part D. Continuous sections and the residue completion}

All cohomology in this part is computed by continuous cochains. In
particular, a cokernel of a homomorphism of integral modules is not
silently replaced by its separated quotient. We write
\(C=c_0(\{q:q\ne5\},\mathbb Z_5)\), with its uniform \(5\)-adic
topology, and \(D_j(a)_q=(q^{j-1}-1)a_q\).

\begin{statement}{D.1}{Proved}{Weights three, four and five with unrestricted ramification}
For \(j=3,4,5\), the continuous cohomology groups of the absolute
Galois group of \(\mathbb Q\) are
\[
\begin{array}{c|ccc}
j&3&4&5\\\hline
H^0(\mathbb Q,\mathbb Z_5(j))&0&0&0\\
H^1(\mathbb Q,\mathbb Z_5(j))&\mathbb Z_5&\mathbb Z/5&\mathbb Z_5\\
H^2(\mathbb Q,\mathbb Z_5(j))&C/D_3C&C/D_4C&C/D_5C.
\end{array}
\]
The last row denotes algebraic cokernels. The first-cohomology
classes are unramified away from \(5\). More precisely, for
\(X=\operatorname{Spec}\mathbb Z[1/5]\), localization gives
\[
R\Gamma(X,\mathbb Z_5(j))\longrightarrow
R\Gamma(\mathbb Q,\mathbb Z_5(j))\longrightarrow
[C\xrightarrow{D_j}C][-1]\longrightarrow R\Gamma(X,\mathbb Z_5(j))[1],
\]
where the bracket is in degrees zero and one. The Jannsen kernel is
\[
\varprojlim{}^1_n H^1(\mathbb Q,\mathbb Z/5^n(j))
=\frac{C\cap\prod_{q\ne5}(q^{j-1}-1)\mathbb Z_5}{D_jC}.
\]
It is nonzero in all three weights. In particular Lead 12 is true.
\end{statement}
\begin{proof}
We give the integral calculation, including the small torsion in
weight four. Let \(L=\mathbb Q(\zeta_5)\),
\(A=\operatorname{Gal}(L/\mathbb Q)\), and let \(\omega\) be its
modulo-five cyclotomic character. Minkowski's bound is
\(24\cdot256^{-1}(4/\pi)^2\sqrt{125}<2\), so every ideal class
of \(L\) is principal. There is one prime above five. The
Brauer invariant-sum sequence therefore gives
\(\operatorname{Br}(\mathcal O_L[1/5])[5]=0\): only that prime
can support a five-primary invariant, and its sum must vanish.
Kummer theory now gives
\[
H^2(G_{L,\{5\}},\mathbb F_5(1))=0,\qquad
H^1(G_{L,\{5\}},\mathbb F_5(1))=
\mathcal O_L[1/5]^\times/(\mathcal O_L[1/5]^\times)^5.
\]
As an \(A\)-module the latter is
\(\mathbf1\oplus\omega\oplus\omega^2\). The valuation at the
unique prime gives \(\mathbf1\); roots of unity give \(\omega\);
the free rank-one unit group gives \(\omega^2\), since a generator
of \(A\) sends \((1+\sqrt5)/2\) to the negative of its inverse.
Averaging is exact because \(|A|=4\) is prime to five. Hence, for
\(j=3,4,5\), the dimensions of \(H^1(G_{\mathbb Q,\{5\}},
\mathbb F_5(j))\) are all one, and \(H^2\) is zero.

The finite-coefficient exact sequences imply vanishing of
\(H^2(G_{\mathbb Q,\{5\}},\mathbb Z/5^n(j))\) for every \(n\).
Finite ramification makes the finite-coefficient cohomology finite
\cite[I, Corollary 4.15]{DE_MilneADT}; its inverse systems are
Mittag--Leffler. The continuous comparison
\cite[Proposition 1.6 and Theorem 2.2]{DE_Jannsen1988} gives
\(H^2(X,\mathbb Z_5(j))=0\). It also gives compactness of the
integral \(H^1\). Its quotient modulo five has dimension one, so
compact Nakayama makes it a cyclic finitely generated
\(\mathbb Z_5\)-module. Its torsion is exactly
\(H^0(\mathbb Q,\mathbb Q_5/\mathbb Z_5(j))\), by the coefficient
sequence and the vanishing of rational invariants. The cyclotomic
image is \(\mathbb Z_5^\times\). This invariant group is zero for
\(j=3,5\), and has order five for \(j=4\), since
\(v_5(2^4-1)=1\). The first-cohomology row follows. The
higher cohomology of \(X\) vanishes here; finite-coefficient
vanishing above degree two follows from
\cite[I, Theorem 4.10(c)]{DE_MilneADT}, and the same inverse-limit
argument applies.

We check the full localization map rather than its projections
at a single finite level. For a finite set of primes removed from
\(X\), purity and localization have residue term
\(R\Gamma(\mathbb F_q,\mathbb Z/5^n(j-1))[-1]\).
Wild inertia has order prime to five, and the five-primary tame
inertia is \(\mathbb Z_5(1)\); these facts give the same residue
complex directly. Arithmetic Frobenius acts on its coefficient
by \(q^{j-1}\), so that complex is
\([\mathbb Z/5^n\xrightarrow{q^{j-1}-1}\mathbb Z/5^n]\).
The comparison with Galois cohomology is
\cite[II, Proposition 2.9]{DE_MilneADT}. Taking the filtered union
of finite ramification sets is legitimate: every finite-valued
continuous cochain factors through a finite Galois quotient.
Thus the residue term at level \(n\) is the direct sum over all
\(q\ne5\). Both of its transition maps are surjective. Their
inverse limits are exactly \(C\), not the unrestricted product.
This proves the displayed triangle. Since every
\(q^{j-1}-1\ne0\), its degree-one differential is injective.
The already computed global term now gives both the claimed
\(H^1\) and \(H^2\).

For the derived-limit formula, reduction modulo \(5^n\) of the
residue complex has first cohomology
\(P_n=\bigoplus_q\ker(D_j:\mathbb Z/5^n\to\mathbb Z/5^n)\).
Because the global \(H^2\) at finite ramification is zero,
localization and finite global \(H^1\) identify
\(\varprojlim{}^1H^1(\mathbb Q,\mathbb Z/5^n(j))\) with
\(\varprojlim{}^1P_n\). The image towers
\(D_j\bigoplus_q\mathbb Z/5^n\) are surjective. Applying inverse
limits to their kernel sequences identifies their inverse limit
with \(C\cap\prod_q D_j\mathbb Z_5\), giving exactly the formula
in the statement. This is also the kernel of reduction to
\(\varprojlim_nH^2(\mathbb Q,\mathbb Z/5^n(j))\).
Choose distinct primes \(q_r\equiv1\pmod{5^r}\). The family
\(z_{q_r}=q_r^{j-1}-1\), with all other coordinates zero,
belongs to \(C\), is locally divisible by the displayed
multipliers, and has unique divided coordinates equal to one
at all \(q_r\). Those coordinates do not tend to zero.
Thus its class is nonzero. Dirichlet's theorem supplies the
primes; the failure is the convergence condition itself.
\end{proof}



\begin{statement}{D.20}{Proved}{Full metabelian Fermat families: the primewise cocycles really embed}
Let \(\ell\in\{3,5\}\), \(n=\ell^r\), \(r\ge1\), and
let \(U_n\) be the complement of the coordinate axes in
\(X_n:X^n+Y^n+Z^n=0\). Write
\(N=\pi_1(U_{n,\overline{\mathbb Q}})\), and use a rational
unit tangent at \(O=[0:-1:1]\) to specify the action.
There is a continuous injective map
\[
 \mathbb Z_\ell^{d_n}\longrightarrow
 \operatorname{Sec}(N/N''),\qquad
 d_n=\frac{(n^2-1)(n^2+3)}4,                         \tag{D.20.1}
\]
whose members have zero abelian cocycle, whose zero member
is the chosen cusp section, and whose every nonzero member
is noncuspidal. This includes the sections of the cusp
normalizers in the full metabelian quotient itself.
In particular \(d_3=24\) and \(d_5=168\).
Thus the full metabelian existence assertions LM6 E.33
and E.34 are valid, without identifying a product of
primewise quotients with the full metabelian group.
\end{statement}
\begin{proof}
We first prove the comparison of derived modules, including
its arithmetic equivariance. For any free profinite group
\(N\) of finite rank \(h\), put
\[
 A=N^{\mathrm{ab}},\quad D=N'/N'',\quad
 P=N^{(\ell)},\quad D_\ell=P'/P'',\quad
 R=\widehat{\mathbb Z}[[A]],\quad
 R_\ell=\mathbb Z_\ell[[P^{\mathrm{ab}}]].
\]
In a free basis with abelian images \(X_i\), the Alexander
sequences are
\[
 0\longrightarrow D\longrightarrow R^h
   \xrightarrow{(r_i)\mapsto\sum r_i(X_i-1)}I_A
   \longrightarrow0,\qquad
 0\longrightarrow D_\ell\longrightarrow R_\ell^h
   \longrightarrow I_{A_\ell}\longrightarrow0.       \tag{D.20.2}
\]
Here is a direct justification of the completed form.
For a finite abelian quotient of the free group, the
covering graph has one vertex for each quotient element
and \(h\) edges at each vertex. Its chain boundary is
the displayed Fox boundary. Pass to all finite abelian
quotients and finite coefficient rings. Its inverse-limit
first homology is \(D\): a finite continuous character
of \(N'\) has kernel containing \(N'\cap V\) for some
open normal \(V\lhd N\), and extends to \(N'V\) by
being trivial on \(V\). This is well-defined because
\([V,N']\subset N'\cap V\). Conversely a finite character
of \(N'V\) trivial on \(N'\) vanishes on a smaller
\(N'W\). Duality for finite abelian groups therefore
identifies the limit with \(N'^{\mathrm{ab}}\).
All chain groups are compact and all finite equations
have compatible solutions, so the inverse limit is exact.
The same proof in the pro-\(\ell\) category gives the
second row of (D.20.2).

In the coefficient-\(\ell\) factor \(\mathbb Z_\ell[[A]]\),
let \(e_\ell\) be the normalized Haar idempotent of
\(A^{(\ell')}\). At a finite quotient this is the
sum of that subgroup divided by its order; the order
is an \(\ell\)-adic unit. These finite sums are compatible.
Extend by zero in the other coefficient factors of \(R\).
Then
\[
 e_\ell R\simeq R_\ell,\qquad
 e_\ell I_A\simeq I_{A_\ell},\qquad
 e_\ell D\xrightarrow[\ q\ ]{\ \sim\ }D_\ell.       \tag{D.20.3}
\]
The last isomorphism follows from the kernels in
(D.20.2), since multiplication by a central idempotent
is exact. Its inverse gives
\(j:D_\ell\hookrightarrow D\) with \(qj=1\).
The idempotent is characteristic under all automorphisms
of \(A\), and \(q\) is natural. Thus \(j\) commutes
with the actual arithmetic action. This proves an
embedding of derived modules, not an embedding of
\(P\) into \(N\).

We next check the arithmetic source used to build the
cocycles. The pair \((X_n,X_n\setminus U_n)\) is smooth
with etale boundary over \(\mathbb Z[1/\ell]\).
Specialization for smooth pairs therefore makes the
tangential action on \(P\) unramified outside \(\ell\);
the tangent is a unit at every other prime. These are
the hypotheses of the tame specialization theorem
\cite[Expos\'e XIII, Corollary 2.9 and \S2.10]{DE_SGA1}.
The Fermat map identifies \(P\) with the kernel of
the map from the pro-\(\ell\) tripod group to
\((\mathbb Z/n)^2\). Indeed a subgroup of finite
\(\ell\)-power index in a free pro-\(\ell\) group
has the pro-\(\ell\) topology induced by that of its
full profinite counterpart: take the core of a finite
\(\ell\)-group quotient, whose finitely many conjugates
again have \(\ell\)-power index. The rational tangents
can be chosen compatibly with the power map. A change
between the two signs of a unit tangent has no
\(\ell\)-primary Kummer class, since \(\ell\) is odd.

On the tripod's Frattini quotient, the action of
\(G_{\mathbb Q(\zeta_\ell)}\) is trivial. The kernel
of the action on that Frattini quotient is pro-\(\ell\),
as in D.16; restriction to the Fermat subgroup has
pro-\(\ell\) image as well. Consequently the action
on \(P\) factors through
\[
 \Gamma_\ell=P_{\rm ar}\rtimes C_{\ell-1},\qquad
 P_{\rm ar}=G_{\mathbb Q(\zeta_\ell),S_\ell}^{(\ell)},
 \quad S_\ell=\{v\mid\ell\}.                         \tag{D.20.4}
\]
To see that \(P_{\rm ar}\) is free, note that
\(\mathbb Q(\zeta_3)\) and \(\mathbb Q(\zeta_5)\)
have class number one and one prime above \(\ell\).
For both fields the Minkowski bound is less than two.
Kummer theory and
the local-invariant description of the Brauer group
give
\(H^2(G_{\mathbb Q(\zeta_\ell),S_\ell},\mathbb F_\ell)=0\):
the class-group term is zero, and the sum map on
the invariants at the unique finite prime is injective.
The maximal pro-\(\ell\) quotient has vanishing
\(H^2(-,\mathbb F_\ell)\), by the kernel-invariant
argument used in D.16. The minimal-presentation
argument there proves that it is free. Its generator
rank, from the \(S_\ell\)-units modulo \(\ell\), is
\(2\) for \(\ell=3\) and \(3\) for \(\ell=5\).
The extension by \(C_{\ell-1}\) splits by the
profinite coprime splitting theorem, applied to its
finite quotients and then compactness. This is a
complement in the arithmetic quotient, not in
\(G_{\mathbb Q}\).

Set \(V=P^{\mathrm{ab}}\) and
\(B=\gamma_2P/\gamma_3P=\bigwedge^2 V\).
The map
\[
 H^1_{\rm cts}(\Gamma_\ell,D_\ell)
       \longrightarrow H^1_{\rm cts}(\Gamma_\ell,B)
                                                               \tag{D.20.5}
\]
is surjective on actual cocycles after normalization
to zero on \(C_{\ell-1}\). For a normalized cocycle
on the right, lift its values on a free basis of
\(P_{\rm ar}\) arbitrarily to \(D_\ell\).
The resulting homomorphism into the pro-\(\ell\)
semidirect product gives a continuous cocycle.
Average this cocycle under \(C_{\ell-1}\), dividing
by \(\ell-1\). Its projection is unchanged, because
the original cocycle was invariant. Extend the
averaged cocycle by zero on the complement.
This proves (D.20.5) without a nonabelian averaging
operation or a choice of unrelated finite-level lifts.

Inflation identifies the target in (D.20.5) with
\(H^1(G_{\mathbb Q,\{\ell,\infty\}},B)\).
Indeed over \(\mathbb Q(\zeta_\ell)\) the graph of
any cocycle is in a pro-\(\ell\) semidirect product,
so it factors through \(P_{\rm ar}\). Every finite
\(\ell\)-primary module occurring here has, over
that field, a filtration with trivial residual
\(\mathbb F_\ell\) factors. The preceding vanishing
of \(H^2\), induction along the filtration, and
averaging by \(C_{\ell-1}\) give zero finite-coefficient
\(H^2\). Finite \(S\)-ramified \(H^1\) groups are
finite, so their inverse systems are Mittag--Leffler;
Jannsen's sequence then gives zero continuous
\(H^2\) with the lattice \(B\). Its \(H^0\) is zero:
at a prime of good reduction different from \(\ell\),
the compact and boundary eigenvalues in \(V\) have
weights one and two, so those in \(B\) have positive
weights. The global Euler formula
\cite[Chapter I, \S4]{DE_MilneADT} now yields
\[
 \operatorname{rank}_{\mathbb Z_\ell}H^1(\Gamma_\ell,B)
       =\dim_{\mathbb Q_\ell}(B\otimes\mathbb Q_\ell)^-.
                                                               \tag{D.20.6}
\]

Here the genus is \(g=(n-1)(n-2)/2\), and the
boundary has \(3n\) points, exactly three real.
Complex conjugation is antisymplectic on compact
homology, so its plus and minus dimensions there
are both \(g\). On the boundary it is minus the
permutation of the cusps; quotienting by the
diagonal removes one minus eigenvalue. Hence
\[
 \dim V^+=\frac{n^2-1}{2},\qquad
 \dim V^-=\frac{n^2+3}{2},\qquad
 \dim(\bigwedge^2V)^-=d_n.                            \tag{D.20.7}
\]
These dimensions are also checked by the executed
integer-matrix certificate
\path{code/de_fermat_ranks.py} at \(n=3,5,9,25\).

Choose a free submodule of rank \(d_n\) in the
cohomology group (D.20.6), and representative
cocycles. Lift its basis by (D.20.5), obtaining
\(b_1,\ldots,b_{d_n}\). For
\(a=(a_i)\in\mathbb Z_\ell^{d_n}\), put
\[
 s_a(g)=j\left(\sum_i a_i b_i(g)\right)s_O(g).        \tag{D.20.8}
\]
The derived module is abelian, so this is an actual
continuous section. Dependence on \((a,g)\) is
jointly continuous. A conjugator between two such
sections has abelian image in
\(H^0(G_{\mathbb Q},N^{\mathrm{ab}})=0\), again by
the positive weights, coefficient prime by
coefficient prime. It therefore lies in \(D\).
Projecting it to \(B\) shows that the difference
of the parameter classes is a coboundary.
The chosen free submodule proves injectivity.

It remains to verify the word ``cuspidal'' in the
full quotient. A cusp loop \(x\) is primitive
in a free basis, obtained by omitting a different
boundary loop from the surface presentation.
If an element of \(N/N''\) centralizes \(x\),
its abelian image, by the class-two commutator
pairing, is a multiple of \(x\). Subtract that
power. The remaining element is in \(D\), is
annihilated by \(X-1\), and embeds in \(R^h\)
by (D.20.2). Multiplication by \(X-1\) is injective
on \(R\): in its coefficient-\(q\) factor, refine
a finite abelian quotient by arbitrarily large
\(q\)-power order in the primitive \(X\) direction.
An \(X\)-invariant coefficient pushes forward as
a multiple of that \(q\)-power, hence becomes
zero modulo every fixed coefficient modulus.
Thus the centralizer is precisely \(\langle x\rangle\).
Its normalizer is the same, since a conjugator
cannot change the primitive abelian vector.
The arithmetic normalizer of a rational cusp
therefore has exactly its usual Kummer packet.
For a nonrational cusp, the arithmetic image
stabilizes its boundary line, so is a proper
subgroup of \(G_{\mathbb Q}\).

For a rational cusp \(d\ne O\), the ratio of
the tangent-line equations has divisor \(n(d-O)\).
Its Abel--Jacobi class is nonzero: a function
with divisor \(d-O\) would give a degree-one
map to the projective line, contrary to \(g\ge1\).
It is therefore nonzero \(\ell\)-primary torsion,
and its \(\ell\)-adic Kummer class is nonzero
by Mordell--Weil. No section in that cusp packet
has zero abelian class. At \(O\), the cusp line
\(I_O\simeq\widehat{\mathbb Z}(1)\) is primitive,
and \(H^0(G_{\mathbb Q},A/I_O)=0\) by the same
weights. Thus \(H^1(G_{\mathbb Q},I_O)\to
H^1(G_{\mathbb Q},A)\) is injective: the only
zero-abelian member of its packet is \(s_O\).
Every nonzero member of (D.20.8) is consequently
noncuspidal, as asserted.
\end{proof}

\begin{statement}{D.21}{Proved}{Nongeometric sections of the full metabelian tripod}
For the tripod over \(\mathbb Q\), the full profinite
metabelian arithmetic quotient admits a continuous
injective family of section classes parametrized by
\(\mathbb Z_5^2\). Every member has zero abelian and
zero class-two cocycle. No member is a rational-point
class, and all but at most three are noncuspidal.
In particular the full metabelian quotient, as well
as its pro-five quotient, has sections that are
neither cuspidal nor induced by rational points.
\end{statement}
\begin{proof}
Write \(P=F^{(5)}\), \(D_5=P'/P''\), and
\(C=\gamma_3P/P''\subset D_5\). This is a closed
abelian pro-five module with a surjection
\[
 C\longrightarrow\gamma_3P/\gamma_4P
             \simeq\mathbb Z_5(3)^2.                 \tag{D.21.1}
\]
The actual action factors through the group
\(\Gamma_5\) of (D.20.4). The free-generator
lifting and averaging proof of (D.20.5) works
word for word for (D.21.1). Its target cohomology
is \(\mathbb Z_5^2\) by D.1, since all the
weight-three classes are unramified outside five.
Choose lifts of a basis to continuous \(C\)-cocycles.
The equivariant Haar embedding
\(j:D_5\hookrightarrow F'/F''\) from (D.20.3)
turns their linear combinations into sections
of the full metabelian quotient. Their class-two
images and their abelian images are zero.

To check conjugacy, a conjugator first has
abelian image in \(H^0(G_{\mathbb Q},F^{\rm ab})=0\),
so lies in \(F'/F''\). Its pro-five degree-two
coordinate is in \(H^0(G_{\mathbb Q},\mathbb Z_5(2))=0\),
because the two sections agree through class two.
It therefore has its pro-five image in \(C\).
Projection to degree three then changes the
parameter only by a coboundary. The chosen
cohomology basis proves injectivity.

A rational point has abelian Kummer coordinates
\((x,1-x)\). The map
\(\mathbb Q^\times\to\varprojlim_m
\mathbb Q^\times/(\mathbb Q^\times)^m\) is injective:
valuations detect a nonunit rational number,
and the only remaining nonidentity element
\(-1\) is detected at the powers of two.
Thus both abelian coordinates cannot vanish
for a rational point of the tripod.

Each of the three cusp lines is primitive
in \(F^{\rm ab}=\widehat{\mathbb Z}(1)^2\),
and its quotient has zero \(G_{\mathbb Q}\)-invariants.
Consequently each cusp packet has at most one
zero-abelian class. The primitive-loop
normalizer calculation in D.20 applies in
rank two as well, so it misses no extra
normalizer section in this quotient. At most
three members of the injective family can
be cuspidal. All the other members give the
claimed nongeometric sections.
\end{proof}

\begin{statement}{D.11}{Proved}{The finite-image criterion of Lead 20}
Let \(L\) be a number field, \(S\) finite containing the infinite
places and those above \(\ell\), and \(M\) a finite free
\(\mathbb Z_\ell\)-module unramified outside \(S\). Then
\[
\varprojlim{}^1_nH^1(G_L,M/\ell^nM)=0
\quad\Longleftrightarrow\quad
\rho_{M(-1)}(G_L)\text{ is finite}.
\]
Equivalently, the valuations of the nonzero Smith entries of
\(\operatorname{Frob}_v-1\) on \(M(-1)\) are uniformly bounded
for \(v\notin S\). Zero entries are excluded from this test.
\end{statement}
\begin{proof}
Write \(A_v=M(-1)\), \(D_v=\operatorname{Frob}_v-1\), and
\(R=c_0(\{v\notin S\},A_v)\). The same finite-coefficient
localization proof as above gives a residue complex
\([R\xrightarrow D R][-1]\). At level \(n\) put
\(P_n=\bigoplus_v\ker(D_v:A_v/\ell^n\to A_v/\ell^n)\), and
let \(E_n\) be its Gysin image in \(H^2(G_{L,S},M/\ell^n)\).
The finite groups at fixed ramification satisfy the
Mittag--Leffler condition. The two kernel sequences therefore
produce
\[
0\longrightarrow
\operatorname{coker}(\varprojlim P_n\to\varprojlim E_n)
\longrightarrow\varprojlim{}^1H^1(G_L,M/\ell^n)
\longrightarrow
\frac{R\cap\prod_vD_vA_v}{DR}\longrightarrow0.                 \tag{D.11.1}
\]
In particular one must check the left term; finiteness of each
\(E_n\) alone does not eliminate it.

A nonzero Smith entry \(\ell^e\) contributes the kernel
\(\ell^{\max(0,n-e)}\mathbb Z/\ell^n\). Its image from level
\(m\) into level \(n\le m\) is generated by
\(\ell^{\max(0,m-e)}\). A zero Smith entry contributes a
surjective tower. A uniform bound \(e\le b\) therefore makes
all nonzero-entry contributions vanish after a shift by \(b\).
If instead the valuations are unbounded, choose distinct places
\(v_r\) with a Smith entry of valuation at least \(r\), and
choose an integral source basis vector \(u_r\) for that entry.
The family \(D_{v_r}u_r\) belongs to \(R\); any preimage has
nonzero reduction in the corresponding source coordinate modulo
\(\ker D_{v_r}+\ell A_{v_r}\). No preimage lies in \(R\).
The right term of (D.11.1) is then nonzero.

Suppose the bound \(b\) exists. Put
\(P_n^0=\bigoplus_v\ker D_v/\ell^n\ker D_v\).
The kernels are saturated, and the tower \(P_n/P_n^0\) is zero
after a shift by \(b\). The integral global module
\(C_S=H^2_{\rm cts}(G_{L,S},M)\) is finitely generated: it is
the compact inverse limit of finite cohomology groups, its
quotient modulo \(\ell\) is finite by the coefficient sequence,
and compact Nakayama applies. The submodule \(J\subset C_S\)
generated by all integral Gysin images is finitely generated
and closed. A finite collection of these images generates it.
The images of \(P_n^0\) are precisely the reductions of \(J\).
Since the difference from \(E_n\) is zero after the fixed shift,
\(\varprojlim E_n=J\). Lifting the finitely many chosen
integral generators to \(\varprojlim P_n\) proves surjectivity
onto \(J\), so the left term of (D.11.1) vanishes as well.
This proves the equivalence with the Smith test, including
\(\ell=2\); possible real-place groups were retained in \(C_S\).

A finite image gives only finitely many possible matrices
\(D_v\), so the test is bounded. Conversely let
\(H=\rho_{M(-1)}(G_L)\subset\operatorname{GL}_d(\mathbb Z_\ell)\)
be infinite. For every \(r\) its congruence kernel
\(H\cap(1+\ell^r\operatorname{Mat}_d(\mathbb Z_\ell))\)
contains a nonidentity element \(g\); otherwise \(H\) would
embed into a finite matrix group. Choose \(k\ge r\) for which
\(g\not\equiv1\pmod{\ell^{k+1}}\). Chebotarev in the finite
extension cut out modulo \(\ell^{k+1}\) supplies a prime with
Frobenius conjugate to \(g\) at that precision. Its difference
from one is nonzero and all its entries are divisible by
\(\ell^r\). Every nonzero Smith entry is therefore divisible
by \(\ell^r\). This proves unboundedness and the assertion.
\end{proof}

\begin{statement}{D.8}{Proved}{All section fibres of the group of order seventy-five}
Let \(E\) be the smooth projective curve \(y^3=x(1-x)\), with
origin the point over \(x=0\), and use the unit tangent of the
tripod at zero. The quotient \(Q_{75}\) in Lead 18 has
\[
Q_{75}\simeq E[5]\rtimes\mu_3,
\]
as a geometric group with arithmetic action. The \(\mu_3\)-action
on the two-dimensional \(\mathbb F_5\)-module \(E[5]\) has
minimal polynomial \(T^2+T+1\). The image in \(Q_{75}\) of the
kernel of all maximal pro-prime quotients of the full geometric
tripod group is exactly \(E[5]\).

The map on continuous section classes is onto
\[
\operatorname{Sec}(Q_{75})\longrightarrow
H^1(\mathbb Q,\mu_3)=\mathbb Q^\times/\mathbb Q^{\times3}.
\]
For a representative \(c\in\mathbb Q^\times\), write
\(E_c:y^3=cx(1-x)\), with its point over zero as origin.
The fibre over \([c]\) is a torsor under
\(H^1(\mathbb Q,E_c[5])\); the origin section gives a preferred
zero after the twist identification has been chosen. Its
obstruction in \(H^2(\mathbb Q,E_c[5])\) is zero.
The group \(H^1(\mathbb Q,E_c[5])\) is nonzero, indeed infinite.
\end{statement}
\begin{proof}
The cyclic covering obtained by adjoining a cube root of
\(x(1-x)\) is totally ramified at \(0,1,\infty\).
Riemann--Hurwitz gives genus one, and all three ramification
points are rational. Killing their inertia gives the geometric
fundamental group of \(E\); quotienting it modulo five gives
\(E[5]\). The order-three deck transformation fixes the origin
and is the multiplication automorphism by a primitive cube
root of unity. As \(5\equiv-1\pmod3\), its action on \(E[5]\)
has the stated irreducible polynomial. The geometric orbifold
cover \([E/\mu_3]\) now gives the semidirect product in the
statement. Its rational fixed origin gives the corresponding
arithmetic splitting. The unit tangent at zero has the same
image: locally \(y^3=x+O(x^2)\), whose leading coefficient is
one, so there is no extra Kummer twist. A different cubic-root
choice conjugates the comparison by \(\mu_3\).

Every nilpotent quotient of \(Q_{75}\) kills \(E[5]\). Indeed
Sylow subgroups in a finite nilpotent group commute, whereas
\(\tau-1\) is invertible on \(E[5]\). Its maximal nilpotent
quotient is consequently \(\mu_3\). If \(N\) is the kernel of
the product of all maximal pro-prime quotients of the free
profinite group, its image has quotient nilpotent, and so
contains \(E[5]\). Conversely its image lies in \(E[5]\),
because the quotient \(\mu_3\) is a pro-prime quotient. This
proves equality.

Represent \([c]\) by the Kummer cocycle
\(a_g=g(c^{1/3})/c^{1/3}\). Twisting the action on \(E[5]\) by
\(a_g\) is precisely the action on \(E_c[5]\), using the
isomorphism \(E\to E_c\) over \(\overline{\mathbb Q}\) given
by \((x,y)\mapsto(x,c^{1/3}y)\). Its origin is rational.
The zero translation together with \(a_g\) is therefore a
section of the semidirect product, so the obstruction vanishes.
All other lifts are its products with continuous
\(E_c[5]\)-cocycles, and translation conjugacy divides by
coboundaries. The additional stabilizer that occurs in a
general nonabelian fibre formula is
\(H^0(\mathbb Q,\mu_3)=\mu_3(\mathbb Q)=1\). It causes no
further quotient here. This proves the torsor assertion for
actual global cohomology, without Selmer local conditions.

For infinitude, let \(L/\mathbb Q\) be a finite Galois
extension splitting \(E_c[5]\) and \(\mu_5\), and let \(S\)
contain all its ramified primes and five. At every rational
prime splitting completely in \(L\), the residue group
\(H^0(\mathbb F_q,E_c[5](-1))\) has dimension two.
There are infinitely many such primes by Chebotarev.
Localization has an exact segment
\[
H^1(\mathbb Q,E_c[5])\longrightarrow
\bigoplus_{q\notin S}H^0(\mathbb F_q,E_c[5](-1))
\longrightarrow H^2(G_{\mathbb Q,S},E_c[5]).
\]
The last group is finite. The middle group has infinite
dimension over \(\mathbb F_5\), so the first group is infinite.
Finally \(H^i(\mu_3,E[5])=0\) for \(i>0\), by averaging over
three. Thus the automorphism complement problem uses coprime
cohomology, whereas the section fibre uses the nonzero global
cohomology just computed. No assertion about lifting these
sections to the full geometric group follows from this finite
calculation.
\end{proof}

\begin{statement}{D.12}{Proved}{Nongeometric sections in the pro-five metabelian tripod}
Let \(G_5=G_{\mathbb Q,\{5,\infty\}}\), acting through the
unit tangent at zero. Every continuous section of \(T_n\)
over \(G_5\) lifts compatibly to all \(T_m\), \(m\ge n\),
and to \(F^{(5)}\). Over the zero abelian section the
class-two section is unique and its class-three fibre is
\(\mathbb Z_5^2\). Its entire fibre lifts to
\(F^{(5)}/(F^{(5)})''\). Infinitely many of these metabelian
section classes are neither cuspidal nor induced by a
rational point of the tripod.
\end{statement}
\begin{proof}
The calculation with \(\mathbb Q(\zeta_5)\) in D.1 proves
\(H^2(G_5,\mathbb F_5(j))=0\) for every integer \(j\),
not only the three displayed there: over that field all
modulo-five Tate twists have the same zero second cohomology,
and averaging over its order-four Galois group is exact.
Coefficient induction, finite cohomology and Jannsen's
sequence therefore prove
\[
 H^2_{\rm cts}(G_5,\mathbb Z_5(j))=0
 \quad\text{for every }j.                              \tag{D.12.1}
\]
The same unit-character calculation gives
\(H^1(G_5,\mathbb Z_5(2))=0\) and
\(H^1(G_5,\mathbb Z_5(3))=\mathbb Z_5\).

The tripod and its unit tangential base point have good
integral reduction outside five. Prime-to-residue-characteristic
specialization consequently factors the pro-five arithmetic
action through \(G_5\). The lower-central layer
\(\gamma_mF^{(5)}/\gamma_{m+1}F^{(5)}\) is finite free
of Witt rank, and its arithmetic action is \(\chi_5^m\):
this follows by taking brackets from the two cyclotomic
abelian generators. Each central lifting obstruction is
therefore in a finite direct sum of groups (D.12.1), and
vanishes. Each nonempty fibre is a torsor under the
corresponding \(H^1\), with no stabilizer quotient. Indeed
a fixed element for any cocycle-twisted action has a first
nonzero central graded term, and that term would be fixed
by \(\chi_5^m\). Some element of \(G_5\) has cyclotomic
value in \(1+5\mathbb Z_5\) of infinite order, so no such
nonzero term is fixed in a torsion-free lattice.
This proves the class-two and class-three assertions.

To pass to the limit, choose successive cocycle lifts, not
just unrelated conjugacy classes. A lift whose projection
is conjugate to the chosen preceding cocycle can be
conjugated to project to it exactly, because the quotient
map on the geometric groups is onto. The resulting
compatible continuous cocycles define a continuous cocycle
into the inverse limit. Every finite quotient of a pro-five
group is nilpotent. Hence
\(F^{(5)}=\varprojlim_nT_n\), with its profinite topology.
This proves existence of the full pro-five lift. Its
metabelian image still records the class-three parameter,
because \((F^{(5)})''\subset\gamma_4F^{(5)}\).

For a rational point \(x\), the abelian coordinates are
the Kummer classes of \(x\) and \(1-x\).
The kernel of
\(\mathbb Q^\times\to H^1(\mathbb Q,\mathbb Z_5(1))\)
is \(\{1,-1\}\), as follows by taking all prime valuations.
There is no \(x\ne0,1\) for which both \(x\) and \(1-x\)
belong to this set. Thus no rational-point section has
zero abelian coordinates. In each of the three cusp
packets the parameter is
\(H^1(\mathbb Q,\mathbb Z_5(1))\). Its abelian image is
multiplication by the primitive vectors
\((1,0),(0,1),(-1,-1)\), respectively; the unit tangent
has zero image. Each packet therefore contributes at most
one zero-abelian section class. Removing these at most
three points from the infinite class-three fibre leaves
infinitely many classes, and all of them have the required
full pro-five and metabelian lifts. This assertion concerns
these pro-five quotients; it asserts no full profinite lift.
\end{proof}

\begin{statement}{D.15}{Proved}{Cyclotomic descent and the nonseparated coinvariants}
Put \(F_\infty=\mathbb Q(\mu_{5^\infty})\),
\(\Delta=G_{F_\infty}\), and
\(\Gamma=\operatorname{Gal}(F_\infty/\mathbb Q)
=C_4\times\Gamma_1\). Then \(\operatorname{cd}_5\Delta\le1\).
For a finite five-primary \(G_{\mathbb Q}\)-module \(A\),
\[
 H^2(G_{\mathbb Q},A)
 =H^1(\Gamma,H^1(\Delta,A)).                            \tag{D.15.1}
\]
For a finite free continuous \(\mathbb Z_5\)-representation
\(M\), set
\(J=\varprojlim_n H^1(\Delta,M/5^nM)\), with the inverse
limit of the discrete topologies, and choose a topological
generator \(\gamma\) of \(\Gamma_1\). Then
\[
 H^2_{\rm cts}(G_{\mathbb Q},M)
 =\frac{J^{C_4}}{(\gamma-1)J^{C_4}},\qquad
 \varprojlim{}^1_nH^1(G_{\mathbb Q},M/5^nM)
 =\frac{\overline{(\gamma-1)J^{C_4}}}
 {(\gamma-1)J^{C_4}}.                                  \tag{D.15.2}
\]
In particular the unclosed image in this formula is essential.
All lower-central lifting problems for \(T_n\) over
\(F_\infty\) are unobstructed.
\end{statement}
\begin{proof}
Let \(E/F_\infty\) be finite. A five-primary Brauer class
over \(E\) descends to a finite number field \(E_0\)
inside \(E\), and its local invariants have finite support.
In the extensions \(E_0\mathbb Q(\mu_{5^r})/E_0\), the
local degrees above every finite prime are unbounded powers
of five, after an initial finite stage. At five this is
cyclotomic ramification. At a prime of residue characteristic
\(q\ne5\), the local extension is unramified and its
Frobenius is a power of the positive integer \(q\) in
\(\mathbb Z_5^\times\). This element has infinite order:
a positive rational integer greater than one cannot be a
root of unity. These unbounded degrees kill all the
finitely many five-primary local invariants at a single
large stage. There are no real places. The global Brauer
invariant sequence proves \(\operatorname{Br}(E)[5]=0\).

Here is the cohomological criterion needed to turn that
calculation into the asserted dimension bound. Since
\(\mu_5\subset F_\infty\), every open subgroup \(U\) of
\(\Delta\) has \(H^2(U,\mathbb F_5)=0\), by Kummer
theory. For any finite \(\Delta\)-module \(A\) killed
by five, restrict to the inverse image of a Sylow-five
subgroup of its finite action group. This has index prime
to five, so restriction on cohomology is injective by
corestriction. On that subgroup \(A\) has a finite
filtration with trivial \(\mathbb F_5\) quotients.
The long exact sequence gives \(H^2(U,A)=0\) by induction.
Induction on the exponent treats every finite five-primary
module. The usual dimension-shifting criterion for
cohomological dimension now gives \(\operatorname{cd}_5
\Delta\le1\). Since \(\operatorname{cd}_5\Gamma=1\),
the finite-coefficient Hochschild--Serre spectral sequence
has only two columns and two rows. Its total degree two
is exactly (D.15.1).

We justify the integral passage rather than invoking an
unqualified Hochschild--Serre theorem for compact modules.
Write \(A_n=M/5^nM\), \(I_n=H^0(\Delta,A_n)\), and
\(J_n=H^1(\Delta,A_n)\). The \(I_n\) are finite.
The transition maps \(J_{n+1}\to J_n\) are onto, because
their obstruction lies in \(H^2(\Delta,M/5M)=0\).
For a discrete five-primary \(\Gamma\)-module \(B\),
continuous cohomology is represented by the two-term complex
\[
 [B^{C_4}\xrightarrow{\gamma-1}B^{C_4}].                \tag{D.15.3}
\]
Taking invariants under \(C_4\) is exact. Thus the towers
of both terms in (D.15.3) for \(B=J_n\) are onto, and
their derived inverse limit is the ordinary two-term
complex with terms \(J^{C_4}\).

At each finite level the truncation of
\(R\Gamma(\Delta,A_n)\) gives a triangle with terms
\(I_n\) and \(J_n[-1]\). Apply \(R\Gamma(\Gamma,-)\)
and then \(R\varprojlim\). The cohomology of
\(R\Gamma(\Gamma,I_n)\) is finite and lies in degrees
zero and one, so its derived inverse limit has no
cohomology above degree one. Consequently total degree
two is exactly the cokernel in (D.15.2). Jannsen's
comparison applies to the original coefficient tower
\cite[Theorem 2.2]{DE_Jannsen1988}.

For the last assertion in (D.15.2), put
\(K_n=(\gamma-1)J_n^{C_4}\). The transition maps on
\(K_n\) are onto. Therefore
\[
 J^{C_4}/(\gamma-1)J^{C_4}
 \longrightarrow\varprojlim_n J_n^{C_4}/K_n
\]
is onto. Its kernel is
\(\varprojlim K_n/(\gamma-1)J^{C_4}\).
Surjectivity of \(J^{C_4}\to J_n^{C_4}\) shows that
\(\varprojlim K_n\) is precisely the closure of the
image. By (D.15.1), the target is
\(\varprojlim_n H^2(G_{\mathbb Q},A_n)\). Jannsen's
exact sequence identifies its kernel with the indicated
\(\varprojlim{}^1\). Finally the same coefficient
argument gives \(H^2_{\rm cts}(\Delta,M)=0\): the finite
second cohomology vanishes and the first-cohomology tower
is onto. Apply this to the finite free central layers
of every \(T_n\).
\end{proof}

\begin{statement}{D.13}{Proved}{All nonempty central lifting fibres through class five}
For unrestricted \(G_{\mathbb Q}\), every nonempty fibre of
\(\operatorname{Sec}(T_n)\to\operatorname{Sec}(T_{n-1})\)
is a torsor under the following group:
\[
\begin{array}{c|ccc}
n&3&4&5\\\hline
\text{fibre group}&\mathbb Z_5^2&(\mathbb Z/5)^3&\mathbb Z_5^6.
\end{array}                                                    \tag{D.13.1}
\]
For a section \(s\) at the preceding stage, its obstruction
lies in \((C/D_nC)^{r_n}\), where
\((r_3,r_4,r_5)=(2,3,6)\). For any representative
\(z=(z_{q,a})\in C^{r_n}\), it vanishes if and only if
\[
 q^{n-1}-1\mid z_{q,a}\quad\hbox{for all }q,a,
 \qquad
 \left(\frac{z_{q,a}}{q^{n-1}-1}\right)_{q\ne5}\in C
                        \quad\hbox{for every }a.               \tag{D.13.2}
\]
The convergence condition is independent of the chosen
representative. This computes the fibres and the exact
form of the convergence test, but not the residue families
of the obstruction for arbitrary input sections.
\end{statement}
\begin{proof}
The central kernel of \(T_n\to T_{n-1}\) is
\(L_n\simeq\mathbb Z_5(n)^{r_n}\), with ranks from the
Witt formula. The Galois action on this graded module is
scalar cyclotomic, and inner twisting by a section has
no effect on it. To see that the action of \(H^1(G,L_n)\)
on a nonempty lifting fibre has no further quotient, it
suffices to check the stabilizer of the underlying
section in \(T_{n-1}\). A nontrivial fixed element has a
first nonzero lower-central coordinate in some free
\(\mathbb Z_5(j)\). The cyclotomic image contains an
integer \(u\in\mathbb Z_5^\times\) with \(u^j\ne1\);
its action cannot fix a nonzero element of this free
module. Thus the stabilizer is trivial. The standard
central cocycle calculation now makes the fibre a
principal homogeneous space under \(H^1(G,L_n)\).
D.1 gives exactly (D.13.1). The obstruction is the class
of \(\widetilde s(g)g(\widetilde s(h))\widetilde s(gh)^{-1}\)
for a continuous Hall-coordinate lift; associativity
makes this a continuous two-cocycle. D.1 identifies its
cohomology with the displayed residue quotient. Each
multiplier is nonzero, so it has a unique coordinatewise
division in \(\mathbb Q_5\). Membership of that divided
family in \(C\) is precisely membership of \(z\) in
\(D_nC^{r_n}\), proving (D.13.2). Changing \(z\) by
\(D_nw\), \(w\in C^{r_n}\), adds \(w\) to the divided
family and preserves its convergence.
\end{proof}

\begin{statement}{D.14}{Proved}{The induced Fermat-kernel section sets and a strict group comparison}
Let \(r\ge1\), \(N=5^r\), and
\[
 K_{n,r}=\ker\bigl(T_n\longrightarrow(\mathbb Z/N(1))^2\bigr).
\]
Its section classes inject into \(\operatorname{Sec}(T_n)\).
Their image is exactly the classes whose two abelian
Kummer coordinates belong to
\(N H^1(\mathbb Q,\mathbb Z_5(1))\).
In contrast, if \(F_5=F^{(5)}\) and
\(K_r=\ker(F_5\to(\mathbb Z/N)^2)\), the natural map
from the full pro-five metabelian completion of the
cover to the subgroup induced from the ambient
metabelian group has nonzero kernel:
\[
 K_r/K_r''\longrightarrow K_r/F_5'',\qquad
 \ker=F_5''/K_r''\ne1.                                 \tag{D.14.1}
\]
Thus the two groups in the comparison of Task D.6 are
different. Equality of their arithmetic section sets
does not follow from identification of the induced
Fermat kernels.
\end{statement}
\begin{proof}
The cyclotomic image modulo \(N\) contains two. Hence
\(H^0(\mathbb Q,(\mathbb Z/N(1))^2)=0\).
In the nonabelian cocycle sequence of
\(1\to K_{n,r}\to T_n\to(\mathbb Z/N(1))^2\to1\),
this gives injectivity of the map on \(H^1\), and its
image is the inverse image of zero in the abelian
\(H^1\). The coefficient exact sequence identifies
that kernel with divisibility by \(N\) of both
integral Kummer coordinates. This proves the first
assertion without a finite-ramification assumption.

For strictness of (D.14.1), let \(c=[x,y]\) and
\(c^x=xcx^{-1}\). Both lie in \(K_r\). Their classes
in \(K_r^{\rm ab}\) are linearly independent over
\(\mathbb Z_5\). One can see this in the finite cover
of the two-petal graph: with deck group
\(A=(\mathbb Z/N)^2\), its edge-chain module is
\(\mathbb Z_5[A]^2\), and the cycle of \(c\) is
\((1-Y,X-1)\), up to a common sign and monomial.
The cycle of \(c^x\) is its translate by \(X\).
The first coordinates of a putative relation are
\((a+bX)(1-Y)\). Since \(N\ge5\), their four support
positions are distinct, forcing \(a=b=0\).
The group \(K_r\) is free pro-five by the Schreier
theorem. Its degree-two lower-central quotient is
\(\bigwedge^2K_r^{\rm ab}\), so
\([c,c^x]\) has nonzero class in
\(\gamma_2K_r/\gamma_3K_r\). But
\(K_r''\subset\gamma_4K_r\subset\gamma_3K_r\),
whereas \([c,c^x]\in F_5''\). This gives a concrete
nontrivial element of the kernel in (D.14.1).
\end{proof}

\begin{statement}{D.17}{Proved}{A finite necessary filter on the infinite mixed-prime section fibres}
Fix \(c\in\mathbb Q^\times\), and let \(S_c\) contain
five, the infinite place, and every prime of bad
reduction of \(E_c\). A class in
\(H^1(\mathbb Q,E_c[5])\) that extends compatibly
through \(E_c[5^n]\) for all \(n\) belongs to the
finite subgroup
\[
 H^1(G_{\mathbb Q,S_c},E_c[5])
       \subset H^1(\mathbb Q,E_c[5]).                    \tag{D.17.1}
\]
In particular every section of the relative pro-five
quotient in Task D.9, and every section of the full
geometric group lying above the cubic class \(c\),
satisfies (D.17.1). Infinitely many classes in each
fibre of D.8 fail this necessary condition and hence
have no such lift.
\end{statement}
\begin{proof}
Compatible finite classes define a continuous class in
\(H^1(\mathbb Q,T_5E_c)\). Indeed the possible
\(\varprojlim^1H^0(\mathbb Q,E_c[5^n])\) term vanishes,
since all the groups in that tower are finite. At a
good prime \(q\ne5\), localization identifies the
integral tame residue of such a class with
\[
 H^0(\mathbb F_q,T_5E_c(-1)).                            \tag{D.17.2}
\]
This group is zero. Frobenius on \(T_5E_c\) has
algebraic eigenvalues of complex absolute value
\(q^{1/2}\); after twisting by \(-1\) their absolute
value is \(q^{-1/2}\), so neither eigenvalue is one.
The lattice is torsion-free, which proves the claim
integrally. Wild inertia at \(q\) has no positive
five-primary cohomology. Thus the entire integral
class, and consequently its reduction modulo five,
is unramified at every \(q\notin S_c\). This proves
(D.17.1); finiteness follows from finite ramification
and finite coefficients.

The geometric map from the punctured cubic cover
to its compactification induces the quotient
\(T_5E_c\) of its geometric pro-five group. Every
relative section therefore induces the compatible
classes just considered. A section of the full group
also induces a relative section. The infinite fibre
in D.8 cannot be contained in the finite group
(D.17.1), proving the final assertion. More explicitly,
at primes splitting in
\(\mathbb Q(E_c[5],\mu_5)\) the finite tame-residue
space is \(\mathbb F_5^2\); localization has only a
finite-dimensional global obstruction at a fixed
ramification set. Arbitrarily many independent such
prime residues therefore give classes excluded by
(D.17.1). This exclusion does not determine which
of the finitely many remaining classes pass the
conditions at three, five, the cusps, and all higher
nonabelian stages.
\end{proof}

\begin{statement}{D.18}{Proved}{Which cohomology the three open rows require}
In the mixed quotient of D.8, finite outer-action lifting
and arithmetic section lifting are different problems.
For the finite groups of orders two and three occurring
in the involution and cyclic-deck calculations,
\(H^i(C,E[5])=0\) for \(i>0\). This concerns the
finite-action reading of LM8 A.12. LM8 F.12 concerns
\(H^1\) of an infinite arithmetic Galois group, and
D.8 exhibits its nonzero mixed-prime fibres. LM8 G.20
concerns arithmetic fixed elements and normalizers;
its fixed-element term is an arithmetic \(H^0\), not
finite prime-to-five group cohomology. The averaging
argument for A.12 therefore gives no vanishing for
F.12 or G.20. Conversely the nonzero fibre of D.8
is not a counterexample to G.20.
\end{statement}
\begin{proof}
If \(|C|\) is prime to five, the averaging idempotent
\(|C|^{-1}\sum_{g\in C}g\) makes invariants exact on
\(\mathbb F_5[C]\)-modules, proving the asserted
cohomology vanishing. The mixed information can
nevertheless persist in degree zero, as the
nontrivial idempotent in the order-seventy-five
involution shows. The finite certificate
\texttt{code/de\_q75.py} checks its multiplication,
its peripheral conjugators and its nontrivial action
on the normal subgroup of order twenty-five.
The section comparison is D.8. A fixed element of
an arithmetic kernel is by definition in
\(H^0(\Delta,N)\), with \(\Delta\) the relevant
infinite arithmetic subgroup. Such a group is
neither \(C_2\) nor \(C_3\), and its action has no
finite averaging operator of the above kind.
These identifications prove the distinctions.
For general finite surface actions in A.12 one
must also check that the coefficient prime is
prime to the acting group's order; no universal
averaging vanishing is asserted there.
\end{proof}

\begin{statement}{D.2}{Proved}{The two unrestricted class-three residue families}
D.30 computes the two actual integral residue
families of an arbitrary class-two section:
\[
 R_{A,q}=\frac{q+1}{2}c_qu_q-d_q\alpha_q,
 \qquad
 R_{B,q}=\frac{q+1}{2}c_qv_q-d_q\beta_q.
\]
Here \(u_q,v_q,d_q\) are its normalized
Frobenius coordinates, and
\(c_q=-(u_q\beta_q-v_q\alpha_q)/(q-1)\)
is its integral central inertia coordinate.
D.30 specifies these uniquely from the
Kummer classes by one fixed normalization.
These families represent the global
obstruction in \((C/D_3C)^2\), including
its Jannsen-kernel part. A lift exists
exactly when both families are divisible
by \(q^2-1\) at every prime and their
divided coordinates tend to zero.
Every nonempty fibre is a
\(\mathbb Z_5^2\)-torsor.
\end{statement}

\begin{statement}{D.3}{Proved}{Realizing the nonzero Jannsen kernel as a class-three obstruction}
D.28 constructs an actual continuous class-two
section, with abelian coordinates
\[
 \left(\kappa(3),\ \kappa(-2)+
                 \sum_{r\ge1}5^r\kappa(q_r)\right),
\]
whose obstruction to class three is a nonzero
element of
\(\varprojlim{}^1_mH^1(\mathbb Q,\mathbb Z/5^m(3))^2\).
The primes \(q_r\) are chosen in explicitly
specified nonempty Chebotarev classes in the
actual finite representation of the point
\(3\), with \(v_5(q_r+1)=r\).
Every finite-coefficient central obstruction
vanishes, but the integral section cannot lift.
The proof establishes global nonvanishing by
continuity, not by assigning an unverified
family of local representatives to the global
cohomology class. This proves Lead 13.
\end{statement}

\begin{statement}{D.4}{Proved}{The unrestricted class-four and class-five images}
D.31 fixes the actual arithmetic Frobenius action
through class five, including its normalized Soul\'e
coefficient. D.32 evaluates the three class-four
and six class-five global residue families in
integral Hall coordinates at every prime other
than five. All nine polynomials are expanded in
the appendix to D.32. A class lifts exactly when
each residue is divisible by \(q^{n-1}-1\) and
the divided family tends to zero. Nonempty fibres
are \((\mathbb Z/5)^3\) at class four and
\(\mathbb Z_5^6\) at class five. The three
order-five parameters are retained explicitly,
and the six R17 congruences at each of two and
three are recovered by an independent integral
Magnus calculation. Other primes are governed by
the complete formulas, with their full five-adic
precision and convergence condition.
\end{statement}

\begin{statement}{D.5}{Open}{The full metabelian set and the higher-derived Fermat comparison}
D.39--D.41 give the actual all-weight arithmetic
action and the complete residue-and-branch
description of the unrestricted pro-five
metabelian section set. D.35 gives all three
cusp packet maps and proves torsion-freeness
of every nonempty Alexander fibre.
D.36 and D.44 compute the unrestricted global
Alexander cohomology above every abelian pair
unramified outside five, retaining the nonzero
derived-limit term. D.42 and D.43 give the
actual local cusp constraint and ramified
local cohomology. The corresponding global
first cohomology for every abelian twist
is proved compact and torsion-free in D.47,
hence a product of copies of \(\mathbb Z_5\);
D.60 gives an infinitely ramified cusp
tangent for which that rank is zero and
the metabelian fibre is a singleton. D.61
strengthens this to rank-zero modules dense
in every five-adic neighborhood of every
abelian pair, even with its second coordinate
held fixed. D.66 strengthens this to a comeagre
zero-cohomology locus, globally and in every
fixed-second-coordinate parameter space. In
each completed cusp packet a comeagre set
of parameters has a singleton metabelian
fibre, with existence supplied by the full
cusp section. D.67 gives a comeagre subset
of these parameters for which the full free
pro-five fibre itself is a singleton, and hence
so is its set of compatible nilpotent section
towers. D.68 proves compactness of the exact
fixed-abelian fibres at every nilpotent stage,
for both the full and metabelian quotients.
Thus individual liftability at every higher
stage does imply a compatible actual limit
in these fibres. Each noncuspidal finite-stage
class above a generic cusp parameter fails
at a finite stage, uniformly outside any
fixed cusp neighborhood; that stage and its
first obstruction are not yet evaluated.
D.69 extends the generic singleton cusp-fibre
assertion to every coefficient prime, including
the sign component at two, and then to the
actual full profinite section fibre by D.24.
Its parameter space is the product of the
individual completed Kummer spaces; it makes
no classification of arbitrary abelian pairs.
D.70 proves the corresponding generic full pro-prime
and full profinite cusp-fibre assertion for every
affine number-field curve with at least three geometric
cusps and one chosen rational cusp. It applies to
the actual Fermat curves at every coefficient prime,
including primes not dividing the covering exponent;
it does not identify their metabelian fibres.
D.71 treats the requested once-punctured elliptic
curve itself: its full cusp packet is already
injected at class two, and a comeagre set of
these class-two classes has a singleton full
pro-prime or full profinite fibre. Its abelian
fibre contains the entire cusp packet; no
singleton abelian assertion is made there.
D.72 proves compactness of exact fixed-class-two
fibres on every affine number-field curve with
a chosen actual section. Individual lifts at
every nilpotent stage can therefore be made
compatible. For D.70 and D.71 generic cusp
parameters, a noncusp finite-stage class fails
at some finite nilpotent stage, uniformly
outside a fixed cusp neighborhood. Neither
the first stage nor its obstruction is evaluated.
D.73 extends generic class-two cusp rigidity
to every affine hyperbolic number-field curve
with one chosen rational cusp, including one
or two geometric punctures. It separates the
two-cusp inertia groups by their degree-two
surface term; it makes no singleton abelian
claim in those cases.
Global pair existence itself lies in a closed
nowhere-dense local slope filter; no generic
existence is asserted.
D.74 computes the first cohomology for every finitely
ramified pair: it is a countably infinite compact product
of copies of \(\mathbb Z_5\). The proof constructs
actual integral cocycles by Ore clearing in the arithmetic
operator ring; it does not infer an integral limit from
growing rational ranks. Every finitely ramified cusp fibre
therefore contains uncountably many noncuspidal, non-point
metabelian classes. Arbitrary infinitely ramified ranks
remain undetermined. D.48 isolates the
unrestricted second-cohomology residue
term and proves a nonzero Jannsen
kernel for every finitely ramified
twist, while retaining the unevaluated
compact base cohomology. D.62 proves that
for every twist, including infinite ramification,
the unrestricted Jannsen kernel is nonzero,
is exactly the closure of zero in the uniform
cochain quotient topology, and is invisible
at every local place. D.63 computes the
separated quotient up to a precisely identified
compact global kernel: it is an open extension
of an explicit convergent family of actual
local operator cokernels. This separated
quotient is never locally compact. D.64 now
evaluates the family supported at five and
one prime of order twenty modulo twenty-five:
its first-cohomology rank is countably infinite,
its compact base second cohomology is the
actual local group at that prime, its separated
global compact kernel vanishes, and its
metabelian fibre exists exactly in the three
local cusp directions. D.65 computes the abstract
residue part of every finitely ramified Jannsen
kernel as a countable product of \(\mathbb Z_5\),
and splits off its retained compact term algebraically.
D.75 proves that the separated compact term is sometimes
large: a cusp parameter gains a split free factor of rank
its extra local divisibility at five, and an explicit
two-prime logarithm cancellation gives a split countable
product of \(\mathbb Z_5\). The complementary compact
factor remains undetermined.
D.76 proves that the contributions of all three peripheral
axes add for every abelian pair. D.77 computes the entire
local-at-five group and global invariant dual, including their
actual map: their kernel is exactly this three-axis factor.
The separated compact kernel is therefore continuously
\(\mathcal S_s\times Z(d(s))\), with only the inverse limit
\(\mathcal S_s\) of finite global localization kernels
retained in this compact calculation.
D.65 gives the whole Jannsen kernel for D.44 and D.64,
with the intrinsic uniform five-adic topology
explicitly separated from the nonseparated
cohomology topology. The exact rank for every
twist, the remaining compact terms, and the
abstract kernel type for general infinite
ramification remain uncomputed.

D.12 proves the existence of nonpoint,
noncuspidal pro-five metabelian classes,
with actual compatible lifts of its constructed
classes to the full free pro-five group.
D.20 independently proves the full profinite
metabelian Fermat assertions LM6 E.33--E.34,
and D.21 gives the corresponding full
metabelian tripod families. D.52 proves the
nilpotent part of the hunch for both the
three-power and five-power Fermat towers,
including compatible actual sections of
the full pronilpotent quotient. D.54 proves
the local prosolvable section theorem, and
D.55 proves the Legendre-relative prosolvable
hunch at every five-power level, as well as
cuspidality for the plain prosolvable
three-power tower. Thus the constructed
noncuspidal families cannot have compatible
lifts to those inverse limits. D.57 strengthens
this for both towers: every nonzero D.20
parameter is already locally nongeometric
at the tower prime, and fails to lift at
some finite derived length. A uniform
finite length exists away from every
fixed neighborhood of the zero parameter.
D.59 retains this obstruction in every
two-prime quotient \(\{\ell,q\}\),
\(q\ne\ell\), including \(q=2\).
It proves the needed local Fermat point
separation by the single \(\ell\)-adic
proper Kummer map using explicit cusp disks.
The least obstructed length and its
obstruction class remain undetermined;
so does the hunch about all section
classes at a fixed derived length. These precise
global-cohomology and higher-lifting
obligations keep the whole of Task D.5 open.
D.78 computes the remaining compact kernel for the actual
point pair \(\kappa(2),\kappa(3)\), realized by \(t=-2\):
it is exactly \(\mathbb Z/5\). D.79 classifies the whole
compact kernel for every invertible matrix recombination of
these two Kummer classes, including all finite and exact
local boundary depths. D.80 extends this classification to
arbitrarily infinitely ramified pairs that are independent
modulo five, have zero five-valuations, and have a primitive
local unit-logarithm vector. Its explicit ordinary rational
points realize \(\mathbb Z_5^d\) for every positive finite
rank \(d\). D.81 removes the local-primitivity restriction
for all globally independent pairs with zero five-valuations.
The entire compact kernel is an explicitly determined product
of a possible order-five term and the three local-depth free
factors. Its finite-support examples include
\(\mathbb Z/5\times\mathbb Z_5^n\) for every \(n\ge3\).
D.82 allows arbitrary five-valuations whenever two independent
residual inertia directions occur away from five. It computes
the full compact kernel from the span of the local valuation
and unit vectors; rank two gives zero, realized by the actual
point \(t=10\). D.83 proves that every finite-coefficient
localization away from five is an isomorphism when the closed
away-five valuation lattice is generated by its local unit
vector. This includes pure five-valuations, infinite ramification,
and the actual point \(t=5\), all with zero separated compact
kernel. D.84 and D.85 determine further dependent-away and
globally independent residual families with their actual maps.
D.87 identifies the compact finite-localization term for every
completed pair with the kernel on its actual compact operator
images, including one common finite coefficient bound.
D.88--D.91 evaluate this term and the entire separated compact
kernel for every parameter in all three cusp packets, including
zero parameters, pure five-valuations, nonprimitive ramification,
finite local indices, and the infinite endpoints. The smaller
cyclic elementary divisors and the natural free boundary quotient
are retained; the product decompositions are noncanonical.
D.92 determines the exact abelian image on every global
valuation line: a nonempty fibre must be peripheral or
supported only at five, and both possibilities occur.
The only ordinary rational points in this locus are
\(-1,2,1/2\), at the level of their abelian images.
D.91 and D.83 therefore compute the compact kernel of every
nonempty fibre on this entire rank-at-most-one locus.
D.93 computes the complete finite compact kernel on the actual
ordinary-point family \(t=2^{5^r}\), \(r\ge1\), with cyclic
exponents \(1,\ldots,r-1,r+1\) and unbounded order. Its
all-depth Smith factors are proved by an integral norm extension.
Arbitrary pairs outside the evaluated families, general infinitely
ramified first-cohomology ranks, and the stated higher-derived
obligations remain Open.
\end{statement}

\begin{statement}{D.6}{Open}{The induced Fermat sections and the remaining full higher stages}
D.45 determines every induced metabelian
Fermat stage and its limit by explicit
arithmetic residues, with unrestricted
ramification, all transition maps, and
the complete all-level intersection.
The full nonmetabelian ambient stages are
evaluated through class five by D.30 and
D.32. D.29 gives their exact general
central criterion, but the all-weight
nonmetabelian arithmetic action required
to evaluate it is not computed here.
LM8 F.22--F.23 treats both versions
explicitly; the broader full-stage
reading of Task D.6 therefore remains open.

The suggested natural equality with
the covering's own metabelian section
set is refuted at every five-power
level by D.37 and D.38.
D.46 strengthens this comparison:
it is surjective on classes unramified
outside five, and every such fibre
contains an actual integral family
of rank \((5^{2r}-1)^2/4\).
Thus its failure of injectivity is
arithmetically substantial, while
the remaining open assertion is
the specified full higher-stage
determination.
\end{statement}

\begin{statement}{D.7}{Proved}{Cyclotomic section sets and exact descent}
D.15 proves the cohomological-dimension assertion
of Lead 16, every central lifting surjectivity
over the cyclotomic field, and the precise
non-Hausdorff coinvariant interpretation of
the derived-limit term. D.16 determines every cyclotomic section set by
convergent free-basis parameters. D.34 gives the
actual affine action of \(\Gamma\) on every
central lifting torsor, its exact descent
equation, and the finite-orbit residue comparison.
It derives the two, three, and six families of
D.30 and D.32 again from this descent equation,
with the same unclosed image and the same
fibres, including the three order-five
parameters arising from \(\Gamma\) itself.
\end{statement}

\begin{statement}{D.9}{Open}{The surviving mixed-prime classes and the relative section conjecture}
D.17 excludes all but finitely many classes in
each infinite fibre \(H^1(\mathbb Q,E_c[5])\)
from even the compatible Tate-module lifting
problem. Which of the remaining classes meet
the local conditions at three, five and the
cusps and lift through every relative pro-five
nonabelian stage has not been determined.
The relative pro-five equality with the Kummer
images of rational points and cusps remains open.
For the full geometric group, D.58 proves that
its first-five fibre is exactly the rational-point
Kummer image whenever the cubic twist has
\(T_5\operatorname{Sha}=0\), in particular in analytic rank
zero or one (respectively one or five classes).
For a possible nonzero \(T_5\operatorname{Sha}\), the
non-Mordell--Weil part of the finite inverse-Selmer
image is not yet decided. D.53 also settles the full-group
question above the trivial cubic class:
exactly its zero \(E_1[5]\)-cohomology class
extends, and every full section there is
cuspidal. In particular
D.8's finite semidirect-product splitting is
not a solution to either lifting problem.
\end{statement}

\begin{statement}{D.10}{Proved}{Primewise geometric coherence and rigidity}
Both M0 and M1 of Lead 19 hold for every
hyperbolic curve over a number field.
M1 is D.24. M0 is D.51: even if the
ordinary point or cusp is initially
allowed to depend on the prime, all
primewise geometric images of one full
section have one global geometric origin.
In particular this settles the specified
punctured CM curve and the unrestricted
tripod case.

The proof of the remaining mixed cases
uses the new finite-field cusp/ordinary
separation in D.50, followed by D.49's
uniform integrality and Siegel--Mahler
argument and D.25's finite-support theorem.
It does not infer agreement on the
arithmetic subgroup of Part A merely
from agreement in every pro-prime quotient.
Equivalently, after a geometric reference
section is fixed, the fibre over its image
modulo the common pro-prime kernel \(N\)
is a singleton: the corresponding orbit
set of nonabelian \(H^1(G_k,N)\) under
the fixed elements of \(\overline\pi/N\)
has one element, by D.24.
\end{statement}

\begin{statement}{D.16}{Proved}{All section classes over the cyclotomic field}
Let \(\Delta_5=G_{\mathbb Q(\mu_{5^\infty})}\), and let
\(P\) be its maximal pro-five quotient. It is a free
pro-five group with a countable basis
\(g_1,g_2,\ldots\) converging to one. The arithmetic
action on \(T_n\) factors through \(P\); denote it
by \(\rho_n\). For every \(n\ge1\), there is a
complete parametrization
\[
 \operatorname{Sec}_{\mathbb Q(\mu_{5^\infty})}(T_n)
 =\{(u_i)_{i\ge1}\in\prod_iT_n:u_i\to1\}/\sim,
 \qquad
 (u_i)\sim\bigl(a^{-1}u_i\rho_n(g_i)(a)\bigr)_i,
 \quad a\in T_n.                                      \tag{D.16.1}
\]
Every such class lifts compatibly through every
higher class and to \(F^{(5)}\). Formula (D.16.1)
is noncanonical through its choice of free basis,
but imposes no equations on the parameters other
than convergence and the displayed conjugacy
relation.
\end{statement}
\begin{proof}
D.15 gives \(H^2(\Delta_5,\mathbb F_5)=0\).
Write \(N\) for the kernel of \(\Delta_5\to P\).
There is no nonzero invariant homomorphism
\(N\to\mathbb F_5\): its kernel would be normal
in \(\Delta_5\), and would give a pro-five
quotient of \(\Delta_5\) in which \(N\) has a
nontrivial image, contradicting maximality of
\(P\). The Hochschild--Serre five-term sequence
therefore makes inflation
\(H^2(P,\mathbb F_5)\to H^2(\Delta_5,\mathbb F_5)\)
injective. Hence \(H^2(P,\mathbb F_5)=0\).

For clarity, the free pro-five criterion here can
be proved by its relation argument. Choose a
minimal convergent generating family of the
countably based pro-five group \(P\), and map
the free pro-five group \(F_0\) on that family
onto \(P\). If \(R\) is its kernel, the map on
\(H^1(-,\mathbb F_5)\) is an isomorphism by
minimality. Since \(H^2(F_0,\mathbb F_5)=0\),
the five-term sequence gives
\(H^1(R,\mathbb F_5)^P=0\). A nontrivial \(R\)
has nontrivial image \(\overline R\) in a finite
five-group quotient \(Q\) of \(F_0\). The
nonzero module \(\overline R/\Phi(\overline R)\)
has nonzero \(Q\)-coinvariants, because the
augmentation ideal of a finite five-group over
\(\mathbb F_5\) is nilpotent. A linear functional
on those coinvariants gives a nonzero
\(F_0\)-invariant homomorphism \(R\to\mathbb F_5\),
a contradiction. Thus \(R=1\).

The basis is countable because the field
\(\mathbb Q(\mu_{5^\infty})\) and its algebraic
closure are countable; its finite extensions
and hence its finite Galois quotients form a
countable family. The rank is infinite:
Kummer theory gives
\(H^1(\Delta_5,\mathbb F_5)
 =\mathbb Q(\mu_{5^\infty})^\times/
   \mathbb Q(\mu_{5^\infty})^{\times5}\).
The rational primes other than five give
infinitely many independent classes. At any
one of them the cyclotomic tower is unramified,
so its normalized valuation takes integer
values and detects its exponent modulo five.
This also shows directly that the discrete
\(H^1\) has countably infinite dimension.
A basis compatible with a countable
neighborhood filtration of the Frattini
quotient has convergent lifts, as used above.

The action of \(\Delta_5\) on the Frattini
quotient of \(F^{(5)}\) is trivial, since its
cyclotomic character is one. The kernel of
\(\operatorname{Aut}(F^{(5)})\to\operatorname{GL}_2(\mathbb F_5)\)
is pro-five: on successive finite Frattini
quotients its congruence kernels are subgroups
of additive homomorphism groups over
\(\mathbb F_5\). Hence the arithmetic action
factors through \(P\). A crossed homomorphism
\(u:\Delta_5\to T_n\), together with this
action, defines a homomorphism into a pro-five
semidirect product, and therefore also factors
through \(P\).

By the universal property of a free pro-five
group on a convergent basis, the assignments
\(g_i\mapsto(u_i,\rho_n(g_i))\) extend to a
continuous homomorphism precisely when
\(u_i\to1\). Its second projection is the
prescribed action, so its first component is
exactly a continuous crossed homomorphism.
Conjugating by \(a\in T_n\) gives exactly
the equivalence relation in (D.16.1).
Finally an open continuous surjection of
compact groups allows a null sequence to be
lifted to a null sequence: lift its sufficiently
late terms inside successively smaller open
neighborhoods of one. Apply this to
\(T_{n+1}\to T_n\), successively and compatibly.
The resulting inverse-limit cocycle has values
in \(\varprojlim_nT_n=F^{(5)}\), as proved in
D.12. This proves all assertions.
\end{proof}

\begin{statement}{D.19}{Proved}{The geometric candidates inside each mixed-prime fibre}
For \(E_c:y^3=cx(1-x)\), with origin the point over
zero, the three cusp points are rational
three-torsion points. Their images in
\(H^1(\mathbb Q,E_c[5])\) are all zero.
The union of the images of rational noncuspidal
points and cusp sections is exactly the Kummer
subspace
\[
 E_c(\mathbb Q)/5E_c(\mathbb Q)
          \hookrightarrow H^1(\mathbb Q,E_c[5]).       \tag{D.19.1}
\]
It has \(\mathbb F_5\)-dimension
\(\operatorname{rank}E_c(\mathbb Q)\); in particular
there is no rational five-torsion correction.
Thus the proposed equality in Task D.9 is an
explicit equality between the as yet unknown
set of liftable classes in the finite filter
D.17 and the finite subspace (D.19.1).
\end{statement}
\begin{proof}
The order-three deck automorphism fixes all three
cusps and the chosen origin. It is multiplication
by a primitive cube root of unity on the elliptic
curve. Its fixed kernel has order three, since
\(N_{\mathbb Q(\zeta_3)/\mathbb Q}(1-\zeta_3)=3\).
It consists exactly of the three displayed
rational points, so their Kummer classes modulo
five vanish. A cuspidal section maps under
compactification to the section of its cusp
point; its tangent parameter lies in cusp
inertia, which compactification kills. This
proves the assertion for every tangent vector
and every section in the cusp packet.

The Kummer sequence identifies the point image
with \(E_c(\mathbb Q)/5E_c(\mathbb Q)\).
Removing the three cusp points removes only
representatives of its zero element, which is
restored by the cusp sections. This proves
(D.19.1). To exclude rational five-torsion,
work over \(k_3=\mathbb Q(\zeta_3)\), where all CM
endomorphisms are defined. Because five is
inert in \(k_3\), \(E_c[5]\) is a one-dimensional
vector space over \(\mathbb F_{25}\), and
\(G_{k_3}\) acts by \(\mathbb F_{25}\)-linear
scalars. A nonzero invariant vector would
force all of \(E_c[5]\) to be invariant.
The Weil pairing would then put \(\mu_5\)
inside \(k_3\), impossible since its degree
over \(\mathbb Q\) is four whereas
\([k_3:\mathbb Q]=2\). Hence
\(E_c[5](\mathbb Q)=0\). Mordell--Weil and
the injectivity of the Kummer map now give
the stated dimension.
\end{proof}

\begin{statement}{D.22}{Refuted}{Ramification alone does not make the Tate parameter divisible}
The divisibility assertion in Steklov's Proposition 4.6,
and its use to split the residual representation in
Proposition 5.10, are false as written. The same
divisibility sentence in LM6 E.8 is false.
For the Fermat cusp over zero, with local coordinate
\(t=X/Z\) and unit tangent, the induced Legendre
Tate parameter has Kummer class
\[
       q(s_O)=\kappa(256^{-1})
          \in H^1(\mathbb Q,\widehat{\mathbb Z}(1)).  \tag{D.22.1}
\]
Its valuation at two is \(-8\), so it is not divisible
by any odd exponent \(n\ge3\), although the ramification
index of the map to the \(j\)-line is \(2n\).
The corrected formal rule is an affine coset:
if \(q=a t^e h(t)\), \(h(0)=1\), then
\[
          q(s)=\kappa(a)+e\,\kappa_t(s).              \tag{D.22.2}
\]
The omitted constant does not invalidate the local
conductor and finite-flat conditions needed for the
Fermat Selmer theorem; those are proved in D.23.
\end{statement}
\begin{proof}
The primary text \cite[Propositions 4.6 and 5.10]{DE_Steklov2026}
replaces a ramified formal extension by
\(k((j^{-1/e}))/k((j^{-1}))\). This replacement requires
the leading coefficient to be an \(e\)-th power in
\(k\), which need not hold. In characteristic zero,
every element of \(1+t k[[t]]\) has a unique \(m\)-th
root there for each positive integer \(m\), by the
binomial series. Such principal units have zero
Kummer class at every finite coefficient level.
Pulling back \([q]=[a]+e[t]\) along a section proves
(D.22.2), including the inverse limit.

For the Legendre family,
\[
 j(\lambda)=256\frac{(\lambda^2-\lambda+1)^3}
                         {\lambda^2(1-\lambda)^2},
 \qquad q=j^{-1}+O(j^{-2}).
\]
At the Fermat cusp, \(\lambda=-t^n\), so
\[
       q=256^{-1}t^{2n}(1+O(t^n)).                   \tag{D.22.3}
\]
The unit tangent has \(\kappa_t(s_O)=0\). Formula
(D.22.1) follows. Applying the valuation map
\(H^1(\mathbb Q,\widehat{\mathbb Z}(1))
\to\widehat{\mathbb Z}\) at two proves the failure
of divisibility. The same counterexample is already
visible over \(\mathbb Q_2\), and modulo \(n\).
The formal leading coefficient is checked exactly
at all three cusps in \texttt{de\_fermat\_local.py}.
\end{proof}

\begin{statement}{D.23}{Proved}{The full Fermat section theorem with the cusp constant repaired}
Every Selmer section of the affine Fermat curve
\(U_n/\mathbb Q\), for odd \(n\ge3\), is cuspidal.
D.56 with the class of all finite groups,
together with the real section theorem, makes
every full section Selmer and hence cuspidal. The complete full section set is the
disjoint union of the packets at the three rational
cusps; each is a torsor under
\(H^1(\mathbb Q,\widehat{\mathbb Z}(1))\).
This repairs the proof of LM6 E.12 and proves
the full-group side of the comparison with D.20.
The full local point-or-cusp theorem used here,
including its affine form, is proved in D.56.
\end{statement}
\begin{proof}
We use ``Selmer'' to mean point or cusp induced at
every completion, with the usual real component
interpretation. First record the finite-support
argument, including its compatibility requirement.
If a Selmer section of an affine curve has, at
all finite places, underlying compactification
points in a finite reduced subscheme \(Z\), take
any characteristic geometric section neighborhood
with compactification of genus at least one.
The lifted section is still Selmer: a point
section fixes an element of the finite etale
fibre, and a cusp section has a rational cusp
above it because its intersection with the
cover group still surjects onto the local
Galois group. The compactification points
survive every finite torsor: evaluate the torsor
on the global section and restrict that one
cohomology class at every place. Theorem 8.2
of \cite{DE_Stoll2007}, applied to the finite
inverse image of \(Z\), supplies a rational
point in that inverse image. Its hypothesis
uses a set of finite places of density one;
no assertion about a scheme-theoretic fibre
of a real component is needed, as clarified
in \cite[\S3]{DE_StollErratum}.

These finite nonempty sets over section
neighborhoods have the finite intersection
property: a common refinement maps into
every prescribed finite collection of levels.
Their inverse limit supplies a horizontal
point fixed by the section in the full cover
tower. Its stabilizer is an ordinary or cusp
decomposition group. Consequently the section
is a rational-point or cusp section. Characteristic
neighborhoods of the required genus are cofinal:
make the orders of all peripheral inertia
images large in a finite quotient and use
Riemann--Hurwitz, then take a characteristic
core. It follows in particular that a nonconstant
map from a compactification to a proper curve
\(W\) with \(W(\mathbb A_\mathbb Q)_\bullet^{\rm f-cov}
=W(\mathbb Q)\) forces every Selmer section
upstairs to be point or cusp induced. The
finite fibre above the resulting rational
point of \(W\) is the required \(Z\).

Now let \(\ell\ge5\) be prime and let \(s\) be
a Selmer section of \(U_\ell\). Compose with the
Legendre family's residual representation to
obtain
\(\bar\rho_s:G_{\mathbb Q}\to\operatorname{GL}_2(\mathbb F_\ell)\),
of determinant \(\bar\chi_\ell\). We check every
finite local condition used in the modularity
argument. At a local ordinary point, write
\(a^\ell+b^\ell+1=0\) and \(\lambda=-a^\ell\).
If \(v_p(a)\ne0\), the exact \(j\)-formula gives
\[
             v_p(j)=8v_p(2)-2\ell|v_p(a)|.          \tag{D.23.1}
\]
If \(v_p(a)=0\) and \(\lambda\equiv1\pmod p\),
use \(1-\lambda=-b^\ell\) instead. At two,
one of these alternatives always occurs and
\(8-2\ell<0\). At an odd prime the remaining
alternative is a Legendre integral equation
with \(\lambda\) and \(1-\lambda\) both units,
hence good reduction. Thus, after a local
quadratic twist, the curve has good reduction
or is a Tate curve, and at odd primes in the
Tate case \(\ell\mid v_p(q)=-v_p(j)\).
These calculations also independently verify
the local point input of
\cite[Lemma 5.12 and Proposition 5.13]{DE_Steklov2026}.

At a local cusp, choose the appropriate one
of \(X/Z,Y/Z,Z/X\) as parameter. In each case
\(j^{-1}\), and hence \(q\), has leading term
\(256^{-1}t^{2\ell}\). D.22 therefore gives,
for every section in that cusp packet,
\[
       [q(s)]_\ell=[256^{-1}]_\ell
                      \in H^1(\mathbb Q_p,\mu_\ell).
                                                               \tag{D.23.2}
\]
The residual representation is a quadratic
twist of the Tate extension represented by
this class. It need not split.

Here are the exact consequences, both for
ordinary points and for cusps. At odd
\(p\ne\ell\), a class with valuation divisible
by \(\ell\) can be represented by a unit.
Its Kummer extension is unramified: raising
to the \(\ell\)-th power is an automorphism
of the principal-unit group, and a root of
the residue-field unit is obtained in an
unramified extension. The cyclotomic character
is also unramified there. At \(p=\ell\), the
same valuation condition makes the Tate
class a unit class. Its extension of
\(\mathbb Z/\ell\) by \(\mu_\ell\) is finite
flat over \(\mathbb Z_\ell\): for a unit \(u\),
the components \(z^\ell=u^a\), \(0\le a<\ell\),
have multiplication
\[
 (a,z)(b,w)=\left(a+b-\ell\left\lfloor\frac{a+b}{\ell}\right\rfloor,
    zw/u^{\lfloor(a+b)/\ell\rfloor}\right).
\]
These finite free component algebras, with
this associative multiplication, give a
finite flat group scheme extending that generic fibre.
The good-reduction case is finite flat by
the finite flat group scheme \(E[\ell]\).
Thus the residual representation has Serre
weight two by \cite[\S2.8, Proposition 4]{DE_Serre1987},
since its determinant on inertia is cyclotomic.
At two, after the quadratic
twist, the inertia action is unipotent with
at most a one-dimensional nontrivial quotient.
Wild inertia is pro-two and cannot map
nontrivially to its additive group of order
\(\ell\). Hence the action is tame and the
Artin conductor exponent is at most one.
In (D.23.2), the constant \(256^{-1}\) is a
unit at every odd prime, so every stated
condition holds despite the error in D.22.

The finitely many local quadratic characters
on inertia patch to one global quadratic
character. Indeed \(G_{\mathbb Q}^{\rm ab}
\simeq\prod_p\mathbb Z_p^\times\); put the
prescribed character on each relevant factor
and the trivial character on the others.
After this twist the representation is
unramified outside \(\{2,\ell\}\), has
conductor dividing two away from \(\ell\),
weight two at \(\ell\), and determinant
\(\bar\chi_\ell\). If it were irreducible,
it would be absolutely irreducible: any
invariant line over \(\overline{\mathbb F}_\ell\)
is an eigenspace of complex conjugation,
whose distinct eigenvalues are \(1,-1\),
so the line is already defined over
\(\mathbb F_\ell\). The precise Serre
modularity theorem
\cite[Theorem 9.1]{DE_KhareWintenberger2009},
with its hypothesis (H) supplied by
\cite[Theorem 0.1 and Corollary 0.2]{DE_Kisin2009},
would give a weight-two cusp form of level
one or two, and trivial nebentypus. Both
spaces vanish since \(X_0(1)\) and
\(X_0(2)\) have genus zero. This proves
reducibility.

There is therefore a section lift to the
pullback over \(U_\ell\) of the finite etale
cover \(Z_\ell\) of the tripod parametrizing
a cyclic subgroup of order \(\ell\) of the
Legendre curve. Work with the connected
component of \(U_\ell\times_{\mathbb P^1-\{0,1,\infty\}}Z_\ell\)
supporting this lift. It has constant
field \(\mathbb Q\), since its arithmetic
image is all of \(G_{\mathbb Q}\). This
cover has a nonconstant map to \(Y_0(4\ell)\),
as follows. Let \(P=(0,0)\), \(Q=(1,0)\),
\(\phi:E\to E'=E/\langle P\rangle\), and
\(\phi':E'\to E\) be its dual. The subgroup
\((\phi')^{-1}(\langle Q\rangle)\) is cyclic
of order four: a lift \(R\) of \(Q\) satisfies
\(2R=\phi(Q)\ne0\). Its sum with \(\phi(C)\),
for the tautological order-\(\ell\) subgroup
\(C\), is cyclic of order \(4\ell\).
The resulting modular map is nonconstant,
since a fixed elliptic curve has only
finitely many degree-two isogenous curves
whereas \(j(E_\lambda)\) is nonconstant.
It extends to a finite morphism of smooth
compactifications. The index, cusp count,
and elliptic-point counts for \(\Gamma_0(4\ell)\)
are respectively \(6(\ell+1),6,0,0\); hence
\[
                 g(X_0(4\ell))=(\ell-3)/2>0.         \tag{D.23.3}
\]
Stoll's Corollary 8.8 \cite{DE_Stoll2007}
applies to every positive-genus \(X_0(N)/\mathbb Q\);
its conclusion is equality with the rational
points, not merely their adelic closure. It
and the finite-support argument show that
the lifted section, and then the original
section, is a point or cusp section.

To exclude the ordinary point case, retain
the full rational two-torsion in the
Legendre curve. For a reducible representation,
the same local quadratic twists make both
diagonal characters unramified away from
\(\ell\). At \(\ell\), the good supersingular
case is impossible because its inertial
representation is irreducible; in the
ordinary and Tate cases the diagonal
characters on inertia are \(1,\bar\chi_\ell\).
One global diagonal character is thus
unramified at all finite primes, hence
trivial by Minkowski's theorem. The two
global characters are \(1\) and
\(\bar\chi_\ell\). A rational elliptic
curve after that quadratic twist therefore
has a subgroup \(\mathbb Z/\ell\) or
\(\mu_\ell\). In the second case quotient
by \(\mu_\ell\); the dual isogeny has
constant kernel \(\mathbb Z/\ell\).
Quadratic twisting and an odd isogeny both
preserve full rational two-torsion. A
rational elliptic curve would consequently
have torsion containing
\(\mathbb Z/2\oplus\mathbb Z/(2\ell)\),
contrary to Mazur's torsion theorem
\cite[Theorem (8)]{DE_Mazur1977} for \(\ell\ge5\).
This proves the prime-exponent assertion,
including five. It also avoids the incorrect
claim that \(Y_1(7)(\mathbb Q)\) is empty.

For exponent three, put
\[
 u=-\frac{3Z}{X+Y},\qquad
 v=\frac{9(X-Y)/(X+Y)-1}{2}.
\]
Substitution gives the smooth cubic
\(v^2+v=u^3-7\). Under
\((x,y)=(4u,8v+4)\) it is
\(y^2=x^3-432=X_0(27)\), with this modular
identification recorded in
\cite[\S10.1]{DE_KrizLi2019}.
Stoll's Corollary 8.8 gives the finite-descent
equality and implies that its rational
group is finite. For clarity, this last
finiteness also follows directly from that
equality: the finite-descent set is compact,
so its projection to \(E(\mathbb Q_p)\) is
a compact countable subgroup; an infinite
compact metrizable group is uncountable.
Good reduction counts at five and seven
are six and nine. Prime-to-residue-characteristic
torsion injects under reduction, so the
rational torsion has order dividing three.
The origin and \((u,v)=(3,-5),(3,4)\) give
all three points, and correspond to the
three rational cusps. The finite-support
argument again makes every Selmer section
cuspidal. The identities, point counts,
and modular genera are independently checked
in \texttt{de\_fermat\_local.py}.

For odd \(n\), choose a prime \(\ell\mid n\).
The coordinate power map \(U_n\to U_\ell\)
is finite etale and preserves local point
and cusp sections. Its section image is
cuspidal by the prime cases. Intersecting
with that cusp decomposition group gives
a rational cusp upstairs and a cusp
section there. Finally D.56, with the
formation of all finite groups, makes every
full affine section Selmer at every finite place. The real
section theorem supplies the real component.
The three rational cusps have tangents over
\(\mathbb Q\); the cusp decomposition sequence
and \cite[Proposition 8 and Theorem 14]{fghStix}
give the asserted disjoint torsors. The
injectivity hypothesis in that theorem holds
over every number field by Mordell--Weil,
as explained in its Remark 15 and Appendix B. No
global density step of the supplied local
preprint is used.
\end{proof}

\begin{statement}{D.24}{Proved}{Primewise agreement with one geometric section implies full agreement}
Let \(U/k\) be any smooth geometrically connected
hyperbolic curve over a number field, let
\(\pi=\pi_1(U_{\bar k})\), and put
\[
 N=\ker\left(\pi\longrightarrow\prod_\ell\pi^{(\ell)}\right).
\]
Let \(s\) be a full section and let \(s_x\) be
the section of a rational point or any fixed
section in the packet of a rational cusp.
If \(s\) and \(s_x\) have the same class in
the quotient by \(N\), then they have the
same full section class. The same conclusion
holds if their classes are specified to agree
separately in every geometric pro-\(\ell\)
quotient. This proves M1 of Lead 19 for all
three requested classes of curves. D.56, with the class of all finite groups,
provides the full local point-or-cusp theorem
for the curve and every full section
neighborhood. No global or pro-\(p\)
section conjecture is assumed.
\end{statement}
\begin{proof}
Write \(X\) for the smooth proper
compactification and \(D=X-U\). We give the
local separation and the final cusp-parameter
argument explicitly, since abelianization
alone does not distinguish all cusp packets.

First, \(\pi\to\prod_\ell\pi^{(\ell)}\) is
surjective. Its image surjects onto each
finite product of finite prime-power
quotients: a proper subdirect product would,
by Goursat's lemma inductively, give a common
nontrivial quotient of groups of coprime
orders. Compactness gives surjectivity onto
the full product. Consequently individual
primewise conjugators lift to one element
of \(\pi\). We may therefore choose
representatives whose images agree as
sections in the quotient by \(N\).

For a finite extension \(F/\mathbb Q_p\) and
an abelian variety \(A/F\), the kernels of
the Kummer maps are
\[
 \ker\bigl(A(F)\to H^1(F,T_pA)\bigr)
       =A(F)_{\rm tors,p'},\qquad
 \ker\bigl(A(F)\to H^1(F,TA)\bigr)=0.             \tag{D.24.1}
\]
Indeed the finite-level Kummer maps
\(A(F)/mA(F)\to H^1(F,A[m])\) are injective.
The finite groups \(H^0(F,A[m])\)
are Mittag--Leffler along a divisibility-cofinal
sequence, so the continuous \(H^1\) Kummer map
agrees with the finite-level inverse limit.
The compact group \(A(F)\) has an open
subgroup isomorphic to \(\mathbb Z_p^{[F:\mathbb Q_p]\dim A}\),
by the formal logarithm. Its intersection
\(\bigcap_r p^rA(F)\) is exactly its finite
prime-to-\(p\) torsion subgroup, and
\(\bigcap_m mA(F)=0\). This also follows
from its decomposition into a
\(\mathbb Z_p\)-lattice and finite primary
factors. The local multiplicative Kummer map
\[
       F^\times\longrightarrow
       H^1(F,\widehat{\mathbb Z}(1))              \tag{D.24.2}
\]
is injective as well: use the valuation
factor \(\mathbb Z\), the finite group of
roots of unity, and the open unit subgroup
with logarithm lattice \(\mathbb Z_p^{[F:\mathbb Q_p]}\).
No injectivity for a single \(T_\ell\) at
\(\ell\ne p\) is asserted.

By A.17, the restriction of \(s\) at every
finite place is a point or cusp section.
Let \(y_v\in X(k_v)\) be its underlying
compactification point. If \(g(X)\ge1\),
the Abel--Jacobi map based at the rational
point \(x\in X(k)\) induces, in every
prime component, the Kummer class of
\([y_v-x]\in J(k_v)\). This holds for a
cusp section too: its inertia is killed
by \(\pi_1(U)\to\pi_1(X)\). Primewise
agreement and (D.24.1) give \([y_v-x]=0\).
The Abel--Jacobi map of a positive-genus
smooth curve is injective: a function
whose divisor is the difference of two
distinct degree-one points would have
degree one, forcing genus zero. Thus
\[
                         y_v=x                  \tag{D.24.3}
\]
at every finite place.

If \(g(X)=0\), pass to a finite extension
splitting the punctures, so that \(X\)
is a projective line. Equality of local
points can be checked after this extension.
There are at least three punctures. For
an ordinary reference point \(x\), a ratio
\((t-b)/(t-c)\) of two distinct puncture
coordinates separates \(x\) from any
other ordinary point. Its Kummer values
agree, and (D.24.2) implies equality of
the points. If \(y_v\) is a cusp, choose
two other punctures \(b,c\). The same
function is regular and nonzero at the
cusp, so its value on every section in
that packet is its ordinary value there.
It again separates that cusp from \(x\).
This proves (D.24.3) for an ordinary
reference point in genus zero.

It remains to separate genus-zero cusps.
Put the reference cusp at zero and two
other punctures at one and infinity. The
open immersion into the tripod preserves
the section agreements. The unit \(1-t\)
has value one at the reference cusp. If
\(y_v\) maps to an ordinary tripod point,
equality of full Kummer classes gives
\(1-y_v=1\), an impossibility. Thus the
only possible competing images are the
cusps one and infinity.

Let the residue characteristic of \(v\)
be \(p\), and choose \(n=p\) for odd
\(p\), and \(n=4\) for \(p=2\). The
reference cusp's \(t\)-Kummer parameter
modulo \(n\) is represented by
\(a\in F^\times\), where \(F\) is the
local field now in use. Consider the cover
\[
       a z^n=t,\qquad w^n=1-t,
       \qquad a z^n+w^n=1.                        \tag{D.24.4}
\]
The two root torsors pull back trivially
along the reference section and along
\(s\), since their pro-\(p\) images agree.
Use representatives with exactly equal
pro-\(p\) images and the same fixed fibre
point. They lift to the same section
neighborhood of this finite etale cover.
Its geometric subgroup \(H\) is the
kernel of the rank-two mod-\(n\) abelian
quotient. Its index is a power of \(p\),
and
\[
 H^{(p)}\ \simeq\
 \text{the corresponding open subgroup of }
 \pi_1(\text{tripod}_{\bar F})^{(p)}.            \tag{D.24.5}
\]
Indeed any normal open subgroup of \(H\)
of \(p\)-power index has a core in the
full geometric group still of \(p\)-power
index: intersect its finitely many
conjugates, and note that both the kernel
quotient and the quotient by \(H\) are
\(p\)-groups. This proves (D.24.5) by
the universal property of completion.
It also shows that equality of the
original pro-\(p\) sections implies equality
of their lifted pro-\(p\) classes. In
particular their proper abelian Kummer
classes upstairs agree.

The compactification of (D.24.4) is a
diagonal Fermat curve of genus
\((n-1)(n-2)/2\ge1\). Its boundary points
are hyperflexes: the tangent line at any
point with one homogeneous coordinate
zero has intersection divisor \(nP\).
Ratios of tangent lines give
\[
                   n[P-Q]=0\quad\text{in }J       \tag{D.24.6}
\]
for any two boundary points. Each lifted
local cusp section has a rational underlying
cusp; one may also extend \(F\) finitely
to split all the boundary points. If
the original cusp images were zero and
one or infinity, their lifts \(P,Q\)
would be distinct, so \([P-Q]\ne0\)
by Abel--Jacobi injectivity. By (D.24.6)
this class is \(p\)-primary and is
detected by (D.24.1), contradicting the
equality of lifted pro-\(p\) Kummer classes.
Hence (D.24.3) also holds for cusps.
Only actual \(p\)-group covers have been
used in this change of neighborhood.

Apply the finite-support argument at
the start of D.23, now over \(k\), to
the one-point subscheme \(\{x\}\subset X\).
Stoll's Theorem 8.2 applies at all finite
places. The real section theorem supplies
the real components for the adelic
finite-descent condition, as in D.23;
support is tested only at finite places.
Cofinal section neighborhoods
produce a horizontal point fixed by the
full section. Thus \(s\) is the section
of \(x\) if \(x\) is ordinary, or lies
in its cusp packet. In the latter case
we must still compare packet parameters.

Here is primewise injectivity within
one cusp packet over a number field.
After a finite extension all \(r=\#D\)
punctures are rational. Put
\(V=\pi^{(\ell),\mathrm{ab}}\). If
\(r\ge2\), cusp inertia
\(I=\mathbb Z_\ell(1)\) embeds primitively
in \(V\), whose boundary sequence is
\[
 0\longrightarrow\mathbb Z_\ell(1)^{r-1}
 \longrightarrow V\longrightarrow T_\ell J
 \longrightarrow0.
\]
The quotient \(V/I\) is torsion free.
It has no Galois invariants: at a place
of good reduction away from \(\ell\),
the Frobenius eigenvalues on its
rationalization have complex absolute
values \(q\) or \(q^{1/2}\), never one.
Thus \(H^1(k,I)\to H^1(k,V)\) is
injective, and equality of pro-\(\ell\)
section classes gives equality of the
cusp parameters.

If \(r=1\), then \(g\ge1\), the pro-\(\ell\)
group is free of rank \(2g\), and its
boundary word has degree-two image
\[
 \omega=\sum_{i=1}^g a_i\wedge b_i
 \quad\text{in}\quad
 W=\gamma_2\pi^{(\ell)}/\gamma_3\pi^{(\ell)}
        =\bigwedge^2 T_\ell J.                   \tag{D.24.7}
\]
This is primitive, its inertia line is
\(I=\mathbb Z_\ell(1)\), and \(W/I\)
is torsion free. Both \(H^0(k,V)\)
and \(H^0(k,W/I)\) vanish by the same
Frobenius argument. A conjugator between
two cusp sections has invariant image
in \(V\), since their abelian images
coincide. It therefore belongs to
\(\gamma_2\). In \(W\), the difference
of the two inertia cocycles is its
coboundary. The injection
\(H^1(k,I)\to H^1(k,W)\), from
\(H^0(k,W/I)=0\), proves equality of
the two parameters.

Restriction to the finite extension
used to split punctures loses no inertia
cohomology: \(I^{G_L}=0\) for a number
field \(L\), so inflation--restriction
injects \(H^1(k,I)\) into \(H^1(L,I)\).
The argument holds for every prime,
including two. Cocycles and coboundaries
in a product are products of the
corresponding cocycles and coboundaries;
continuity is coordinatewise. Therefore
\[
 H^1(k,\widehat{\mathbb Z}(1))
         =\prod_\ell H^1(k,\mathbb Z_\ell(1)).
\]
The full cusp parameters coincide,
completing the proof of M1.
\end{proof}

\begin{statement}{D.25}{Proved}{Primewise coherence for finite sets, including the given CM curve}
In M0 of Lead 19, suppose that the set of
underlying rational points and cusps
\(\{x_\ell\}_\ell\) is finite. Then the full
section is geometric and every \(x_\ell\)
is the same point or the same cusp.
Consequently M0 holds for every hyperbolic
curve whose compactification has genus at
least two, for every curve with finite
\(X(k)\), and whenever all primewise
sections are cuspidal. It also holds for
the specified \(U=\{y^2=x^3-x\}-\{O\}\)
over \(K=\mathbb Q(i)\); in fact
\[
 E(K)=\{O,(0,0),(1,0),(-1,0),
          (i,\pm(1-i)),(-i,\pm(1+i))\}
       \simeq\mathbb Z/4\oplus\mathbb Z/2.        \tag{D.25.1}
\]
The local geometricity used in the cohomological
conclusion is the full-group case of D.56.
The Mordell--Weil calculation is given directly below.
\end{statement}
\begin{proof}
We first check the single-prime separation
input, with its precise scope. Mochizuki's
Theorem 19.1 \cite{DE_Mochizuki1999} states
that the pro-\(p\) point-section map for a
hyperbolic curve is injective over a
sub-\(p\)-adic field. Its Definition 15.4
includes number fields and finite extensions
of \(\mathbb Q_p\), for every prime \(p\),
including two. This is an injectivity
theorem and asserts no pro-\(p\) section
conjecture.

It also separates the underlying points
of point and cusp sections in our setting.
Indeed, if two such pro-\(p\) classes agree,
choose a common pro-\(p\) section neighborhood
whose compactification has genus at least
two. Such neighborhoods are cofinal. To
see this in the affine case, put
\(\epsilon=2g-2+r>0\). Each peripheral
element has infinite order in the free
pro-\(p\) group: for \(r\ge2\) its
abelian image is primitive; for \(r=1\)
its degree-two symplectic image is nonzero.
Choose a finite characteristic \(p\)-group
quotient in which all peripheral orders
\(e_j\) exceed \(2r/\epsilon\). For the
associated degree-\(d\) cover,
\[
 2g_Y-2=d\left(2g-2+r-\sum_{j=1}^r e_j^{-1}\right)>0.
                                                               \tag{D.25.2}
\]
Intersecting with any prescribed pro-\(p\)
neighborhood proves cofinality. The open
pro-\(p\) completion identity (D.24.5)
shows that equality of section classes
persists in the cover. The two full
geometric sections lift to rational
compactification points there, and these
remain distinct if their original
underlying point or cusp was distinct.
Theorem 19.1 on that proper hyperbolic
compactification contradicts equality.
This proves the asserted separation of
underlying points; it does not confuse
cusp parameters with their underlying cusp.

Apply this at every finite place \(v\)
of the number field, taking \(p\) to be
its residue characteristic. D.56, applied to the full fundamental group, says
that the full section has local underlying
point \(y_v\in X(k_v)\). Its pro-\(p\)
class is that of \(x_p\), so the separation
just proved gives
\[
                         y_v=x_p.                \tag{D.25.3}
\]
If \(\{x_\ell\}\) is finite, form its
finite reduced union \(Z\subset X\) over
\(k\). The finite-support argument of
D.23, including real components and
cofinal full section neighborhoods,
proves that the full section is geometric
at some point of \(Z\). Applying the
same single-prime separation theorem over
the number field, once for each \(\ell\),
shows that every \(x_\ell\) equals this
point. This proves the finite-set criterion.
If \(g(X)\ge2\), its hypothesis follows
from Faltings's theorem
\cite[Satz 7]{DE_Faltings1983}. If all
primewise sections are cuspidal, it follows
from finiteness of the boundary.

We give an independent descent calculation
for the asserted CM example. Set
\[
 E:y^2=x^3-x,\qquad E':v^2=u^3+4u,
\]
with dual two-isogenies
\[
 \phi(x,y)=(x-x^{-1},y(1+x^{-2})),\qquad
 \psi(u,v)=((u+4u^{-1})/4,v(1-4u^{-2})/8).
                                                               \tag{D.25.4}
\]
The formulas extend over the omitted
points, have the rational kernels generated
by \((0,0)\), and satisfy \(\psi\phi=[2]\).
The descent maps to \(\mathbb Q^\times/
\mathbb Q^{\times2}\) are \(x\) and \(u\)
away from their zero and pole; their
values at infinity are one, and at
\((0,0)\) are respectively \(-1\) and
\(4\). The chord formula makes them
homomorphisms: the product of the three
nonzero intersection abscissae of a line
with either cubic is a square. Repeated
intersection points use the tangent formula.
For the exceptional point \(T=(0,0)\),
the addition formula is
\(x(P+T)=b/x(P)\) on
\(y^2=x(x^2+ax+b)\); this verifies the
homomorphism identity with \(\alpha(T)=b\).
Opposite points have the same abscissa,
and the identity point has class one.

Their kernels are respectively
\(\psi E'(\mathbb Q)\) and
\(\phi E(\mathbb Q)\). Here this can be
verified without a black-box descent
algorithm. The first coordinate of
\(\psi(u,v)\) is \((v/(2u))^2\).
If \(x=b^2\ne0\), the point with
\(u=2x+2y/b\), \(v=2bu\) maps to
\((x,y)\). Likewise the first coordinate
of \(\phi(x,y)\) is \((y/x)^2\).
If \(u=b^2\ne0\), set
\(x=(u+v/b)/2\), \(y=bx\); this is a
preimage under \(\phi\). The exceptional
kernel assertions follow from (D.25.4);
in particular \((0,0)\in E'\) is the
image of \((1,0)\in E\).

For a curve \(y^2=x(x^2+ax+b)\), writing
\(x=A/C^2\), \(y=B/C^3\) in coprime
integers shows that the square class of
\(x\) has a squarefree representative
dividing \(b\): at any prime dividing
\(A\) but not \(b\), the valuation of
\(A(A^2+aAC^2+bC^4)\) forces the
valuation of \(A\) to be even. Thus the
first descent image is exactly
\(\{1,-1\}\). The second is contained
in \(\{1,2,-1,-2\}\); the equation of
\(E'\) makes \(u\ge0\) at a real point,
so it is \(\{1,2\}\), both values being
attained by infinity and \((2,4)\).
Because \(\ker\psi\subset\phi E(\mathbb Q)\),
the quotient sequence for \([2]=\psi\phi\)
now gives
\[
       \#(E(\mathbb Q)/2E(\mathbb Q))=2\cdot2=4.
\]
There are already four rational points
of order dividing two. Mordell--Weil
therefore gives rank zero. Good reduction
counts \(\#E(\mathbb F_3)=4\) and
\(\#E(\mathbb F_5)=8\) then imply
\(E(\mathbb Q)=E[2]\), by injectivity
of prime-to-residue-characteristic torsion.

For \(P\in E(\mathbb Q(i))\), both
\(P+\bar P\) and \([i](P-\bar P)\)
are rational, where
\([i](x,y)=(-x,iy)\). They are consequently
killed by two, and \(4P=0\). Thus
\(E(K)\) is finite. Its reduction at
a prime of \(K\) above five has at most
eight points and is injective on this
two-primary group. The eight points
in (D.25.1) all satisfy the equation;
the four with abscissa \(\pm i\) double
to \((0,0)\). They exhaust \(E(K)\)
and give its indicated group structure.
The identities, doubles, and counts over
\(\mathbb F_3,\mathbb F_5,\mathbb F_9\)
are independently executed in
\texttt{de\_cm\_rational.py}. Finite-set
coherence now proves M0 for \(U/K\).
\end{proof}

\begin{statement}{D.26}{Proved}{Genus-zero coherence outside the infinitely mixed case}
Let \(U/k\) be hyperbolic with genus-zero
compactification, and assume the hypotheses
of M0. If only finitely many primes have
cuspidal primewise image, or if only
finitely many have ordinary-point image,
then M0 holds: all primewise underlying
points or cusps coincide. This includes
the all-ordinary case of the tripod over
\(\mathbb Q\). A remaining counterexample
in genus zero would require infinitely
many ordinary primes, infinitely many
cuspidal primes, and infinitely many
distinct ordinary rational points.
\end{statement}
\begin{proof}
If only finitely many ordinary primes
occur, the set of underlying points is
finite, since the boundary is finite;
apply D.25. Now let \(B\) be the finite
set of cuspidal primes. Enlarge the
number field once to split all punctures
and identify three of them with
\(0,1,\infty\). This identifies the
compactification with a projective line;
the coordinate \(t\) is injective on
its points. Fix an ordinary prime
\(\ell_0\) and put \(a=t(x_{\ell_0})\).

At a finite place \(v\) of residue
characteristic \(p\notin B\), D.25's
local separation gives \(y_v=x_p\).
For any ordinary prime \(\ell\), the
units \(t\) and \(1-t\) identify the
two primewise abelian Kummer coordinates.
Taking the local valuation yields
\[
 v(t(x_\ell))=v(t(x_p))=v(a),\qquad
 v(1-t(x_\ell))=v(1-t(x_p))=v(1-a).              \tag{D.26.1}
\]
These are equalities of integers: each
equality first holds in \(\mathbb Z_\ell\)
or \(\mathbb Z_{\ell_0}\), into which
\(\mathbb Z\) injects. It is essential
here that \(y_v\) is an ordinary point;
the valuation of a completed cusp tangent
would instead lie in \(\widehat{\mathbb Z}\).

Take \(S\) to contain the archimedean
places, the places above \(B\), and the
support of the divisors of \(a\) and
\(1-a\). Equation (D.26.1) says that
\(t(x_\ell)\) and \(1-t(x_\ell)\)
are both \(S\)-units for every ordinary
\(\ell\). The two-variable \(S\)-unit
theorem \cite[Theorem 1, with \(n=2,c=1,d=0\)]{DE_Evertse1984}
makes this set finite. Its projective
tuple is \((t,1-t,-1)\); all entries
and all proper nonempty subsums are
nonzero because \(t\ne0,1\).
The underlying ordinary points are
therefore finite, and adding the finitely
many cusps still gives a finite set of
points of the original curve over \(k\).
D.25 applies over \(k\); no descent of
section conjugacy from the auxiliary
splitting field is required.

For the all-ordinary tripod over
\(\mathbb Q\), an elementary shortcut
is available: (D.26.1) holds at every
rational prime. Ratios of any two
\(t(x_\ell)\), and of their complements,
are therefore signs. The equations
\(t'=\pm t\), \(1-t'=\pm(1-t)\)
force \(t'=t\) when \(t,t'\ne0,1\).
\end{proof}

\begin{statement}{D.27}{Proved}{Mixed coherence with omitted punctures or one cusp over a finite-unit field}
Let \(U/k\) have genus-zero compactification
and satisfy the hypotheses of M0. Let
\(C\) be the set of cusps occurring among
the primewise underlying points \(x_\ell\).
If at least three geometric punctures
do not belong to \(C\), then M0 holds.
If \(k=\mathbb Q\) or an imaginary
quadratic field, at least three punctures
are \(k\)-rational, and \(C\) consists
of one cusp, then M0 also holds.
Consequently a remaining counterexample
for the tripod over \(\mathbb Q\) must
use at least two distinct cusps at its
cuspidal primes.
\end{statement}
\begin{proof}
The case \(C=\varnothing\) is D.26.
Otherwise choose a prime \(\ell_0\)
whose underlying point is a cusp
\(c_0\in C\). In the first assertion,
pass to one finite extension \(L/k\)
over which the boundary splits, and
put three punctures outside \(C\) at
\(0,1,\infty\). The coordinate \(t\)
then has the property that both \(t\)
and \(1-t\) are invertible regular
functions on \(U_L\), regular and
nonzero at every member of \(C\).

Here is the valuation argument, including
the role of the fixed cuspidal prime.
Write \(f=t\) or \(f=1-t\). Since
\(f\) is regular and nonzero at \(c_0\),
the pushforward of the \(\ell_0\)-adic
section under \(f:U_L\to\mathbb G_m\)
is exactly \(\kappa_{\ell_0}(f(c_0))\).
There is no tangent contribution: the
order of \(f\) at this cusp is zero.
At a finite place \(v\) of \(L\) whose
residue characteristic \(p\) has ordinary
primewise image, D.25 identifies the local
underlying point with \(x_p\). Thus
\[
              v(f(x_p))=v(f(c_0)).              \tag{D.27.1}
\]
This is an integer equality, obtained
first in \(\mathbb Z_{\ell_0}\).
If instead \(p\) has cuspidal image
at \(c_p\in C\), the full local section
has underlying cusp \(c_p\), again by
D.25. Its image under \(f\), in every
primewise Kummer coordinate, is
\(\kappa(f(c_p))\), independent of its
completed tangent parameter.

Let \(\ell\) be any ordinary prime.
Compare its global point section
\(x_\ell\) with these local images of
the full section. Taking valuations
in the \(\ell\)-adic coordinate gives
\[
 v(f(x_\ell))=
 \begin{cases}
 v(f(c_0)),&p\text{ has ordinary image},\\
 v(f(c_p)),&p\text{ has cuspidal image}.
 \end{cases}                                   \tag{D.27.2}
\]
Choose a finite set \(S\) of places of
\(L\) containing the archimedean places
and every divisor of \(t(c)\) and
\(1-t(c)\), for \(c\in C\). Equations
(D.27.1)--(D.27.2) imply that all
\(t(x_\ell)\) and \(1-t(x_\ell)\)
are \(S\)-units. The same two-variable
\(S\)-unit theorem used in D.26 makes
the set of ordinary \(x_\ell\) finite.
Adding \(C\) gives a finite set over
the original field \(k\), so D.25
applies there. The auxiliary extension
has only been used to prove finiteness.

For the second assertion, put the sole
occurring cusp at \(0\), and two other
\(k\)-rational punctures at \(1,\infty\).
Now use just \(f=1-t\). It has value
one at the sole cusp. The argument
above therefore gives
\[
                   v(1-t(x_\ell))=0
\]
at every finite place, for every ordinary
\(\ell\). Hence \(1-t(x_\ell)\) belongs
to \(\mathcal O_k^\times\). This group
is finite for the stated fields: over
\(\mathbb Q\) it is \(\{1,-1\}\);
for an imaginary quadratic field its
elements have complex absolute value
one and lie in the discrete lattice
\(\mathcal O_k\subset\mathbb C\).
Again the point set is finite and D.25
finishes. For the tripod over
\(\mathbb Q\), the only possible
ordinary coordinate is \(t=2\), since
\(1-t=1\) would give the excluded
cusp \(t=0\). Finite-set coherence
then excludes an actual mixture of
that ordinary point with the cusp.
\end{proof}

\begin{statement}{D.28}{Proved}{An actual class-three obstruction invisible at every finite coefficient level}
There are distinct primes \(q_r\notin\{2,3,5\}\),
with \(v_5(q_r+1)=r\), and a continuous
section \(s\) of \(T_2\rtimes G_{\mathbb Q}\)
whose abelian coordinates are
\[
 a=\kappa(3),\qquad
 b=\kappa(-2)+\sum_{r\ge1}5^r\kappa(q_r),        \tag{D.28.1}
\]
such that its obstruction
\[
  o_3(s)\in H^2_{\mathrm{cts}}
       (\mathbb Q,\mathbb Z_5(3)^{\oplus2})
\]
is nonzero and lies in the Jannsen kernel.
In particular \(s\) lifts through every
central extension with kernel
\((\mathbb Z/5^m(3))^{\oplus2}\) obtained
from \(T_3\to T_2\), but it has no
continuous lift to \(T_3\). No local or
global section conjecture is used.
More generally, the same primes give
an injectively parametrized family of
class-two sections indexed by
\(\prod_{r\ge1}\mathbb Z_5\), by replacing
\(5^r\) in (D.28.1) by \(5^ru_r\).
Every member lifts with every finite
central coefficient module; every member
with \((u_r)\notin c_0\) has a nonzero
integral obstruction. Thus there are
continuum many such section classes.
\end{statement}
\begin{proof}
We give the global cocycle construction,
the actual arithmetic Frobenius conditions,
and the continuity obstruction separately.

\emph{Integral coordinates and the arithmetic action.}
Put \(X=\log x\), \(Y=\log y\),
\(C=[X,Y]\), \(A=[X,C]\), \(B=[Y,C]\).
The free nilpotent pro-five group of class
at most three is the BCH group of the free
nilpotent \(\mathbb Z_5\)-Lie algebra on
\(X,Y\). This assertion is integral: the
only BCH denominators are two and twelve,
both units; logarithms of the ordinary Hall
group commutators are related to this Lie
basis by a triangular integral matrix with
diagonal entries one. The usual unit-tangent
action satisfies
\[
  g(x)=x^{\chi(g)},\qquad
  g(y)=f_g^{-1}y^{\chi(g)}f_g,\qquad f_g\in F',
                                                               \tag{D.28.2}
\]
as in Ihara's primary formulation
\cite[\S1, first paragraph]{DE_Ihara2000}.
It is therefore diagonal in class two.

We will use a diagonal integral action in
class three as well, and justify that
change of coordinates. The proof of D.1,
in weight two, gives
\[
 H^1(G_{\mathbb Q,\{5\}},\mathbb Z_5(2))=0,
 \qquad
 H^2(G_{\mathbb Q,\{5\}},\mathbb Z/5^m(2))=0
 \quad(m\ge1).                                 \tag{D.28.3}
\]
Indeed the residual unit characters are
\(1,\omega,\omega^2\), and twisting by
\(\omega\) leaves no invariant vector;
the residual second cohomology is zero
as in D.1. Induction in \(m\), followed
by the finite-group Mittag--Leffler argument,
proves (D.28.3). The integral localization
complex has injective differential
\(q-1\) at every \(q\ne5\), so also
\(H^1(G_{\mathbb Q},\mathbb Z_5(2))=0\).

By (D.28.2), the possible nondiagonal
term is
\(g(Y)=\chi(g)Y+d(g)B\), with
\(g(X)=\chi(g)X\). Compatibility says
that \(d(g)/\chi(g)\) is a continuous
one-cocycle of weight two. Thus there is
an integral \(u\in\mathbb Z_5\) with
\(d(g)=u\chi(g)(\chi(g)^2-1)\).
Replacing \(Y\) by \(Y-uB\) makes the
action exactly
\[
       D_\lambda(X,Y,C,A,B)
       =(\lambda X,\lambda Y,\lambda^2 C,
                       \lambda^3 A,\lambda^3 B).
                                                               \tag{D.28.4}
\]
This replacement is integral, leaves class
two unchanged, and leaves the displayed
brackets unchanged. We henceforth use it.

\emph{The Frobenius classes occur in an actual arithmetic image.}
Let \(s_0\) be the genuine section defined
by the rational point \(3\) of the tripod.
Write its class-two logarithmic cocycle as
\(aX+b_0Y+c_0C\), where
\(a=\kappa(3)\), \(b_0=\kappa(-2)\).
It is unramified outside \(2,3,5\):
both point and unit tangent extend in the
smooth pair away from these primes, so
tame specialization applies as in D.20.

Let \(H_5=\exp(\mathbb F_5X\oplus
\mathbb F_5Y\oplus\mathbb F_5C)\).
For every \(N\ge1\), the restriction
of \(s_0\bmod5\) to
\(G_{\mathbb Q(\zeta_{5^N})}\) is a
homomorphism to \(H_5\), since the
class-two action is trivial there.
Its abelian image is all \(\mathbb F_5^2\).
To verify this, the two Kummer classes
\(3,-2\) are independent over
\(\mathbb Q(\zeta_{5^N})\): valuations
at primes above three and two are
respectively \((1,0)\) and \((0,1)\),
as the cyclotomic extension is unramified
at both primes. A subgroup of \(H_5\)
surjecting on its abelianization is
\(H_5\) itself. The commutator of two
independent abelian vectors generates
its center; adjoining that center to
the abelian surjection gives all 125
elements.

Consequently the genuine finite
representation
\[
 G_{\mathbb Q}\longrightarrow
 H_5\rtimes(\mathbb Z/5^N)^\times,
 \qquad g\longmapsto(s_0(g)\bmod5,\chi(g)\bmod5^N)
                                                               \tag{D.28.5}
\]
is surjective. Its kernel over the
cyclotomic factor already surjects onto
\(H_5\), and the cyclotomic factor is
surjective. Take \(N=r+1\) and the
element
\[
                  (\exp C,-1+5^r).
                                                               \tag{D.28.6}
\]
Chebotarev supplies a prime \(q_r\)
in its conjugacy class, avoiding
\(2,3,5,q_1,\ldots,q_{r-1}\).
Thus \(v_5(q_r+1)=r\). In the local
class-two section, an integral conjugator
kills the degree-one Frobenius coordinate,
because \(q_r-1\) is a five-adic unit.
The remaining central Frobenius coordinate
is a unit modulo five. This property is
independent of the representative of the
Chebotarev conjugacy class: conjugation
by \(H_5\) changes the degree-one part
and leaves the normalized central part
unchanged, while a cyclotomic conjugator
multiplies the latter by a nonzero square.
The finite algebra is checked exhaustively
in \texttt{de\_lim1\_class3.py}.

\emph{An actual continuous class-two perturbation.}
For a weight-one cocycle \(k\), define
the continuous two-cocycle
\[
 \Omega(a,k)(g,h)=\tfrac12\chi(g)
                    \bigl(a(g)k(h)-k(g)a(h)\bigr).
                                                               \tag{D.28.7}
\]
With differential
\(\delta c(g,h)=\chi(g)^2c(h)-c(gh)+c(g)\),
the class-two BCH equation is
\(\delta c=-\Omega(a,b)\).
For \(k_r=\kappa(q_r)\), this cocycle
is defined over
\(G_{\mathbb Q,S_r}\),
\(S_r=\{3,5,q_r\}\). Its integral
second cohomology in weight two is zero.
Indeed (D.28.3) gives the base vanishing
at five, and localization at three and
\(q_r\) adds two acyclic residue complexes:
their differentials \(3-1\) and
\(q_r-1\) are units in \(\mathbb Z_5\).
Equivalently one first obtains zero with
every finite coefficient module and then
uses finiteness of the finite-ramification
\(H^1\)'s to eliminate \(\varprojlim^1\).
There is therefore a normalized continuous
cochain \(e_r:G_{\mathbb Q,S_r}\to\mathbb Z_5(2)\)
such that \(\delta e_r=-\Omega(a,k_r)\).

Choose normalized integral Kummer cocycles
and inflate these cochains to
\(G_{\mathbb Q}\). The sums
\[
 b=b_0+\sum_{r\ge1}5^rk_r,
 \qquad c=c_0+\sum_{r\ge1}5^re_r             \tag{D.28.8}
\]
converge uniformly: every cochain takes
values in \(\mathbb Z_5\), so its
\(r\)-th summand is bounded by
\(5^{-r}\). Their limits are continuous.
The differential is continuous and
\(\Omega\) is linear in its second
argument. Hence
\(\exp(aX+bY+cC)\) is an actual
class-two section \(s\). Moreover
\[
                         s\bmod5=s_0\bmod5.    \tag{D.28.9}
\]
It is unramified outside
\(\{2,3,5,q_1,q_2,\ldots\}\). At
\(q_r\), its inertia has degree-one
logarithm \(5^rY\), after choosing the
five-primary tame generator so that
\(\kappa(q_r)(\tau_r)=1\). Every
other Kummer input is a local unit and
has zero inertia coordinate there.

\emph{The exact local class-three equation.}
Fix \(q=q_r\), let \(t=5^r\), and
write \(\sigma\tau\sigma^{-1}=\tau^q\).
The geometric action is unramified at
\(q\), and wild and prime-to-five tame
inertia have trivial geometric image.
Suppose first that a local class-three
lift is given. Write its Frobenius
logarithm in degree one as \(w_1\).
Conjugate by \(\exp(w_1/(q-1))\).
This is integral and kills \(w_1\)
exactly. The class-two tame equation
then forces the inertia degree-two
coordinate to be zero, since its
coefficient is \(q^2-q\), a unit.
The Frobenius central coordinate
\(\beta C\) is a unit by (D.28.9)
and the choice (D.28.6). Thus the
normalized logarithms have the form
\[
       \log s(\sigma)=\beta C+w_3,
       \log s(\tau)=tY+z_AA+z_BB,
       \qquad \beta\in\mathbb Z_5^\times.       \tag{D.28.10}
\]
The same normalization is available for
the prescribed class-two section before
choosing a lift.

By (D.28.4), the logarithm of the tame
relation is
\[
 \exp(\operatorname{ad}(\beta C+w_3))
       D_q(tY+z_AA+z_BB)
       =q(tY+z_AA+z_BB).
\]
There are no omitted terms of degree
at most three. Since \([C,Y]=-B\),
this is exactly
\[
       (q^2-1)z_A=0,\qquad
       (q^2-1)z_B=t\beta.                      \tag{D.28.11}
\]
All signs, normalizations, and BCH terms
in this calculation are independently
verified symbolically by
\texttt{de\_lim1\_class3.py}. In particular
\[
         z_A=0,\qquad
         z_B=\frac{5^r\beta}{q_r^2-1}
                    \in\mathbb Z_5^\times.     \tag{D.28.12}
\]
Here \(v_5(q_r^2-1)=r\). Conversely
(D.28.12), with any degree-three
Frobenius coordinate, constructs a local
lift: it satisfies the single tame
relation and hence defines a continuous
homomorphism from the tame local quotient.
The unramified procyclic factor causes
no additional condition, since the
closed subgroup generated by any chosen
Frobenius lift is a continuous quotient
of \(\widehat{\mathbb Z}\).

\emph{Finite-level vanishing and integral nonvanishing.}
The section lifts locally at every
\(q_r\) by the preceding construction.
At primes outside the displayed
ramification set it has an unramified
lift. At two and three the local
weight-three second cohomology is zero,
since \(2^2-1\) and \(3^2-1\) are
five-adic units. At five it is zero
by local duality: the residual character
on the dual \(\mathbb F_5(-2)\) is
nontrivial. The local finite cohomology
groups are finite, so no local derived-limit
term is being suppressed here. At the
real place, positive-degree five-primary
cohomology vanishes.

For every \(m\), the finite-coefficient
localization calculation in D.1 is
\[
 H^2(\mathbb Q,\mathbb Z/5^m(3))
   \simeq\bigoplus_{q\ne5}
       (\mathbb Z/5^m)/(q^2-1)(\mathbb Z/5^m).
                                                               \tag{D.28.13}
\]
The base \(H^2\) and \(H^3\) at five
are both zero, exactly as proved there.
The components in (D.28.13) are the
local restriction obstructions. The
local lifts thus prove
\(o_3(s)\bmod5^m=0\) for every \(m\).

If an integral continuous lift existed,
its reduction modulo five in the BCH
lattice, together with the residual
cyclotomic action, would be a finite
continuous Galois representation. It
would be unramified outside finitely
many rational primes. But (D.28.12)
forces its inertia at every \(q_r\)
to have a nonzero degree-three
coordinate modulo five. The local
conjugations are integral; they preserve
the property of being the identity
modulo five. This is a contradiction.
A lift whose projection is merely
conjugate to \(s\) can be conjugated
by a lift of that class-two conjugator,
so it gives the same contradiction.
Thus \(o_3(s)\ne0\). Jannsen's exact
sequence, with (D.28.13), identifies
this nonzero obstruction with an element
of the claimed \(\varprojlim^1\) kernel.
The integral failure is proved for
the actual global section, without
deducing global nonvanishing from
zero local cohomology classes.

For the final assertion, replace the
coefficients in (D.28.8) by \(5^ru_r\).
All sums still converge uniformly and
their reductions modulo five still
equal the seed. At \(q_r\), formula
(D.28.12) becomes
\[
                 z_B=\frac{5^ru_r\beta_r}{q_r^2-1},
                 \qquad v_5(z_B)=v_5(u_r),      \tag{D.28.14}
\]
with the equality understood also for
zero coordinates. It is always integral,
so all local and finite-coefficient
obstructions still vanish. If \((u_r)\)
does not tend to zero, choose \(M\)
such that infinitely many \(u_r\) are
nonzero modulo \(5^M\). An integral
lift would then ramify modulo \(5^M\)
at infinitely many \(q_r\), the same
contradiction. Distinct sequences give
distinct global abelian Kummer classes,
as their valuations at \(q_r\) differ
by \(5^r(u_r-u'_r)\); hence their
class-two section classes are distinct.
For example the sequences in
\(\{1,2\}^{\mathbb N}\) already give
continuum many nonliftable classes.
No sufficiency assertion for
\((u_r)\in c_0\) is needed or made.
\end{proof}

\begin{statement}{D.29}{Proved}{The integral tame-relator criterion with its full convergence condition}
Let \(n\ge2\), \(L_n=\gamma_nF^{(5)}/
\gamma_{n+1}F^{(5)}\simeq\mathbb Z_5(n)^{d_n}\),
and let \(s\) be an actual continuous
cocycle in \(T_{n-1}\). Choose the
continuous Hall-coordinate set-section
\(j:T_{n-1}\to T_n\) with highest
coordinates zero and \(j(1)=1\).
For every \(q\ne5\), choose a five-primary
tame inertia generator \(\tau_q\) and
arithmetic Frobenius \(\sigma_q\), with
\(\sigma_q\tau_q\sigma_q^{-1}=\tau_q^q\)
in the tame quotient. Write
\[
  v_q=s(\tau_q),\quad w_q=s(\sigma_q),\quad
  \Delta_q=j(w_q)\,\sigma_q(j(v_q))\,
                    j(w_q)^{-1}j(v_q)^{-q}\in L_n,
  \qquad R_q=q^{-1}\Delta_q.                    \tag{D.29.1}
\]
Central elements in this formula are
written additively after the indicated
group product has been formed. Then
\((R_q)_{q\ne5}\in c_0^{d_n}\), and,
in the residue convention of D.1,
\[
 o_n(s)=[(R_q)_q]\in
 \left(c_0/\operatorname{diag}(q^{n-1}-1)c_0\right)^{d_n}.
                                                               \tag{D.29.2}
\]
The sign of this convention gives the
class-two residue
\(u_q\beta_q-v_q\alpha_q\), as in R16.
A continuous lift exists if and only if
\[
 R_q\in(q^{n-1}-1)L_n\quad\text{for every }q\ne5,
 \qquad
 \left(\frac{R_q}{q^{n-1}-1}\right)_q\in c_0^{d_n}.
                                                               \tag{D.29.3}
\]
There is no additional condition at five.
Every nonempty fibre of section classes
is a torsor under \(H^1(\mathbb Q,L_n)\).
All statements use the actual arithmetic
action and integral Hall powers; no
class-five BCH integrality is assumed.
\end{statement}
\begin{proof}
First extend the finite-coefficient
calculation of D.1 to every Tate weight.
Over \(\mathbb Q(\zeta_5)\), the
residual Tate characters are trivial,
the class group is trivial, and the
single prime above five leaves no
five-primary Brauer invariant. Averaging
by its group of order four gives
\(H^2(G_{\mathbb Q,\{5\}},\mathbb F_5(n))=0\)
for every integer \(n\). Induction in
the coefficient exponent and the
cohomological-dimension bound used in
D.1 give
\[
 H^2(G_{\mathbb Q,\{5\}},L_n/5^m)=
 H^3(G_{\mathbb Q,\{5\}},L_n/5^m)=0.            \tag{D.29.4}
\]
In particular this also holds when
\(n\equiv1\pmod4\). The finite-ramification
\(H^1\) is finite. Localization therefore
gives, with \(M_m=L_n/5^m\),
\[
\begin{split}
 H^2(\mathbb Q,M_m)&\simeq
       \bigoplus_{q\ne5}M_m/(q^{n-1}-1)M_m,\\
 H^1(\mathbb Q,M_m)&\longrightarrow
       \bigoplus_{q\ne5}\ker(q^{n-1}-1:M_m\to M_m)
                     \longrightarrow0.
\end{split}
\tag{D.29.5}
\]
The kernel of the second arrow is
\(H^1(G_{\mathbb Q,\{5\}},M_m)\).

Integral Hall coordinates give the
specified set-section \(j\). To be
explicit about the integrality issue,
each Hall generator is an actual iterated
group commutator, and normal forms are
ordered products of its \(\mathbb Z_5\)-powers.
Collection and powers are integer-valued
Hall polynomials. Such polynomials extend
continuously from integers to \(\mathbb Z_5\)
and remain integral, including terms
\(\binom{a}{5}\) in class five. Unique
normal forms follow by comparing successive
lower-central quotients with the integral
Hall basis; existence follows by successive
collection and completeness. Thus appending zero
top coordinates to the lower coordinates
gives a continuous set-section, even when logarithmic
coordinates would require division by five.

At \(q\ne5\), the tangential geometric
action is unramified, by specialization
of the tripod's smooth pair over
\(\mathbb Z[1/5]\). On inertia the cocycle
is a homomorphism into a pro-five group.
It kills wild inertia and the prime-to-five
tame factors. Its lower-stage tame
relation is exactly the projection of
(D.29.1), proving \(\Delta_q\in L_n\).
If the chosen lift of \(v_q\) is changed
to \(j(v_q)z_q\), with \(z_q\in L_n\),
then
\[
          \Delta_q\longmapsto
            \Delta_q+(q^n-q)z_q.               \tag{D.29.6}
\]
Changing the lift of Frobenius by a
central element has no effect. Local
lifting is therefore equivalent to
the first condition in (D.29.3), and
the unique integral inertia correction
is \(z_q=-R_q/(q^{n-1}-1)\).

The sequence \(v_q\) tends to the
identity. Indeed, every finite quotient
of a continuous Galois section is
unramified outside finitely many primes.
The same statement follows directly
from \(\tau_q\to1\) in \(G_{\mathbb Q}\)
and continuity of the cocycle. The
Hall section sends \(v_q\) to elements
tending to one. Its conjugators range
over a compact group, and the arithmetic
automorphisms range over a compact
equicontinuous group. Formula (D.29.1)
thus gives \(R_q\to0\). The scalar
\(q^{-1}\) is always a five-adic unit.

We verify that this is the integral
global residue class, rather than only
its family of local reductions. The
global continuous set-lift \(j\circ s\)
defines the obstruction two-cocycle.
On five-primary tame inertia, replace
its set-lift by the homomorphism
\[
          \tau_q^a\longmapsto j(v_q)^a,
                    \qquad a\in\mathbb Z_5.
\]
The correcting one-cochain is the
central element
\[
              j(v_q)^a j(v_q^a)^{-1}.           \tag{D.29.7}
\]
It tends uniformly to zero as \(q\)
tends to infinity: the maps involved
are continuous on the compact space
\(\mathbb Z_5\times T_{n-1}\), and
\(v_q\to1\). Hence these homotopies
have finite support after reduction
modulo any \(5^m\), and are compatible
as \(m\) varies. They may therefore
be used in the direct-sum localization
complex and in its derived inverse limit.

In that complex, the inertia cohomology
\(H^1(I_q,L_n)\) is identified with
\(L_n(-1)\) by evaluation on \(\tau_q\).
The defect of Frobenius equivariance
of the displayed homomorphic inertia
lift is evaluated on \(\tau_q\) as
\[
 \widetilde\sigma_q\,
       \widetilde\tau_q^{\,1/q}\,
       \widetilde\sigma_q^{-1}
       \widetilde\tau_q^{-1}
              =q^{-1}\Delta_q.                 \tag{D.29.8}
\]
This is precisely the inhomogeneous
degree-one residue cochain for Frobenius.
Its differential on a change of inertia
lift is \(q^{n-1}-1\), as also follows
from (D.29.6). The compatible, uniformly
vanishing homotopies (D.29.7) identify
the global obstruction with the actual
\(c_0\) family (D.29.8) in the
localization triangle of D.1. This
proves (D.29.2), including the
nonseparated part. Replacing \(j\) by
another continuous identity-preserving
set-section changes the family by
\((q^{n-1}-1)h(v_q)\), with
\(h(v_q)\to0\), so the displayed
cokernel class is unchanged.

For completeness there is an independent
proof of the sufficiency of (D.29.3)
that does not depend on naming the
global residue representative. The
first line of (D.29.5) shows that
finite central lifts of \(s\) exist
for every \(m\). In any such lift,
the difference of its inertia correction
from the prescribed \(z_q\bmod5^m\)
belongs to
\(\ker(q^{n-1}-1:M_m\to M_m)\).
This family has finite support. For
the given lift, that follows from
continuity of its finite central
coordinate; for the prescribed corrections,
it follows from \(z_q\to0\). The
surjection in the second line of
(D.29.5) supplies one global one-cocycle
which corrects all those inertia
coordinates simultaneously.

Thus the set of exact finite central
lifts whose inertia corrections equal
\(z_q\bmod5^m\) is nonempty. It is
finite: differences of any two such
lifts are unramified outside five,
and their classes form a torsor under
the finite group
\(H^1(G_{\mathbb Q,\{5\}},M_m)\);
the set of coboundaries is finite as
well, being an image of \(M_m\).
Reduction maps preserve these prescribed
inertia coordinates. An inverse system
of finite nonempty sets indexed by
the positive integers has nonempty
inverse limit: every finite collection
of compatibility equations is solved
by a point at its largest level, and
compactness applies. A compatible
sequence of exact lifts is a continuous
lift to \(T_n\), since continuity into
the inverse limit is coordinatewise.
This proves sufficiency and explicitly
retains the convergence requirement.
Necessity follows either from (D.29.2)
or from continuity of an integral
lift's central inertia corrections.

There is no missing condition at five:
the first line of (D.29.5) is a global
isomorphism, so vanishing of all the
displayed finite residues already
annihilates the global obstruction.
In weights congruent to one modulo
four, local second cohomology at five
need not vanish; this argument does
not assert that it does.

Finally the usual central lifting
torsor has no further conjugacy
stabilizer. A fixed geometric element
of \(T_{n-1}\) for the section-twisted
action would have a first nonzero
lower-central image fixed by a positive
power of \(\chi\). That graded module
is torsion free and has no invariants.
Thus the element must be the identity.
The fibre is consequently exactly an
\(H^1(\mathbb Q,L_n)\)-torsor.
\end{proof}

\begin{statement}{D.30}{Proved}{The evaluated class-three obstruction of an arbitrary Kummer pair}
Let \((\alpha,\beta)\) be an admissible
pair in R16. Fix once an element
\(g_*\in G_{\mathbb Q}\) with
\(\chi_5(g_*)=2\). Represent the two
Kummer classes by their unique cocycles
\(a,b\) satisfying
\(a(g_*)=b(g_*)=0\), and let \(c\)
be the unique continuous integral
cochain satisfying
\[
 \delta c=-\tfrac12(a\smile b-b\smile a),
                   \qquad c(g_*)=0.             \tag{D.30.1}
\]
For \(q\ne5\), put
\[
\begin{array}{lll}
 a_q=\alpha_q=a(\tau_q),&
 b_q=\beta_q=b(\tau_q),&c_q=c(\tau_q),\\
 u_q=a(\sigma_q),&v_q=b(\sigma_q),&d_q=c(\sigma_q).
\end{array}
\]
Then
\[
 c_q=-\frac{u_q\beta_q-v_q\alpha_q}{q-1},
 \quad
 R_{A,q}=\frac{q+1}{2}c_qu_q-d_q\alpha_q,
 \quad
 R_{B,q}=\frac{q+1}{2}c_qv_q-d_q\beta_q.          \tag{D.30.2}
\]
The two \(c_0\) families \((R_{A,q})\),
\((R_{B,q})\) represent the actual
class-three obstruction in \((C/D_3C)^2\).
The lift criterion is their pointwise
divisibility by \(q^2-1\), with both
divided families tending to zero.
Every nonempty fibre is a
\(\mathbb Z_5^2\)-torsor. These formulas
settle Task D.2, including its integral
convergence part.
\end{statement}
\begin{proof}
The character is surjective onto
\(\mathbb Z_5^\times\), so \(g_*\)
exists. A change of representative of
a given Kummer class is the coboundary
\((\chi-1)t\); since \(\chi(g_*)-1=1\),
the required normalized representatives
exist uniquely. Admissibility in R16
is precisely existence of a continuous
integral solution to (D.30.1) before
normalization. By D.28.3 and its
unrestricted consequence,
\(H^1(\mathbb Q,\mathbb Z_5(2))=0\).
The difference of two such solutions
is therefore \((\chi^2-1)t\).
Since \(\chi(g_*)^2-1=3\) is a unit,
the normalization makes \(c\) unique.
Thus (D.30.1) specifies one actual
cochain from the input Kummer pair,
without an unchosen class-two torsor.

Use the integral diagonal class-three
coordinates of D.28.4. At a fixed
prime suppress the subscript and write
\[
        V=a_qX+b_qY+c_qC,
        W=u_qX+v_qY+d_qC.
\]
The class-two tame relation is
\[
       (q-1)c_qC=-[u_qX+v_qY,a_qX+b_qY],
                                                               \tag{D.30.3}
\]
proving the first formula of (D.30.2).
Here \(q-1\ne0\), and the quotient is
integral because the input is an actual
class-two section. The BCH set-lift
with zero degree-three logarithmic
coordinate preserves the powers of
each inertia element. Its relator is
\[
 \Delta_q=
 \exp\bigl(\exp(\operatorname{ad}W)D_qV\bigr)
                       \exp(-qV).
\]
By (D.30.3) its degree-one and
degree-two logarithms vanish, so its
degree-three logarithm is simply the
degree-three part of
\(\exp(\operatorname{ad}W)D_qV-qV\).
After dividing by \(q\), this is
\[
 q[W_1,c_qC]+[d_qC,V_1]
                    +\tfrac12[W_1,[W_1,V_1]].
\]
Substituting (D.30.3) reduces it to
\[
             \frac{q+1}{2}c_q[W_1,C]
                         -d_q[V_1,C],          \tag{D.30.4}
\]
whose \(A,B\) coefficients are exactly
(D.30.2). The extended exact symbolic
certificate \texttt{de\_lim1\_class3.py}
checks (D.30.3)--(D.30.4), including
the final BCH multiplication, without
discarding a noncentral error term.

Because \(a_q,b_q,c_q\to0\) and
\(u_q,v_q,d_q\in\mathbb Z_5\), both
families in (D.30.2) belong to \(c_0\).
D.29 now identifies their global
cokernel classes and proves both
directions of the criterion. Its
torsor statement and
\(H^1(\mathbb Q,\mathbb Z_5(3))=\mathbb Z_5\)
give the claimed fibre. Changing the
fixed normalization changes the
families by \(D_3c_0\), as follows
also from D.29; the criterion is a
property of the original section class.
\end{proof}

\begin{statement}{D.31}{Proved}{The actual arithmetic constants needed through class five}
Use the standard interval path between the unit
tangents at zero and one. Write
\(X=\log x\), \(Y=\log y\), and set
\[
 C=[X,Y],\quad U=[X,C],\quad V=[Y,C],\qquad
 E_1=[X,U],\quad E_2=[Y,U],\quad E_3=[Y,V].
\]
For the actual arithmetic Frobenius \(g=\sigma_q\)
at any rational prime \(q\ne5\), the Ihara action
\(g(x)=x^q\), \(g(y)=f_g^{-1}y^qf_g\) has
\[
 \log f_g\equiv r_qC+h_q(U+V)
             +c_q(E_1+E_3)+d_qE_2\pmod{\text{degree}\ge5},
                                                        \tag{D.31.1}
\]
where the constants, in these precise conventions, are
\[
 r_q=\frac{q^2-1}{24},\qquad
 c_q=-\frac{q^4-1}{1440},\qquad
 d_q=-\frac{(q^2-1)(q^2+11)}{5760},\qquad
 h_q=\frac{\kappa_3([\sigma_q,c_\infty])}{96}.
                                                        \tag{D.31.2}
\]
Here \(c_\infty\) is the complex conjugation
determined by the same real interval, and
\([a,b]=aba^{-1}b^{-1}\). The raw Soul\'e character
\(\kappa_3\) on \(\ker\chi_5\) is normalized by
\[
 u_m=\prod_{\substack{1\le a<5^m\\5\nmid a}}
                  (1-\zeta_{5^m}^{a})^{a^2},\qquad
 \frac{\delta(u_m^{1/5^m})}{u_m^{1/5^m}}
       =\zeta_{5^m}^{\,\kappa_3(\delta)\bmod5^m}.
                                                        \tag{D.31.3}
\]
All four coefficients in (D.31.2) belong to
\(\mathbb Z_5\). Thus (D.31.1) determines the
actual integral arithmetic action through class
five; its parameter \(h_q\) is not a freely
chosen local constant.
\end{statement}
\begin{proof}
The unit-tangent Galois action belongs to the
Grothendieck--Teichm\"uller group and has
\(f_g\in F'\), with the two-cycle, three-cycle,
and five-cycle relations
\cite[\S1]{DE_Ihara2000}. The two-cycle identity
\(f(y,x)f(x,y)=1\) gives precisely the shape
in (D.31.1): interchange of \(X,Y\) sends
\(C\) to \(-C\), interchanges \(U,V\) with
minus signs, interchanges \(E_1,E_3\) with
minus signs, and sends \(E_2\) to \(-E_2\).

Here is an explicit computation of the even
coefficients, independent of the supplied R17
program. Put \(z=(xy)^{-1}\), \(m=(q-1)/2\).
The coefficient of \(XY\) in
\[
 x^m f(z,x)z^m f(y,z)y^m f(x,y)=1
\]
is \(3r-(q^2-1)/8\); hence
\(r=(q^2-1)/24\). For degree four use the
group pentagon in the pure braid group on four
strings, with \(A=A_{12}\), \(B=A_{23}\),
\(D=A_{13}\), \(V_0=A_{24}\), \(W=A_{34}\):
\[
 f(A,BV_0)f(DB,W)
       =f(B,W)f(AD,V_0W)f(A,B).                 \tag{D.31.4}
\]
This is precisely relation (III) of
\cite[\S1, manuscript p.4]{DE_LochakSchneps1994};
the pure braid generators used below are
defined there on manuscript p.5.
The following specifies a direct, checkable
coefficient certificate for this substitution.
On the free group with generators \(z_0,\ldots,z_3\),
the Artin generator \(s_i\) sends
\[
 z_i\longmapsto z_i z_{i+1}z_i^{-1},\qquad
 z_{i+1}\longmapsto z_i
\]
and fixes the other generators. Its inverse sends
\(z_i\) to \(z_{i+1}\) and \(z_{i+1}\) to
\(z_{i+1}^{-1}z_i z_{i+1}\). Define
\(A_{ij}=s_{j-1}\cdots s_{i+1}s_i^2
s_{i+1}^{-1}\cdots s_{j-1}^{-1}\), with
zero-based indices for this calculation. Products
act on the right. These formulas satisfy the
braid relations by direct substitution.

Substitute \(z_i=1+Z_i\) and discard words of
length greater than five. In the image of \(Z_0\),
the coefficients of the words \(Z_0^2Z_1Z_2Z_3\)
and \(Z_0^2Z_1Z_3Z_2\) in the difference of the
two sides of (D.31.4), with the right side
subtracted from the left side or vice versa
consistently, give the equivalent equations
\[
 \begin{pmatrix}-3&2\\1&1\end{pmatrix}
       \binom{c}{d}
       +\binom{-r^2}{r/12+r^2/2}=0.             \tag{D.31.5}
\]
The determinant is \(-5\), which is nonzero
over \(\mathbb Q_5\); no reduction of this
matrix modulo five is made. Solving gives
\[
       c=-r/30-2r^2/5,\qquad
       d=-r/20-r^2/10.
\]
Substitution of \(r=(q^2-1)/24\) gives
(D.31.2). The exact calculation
\texttt{code/de\_gt\_low.py} records these two
Artin rows, checks every other row through
degree five, and verifies that the column of
the arbitrary coefficient \(h\) is zero.
It also checks the full hexagon through degree
four after substitution. The calculation uses
the geometric Artin action itself; it needs no
faithfulness assertion about a separately
presented free-fibre quotient.

We next determine the odd coefficient, including
its normalization. The path cocycle identity is
\(f_{gh}=f_g\,g(f_h)\). Its degree-three part
shows that \(h\) is a genuine
\(\mathbb Z_5(3)\)-valued one-cocycle: the
first non-diagonal term of the action on
\(C\) has degree four, so no degree-three
cross term intervenes. Complex conjugation
fixes the real defining path and satisfies
\(f_{c_\infty}=1\). Consequently
\[
        h([g,c_\infty])=2h(g),\qquad
        [g,c_\infty]\in\ker\chi_5.             \tag{D.31.6}
\]
On \(\ker\chi_5\), formula (2.4),
Proposition 2.2, and Lemma 2.4 of
\cite{DE_NakamuraSakugawaWojtkowiak2017}
give, at the endpoint \(1\) with its interval
tangent,
\[
 \operatorname{coeff}_{U}(\log f_\delta)
       =-\ell i_3(1)(\delta)
       =-\frac{\widetilde\kappa_3(\delta)}{2},
 \qquad
 \widetilde\kappa_3
       =\frac{\kappa_3}{1-5^2}.
\]
The raw character in that paper's (1.2) and
(2.11) is exactly (D.31.3). Their path convention
\(\gamma=f_g\,g(\gamma)\) is the one giving
\(g(y)=f_g^{-1}y^{\chi_5(g)}f_g\) here.
The same \S2.1 explicitly permits number fields;
this is not an application of its later local
Coleman formula at a Frobenius outside five.
Thus \(h(\delta)=\kappa_3(\delta)/48\).
Equation (D.31.6) proves the final formula
in (D.31.2).

The only nonunit denominator in \(c_q,d_q\)
is a single factor of five. For a five-adic
unit \(q\), \(q^4-1\) is divisible by five.
Also \((q^2-1)(q^2+11)\) vanishes modulo five:
the two possibilities for \(q^2\) are one
and four. Hence these coefficients are integral.
The coefficient \(h\) is integral since logarithms
through degree three have five-integral
denominators; (D.31.3) also makes this clear.
Finally, a term of degree at least five in
\(\log f_g\) can affect the conjugation of
\(y^q\) only in degree at least six. No unknown
coefficient of \(f_g\) remains in the action
needed for the class-five calculation.
\end{proof}

\begin{statement}{D.32}{Proved}{Nine explicit global residue polynomials and all lifting fibres through class five}
Order the standard Lyndon words by weight and
then lexicographically, with \(X<Y\):
\[
 \begin{array}{c|l}
 1&X,Y\\
 2&XY\\
 3&XXY,XYY\\
 4&XXXY,XXYY,XYYY\\
 5&XXXXY,XXXYY,XXYXY,XXYYY,XYXYY,XYYYY.
 \end{array}                                             \tag{D.32.1}
\]
Let \(H_X=x,H_Y=y\), and recursively let
\(H_w=[H_u,H_v]\), where \(w=uv\) is its
standard Lyndon factorization with longest
proper Lyndon suffix \(v\). Products of the
\(H_w\)'s in the displayed order define
integral Hall coordinates. For an actual
\(T_{n-1}\)-section, \(n=4\) or \(5\), write
its inertia and Frobenius coordinates at \(q\ne5\)
as \((a_1,\ldots,a_k)\),
\((b_1,\ldots,b_k)\), where \(k=5\) or \(8\),
respectively. Use the original interval coordinates
of D.31, not the diagonalized class-three
coordinates of D.30.

The following finite formulas evaluate every
residue. Work in
\[
 \mathcal A=\mathbb Q_5\langle\mathsf X,\mathsf Y\rangle/
                 (\text{words of length}>5),
\]
put
\(x=1+\mathsf X,y=1+\mathsf Y\), and use
\[
 G^t=\sum_{j=0}^{5}\binom tj(G-1)^j,
 \quad X=\sum_{j=1}^{4}(-1)^{j+1}\mathsf X^j/j,
 \quad Y=\sum_{j=1}^{4}(-1)^{j+1}\mathsf Y^j/j.
                                                        \tag{D.32.2}
\]
Form the expression \(F\), truncated in degree
four, by exponentiating the explicit right side
of (D.31.1), with \(h=h_q\) from (D.31.2).
The exponential only involves \(1+L+L^2/2\),
since \(L\) starts in degree two. Substitute
\[
 \alpha_q(x)=x^q,\qquad
 \alpha_q(y)=F^{-1}y^qF,
 \quad A=\prod_{|w|<n}H_w^{a_w},\quad
 B=\prod_{|w|<n}H_w^{b_w},\quad
 P=B\alpha_q(A)B^{-1}A^{-q}.                    \tag{D.32.3}
\]
For a word \(w\), write \(p_w=[w](P-1)\).
Then the three and six residue polynomials are
\[
 \boldsymbol R_4=q^{-1}
       (p_{XXXY},p_{XXYY},p_{XYYY}),             \tag{D.32.4}
\]
\[
 \begin{aligned}
 \boldsymbol R_5=q^{-1}\bigl(&
 p_{XXXXY},p_{XXXYY},p_{XXYXY}+2p_{XXXYY},\\
 &p_{XXYYY},p_{XYXYY}+3p_{XXYYY},p_{XYYYY}\bigr).
 \end{aligned}                                          \tag{D.32.5}
\]
Here the words mean words in
\(\mathsf X,\mathsf Y\) for coefficient extraction.
The appendix gives all nine expressions expanded
as rational polynomials in
\(q,h,a_1,\ldots,a_8,b_1,\ldots,b_8\).
They take integral values at the stated inputs.
They are the actual global obstruction families
in \((c_0/D_nc_0)^{r_n}\), with
\(D_n=(q^{n-1}-1)_q\) and \((r_4,r_5)=(3,6)\).
Consequently
\[
 s\text{ lifts to }T_n
 \quad\Longleftrightarrow\quad
 \left\{
 \begin{array}{l}
  \boldsymbol R_{n,q}\in(q^{n-1}-1)\mathbb Z_5^{r_n}
                              \quad\hbox{for every }q\ne5,\\
  \boldsymbol R_{n,q}/(q^{n-1}-1)\longrightarrow0.
 \end{array}\right.                            \tag{D.32.6}
\]
Every nonempty fibre at class four is a torsor
under \((\mathbb Z/5)^3\), and every nonempty
fibre at class five is a torsor under
\(\mathbb Z_5^6\).
\end{statement}
\begin{proof}
The truncated Magnus algebra uses actual group
multiplication and binomial powers. In particular,
it does not use an integral Baker--Campbell--Hausdorff
law in class five. The Hall basis theorem gives
unique coordinates in the free nilpotent
pro-five group. One can check the needed basis
directly: the leading homogeneous coefficients
of the recursive commutators in (D.32.1),
evaluated at the Lyndon monomials of that weight,
form the matrices
\[
 M_4=I_3,\qquad
 M_5=\begin{pmatrix}
 1&0&0&0&0&0\\
 0&1&0&0&0&0\\
 0&-2&1&0&0&0\\
 0&0&0&1&0&0\\
 0&0&0&-3&1&0\\
 0&0&0&0&0&1
 \end{pmatrix}.                                \tag{D.32.7}
\]
These are unitriangular over \(\mathbb Z\).
Successively removing the corresponding
commutator power recovers every coordinate;
completion gives the same procedure for
\(\mathbb Z_5\)-valued coordinates. This also
proves continuity of the set-section setting
the top Hall coordinates to zero.

For an actual lower-stage section, all
homogeneous coefficients of \(P-1\) below
weight \(n\) vanish, by the tame relation.
Thus its degree-\(n\) part lies in the free
Lie lattice of that degree. Multiplying the
vector of its Lyndon coefficients by
\(M_n^{-1}\) gives its Hall coordinates.
The inverses of (D.32.7) give exactly
(D.32.4)--(D.32.5), including the factors
two and three. D.31 supplies the actual
arithmetic action in (D.32.3). The resulting
Hall vector is therefore \(q^{-1}\Delta_q\)
of D.29, with the same continuous set-section.
D.29 proves both its global interpretation,
including the nonseparated term, and the
equivalence (D.32.6).

The displayed rational polynomials are integral
because they compute these integral group
coordinates. Denominators are evaluated before
reduction. This point matters at class five:
\(x^5=(1+\mathsf X)^5\) has coefficient one
at \(\mathsf X^5\), whereas \(x^0\) has
coefficient zero there, even though the two
exponents agree modulo five. Reducing each
Hall exponent modulo five before evaluating
the degree-five binomial polynomial is invalid.

The independent program
\texttt{code/de\_tame\_hall.py} implements
(D.32.2)--(D.32.3) by truncated noncommutative
polynomials, and extracts the coordinates by
the recursive commutators themselves. It gives
the matrices (D.32.7) and all nine expanded
polynomials, with respective term counts
\[
        (45,66,48),\qquad
        (91,169,170,168,190,98).
\]
The program also checks thirty exact Hall
collection round trips, the above precision
example, and one hundred local normalized
rows used below. The JSON records exact rational
numbers and expanded polynomial expressions;
\texttt{code/de\_hall\_tex.py} produces the
following complete polynomial appendix.
These are symbolic finite polynomial identities,
not numerical approximations to an unknown
Frobenius action.

We make the three class-four parameters explicit.
The one-cocycle
\[
                  z_4(g)=\frac{\chi_5(g)^4-1}{5}
                                                        \tag{D.32.8}
\]
represents a generator of
\(H^1(\mathbb Q,\mathbb Z_5(4))=\mathbb Z/5\).
It is integral, is killed by five as a
cohomology class, and is nonzero since no
integral cobounding constant can be \(1/5\).
After choosing any one lift to class four,
all its lifts over the given class-three
section are obtained by adding
\(z_4(g)(t_1,t_2,t_3)\),
\(t_i\in\{0,1,2,3,4\}\), in the central
Hall basis \(H_{XXXY},H_{XXYY},H_{XYYY}\).
At a prime \(q\ne5\) this leaves the
inertia coordinates unchanged and changes
\[
      (b_6,b_7,b_8)\longmapsto
      (b_6,b_7,b_8)+\frac{q^4-1}{5}(t_1,t_2,t_3).
                                                        \tag{D.32.9}
\]
Substitute all 125 choices into the six
polynomials (D.32.5). This specifies which
class-four lifts continue, with no quotient
by an uncomputed centralizer. The two fibre
groups now follow from D.13 and D.29.

Here is the precise comparison with R17. At
\(q=2\) or \(3\), a unique integral geometric
conjugation kills Frobenius coordinates of
weights one through three, successively,
because \(q^j-1\) is a unit for \(j=1,2,3\).
The actual coefficient \(h_q\) can also be
removed for this local calculation: use the
integral change of generators fixing \(x\)
and sending \(y\) to \(k^{-1}yk\), where
\[
             \log k=-\frac{h_q}{q^3-1}(U+V).
\]
This conjugates the arithmetic action through
class five to the action with \(h=0\).
The denominator is a unit, and all exponential
denominators involved are units because
\(\log k\) starts in degree three. The code
checks this conjugation through degree five
at \(q=2,3,7,13\), retaining the actual
degree-two and degree-four arithmetic terms.

Write the remaining Frobenius coordinate in
the Lie basis \(E_1,E_2,E_3\) as
\(W=w_1E_1+w_2E_2+w_3E_3\), and put
\[
        v_i=\left(\frac{q^4-1}{5}\right)^{-1}w_i
                     \pmod5,
        \qquad (a,b)=(a_1,a_2)\pmod5.
\]
At two and three these inverses are integral
units. The first three tame relations uniquely
determine the inertia coordinates of weights
two through four from \(a_1,a_2\). In the
degree-five basis
\[
 [X,E_1],\ [Y,E_1],\ [Y,E_2],\ [Y,E_3],\ [C,U],\ [C,V],
\]
the obstruction to an integral inertia lift
is the following six equations over
\(\mathbb F_5\):
\[
 \begin{pmatrix}
 a^4b\\3b-a^3b^2\\2b+a^2b^3\\3b-ab^4\\a^3b^2\\2a^2b^3
 \end{pmatrix}
 =
 \begin{pmatrix}
 av_1\\bv_1+av_2\\bv_2+av_3\\bv_3\\av_2\\av_3
 \end{pmatrix}.                                \tag{D.32.10}
\]
For transparency, the change from this basis
to the degree-five Lyndon basis in (D.32.1)
is the unimodular matrix
\[
 \begin{pmatrix}
 1&0&0&0&0&0\\0&-1&0&0&0&0\\0&1&0&0&-1&0\\
 0&0&1&0&0&0\\0&0&1&0&0&-1\\0&0&0&-1&0&0
 \end{pmatrix}.                                \tag{D.32.11}
\]
The program independently obtains (D.32.10)
at both primes and enumerates all its solutions:
\[
 \begin{array}{c|c}
 (a,b)&(v_1,v_2,v_3)\\\hline
 (0,0)&\text{arbitrary}\\
 a\ne0,\ b=0&(0,0,0)\\
 a=0,\ b\ne0&(3,2,3)\\
 a=b\ne0&(1,1,2).
 \end{array}                                             \tag{D.32.12}
\]
There are 137 solutions among 3125 tuples.
The same count holds after discarding the last
two rows for the metabelian quotient: if
\(a,b\ne0\), the first and fourth rows give
\(v_1=a^3b\) and \(v_3=3-a b^3\), and the
middle rows imply \(a=b\); substitution gives
the displayed triple. The cases where either
\(a\) or \(b\) is zero are immediate from
the remaining rows. Thus the pair
\((\kappa(2),2\kappa(2))\), whose inertia
at two is \((1,2)\), cannot lift to class
five for any of its class-four choices.

Finally, equations (D.32.10) at two and three
are only the indicated local simplifications.
They are not substituted for (D.32.6) at other
primes. For example \(v_5(7^4-1)=2\); the
weight-four Frobenius normalization there need
not permit division by \((7^4-1)/5\) inside
\(\mathbb Z_5\). Primes congruent to one or
minus one modulo five also have nonunit lower
normalization factors. Formulas
(D.32.2)--(D.32.6), with the fixed constants
(D.31.2), cover all those primes without a
local normalization or discarded parameter.
They supply both their full congruences and
the divided-family convergence that the
twelve R17 congruences alone cannot detect.
\end{proof}

\subsection*{Expanded integral Hall residue polynomials for D.32}
This appendix is generated by \texttt{de\_hall\_tex.py}.\par
It uses the exact output of \texttt{de\_tame\_hall.py}.
Both scripts are in the \texttt{code/} directory.
The variables and their order are defined in D.32.
For each prime substitute its actual $q,h_q,a_i,b_i$.
These are rational presentations of integer-valued polynomials;
denominators divisible by five must be cancelled before reduction.
No lower-stage relation has been used to shorten these expressions.
\begingroup\small\allowdisplaybreaks[4]
\begin{align*}
R_{4,1}={}& \frac{a_{1}^{3} a_{2} q^{3}}{8} + \frac{a_{1}^{2} a_{3} q^{3}}{2} + \frac{a_{1} a_{3} q^{3}}{2}\\
&{}+ a_{1} a_{4} q^{3} + \frac{a_{3} q^{3}}{6} + a_{4} q^{3}\\
&{}- \frac{a_{1}^{3} a_{2} q^{2}}{12} - \frac{a_{1}^{3} b_{2} q^{2}}{6} + \frac{a_{1}^{2} a_{2} b_{1} q^{2}}{2}\\
&{}- \frac{a_{1}^{2} a_{2} q^{2}}{3} - \frac{a_{1}^{2} a_{3} q^{2}}{6} + a_{1} a_{3} b_{1} q^{2}\\
&{}- a_{1} a_{3} q^{2} + \frac{a_{3} b_{1} q^{2}}{2} - \frac{a_{3} q^{2}}{2}\\
&{}+ a_{4} b_{1} q^{2} - a_{4} q^{2} - \frac{a_{1}^{3} a_{2} q}{24}\\
&{}+ \frac{a_{1}^{2} a_{2} q}{4} - \frac{a_{1}^{2} a_{3} q}{4} - \frac{a_{1}^{2} b_{1} b_{2} q}{2}\\
&{}+ \frac{a_{1}^{2} b_{2} q}{2} - \frac{a_{1}^{2} b_{3} q}{2} + \frac{a_{1} a_{2} b_{1}^{2} q}{2}\\
&{}- a_{1} a_{2} b_{1} q + \frac{a_{1} a_{2} q}{6} + \frac{a_{1} a_{3} q}{4}\\
&{}- \frac{a_{1} a_{4} q}{2} + \frac{a_{3} b_{1}^{2} q}{2} - a_{3} b_{1} q\\
&{}+ \frac{a_{3} q}{3} + \frac{a_{1}^{2} a_{2}}{12} - \frac{a_{1}^{2} a_{3}}{12}\\
&{}- \frac{a_{1} a_{2}}{6} + \frac{a_{1} a_{3}}{4} - \frac{a_{1} a_{4}}{2}\\
&{}- \frac{a_{1} b_{1}^{2} b_{2}}{2} + a_{1} b_{1} b_{2} - a_{1} b_{1} b_{3}\\
&{}- \frac{a_{1} b_{2}}{3} + \frac{a_{1} b_{3}}{2} - a_{1} b_{4}\\
&{}+ \frac{a_{2} b_{1}^{3}}{6} - \frac{a_{2} b_{1}^{2}}{2} + \frac{a_{2} b_{1}}{3}.
\end{align*}
\begin{align*}
R_{4,2}={}&- \frac{a_{1}^{2} a_{2}^{2} q^{3}}{8} - a_{1} a_{2} a_{3} q^{3} - \frac{a_{1} a_{2} q^{3}}{24}\\
&{}- \frac{a_{1} a_{3} q^{3}}{2} + a_{1} a_{5} q^{3} - \frac{a_{2} a_{3} q^{3}}{2}\\
&{}- a_{2} a_{4} q^{3} - \frac{7 a_{3} q^{3}}{24} - \frac{a_{4} q^{3}}{2}\\
&{}+ \frac{a_{5} q^{3}}{2} + \frac{a_{1}^{2} a_{2}^{2} q^{2}}{6} - \frac{a_{1}^{2} a_{2} q^{2}}{6}\\
&{}- \frac{a_{1} a_{2}^{2} b_{1} q^{2}}{2} + \frac{a_{1} a_{2}^{2} q^{2}}{12} + \frac{a_{1} a_{2} a_{3} q^{2}}{3}\\
&{}- a_{1} a_{3} b_{2} q^{2} - \frac{a_{1} a_{3} q^{2}}{2} - a_{2} a_{3} b_{1} q^{2}\\
&{}+ \frac{a_{2} a_{3} q^{2}}{2} - \frac{a_{2} b_{1} q^{2}}{24} + \frac{a_{2} q^{2}}{48}\\
&{}- \frac{a_{3} b_{1} q^{2}}{2} - \frac{a_{3} b_{2} q^{2}}{2} - a_{4} b_{2} q^{2}\\
&{}- \frac{a_{4} q^{2}}{2} + a_{5} b_{1} q^{2} - \frac{a_{5} q^{2}}{2}\\
&{}+ \frac{a_{1}^{2} a_{2} q}{8} + \frac{a_{1}^{2} b_{2}^{2} q}{4} + \frac{a_{1}^{2} b_{2} q}{4}\\
&{}- \frac{a_{1} a_{2}^{2} q}{8} + \frac{a_{1} a_{2} a_{3} q}{2} - \frac{a_{1} a_{2} b_{1} q}{2}\\
&{}+ a_{1} a_{2} b_{3} q + \frac{a_{1} a_{2} q}{6} + \frac{a_{1} a_{3} q}{2}\\
&{}- \frac{a_{1} a_{5} q}{2} - \frac{a_{2}^{2} b_{1}^{2} q}{4} + \frac{a_{2}^{2} b_{1} q}{4}\\
&{}+ \frac{a_{2} a_{4} q}{2} - a_{3} b_{1} b_{2} q - \frac{a_{3} b_{1} q}{2}\\
&{}+ \frac{a_{3} b_{2} q}{2} + \frac{7 a_{3} q}{24} - a_{2} h\\
&{}- \frac{a_{1}^{2} a_{2}^{2}}{24} + \frac{a_{1}^{2} a_{2}}{24} + \frac{a_{1} a_{2}^{2}}{24}\\
&{}+ \frac{a_{1} a_{2} a_{3}}{6} - \frac{a_{1} a_{2}}{8} + \frac{a_{1} a_{3}}{2}\\
&{}- \frac{a_{1} a_{5}}{2} + \frac{a_{1} b_{1} b_{2}^{2}}{2} + \frac{a_{1} b_{1} b_{2}}{2}\\
&{}- \frac{a_{1} b_{2}^{2}}{4} + a_{1} b_{2} b_{3} - \frac{a_{1} b_{2}}{4}\\
&{}+ a_{1} b_{3} - a_{1} b_{5} + \frac{a_{2} a_{4}}{2}\\
&{}- \frac{a_{2} b_{1}^{2}}{4} + a_{2} b_{1} b_{3} + \frac{7 a_{2} b_{1}}{24}\\
&{}+ a_{2} b_{4} - \frac{a_{2}}{48} + a_{4}.
\end{align*}
\begin{align*}
R_{4,3}={}& \frac{a_{1} a_{2}^{3} q^{3}}{24} + \frac{a_{2}^{2} a_{3} q^{3}}{2} + \frac{a_{2}^{2} q^{3}}{48}\\
&{}+ \frac{a_{2} a_{3} q^{3}}{2} - a_{2} a_{5} q^{3} + \frac{a_{3} q^{3}}{6}\\
&{}- a_{5} q^{3} - \frac{a_{1} a_{2}^{3} q^{2}}{12} + \frac{a_{1} a_{2}^{2} q^{2}}{6}\\
&{}+ \frac{a_{2}^{3} b_{1} q^{2}}{6} - \frac{a_{2}^{2} a_{3} q^{2}}{6} + a_{2} a_{3} b_{2} q^{2}\\
&{}+ a_{2} a_{3} q^{2} + \frac{a_{2} b_{2} q^{2}}{24} + \frac{a_{2} q^{2}}{24}\\
&{}+ \frac{a_{3} b_{2} q^{2}}{2} + \frac{a_{3} q^{2}}{2} - a_{5} b_{2} q^{2}\\
&{}- a_{5} q^{2} + \frac{a_{1} a_{2}^{3} q}{24} - \frac{a_{1} a_{2}^{2} q}{4}\\
&{}+ \frac{a_{1} a_{2} q}{6} - \frac{a_{2}^{2} a_{3} q}{4} + \frac{a_{2}^{2} b_{1} q}{2}\\
&{}- \frac{a_{2}^{2} b_{3} q}{2} - \frac{a_{2}^{2} q}{48} - \frac{3 a_{2} a_{3} q}{4}\\
&{}+ \frac{a_{2} a_{5} q}{2} + \frac{a_{3} b_{2}^{2} q}{2} + a_{3} b_{2} q\\
&{}+ \frac{a_{3} q}{3} + a_{2} h + \frac{a_{1} a_{2}^{2}}{12}\\
&{}- \frac{a_{1} a_{2}}{6} - \frac{a_{1} b_{2}^{3}}{6} - \frac{a_{1} b_{2}^{2}}{2}\\
&{}- \frac{a_{1} b_{2}}{3} - \frac{a_{2}^{2} a_{3}}{12} - \frac{3 a_{2} a_{3}}{4}\\
&{}+ \frac{a_{2} a_{5}}{2} + \frac{a_{2} b_{1}}{3} - a_{2} b_{2} b_{3}\\
&{}- \frac{a_{2} b_{2}}{24} - \frac{3 a_{2} b_{3}}{2} + a_{2} b_{5}\\
&{}- \frac{a_{2}}{24} - a_{3} + 2 a_{5}.
\end{align*}
\begin{align*}
R_{5,1}={}& \frac{a_{1}^{4} a_{2} q^{4}}{30} + \frac{a_{1}^{3} a_{3} q^{4}}{6} + \frac{a_{1}^{2} a_{3} q^{4}}{4}\\
&{}+ \frac{a_{1}^{2} a_{4} q^{4}}{2} + \frac{a_{1} a_{3} q^{4}}{6} + a_{1} a_{4} q^{4}\\
&{}+ a_{1} a_{6} q^{4} + \frac{a_{3} q^{4}}{24} + \frac{7 a_{4} q^{4}}{12}\\
&{}+ \frac{3 a_{6} q^{4}}{2} - \frac{a_{1}^{4} a_{2} q^{3}}{48} - \frac{a_{1}^{4} b_{2} q^{3}}{24}\\
&{}+ \frac{a_{1}^{3} a_{2} b_{1} q^{3}}{6} - \frac{3 a_{1}^{3} a_{2} q^{3}}{16} - \frac{a_{1}^{3} a_{3} q^{3}}{24}\\
&{}+ \frac{a_{1}^{2} a_{3} b_{1} q^{3}}{2} - \frac{3 a_{1}^{2} a_{3} q^{3}}{4} + \frac{a_{1} a_{3} b_{1} q^{3}}{2}\\
&{}- \frac{3 a_{1} a_{3} q^{3}}{4} + a_{1} a_{4} b_{1} q^{3} - \frac{3 a_{1} a_{4} q^{3}}{2}\\
&{}+ \frac{a_{3} b_{1} q^{3}}{6} - \frac{a_{3} q^{3}}{4} + a_{4} b_{1} q^{3}\\
&{}- \frac{3 a_{4} q^{3}}{2} + a_{6} b_{1} q^{3} - \frac{3 a_{6} q^{3}}{2}\\
&{}- \frac{a_{1}^{4} a_{2} q^{2}}{72} + \frac{a_{1}^{3} a_{2} q^{2}}{8} - \frac{a_{1}^{3} a_{3} q^{2}}{12}\\
&{}- \frac{a_{1}^{3} b_{1} b_{2} q^{2}}{6} + \frac{a_{1}^{3} b_{2} q^{2}}{4} - \frac{a_{1}^{3} b_{3} q^{2}}{6}\\
&{}+ \frac{a_{1}^{2} a_{2} b_{1}^{2} q^{2}}{4} - \frac{3 a_{1}^{2} a_{2} b_{1} q^{2}}{4} + \frac{11 a_{1}^{2} a_{2} q^{2}}{36}\\
&{}+ \frac{a_{1}^{2} a_{3} q^{2}}{6} - \frac{a_{1}^{2} a_{4} q^{2}}{6} + \frac{a_{1} a_{3} b_{1}^{2} q^{2}}{2}\\
&{}- \frac{3 a_{1} a_{3} b_{1} q^{2}}{2} + \frac{11 a_{1} a_{3} q^{2}}{12} + \frac{a_{3} b_{1}^{2} q^{2}}{4}\\
&{}- \frac{3 a_{3} b_{1} q^{2}}{4} + \frac{11 a_{3} q^{2}}{24} + \frac{a_{4} b_{1}^{2} q^{2}}{2}\\
&{}- \frac{3 a_{4} b_{1} q^{2}}{2} + \frac{11 a_{4} q^{2}}{12} + \frac{a_{1}^{3} a_{2} q}{16}\\
&{}- \frac{a_{1}^{3} a_{3} q}{24} - \frac{11 a_{1}^{2} a_{2} q}{48} + \frac{a_{1}^{2} a_{3} q}{4}\\
&{}- \frac{a_{1}^{2} a_{4} q}{4} - \frac{a_{1}^{2} b_{1}^{2} b_{2} q}{4} + \frac{3 a_{1}^{2} b_{1} b_{2} q}{4}\\
&{}- \frac{a_{1}^{2} b_{1} b_{3} q}{2} - \frac{11 a_{1}^{2} b_{2} q}{24} + \frac{a_{1}^{2} b_{3} q}{2}\\
&{}- \frac{a_{1}^{2} b_{4} q}{2} + \frac{a_{1} a_{2} b_{1}^{3} q}{6} - \frac{3 a_{1} a_{2} b_{1}^{2} q}{4}\\
&{}+ \frac{11 a_{1} a_{2} b_{1} q}{12} - \frac{a_{1} a_{2} q}{8} - \frac{a_{1} a_{3} q}{6}\\
&{}+ \frac{a_{1} a_{4} q}{4} - \frac{a_{1} a_{6} q}{2} + \frac{a_{3} b_{1}^{3} q}{6}\\
&{}- \frac{3 a_{3} b_{1}^{2} q}{4} + \frac{11 a_{3} b_{1} q}{12} - \frac{a_{3} q}{4}\\
&{}+ \frac{a_{1}^{4} a_{2}}{720} - \frac{11 a_{1}^{2} a_{2}}{144} + \frac{a_{1}^{2} a_{3}}{12}\\
&{}- \frac{a_{1}^{2} a_{4}}{12} + \frac{a_{1} a_{2}}{8} - \frac{a_{1} a_{3}}{6}\\
&{}+ \frac{a_{1} a_{4}}{4} - \frac{a_{1} a_{6}}{2} - \frac{a_{1} b_{1}^{3} b_{2}}{6}\\
&{}+ \frac{3 a_{1} b_{1}^{2} b_{2}}{4} - \frac{a_{1} b_{1}^{2} b_{3}}{2} - \frac{11 a_{1} b_{1} b_{2}}{12}\\
&{}+ a_{1} b_{1} b_{3} - a_{1} b_{1} b_{4} + \frac{a_{1} b_{2}}{4}\\
&{}- \frac{a_{1} b_{3}}{3} + \frac{a_{1} b_{4}}{2} - a_{1} b_{6}\\
&{}+ \frac{a_{2} b_{1}^{4}}{24} - \frac{a_{2} b_{1}^{3}}{4} + \frac{11 a_{2} b_{1}^{2}}{24}\\
&{}- \frac{a_{2} b_{1}}{4}.
\end{align*}
\begin{align*}
R_{5,2}={}&- \frac{a_{1}^{3} a_{2}^{2} q^{4}}{20} - \frac{a_{1}^{2} a_{2} a_{3} q^{4}}{2} - \frac{a_{1}^{2} a_{2} q^{4}}{48}\\
&{}- \frac{a_{1}^{2} a_{3} q^{4}}{4} + \frac{a_{1}^{2} a_{5} q^{4}}{2} - \frac{a_{1} a_{2} a_{3} q^{4}}{2}\\
&{}- a_{1} a_{2} a_{4} q^{4} - \frac{7 a_{1} a_{3} q^{4}}{24} - \frac{a_{1} a_{4} q^{4}}{2}\\
&{}+ \frac{a_{1} a_{5} q^{4}}{2} + a_{1} a_{7} q^{4} - \frac{a_{2} a_{3} q^{4}}{6}\\
&{}- a_{2} a_{4} q^{4} - a_{2} a_{6} q^{4} + \frac{a_{2} q^{4}}{1440}\\
&{}- \frac{5 a_{3} q^{4}}{48} - \frac{13 a_{4} q^{4}}{24} + \frac{a_{5} q^{4}}{6}\\
&{}- \frac{a_{6} q^{4}}{2} + a_{7} q^{4} + \frac{a_{1}^{3} a_{2}^{2} q^{3}}{16}\\
&{}- \frac{a_{1}^{3} a_{2} q^{3}}{16} - \frac{a_{1}^{2} a_{2}^{2} b_{1} q^{3}}{4} + \frac{a_{1}^{2} a_{2}^{2} q^{3}}{8}\\
&{}+ \frac{a_{1}^{2} a_{2} a_{3} q^{3}}{8} - \frac{a_{1}^{2} a_{3} b_{2} q^{3}}{2} - \frac{a_{1}^{2} a_{3} q^{3}}{4}\\
&{}- a_{1} a_{2} a_{3} b_{1} q^{3} + a_{1} a_{2} a_{3} q^{3} - \frac{a_{1} a_{2} b_{1} q^{3}}{24}\\
&{}+ \frac{a_{1} a_{2} q^{3}}{24} - \frac{a_{1} a_{3} b_{1} q^{3}}{2} - \frac{a_{1} a_{3} b_{2} q^{3}}{2}\\
&{}+ \frac{a_{1} a_{3} q^{3}}{4} - a_{1} a_{4} b_{2} q^{3} - \frac{a_{1} a_{4} q^{3}}{2}\\
&{}+ a_{1} a_{5} b_{1} q^{3} - a_{1} a_{5} q^{3} - \frac{a_{2} a_{3} b_{1} q^{3}}{2}\\
&{}+ \frac{a_{2} a_{3} q^{3}}{2} - a_{2} a_{4} b_{1} q^{3} + a_{2} a_{4} q^{3}\\
&{}- \frac{7 a_{3} b_{1} q^{3}}{24} - \frac{a_{3} b_{2} q^{3}}{6} + \frac{5 a_{3} q^{3}}{24}\\
&{}- \frac{a_{4} b_{1} q^{3}}{2} - a_{4} b_{2} q^{3} + \frac{a_{5} b_{1} q^{3}}{2}\\
&{}- \frac{a_{5} q^{3}}{2} - a_{6} b_{2} q^{3} - \frac{a_{6} q^{3}}{2}\\
&{}+ a_{7} b_{1} q^{3} - a_{7} q^{3} + \frac{a_{1}^{3} a_{2}^{2} q^{2}}{72}\\
&{}+ \frac{a_{1}^{3} a_{2} q^{2}}{24} + \frac{a_{1}^{3} b_{2}^{2} q^{2}}{12} + \frac{a_{1}^{3} b_{2} q^{2}}{12}\\
&{}- \frac{a_{1}^{2} a_{2}^{2} q^{2}}{6} + \frac{a_{1}^{2} a_{2} a_{3} q^{2}}{4} - \frac{a_{1}^{2} a_{2} b_{1} q^{2}}{4}\\
&{}+ \frac{a_{1}^{2} a_{2} b_{3} q^{2}}{2} + \frac{3 a_{1}^{2} a_{2} q^{2}}{16} + \frac{a_{1}^{2} a_{3} q^{2}}{6}\\
&{}- \frac{a_{1}^{2} a_{5} q^{2}}{6} - \frac{a_{1} a_{2}^{2} b_{1}^{2} q^{2}}{4} + \frac{a_{1} a_{2}^{2} b_{1} q^{2}}{2}\\
&{}- \frac{a_{1} a_{2}^{2} q^{2}}{18} - \frac{a_{1} a_{2} a_{3} q^{2}}{6} + \frac{a_{1} a_{2} a_{4} q^{2}}{3}\\
&{}- a_{1} a_{3} b_{1} b_{2} q^{2} - \frac{a_{1} a_{3} b_{1} q^{2}}{2} + a_{1} a_{3} b_{2} q^{2}\\
&{}+ \frac{13 a_{1} a_{3} q^{2}}{24} - \frac{a_{2} a_{3} b_{1}^{2} q^{2}}{2} + a_{2} a_{3} b_{1} q^{2}\\
&{}- \frac{a_{2} a_{3} q^{2}}{3} - \frac{a_{2} b_{1}^{2} q^{2}}{48} + \frac{a_{2} b_{1} q^{2}}{24}\\
&{}- \frac{a_{2} q^{2}}{72} - \frac{a_{3} b_{1}^{2} q^{2}}{4} - \frac{a_{3} b_{1} b_{2} q^{2}}{2}\\
&{}+ \frac{a_{3} b_{1} q^{2}}{4} + \frac{a_{3} b_{2} q^{2}}{2} + \frac{5 a_{3} q^{2}}{48}\\
&{}- a_{4} b_{1} b_{2} q^{2} - \frac{a_{4} b_{1} q^{2}}{2} + a_{4} b_{2} q^{2}\\
&{}+ \frac{13 a_{4} q^{2}}{24} + \frac{a_{5} b_{1}^{2} q^{2}}{2} - a_{5} b_{1} q^{2}\\
&{}+ \frac{a_{5} q^{2}}{3} - a_{1} a_{2} h q - a_{3} h q\\
&{}- \frac{a_{1}^{3} a_{2}^{2} q}{48} + \frac{a_{1}^{3} a_{2} q}{48} + \frac{a_{1}^{2} a_{2} a_{3} q}{8}\\
&{}- \frac{a_{1}^{2} a_{2} q}{8} + \frac{a_{1}^{2} a_{3} q}{4} - \frac{a_{1}^{2} a_{5} q}{4}\\
&{}+ \frac{a_{1}^{2} b_{1} b_{2}^{2} q}{4} + \frac{a_{1}^{2} b_{1} b_{2} q}{4} - \frac{a_{1}^{2} b_{2}^{2} q}{4}\\
&{}+ \frac{a_{1}^{2} b_{2} b_{3} q}{2} - \frac{a_{1}^{2} b_{2} q}{4} + \frac{a_{1}^{2} b_{3} q}{2}\\
&{}- \frac{a_{1}^{2} b_{5} q}{2} + \frac{a_{1} a_{2}^{2} q}{12} - \frac{a_{1} a_{2} a_{3} q}{4}\\
&{}+ \frac{a_{1} a_{2} a_{4} q}{2} - \frac{a_{1} a_{2} b_{1}^{2} q}{4} + a_{1} a_{2} b_{1} b_{3} q\\
&{}+ \frac{13 a_{1} a_{2} b_{1} q}{24} - \frac{a_{1} a_{2} b_{3} q}{2} + a_{1} a_{2} b_{4} q\\
&{}- \frac{a_{1} a_{2} q}{8} - \frac{a_{1} a_{3} q}{4} + \frac{a_{1} a_{4} q}{2}\\
&{}+ \frac{a_{1} a_{5} q}{4} - \frac{a_{1} a_{7} q}{2} - \frac{a_{2}^{2} b_{1}^{3} q}{12}\\
&{}+ \frac{a_{2}^{2} b_{1}^{2} q}{4} - \frac{a_{2}^{2} b_{1} q}{6} + \frac{a_{2} a_{6} q}{2}\\
&{}- \frac{a_{3} b_{1}^{2} b_{2} q}{2} - \frac{a_{3} b_{1}^{2} q}{4} + a_{3} b_{1} b_{2} q\\
&{}+ \frac{13 a_{3} b_{1} q}{24} - \frac{a_{3} b_{2} q}{3} - \frac{5 a_{3} q}{24}\\
&{}- a_{2} b_{1} h + a_{2} h - \frac{a_{1}^{3} a_{2}^{2}}{180}\\
&{}+ \frac{a_{1}^{2} a_{2}^{2}}{24} - \frac{a_{1}^{2} a_{2}}{24} + \frac{a_{1}^{2} a_{3}}{12}\\
&{}- \frac{a_{1}^{2} a_{5}}{12} - \frac{a_{1} a_{2}^{2}}{36} - \frac{a_{1} a_{2} a_{3}}{12}\\
&{}+ \frac{a_{1} a_{2} a_{4}}{6} + \frac{a_{1} a_{2}}{12} - \frac{a_{1} a_{3}}{4}\\
&{}+ \frac{a_{1} a_{4}}{2} + \frac{a_{1} a_{5}}{4} - \frac{a_{1} a_{7}}{2}\\
&{}+ \frac{a_{1} b_{1}^{2} b_{2}^{2}}{4} + \frac{a_{1} b_{1}^{2} b_{2}}{4} - \frac{a_{1} b_{1} b_{2}^{2}}{2}\\
&{}+ a_{1} b_{1} b_{2} b_{3} - \frac{a_{1} b_{1} b_{2}}{2} + a_{1} b_{1} b_{3}\\
&{}- a_{1} b_{1} b_{5} + \frac{a_{1} b_{2}^{2}}{6} - \frac{a_{1} b_{2} b_{3}}{2}\\
&{}+ a_{1} b_{2} b_{4} + \frac{a_{1} b_{2}}{6} - \frac{a_{1} b_{3}}{2}\\
&{}+ a_{1} b_{4} + \frac{a_{1} b_{5}}{2} - a_{1} b_{7}\\
&{}+ \frac{a_{2} a_{6}}{2} - \frac{a_{2} b_{1}^{3}}{12} + \frac{a_{2} b_{1}^{2} b_{3}}{2}\\
&{}+ \frac{13 a_{2} b_{1}^{2}}{48} - \frac{a_{2} b_{1} b_{3}}{2} + a_{2} b_{1} b_{4}\\
&{}- \frac{5 a_{2} b_{1}}{24} + a_{2} b_{6} + \frac{19 a_{2}}{1440}\\
&{}+ a_{6}.
\end{align*}
\begin{align*}
R_{5,3}={}& \frac{a_{1}^{3} a_{2}^{2} q^{4}}{10} + \frac{a_{1}^{2} a_{2} a_{3} q^{4}}{2} + \frac{a_{1} a_{2} a_{3} q^{4}}{4}\\
&{}+ \frac{a_{1} a_{2} a_{4} q^{4}}{2} + \frac{a_{1} a_{3}^{2} q^{4}}{2} + \frac{a_{2} a_{3} q^{4}}{6}\\
&{}+ a_{2} a_{4} q^{4} + a_{2} a_{6} q^{4} - \frac{a_{2} q^{4}}{1440}\\
&{}+ \frac{a_{3}^{2} q^{4}}{4} + \frac{a_{3} q^{4}}{12} - \frac{a_{4} q^{4}}{2}\\
&{}- \frac{a_{5} q^{4}}{6} - \frac{3 a_{1}^{3} a_{2}^{2} q^{3}}{16} - \frac{a_{1}^{3} a_{2} b_{2} q^{3}}{4}\\
&{}- \frac{a_{1}^{3} a_{2} q^{3}}{8} + \frac{a_{1}^{2} a_{2}^{2} b_{1} q^{3}}{2} + \frac{a_{1}^{2} a_{2}^{2} q^{3}}{16}\\
&{}- \frac{7 a_{1}^{2} a_{2} a_{3} q^{3}}{8} - \frac{a_{1}^{2} a_{3} b_{2} q^{3}}{2} - \frac{a_{1}^{2} a_{3} q^{3}}{2}\\
&{}+ \frac{3 a_{1} a_{2} a_{3} b_{1} q^{3}}{2} + \frac{a_{1} a_{2} a_{3} q^{3}}{2} - \frac{a_{1} a_{2} a_{4} q^{3}}{2}\\
&{}+ \frac{a_{1} a_{2} q^{3}}{24} - a_{1} a_{3}^{2} q^{3} - a_{1} a_{4} q^{3}\\
&{}- a_{1} a_{5} q^{3} + \frac{a_{2} a_{3} b_{1} q^{3}}{2} + \frac{a_{2} a_{3} q^{3}}{2}\\
&{}+ a_{2} a_{4} b_{1} q^{3} + a_{2} a_{4} q^{3} + \frac{a_{3}^{2} b_{1} q^{3}}{2}\\
&{}- \frac{3 a_{3}^{2} q^{3}}{4} - a_{3} a_{4} q^{3} + \frac{a_{3} b_{2} q^{3}}{6}\\
&{}+ \frac{a_{3} q^{3}}{8} + a_{4} b_{2} q^{3} - \frac{a_{4} q^{3}}{2}\\
&{}- \frac{a_{5} q^{3}}{2} + a_{6} b_{2} q^{3} - a_{6} q^{3}\\
&{}- a_{7} q^{3} + \frac{5 a_{1}^{3} a_{2}^{2} q^{2}}{72} + \frac{a_{1}^{3} a_{2} b_{2} q^{2}}{4}\\
&{}+ \frac{a_{1}^{3} a_{2} q^{2}}{12} + \frac{a_{1}^{3} b_{2}^{2} q^{2}}{6} + \frac{a_{1}^{3} b_{2} q^{2}}{6}\\
&{}- \frac{a_{1}^{2} a_{2}^{2} b_{1} q^{2}}{2} - \frac{a_{1}^{2} a_{2}^{2} q^{2}}{24} + \frac{a_{1}^{2} a_{2} a_{3} q^{2}}{4}\\
&{}- a_{1}^{2} a_{2} b_{1} b_{2} q^{2} - \frac{a_{1}^{2} a_{2} b_{1} q^{2}}{2} + \frac{a_{1}^{2} a_{2} b_{2} q^{2}}{4}\\
&{}- \frac{a_{1}^{2} a_{2} b_{3} q^{2}}{2} + \frac{a_{1}^{2} a_{2} q^{2}}{3} + \frac{a_{1}^{2} a_{3} b_{2} q^{2}}{2}\\
&{}+ \frac{a_{1}^{2} a_{3} q^{2}}{6} + \frac{3 a_{1} a_{2}^{2} b_{1}^{2} q^{2}}{4} + \frac{a_{1} a_{2}^{2} b_{1} q^{2}}{4}\\
&{}- \frac{a_{1} a_{2}^{2} q^{2}}{9} - \frac{3 a_{1} a_{2} a_{3} b_{1} q^{2}}{2} + \frac{a_{1} a_{2} a_{3} q^{2}}{12}\\
&{}- \frac{a_{1} a_{2} a_{4} q^{2}}{6} + \frac{a_{1} a_{3}^{2} q^{2}}{3} - a_{1} a_{3} b_{1} b_{2} q^{2}\\
&{}- a_{1} a_{3} b_{1} q^{2} + \frac{3 a_{1} a_{3} b_{2} q^{2}}{2} - a_{1} a_{3} b_{3} q^{2}\\
&{}+ a_{1} a_{3} q^{2} + a_{2} a_{3} b_{1}^{2} q^{2} + \frac{a_{2} a_{3} b_{1} q^{2}}{2}\\
&{}- \frac{2 a_{2} a_{3} q^{2}}{3} + \frac{a_{2} b_{1} q^{2}}{24} - \frac{a_{2} q^{2}}{36}\\
&{}- a_{3}^{2} b_{1} q^{2} + \frac{a_{3}^{2} q^{2}}{2} + \frac{a_{3} b_{2} q^{2}}{2}\\
&{}- \frac{a_{3} b_{3} q^{2}}{2} + \frac{a_{3} q^{2}}{6} - a_{4} b_{1} q^{2}\\
&{}+ a_{4} b_{2} q^{2} - a_{4} b_{3} q^{2} + a_{4} q^{2}\\
&{}- a_{5} b_{1} q^{2} + \frac{2 a_{5} q^{2}}{3} + \frac{a_{1}^{3} a_{2}^{2} q}{48}\\
&{}+ \frac{a_{1}^{3} a_{2} q}{24} - \frac{a_{1}^{2} a_{2}^{2} q}{16} + \frac{a_{1}^{2} a_{2} a_{3} q}{8}\\
&{}+ \frac{a_{1}^{2} a_{2} b_{1} b_{2} q}{2} - \frac{a_{1}^{2} a_{2} b_{2} q}{4} + \frac{a_{1}^{2} a_{2} b_{3} q}{2}\\
&{}- \frac{a_{1}^{2} a_{2} q}{4} + \frac{a_{1}^{2} a_{3} q}{4} + \frac{a_{1}^{2} b_{1} b_{2}^{2} q}{2}\\
&{}+ \frac{a_{1}^{2} b_{1} b_{2} q}{2} - \frac{a_{1}^{2} b_{2}^{2} q}{2} + \frac{a_{1}^{2} b_{2} b_{3} q}{2}\\
&{}- \frac{a_{1}^{2} b_{2} q}{2} + \frac{a_{1}^{2} b_{3} q}{2} - \frac{a_{1} a_{2}^{2} b_{1}^{2} q}{4}\\
&{}+ \frac{a_{1} a_{2}^{2} b_{1} q}{4} + \frac{a_{1} a_{2}^{2} q}{6} - \frac{3 a_{1} a_{2} a_{3} q}{4}\\
&{}- a_{1} a_{2} b_{1}^{2} b_{2} q - \frac{a_{1} a_{2} b_{1}^{2} q}{2} + \frac{a_{1} a_{2} b_{1} b_{2} q}{2}\\
&{}- a_{1} a_{2} b_{1} b_{3} q + a_{1} a_{2} b_{1} q - a_{1} a_{2} b_{3} q\\
&{}- \frac{5 a_{1} a_{2} q}{24} + \frac{a_{1} a_{3}^{2} q}{4} + a_{1} a_{3} b_{1} b_{2} q\\
&{}- \frac{a_{1} a_{3} b_{2} q}{2} + a_{1} a_{3} b_{3} q - \frac{a_{1} a_{3} q}{2}\\
&{}+ \frac{a_{1} a_{4} q}{2} + \frac{a_{1} a_{5} q}{2} + \frac{a_{2}^{2} b_{1}^{3} q}{3}\\
&{}- \frac{a_{2}^{2} b_{1} q}{3} - \frac{a_{2} a_{3} b_{1}^{2} q}{2} + \frac{a_{2} a_{3} b_{1} q}{2}\\
&{}- a_{2} a_{4} q - \frac{a_{2} a_{6} q}{2} + \frac{a_{3} a_{4} q}{2}\\
&{}- \frac{a_{3} b_{1}^{2} q}{2} + a_{3} b_{1} b_{2} q + a_{3} b_{1} q\\
&{}- \frac{2 a_{3} b_{2} q}{3} + \frac{a_{3} b_{3} q}{2} + a_{3} b_{4} q\\
&{}- \frac{3 a_{3} q}{8} + a_{2} h - \frac{a_{1}^{3} a_{2}^{2}}{360}\\
&{}+ \frac{a_{1}^{2} a_{2}^{2}}{24} - \frac{a_{1}^{2} a_{2}}{12} + \frac{a_{1}^{2} a_{3}}{12}\\
&{}- \frac{a_{1} a_{2}^{2}}{18} - \frac{a_{1} a_{2} a_{3}}{12} + \frac{a_{1} a_{2} a_{4}}{6}\\
&{}+ \frac{a_{1} a_{2}}{6} - \frac{a_{1} a_{3}^{2}}{12} - \frac{a_{1} a_{3}}{2}\\
&{}+ \frac{a_{1} a_{4}}{2} + \frac{a_{1} a_{5}}{2} + \frac{a_{1} b_{1}^{2} b_{2}}{2}\\
&{}- \frac{a_{1} b_{1} b_{2}^{2}}{2} - a_{1} b_{1} b_{2} + a_{1} b_{1} b_{3}\\
&{}+ \frac{a_{1} b_{2}^{2}}{3} - \frac{3 a_{1} b_{2} b_{3}}{2} - a_{1} b_{2} b_{4}\\
&{}+ \frac{a_{1} b_{2}}{3} + \frac{a_{1} b_{3}^{2}}{2} - a_{1} b_{3}\\
&{}+ a_{1} b_{4} + a_{1} b_{5} - a_{2} a_{4}\\
&{}- \frac{a_{2} a_{6}}{2} - \frac{a_{2} b_{1}^{3}}{6} + \frac{a_{2} b_{1}^{2}}{2}\\
&{}- a_{2} b_{1} b_{3} - \frac{3 a_{2} b_{1}}{8} - 2 a_{2} b_{4}\\
&{}- a_{2} b_{6} + \frac{41 a_{2}}{1440} + \frac{a_{3} a_{4}}{2}\\
&{}+ a_{6} + a_{7}.
\end{align*}
\begin{align*}
R_{5,4}={}& \frac{a_{1}^{2} a_{2}^{3} q^{4}}{30} + \frac{a_{1} a_{2}^{2} a_{3} q^{4}}{2} + \frac{a_{1} a_{2}^{2} q^{4}}{48}\\
&{}+ \frac{a_{1} a_{2} a_{3} q^{4}}{2} - a_{1} a_{2} a_{5} q^{4} + \frac{a_{1} a_{3} q^{4}}{6}\\
&{}- a_{1} a_{5} q^{4} + a_{1} a_{8} q^{4} + \frac{a_{2}^{2} a_{3} q^{4}}{4}\\
&{}+ \frac{a_{2}^{2} a_{4} q^{4}}{2} + \frac{7 a_{2} a_{3} q^{4}}{24} + \frac{a_{2} a_{4} q^{4}}{2}\\
&{}- \frac{a_{2} a_{5} q^{4}}{2} - a_{2} a_{7} q^{4} - \frac{a_{2} q^{4}}{5760}\\
&{}+ \frac{5 a_{3} q^{4}}{48} + \frac{a_{4} q^{4}}{6} - \frac{13 a_{5} q^{4}}{24}\\
&{}- a_{7} q^{4} + \frac{a_{8} q^{4}}{2} - \frac{a_{1}^{2} a_{2}^{3} q^{3}}{16}\\
&{}+ \frac{a_{1}^{2} a_{2}^{2} q^{3}}{8} + \frac{a_{1} a_{2}^{3} b_{1} q^{3}}{6} - \frac{a_{1} a_{2}^{3} q^{3}}{48}\\
&{}- \frac{a_{1} a_{2}^{2} a_{3} q^{3}}{8} + a_{1} a_{2} a_{3} b_{2} q^{3} + a_{1} a_{2} a_{3} q^{3}\\
&{}+ \frac{a_{1} a_{2} b_{2} q^{3}}{24} + \frac{a_{1} a_{2} q^{3}}{24} + \frac{a_{1} a_{3} b_{2} q^{3}}{2}\\
&{}+ \frac{a_{1} a_{3} q^{3}}{2} - a_{1} a_{5} b_{2} q^{3} - a_{1} a_{5} q^{3}\\
&{}+ \frac{a_{2}^{2} a_{3} b_{1} q^{3}}{2} - \frac{a_{2}^{2} a_{3} q^{3}}{4} + \frac{a_{2}^{2} b_{1} q^{3}}{48}\\
&{}- \frac{a_{2}^{2} q^{3}}{96} + \frac{a_{2} a_{3} b_{1} q^{3}}{2} + \frac{a_{2} a_{3} b_{2} q^{3}}{2}\\
&{}+ \frac{a_{2} a_{3} q^{3}}{4} + a_{2} a_{4} b_{2} q^{3} + a_{2} a_{4} q^{3}\\
&{}- a_{2} a_{5} b_{1} q^{3} + \frac{a_{2} a_{5} q^{3}}{2} + \frac{a_{3} b_{1} q^{3}}{6}\\
&{}+ \frac{7 a_{3} b_{2} q^{3}}{24} + \frac{5 a_{3} q^{3}}{24} + \frac{a_{4} b_{2} q^{3}}{2}\\
&{}+ \frac{a_{4} q^{3}}{2} - a_{5} b_{1} q^{3} - \frac{a_{5} b_{2} q^{3}}{2}\\
&{}- a_{7} b_{2} q^{3} - a_{7} q^{3} + a_{8} b_{1} q^{3}\\
&{}- \frac{a_{8} q^{3}}{2} + \frac{a_{1}^{2} a_{2}^{3} q^{2}}{72} - \frac{a_{1}^{2} a_{2}^{2} q^{2}}{6}\\
&{}+ \frac{a_{1}^{2} a_{2} q^{2}}{9} + \frac{a_{1} a_{2}^{3} q^{2}}{24} - \frac{a_{1} a_{2}^{2} a_{3} q^{2}}{4}\\
&{}+ \frac{a_{1} a_{2}^{2} b_{1} q^{2}}{2} - \frac{a_{1} a_{2}^{2} b_{3} q^{2}}{2} - \frac{5 a_{1} a_{2}^{2} q^{2}}{48}\\
&{}- \frac{a_{1} a_{2} a_{3} q^{2}}{2} + \frac{a_{1} a_{2} a_{5} q^{2}}{3} + \frac{a_{1} a_{3} b_{2}^{2} q^{2}}{2}\\
&{}+ a_{1} a_{3} b_{2} q^{2} + \frac{a_{1} a_{3} q^{2}}{3} + \frac{a_{2}^{3} b_{1}^{2} q^{2}}{12}\\
&{}- \frac{a_{2}^{3} b_{1} q^{2}}{12} - \frac{a_{2}^{2} a_{4} q^{2}}{6} + a_{2} a_{3} b_{1} b_{2} q^{2}\\
&{}+ a_{2} a_{3} b_{1} q^{2} - \frac{a_{2} a_{3} b_{2} q^{2}}{2} - \frac{13 a_{2} a_{3} q^{2}}{24}\\
&{}+ \frac{a_{2} b_{1} b_{2} q^{2}}{24} + \frac{a_{2} b_{1} q^{2}}{24} - \frac{a_{2} b_{2} q^{2}}{48}\\
&{}- \frac{13 a_{2} q^{2}}{576} + \frac{a_{3} b_{1} b_{2} q^{2}}{2} + \frac{a_{3} b_{1} q^{2}}{2}\\
&{}+ \frac{a_{3} b_{2}^{2} q^{2}}{4} + \frac{a_{3} b_{2} q^{2}}{4} - \frac{5 a_{3} q^{2}}{48}\\
&{}+ \frac{a_{4} b_{2}^{2} q^{2}}{2} + a_{4} b_{2} q^{2} + \frac{a_{4} q^{2}}{3}\\
&{}- a_{5} b_{1} b_{2} q^{2} - a_{5} b_{1} q^{2} + \frac{a_{5} b_{2} q^{2}}{2}\\
&{}+ \frac{13 a_{5} q^{2}}{24} + a_{1} a_{2} h q + \frac{a_{2}^{2} h q}{2}\\
&{}+ a_{3} h q + \frac{a_{1}^{2} a_{2}^{3} q}{48} - \frac{a_{1}^{2} a_{2} q}{12}\\
&{}- \frac{a_{1}^{2} b_{2}^{3} q}{12} - \frac{a_{1}^{2} b_{2}^{2} q}{4} - \frac{a_{1}^{2} b_{2} q}{6}\\
&{}- \frac{a_{1} a_{2}^{3} q}{48} - \frac{a_{1} a_{2}^{2} a_{3} q}{8} + \frac{a_{1} a_{2}^{2} q}{8}\\
&{}- \frac{3 a_{1} a_{2} a_{3} q}{4} + \frac{a_{1} a_{2} a_{5} q}{2} + \frac{a_{1} a_{2} b_{1} q}{3}\\
&{}- a_{1} a_{2} b_{2} b_{3} q - \frac{a_{1} a_{2} b_{2} q}{24} - \frac{3 a_{1} a_{2} b_{3} q}{2}\\
&{}+ a_{1} a_{2} b_{5} q - \frac{a_{1} a_{2} q}{8} - \frac{a_{1} a_{3} q}{2}\\
&{}+ a_{1} a_{5} q - \frac{a_{1} a_{8} q}{2} - \frac{a_{2}^{2} a_{4} q}{4}\\
&{}+ \frac{a_{2}^{2} b_{1}^{2} q}{4} - \frac{a_{2}^{2} b_{1} b_{3} q}{2} - \frac{13 a_{2}^{2} b_{1} q}{48}\\
&{}- \frac{a_{2}^{2} b_{4} q}{2} + \frac{a_{2}^{2} q}{96} - \frac{3 a_{2} a_{4} q}{4}\\
&{}+ \frac{a_{2} a_{7} q}{2} + \frac{a_{3} b_{1} b_{2}^{2} q}{2} + a_{3} b_{1} b_{2} q\\
&{}+ \frac{a_{3} b_{1} q}{3} - \frac{a_{3} b_{2}^{2} q}{4} - \frac{13 a_{3} b_{2} q}{24}\\
&{}- \frac{5 a_{3} q}{24} + a_{2} b_{1} h + a_{2} b_{2} h\\
&{}+ \frac{a_{2} h}{2} - \frac{a_{1}^{2} a_{2}^{3}}{180} + \frac{a_{1}^{2} a_{2}^{2}}{24}\\
&{}- \frac{a_{1}^{2} a_{2}}{36} - \frac{a_{1} a_{2}^{2}}{24} - \frac{a_{1} a_{2} a_{3}}{4}\\
&{}+ \frac{a_{1} a_{2} a_{5}}{6} + \frac{a_{1} a_{2}}{12} - \frac{a_{1} a_{3}}{2}\\
&{}+ a_{1} a_{5} - \frac{a_{1} a_{8}}{2} - \frac{a_{1} b_{1} b_{2}^{3}}{6}\\
&{}- \frac{a_{1} b_{1} b_{2}^{2}}{2} - \frac{a_{1} b_{1} b_{2}}{3} + \frac{a_{1} b_{2}^{3}}{12}\\
&{}- \frac{a_{1} b_{2}^{2} b_{3}}{2} + \frac{a_{1} b_{2}^{2}}{4} - \frac{3 a_{1} b_{2} b_{3}}{2}\\
&{}+ a_{1} b_{2} b_{5} + \frac{a_{1} b_{2}}{6} - a_{1} b_{3}\\
&{}+ 2 a_{1} b_{5} - a_{1} b_{8} - \frac{a_{2}^{2} a_{4}}{12}\\
&{}- \frac{3 a_{2} a_{4}}{4} + \frac{a_{2} a_{7}}{2} + \frac{a_{2} b_{1}^{2}}{6}\\
&{}- a_{2} b_{1} b_{2} b_{3} - \frac{a_{2} b_{1} b_{2}}{24} - \frac{3 a_{2} b_{1} b_{3}}{2}\\
&{}+ a_{2} b_{1} b_{5} - \frac{5 a_{2} b_{1}}{24} - a_{2} b_{2} b_{4}\\
&{}+ \frac{a_{2} b_{2}}{48} - \frac{3 a_{2} b_{4}}{2} + a_{2} b_{7}\\
&{}+ \frac{131 a_{2}}{5760} - a_{4} + 2 a_{7}.
\end{align*}
\begin{align*}
R_{5,5}={}&- \frac{a_{1}^{2} a_{2}^{3} q^{4}}{30} - \frac{a_{1} a_{2}^{2} a_{3} q^{4}}{6} + \frac{a_{2}^{2} a_{3} q^{4}}{4}\\
&{}+ \frac{a_{2}^{2} a_{4} q^{4}}{2} - \frac{a_{2} a_{3}^{2} q^{4}}{2} + \frac{7 a_{2} a_{3} q^{4}}{24}\\
&{}+ \frac{a_{2} a_{4} q^{4}}{2} - \frac{a_{2} a_{5} q^{4}}{2} - a_{2} a_{7} q^{4}\\
&{}+ \frac{a_{2} q^{4}}{1440} - \frac{a_{3}^{2} q^{4}}{4} + a_{3} a_{5} q^{4}\\
&{}+ \frac{a_{3} q^{4}}{24} - \frac{a_{5} q^{4}}{12} + \frac{a_{8} q^{4}}{2}\\
&{}+ \frac{a_{1}^{2} a_{2}^{3} q^{3}}{12} + \frac{a_{1}^{2} a_{2}^{2} b_{2} q^{3}}{6} + \frac{3 a_{1}^{2} a_{2}^{2} q^{3}}{16}\\
&{}- \frac{a_{1} a_{2}^{3} b_{1} q^{3}}{6} - \frac{a_{1} a_{2}^{3} q^{3}}{16} + \frac{5 a_{1} a_{2}^{2} a_{3} q^{3}}{12}\\
&{}+ a_{1} a_{2} a_{3} b_{2} q^{3} + \frac{7 a_{1} a_{2} a_{3} q^{3}}{4} + \frac{a_{1} a_{2} b_{2} q^{3}}{24}\\
&{}+ \frac{a_{1} a_{2} q^{3}}{12} + \frac{a_{1} a_{3} b_{2} q^{3}}{2} + a_{1} a_{3} q^{3}\\
&{}- a_{1} a_{5} b_{2} q^{3} - 2 a_{1} a_{5} q^{3} - \frac{a_{2}^{2} a_{3} b_{1} q^{3}}{2}\\
&{}- \frac{3 a_{2}^{2} a_{3} q^{3}}{4} - \frac{a_{2}^{2} q^{3}}{32} + \frac{a_{2} a_{3}^{2} q^{3}}{2}\\
&{}+ \frac{a_{2} a_{3} b_{2} q^{3}}{2} + \frac{a_{2} a_{3} q^{3}}{4} + a_{2} a_{4} b_{2} q^{3}\\
&{}+ 2 a_{2} a_{4} q^{3} + \frac{3 a_{2} a_{5} q^{3}}{2} - \frac{a_{3}^{2} b_{2} q^{3}}{2}\\
&{}- \frac{a_{3}^{2} q^{3}}{4} + \frac{7 a_{3} b_{2} q^{3}}{24} + \frac{a_{3} q^{3}}{3}\\
&{}+ \frac{a_{4} b_{2} q^{3}}{2} + a_{4} q^{3} - \frac{a_{5} b_{2} q^{3}}{2}\\
&{}+ \frac{a_{5} q^{3}}{2} - a_{7} b_{2} q^{3} - 2 a_{7} q^{3}\\
&{}- \frac{3 a_{8} q^{3}}{2} - \frac{a_{1}^{2} a_{2}^{3} q^{2}}{12} - \frac{a_{1}^{2} a_{2}^{2} b_{2} q^{2}}{4}\\
&{}- \frac{5 a_{1}^{2} a_{2}^{2} q^{2}}{24} + \frac{a_{1}^{2} a_{2} b_{2} q^{2}}{4} + \frac{a_{1}^{2} a_{2} q^{2}}{3}\\
&{}+ \frac{a_{1} a_{2}^{3} b_{1} q^{2}}{4} + \frac{a_{1} a_{2}^{3} q^{2}}{8} - \frac{a_{1} a_{2}^{2} a_{3} q^{2}}{6}\\
&{}+ \frac{a_{1} a_{2}^{2} b_{1} b_{2} q^{2}}{2} + \frac{3 a_{1} a_{2}^{2} b_{1} q^{2}}{4} - \frac{a_{1} a_{2}^{2} q^{2}}{4}\\
&{}- \frac{a_{1} a_{2} a_{3} b_{2} q^{2}}{2} - \frac{a_{1} a_{2} a_{5} q^{2}}{3} + a_{1} a_{3} b_{2}^{2} q^{2}\\
&{}+ \frac{5 a_{1} a_{3} b_{2} q^{2}}{2} + a_{1} a_{3} q^{2} - \frac{a_{2}^{3} b_{1}^{2} q^{2}}{4}\\
&{}- \frac{a_{2}^{3} b_{1} q^{2}}{4} + \frac{a_{2}^{2} a_{3} b_{1} q^{2}}{2} + \frac{a_{2}^{2} a_{3} q^{2}}{6}\\
&{}- \frac{a_{2}^{2} a_{4} q^{2}}{6} + \frac{a_{2} a_{3}^{2} q^{2}}{3} + \frac{3 a_{2} a_{3} b_{1} q^{2}}{2}\\
&{}- \frac{3 a_{2} a_{3} b_{2} q^{2}}{2} - \frac{37 a_{2} a_{3} q^{2}}{24} + \frac{a_{2} b_{1} q^{2}}{12}\\
&{}- \frac{a_{2} b_{2} q^{2}}{16} - \frac{a_{2} b_{3} q^{2}}{24} - \frac{19 a_{2} q^{2}}{288}\\
&{}+ a_{3}^{2} q^{2} - a_{3} a_{5} q^{2} + a_{3} b_{1} q^{2}\\
&{}+ \frac{a_{3} b_{2}^{2} q^{2}}{4} + \frac{a_{3} b_{2} q^{2}}{4} - \frac{a_{3} b_{3} q^{2}}{2}\\
&{}- \frac{7 a_{3} q^{2}}{24} + \frac{a_{4} b_{2}^{2} q^{2}}{2} + 2 a_{4} b_{2} q^{2}\\
&{}+ a_{4} q^{2} - 2 a_{5} b_{1} q^{2} + \frac{3 a_{5} b_{2} q^{2}}{2}\\
&{}+ a_{5} b_{3} q^{2} + \frac{19 a_{5} q^{2}}{12} + \frac{a_{2}^{2} h q}{2}\\
&{}+ \frac{a_{1}^{2} a_{2}^{3} q}{24} + \frac{a_{1}^{2} a_{2}^{2} b_{2} q}{12} - \frac{a_{1}^{2} a_{2}^{2} q}{16}\\
&{}- \frac{a_{1}^{2} a_{2} b_{2} q}{4} - \frac{a_{1}^{2} a_{2} q}{4} - \frac{a_{1}^{2} b_{2}^{3} q}{4}\\
&{}- \frac{3 a_{1}^{2} b_{2}^{2} q}{4} - \frac{a_{1}^{2} b_{2} q}{2} - \frac{a_{1} a_{2}^{3} b_{1} q}{12}\\
&{}- \frac{a_{1} a_{2}^{3} q}{16} - \frac{a_{1} a_{2}^{2} a_{3} q}{6} + \frac{a_{1} a_{2}^{2} b_{1} q}{4}\\
&{}+ \frac{3 a_{1} a_{2}^{2} q}{8} - \frac{a_{1} a_{2} a_{3} b_{2} q}{2} - \frac{3 a_{1} a_{2} a_{3} q}{2}\\
&{}+ \frac{a_{1} a_{2} a_{5} q}{2} + \frac{a_{1} a_{2} b_{1} b_{2} q}{2} + a_{1} a_{2} b_{1} q\\
&{}- a_{1} a_{2} b_{2} b_{3} q - \frac{a_{1} a_{2} b_{2} q}{24} - 2 a_{1} a_{2} b_{3} q\\
&{}- \frac{a_{1} a_{2} q}{3} - a_{1} a_{3} b_{2} q - a_{1} a_{3} q\\
&{}+ a_{1} a_{5} b_{2} q + a_{1} a_{5} q + \frac{a_{2}^{2} a_{3} b_{1} q}{2}\\
&{}+ \frac{a_{2}^{2} a_{3} q}{4} - \frac{a_{2}^{2} a_{4} q}{4} + \frac{a_{2}^{2} b_{1}^{2} q}{4}\\
&{}- \frac{3 a_{2}^{2} b_{1} q}{4} + \frac{a_{2}^{2} b_{3} q}{2} - \frac{a_{2}^{2} b_{4} q}{2}\\
&{}+ \frac{a_{2}^{2} q}{32} - \frac{a_{2} a_{3}^{2} q}{4} + a_{2} a_{3} b_{1} q\\
&{}+ \frac{a_{2} a_{3} q}{2} - \frac{5 a_{2} a_{4} q}{4} - a_{2} a_{5} b_{1} q\\
&{}- \frac{a_{2} a_{5} q}{2} + \frac{a_{2} a_{7} q}{2} - \frac{a_{3}^{2} q}{2}\\
&{}+ \frac{a_{3} a_{5} q}{2} + 2 a_{3} b_{1} b_{2} q + a_{3} b_{1} q\\
&{}- \frac{3 a_{3} b_{2}^{2} q}{4} - \frac{37 a_{3} b_{2} q}{24} + \frac{a_{3} b_{3} q}{2}\\
&{}- a_{3} b_{5} q - \frac{7 a_{3} q}{12} + a_{2} b_{2} h\\
&{}+ \frac{a_{2} h}{2} - \frac{a_{1}^{2} a_{2}^{3}}{120} + \frac{a_{1}^{2} a_{2}^{2}}{12}\\
&{}- \frac{a_{1}^{2} a_{2}}{12} + \frac{a_{1} a_{2}^{2} a_{3}}{12} - \frac{a_{1} a_{2}^{2}}{8}\\
&{}- \frac{a_{1} a_{2} a_{3}}{4} - \frac{a_{1} a_{2} a_{5}}{6} + \frac{a_{1} a_{2}}{4}\\
&{}- a_{1} a_{3} + a_{1} a_{5} - a_{1} b_{1} b_{2}^{2}\\
&{}- a_{1} b_{1} b_{2} + \frac{a_{1} b_{2}^{3}}{4} - \frac{a_{1} b_{2}^{2} b_{3}}{2}\\
&{}+ \frac{3 a_{1} b_{2}^{2}}{4} - \frac{5 a_{1} b_{2} b_{3}}{2} + a_{1} b_{2} b_{5}\\
&{}+ \frac{a_{1} b_{2}}{2} - 2 a_{1} b_{3} + 2 a_{1} b_{5}\\
&{}+ \frac{a_{2}^{2} a_{3}}{12} - \frac{a_{2}^{2} a_{4}}{12} - \frac{a_{2} a_{3}^{2}}{12}\\
&{}+ \frac{a_{2} a_{3}}{2} - \frac{5 a_{2} a_{4}}{4} - \frac{a_{2} a_{5}}{2}\\
&{}+ \frac{a_{2} a_{7}}{2} + \frac{a_{2} b_{1}^{2}}{2} - 2 a_{2} b_{1} b_{3}\\
&{}- \frac{7 a_{2} b_{1}}{12} + a_{2} b_{2} b_{3} - a_{2} b_{2} b_{4}\\
&{}+ \frac{a_{2} b_{2}}{16} + \frac{a_{2} b_{3}^{2}}{2} + \frac{25 a_{2} b_{3}}{24}\\
&{}- \frac{5 a_{2} b_{4}}{2} - a_{2} b_{5} + a_{2} b_{7}\\
&{}+ \frac{47 a_{2}}{720} - \frac{a_{3} a_{5}}{2} + \frac{a_{3}}{2}\\
&{}- 2 a_{4} - 2 a_{5} + 2 a_{7}\\
&{}+ a_{8}.
\end{align*}
\begin{align*}
R_{5,6}={}&- \frac{a_{1} a_{2}^{4} q^{4}}{120} - \frac{a_{2}^{3} a_{3} q^{4}}{6} - \frac{a_{2}^{3} q^{4}}{144}\\
&{}- \frac{a_{2}^{2} a_{3} q^{4}}{4} + \frac{a_{2}^{2} a_{5} q^{4}}{2} - \frac{a_{2} a_{3} q^{4}}{6}\\
&{}+ a_{2} a_{5} q^{4} - a_{2} a_{8} q^{4} + \frac{a_{2} q^{4}}{1440}\\
&{}- \frac{a_{3} q^{4}}{24} + \frac{7 a_{5} q^{4}}{12} - \frac{3 a_{8} q^{4}}{2}\\
&{}+ \frac{a_{1} a_{2}^{4} q^{3}}{48} - \frac{a_{1} a_{2}^{3} q^{3}}{16} - \frac{a_{2}^{4} b_{1} q^{3}}{24}\\
&{}+ \frac{a_{2}^{3} a_{3} q^{3}}{24} - \frac{a_{2}^{2} a_{3} b_{2} q^{3}}{2} - \frac{3 a_{2}^{2} a_{3} q^{3}}{4}\\
&{}- \frac{a_{2}^{2} b_{2} q^{3}}{48} - \frac{a_{2}^{2} q^{3}}{32} - \frac{a_{2} a_{3} b_{2} q^{3}}{2}\\
&{}- \frac{3 a_{2} a_{3} q^{3}}{4} + a_{2} a_{5} b_{2} q^{3} + \frac{3 a_{2} a_{5} q^{3}}{2}\\
&{}- \frac{a_{3} b_{2} q^{3}}{6} - \frac{a_{3} q^{3}}{4} + a_{5} b_{2} q^{3}\\
&{}+ \frac{3 a_{5} q^{3}}{2} - a_{8} b_{2} q^{3} - \frac{3 a_{8} q^{3}}{2}\\
&{}- \frac{a_{1} a_{2}^{4} q^{2}}{72} + \frac{a_{1} a_{2}^{3} q^{2}}{8} - \frac{11 a_{1} a_{2}^{2} q^{2}}{72}\\
&{}+ \frac{a_{2}^{3} a_{3} q^{2}}{12} - \frac{a_{2}^{3} b_{1} q^{2}}{4} + \frac{a_{2}^{3} b_{3} q^{2}}{6}\\
&{}+ \frac{a_{2}^{3} q^{2}}{144} + \frac{a_{2}^{2} a_{3} q^{2}}{3} - \frac{a_{2}^{2} a_{5} q^{2}}{6}\\
&{}- \frac{a_{2} a_{3} b_{2}^{2} q^{2}}{2} - \frac{3 a_{2} a_{3} b_{2} q^{2}}{2} - \frac{11 a_{2} a_{3} q^{2}}{12}\\
&{}- \frac{a_{2} b_{2}^{2} q^{2}}{48} - \frac{a_{2} b_{2} q^{2}}{16} - \frac{11 a_{2} q^{2}}{288}\\
&{}- \frac{a_{3} b_{2}^{2} q^{2}}{4} - \frac{3 a_{3} b_{2} q^{2}}{4} - \frac{11 a_{3} q^{2}}{24}\\
&{}+ \frac{a_{5} b_{2}^{2} q^{2}}{2} + \frac{3 a_{5} b_{2} q^{2}}{2} + \frac{11 a_{5} q^{2}}{12}\\
&{}- \frac{a_{2}^{2} h q}{2} - \frac{a_{1} a_{2}^{3} q}{16} + \frac{11 a_{1} a_{2}^{2} q}{48}\\
&{}- \frac{a_{1} a_{2} q}{8} + \frac{a_{2}^{3} a_{3} q}{24} + \frac{a_{2}^{2} a_{3} q}{2}\\
&{}- \frac{a_{2}^{2} a_{5} q}{4} - \frac{11 a_{2}^{2} b_{1} q}{24} + \frac{a_{2}^{2} b_{2} b_{3} q}{2}\\
&{}+ \frac{a_{2}^{2} b_{2} q}{48} + a_{2}^{2} b_{3} q - \frac{a_{2}^{2} b_{5} q}{2}\\
&{}+ \frac{a_{2}^{2} q}{32} + \frac{11 a_{2} a_{3} q}{12} - \frac{5 a_{2} a_{5} q}{4}\\
&{}+ \frac{a_{2} a_{8} q}{2} - \frac{a_{3} b_{2}^{3} q}{6} - \frac{3 a_{3} b_{2}^{2} q}{4}\\
&{}- \frac{11 a_{3} b_{2} q}{12} - \frac{a_{3} q}{4} - a_{2} b_{2} h\\
&{}- \frac{3 a_{2} h}{2} + \frac{a_{1} a_{2}^{4}}{720} - \frac{11 a_{1} a_{2}^{2}}{144}\\
&{}+ \frac{a_{1} a_{2}}{8} + \frac{a_{1} b_{2}^{4}}{24} + \frac{a_{1} b_{2}^{3}}{4}\\
&{}+ \frac{11 a_{1} b_{2}^{2}}{24} + \frac{a_{1} b_{2}}{4} + \frac{a_{2}^{2} a_{3}}{6}\\
&{}- \frac{a_{2}^{2} a_{5}}{12} + \frac{11 a_{2} a_{3}}{12} - \frac{5 a_{2} a_{5}}{4}\\
&{}+ \frac{a_{2} a_{8}}{2} - \frac{a_{2} b_{1}}{4} + \frac{a_{2} b_{2}^{2} b_{3}}{2}\\
&{}+ \frac{a_{2} b_{2}^{2}}{48} + 2 a_{2} b_{2} b_{3} - a_{2} b_{2} b_{5}\\
&{}+ \frac{a_{2} b_{2}}{16} + \frac{11 a_{2} b_{3}}{6} - \frac{5 a_{2} b_{5}}{2}\\
&{}+ a_{2} b_{8} + \frac{3 a_{2}}{80} + a_{3}\\
&{}- 3 a_{5} + 3 a_{8}.
\end{align*}
\endgroup


\begin{statement}{D.33}{Proved}{An exact nonclosed lifting locus}
In the family of D.28 obtained by replacing
\(5^r\) with \(5^ru_r\),
\(u=(u_r)\in\prod_{r\ge1}\mathbb Z_5\),
the resulting actual class-two section \(s_u\)
lifts to class three if and only if
\[
                         u_r\longrightarrow0.             \tag{D.33.1}
\]
Every member has vanishing obstruction with
all finite coefficient modules. The map
\(u\mapsto[s_u]\) is a continuous injection
for the product topology on the parameters
and the quotient of the uniform topology on
continuous cocycles. In particular, the image
of the integral lifting map from class three
to class two is not closed: the finite-support
truncations of \((1,1,1,\ldots)\) give
liftable classes converging to the nonliftable
class of D.28.
\end{statement}
\begin{proof}
All ingredients of the constructed cocycle are
unramified outside
\(\{2,3,5,q_1,q_2,\ldots\}\). At each
\(q_r\), its degree-one inertia in the
diagonal coordinates of D.28 is
\(5^ru_rY\), and its normalized degree-two
Frobenius coordinate \(d'_r\) is a unit:
the entire perturbation is divisible by five,
whereas the actual seed has the chosen central
unit coordinate modulo five. The exact local
equation in D.28 gives the degree-three
inertia correction
\[
              \frac{5^ru_r d'_r}{q_r^2-1}[C,Y].             \tag{D.33.2}
\]
Its scalar coefficient is a unit times \(u_r\),
since \(v_5(q_r+1)=r\) and \(q_r-1\) is
a unit. Passing between this local normalization
and the fixed global BCH set-lift changes the
correction by a family tending to zero. Indeed,
the normalizing degree-one conjugators divide
only by \(q_r-1\), hence lie in a fixed
compact integral group, and inertia tends to
one uniformly. Conjugation and the two
set-lifts are continuous, so their central
differences tend uniformly to zero on that
compact set. Equivalently this follows by
substitution in the explicit formulas of D.30.

At two and three there are only finitely many
coordinates, and their weight-three local
obstruction groups vanish because \(q^2-1\)
is a five-adic unit. At every other prime
different from five and the \(q_r\)'s the
inertia is trivial, so the tame residue is
zero. Thus the pointwise divisibility holds
for every parameter sequence, and the
divided-residue convergence condition of D.29
holds precisely when (D.33.1) does. This proves
the equivalence, including sufficiency.
The finite-coefficient vanishing was proved
uniformly for this family in D.28.

Both the Kummer perturbations and their
chosen integral primitives have norm at most
one. The tails of the series defining \(s_u\)
therefore have norm at most \(5^{-N}\),
uniformly in all parameters, after the first
\(N\) terms. Dependence on each finite
parameter block is continuous. This proves
continuity into the space of continuous
cocycles with its uniform topology.
Distinct parameter sequences have distinct
second abelian Kummer classes, since the
valuation at \(q_r\) is exactly \(5^ru_r\).
They cannot become conjugate nonabelian
classes. Each finite-support truncation
satisfies (D.33.1), and the constant sequence
one does not. The stated convergence and
nonclosedness follow.
\end{proof}

\begin{statement}{D.34}{Proved}{The explicit cyclotomic descent equation for every central stage}
Let \(\Delta=G_{\mathbb Q(\mu_{5^\infty})}\),
\(\Gamma=G_{\mathbb Q}/\Delta=C_4\times\Gamma_1\),
and choose \(\gamma\in\Gamma_1\) with
\(\chi_5(\gamma)=1+5\). Fix an actual
lower-stage cocycle \(s:G_{\mathbb Q}\to T_{n-1}\),
and put \(L=L_n\). The exact continuous lifts
of \(s|_\Delta\) to \(T_n\) form a nonempty
affine space under
\[
          J=\operatorname{Hom}_{\rm cts}(\Delta,L)
            =\varprojlim_m H^1(\Delta,L/5^m).
                                                        \tag{D.34.1}
\]
For any such lift \(t\), and any \(g\in G_{\mathbb Q}\),
choose a lift \(\widetilde S_g\) of the arithmetic
element \((s(g),g)\) in \(T_n\rtimes G_{\mathbb Q}\).
Writing \(\widetilde t(\delta)=(t(\delta),\delta)\),
the formula
\[
 (g*t)(\delta)=
 \widetilde S_g\,\widetilde t(g^{-1}\delta g)
       \widetilde S_g^{-1}
 \quad\hbox{in the arithmetic fibre above }\delta
                                                        \tag{D.34.2}
\]
defines an affine action of \(\Gamma\).
After averaging over \(C_4\), choose \(t\) fixed
by \(C_4\). Then
\[
        d_t=\gamma*t-t\in J^{C_4},\qquad
 s\text{ lifts over }\mathbb Q
   \Longleftrightarrow
       d_t\in(\gamma-1)J^{C_4}.                 \tag{D.34.3}
\]
The image in (D.34.3) is not replaced by its
closure. Under the finite-orbit residue
comparison, the class of \(d_t\) is exactly
the family \(q^{-1}\Delta_q\) of D.29.
Thus (D.34.3), evaluated through weights three,
four and five, is precisely D.30 and D.32,
with all their divisibility and convergence
conditions. The \(\Gamma\)-descent contribution
to the fibre is
\[
 \begin{array}{c|ccc}
       n&3&4&5\\\hline
 H^1(\Gamma,L_n)&0&(\mathbb Z/5)^3&0\\
 J^\Gamma&\mathbb Z_5^2&0&\mathbb Z_5^6 .
 \end{array}                                             \tag{D.34.4}
\]
\end{statement}
\begin{proof}
On \(L_n\) the action of \(\Delta\) is trivial.
D.15 proves \(H^2_{\rm cts}(\Delta,L_n)=0\),
so exact lifts \(t\) exist. Their differences
are the homomorphisms (D.34.1); there are no
central coboundaries over \(\Delta\).
The topology on \(J\) is the inverse limit of
the discrete topologies on the finite-coefficient
homomorphism groups, equivalently uniform
convergence modulo every \(5^m\).

In (D.34.2), replacing \(\widetilde S_g\) by
another lift multiplies it by an element of
\(L\). That element commutes with
\(\widetilde t(\Delta)\), since \(\Delta\)
acts trivially on \(L\). Hence the action
does not depend on this choice. Products of
two chosen lifts differ from the chosen lift
of their product by an element of \(L\),
so (D.34.2) is an actual group action.
For \(g\in\Delta\), choose
\(\widetilde S_g=\widetilde t(g)\); the
homomorphism identity for \(\widetilde t\)
then shows that \(g*t=t\).
Thus the action factors through \(\Gamma\).
Its linear part on \(J\) is the usual
conjugation action on \(\Delta\), combined
with multiplication by \(\chi_5(g)^n\)
on \(L\). All the maps are continuous:
uniform continuity on the compact groups
involved gives continuity into each discrete
finite-coefficient homomorphism group.

The translation \(g*t-t\) is a continuous
one-cocycle of \(\Gamma\) with values in
\(J\). Averaging an affine \(C_4\)-orbit is
legitimate because four is invertible on
\(J\); it supplies a \(C_4\)-fixed \(t\).
For this choice the value \(d_t\) lies in
\(J^{C_4}\). Replacing \(t\) by \(t+b\)
changes it to \(d_t+(\gamma-1)b\).
Consequently a \(\Gamma\)-fixed exact
\(\Delta\)-lift exists exactly when the
right side of (D.34.3) holds.

Such a fixed lift descends to an actual
global lift, not merely to an invariant
conjugacy class. To see this, let
\(\mathcal E\) be the pullback of
\(T_n\rtimes G_{\mathbb Q}\) above the
section \(s\), an extension of
\(G_{\mathbb Q}\) by \(L\).
The subgroup \(\widetilde t(\Delta)\) is
closed and normal in \(\mathcal E\) by
(D.34.2), and has trivial intersection
with \(L\). Its quotient gives
\[
 1\longrightarrow L\longrightarrow
 \mathcal E/\widetilde t(\Delta)
       \longrightarrow\Gamma\longrightarrow1.
                                                        \tag{D.34.5}
\]
Continuous \(H^2(\Gamma,L)\) is zero.
Indeed \(\Gamma\) has five-cohomological
dimension one on discrete finite coefficients;
its first cohomology with \(L/5^m\) is finite,
so the derived-limit term vanishes as well.
Thus (D.34.5) has a continuous splitting.
Moreover
\(\mathcal E\simeq
G_{\mathbb Q}\times_\Gamma
(\mathcal E/\widetilde t(\Delta))\):
injectivity follows from the trivial
intersection, and surjectivity follows by
adjusting any lift by an element of
\(\widetilde t(\Delta)\).
Pullback of the splitting proves descent.
The converse is immediate by restriction.
This also proves directly the central
obstruction interpretation of (D.34.3),
in agreement with D.15.

We spell out how the residue expressions arise
from this descent calculation. Above a rational
prime \(q\ne5\), the cyclotomic decomposition
subgroup is \(H_q=\overline{\langle q\rangle}
\subset\Gamma\). It has finite index:
\(q\), a positive rational integer greater
than one, has nonzero five-adic logarithm
after its Teichm\"uller part is removed.
Thus the primes of the cyclotomic field
above \(q\) form the finite orbit
\(\Gamma/H_q\). The corresponding residue
module is
\[
       \operatorname{Ind}_{H_q}^{\Gamma}L(-1).
                                                        \tag{D.34.6}
\]
At finite coefficients, its derived
\(\Gamma\)-invariants are, by the explicit
finite-orbit Shapiro map, the two-term
complex
\[
             [\,L/5^m
                 \xrightarrow{q^{n-1}-1}L/5^m\,].
                                                        \tag{D.34.7}
\]
One may construct this map by choosing an
orbit representative and transporting around
the orbit. The inverse transports a value to
all orbit representatives; the homotopy uses
partial sums of transported values. All
coefficients are five-adic units or sums
of such coefficients, hence have norm at
most one, uniformly in the orbit length.
These maps therefore pass to direct sums
over \(q\), and then to their derived
inverse limits without losing the \(c_0\)
condition. The resulting residue complex
is exactly
\[
                [\,c_0(L)
                     \xrightarrow{D_n}c_0(L)\,].
                                                        \tag{D.34.8}
\]
This is also obtained by applying derived
\(\Gamma\)-invariants to the localization
triangle over the cyclotomic field.
The arithmetic base term is
\(R\Gamma(G_{\mathbb Q,\{5\}},L)\);
its cohomology in degrees two and three
vanishes by D.29. Hence (D.34.8) computes
the complete global degree-two obstruction,
including the unclosed image.

There is a direct check of the representative
which fixes its normalization. Put
\(v_q=s(\tau_q)\), \(w_q=s(\sigma_q)\), and
let \(j\) be the global Hall set-section.
The restricted lift \(t\) sends tame inertia
to a homomorphism, whose generator has the
form \(j(v_q)u_q\), \(u_q\in L\).
The family \(u_q\) tends to zero:
\(\tau_q\to1\) in \(\Delta\), and both
the continuous \(t\) and \(j\circ s\)
send these elements to one.
Transport around the Frobenius orbit in
(D.34.6) evaluates the affine descent defect
on \(\tau_q\) by
\[
 \widetilde S_{\sigma_q}
     \widetilde t(\tau_q^{1/q})
     \widetilde S_{\sigma_q}^{-1}
     \widetilde t(\tau_q)^{-1}
      =q^{-1}\Delta_q+(q^{n-1}-1)u_q
           \quad\hbox{in }L.                  \tag{D.34.9}
\]
Here \(\Delta_q\) is the actual group
relator (D.29.1). Its factor \(q^{-1}\)
comes from evaluating on
\(\sigma_q^{-1}\tau_q\sigma_q=\tau_q^{1/q}\).
Thus the cocycle \(d_t\), followed by the
residue map and finite-orbit Shapiro, gives
exactly the class of \(q^{-1}\Delta_q\)
in (D.34.8); the correcting family is
in \(c_0(L)\), not merely in the full
product. Evaluating this same relator
in the integral class-three BCH coordinates
gives (D.30.2). Evaluating it in the
integral Hall coordinates gives the
three and six polynomials of D.32.
This derives all three lifting criteria
as \(\Gamma\)-descent equations and
retains the derived-limit term.

Finally, restriction and inflation give
the exact sequence
\[
 0\longrightarrow H^1(\Gamma,L_n)
 \longrightarrow H^1(G_{\mathbb Q},L_n)
 \longrightarrow J^\Gamma\longrightarrow0.
                                                        \tag{D.34.10}
\]
The last zero follows from
\(H^2_{\rm cts}(\Gamma,L_n)=0\), just
proved. If \(n=3\) or \(5\), the
nontrivial residual \(C_4\)-character
kills \(L_n^{C_4}\), so
\(H^1(\Gamma,L_n)=0\). If \(n=4\),
that character is trivial and
\[
 H^1(\Gamma,\mathbb Z_5(4))
  =\mathbb Z_5/((1+5)^4-1)\mathbb Z_5
  =\mathbb Z/5.
 \]
Its generator is exactly \(z_4\) of
(D.32.8). Inserting the global groups
computed in D.1 proves (D.34.4).
Together with the complete convergent
free-basis parameterization of D.16,
this supplies the requested cyclotomic
section sets and their global descent.
\end{proof}

\begin{statement}{D.35}{Proved}{The full Alexander action, its free resolution, and a torsion-free metabelian fibre}
Put \(M=F^{(5)}/(F^{(5)})''\), and
\[
 A=M'=\Lambda C,\qquad
 \Lambda=\mathbb Z_5[[u,v]],\quad
 t_1=1+u,\quad t_2=1+v,\quad C=[x,y].
\]
Write \(f_g=B_gC\in A\), additively, for the
actual interval Ihara cocycle. For any abelian
Kummer cocycle \(s=(a,b)\), its twisted
Alexander module \(A_s\) has the explicit action
\[
 \rho_s(g)m=
 t_1^{a(g)}t_2^{b(g)}
 S_{\chi_5(g)}(t_1)S_{\chi_5(g)}(t_2)
       (1+uvB_g)\,
 m(t_1^{\chi_5(g)}-1,t_2^{\chi_5(g)}-1),
 \quad S_z(t)=\frac{t^z-1}{t-1}.               \tag{D.35.1}
\]
Here \(m\) is written as a power series in
\(u,v\); all substitutions and binomial series
are integral and continuous.

For every \(s\), including those with
unrestricted ramification,
\[
 H^0(\Delta,A_s/5)=H^0(\Delta,A_s)=0,
 \qquad H^1_{\rm cts}(G_{\mathbb Q},A_s)
                    \text{ is torsion free}.              \tag{D.35.2}
\]
More precisely, let \(d_1,d_2,\ldots\to1\)
be a free basis of \(\Delta^{(5)}\), as in
D.16, put
\(\mu_i=t_1^{a(d_i)}t_2^{b(d_i)}(1+uvB_{d_i})\),
and let \(c_0(\Lambda)\) denote sequences
tending to zero in the \((5,u,v)\)-adic
topology. Then
\[
 0\longrightarrow\Lambda
   \xrightarrow{\,m\mapsto((\mu_i-1)m)_i\,}
   c_0(\Lambda)\longrightarrow J_s
   \longrightarrow0,\qquad
 J_s=H^1_{\rm cts}(\Delta,A_s),                 \tag{D.35.3}
\]
is an exact sequence of topological modules
with closed first image. With the actual
induced \(\Gamma\)-action specified below,
\[
 H^0(G_{\mathbb Q},A_s)=0,\qquad
 H^1(G_{\mathbb Q},A_s)=J_s^\Gamma,\qquad
 H^2(G_{\mathbb Q},A_s)=
       J_s^{C_4}/(\gamma-1)J_s^{C_4}.           \tag{D.35.4}
\]
The last image remains unclosed. Every
nonempty metabelian section fibre over \(s\)
is an \(H^1(G_{\mathbb Q},A_s)\)-torsor,
whose translation group is torsion free.
The three complete cusp packets in \(M\)
are given explicitly in (D.35.8) below.
\end{statement}
\begin{proof}
The Fox--Magnus description, proved in D.20,
realizes \(M\) as the pairs
\((T,\mathbf z)\), where
\(T=t_1^\alpha t_2^\beta\),
\(\mathbf z=z_1e_1+z_2e_2\), and
\[
      uz_1+vz_2=T-1,\qquad
 (T,\mathbf z)(T',\mathbf z')
       =(TT',\mathbf z+T\mathbf z').
                                                        \tag{D.35.5}
\]
Here \(x=(t_1,e_1)\), \(y=(t_2,e_2)\),
and \(C=(1,-ve_1+ue_2)\). The kernel of
the row \((u,v)\) is freely generated by
\((-v,u)\), since \(u,v\) is a regular
sequence in \(\Lambda\). Thus the derived
module is exactly \(\Lambda C\).

For \(\chi=\chi_5(g)\), the images of
the two vectors \(e_1,e_2\) under the
arithmetic action are
\[
 e_1\longmapsto S_\chi(t_1)e_1,\qquad
 e_2\longmapsto S_\chi(t_2)e_2
                       +(t_2^\chi-1)B_g C.
                                                        \tag{D.35.6}
\]
The coefficients are transformed by
\(t_i\mapsto t_i^\chi\). The commutator
of these two group elements has vector
\[
 (1-t_2^\chi)S_\chi(t_1)e_1
 +(t_1^\chi-1)
       \{S_\chi(t_2)e_2+(t_2^\chi-1)B_gC\}.
 \]
Collecting its \(C\)-coefficient gives
\(S_\chi(t_1)S_\chi(t_2)(1+uvB_g)\).
Conjugation by any lift of the abelian
element \((a(g),b(g))\) multiplies \(A\)
by \(t_1^{a(g)}t_2^{b(g)}\), proving
(D.35.1). Different lifts differ by
elements of the abelian group \(A\), so
the formula is independent of that
choice and is an actual group action.
The generic polynomial identity is
checked in \texttt{code/de\_alexander.py}.

Here is a primitive arithmetic element that
rules out invariants for every twist.
Let \(\zeta=\zeta_5\) and
\(\phi=-\zeta^2-\zeta^3=(1+\sqrt5)/2\).
The raw Soul\'e unit (D.31.3) at \(m=1\)
has the exact factorization
\[
 \prod_{a=1}^4(1-\zeta^a)^{a^2}
       =(1-\zeta)^{30}\phi^{13}
       =\phi^3\bigl((1-\zeta)^6\phi^2\bigr)^5.
                                                        \tag{D.35.7}
\]
For example,
\(1+\zeta=-\phi\zeta^3\),
\(1+\zeta+\zeta^2=\phi\zeta\), and
\(1+\zeta+\zeta^2+\zeta^3=-\zeta^4\)
prove the identity immediately; the code
also verifies it exactly modulo
\(\Phi_5(\zeta)\). In the units modulo
fifth powers of \(\mathbb Q(\zeta_5)\),
\(\phi\) and \(\zeta_5\) are independent,
by the unit calculation in D.1.
The unique degree-five subextension of
the cyclotomic \(\mathbb Z_5\)-extension
has Kummer class \(\zeta_5\).
Thus the Kummer character of \(\phi\),
and hence \(\kappa_3\bmod5\), is nonzero
on \(\Delta\). Choose
\(\delta\in\Delta\) with
\(h_\delta=\kappa_3(\delta)/48\ne0\bmod5\).

As \(\chi_5(\delta)=1\), D.31 gives
\(\log f_\delta=h_\delta(U+V)+O(4)\).
Consequently the first nonzero homogeneous
term of \(uvB_\delta\bmod5\) is
\(h_\delta uv(u+v)\), of degree three.
The multiplier of \(\delta\) on
\(A_s/5=\mathbb F_5[[u,v]]C\) is
\[
 t_1^{a(\delta)}t_2^{b(\delta)}
                       (1+uvB_\delta).
\]
If either \(a(\delta)\) or \(b(\delta)\)
is nonzero modulo five, this multiplier
has a nonzero degree-one term.
If both are divisible by five, their
monomial has no nonconstant term below
degree five, by the characteristic-five
identity \((1+u)^5=1+u^5\). Its product
therefore retains the nonzero degree-three
term just exhibited. In all cases the
multiplier is different from one.
The power-series ring over \(\mathbb F_5\)
is a domain, so this single element
annihilates all possible invariants.
This proves \(H^0(\Delta,A_s/5)=0\).
It also proves \(H^0(\Delta,A_s)=0\)
by successive division by five and
five-adic separation. The long exact
sequence for multiplication by five on
\(A_s\) now proves that
\(H^1(G_{\mathbb Q},A_s)\) has no
five-torsion; all its finite-order
elements would have five-power order.

Over \(\Delta\), formula (D.35.1) is
multiplication by a continuous character
with values in \(1+(u,v)\Lambda\),
a pro-five group. Thus it factors through
the free pro-five group
\(\Delta^{(5)}\). A continuous crossed
homomorphism is uniquely specified by
arbitrary values on \(d_i\) tending to
zero: apply the universal property to
the semidirect product
\(A_s\rtimes\operatorname{im}(\rho_s|_\Delta)\).
Principal crossed homomorphisms have
values \(((\mu_i-1)m)_i\). Their map is
injective by the invariant calculation.
Its image is compact, hence closed in
the Hausdorff space \(c_0(\Lambda)\).
This proves (D.35.3), including its
topology; equivalently one can pass
through the finite \((5,u,v)\)-adic
quotients and choose compatible
representatives, since the coboundary
sets at every finite level are finite.

For completeness, the operator on this
resolution is completely specified as
follows. A tuple \(\mathbf z\) represents
a crossed homomorphism \(z\). For a
lift \(g\in G_{\mathbb Q}\) of an
element of \(\Gamma\), the transformed
tuple has coordinates
\[
       \bigl(\rho_s(g)\,
             z(g^{-1}d_i g)\bigr)_{i\ge1}.
\]
The values on the words \(g^{-1}d_i g\)
are computed by
\[
 z(w)=\sum_j
       \left(\frac{\partial w}{\partial d_j}\right)_{\mu}
                         z(d_j),
\]
the continuous Fox derivative evaluated
at \(d_j\mapsto\mu_j\).
One may equivalently evaluate successive
finite words modulo every power of
\((5,u,v)\); the answers stabilize,
so no unjustified infinite sum is
introduced. Changing the lift \(g\)
changes this tuple by a principal
crossed homomorphism. Hence it defines
the actual \(\Gamma\)-action on \(J_s\).

Apply the argument of D.15 to the finite
coefficient tower \(A_s/(5,u,v)^mA_s\).
Its degree-one cohomology over \(\Delta\)
has surjective transition maps, because
\(\operatorname{cd}_5\Delta\le1\).
Its degree-zero groups are finite.
Their inverse limit is
\(H^0(\Delta,A_s)=0\), so that finite
tower is pro-zero. It contributes
nothing after derived inverse limit.
At each finite coefficient level, apply
the standard procyclic resolution to
the discrete module \(H^1(\Delta,A_s/(5,u,v)^mA_s)\);
this discrete group need not be finite.
Its transition maps, and those after
taking \(C_4\)-invariants, are onto.
Thus the derived inverse limit of
these two-term complexes is their
ordinary termwise inverse limit.
The resulting two-term complex for
\(\Gamma\) is
\([J_s^{C_4}\xrightarrow{\gamma-1}J_s^{C_4}]\),
shifted by one. This proves (D.35.4)
for continuous cohomology, without
discarding a derived-limit term.
The ordinary abelian lifting torsor
has no stabilizer quotient:
\(H^0(G_{\mathbb Q},\mathbb Z_5(1)^2)=0\).
This proves the assertion about every
nonempty metabelian fibre.

We finally give the cusp maps in these
coordinates. Put
\(d=y^{-1/2}x^{-1}y^{-1/2}\), a conjugate
of the third boundary element, and let
\((T,\mathbf v)\) be its pair in (D.35.5).
Explicitly
\[
 T=(t_1t_2)^{-1},\quad
 \mathbf v=-t_1^{-1}t_2^{-1/2}e_1
       +S_{-1/2}(t_2)(1+t_1^{-1}t_2^{-1/2})e_2.
 \]
Let \(\mathbf v_g\) be its arithmetic
transform using (D.35.6). There is a
unique \(K_g\in\Lambda\) satisfying
\[
       (1-T^{\chi_5(g)})K_g C
           = S_{\chi_5(g)}(T)\mathbf v-\mathbf v_g.
 \]
The division asserted here is integral.
To verify it, start with the actual
cuspidal conjugator for \(d\).
The logarithm of \(d\) has zero term
in degree two, so comparison in degree
two shows that the abelianization of
this conjugator is parallel to \((-1,-1)\).
Multiplying it by a power of \(d\)
therefore makes its abelianization zero.
It then lies in \(A\) and supplies
the required \(K_g\). Uniqueness holds
since \(\Lambda\) is a domain and
\(1-T^{\chi_5(g)}\ne0\).
It also proves the cocycle identity
for \(k_g=(1,K_g C)\).
This is the zero-abelian representative
of the unit tangent packet at infinity;
the unit tangent has zero abelian
Kummer class because the possible
sign is \(-1\), which is five-divisible
over \(\mathbb Q\).

For arbitrary continuous Kummer cocycles
\(\eta:G_{\mathbb Q}\to\mathbb Z_5(1)\),
the complete three packets are
\[
 \begin{array}{c|c|c}
 \text{cusp}&\text{metabelian cocycle}
                         &\text{abelian class}\\\hline
 0&x^{\eta(g)}&(\eta,0)\\
 1&y^{\eta(g)}f_g&(0,\eta)\\
 \infty&d^{\eta(g)}k_g&(-\eta,-\eta).
 \end{array}                                             \tag{D.35.8}
\]
The unit tangents correspond to \(\eta=0\).
The twisted arithmetic action fixes
the corresponding boundary element
up to its cyclotomic power, which
proves the cocycle identity for every
entry. Conversely the decomposition
group of a rational cusp is the
extension by its procyclic inertia,
whose sections differ exactly by
these Kummer cocycles. The primitive
abelian vectors in (D.35.8) show that
different packet parameters remain
different metabelian classes.
\end{proof}

\paragraph{What remains in Task D.5.}
Formulas (D.35.1)--(D.35.4) determine the
coefficient action and give a concrete
free-resolution operator. They do not
evaluate the kernel and unclosed
cokernel of its actual \(\Gamma\)-action
for arbitrary \(a,b\). That evaluation,
and the full obstruction represented
by \(x^{a(g)}y^{b(g)}\,
g(x^{a(h)}y^{b(h)})\,
(x^{a(gh)}y^{b(gh)})^{-1}\) in \(A_s\),
is the remaining Alexander calculation.
The torsion-free assertion is already
a restriction on any proposed answer:
the three order-five fibres at class
four cannot be assembled into an
independent finite torsion factor of
the full metabelian fibre. Their
extensions into higher weights must
be retained. The higher-derived
Fermat lifting has the finite-stage
nonexistence theorem D.57, while the
first obstructed length and its class
remain open. The
Legendre-relative prosolvable assertion
is proved in D.55, as recorded in D.5.

\begin{statement}{D.36}{Proved}{The complete unramified-outside-five metabelian section fibres}
Let \(G_5=G_{\mathbb Q,\{5,\infty\}}\).
For every \(G_5\)-abelian section
\[
            (a,b)=(A\kappa(5),B\kappa(5)),
                  \qquad (A,B)\in\mathbb Z_5^2,
\]
its full Alexander module satisfies
\[
 H^0(G_5,A_s)=H^2_{\rm cts}(G_5,A_s)=0,\qquad
 H^1_{\rm cts}(G_5,A_s)
       \simeq\prod_{j=1}^{\infty}\mathbb Z_5
       \quad\text{as a compact }\mathbb Z_5\text{-module}.
                                                        \tag{D.36.1}
\]
The isomorphism is noncanonical. Moreover
\[
       H^1(G_{\mathbb Q},A_s)=H^1(G_5,A_s).
\]
Consequently the entire unrestricted
metabelian section fibre over each of
these abelian pairs is represented over
\(G_5\), and
\[
       \operatorname{Sec}_{G_5}(M)
                  \longrightarrow\mathbb Z_5^2
                                                        \tag{D.36.2}
\]
is onto, and every fibre is a torsor under
the displayed product. No class in this
section set is the class of an ordinary
rational point of the tripod.
The cusp packets are the three maps of
D.35 with \(\eta=c\kappa(5)\).
Thus there are no cusp classes over
\((A,B)\) off the two coordinate axes
and the diagonal. On the nonzero part
of each of these three lines there
is exactly one cusp class per fibre;
over \((0,0)\) there are exactly three.
\end{statement}
\begin{proof}
The arithmetic action and the specified
abelian cocycle factor through \(G_5\).
The \(I=(u,v)\)-adic graded pieces of
\(A_s\), independently of \(A,B\), are
\[
             I^jA_s/I^{j+1}A_s
                  \simeq\mathbb Z_5(j+2)^{j+1}.
                                                        \tag{D.36.3}
\]
Indeed every Kummer monomial and every
arithmetic factor \(1+uvB_g\) has
constant term one, while the leading
term of (D.35.1) multiplies \(C\) by
\(\chi_5^2\) and each \(u,v\) by
\(\chi_5\). Every finite quotient
of \(I^jA_s\) by powers of five and
of \(I\) therefore has a finite
filtration whose factors are residual
Tate characters. Their \(H^2(G_5,-)\)
vanishes by D.1, extended to all
weights in D.12. The same holds for
the kernels of the transition maps.
Thus the \(H^1\) transition maps of
these finite towers are onto.
Jannsen's sequence gives
\[
          H^2_{\rm cts}(G_5,I^jA_s)=0
                   \quad(j\ge0).
                                                        \tag{D.36.4}
\]
The zero \(H^0\) follows either from
the first nonzero positive-weight
graded piece or from D.35.
For the torsion-free assertion, the
primitive element \(\delta\) in the
proof of D.35 can be chosen in the
cyclotomic kernel of \(G_5\):
the Kummer extension of the unit
\(\phi\) is unramified outside five.
Thus the same proof applies to
\(H^1(G_5,A_s)\).

This \(H^1\) is compact and metrizable.
The maximal pro-five quotient over
\(\mathbb Q(\zeta_5)\), unramified
outside five, is the free pro-five
group on three generators used in
D.20. The action on \(A_s\), and
any crossed homomorphism to it,
factor through its extension by
the group of order four: the
corresponding semidirect image is
pro-five over that finite group.
Cocycles are therefore specified by
finitely many elements of the compact
metrizable module \(A_s\), subject
to closed relations. Coboundaries
form a compact image as well.
Their quotient is compact metrizable.

For every \(N\), (D.36.4) gives a
surjection
\[
 H^1(G_5,A_s)\longrightarrow
                 H^1(G_5,A_s/I^NA_s).
\]
All the integral positive-weight
graded quotients have zero \(H^0\)
and zero \(H^2\), so their \(H^1\)
ranks add in this filtration.
The residual unit calculation of
D.1 gives, for all \(j\ge2\),
\[
 \operatorname{rank}_{\mathbb Z_5}
            H^1(G_5,\mathbb Z_5(j))
      =\begin{cases}1&j\text{ odd},\\0&j\text{ even}.
        \end{cases}
\]
One can read this directly from the
four residual characters: their
first-cohomology dimensions are
one, one, zero, one in weights
zero, one, two, three modulo four;
subtract the residual invariant
dimension, one only in weight zero.
The coefficient long exact sequence,
using integral \(H^2=0\), gives
the asserted integral ranks.
Consequently
\[
 \operatorname{rank}_{\mathbb Z_5}
       H^1(G_5,A_s/I^NA_s)
       =r(r+1),\qquad r=\lfloor N/2\rfloor .
                                                        \tag{D.36.5}
\]
These ranks are unbounded.
A compact metrizable torsion-free
\(\mathbb Z_5\)-module is a finite
or countable product of copies of
\(\mathbb Z_5\): its Pontryagin
dual is a countable divisible
five-primary torsion group, hence
a direct sum of copies of
\(\mathbb Q_5/\mathbb Z_5\).
The unbounded ranks (D.36.5)
exclude a finite product, proving
(D.36.1).

The Kummer calculation over
\(\mathbb Z[1/5]\) identifies the
abelian section set with
\(\mathbb Z_5^2\), generated in
each coordinate by \(\kappa(5)\).
The obstruction to a metabelian
lift lies in the zero group
\(H^2(G_5,A_s)\), and there is
no abelian fixed-element quotient,
as in D.35. This proves the
surjectivity and all the fibres
of (D.36.2).

To identify these with the entire
unrestricted fibres, consider any
global metabelian section over the
same abelian cocycle. At \(q\ne5\)
its inertia has zero abelian image,
so lies in \(A\). Wild inertia and
the prime-to-five tame factors act
trivially: the geometric action is
unramified, and their homomorphisms
to the pro-five group are zero.
If a five-primary tame generator
had image \(m\ne0\), choose its
first nonzero \(I\)-adic grade
\(I^jA/I^{j+1}A\). The tame
relation would give
\[
              \rho_s(\sigma_q)m=q m.
\]
On that first grade the left
side is \(q^{j+2}m\), by
(D.36.3). The nonzero integer
\(q^{j+2}-q\) cannot annihilate
an element of this torsion-free
lattice. This contradiction shows
that the section is unramified
at every \(q\ne5\), and hence
factors through \(G_5\).
The identical argument for a
crossed homomorphism into \(A_s\)
shows that every such global
crossed homomorphism is unramified
outside five. This proves the
asserted equality of \(H^1\)
and the unrestricted-fibre claim.
More generally the argument excludes
new ramification at any prime
outside five where the abelian
cocycle is unramified; it does
not remove the convergence problem
when infinitely many abelian
inertia coordinates are nonzero.

If a rational point had a section
in this set, its abelian Kummer
classes would be unramified outside
five. The valuation maps then give
\(x,1-x\in\mathbb Z[1/5]^\times
=\{\pm5^n:n\in\mathbb Z\}\).
Both of these units reduce to one
in \(\mathbb F_2\), contradicting
\(x+(1-x)=1\). Thus no ordinary
rational-point class occurs.
The cusp maps and their primitive
abelian vectors were proved in
D.35. Their allowed \(G_5\)-parameters
are exactly \(c\kappa(5)\), giving
the three indicated lines.

It remains only to distinguish the
three unit classes over the origin.
Restrict to the primitive
\(\delta\in\ker(\chi_5|_{G_5})\)
used above. Through weight three,
the coefficients of their Alexander
cocycles, in the basis \(U,V\), are
respectively
\[
                   (0,0),\qquad
                   h_\delta(1,1),\qquad
                   h_\delta(0,1).
\]
The middle value is D.31.
For the last, use
\(\log d=-X-Y+O(3)\) and the
weight-four action
\(\delta(Y)=Y+[Y,h_\delta(U+V)]+O(5)\).
Conjugation by \(aU+bV\) changes
\(-X-Y\) by
\(a(E_1+E_2)+b(E_2+E_3)\).
The equation defining \(k_\delta\)
therefore gives \(a=0,b=h_\delta\).
These three values are distinct.
A geometric conjugator between
zero-abelian global sections must
lie in \(A\), because the abelian
arithmetic action has no invariants.
Over this \(\delta\), conjugation
by \(A\) changes the Alexander
coefficient only by a multiple
of \(uvB_\delta\), which starts
in degree three. It cannot change
these linear coefficients.
Thus the three unit classes are
distinct, completing the count.
\end{proof}

\begin{statement}{D.37}{Proved}{The arithmetic Fermat metabelian comparison is not injective at any five-power level}
For every \(r\ge1\), let
\[
 K_r=\ker\{F_5\longrightarrow(\mathbb Z/5^r)^2\},
 \qquad A_r=F_5''/K_r'',\qquad F_5=F^{(5)}.
\]
The natural map
\[
 \operatorname{Sec}_{\mathbb Q}(K_r/K_r'')
       \longrightarrow
 \operatorname{Sec}_{\mathbb Q}(K_r/F_5'')
                                                        \tag{D.37.1}
\]
has at least 25 distinct elements in its
fibre over the unit tangential section
at zero. These classes are represented
by cocycles unramified outside five.
Thus the section-set equality proposed
in Task D.6, under its natural comparison,
is false at every five-power level.
\end{statement}
\begin{proof}
The group \(A_r\) is abelian:
\(F_5'\subset K_r\), hence
\(F_5''\subset K_r'\). We first exhibit
the finite arithmetic quotient which
detects different section classes.
Write \(C=[X,Y]\), \(U=[X,C]\),
\(V=[Y,C]\) in the graded free Lie
algebra of \(F_5\). The degree-four
part of \([F_5',F_5']\) is zero,
because its degree-two parts are
multiples of the same generator \(C\).
Thus \(F_5''\subset\gamma_5F_5\).
Its image in degree five is exactly
the free rank-two module generated
by \([C,U]\) and \([C,V]\). These
two elements are independent in the
degree-five Hall basis; see the last
two columns of (D.32.11).

Use the characteristic-five Magnus
embedding, with augmentation ideal
\(\mathfrak a=(\mathsf X,\mathsf Y)\).
Every element of \(K_r\) has trivial
abelianization modulo five, so has
Magnus series in \(1+\mathfrak a^2\).
Consequently
\[
 K_r'\subset1+\mathfrak a^4,\qquad
 K_r''\subset1+\mathfrak a^8
               \quad\hbox{in characteristic five}.
 \]
Therefore the degree-five map modulo
five kills \(K_r''\) and gives an
arithmetic equivariant surjection
\[
         A_r\longrightarrow
             \mathbb F_5(5)^2.                \tag{D.37.2}
\]
The arithmetic action on this leading
degree is exactly \(\chi_5^5\);
inner conjugations do not change it.
No group-level kernel alone is being
used to infer a section difference.

We now realize every cohomology class
of (D.37.2) by an actual arithmetic
cocycle. Let \(Q_5\) be the quotient
of \(G_5\) whose subgroup over
\(\mathbb Q(\zeta_5)\) is the maximal
pro-five extension unramified outside
five. As established in D.20, that
subgroup \(P\) is free pro-five of
rank three. The group of order four
has a complement, so \(Q_5=P\rtimes C_4\),
and \(\operatorname{cd}_5 Q_5=1\).
The actual action on \(F_5\) factors
through \(Q_5\): over \(\mathbb Q(\zeta_5)\)
its image lies in the pro-five kernel
of \(\operatorname{Aut}(F_5)\to
\operatorname{GL}(F_5/\Phi(F_5))\).
The same holds for the characteristic
subgroups \(K_r,K_r''\) and the
module \(A_r\).

If \(N_r\) is the kernel of (D.37.2),
then \(H^2_{\rm cts}(Q_5,N_r)=0\).
Indeed every finite quotient of
this compact module has zero
second cohomology, and the first
cohomology transition maps are onto
by the same dimension-one assertion
for their finite kernels. Their
derived-limit term is zero.
Hence
\[
 H^1(Q_5,A_r)\longrightarrow
             H^1(Q_5,\mathbb F_5(5)^2)
                    =\mathbb F_5^2
                    \longrightarrow0.         \tag{D.37.3}
\]
The last equality is the residual
weight-one calculation of D.1;
\(\mathbb F_5(5)=\mathbb F_5(1)\).
First cohomology with this module
already factors through \(Q_5\),
since a crossed homomorphism together
with its action has pro-five image
over the group of order four.
It is therefore the same group as
the corresponding \(G_5\)-cohomology.

Inflation of (D.37.3) first to
\(G_5\) and then to \(G_{\mathbb Q}\)
is injective in degree one.
Each of its 25 different images
in the finite target can thus be
represented by an actual distinct
\(A_r\)-cohomology class on
\(G_{\mathbb Q}\), unramified
outside five. Multiplying the unit
section by these cocycles produces
sections of \(K_r/K_r''\) which
all map to the unit section of
\(K_r/F_5''\).

There is no further identification
by geometric conjugation. In the
abelian-kernel fibre formula the
only possible stabilizer is
\(H^0(G_{\mathbb Q},K_r/F_5'')\).
This is trivial: the group embeds
in \(F_5/F_5''\), whose first
nonzero lower-central grade is a
torsion-free positive Tate module.
An arithmetic fixed element would
give an invariant in that grade,
which is impossible. Thus the
fibre in (D.37.1) is the actual
\(H^1(G_{\mathbb Q},A_r)\)-torsor,
and the 25 classes remain distinct.
This proves failure of the natural
arithmetic comparison, strengthening
the geometric kernel witness in D.14.
\end{proof}

\begin{statement}{D.38}{Refuted}{Equality of the two Fermat metabelian section sets under the natural comparison}
For no \(r\ge1\) is the natural comparison
from the section set of \(K_r/K_r''\)
to that of \(K_r/F_5''\) a bijection.
The claimed equality in Task D.6 is
therefore false under this comparison.
\end{statement}
\begin{proof}
D.37 exhibits at least 25 distinct
actual section classes over its
single unit tangential class, for
every \(r\). Thus the map is not
injective. This refutes the proposed
natural equality; it does not purport
to classify all classes in either
section set or to decide whether
the comparison is onto for arbitrary
ramification.
\end{proof}

\begin{statement}{D.39}{Proved}{All arithmetic coefficients of the metabelian action}
Retain the generators and Alexander coefficient \(B_g\) of D.35,
and write
\[
 {\cal B}_g(t_1,t_2)
 =1+(t_1-1)(t_2-1)B_g(t_1-1,t_2-1),\qquad
 D_m(X,Y)=(X+Y)^m-X^m-Y^m.
\]
Let \(c_\infty\) be the usual real conjugation, for which
\(\chi_5(c_\infty)=-1\) and \(f_{c_\infty}=1\).
Let \(\mathrm B_m\) denote the Bernoulli numbers, with
\(\mathrm B_2=1/6\).
For odd \(m\ge3\) define the raw Soul\'e character on
\(\Delta=\ker\chi_5\) by
\[
 u_{m,n}=\prod_{\substack{1\le a<5^n\\5\nmid a}}
                   (1-\zeta_{5^n}^{\,a})^{a^{m-1}},\qquad
 \frac{\delta(u_{m,n}^{1/5^n})}{u_{m,n}^{1/5^n}}
             =\zeta_{5^n}^{\,\kappa_m(\delta)}.             \tag{D.39.1}
\]
The roots in this formula have the compatible convention of
\cite[(1.1)--(1.2)]{DE_NakamuraSakugawaWojtkowiak2017}.
Then, for every \(g\in G_{\mathbb Q}\), with \(\chi=\chi_5(g)\),
\[
\begin{split}
 \log {\cal B}_g(e^X,e^Y)
 ={}&
 \sum_{\substack{m\ge2\\m\ {\rm even}}}
 \frac{\mathrm B_m(\chi^m-1)}{2m\,m!}\,D_m(X,Y)\\
 &+\sum_{\substack{m\ge3\\m\ {\rm odd}}}
 \frac{\kappa_m([g,c_\infty])}
      {2(5^{m-1}-1)m!}\,D_m(X,Y).
\end{split}
\tag{D.39.2}
\]
This is an identity of formal power series over \(\mathbb Q_5\).
On substituting \(X=\log(1+u)\), \(Y=\log(1+v)\), exponentiating,
subtracting one and dividing by \(uv\), it gives the actual
integral coefficient \(B_g\in\mathbb Z_5[[u,v]]\).
Thus (D.39.2) specifies the entire arithmetic action in D.35;
no unspecified Frobenius coefficient remains in that formula.
The raw characters on the right are concrete Kummer characters,
not freely assignable parameters.
\end{statement}
\begin{proof}
We first justify the shape of the logarithm, including its
normalization. The normalized adelic beta function in
\cite[\S2.3, p.~101]{PBC_Furusho2003} is defined by
\[
 g(C)=S_\chi(t_1)S_\chi(t_2)\,{\cal B}_g(t_1,t_2)\,C.
\]
This agrees exactly with D.35. Ihara's original
\cite[introduction, p.~1, and \S2, display (6), preprint p.~9]{DE_Ihara2000}
states the beta factorization for every element of
\(\widehat{\rm GT}\), and specifies the coefficient
ring and its faithful evaluation map. In his notation
the gamma element lies in
\(\widetilde A_1=W[[\widehat{\mathbb Z}]]^\times\), and
its function takes values in \(W\otimes\mathbb Q^{\rm ab}\),
not merely in \(W=\prod_pW(\overline{\mathbb F}_p)\).
This distinction matters at five-power torsion.
The factorization therefore gives, at every torsion
character, the identity
\[
 {\cal B}_g(t_1,t_2){\cal B}_g(t_1t_2,t_3)
 ={\cal B}_g(t_2,t_3){\cal B}_g(t_1,t_2t_3).               \tag{D.39.3}
\]
There is no need to choose a formal gamma function.
Immediately after display (6), Ihara explicitly proves
injectivity by the invertibility of the Fourier matrix
\((1\otimes e^{-2\pi i ab/N})_{a,b\in\mathbb Z/N}\)
in \(W\otimes\mathbb Q^{\rm ab}\) for every \(N\).
In particular this includes \(N=5^r\); it is not
restricted to torsion of order prime to five.
Tensoring this matrix argument in three variables
proves (D.39.3) at every finite quotient of the
completed group ring. Project to the coefficient-five,
pro-five factor, whose finite quotients are
\(\mathbb Z_5[(\mathbb Z/5^r)^3]\). Equivalently,
their Fourier maps become injective over
\(\mathbb Q_5(\mu_{5^r})\); integral injectivity follows
from their torsion-freeness. Thus (D.39.3) holds in
the full formal ring, and the logarithmic substitution
used next is justified.

Put \(L(X,Y)=\log{\cal B}_g(e^X,e^Y)\). It vanishes if either
variable is zero. Taking logarithms in (D.39.3), then
differentiating with respect to its third variable at zero,
gives
\[
 \partial_2L(X,Y)=\ell(X+Y)-\ell(Y),
 \qquad \ell(T)=\partial_2L(T,0).
\]
Integration in \(Y\), with \(L(X,0)=0\), gives
\[
               L(X,Y)=\sum_{m\ge2}b_m(g)D_m(X,Y).         \tag{D.39.4}
\]
This elementary calculation also removes the linear
ambiguity in a possible gamma factor.

Next we compute the even coefficients directly from the
third peripheral class. Put \(z=(xy)^{-1}\). There is a
conjugator \(k\) with \(k\,g(z)\,k^{-1}=z^\chi\).
If \(k^{\rm ab}=(a,b)\), comparison in the class-two Lie
quotient gives
\[
 -\frac{\chi^2}{2}+\chi(b-a)=-\frac{\chi}{2},
 \qquad b-a=\frac{\chi-1}{2}.                            \tag{D.39.5}
\]
This comparison is integral at five, since only two is
inverted. We can also check the complete specialization,
not just its leading term. In the Fox module specialize
\(t_1=t,\ t_2=t^{-1}\). Then
\[
 C=\frac{t-1}{t}(e_1+t e_2),\qquad
 v\bigl(g(z)\bigr)
   =-S_\chi(t){\cal B}_g(t,t^{-1})(e_1+t e_2),
 \qquad v(z^\chi)=-\chi(e_1+t e_2).
\]
Here \(v\) denotes the Fox vector. Conjugation by \(k\)
multiplies the first vector by \(t^{a-b}\), since the
abelian part of \(g(z)\) is one after this specialization.
Using (D.39.5) therefore proves the exact reflection formula
\[
 {\cal B}_g(t,t^{-1})
       =\frac{\chi\,t^{(\chi-1)/2}}{S_\chi(t)},\qquad
 {\cal B}_g(e^X,e^{-X})
       =\frac{\chi\sinh(X/2)}{\sinh(\chi X/2)}.            \tag{D.39.6}
\]
Both displayed quotients have constant term one.
The usual generating identity for Bernoulli numbers, or
its formal derivative followed by integration, gives
\[
 \log\frac{\chi\sinh(X/2)}{\sinh(\chi X/2)}
 =\sum_{\substack{m\ge2\\m\ {\rm even}}}
       \frac{\mathrm B_m(1-\chi^m)}{m\,m!}\,X^m.
\]
Since \(D_m(X,-X)=-2X^m\) for even \(m\) and is zero
for odd \(m\), (D.39.4) now gives the first sum of
(D.39.2), including its sign.

For the odd coefficients, the semilinear action on \(C\)
and the identity
\(S_{\chi\chi'}(t)=S_\chi(t)S_{\chi'}(t^\chi)\)
give
\[
 {\cal B}_{gh}(t_1,t_2)
  ={\cal B}_g(t_1,t_2)
       {\cal B}_h(t_1^{\chi_5(g)},t_2^{\chi_5(g)}).
\]
Consequently \(b_m\) is a continuous \(1\)-cocycle with
values in \(\mathbb Q_5(m)\). On \(\Delta\), the explicit
formula \cite[(1.1)--(1.2)]{DE_NakamuraSakugawaWojtkowiak2017}
reads
\[
 b_m(\delta)=\frac{\kappa_m(\delta)}
                   {(5^{m-1}-1)m!}
 \quad(m\ge3\ {\rm odd}).
\]
There is no cyclotomic correction here, since the two
\(S\)-factors are one on \(\Delta\). Since
\(b_m(c_\infty)=0\), the cocycle identity gives
\[
 b_m([g,c_\infty])
       =(1-(-1)^m)b_m(g)=2b_m(g)
                  \quad(m\ {\rm odd}).
\]
The commutator lies in \(\Delta\). This proves the
second sum and fixes its extension off \(\Delta\)
without choosing a lift of a cyclotomic generator.

Finally the formal transformations used above are
invertible over \(\mathbb Q_5\). They were applied to
the actual integral beta series, so their inverse gives
that same integral series. In particular the assertion
of integrality does not claim that arbitrary choices
of the numbers in (D.39.1) give integral coefficients.
For example the even weight-two coefficient is
\((\chi^2-1)/24\), and the coefficient of \(X^2Y\)
in weight three is
\(\kappa_3([g,c_\infty])/96\), exactly as in D.31.
\end{proof}

\begin{statement}{D.40}{Proved}{All positive Tate layers and explicit generators}
For every integer \(n\ge2\), put
\[
 C=c_0(\{q:q\ne5\},\mathbb Z_5),\qquad
 (D_na)_q=(q^{n-1}-1)a_q,\qquad
 \epsilon_n=1+v_5(n).
\]
Then
\[
\begin{split}
 H^0(G_{\mathbb Q},\mathbb Z_5(n))&=0,\\
 H^1(G_{\mathbb Q},\mathbb Z_5(n))&=
 \begin{cases}
  \mathbb Z_5,&n\ge3\text{ odd},\\
  \mathbb Z/5^{\epsilon_n},&n\equiv0\pmod4,\\
  0,&n\equiv2\pmod4,
 \end{cases}\\
 H^2(G_{\mathbb Q},\mathbb Z_5(n))&=C/D_nC.
\end{split}
\tag{D.40.1}
\]
The last cokernel is algebraic. Its kernel of passage
to all finite-coefficient second cohomology groups is
\[
 \frac{C\cap\prod_{q\ne5}(q^{n-1}-1)\mathbb Z_5}{D_nC}.
                                                               \tag{D.40.2}
\]
Every integral first-cohomology class is unramified outside
five. In the odd case an explicit generator is the cocycle
\[
          \widehat\kappa_n(g)
                   =\frac12\kappa_n([g,c_\infty]).        \tag{D.40.3}
\]
In the case \(n\equiv0\pmod4\), a generator is
\[
           z_n(g)=\frac{\chi_5(g)^n-1}{5^{\epsilon_n}}.
                                                               \tag{D.40.4}
\]
\end{statement}
\begin{proof}
The residual calculation in D.1 was made over
\(L=\mathbb Q(\zeta_5)\) before specifying the weight.
It gives \(H^2(G_5,\mathbb F_5(n))=0\) in every weight.
The dimensions of \(H^1(G_5,\mathbb F_5(n))\) for
\(n\equiv0,1,2,3\pmod4\) are respectively \(1,1,0,1\):
these are the invariant parts of
\((\mathbf1\oplus\omega\oplus\omega^2)\otimes\omega^{n-1}\).
The same finite-coefficient and Mittag--Leffler argument
as in D.1 gives
\[
 H^2(G_5,\mathbb Z_5(n))=0.
\]
It also makes the integral \(H^1\) a cyclic compact
\(\mathbb Z_5\)-module, or zero in residual weight two.
Its torsion is
\(H^0(G_{\mathbb Q},\mathbb Q_5/\mathbb Z_5(n))\).
The full cyclotomic image \(\mathbb Z_5^\times\) makes
this group zero unless \(4\mid n\); in that case the
lifting-the-exponent formula gives its order
\(5^{1+v_5(n)}\). A nonzero cyclic \(\mathbb Z_5\)-module
with this nonzero torsion must be exactly the indicated
finite cyclic group. For odd \(n\), residual dimension
one and absence of torsion give \(\mathbb Z_5\).

The localization proof of D.1 applies verbatim for
every \(n\ge2\), with residue complex
\([C\xrightarrow{D_n}C][-1]\).
Its first differential is injective because the nonzero
integer \(q^{n-1}-1\) acts injectively on \(\mathbb Z_5\).
It gives both the first-cohomology unramification and
the algebraic second-cohomology cokernel, including
(D.40.2). Thus no separated quotient is taken in
passing to these additional weights.

For (D.40.3), D.39 shows that
\(\kappa_n([g,c_\infty])/2\) is an integral continuous
\(1\)-cocycle: it is the coefficient \(b_n(g)\) there
multiplied by the fixed scalar
\((5^{n-1}-1)n!\). Its restriction to \(\Delta\) is
\(\kappa_n\), since conjugation by \(c_\infty\) acts
by \((-1)^n=-1\) on the character.
We verify that it is primitive; a nonzero rational
character alone would not suffice.

At level one the exponents \(a^{n-1}\) modulo five
depend only on \(n-1\) modulo four. If \(n\equiv1\pmod4\),
the Kummer element \(u_{n,1}\) has the same class as
\(\prod_{a=1}^4(1-\zeta_5^a)=5\) in \(L^\times/L^{\times5}\).
If \(n\equiv3\pmod4\), it has the same class as \(u_{3,1}\),
which D.35 computes exactly as
\[
 u_{3,1}=(1-\zeta_5)^{30}\varphi^{13}
         =\varphi^3\bigl((1-\zeta_5)^6\varphi^2\bigr)^5,
 \qquad \varphi=(1+\sqrt5)/2.
\]
Neither class is in the line generated by \(\zeta_5\).
For \(5\), the valuation at \(1-\zeta_5\) is four,
not divisible by five. For \(\varphi^3\), the free
unit generator has exponent three, whereas \(\zeta_5\)
is torsion. The field \(L(\mu_{5^\infty})\) has exactly
one degree-five extension of \(L\), namely
\(L(\zeta_{25})\), whose Kummer line is generated
by \(\zeta_5\). Thus neither of the two classes
becomes a fifth power in that cyclotomic tower.
The character \(\kappa_n\) is consequently nonzero
modulo five on \(\Delta\). This proves primitivity
of (D.40.3) in the already computed rank-one group.

Finally \(5^{-\epsilon_n}\bmod\mathbb Z_5\) is a
generator of the invariant torsion group when \(4\mid n\).
Its connecting cocycle is precisely (D.40.4), so that
cocycle generates the finite first-cohomology group,
with the stated exact order.
\end{proof}

\begin{statement}{D.41}{Proved}{An explicit all-weight metabelian lifting criterion}
Put \(M=F_5/F_5''\), \(A=M'=\Lambda C\),
\(\Lambda=\mathbb Z_5[[u,v]]\), \(t_1=1+u,t_2=1+v\),
and \(I=(u,v)\). Every element of \(M\) has a unique
continuous normal form
\[
                    C^\lambda x^a y^b,
       \qquad (a,b)\in\mathbb Z_5^2,\quad\lambda\in\Lambda.
\]
Here \(C^\lambda\) means the Alexander-module element
\(\lambda C\), not a scalar power when \(\lambda\notin\mathbb Z_5\).
Writing such an element as \((a,b,\lambda)\), its group law is
\[
 (a,b,\lambda)(c,d,\eta)=
 \bigl(a+c,b+d,\lambda+t_1^at_2^b\eta
                   -t_1^aS_b(t_2)S_c(t_1)\bigr).         \tag{D.41.1}
\]
For \(g\in G_{\mathbb Q}\), put \(\chi=\chi_5(g)\),
\(\phi_\chi(t_i)=t_i^\chi\), and
\[
 U_g=S_\chi(t_1)S_\chi(t_2){\cal B}_g(t_1,t_2).
\]
With \(B_g\) given explicitly by D.39, its arithmetic action is
\[
 g(a,b,\lambda)=
 \bigl(\chi a,\chi b,\,
 U_g\phi_\chi(\lambda)
       +t_1^{\chi a}(t_2^{\chi b}-1)B_g\bigr).            \tag{D.41.2}
\]
Thus an abelian cocycle \((a_g,b_g)\) has a metabelian
lift precisely when the following integral cochain equation
has a continuous solution \(\lambda:G_{\mathbb Q}\to\Lambda\):
\[
\begin{split}
 \lambda_{gh}={}&\lambda_g+
       t_1^{a_g}t_2^{b_g}U_g\phi_{\chi_5(g)}(\lambda_h)\\
 &+t_1^{a_g}t_2^{b_g}t_1^{\chi_5(g)a_h}
              (t_2^{\chi_5(g)b_h}-1)B_g\\
 &-t_1^{a_g}S_{b_g}(t_2)S_{\chi_5(g)a_h}(t_1).
\end{split}
\tag{D.41.3}
\]
The solvability test for (D.41.3) is the following explicit
successive residue test, including its convergence condition.

Let \(M_n=M/I^{n-1}A\), with \(M_1=M/A\), and suppose an
actual section of \(M_{n-1}\), \(n\ge2\), has already been
chosen. At \(q\ne5\), write its tame inertia and Frobenius
normal forms as
\[
 \tau_q=(\alpha,\beta,c),\qquad
 \sigma_q=(a,b,d),\qquad
 \mu=t_1^\alpha t_2^\beta,\quad \nu=t_1^at_2^b.
\]
Choose for \(c,d\) their polynomial representatives of
total degree \(<n-2\), taking both to be zero when \(n=2\).
Use the actual arithmetic \(B_{\sigma_q}\) of D.39 and put
\[
\begin{split}
 P_q(\alpha,\beta)
  &=-\sum_{j=0}^{q-1}
           t_1^{j\alpha}S_{j\beta}(t_2)S_\alpha(t_1),\\
 K_q(a,b;\alpha,\beta)
  &=t_1^{q\alpha}S_{q\beta}(t_2)S_a(t_1)
                 -t_1^aS_b(t_2)S_{q\alpha}(t_1),\\
 {\cal R}_q
  &=\nu U_{\sigma_q}\phi_q(c)-S_q(\mu)c
          +(1-\mu^q)d\\
  &\hspace{8mm}
       +\nu t_1^{q\alpha}(t_2^{q\beta}-1)B_{\sigma_q}
           +K_q(a,b;\alpha,\beta)-P_q(\alpha,\beta).
\end{split}
\tag{D.41.4}
\]
For \(0\le i\le n-2\), define
\[
 R_{n,q,i}=q^{-1}
              [u^iv^{\,n-2-i}]\,{\cal R}_q.              \tag{D.41.5}
\]
Each family \(R_{n,\bullet,i}\) is in \(C\).
The chosen section lifts to \(M_n\) if and only if
\[
\begin{gathered}
 R_{n,q,i}\in(q^{n-1}-1)\mathbb Z_5
                         \quad\hbox{for every }q\ne5,\ i,\\
 \left(\frac{R_{n,q,i}}{q^{n-1}-1}\right)_{q\ne5}
                  \in c_0(\{q:q\ne5\},\mathbb Z_5)
                         \quad\hbox{for every }i.
\end{gathered}
\tag{D.41.6}
\]
When it lifts, its fibre on section classes is a torsor under
\[
 \begin{cases}
    \mathbb Z_5^{\,n-1},&n\ge3\text{ odd},\\
    (\mathbb Z/5^{1+v_5(n)})^{\,n-1},&n\equiv0\pmod4,\\
    0,&n\equiv2\pmod4.
 \end{cases}                                           \tag{D.41.7}
\]
The translations can be taken to be the explicit cocycles
of D.40 times \(u^iv^{n-2-i}C\).

Finally, \(\operatorname{Sec}(M)\) is exactly the set of
infinite compatible branches obtained by starting with
an abelian section and applying (D.41.4)--(D.41.7)
at every weight. No uniform extra condition across the
weights is required; the \(c_0\) condition at each weight
is required. This specifies an explicit obstruction
system with fixed arithmetic coefficients, rather than
leaving the finite-stage obstructions unspecified.
D.68 further proves that, after fixing the
abelian class or any lower-stage class, individual
nonemptiness at every higher stage suffices for
a compatible branch: the actual finite-class
fibres are compact. This conclusion retains
each weight's unrestricted convergence condition.
It does not replace the unrestricted Alexander
cohomology calculation still sought in Task D.5
by an asserted closed formula for its answer.
\end{statement}
\begin{proof}
The normal form follows from the exact Alexander
sequence in D.35 and the continuous set-theoretic
lift \((a,b)\mapsto x^ay^b\) of the abelianization.
It is unique because \(A=\Lambda C\).
In the Fox embedding its vector is
\[
 \lambda C+S_a(t_1)e_1+t_1^aS_b(t_2)e_2,\qquad
 C=(1-t_2)e_1+(t_1-1)e_2.
\]
Multiplication of Fox vectors is
\((T,v)(T',v')=(TT',v+Tv')\).
Subtracting the vector of \(x^{a+c}y^{b+d}\) gives
exactly \(-t_1^aS_b(t_2)S_c(t_1)C\); this proves
(D.41.1). It is valid for all five-adic exponents
by continuity, not only for integers.
The identities
\[
 g(x^a)=x^{\chi a},\qquad
 g(y^b)=f_g^{-1}y^{\chi b}f_g,\qquad
 g(\lambda C)=U_g\phi_\chi(\lambda)C
\]
then prove (D.41.2), again with integral coefficients.
Combining these two formulas with the cocycle identity
proves (D.41.3). Its constant term is the usual
class-two cup-product equation.

We verify the entire tame expression before taking
its leading coefficient. For \(J=(\alpha,\beta,c)\),
repeated use of (D.41.1) gives
\[
 J^q=(q\alpha,q\beta,S_q(\mu)c+P_q(\alpha,\beta)).
\]
For \(F=(a,b,d)\) and \(J'=(q\alpha,q\beta,c')\),
the same law gives
\[
 FJ'F^{-1}=
 (q\alpha,q\beta,(1-\mu^q)d+\nu c'
                         +K_q(a,b;\alpha,\beta)).
\]
Substitute
\(c'=U_{\sigma_q}\phi_q(c)
             +t_1^{q\alpha}(t_2^{q\beta}-1)B_{\sigma_q}\).
Two elements with the same abelian normal form have
quotient in \(A\) equal to the difference of their
Alexander coefficients. Hence their difference is
exactly \({\cal R}_q C\):
\[
 F\,\sigma_q(J)\,F^{-1}J^{-q}
                                  =C^{{\cal R}_q}.      \tag{D.41.8}
\]
The arithmetic action is unramified at \(q\ne5\);
the tame relation is thus the appropriate actual
local Galois relation. Prime-to-five inertia has
trivial geometric image and contributes nothing.

For the metabelian free group,
\[
 \gamma_{n}M=I^{n-2}A\quad(n\ge2),\qquad
 \gamma_nM/\gamma_{n+1}M
       =\bigoplus_{i=0}^{n-2}
                \mathbb Z_5(n)\,u^iv^{n-2-i}C.          \tag{D.41.9}
\]
The first identity follows by commuting \(A\)
with \(x\) and \(y\), which multiplies by \(u\)
and \(v\), and then by induction. The second
follows from (D.41.2), whose action on that grade
is \(\chi_5^n\).
Because the input was an actual section of
\(M_{n-1}\), all lower-degree terms in
(D.41.8) vanish. Its degree \(n-2\) is therefore
exactly the central extension's relator defect.

The coefficient representatives used here give a
continuous integral set lift. If an inertia normal
form is zero modulo \(5^e\), so is its chosen
lift, and (D.41.8) is zero modulo \(5^e\).
Continuity of the lower section consequently
makes the families (D.41.5) tend to zero.
The normalization by \(q^{-1}\) is the same
normalization checked in D.29 and D.34.
Changing a highest-weight inertia coordinate by
\(z_i\) changes \(R_{n,q,i}\) by
\((q^{n-1}-1)z_i\): in (D.41.4) the two
leading actions are \(q^n\) and \(q\).
Thus the displayed residues represent the
obstruction in the actual algebraic cokernel
\((C/D_nC)^{n-1}\).

The central lifting theorem D.29 applies to
(D.41.9). Its proof uses the zero second
cohomology over \(G_5\), which holds in every
weight by D.40, and the explicit continuous
residue complex; it therefore applies without
a fixed bound on the class. It proves both
necessity and sufficiency of (D.41.6),
including its divided-family convergence.
The fibre is \(H^1(G_{\mathbb Q},\mathbb Z_5(n))^{n-1}\).
There is no stabilizer quotient: a geometric
fixed element of any of these twisted nilpotent
groups would have a first nonzero positive
Tate grade, which has no invariants. D.40 now
gives (D.41.7) and its stated generators.
The identities (D.41.1), (D.41.2), and the
complete tame expression, with negative as
well as positive exponents, are independently
checked by the executed exact-arithmetic script
\texttt{de\_metabelian.py}.

For completeness, compatibility in the last
assertion is essential. Given a compatible
sequence of section classes, choose representatives
successively. At a new level, lift the geometric
conjugator relating its projection to the previous
representative and conjugate it; the representatives
are then exactly compatible. Their inverse limit
is a continuous section, since the topology of
\(\Lambda\) is the product topology of its
\(\mathbb Z_5\)-coefficients. Conversely every
section gives such a branch. If two sections
give the same branch of classes, the sets of
elements of the compact group \(M\) conjugating
them through level \(n\) are nested nonempty
closed sets. Their intersection gives an
actual conjugator. This proves the asserted
bijection. It does not infer the existence
of a branch from unrelated nonempty sets
of lifts at arbitrarily high levels.
\end{proof}

\begin{statement}{D.42}{Proved}{Every ramified local metabelian section is cuspidal}
Let \(k/\mathbb Q_q\) be finite, \(q\ne5\), with residue
field of cardinality \(Q\), and retain the actual unit-tangent
action on \(M=F_5/F_5''\). Every continuous section
\[
             G_k\longrightarrow M\rtimes G_k
\]
either is unramified on its geometric coordinate, or belongs,
up to \(M\)-conjugation, to a cusp packet.
More precisely, if a five-primary tame inertia generator
has image \(I\ne1\), then \(I\) is conjugate to a nonzero
five-adic power of one of the three peripheral generators,
and the whole local section lies in that cusp decomposition
group.

In particular, for any global metabelian section and any
prime \(q\ne5\), its abelian inertia vector \((\alpha_q,\beta_q)\)
satisfies
\[
       \alpha_q\beta_q(\alpha_q-\beta_q)=0.                \tag{D.42.1}
\]
If this vector is zero, its entire inertia image is zero.
This is an unconditional statement about the actual
metabelian quotient; it does not use the supplied local
section conjecture for full fundamental groups.
\end{statement}
\begin{proof}
We use two integral facts, which we spell out because
rationalizing at the augmentation alone would miss
the nonidentity five-power characters.

First consider a finite character
\(\lambda:(\mathbb Z_5)^2\to\mu_{5^\infty}\), with
\(\lambda(t_1)=\zeta_1,\lambda(t_2)=\zeta_2\), such that
\[
             \zeta_1\ne1,\quad\zeta_2\ne1,\quad
                                      \zeta_1\zeta_2\ne1. \tag{D.42.2}
\]
Choose \(N=5^r\) containing both character orders.
The subgroup \(K_N=\ker(F_5\to(\mathbb Z/N)^2)\)
is the pro-five geometric group of the open Fermat
cover \(X^N+Y^N=Z^N\).
Since \(F_5'\subset K_N\), we have
\(F_5''\subset K_N'\). Thus an element of \(M\)
whose abelian image is divisible by \(N\)
has a well-defined image in the actual
homology \(K_N/K_N'\); this is the comparison
used below for \(I^N\).
Its Fox chain complex,
after taking the nontrivial character \(\lambda\),
has differential
\[
 (z_1,z_2)\longmapsto
       (\zeta_1-1)z_1+(\zeta_2-1)z_2.
\]
The kernel is the line generated by
\((1-\zeta_2,\zeta_1-1)\), the specialization of \(C\).
The boundary part of the homology of this open curve
has characters only where one of
\(\zeta_1,\zeta_2,\zeta_1\zeta_2\) is one.
Indeed its three permutation modules are induced
from the stabilizers generated respectively by
\((1,0),(0,1),(1,1)\). Thus the line under (D.42.2)
is a character component of the Tate module of the
proper Fermat Jacobian.

The proper curve is smooth over the residue field
because \(q\nmid N\). If \(h\) is such that
\(Q^h\equiv1\pmod N\), arithmetic Frobenius to the
power \(h\) preserves this character line, and its
eigenvalue has every complex absolute value
\(Q^{h/2}\). We use the arithmetic Frobenius on the
covariant Tate module here: it is the endomorphism
raising coordinates to their \(Q^h\)-th powers.
The precise eigenvalue assertion and this Tate-module
identification are
\cite[Theorem 14.1(b), and (9.20)]{DE_MilneAV}.
Multiplying Frobenius by a deck transformation
only multiplies the eigenvalue by a root of unity.
Consequently it cannot acquire eigenvalue \(Q^h\).

The second fact is an elementary bounded-measure
lemma. For \(b\in\mathbb Z_5\), the solutions
\(f\in\mathbb Z_5[[t-1]]=\mathbb Z_5[[\mathbb Z_5]]\)
of
\[
                       f(t)=t^b f(t^Q)                  \tag{D.42.3}
\]
are exactly the multiples of \(t^{z_0}\) if
\(z_0=-b/(Q-1)\) belongs to \(\mathbb Z_5\), and
are zero otherwise. To prove this, interpret \(f\)
as an integral measure on \(\mathbb Z_5\).
Equation (D.42.3) says that the measure is invariant
under \(T(z)=Qz+b\). Take a residue ball \(U\)
not containing the possible fixed point \(z_0\),
small enough that \(v_5(z-z_0)\) is constant on it.
Some positive power \(T^d\) preserves \(U\);
we may further make \(Q^d\equiv1\pmod5\).
The integer \(Q^d-1\) is nonzero.
On refinements of \(U\) modulo \(5^m\), the
orbits of \(T^d\) have lengths tending through
arbitrarily large powers of five, by
\[
 v_5(Q^{d5^j}-1)=v_5(Q^d-1)+j.
\]
Invariance makes the measure constant on every
such orbit of refined balls. Hence its value
on \(U\) is divisible by arbitrarily large
powers of five, so is zero. This argument
also applies when \(z_0\notin\mathbb Z_5\).
The measure is therefore supported on the
single fixed point, if it exists. Such a
measure is precisely a scalar Dirac measure.
In particular no integral Haar measure on
an infinite orbit survives when \(Q\equiv1\pmod5\).

We now apply these facts to the section. The
arithmetic action on \(M\) is unramified at
\(q\ne5\). Wild inertia and the prime-to-five
tame factors therefore have zero geometric
image. Let \(I=(\alpha,\beta,c)\) be the image
of a five-primary tame generator, and let
\(F\) be the geometric coordinate at an
arithmetic Frobenius \(\sigma\). They satisfy
\[
                      F\,\sigma(I)\,F^{-1}=I^Q.          \tag{D.42.4}
\]
If \((\alpha,\beta)=0\), the first nonzero
\(I\)-adic grade of a nonzero \(c\in\Lambda C\)
would have Frobenius eigenvalue \(Q^{j+2}\)
and also \(Q\), as in D.36. Since that grade
is torsion-free, this is impossible. Thus
zero abelian inertia gives \(I=1\).

Suppose now that \((\alpha,\beta)\ne0\), and put
\(\mu=t_1^\alpha t_2^\beta\).
For every character satisfying \(\lambda(\mu)=1\)
and (D.42.2), the Fox vector \(v(I)\) vanishes
at that character. To see this rigorously,
\(I^N\in K_N\), and its specialized Fox vector
is \(N v(I)\), since \(\lambda(\mu)=1\).
Iterating (D.42.4) through the above power \(h\)
would make that vector an eigenvector with
eigenvalue \(Q^h\) of Frobenius times a deck
transformation. The first fact proves it is
zero; \(N\) is nonzero in \(\mathbb Q_5\), so
\(v(I)\) is zero there as well.

This vanishing forces a peripheral direction,
including when both entries are divisible
by five. Write
\[
 (\alpha,\beta)=5^e(\alpha_0,\beta_0),
 \qquad(\alpha_0,\beta_0)\text{ primitive}.
\]
Complete \((\alpha_0,\beta_0)\) to a
\(\operatorname{GL}_2(\mathbb Z_5)\)-matrix,
and use the corresponding multiplicative
coordinates \(r=t_1^{\alpha_0}t_2^{\beta_0},s\).
Then
\[
 R:=\Lambda/(\mu-1)
   =\mathbb Z_5[[r-1,s-1]]/(r^{5^e}-1).
\]
This ring is torsion-free and, after adjoining
the \(5^e\)-th roots of unity, injects into
the product of the rings
\(\mathcal O[[s-1]]\) given by \(r=\zeta\),
\(\zeta^{5^e}=1\). Injectivity follows by
writing representatives of degree \(<5^e\)
in \(r\) and using the nonzero Vandermonde
determinant over the fraction field.

In the component \(r=1\), none of
\(t_1,t_2,t_1t_2\) is identically one unless
\(\beta=0,\alpha=0,\alpha=\beta\), respectively.
If none of these equalities holds, only
finitely many five-power values of \(s\)
lie on the three boundary loci. Indeed a
nonzero exponent in \(\mathbb Z_5\) has a
finite kernel on \(\mu_{5^\infty}\).
Both coordinates of \(v(I)\) thus vanish
at infinitely many other five-power values.
A nonzero element of \(\mathcal O[[s-1]]\)
has only finitely many such zeros:
divide by the smallest coefficient valuation
and apply Weierstrass preparation.
Therefore the whole vector vanishes in this
component. Its augmentation would give
\((\alpha,\beta)=0\), a contradiction.
This proves (D.42.1). The identity component
was taken only after retaining all
\(5^e\) components of \(R\); no rational
augmentation completion was substituted
for the integral quotient.

Permuting the punctures by an automorphism
of the tripod over \(\mathbb Q\), and changing
the rational tangential base point along it,
preserves the assertion about cusp decomposition
groups and its tame relation. We may therefore
treat the direction \(\beta=0,\alpha\ne0\).
Use the unit tangent at this cusp, so that
\(\sigma(x)=x^Q\). Write \(F=(a,b,d)\).
On the component \(t_1=1,\ t_2=t\), put
\[
                     f(t)=\alpha-(t-1)c(0,t-1).
\]
The Fox vector of \(I\) is \(f(t)e_1\).
Here (D.35)'s actual multiplier satisfies
\(U_\sigma(1,t)=Q S_Q(t)\), because
\({\cal B}_\sigma(1,t)=1\). Consequently
(D.42.4) gives exactly (D.42.3).
The measure lemma and \(f(1)=\alpha\ne0\)
give an element \(z_0\in\mathbb Z_5\) with
\[
                          f(t)=\alpha t^{z_0}.          \tag{D.42.5}
\]

The element \(I_0=y^{z_0}x^\alpha y^{-z_0}\)
has Alexander coefficient
\[
                   c_0=-S_{z_0}(t_2)S_\alpha(t_1).
\]
Equation (D.42.5) says \(c=c_0\) on the
component \(t_1=1\). On every nonidentity
component \(t_1=\zeta,\ \zeta^\alpha=1\),
the first fact and Weierstrass argument
give \(c=0\): outside the two exceptional
values \(t_2=1,\zeta^{-1}\), the vector
\(c(1-t_2,\zeta-1)\) is zero, and
\(\zeta-1\ne0\). The coefficient \(c_0\)
also vanishes there because \(S_\alpha(\zeta)=0\).
The integral injection of \(R\) into its
components now gives
\[
                          c-c_0\in(\mu-1)\Lambda.
\]
Conjugation by \(C^m\) changes an Alexander
coefficient by \((1-\mu)m\). Hence \(I\)
is conjugate to \(I_0\), and then to
\(x^\alpha\), inside \(M\) itself.

Finally the centralizer of \(x^a\) in
\(M\), for every nonzero \(a\in\mathbb Z_5\),
is exactly \(x^{\mathbb Z_5}\).
For \(H=(h,k,m)\), the equation
\(Hx^aH^{-1}=x^a\), computed by (D.41.1),
is
\[
 S_a(t_1)\bigl((t_1-1)m+t_1^hS_k(t_2)\bigr)=0.
\]
The ring \(\Lambda\) is a domain and
\(S_a(t_1)\ne0\). Setting \(t_1=1\)
first gives \(S_k(t_2)=0\), hence \(k=0\).
It then gives \((t_1-1)m=0\), hence \(m=0\).
After the preceding conjugation of inertia,
(D.42.4) therefore forces \(F\) into
\(x^{\mathbb Z_5}\). Inertia and Frobenius
topologically generate the relevant tame
quotient; the discarded inertia factors
have zero image. The full local section
lies in the cusp decomposition group.
This completes the proof.
\end{proof}

\begin{statement}{D.43}{Proved}{Local Alexander cohomology with nonzero abelian inertia}
Let \(k,Q\) be as in D.42, and let \(s=(a_g,b_g)\)
be any abelian cocycle, whether or not it has a
metabelian lift. Give \(A_s=\Lambda C\) the action
\(\rho_s\) of D.35. If its inertia vector
\((\alpha,\beta)\) is nonzero, put
\[
 \mu=t_1^\alpha t_2^\beta,\quad
 R=\Lambda/(\mu-1),\quad
 T=Q^{-1}\rho_s(\sigma)\ \text{on }R,
\]
where \(\sigma\) is arithmetic Frobenius. Then
\[
 H^0(G_k,A_s)=H^1(G_k,A_s)=0,\qquad
 H^2(G_k,A_s)=R/(T-1)R.                                \tag{D.43.1}
\]
Here the operator is the actual explicit operator
of D.39, and the last image is closed.

There is a complete evaluation in the primitive
axis case \(\beta=0,\alpha\in\mathbb Z_5^\times\).
Write \(b=b_\sigma\) and \(s_Q=v_5(Q-1)\). Then
\[
 H^2(G_k,A_s)\cong
 \begin{cases}
    \displaystyle\prod_{\mathbb N}\mathbb Z_5,
                        &b\in(Q-1)\mathbb Z_5,\\
    \mathbb Z_5^{\,5^{v_5(b)}},
                        &v_5(b)<s_Q.
 \end{cases}                                           \tag{D.43.2}
\]
The second case can occur only when \(Q\equiv1\pmod5\).
The identifications are noncanonical isomorphisms
of compact \(\mathbb Z_5\)-modules.
\end{statement}
\begin{proof}
The action of a five-primary inertia generator
on \(A_s\) is multiplication by \(\mu\).
Since \(\Lambda\) is a domain and \(\mu\ne1\),
its invariants are zero, and its first
cohomology is \(R\). Its higher cohomology
vanishes. These compact-coefficient assertions
can be obtained from the procyclic resolution
at every finite coefficient quotient:
all the local finite-coefficient groups are
finite, so their inverse systems are
Mittag--Leffler. In particular no local
derived-limit term is being discarded.

Frobenius acts on this inertia cohomology by
\(Q^{-1}\rho_s(\sigma)\).
Indeed \(\sigma^{-1}\tau\sigma=\tau^{1/Q}\);
the crossed-homomorphism power formula gives
\(S_{1/Q}(\mu)\), which becomes
\(S_Q(\mu)^{-1}\) after applying \(\rho_s(\sigma)\).
Modulo \(\mu-1\) its value is \(Q^{-1}\).
Hochschild--Serre for inertia and the
procyclic residue-field group therefore gives
\[
 H^1(G_k,A_s)=R^{T=1},\qquad
 H^2(G_k,A_s)=R/(T-1)R.
\]

To show \(R^{T=1}=0\), evaluate an invariant
at a finite character of \(R\) off the three
boundary loci. Iterating Frobenius to fix
that character would give eigenvalue \(Q^h\)
on a proper Fermat Tate-module character,
up to a root of unity from the abelian
Kummer twist. D.42's weight argument
therefore makes the evaluation zero.
The same componentwise Weierstrass argument
makes the element zero on every component
not contained in a boundary locus.
If a component is contained in a boundary
locus, the inertia direction is a cusp
direction. After permuting the punctures
it is the component \(t_1=1\) of
\(\mu=t_1^\alpha\), and on it
\[
                Tm(t)=t^bS_Q(t)m(t^Q).
\]
Multiplication by \(t-1\) turns \(Tm=m\)
into the invariant-measure equation of
D.42 for \(f=(t-1)m\). This measure has
augmentation zero. An invariant integral
measure is supported at a fixed point,
so augmentation zero forces it to vanish.
Thus \(m=0\) on this component as well,
proving (D.43.1). The ring \(R\) is compact,
and \(T-1\) is continuous; its image is
compact and hence closed.

For (D.43.2), primitivity of \(\alpha\)
gives \(R=\mathbb Z_5[[t-1]]\).
Multiplication by \(t-1\) identifies this
module, with its \(T\)-action, with the
augmentation ideal \(J\) in the measure
module \(V=\mathbb Z_5[[\mathbb Z_5]]\),
where \(T\) acts by the affine permutation
\(z\mapsto Qz+b\).
Let \(\mathcal O\) be the compact orbit
space under the closure of this permutation.
Then
\[
                     V/(T-1)V=\mathbb Z_5[[\mathcal O]].
                                                               \tag{D.43.3}
\]
This follows either from its universal
property for continuous maps into profinite
\(\mathbb Z_5\)-modules, or from the finite
permutation quotients followed by inverse
limit. The image of \(T-1\) is closed,
so no extra closure is inserted.

If \(z_0=-b/(Q-1)\) is integral, translation
by \(z_0\) identifies the action with
multiplication by \(Q\). Its orbit space
has the fixed point zero and, at each
valuation \(j\ge0\), finitely many nonzero
orbits indexed by
\(\mathbb Z_5^\times/\overline{\langle Q\rangle}\).
This quotient is finite because the
closure of \(\langle Q\rangle\) is open
in \(\mathbb Z_5^\times\). These isolated
orbits tend to zero as \(j\to\infty\).
Consequently the augmentation kernel of
\(\mathbb Z_5[[\mathcal O]]\) is the
countable product of \(\mathbb Z_5\)'s,
with convergent generators
\(\delta_o-\delta_0\).
The fixed Dirac measure splits
\(V\to\mathbb Z_5\) equivariantly, so
this kernel is exactly \(J/(T-1)J\).

If \(v_5(b)=t<s_Q\), there is no integral
fixed point. Every orbit is dense in one
coset of \(5^t\mathbb Z_5\). In fact
\[
 v_5(T^{5^j}(z)-z)
   =v_5\bigl(S_{5^j}(Q)((Q-1)z+b)\bigr)=t+j
\]
for all \(z\in\mathbb Z_5\).
Thus \(\mathcal O\) consists of exactly
\(5^t\) points and the right side of
(D.43.3) is \(\mathbb Z_5^{5^t}\).
There are no invariant measures in \(V\),
by D.42's lemma. Applying invariants and
coinvariants to
\(0\to J\to V\to\mathbb Z_5\to0\)
therefore gives an exact sequence
\[
 0\longrightarrow\mathbb Z_5
  \longrightarrow J/(T-1)J
  \longrightarrow
       \ker(\mathbb Z_5^{5^t}\xrightarrow{\rm sum}\mathbb Z_5)
  \longrightarrow0.
\]
It splits as compact \(\mathbb Z_5\)-modules,
since the right module is finite free.
Its middle term is \(\mathbb Z_5^{5^t}\),
as asserted.
\end{proof}

\begin{statement}{D.44}{Proved}{Unrestricted global Alexander cohomology over every unramified abelian pair}
Let \(s=(A\kappa(5),B\kappa(5))\), with
\(A,B\in\mathbb Z_5\), and let \(A_s=\Lambda C\)
have its actual action from D.35.
Let \(\mathfrak m=(5,u,v)\) and put
\[
 {\cal C}_\Lambda=
 \left\{(z_q)_{q\ne5}\in\prod_{q\ne5}\Lambda:
  \text{for every }r,\ z_q\in\mathfrak m^r
                     \text{ for all but finitely many }q\right\}.
\]
Give this module the inverse-limit topology of the
discrete direct sums
\(\bigoplus_{q\ne5}\Lambda/\mathfrak m^r\).
This is the uniform \(\mathfrak m\)-adic \(c_0\)
topology, not the subspace topology of the product.
Choose arithmetic Frobenii \(\sigma_q\), and define
the integral operators
\[
 D_{s,q}m=q^{-1}t_1^{a_{\sigma_q}}t_2^{b_{\sigma_q}}
       U_{\sigma_q}\phi_q(m)-m,\qquad
 D_s(z)_q=D_{s,q}z_q,                                   \tag{D.44.1}
\]
where every coefficient of \(U_{\sigma_q}\) is fixed
by D.39. The unrestricted cohomology is
\[
 H^0(G_{\mathbb Q},A_s)=0,\qquad
 H^1(G_{\mathbb Q},A_s)\cong\prod_{\mathbb N}\mathbb Z_5,
 \qquad
 H^2(G_{\mathbb Q},A_s)
              ={\cal C}_\Lambda/D_s{\cal C}_\Lambda.     \tag{D.44.2}
\]
The last cokernel is algebraic. More precisely,
with \(X=\operatorname{Spec}\mathbb Z[1/5]\),
there is a localization triangle
\[
 R\Gamma(X,A_s)\longrightarrow
 R\Gamma(\mathbb Q,A_s)\longrightarrow
 [\,{\cal C}_\Lambda\xrightarrow{D_s}{\cal C}_\Lambda\,][-1]
 \longrightarrow R\Gamma(X,A_s)[1].                     \tag{D.44.3}
\]
For the coefficient tower \(A_s/\mathfrak m^rA_s\),
the Jannsen kernel is
\[
 \varprojlim\nolimits^1_r
       H^1(G_{\mathbb Q},A_s/\mathfrak m^rA_s)
 =
 \frac{{\cal C}_\Lambda\cap\prod_{q\ne5}D_{s,q}\Lambda}
      {D_s{\cal C}_\Lambda}.                            \tag{D.44.4}
\]
It is nonzero for every pair \(A,B\). In particular the
unrestricted second cohomology in (D.44.2) is
non-Hausdorff with the displayed residue-quotient
topology, even though its base cohomology over
\(X\) has \(H^2=0\).
\end{statement}
\begin{proof}
The finite modules \(A_s/\mathfrak m^rA_s\) are
unramified outside five. Their finite filtrations
have successive quotients of the form
\(\mathbb F_5(n)\), \(n\ge2\).
The residual calculation of D.1, which holds
in every weight, gives zero second cohomology
over \(X\) for all these modules. Their first
cohomology groups are finite; the inverse
systems are therefore Mittag--Leffler.
Thus \(H^2(X,A_s)=H^3(X,A_s)=0\).
The full \(H^1(X,A_s)\) was evaluated in
D.36 as the countable product of \(\mathbb Z_5\)'s.

At level \(r\), localization has residue
complex
\[
 \bigoplus_{q\ne5}
 [\,A_s/\mathfrak m^rA_s
        \xrightarrow{\rho_s(\sigma_q)/q-1}
                         A_s/\mathfrak m^rA_s\,][-1].
\]
The twist by \(q^{-1}\) is the tame-inertia
cohomology twist, as directly checked in D.43.
All finite-valued continuous global cochains
factor through a finite Galois quotient, so
the residue term is the direct sum over primes.
The transition maps on the two terms are
surjective. Their inverse limits are exactly
\({\cal C}_\Lambda\), giving (D.44.3).
Each operator in (D.44.1) preserves every
\(\mathfrak m^r\), uniformly in \(q\), so
the displayed family operator is well-defined.

Every \(D_{s,q}\) is injective. If \(m\ne0\),
take its first nonzero \(I\)-adic homogeneous
degree \(j\). On it \(D_{s,q}\) is
multiplication by the nonzero integer
\(q^{j+1}-1\). That degree is torsion-free,
so \(D_{s,q}m\ne0\).
The localization triangle and the base
vanishings now give (D.44.2), including
the equality of unrestricted and base
first cohomology.

For the derived-limit kernel, take the
images of the finite-level differentials.
Writing \(A_r=A_s/\mathfrak m^rA_s\), finite-level
localization and the base second-cohomology
vanishing give the exact sequence
\[
 0\longrightarrow H^1(X,A_r)
 \longrightarrow H^1(\mathbb Q,A_r)
 \longrightarrow\ker D_r\longrightarrow0 .
\]
The first groups are finite and form a
Mittag--Leffler system. Their first derived
limit is zero, so the first derived limits
of the last two systems agree.
This is the comparison with Jannsen's
global first-cohomology term used here.
The image modules of \(D_r\) have
surjective transition maps:
lift any preimage at one finite level to
the next level. The inverse limit of
these image modules is
\[
                 {\cal C}_\Lambda
                     \cap\prod_{q\ne5}D_{s,q}\Lambda.
\]
For each fixed \(q\), the image
\(D_{s,q}\Lambda\) is compact and closed.
Solvability in it is consequently equivalent
to solvability at every finite coefficient
level. Conversely a family in the displayed
intersection has finite support at each
coefficient level, so it belongs to that
finite-level image. Taking the inverse
limits of the kernel-image exact sequences,
as in D.1, proves (D.44.4).
This identifies the actual Jannsen kernel,
not merely the product of its local images.

Finally choose distinct primes \(q_r\ne5\)
which split completely in a finite Galois
extension on which both the action on
\(A_s/\mathfrak m^rA_s\) and
\(\chi_5\bmod5^r\) become trivial.
Chebotarev supplies infinitely many such
primes at each step. Define
\[
 z_{q_r}=D_{s,q_r}(1),\qquad z_q=0
                                  \text{ at other primes}.
\]
Then \(z_{q_r}\in\mathfrak m^r\), so
\(z\in{\cal C}_\Lambda\), and it lies
pointwise in every required image.
Injectivity of all \(D_{s,q}\) makes its
only divided family equal to one at
the \(q_r\)'s and zero elsewhere.
This family does not tend to zero.
Thus \(z\) gives a nonzero class in
(D.44.4). Its finite-support truncations
belong to \(D_s{\cal C}_\Lambda\) and
converge to \(z\). The image is therefore
not closed, proving the final assertion.
The construction uses the actual finite
arithmetic action in Chebotarev, rather
than choosing Frobenius coefficients
independently of that action.
\end{proof}

\begin{statement}{D.45}{Proved}{All induced metabelian Fermat stages and their complete intersection}
Let \(N=5^r\), \(r\ge1\), and put
\[
 \begin{gathered}
 M_n=F_5/(F_5''\gamma_{n+1}F_5),\qquad
 M=F_5/F_5'',\\
 J_{n,r}=\ker\{M_n\longrightarrow(\mathbb Z/N(1))^2\},
 \qquad
 J_r=K_r/F_5''.
 \end{gathered}
\]
The following parameter system determines the unrestricted
section sets of every \(J_{n,r}\) and \(J_r\).
Write
\[
 {\cal H}=c_0(\{\hbox{rational primes}\},\mathbb Z_5)
       =H^1(G_{\mathbb Q},\mathbb Z_5(1)),
\]
using prime-exponent Kummer coordinates. Start with arbitrary
\((A,B)\in{\cal H}^2\), and use the actual abelian
cocycle \((N\kappa_A,N\kappa_B)\). At every weight apply
the fully specified arithmetic expressions (D.39.2) and
(D.41.4)--(D.41.6), with these abelian inputs.
At weight two the test has the particularly explicit form
\[
 \begin{split}
 W_q(A,B)&=
       a_q B_q-b_q A_q+\frac{q-1}{2}A_qB_q,\\
 &\frac{N^2 W_q(A,B)}{q-1}\in\mathbb Z_5
                           \quad(q\ne5),\qquad
 \left(\frac{N^2 W_q(A,B)}{q-1}\right)_{q\ne5}
                                      \in c_0.
 \end{split}                                                   \tag{D.45.1}
\]
Here \(A_q,B_q\) are the prime exponents, and
\(a_q=\kappa_A(\sigma_q)\), \(b_q=\kappa_B(\sigma_q)\)
for the chosen actual tame Frobenius lifts and normalized
Kummer cocycles. At subsequent weights the local inputs
are those of the section already chosen. Each successful
step of weight \(n\ge2\) has precisely the following
translation group:
\[
 \begin{cases}
  \mathbb Z_5^{\,n-1},&n\text{ odd},\\
  (\mathbb Z/5^{\,1+v_5(n)})^{\,n-1},&4\mid n,\\
  0,&n\equiv2\pmod4.
 \end{cases}                                                   \tag{D.45.2}
\]
Its generators are the concrete cocycles (D.40.3) and
(D.40.4), in each indicated Alexander monomial.
There are no additional geometric stabilizers.
Branches through weight \(n\) are exactly
\(\operatorname{Sec}_{\mathbb Q}(J_{n,r})\);
infinite compatible branches are exactly
\(\operatorname{Sec}_{\mathbb Q}(J_r)\).
Every division at every step retains the divided-\(c_0\)
condition; pointwise solvability alone is insufficient.

The level map \(r+1\to r\) is inclusion of the resulting
ambient section classes. In the displayed initial
parameters it sends \((A,B)\) to \((5A,5B)\);
all later coordinates are those of the same ambient
section. In particular
\[
 \begin{split}
 \varprojlim_r\operatorname{Sec}_{\mathbb Q}(J_r)
 &=\bigcap_r
   \operatorname{im}\{\operatorname{Sec}_{\mathbb Q}(J_r)
                         \longrightarrow
                       \operatorname{Sec}_{\mathbb Q}(M)\}\\
 &=H^1(G_{\mathbb Q},A_0)
       \ \simeq\ \prod_{j=1}^{\infty}\mathbb Z_5 ,
 \qquad A_0=M'\text{ with the unit-section action}.
 \end{split}                                                   \tag{D.45.3}
\]
The last isomorphism is noncanonical; the preceding
identifications are canonical relative to the specified
unit section. Exactly the three unit cusp classes in
D.35 lie in this intersection, and none of its classes
is induced by an ordinary rational point.
At a fixed level \(r\), the three induced cusp packets
are precisely the formulas of D.35 with tangent
parameter \(N\eta\), \(\eta\in{\cal H}\).

This statement concerns the explicitly typed metabelian
quotients. The source LM8 F.22--F.23 also writes full
nonmetabelian lower-central kernels, separately from its
metabelian kernels. Their unrestricted versions are
computed here through class five by D.30 and D.32.
No evaluated all-weight nonmetabelian Ihara action is
asserted in this statement.
\end{statement}
\begin{proof}
Kummer theory and the unique factorization of rational
numbers give
\[
 \mathbb Q^\times/(\mathbb Q^\times)^{5^e}
       =\bigoplus_q\mathbb Z/5^e.
\]
The factor \(-1\) disappears because five is odd.
Taking the inverse limit gives \({\cal H}\):
at every finite coefficient level only finitely many
prime exponents are nonzero. In particular this
module is torsion-free, and
\(\bigcap_r5^r{\cal H}=0\).
One may fix its cocycle representatives simultaneously
by choosing an element \(g_*\) with \(\chi_5(g_*)=2\)
and replacing each Kummer cocycle by its unique
cohomologous one vanishing at \(g_*\). Since
\(\chi_5(g_*)-1=1\), this normalization is integral
and commutes with multiplication by \(N\).

For the exact sequence with quotient
\((\mathbb Z/N(1))^2\), that quotient has no arithmetic
invariants: the element \(g_*\) acts by two.
The nonabelian cocycle sequence therefore makes
\[
 H^1(G_{\mathbb Q},J_{n,r})
     \longrightarrow H^1(G_{\mathbb Q},M_n)
\]
injective, with image the inverse image of zero
in \(H^1(G_{\mathbb Q},(\mathbb Z/N(1))^2)\).
The coefficient sequence for multiplication by \(N\)
identifies this image with divisibility by \(N\)
of both abelian Kummer classes. The same argument
applies to \(M\) and \(J_r\). It uses the actual
finite quotient, not an assertion that geometric
conjugacy automatically descends to a subgroup.

We can now evaluate the image by D.41. For the
first step set \(c=d=0\) in (D.41.4). Its constant
term is
\[
 q^{-1}[{\cal R}_q]_{u^0v^0}
   =a\beta-b\alpha+\frac{q-1}{2}\alpha\beta.
\]
Indeed the two constant terms in that formula are
\[
 [K_q]_0=q(a\beta-b\alpha),\qquad
 [P_q]_0=-\frac{q(q-1)}2\alpha\beta;
\]
the beta-series contribution has positive
augmentation degree. Substitution
\((\alpha,\beta,a,b)=N(A_q,B_q,a_q,b_q)\)
gives (D.45.1), including its convergence condition.
All higher substitutions are equally literal in
(D.41.4). D.39 supplies every arithmetic coefficient
there, so this is an arithmetic residue parameter
system, not just a renamed inverse limit of unknown
section sets. D.40 gives the exact groups and
generators (D.45.2). The proof of D.41 gives the
absence of stabilizers and the equivalence between
actual sections and compatible branches.
Thus every finite metabelian stage and its full
limit are accounted for.

The quotient maps at successive Fermat levels
are the inclusions \(J_{r+1}\subset J_r\subset M\).
The injections on section sets just proved make
their transition maps exactly inclusions of
the divisibility loci. Their intersection is
therefore the zero-abelian fibre. A representative
of any class in that fibre can be conjugated
once so that its abelian cocycle is identically
zero: lift to \(M\) the abelian element whose
coboundary is that cocycle. The resulting
geometric cocycle lies in \(A_0\), hence in
every \(J_r\), simultaneously. Two such
representatives are conjugate in \(M\) only
if their conjugator has invariant abelian
image, which is zero. They are consequently
conjugate exactly by \(A_0\). This proves the
middle identification in (D.45.3), including
both its injectivity and its surjectivity.
D.36 evaluates the group as the displayed
countable product and proves that every one
of these classes is unramified outside five.

Finally the three cusp abelian directions are
primitive, and their tangent maps are exactly
those of D.35. Divisibility by \(N\) is thus
equivalent to a tangent parameter \(N\eta\).
Divisibility at every level forces that
parameter to be zero. D.36 proves that the
three resulting unit classes are distinct
and that no ordinary rational-point class
has zero abelian image. This proves all
the cusp and point assertions.
\end{proof}

\begin{statement}{D.46}{Proved}{Large integral fibres of every unramified Fermat comparison}
For \(N=5^r\), \(r\ge1\), the natural comparison
\[
 \operatorname{Sec}_{G_5}(K_r/K_r'')
        \longrightarrow
 \operatorname{Sec}_{G_5}(K_r/F_5'')
                                                               \tag{D.46.1}
\]
is surjective. Above every class on the right
there is a jointly continuous injective family
of section classes with parameter module
\[
              \mathbb Z_5^{\,e_N},\qquad
              e_N=\frac{(N^2-1)^2}{4}.                         \tag{D.46.2}
\]
The classes remain distinct after inflation
to \(G_{\mathbb Q}\). In particular \(e_5=144\)
and \(e_{25}=97\,344\).
Above the unit class at zero, the origin can
be chosen to be its own unit lift; then every
nonzero member of this family is noncuspidal,
and no member is induced by an ordinary
rational point.
\end{statement}
\begin{proof}
Put
\[
 A_r=F_5''/K_r'',\qquad
 V_r=F_5'/K_r'
      =\ker\{K_r^{\rm ab}\longrightarrow N\mathbb Z_5^2\},
 \qquad B_r=\bigwedge\nolimits^2 V_r .
\]
These are actual subquotients of the covering
group. The last exact sequence is integral:
the map on the right is onto, and its
kernel is precisely the image of \(F_5'\).
Schreier's formula gives
\(\operatorname{rank}K_r^{\rm ab}=N^2+1\);
therefore \(V_r\) is a primitive free
sublattice of rank \(N^2-1\).
The derived module \(A_r\) is abelian
because \(F_5'\subset K_r\), hence
\(F_5''\subset K_r'\).

The free pro-five group \(K_r\) has the
integral degree-two quotient
\[
       K_r'/\gamma_3K_r
                     =\bigwedge\nolimits^2K_r^{\rm ab}.
\]
The image of \(F_5''=[F_5',F_5']\) in this
quotient is exactly \(\bigwedge^2V_r\):
commutators map to exterior products, and
every pair of elements of \(V_r\) lifts to
elements of \(F_5'\). Since
\(K_r''\subset\gamma_4K_r\subset\gamma_3K_r\),
this gives an arithmetic-equivariant
surjection of compact modules
\[
                         A_r\twoheadrightarrow B_r.            \tag{D.46.3}
\]
Primitivity of \(V_r\) ensures that its
exterior square embeds integrally in the
displayed degree-two quotient. No
rationalization or dimension count is
being used to claim the surjection.
Conjugation by \(K_r\) acts trivially on
this quotient. Thus twisting \(A_r\) by
any lift of any section of \(K_r/F_5''\)
leaves the induced arithmetic action on
\(B_r\) unchanged.

Every section over \(G_5\), together with
the actual action, factors through the
arithmetic quotient \(Q_5=P\rtimes C_4\)
of D.37. Over \(\mathbb Q(\zeta_5)\)
its graph has pro-five image; this proves
the assertion for the cocycle, as well
as for the original action. Here \(P\)
is free pro-five of rank three and
\(\operatorname{cd}_5 Q_5=1\).
For every countably based compact
five-primary \(Q_5\)-module \(L\),
\[
                         H^2_{\rm cts}(Q_5,L)=0.               \tag{D.46.4}
\]
Indeed choose a cofinal tower of finite
invariant quotients of \(L\). Every finite
second cohomology group is zero.
The maps in their first-cohomology tower
are onto, by the same vanishing for
the finite kernels. Thus its first
derived limit is zero, proving (D.46.4).
The pullback of (D.46.1) by a given
section is an abelian-kernel extension
with kernel \(A_r\). Its continuous
extension class vanishes by (D.46.4),
so the section lifts over \(Q_5\).
For this use of continuous extension
cohomology one may choose a continuous
set section of the quotient: successive
finite clopen coset transversals, refined
along a countable neighborhood basis,
give such a section of this open
surjection of profinite groups.
This proves surjectivity with the
specified ramification bound.

Choose such a lift \(\widetilde s\).
By (D.46.3), (D.46.4) applied to its
kernel, and the observation about
twisting, there is a surjection
\[
 H^1_{\rm cts}(Q_5,(A_r)_{\widetilde s})
                   \longrightarrow H^1_{\rm cts}(Q_5,B_r).
                                                               \tag{D.46.5}
\]
The target is finitely generated over
\(\mathbb Z_5\), since \(P\) has a
finite free generating set.
Its rank is \(e_N\), as follows.
Complex conjugation on \(K_r^{\rm ab}\)
has plus and minus ranks
\((N^2-1)/2\) and \((N^2+3)/2\),
respectively, by the actual compact
and boundary homology calculation in
D.20. The quotient \(N\mathbb Z_5(1)^2\)
has minus rank two. Exactness of the
plus/minus projectors over \(\mathbb Z_5\)
therefore gives
\[
 \operatorname{rank}V_r^+
    =\operatorname{rank}V_r^-=\frac{N^2-1}{2},\qquad
 \dim_{\mathbb Q_5}(B_r\otimes\mathbb Q_5)^-=e_N .
                                                               \tag{D.46.6}
\]
The executed signed-permutation
calculation in \texttt{de\_fermat\_relative.py}
checks the exterior-square ranks
at \(N=5,25\), as well as the
constant-term substitution in
(D.45.1).
At a good prime different from five,
the eigenvalues on \(K_r^{\rm ab}\)
have positive weights one or two;
the same is true on its invariant
subspace \(V_r\), and exterior products
again have strictly positive weight.
Thus \(H^0(Q_5,B_r)=0\).
The second cohomology is zero by
(D.46.4). The integral global Euler
formula used in D.20 now gives
\(\operatorname{rank}H^1(Q_5,B_r)=e_N\).

Choose a free submodule of rank \(e_N\)
in this finitely generated target,
and lift a basis through (D.46.5)
to actual continuous cocycles
\(b_1,\ldots,b_{e_N}\).
For \(z=(z_i)\in\mathbb Z_5^{e_N}\)
multiply \(\widetilde s\) by the
\((A_r)_{\widetilde s}\)-cocycle
\(\sum_i z_i b_i\). The kernel is
abelian, so this is an actual section.
Its dependence on both \(z\) and
the Galois argument is continuous.
Different parameters have different
images in \(H^1(Q_5,B_r)\), hence
different \(A_r\)-cohomology classes.

There is no geometric stabilizer
identifying these classes. The twisted
group \(K_r/F_5''\) embeds in the
twisted metabelian group \(M\).
A nonidentity fixed element would
have a first nonzero lower-central
grade. On that torsion-free grade
the action, even after twisting,
is the positive Tate character
\(\chi_5^j\), which has no invariants.
Thus \(H^0(Q_5,(K_r/F_5'')_s)=1\).
The exact nonabelian fibre sequence
identifies the actual comparison
fibre with its \(A_r\)-cohomology
torsor, proving injectivity in
(D.46.2). Inflation first to \(G_5\)
and then to \(G_{\mathbb Q}\) is
injective on this first cohomology:
the action factors through \(Q_5\),
and inflation in degree one has
zero kernel. The same vanishing
of the geometric stabilizer holds
over \(G_{\mathbb Q}\). All members
therefore remain distinct as full
global section classes.

For the last assertion, choose the
unit lift as the origin. The open
Fermat cover of odd degree has
exactly three rational cusps.
The distinct boundary classes in
\(K_r^{\rm ab}\) distinguish its
geometric cusp inertias; hence
a cuspidal section over \(\mathbb Q\)
must lie over one of these rational
cusps even in this metabelian quotient.
Their inertia generators map to
the \(N\)-th powers of the three
tripod peripheral generators,
with unit leading tangent factors;
a possible sign has zero
five-primary Kummer class.
If a cuspidal class maps to the
unit class at zero, its tangent
parameter must be zero, because
its ambient abelian vector is
\(N\) times a primitive direction
and \({\cal H}\) is torsion-free.
The other two unit cusp classes
have different images by D.36.
There is consequently only one
cuspidal class in this fibre,
namely its chosen origin.
An ordinary rational point would
map to an ordinary tripod point
whose abelian Kummer pair is zero,
which D.36 excludes. This proves
the noncuspidal and nonpoint
claims without appealing to a
full section conjecture.
\end{proof}

\begin{statement}{D.47}{Proved}{Uniform ramification and compactness of every global Alexander fibre}
Let \(s=(\kappa_A,\kappa_B)\) be any abelian Kummer
pair, with \(A,B\in{\cal H}\), and let \(A_s=\Lambda C\)
have the actual twisted action of D.35.
Existence of a metabelian lift of \(s\) is not assumed.
Put \(\mathfrak m=(5,u,v)\), \(W_r=A_s/\mathfrak m^rA_s\),
and, with \(v_5(0)=+\infty\), define the finite set
\[
 S_r=\{5,\infty\}\ \cup\
       \{q\ne5:\min(v_5(A_q),v_5(B_q))<r-1\}.
                                                               \tag{D.47.1}
\]
Every integral continuous \(A_s\)-valued cocycle,
reduced modulo \(\mathfrak m^r\), factors through
\(G_{\mathbb Q,S_r}\). The finite set \(S_r\)
depends only on the abelian pair, uniformly
for all the integral cocycles.

Consequently \(H^1_{\rm cts}(G_{\mathbb Q},A_s)\),
with its continuous-cochain quotient topology,
is a compact, countably based, torsion-free
\(\mathbb Z_5\)-module. In particular there is
a noncanonical topological isomorphism
\[
 H^1_{\rm cts}(G_{\mathbb Q},A_s)
                \simeq\prod_{j\in I_s}\mathbb Z_5,
                    \qquad |I_s|\le\aleph_0 .                  \tag{D.47.2}
\]
Thus every nonempty unrestricted metabelian
section fibre is a torsor under such a compact
module. For \(A,B\) supported only at five,
D.36 gives \(|I_s|=\aleph_0\).
The value of \(|I_s|\) for all arbitrarily
ramified pairs is not determined by this
statement.
\end{statement}
\begin{proof}
Fix \(q\ne5\), and let \(\tau_q\) be its
five-primary tame inertia generator.
Its action on \(A_s\) is multiplication by
\[
                   \mu_q=t_1^{A_q}t_2^{B_q}.
\]
The original arithmetic action is unramified
there. If \((A_q,B_q)\ne(0,0)\), D.43 gives
the integral vanishing
\(H^1(G_{\mathbb Q_q},A_s)=0\).
The restriction of every global cocycle is
therefore an actual integral coboundary:
for some \(h_q\in A_s\), its value on inertia is
\[
                         (\mu_q-1)h_q,                        \tag{D.47.3}
\]
up to the immaterial convention for the
sign of a coboundary. No bound on \(h_q\)
has to be chosen: all its coefficients
are in the fixed integral compact module.

If \(A_q=B_q=0\), inertia acts trivially.
The restriction of a cocycle to inertia
is then a homomorphism with value \(z_q\)
on \(\tau_q\). The tame relation forces
\[
                  \rho_s(\sigma_q)z_q=qz_q.
\]
A nonzero \(z_q\) has a first nonzero
augmentation degree \(j\). On that grade
\(\rho_s(\sigma_q)\) is multiplication
by \(q^{j+2}\). Thus
\((q^{j+2}-q)[z_q]_j=0\), impossible in
the torsion-free graded lattice. Hence
\(z_q=0\). Wild inertia and the
prime-to-five tame factors have zero
cocycle values, since their action is
trivial and they have no nonzero
continuous homomorphism to a pro-five
additive module.

For \(\alpha\in5^e\mathbb Z_5\), the
binomial series satisfies
\[
                  (1+u)^\alpha-1\in\mathfrak m^{e+1}.
                                                               \tag{D.47.4}
\]
For integral exponents this follows
inductively from
\((1+z)^5-1=5z+10z^2+10z^3+5z^4+z^5\);
continuity gives every five-adic exponent.
The same assertion holds for \(v\).
Thus if \(q\notin S_r\),
\(\mu_q-1\in\mathfrak m^r\).
The action on \(W_r\) is unramified
at \(q\), and (D.47.3), or the preceding
zero-residue argument, shows that every
integral cocycle also has zero inertia
value in \(W_r\). Its graph together
with its action is a finite continuous
Galois representation unramified
outside \(S_r\), proving the first
assertion. Finiteness of \(S_r\)
is exactly the \(c_0\) condition on
the two exponent families.

Here is a precise compactness argument
which does not assume that all finite
coefficient cocycles lift integrally.
For each \(r\), the set
\[
                E_r=Z^1(G_{\mathbb Q,S_r},W_r)
\]
is finite: its first cohomology and
its group of coboundaries are finite
for a finite ramification set and
finite coefficients. Inflate its
members to \(G_{\mathbb Q}\).
The space of integral cocycles is
exactly the closed subset of
\(\prod_rE_r\) defined by equality
of the inflated cocycles after
coefficient reduction. Every
compatible family defines a
continuous map into
\(\varprojlim_r W_r=A_s\) and
satisfies the cocycle equation;
conversely the uniform ramification
assertion puts every integral
cocycle in this closed subset.
These are finite compatibility
conditions on a product of finite
discrete sets, so the resulting
space is compact and countably based.

One must compare after inflation:
an arbitrary cocycle in \(E_{r+1}\)
need not reduce to a cocycle
unramified outside the smaller
set \(S_r\). Thus no nonexistent
transition map between all the
finite \(E_r\)'s is being used.
Equivalently one may replace them
by the finite subsets satisfying
all extendability conditions.
The closed-subset description
already proves what is needed,
without asserting surjectivity
of any finite-coefficient lifting
map.

The topology in this description
is precisely uniform
\(\mathfrak m\)-adic convergence
of continuous cocycles. The
coboundaries form the continuous
image of the compact module \(A_s\);
this image is compact and closed.
The first cohomology quotient
is therefore compact, Hausdorff,
and countably based. D.35 proves
that it is torsion-free for every
abelian twist, using the actual
primitive Soul\'e multiplier
modulo five.

Finally the Pontryagin dual of a
compact torsion-free pro-five
\(\mathbb Z_5\)-module is a divisible
five-primary torsion group:
surjectivity of multiplication
by five on the dual is dual to
injectivity on the original
compact module. A divisible
five-primary group is a direct
sum of copies of
\(\mathbb Q_5/\mathbb Z_5\);
choosing a basis of its five-torsion
and successively compatible
five-power roots proves this
classification. Countable
basedness makes that basis
at most countable. Dualizing
gives (D.47.2).

This argument deliberately leaves
the rank question separate.
Unbounded ranks in finite
augmentation quotients alone
cannot settle it: the compact
system \(L_n=\mathbb Z_5^n\)
with transition
\((x_1,\ldots,x_{n+1})
 \mapsto(5x_1,\ldots,5x_n)\)
has increasing finite ranks
and zero inverse limit.
For the ramified Alexander
modules the actual images of
the cohomology transition maps,
including their integral indices,
must still be controlled to
exclude this phenomenon.
\end{proof}

\begin{statement}{D.48}{Proved}{The unrestricted second cohomology of every finitely ramified Alexander twist}
Let \(s=(\kappa_A,\kappa_B)\) be an abelian pair
whose prime-exponent families have finite support.
Choose a finite \(S\) containing that support,
five, and the infinite place. Put
\[
 \begin{gathered}
 W_r=A_s/\mathfrak m^r A_s,\qquad
 {\cal C}_S=
   \{(z_q)_{q\notin S}\in\prod_{q\notin S}\Lambda:
                  z_q\longrightarrow0\text{ in }\mathfrak m\},\\
 D_{s,q}=q^{-1}\rho_s(\sigma_q)-1,\qquad
 D_s=\prod_{q\notin S}D_{s,q}.
 \end{gathered}
\]
Equip \({\cal C}_S\) with the inverse-limit
topology of the discrete direct sums
\(\bigoplus_{q\notin S}\Lambda/\mathfrak m^r\),
as in D.44. Every operator \(D_{s,q}\)
is the explicitly specified arithmetic
operator of D.39 and D.41 and is injective.
There are exact sequences
\[
 \begin{gathered}
 H^1(G_{\mathbb Q,S},A_s)
                  \xrightarrow{\ \sim\ }
                    H^1(G_{\mathbb Q},A_s),\\
 0\longrightarrow H^2(G_{\mathbb Q,S},A_s)
  \longrightarrow H^2(G_{\mathbb Q},A_s)
  \longrightarrow {\cal C}_S/D_s{\cal C}_S
  \longrightarrow0.
 \end{gathered}                                                 \tag{D.48.1}
\]
The first term in the second row is compact
and is the inverse limit of its finite
coefficient groups. In particular it is
retained, rather than discarded because
those finite groups are finite.

More precisely put
\[
 \begin{split}
 {\cal L}_s&=
  \varprojlim_r\ker\{H^2(G_{\mathbb Q,S},W_r)
                              \longrightarrow
                         H^2(G_{\mathbb Q},W_r)\},\\
 {\cal K}_s&=
 \frac{{\cal C}_S\cap
             \prod_{q\notin S}D_{s,q}\Lambda}
      {D_s{\cal C}_S}.
 \end{split}
\]
Then the actual unrestricted Jannsen kernel
fits into the exact sequence
\[
 0\longrightarrow{\cal L}_s
  \longrightarrow
       \varprojlim\nolimits_r^1H^1(G_{\mathbb Q},W_r)
  \longrightarrow{\cal K}_s
  \longrightarrow0.                                           \tag{D.48.2}
\]
Here \({\cal L}_s\) is compact. For every such
abelian pair, \({\cal K}_s\) contains an explicit
algebraic copy of
\[
           \left(\prod_{j\ge1}\mathbb Z_5\right)
                             \big/c_0(\mathbb N,\mathbb Z_5).
                                                               \tag{D.48.3}
\]
In particular the Jannsen kernel is nonzero
for every finitely ramified Alexander twist,
including twists for which no metabelian
section exists. Equations (D.48.1)--(D.48.2)
separate the explicit unrestricted residue
term from the still unevaluated compact
base cohomology for a general ramified pair.
D.64 evaluates this base and proves that
\({\cal L}_s=0\) for the family supported at
five and one prime of order twenty modulo
twenty-five. D.65 determines the residue kernel
in (D.48.2) as an abstract countable product
of \(\mathbb Z_5\), and proves that its extension
by \({\cal L}_s\) splits algebraically, while
retaining the original cohomology topology.
\end{statement}
\begin{proof}
At each finite coefficient level the module
\(W_r\) is unramified outside \(S\).
Localization, taking the union over finite
sets of additional primes, gives the
continuous residue triangle
\[
 R\Gamma(G_{\mathbb Q,S},W_r)
 \longrightarrow R\Gamma(G_{\mathbb Q},W_r)
 \longrightarrow
 [C_r\xrightarrow{D_r}C_r][-1],
 \qquad
 C_r=\bigoplus_{q\notin S}\Lambda/\mathfrak m^r .
                                                               \tag{D.48.4}
\]
This is the same finite-coefficient Gysin
calculation used in D.1 and D.44.
The factor \(q^{-1}\) is the inertia
Tate twist; the action on the other
factor is the actual \(\rho_s(\sigma_q)\).
The transition maps on both terms \(C_r\)
are surjective, as are those on their
image modules \(D_rC_r\): lift a
preimage and apply the next differential.
Consequently derived inverse limit
turns the residue term in (D.48.4)
into exactly
\([{\cal C}_S\xrightarrow{D_s}{\cal C}_S][-1]\).
No closure of the differential's image
is taken.

The finite \(S\)-ramified cohomology
groups in degrees one and two are finite.
Their inverse systems are therefore
Mittag--Leffler, without needing their
individual transition maps to be onto.
The odd-primary cohomological dimension
of \(G_{\mathbb Q,S}\) is at most two.
It follows that
\[
 \begin{split}
 H^2(G_{\mathbb Q,S},A_s)
       &=\varprojlim_rH^2(G_{\mathbb Q,S},W_r),\\
 H^3(G_{\mathbb Q,S},A_s)&=0. 
 \end{split}                                                   \tag{D.48.5}
\]
The first displayed group is compact.
If \(D_{s,q}z=0\), the first nonzero
augmentation degree \(j\) of \(z\)
would satisfy
\((q^{j+1}-1)[z]_j=0\).
The graded lattice is torsion-free
and this integer is nonzero, so
\(z=0\). Thus \(D_s\) is injective.
The long exact sequence of the
inverse-limit triangle now gives
(D.48.1).

We identify precisely the global
derived-limit correction. Write
\[
 B_r=H^2(G_{\mathbb Q,S},W_r),\qquad
 H_r=H^2(G_{\mathbb Q},W_r),\qquad
 \beta_r:B_r\longrightarrow H_r .
 \]
The finite-level sequence of (D.48.4)
ends in the exact row
\[
 0\longrightarrow\operatorname{im}\beta_r
   \longrightarrow H_r
   \longrightarrow\operatorname{coker}D_r
   \longrightarrow0 .
 \]
All groups \(\operatorname{im}\beta_r\)
are finite, so taking the inverse
limit preserves exactness on the right.
Similarly the finite kernels of
\(\beta_r\) have zero first derived
limit. Therefore the natural map
\[
 B:=\varprojlim_rB_r
              \longrightarrow
                \varprojlim_r\operatorname{im}\beta_r
 \]
is onto, with kernel \({\cal L}_s\).
The right-hand vertical map in the
comparison with (D.48.1) is
\[
 {\cal C}_S/D_s{\cal C}_S
              \longrightarrow
                   \varprojlim_r\operatorname{coker}D_r .
 \]
It is onto, because the systems \(C_r\)
and \(D_rC_r\) have surjective maps.
Its kernel is \({\cal K}_s\).
Indeed at each fixed prime \(q\)
the image \(D_{s,q}\Lambda\) is
compact and closed. Solvability
at every finite coefficient level
is equivalent to membership in
that image. Finite support at each
level gives exactly
\[
       \varprojlim_rD_rC_r
            ={\cal C}_S\cap
                  \prod_{q\notin S}D_{s,q}\Lambda .
 \]
Comparing the two exact rows and
using the surjectivity on the left,
the snake lemma gives (D.48.2).
The middle vertical kernel is
exactly
\(\varprojlim^1_rH^1(G_{\mathbb Q},W_r)\)
by Jannsen's continuous-cohomology
sequence. In particular the compact
term \({\cal L}_s\) has not been
silently assumed to vanish.

Finally choose distinct primes
\(q_j\notin S\) splitting completely
in a finite Galois extension which
trivializes both \(W_j\) and
\(\chi_5\bmod5^j\). Chebotarev
gives infinitely many choices at
each step, with the actual finite
arithmetic representation prescribed.
Then
\[
                      D_{s,q_j}(1)\in\mathfrak m^j .
 \]
For \(a=(a_j)\in\prod_j\mathbb Z_5\)
define the residue family
\[
 z(a)_{q_j}=a_jD_{s,q_j}(1),\qquad
 z(a)_q=0\quad\hbox{at all other primes}.
 \]
This lies in \({\cal C}_S\) and
pointwise in every required image.
Since each \(D_{s,q}\) is injective,
its only divided family is
\(a_j\cdot1_\Lambda\) at \(q_j\),
and zero elsewhere. That family
is in \({\cal C}_S\) if and only if
\(a_j\to0\) five-adically; for a
constant in \(\Lambda\),
membership in \(\mathfrak m^e\)
means divisibility by \(5^e\).
Thus \(a\mapsto[z(a)]\) has kernel
exactly \(c_0(\mathbb N,\mathbb Z_5)\),
proving the injection (D.48.3).
The residue quotient has a
nonclosed differential image in
its displayed topology: every
finite-support truncation of
\(z(1,1,\ldots)\) is in that
image, while the full family
is not. This final observation
concerns the displayed residue
quotient topology, without
assigning an unproved topological
splitting to (D.48.1).
\end{proof}

\begin{statement}{D.49}{Proved}{M0 with finitely many cuspidal primewise images in every genus}
Let \(U/K\) be a hyperbolic curve
over a number field, and let \(s\) be a full
section whose image in every geometric pro-\(\ell\)
quotient is geometric, in the sense of Lead 19.
If only finitely many of these primewise images
are cuspidal, then \(s\) is geometric.
In fact it is induced by a single ordinary
rational point, and all its primewise images
are induced by that point.

For affine \(U\), the proof produces a finite
set of places \(S\), independent of \(\ell\),
for which every ordinary point underlying a
primewise image is \(S\)-integral on one
fixed smooth model of \(U\), after a fixed
finite extension of \(K\).
Consequently this includes every genus-one
case with positive Mordell--Weil rank when
the primewise images are all ordinary, or
only finitely many are cuspidal.
It strengthens the genus-zero result D.26.
\end{statement}
\begin{proof}
The proper case has genus at least two and
is already D.25. Suppose \(U\) is affine,
with smooth proper compactification \(X\)
and nonempty reduced boundary \(D=X\setminus U\).
Let \(C\) be the finite set of rational
primes whose primewise images are cuspidal.
For \(\ell\notin C\), write \(P_\ell\in U(K)\)
for the unique point inducing the
pro-\(\ell\) section. Choose one such
prime \(\ell_0\) and point \(P_{\ell_0}\).

Pass to one finite extension \(L/K\)
splitting all the geometric boundary
points. Choose a finite set \(S\) of
places of \(L\) which contains the
Archimedean places, the places over
\(C\cup\{\ell_0\}\), all bad places
of the pair \((X,D)\), and all
places at which \(P_{\ell_0}\) fails
to be integral on \(U\).
Enlarge it if needed so that
\((X,D)\) has a smooth proper model
with disjoint boundary sections over
\(\mathcal O_{L,S}\).
The reference point then extends to
an integral section of its open
complement.

We first spell out the local
integrality test. Let \(v\notin S\)
have residue characteristic \(p\),
and let \(\ell\ne p\).
Use \(P_{\ell_0}\) as geometric
and arithmetic base point.
Good reduction of the pair and
this integral base point make
the action of local inertia
trivial on the geometric
pro-\(\ell\) group. An integral
ordinary point has trivial
geometric inertia coordinate.
These assertions are the
prime-to-\(p\) tame specialization
of \cite[Expos\'e XIII,
Corollary 2.9 and
\S\S2.10--2.12]{DE_SGA1},
applied to this smooth pair.

Conversely, if an ordinary point
\(P\in U(L_v)\) is not integral,
properness extends it to \(X\),
where its reduction is one of the
boundary sections \(d\).
Choose a parameter \(t\) for
that section. Then
\(e=v(t(P))\) is a positive
integer. The punctured formal
disc has pro-\(\ell\) inertia
\(\mathbb Z_\ell(1)\), whose
Kummer coordinate under \(P\)
is \(\kappa_\ell(t(P))\).
Its restriction to tame inertia
is \(e\) times the tame character:
the unit part of \(t(P)\)
contributes no inertia for
\(\ell\ne p\). Inclusion of
the disc therefore makes the
geometric inertia coordinate
of the point section conjugate
to \(c_d^{\,e}\), where \(c_d\)
is the peripheral generator.
This is nonidentity. If
\(\#D\ge2\), the boundary
vector of \(c_d\) is primitive
and nonzero in the torsion-free
abelianization. If \(\#D=1\),
the genus is positive and its
degree-two class is the
primitive nonzero symplectic
sum \(\sum_i[a_i,b_i]\).
In either case multiplication
by the nonzero integer \(e\)
cannot kill this class.
Thus, for ordinary local points,
\[
 \begin{gathered}
 P\text{ is integral at }v\\
 \Longleftrightarrow\
 \text{its pro-}\ell\text{ point section has
 trivial geometric inertia at }v
 \quad(\ell\ne p).
 \end{gathered}                                                 \tag{D.49.1}
\]
Triviality here is preserved
by geometric conjugation,
because arithmetic inertia
acts trivially with the chosen
integral base point.

Now restrict the full section
\(s\) to \(G_{L_v}\).
D.56 with the class of all
finite groups makes this full
local section ordinary or
cuspidal, including for affine curves.
Its pro-\(p\) image is induced
by \(P_p\), since \(p\notin C\).
The local pro-\(p\) point and
cusp separation established
in D.25 forces the underlying
local point to be exactly
\(P_p\in U(L_v)\). This uses
the residue-characteristic
prime \(p\); it does not
assume point-section
injectivity at an unrelated
prime.

The same full local section,
projected instead to
pro-\(\ell_0\), is the
localization of the global
pro-\(\ell_0\) class of
\(P_{\ell_0}\).
Hence the pro-\(\ell_0\)
point sections of \(P_p\)
and \(P_{\ell_0}\) are
geometrically conjugate.
The latter has trivial
inertia, and \(\ell_0\ne p\)
by the definition of \(S\).
Equation (D.49.1) therefore
forces \(P_p\) to be integral
at \(v\).

Finally let \(\ell\notin C\)
be any ordinary prime index.
If \(\ell=p\), its point
is \(P_p\) and is already
integral. If \(\ell\ne p\),
project the same full local
section to pro-\(\ell\).
The point sections of
\(P_\ell\) and \(P_p\)
are geometrically conjugate.
Since \(P_p\) is integral,
its inertia coordinate is
trivial; (D.49.1) forces
\(P_\ell\) to be integral
as well. This proves
simultaneously, for every
ordinary index \(\ell\)
and every \(v\notin S\),
that
\[
                       P_\ell\in U(\mathcal O_{L,S}).
                                                               \tag{D.49.2}
\]

Siegel--Mahler finiteness now
makes the set in (D.49.2)
finite: an affine hyperbolic
curve has positive genus, or
has genus zero with at least
three geometric boundary
points. The precise
\(S\)-integer theorem and
the reduction from positive
genus with one or two
boundary points are proved
in \cite[Theorem 1 and
\S3]{DE_CorvajaZannier2002}.
This is why the argument
also covers a once-punctured
positive-rank genus-one
curve, where there need
not be a nonconstant global
unit function.

The original \(K\)-points
\(\{P_\ell:\ell\notin C\}\)
are consequently finite.
Adding the finite boundary
gives finite support for
all the primewise geometric
images. D.25, including
its full section-neighborhood
and Stoll finite-support
argument, makes \(s\)
geometric over \(K\).
Its real localizations use
the real section theorem
as in D.24; the integrality
argument above concerns
only finite places.
No descent of an arbitrary
section conjugator from
\(L\) to \(K\) is used:
the finite extension only
proves finiteness of the
original set of \(K\)-points.
Since a pro-\(\ell_0\)
image is ordinary, the
resulting global geometric
section is ordinary.
Global pro-\(\ell\)
point/cusp injectivity from
D.25 then identifies every
primewise image with that
same point.
\end{proof}

\begin{statement}{D.50}{Proved}{Finite-field separation of ordinary and cuspidal pro-prime sections}
Let \(k\) be a finite field of characteristic \(p\),
let \(U/k\) be an affine hyperbolic curve, and let
\(\ell\ne p\). An ordinary \(k\)-point section of the
geometrically pro-\(\ell\) fundamental group of \(U\)
is not geometrically conjugate to a section in a
rational cusp packet. This assertion does not require
injectivity of the ordinary point-section map over
a finite field.

Consequently, for a smooth proper curve with disjoint
boundary sections over a complete discrete valuation ring,
whose boundary complement is hyperbolic and whose
finite residue field has characteristic \(p\),
the pro-\(\ell\) section of an integral ordinary
point is not geometrically conjugate to any
cuspidal section, for \(\ell\ne p\).
\end{statement}

\begin{proof}
We give the arithmetic configuration-fibre argument,
including the particular comparison of groups that
is needed. Write
\[
 U_2=(U\times_k U)\setminus\operatorname{diag}(U),
 \qquad
 1\longrightarrow F\longrightarrow\Pi_{U_2}
   \xrightarrow{p_1}\Pi_U\longrightarrow1,
                                                               \tag{D.50.1}
\]
where both arithmetic fundamental groups are
geometrically pro-\(\ell\).
The sequence is exact, and its geometric fibre is
the full pro-\(\ell\) group of a once further punctured
curve, by \cite[Definition 1.3 and
Proposition 1.4(i)]{DE_Wakabayashi2012}.
For a section \(s:G_k\to\Pi_U\), put
\[
 E_s=\Pi_{U_2}\times_{\Pi_U,s}G_k .
                                                               \tag{D.50.2}
\]
Geometrically conjugate sections give isomorphic
extensions \(E_s\) with kernel \(F\): lift the
conjugating element through the surjection in
(D.50.1) and conjugate the pullback.
If \(s=s_x\) is the section of \(x\in U(k)\),
the homotopy exact sequence identifies \(E_{s_x}\)
with \(\Pi_{U\setminus\{x\}}\), over \(G_k\),
and identifies their entire geometric kernels.

For comparison, fix a rational cusp \(d\).
The logarithmic compactification of the family
\(U_2\to U\) has, over \(d\), the stable fibre
obtained by attaching a tripod to \(U\) at \(d\).
This is the explicit blowup construction of
\cite[Section 2, pp.~14--16]{DE_Wakabayashi2012}.
Let \(D_d\subset\Pi_U\) be its cuspidal decomposition
group. The arithmetic logarithmic fibre group
is precisely \(p_1^{-1}(D_d)\).
Here is the group comparison behind this assertion.
Its major vertex, minor vertex and node groups are,
respectively,
\[
 D_d\times_{G_k}\Pi_U,\qquad
 \Pi_{\mathbb P^1\setminus\{0,1,\infty\}}
       \times_{G_k}D_d,\qquad
 D_d\times_{G_k}D_d .                                          \tag{D.50.3}
\]
The maps to the first \(D_d\) are the structural
maps of the logarithmic fibre. These descriptions
and the node inclusions are
\cite[Lemma 4.1]{DE_Wakabayashi2012};
the node's two parameter coordinates differ by
the unimodular change \((a,b)\mapsto(a+b,b)\)
written in that lemma's proof.
The vertex and node groups form an amalgam
in the category of arithmetic groups with pro-\(\ell\)
geometric kernel, by
\cite[Lemma 4.2(ii)]{DE_Wakabayashi2012}.
The node projection to \(D_d\) splits by the
diagonal \(D_d\to D_d\times_{G_k}D_d\);
these compatible splittings turn the arithmetic
vertex amalgam into the semidirect product of
the relative vertex amalgam with \(D_d\).
Thus taking the relative kernel here is justified
by these splittings and the universal property,
without an assertion that arbitrary pushouts
preserve kernels.
On taking the relative geometric kernel over \(D_d\),
this is the actual configuration-fibre amalgam
\[
 F=A*_{C}B,\qquad
 A=\pi_1^{(\ell)}(U_{\bar k}),\quad
 B=\pi_1^{(\ell)}
       ((\mathbb P^1\setminus\{0,1,\infty\})_{\bar k}),
 \quad C=\mathbb Z_\ell(1).                                   \tag{D.50.4}
\]
The edge \(C\) is the old cusp inertia in \(A\)
and one peripheral subgroup in \(B\).
This identifies the whole kernel in (D.50.1),
not a quotient of that kernel: the geometric
van Kampen isomorphism is
\cite[Definition 1.4 and
Proposition 1.5(iii)]{DE_Mochizuki2010},
and its logarithmic specialization is exactly
the amalgam just quoted from Wakabayashi.
Thus the natural map from the logarithmic fibre
to \(p_1^{-1}(D_d)\) is an isomorphism on
the relative geometric kernels and on the quotient
\(D_d\), hence is an isomorphism.

Suppose \(s\) is any section in the packet of \(d\).
Pull (D.50.3) back along \(s:G_k\to D_d\).
The node group contains a lift of \(s\), obtained
from the pair \((s,s)\) in its two \(D_d\) coordinates.
Conjugation by this lift preserves the two
vertex subgroups \(A,B\) and their common edge
\(C\) in (D.50.4).
This observation handles every tangent parameter:
the argument does not replace the given section
by the unit tangent section.

Both \(A/\overline{\langle\!\langle C\rangle\!\rangle}\)
and \(B/\overline{\langle\!\langle C\rangle\!\rangle}\)
have nonzero homomorphisms to \(C_\ell=\mathbb Z/\ell\).
For \(B\) the quotient is \(\mathbb Z_\ell\).
For \(A\), filling in the cusp leaves an abelianization
of rank \(2g+r-2\) if \(r\ge2\), and rank \(2g\)
if \(r=1\); these ranks are positive because
\(2g-2+r>0\).
Choose surjective characters from these two groups
to \(C_\ell\). After a finite extension of \(k\),
the chosen Frobenius lift fixes both characters:
each belongs to the finite set of homomorphisms
from a finitely generated pro-\(\ell\) group to
\(C_\ell\). The universal property of the amalgam
therefore gives a surjection
\[
 F\twoheadrightarrow P=C_\ell*_{\mathrm{pro}\text{-}\ell}C_\ell
                                                               \tag{D.50.5}
\]
on which this Frobenius acts identically.
The quotient kills \(C\); it sends the two vertex
groups onto the two indicated factors.

Let \(R\) be the kernel of the natural map
\(P\to C_\ell^2\), and let \(H\subset F\) be its
inverse image. Then \([F:H]=\ell^2\) and
\(H\twoheadrightarrow R\).
The group \(R\) is free pro-\(\ell\) of rank
\[
                  \ell^2-2\ell+1=(\ell-1)^2.                    \tag{D.50.6}
\]
For completeness, the finite quotient graph has
the vertices \(C_\ell^2/C_\ell\) from each of the
two factors and the edges \(C_\ell^2\), with the
two quotient maps as incidences. This is the
connected complete bipartite graph with \(\ell\)
vertices in each part. The vertex and edge
stabilizers in \(R\) are trivial, so its
fundamental pro-\(\ell\) group is free with
one generator for every edge outside a spanning
tree, proving (D.50.6). Equivalently, the
Reidemeister--Schreier presentation of
\(\langle a,b\mid a^\ell=b^\ell=1\rangle_{\mathrm{pro}\text{-}\ell}\)
has precisely this many free generators.
In particular there is a surjection
\[
 H^{\mathrm{ab}}\otimes\mathbb Q_\ell
       \twoheadrightarrow
       \mathbb Q_\ell^{\,(\ell-1)^2}                            \tag{D.50.7}
\]
on which Frobenius is the identity. This is an
actual Galois-equivariant quotient: the action
on the whole group \(P\) in (D.50.5) was trivial,
so it is trivial on \(R\) and its abelianization.

If the given cusp section were geometrically
conjugate to \(s_x\), the isomorphism of the
pullbacks (D.50.2) would identify \(H\) with
the geometric group of a connected finite
étale cover \(V\to U\setminus\{x\}\), after
the finite constant extension already made.
Indeed \(H\) is invariant under the chosen
Frobenius lift, and \(H\rtimes G_k\) is an
open arithmetic subgroup of the identified
extension. Thus (D.50.7) would be a nonzero
trivial quotient of
\(\pi_1^{(\ell)}(V_{\bar k})^{\mathrm{ab}}
\otimes\mathbb Q_\ell\).
This is impossible. The boundary homology
has arithmetic Frobenius eigenvalues \(q\zeta\),
and the quotient coming from the Jacobian of
the smooth proper compactification has eigenvalues
of complex absolute value \(q^{1/2}\).
Here \(q=\#k>1\), and \(\zeta\) is a root of
unity coming from the permutation of the boundary.
The assertion for the Jacobian is
\cite[Theorem 14.1(b) and (9.20)]{DE_MilneAV};
the boundary assertion follows directly from
the localization sequence.
Neither part has eigenvalue one.
This contradicts (D.50.7) and proves
the finite-field assertion.

Finally consider the discrete valuation ring in
the statement, with fraction field \(K_v\).
Use the integral ordinary point as base point.
Prime-to-\(p\) tame specialization makes local
arithmetic inertia act trivially on the
geometric pro-\(\ell\) group, and the ordinary
point section has trivial geometric inertia.
If a cuspidal section were conjugate to it,
its inertia coordinate would also vanish.
Peripheral inertia is a copy of
\(\mathbb Z_\ell(1)\) in the geometric group,
so this forces its tangent valuation coordinate
to be zero. The section consequently descends,
under tame specialization, to a section in
the corresponding cusp packet over the
finite residue field. The integral point
descends to an ordinary residue point.
Compatibility of point sections, cusp inertia
and these specialization maps is
\cite[Expos\'e XIII, Corollary 2.9 and
\S\S2.10--2.12]{DE_SGA1}, as in D.49.
The finite-field assertion gives the required
contradiction.
\end{proof}

\begin{statement}{D.51}{Proved}{Primewise geometric sections have one global geometric origin}
Let \(U\) be any hyperbolic
curve over a number field \(K\). If a full
section has a geometric image in every
geometrically pro-prime quotient, then
the full section is geometric, and all
these primewise images come from the same
ordinary point or the same cusp packet.
Thus M0 of Lead 19 holds, including the
infinitely mixed genus-zero and genus-one
cases left open before D.50. Together with
D.24, both M0 and M1 are proved; their full
local section input is D.56.
\end{statement}

\begin{proof}
The proper case, and the case of only finitely
many underlying ordinary points, are already
D.25. Hence suppose that \(U\) is affine
and some pro-\(\ell_0\) image is the ordinary
point section of \(P_{\ell_0}\in U(K)\).
Make a finite extension \(L/K\) splitting
the boundary, and choose a finite set \(S\)
of places of \(L\) containing the infinite
places, the places above \(\ell_0\), the
places of bad reduction of the smooth
proper pair, and those where
\(P_{\ell_0}\) is not integral.
Outside \(S\) the pair has good reduction
and the reference point is integral.

Let \(v\notin S\) have residue characteristic
\(p\), so \(p\ne\ell_0\).
The full localization of the section over
\(L_v\) is ordinary or cuspidal by D.56
applied to the class of all finite groups. Its pro-\(\ell_0\) image
is the localization of \(P_{\ell_0}\).
The local assertion of D.50 excludes the
cuspidal possibility. The localization is
therefore an ordinary point section.
If its pro-\(p\) global image had been
cuspidal, its localization would be
cuspidal too, contradicting residue-characteristic
point/cusp separation from D.25.
Consequently every prime index \(p\)
which occurs as a residue characteristic
of a place outside \(S\) has ordinary
primewise image. In particular the set
of cuspidal prime indices is finite,
being contained in the finite set of
rational primes lying below \(S\).

D.49 now applies over the original field
\(K\): it proves that the ordinary primewise
points form a finite set and that the full
section is an ordinary \(K\)-point section.
The extension \(L/K\) was used only to
establish a finite set of exceptional
prime indices; no descent of a conjugating
element is required. The uniqueness of
the common point in every pro-prime
quotient is the global point/cusp
separation proved in D.25.

If there is no ordinary primewise image,
all images are cuspidal. Their underlying
points lie in the finite boundary, so
D.25 applies directly and makes the full
section geometric. Its ordinary alternative
is excluded by point/cusp separation.
The full section therefore belongs to a
single global cusp packet. Its primewise
images are exactly the projections of that
packet section; D.24 also supplies the
corresponding M1 uniqueness.
\end{proof}

\begin{statement}{D.52}{Proved}{Nongeometric sections at every nilpotent Fermat stage}
Let \(p\in\{3,5\}\), \(r\ge0\), \(n=p^r\),
and let \(U_n\) be the affine Fermat curve
\[
             u^n+v^n=1,\qquad uv\ne0 .
\]
For \(r=0\) this is the tripod.
Use the rational unit tangent at \((u,v)=(0,1)\),
and write \(N=\pi_1(U_{n,\overline{\mathbb Q}})\).
There is a section of its full pronilpotent
arithmetic quotient whose image in
\(N/\overline{\gamma_{d+1}N}\) is neither a
rational-point image nor the image of any
cuspidal packet, for every \(d\ge1\).
All these images are compatible.

The same construction gives a section of
the full maximal pro-\(p\) geometric quotient
of \(U_n\); hence its images in every
pro-\(p\) derived quotient also remain
nongeometric. This proves the natural
lower-central, nilpotent assertion in the
Fermat hunch of Task D.5 for both specified
prime-power towers. It does not identify the
pronilpotent quotient with a full profinite
solvable quotient of fixed derived length.
Here a cusp image means the image of an
actual tangential packet; in an abelian
quotient, merely normalizing a peripheral
subgroup would be a weaker condition.
\end{statement}

\begin{proof}
Put \(G_p=G_{\mathbb Q,\{p,\infty\}}\).
We first record the uniform central vanishing
\[
 H^2_{\rm cts}(G_p,\mathbb Z_p(j))=0
                      \qquad(j\in\mathbb Z).                    \tag{D.52.1}
\]
For \(p=5\) this is the all-weight argument
of D.12. For \(p=3\), take
\(L=\mathbb Q(\zeta_3)\).
Its Minkowski bound is
\((2/\pi)\sqrt3<2\), so it has class
number one. There is a single prime above
three and no real place. The invariant-sum
sequence gives
\(\operatorname{Br}(\mathcal O_L[1/3])[3]=0\).
Kummer theory, as in D.1, therefore gives
\[
 H^2(G_{L,\{3\}},\mathbb F_3(1))=0.
 \]
All residual Tate twists are trivial over
\(L\). Restriction to \(L\), followed by
averaging over its degree-two Galois group,
proves the residual vanishing over
\(\mathbb Q\) in every weight.
Coefficient induction proves it modulo
\(3^m\), for every \(m\). Finite-coefficient
cohomology at this fixed ramification set
is finite; its first-cohomology tower is
Mittag--Leffler. Jannsen's comparison then
gives (D.52.1), including its possible
derived-limit term. Thus no regular-prime
assertion is being invoked without its
residual and integral checks.

Let \(F_p\) be the full free pro-\(p\)
geometric tripod group. Its action through
the unit tangent factors through \(G_p\)
by tame specialization. In its abelian
quotient choose the continuous cocycle
\[
              (n\kappa_p(p),\,2n\kappa_p(p)).
                                                               \tag{D.52.2}
\]
Every successive lower-central kernel
of \(F_p\) is finite free of scalar Tate
weight, so (D.52.1) removes every central
lifting obstruction. Choose lifts which
project to the preceding cocycles exactly;
the adjustment by a geometric conjugator
is possible because each group projection
is surjective. Their compatible inverse
limit is an actual continuous \(F_p\)-cocycle:
every finite quotient of a pro-\(p\) group
is nilpotent.

The covering \(U_n\to U_1\) is
\((u,v)\mapsto t=u^n\), with
\(1-t=v^n\). Its geometric subgroup
has quotient \((\mathbb Z/n)^2\).
Consequently the maximal pro-\(p\)
group of \(U_n\) is canonically
\[
 K_r=\ker\bigl(F_p\longrightarrow
                         (\mathbb Z/n)^2\bigr).
                                                               \tag{D.52.3}
\]
To justify the completion comparison,
a normal subgroup of \(p\)-power index
in the covering's geometric group has
a core of \(p\)-power index in the
ambient group: intersect its conjugates
under the finite \(p\)-group of deck
transformations. Thus the induced
pro-\(p\) topology is the full
pro-\(p\) topology of the subgroup.
The log parameter \(t=u^n\), with unit
leading coefficient, identifies the
chosen arithmetic actions as well.
The cocycle in (D.52.2) is zero
modulo \(n\) pointwise. All the exact
lifts above therefore take values in
\(K_r\). They give the promised full
pro-\(p\) section of the covering.

Its image in the ambient tripod
abelianization is (D.52.2).
This lies on none of the three
peripheral lines
\((a,0),(0,a),(-a,-a)\).
It cannot be the image of a rational
point either. The kernel of
\(\mathbb Q^\times\to
H^1(\mathbb Q,\mathbb Z_p(1))\)
is \(\{\pm1\}\) for odd \(p\), by
the prime valuations. Such a point
would therefore require
\[
       t=\epsilon p^n,\qquad
       1-t=\delta p^{2n},\qquad
       \epsilon,\delta\in\{1,-1\}.
 \]
But \(p^{2n}-p^n>1\), so none of
these four possibilities can sum
to one. A rational point or a
cuspidal packet of \(U_n\) maps
to the corresponding point or
packet on the tripod. Hence its
pro-\(p\) section cannot be the
one just constructed. This argument
already takes place in an abelian
quotient, so it survives every
lower-central or derived truncation.

Finally the maximal pronilpotent
quotient of a profinite group is
canonically the product of its
maximal pro-prime quotients. More
precisely, here
\[
 N^{\rm nil}=\prod_\ell N^{(\ell)},
 \qquad
 N/\overline{\gamma_{d+1}N}
   =\prod_\ell
       N^{(\ell)}/\overline{\gamma_{d+1}N^{(\ell)}} .
                                                               \tag{D.52.4}
\]
Every finite nilpotent group is the
product of its characteristic Sylow
subgroups; applying this to all finite
quotients proves both identifications,
including their topology and arithmetic
equivariance. Use the constructed
section in the \(p\)-factor and the
chosen unit cusp section in every
other prime factor. The product is
continuous in the product topology
and is an actual section of the
pronilpotent arithmetic extension.
Projection to the \(p\)-abelian
factor retains (D.52.2), proving
the stated nongeometric property
at every \(d\ge1\).
\end{proof}

\begin{statement}{D.53}{Proved}{The full-group mixed fibre above the trivial cubic class}
In Task D.9, take the trivial cubic class
\(c\in\mathbb Q^\times/\mathbb Q^{\times3}\).
Then the subset of
\(H^1(\mathbb Q,E_1[5])\) which extends to
a section of the full geometric group
\(F\) is exactly \(\{0\}\). This is
exactly the set of Kummer classes of
rational points and cusps in this fibre.
Every such full section is cuspidal.
The assertion concerns the full group;
it does not decide which further classes
extend only to the relative pro-five
quotient.
\end{statement}
\begin{proof}
The cyclic cover corresponding to the
trivial cubic class is
\[
 V:\ z^3=t(1-t),\qquad t(1-t)\ne0 .
\]
Its compactification has the three
rational cusps over \(0,1,\infty\).
The substitution
\[
       X=-4z,\qquad Y=4(2t-1)
\]
identifies its smooth compactification
with
\[
             E:\quad Y^2=X^3+16.
\]
The three cusps become
\(O,(0,4),(0,-4)\).
There is a nonconstant rational map
to the elliptic curve used in D.23:
\[
 \phi:E\longrightarrow E',\qquad
 E':\ Y'^2=X'^3-432=X_0(27),
 \qquad
 X'=\frac{X^3+64}{X^2},\quad
 Y'=\frac{Y(X^3-128)}{X^3}.                         \tag{D.53.1}
\]
Substitution gives the exact identity
\[
 Y'^2-X'^3+432
 =\frac{(X^3-128)^2}{X^6}(Y^2-X^3-16).
\]
Thus the map extends to a morphism of
smooth proper curves. Its \(X'\)-pole
divisor has degree six, so its degree
is three; it maps \(O\) to \(O\).
It is the degree-three isogeny with
kernel \(O,(0,4),(0,-4)\).
Only finiteness and nonconstancy of
the morphism are needed below.

D.23 proves from the checked
\cite[Corollary 8.8]{DE_Stoll2007} that
\[
 E'(\mathbb A_{\mathbb Q})_\bullet^{\rm f\text{-}cov}
       =E'(\mathbb Q),
 \qquad |E'(\mathbb Q)|=3.                         \tag{D.53.2}
\]
In particular \(E(\mathbb Q)\) is finite,
since it maps into that finite set
under a finite morphism. Good reduction
counts for \(Y^2=X^3+16\) give
\[
                |E(\mathbb F_5)|=6,\qquad
                |E(\mathbb F_7)|=9.
\]
Reduction injects prime-to-residue-
characteristic rational torsion. Using
both primes excludes all torsion except
a group of order dividing three. The
three displayed cusps therefore give
all of \(E(\mathbb Q)\), so \(V(\mathbb Q)\)
is empty. The identities and reduction
counts are executed in
\path{code/de_cubic_trivial_fibre.py}.

A full section of \(V/\mathbb Q\) is
Selmer at every finite place by D.56
with the class of all finite groups;
the real section theorem provides
the real component. Its image under
\(\phi\) gives an adelic point of \(E'\)
surviving every finite torsor. Indeed
evaluating any such torsor on the
one global section supplies the same
global twisting class at every place.
Equation (D.53.2) places the local
compactification points in the finite
inverse image of one point of
\(E'(\mathbb Q)\). The full
section-neighborhood finite-support
argument proved in D.23 then makes
the original section point or cusp
induced over \(\mathbb Q\). Since
\(V(\mathbb Q)=\varnothing\), it is
cuspidal.

Finally, a section of the tripod
whose cubic class is trivial lifts
to this actual cyclic cover: after
the coboundary conjugation which
trivializes its finite cubic
cocycle, its arithmetic image lies
in the covering subgroup. Conversely
a section of \(V\) maps to a tripod
section with trivial cubic class.
This is the ordinary finite-etale
covering correspondence, not a
pro-five approximation. D.19 shows
that every cusp packet of \(V\)
has zero \(E[5]\)-Kummer image,
because its underlying cusp point
has order three and five is
coprime to three. At least one such
packet exists, so zero both occurs
and is the only full-group class.
\end{proof}

\begin{statement}{D.54}{Proved}{The local section theorem for the full geometrically prosolvable group}
Let \(k/\mathbb Q_p\) be finite and let \(U/k\) be a smooth
geometrically connected hyperbolic curve.  Every continuous section of
\[
  \pi_1(U)/\ker\bigl(\pi_1(U_{\bar k})\to
              \pi_1(U_{\bar k})^{\mathrm{sol}}\bigr)
                  \longrightarrow G_k
\]
is induced by a \(k\)-rational point of \(U\), or belongs to the
packet of a \(k\)-rational cusp.  If \(U\) is proper, only the first
alternative occurs.  The quotient is the full prosolvable inverse
limit, with no bound on derived length.  The proof uses the checked
analytic constructions of the supplied local preprint and proves
that they retain solvable monodromy; it does not assume its main
section theorem or a prosolvable variant of that theorem.
\end{statement}
\begin{proof}
We first prove the assertion for a proper hyperbolic curve
$X/k$. The affine case is treated in Step 7 by its actual
finite solvable section neighborhoods.
\paragraph{Step 1: exact inputs and fixed two-prime set.}

Fix once and for all an odd prime $\ell\ne p$ and put
$\Sigma=\{p,\ell\}$.  This same set is used on every curve and in every
comparison below.  The primary external statements needed are:

\begin{enumerate}
\item Mochizuki--Tsujimura, \emph{Resolution of Nonsingularities,
Point-theoreticity, and Metric-admissibility for p-adic Hyperbolic
Curves}, 30 March 2024, Theorem A and Proposition 2.4(iv):
geometrically pro-$\Sigma$ finite covers provide stable-model domination \cite{MT2024}.
Theorem F(i), equivalently Proposition 2.4(vii), of \cite{MT2024} gives a unique
section-fixed point in the universal pro-$\Sigma$ pro-Berkovich space
of a proper hyperbolic curve when the inertia and cyclotomic hypotheses
there hold.  Both hold for $H=G_k$: tame inertia surjects onto
$\mathbb Z_\ell$, and the closure of the powers of the residue-field
cardinality is open in $\mathbb Z_\ell^\times$.
\item Amini--Baker--Brugall\'e--Rabinoff, \emph{Lifting harmonic
morphisms I}, arXiv:1303.4812v3, Corollary 6.21 and Theorem 7.4(2):
the automorphisms of a lifted cover are precisely the stabilizer of
its annular gluing data \cite{DE_ABBR2015}.  Theorem 7.7 permits subdivision of loops.
We construct invariant gluing data below; existence of an arbitrary
lift alone would not establish the required Galois group.
\item Hoshi, \emph{Conditional results on the birational section
conjecture over small number fields}, RIMS-1742, Proposition 1.7:
over $k$, every geometrically pro-$\mathcal C$ birational section is
geometric if the full formation $\mathcal C$ contains the $p$-groups.
We apply this to finite solvable groups \cite{DE_Hoshi2012}.  In particular the conclusion
concerns the whole section image in that quotient, not only its
further pro-$p$ image.
\end{enumerate}

The analytic calculations of Sections 3--6 of the supplied local
preprint \cite{A_OpenAILocal} are used with their stated domains: roots exist on the
\emph{entire} inverse collar, and the transfer is of an $H$-valued
point, including its residue-field embedding, as in \cite[Fact 5.1.3]{DE_Temkin2010}.  Steps 3--5 below
check the additional monodromy assertions and recall the relevant
analytic calculations.  No prosolvable version of the supplied
preprint is being quoted.

\paragraph{Step 2: solvable covers and a compatible fixed point.}

Write $\Delta=\pi_1(X_{\bar k})^{\mathrm{sol}}$ and
$\Pi=\pi_1(X)/\ker(\pi_1(X_{\bar k})\to\Delta)$, and fix an actual
section $s:G_k\to\Pi$.  Since a finitely generated profinite group has
only finitely many open subgroups of each prescribed index, the
characteristic open subgroups $V\subset\Delta$ are cofinal.  The open
subgroup
\[
                 H_V=V\,s(G_k)\subset\Pi
\]
defines a connected finite \'etale cover $X_V\to X$ over $k$.
It is geometrically connected, and its geometric cover is Galois
with group $\Delta/V$.

The geometrically prosolvable fundamental group of $X_V$ is exactly
$H_V$.  Here is the needed finite-group argument.  A finite solvable
Galois cover of $X_{V,\bar k}$ has a normal closure over $X_{\bar k}$.
Over the Galois cover $X_{V,\bar k}$ this closure is a compositum of
finitely many conjugate solvable extensions.  Its Galois group over
$X_{V,\bar k}$ is a subgroup of the product of their solvable groups;
its quotient over $X_{\bar k}$ is $\Delta/V$, also solvable.  The total
group is therefore solvable.  This proves the asserted identification
by the classification of finite covers.  The same argument applies to
finite function-field extensions, and will also be used below.

Burnside's two-prime theorem implies that every finite $\Sigma$-group
is solvable.  Thus the section of $H_V$ projects to an actual section
$s_V^\Sigma$ of the geometrically pro-$\Sigma$ group of $X_V$.
By the first external input, it has a unique fixed point $z_V$ in
the universal pro-$\Sigma$ pro-Berkovich space of $X_V$.

These points are compatible, for a precise reason that does not
compare unrelated quotients.  If $V'\subset V$, pullback of a
geometric finite $\Sigma$-Galois cover of $X_V$ to $X_{V'}$ again has
$\Sigma$ monodromy.  With the common geometric base point, this
constructs the natural map from the universal pro-$\Sigma$ tower of
$X_{V'}$ to that of $X_V$.  The map is equivariant for the two sections,
since both sections are induced by the same homomorphism $s$.
The image of $z_{V'}$ is consequently fixed by $s_V^\Sigma$ in the
\emph{target} tower.  Its uniqueness there gives
\[
                         z_{V'}\longmapsto z_V.
\]
In particular their projections
$x_V\in (X_V\otimes_k\mathbb C_p)^{\mathrm{an}}$ are compatible.
The geometric covers $X_{V,\bar k}$ are cofinal in the universal
prosolvable geometric tower.  Hence $(x_V)_V$ is a point of that
whole pro-Berkovich tower, fixed by $s(G_k)$.

If its projection $x_X$ is type 1, every $x_V$ is type 1, since a
finite morphism has type-1 fibres over type-1 points.  The descent
action of $G_k$ fixes $x_V$, so $x_V\in X_V(k)$, using
$\mathbb C_p^{G_k}=k$, as proved in \cite[Theorem 1, p.~176]{DE_Tate1967}.  The compatible rational points in the
geometrically Galois covers show directly that $s(G_k)$ stabilizes
one lift of a point of $X(k)$.  Horizontal inertia in an \'etale
tower is trivial, so this stabilizer is precisely its point-section
subgroup.

If $x_X$ is not type 1, its restriction to the rational function
field is a real-valued rank-one valuation extending the $p$-adic
valuation: on a smooth analytic curve over $\mathbb C_p$, the zeros
and poles of a nonzero rational function are type 1.  The compatible
points give such a valuation on the universal solvable \'etale
field, preserved by $s(G_k)$.  The remaining argument excludes this
second alternative.

\paragraph{Step 3: cyclic roots on stable edges, with invariant gluing.}

Let $C/\mathbb C_p$ be proper of genus at least two, with split stable
skeleton, and let $e$ be one of its node edges.  We claim that there
is a finite \'etale cover of $C$ with solvable Galois closure on whose
entire inverse node annulus every annular coordinate has an analytic
$\ell$-th root.  In fact the edge construction below uses towers
of cyclic degree-$\ell$ covers.

Suppose first that $e$ lies on a simple circuit, oriented cyclically.
At each circuit vertex let $P$ be the outgoing mark and $Q$ the
incoming mark on its smooth residue curve $D$.  Divisibility of
$\operatorname{Pic}^0(D)$ gives a divisor $E$ such that
$P-Q+\ell E$ is principal.  The equation
\[
                          y^\ell=f,
       \qquad \operatorname{div}(f)=P-Q+\ell E
\]
defines a connected cyclic residue cover, totally ramified just at
$P,Q$.  Equip it with the generator $\tau:y\mapsto\zeta_\ell y$.
At every noncircuit vertex take $\ell$ identity copies, cyclically
permuted by $\tau$.  Circuit edges have one lift of expansion factor
$\ell$; every other edge has $\ell$ lifts of factor one.  Match the
unramified marked fibres by $C_\ell$-equivariant bijections.  Thus all
the data carry one coherent $C_\ell$ action.  The covering is
connected, and the local Riemann--Hurwitz equality
\[
  2g(D')-2=\ell(2g(D)-2)+2(\ell-1)
\]
shows that its augmented graph ramification is zero.  There are no
punctures, so a tame lift is \'etale.

The lift can be chosen $C_\ell$-equivariant.  On an orbit of
unramified edges, choose one gluing isomorphism and transport it
under $\tau$; the orbit has length $\ell$, so the last compatibility
is automatic.  For a circuit edge, its two branch coordinates can
be taken to have actions
\[
              u\longmapsto\zeta_\ell u,
       \qquad v\longmapsto\zeta_\ell^{-1}v.
\]
The inverse characters follow from the respective orders congruent
to $1$ and $-1$ modulo $\ell$ of $f$; multiplying a local Kummer
root by an integral power of the base parameter does not change
its deck character.  On the annular overlap the gluing has $uv=c$;
it commutes with these actions.  More intrinsically, the local
degree-$\ell$ annular deck group is $\mu_\ell$, and the two
restrictions are the same element when expressed in the common
oriented annular coordinate.  Thus the gluing data are fixed by
$\tau$.  Corollary 6.21 and Theorem 7.4(2) of the second external
input now lift the action.  Since the connected cover has degree
$\ell$, it is a $C_\ell$-Galois cover.  Subdividing loops inserts
the map $u\mapsto u^\ell$ and preserves these inverse characters.
The expansion factor on every annulus above $e$ is $\ell$.

If $e$ is separating, each side of the cut stable graph contains
a circuit or a positive-genus vertex.  A genus-zero tree side with
$N$ vertices would have total valence $2N-1$, whereas stability
requires at least $3N$.  On each side choose a connected cyclic
degree-$\ell$ unramified cover, using respectively a nonzero
homomorphism of the graph fundamental group to $C_\ell$, or a
nonzero $\ell$-torsion line bundle on the positive-genus residue
component.  Complete the residue and graph data by identity copies
and equivariant matching.  The $\ell$ bridge preimages can be
matched equivariantly between the two connected sides.  All gluing
edges are then free $C_\ell$-orbits, so the preceding transport
argument gives a connected cyclic \'etale cover of $C$.
Every bridge preimage is now nonseparating: the other $\ell-1$
bridges still join the two connected sides.  Apply the previous
construction to the finitely many bridge preimages and take a
connected component of the fibre product.  Its Galois closure has
an $\ell$-group as group, by the finite-group argument in Step 2.

The root assertion holds on whole annuli.  An analytic annular unit
has a factorization $aS^d(1+h)$ with $|h|<1$ everywhere on the open
annulus and a single dominant Laurent monomial.  Expansion factor
$\ell$ makes $d$ divisible by $\ell$; constants have roots in
$\mathbb C_p$ and $(1+h)^{1/\ell}$ converges on every compact
subannulus.  These roots agree on overlaps.  The inverse-skeleton
property of the tame lift identifies the full inverse node annulus
with the indicated node annuli, including their attached open
balls.  This proves the claim for one edge, and finite fibre
products give it for finitely many edges simultaneously.

For small collars around arbitrary algebraic points, perform the
blowups used in the supplied preprint to insert target node annuli.
Apply stable-model domination with our fixed $\Sigma=\{p,\ell\}$,
so its geometric covering group is solvable.  Choose an irrational
level strictly inside each target node annulus.  A source model
vertex gives a rational value to the target parameter; a smooth
disk tube also gives a rational constant value if it maps into that
target annulus.  Thus the full inverse irrational level lies in the
finitely many source node annuli.  Compactness gives an open collar
about the level whose full inverse lies in those annuli.  Apply the
edge construction to these finitely many annuli.  It follows that
the collar-root lemma of the supplied preprint holds with
\emph{solvable Galois closure} throughout.

\paragraph{Step 4: equivariant patching and arbitrary-order auxiliary roots.}

Let $C\to X_{\bar k}$ be a finite Galois cover with specified finite
solvable group $G$.  At every branch point $b$, choose a sufficiently
small disk with parameter $t_b$ such that each inverse disk has the
normal form $t_b=z^{e}$.  This uses characteristic zero and analytic
roots of units after shrinking; it does not require $e$ prime to $p$.
Suppose a finite cover $Y\to X_{\mathbb C_p}$ supplies an analytic
$m_b$-th root of $t_b$ on the entire inverse collar, with each
$e\mid m_b$.  Outside the inner disks the pullback of $C$ is a
$G$-torsor.  On each inverse disk replace it by the trivial
$G$-torsor.  On the overlap the equations split into all linear
factors, so the exterior torsor is trivial there.  A choice of one
sheet on each connected overlap identifies it \emph{equivariantly}
with $G$; the other sheets are its $G$-translates.  Glue through
these $G$-equivariant trivializations.

The result is a finite \'etale $G$-torsor on $Y^{\mathrm{an}}$.
There is no compatibility between distinct branch disks because
they are disjoint.  Proper analytic GAGA \cite[Theorem A.1(ii)]{DE_Poineau2010} algebraizes both the
cover and its $G$ action; \cite[Fact 5.1.5]{DE_Temkin2010} detects etaleness after analytification.  A connected component $V$ is a Galois
cover of $Y$ with group a subgroup of $G$, hence solvable.  If
$Y\to X$ has solvable Galois closure, so does $V\to X$, by
Step 2.  Descent of a finite \'etale cover along the algebraically
closed extension $\bar k\subset\mathbb C_p$ preserves the action
by full faithfulness.  No descent of the exterior analytic map is
needed.

For arbitrary positive $m$, the auxiliary construction of Section 5
of the supplied preprint also has solvable Galois closure.  Its
first curve is
\[
                  Z:\quad y^\ell=\frac{T(T-1)}{T-\alpha}.
\]
This is a cyclic degree-$\ell$ cover, with horizontal indices
$\ell$ at $0,\alpha,1,\infty$.  Its stable skeleton consists of
two rational vertices joined by $\ell$ parallel edges, each of
length $a=v_p(\alpha)$, and $g(Z)=\ell-1$.  Let $J$ be its
Jacobian and $i:Z\to J$ an Abel--Jacobi map.  The cover
\[
                       Z'=Z\times_{i,J,[m]}J
\]
is a torsor under the abelian finite group $J[m]$ over $\bar k$.
Every connected component has abelian Galois group over $Z$.
Its normal closure over the $T$-line therefore has solvable group,
an extension of an abelian subgroup of a product by $C_\ell$.
Its horizontal ramification indices still divide $\ell$.

For completeness, the roots supplied by this cover are the required
global analytic roots on all inverse sheets.  The totally degenerate
Jacobian has toric uniformization $\mathbb T^{\mathrm{an}}/\Lambda$ \cite[Sections 4.1--4.2 and 4.5.2]{DE_BakerRabinoff2015}.
The Abel--Jacobi map on each of the $\ell$ open annuli lifts to
$\mathbb T^{\mathrm{an}}$.  The cycle $e_j-e_{j'}$ gives a character
whose pullback has slope one by \cite[Proposition 6.1 and Corollary 6.6]{DE_BakerRabinoff2015}, hence is
\[
                 u_j=c_j T(1+h_j),\qquad
        |h_j|\le p^{-\min\{v_p(T),a-v_p(T)\}}.
\]
On $Z'$ the character has an analytic $m$-th root.  Indeed local
torus lifts of the division point have $m$-th power equal to the
original lift times a period.  Decompose by its coset in
$\Lambda/\Lambda^m$, choose one period representative, and adjust
local lifts to that representative.  Their differences on overlaps
are $m$-torsion periods, hence trivial because $\Lambda$ is torsion
free.  The adjusted lifts therefore glue on every clopen piece.
If
\[
        D>v_p(m)+\frac1{p-1},\qquad a>2D,
\]
then $\exp(m^{-1}\log(1+h_j))$ converges on
$D<v_p(T)<a-D$.  Dividing the character root by this function and
an $m$-th root of $c_j$ produces a root of $T$ on the whole inverse
middle collar.  Pullback by $T=t_b/\beta_b$ and finite composita
retain both this property and solvability of the Galois closure.

We also use the following horizontal ramification fact.  Over the
algebraically closed characteristic-zero residue field of a
horizontal valuation, the local extensions of the completed
function field have the form $\bar k((t^{1/e}))$.  Consequently
indices dividing $\ell$ remain so under composita, conjugates,
normal closures, normalized pullback, and composition on top with
an \'etale cover.  Thus the normal closures just used do not lose
the horizontal index bound.

\paragraph{Step 5: completion contains every solvable birational extension.}

Put $K=k(X)$ and $K_0=\bar k(X)$.  Inside a fixed $K^{\mathrm{sep}}$,
let $B$ be the maximal solvable extension of $K_0$, and let $L\subset B$
be the universal geometrically solvable \'etale field of $X$.
Both are normal over $K$.  Fix any prolongation of a real rank-one
valuation extending $v_p$, embed their completions in
$E=\widehat{K^{\mathrm{sep}}}$, and put $H=\widehat L\subset E$.
In particular $\mathbb C_p\subset H$.  We prove
\[
                              B\subset H.                 \tag{*}
\]

First let $F/K_0$ be finite Galois with solvable group and with every
horizontal index dividing $\ell$.  Its branch disks can be chosen
to avoid the analytic point defined by $K_0\hookrightarrow H$:
this point projects to the generic point of $X_{\bar k}$, and hence
is distinct from each algebraic branch point.  Step 3 gives an
\'etale cover $Y\to X$ with solvable Galois closure and all the
necessary collar roots.  Equivariant patching in Step 4 gives
an \'etale cover $V\to X$ with solvable Galois closure and an
exterior analytic map $V\to C$, where $C$ has function field $F$.
Its function field embeds in $L$, and hence its $H$-valued generic
point transfers through the exterior analytic map to an $H$-valued
generic point of $C$.  On completed residue fields this is a
$K_0$-embedding $F\hookrightarrow H$.  Normality of $F/K_0$ makes
its image equal to the chosen field $F\subset E$.  This is an
inclusion of the field, not only the existence of an abstract
isomorphic copy.

Let $L_\ell^{\mathrm{sol}}$ be the compositum of all these finite
solvable Galois extensions with indices dividing $\ell$.  The first
stage proves $L_\ell^{\mathrm{sol}}\subset H$.

Now take \emph{any} finite solvable Galois extension $F/K_0$.
At a branch point $b$ let $m_b$ be the least common multiple of its
indices.  The auxiliary covers of Step 4, pulled back by the
local rational parameters and combined by a compositum, give
$W\to X_{\bar k}$ with solvable Galois closure, horizontal indices
dividing $\ell$, and all analytic $m_b$-th roots on the full inverse
collars.  Equivariant patching produces $V\to W$ \'etale with
solvable group on each connected component.  Thus $V\to X$ has
solvable Galois closure and horizontal indices dividing $\ell$.
Its field embeds in $L_\ell^{\mathrm{sol}}\subset H$.  The same
$H$-valued point transfer and normality again give $F\subset H$.
Every element of $B$ lies in one of these finite solvable Galois
extensions.  This proves (*), with no bound on the ramification
orders or derived lengths of $F$.

Restriction of decomposition groups is consequently an isomorphism
\[
                  D_w(B/K)\xrightarrow{\sim}D_w(L/K).       \tag{**}
\]
Surjectivity follows by lifting an automorphism and adjusting it by
$\operatorname{Gal}(B/L)$, which acts transitively on prolongations;
the infinite case follows from finite-level transitivity and
compactness.  An element of the kernel is an isometry after
normalization $w(p)=1$ and extends to $\widehat B$.  It fixes
$H=\widehat L$ pointwise.  But (*) implies
$B\subset H\subset\widehat B$, so it fixes $B$.  This proves
injectivity; compactness makes the inverse continuous.

\paragraph{Step 6: excluding a section in a rank-one decomposition group.}

If $s(G_k)$ fixes a rank-one valuation in the \'etale tower,
(**) gives a continuous lift
\[
             \widehat s:G_k\longrightarrow\operatorname{Gal}(B/K)
\]
into its birational decomposition group.  By Hoshi's Proposition 1.7
the whole image also lies in a horizontal decomposition group at
some $a\in X(k)$.  Let
$M=B^{\widehat s(G_k)}$.  The restrictions $v_1,v_2$ of the vertical
and horizontal valuations to $M$ have rank one, and
\[
                         v_1(p)=1,\qquad v_2|_k=0.
\]
Rank one is preserved here because all the field extensions from
$K$ are algebraic, so their value-group quotients are torsion.

The required Hensel argument does not presume that $M$ is henselian.
The field $B$ is quadratically closed.  In fact for any $b\in B$,
choose a finite solvable Galois extension $F/K_0$ containing $b$;
adjoining square roots of all its conjugates is an abelian extension
of $F$ whose normal closure over $K_0$ is solvable.  Thus $\sqrt b\in B$.
Every quadratic extension of $M$ is therefore contained in $B$.
Since all of $\operatorname{Gal}(B/M)=\widehat s(G_k)$ stabilizes
each selected valuation, every such quadratic extension has a
unique prolongation of $v_i$.  Equivalently these valuations are
2-henselian.  More concretely, if $u\in M$ satisfies
\[
                             v_i(u-1)>2v_i(2),
\]
Hensel's lemma gives a square root in the henselization of $M$ at
$v_i$.  It also belongs to $B$.  On the quadratic extension
$M(\sqrt u)/M$, the full Galois group is its decomposition group,
so its intersection with that henselization is $M$.  Hence
$\sqrt u\in M$.  This proves exactly the square-root assertion
used next, including $p=2$ and nondiscrete valuations.

Choose a nonsquare $d\in k^\times$.  It remains nonsquare in $M$,
since the projection $\widehat s(G_k)\to G_k$ is onto.  Let $t\in K$
be a uniformizer at $a$, choose an integer $N>v_1(t)$ and set
$z=p^N/t$.  Then $v_1(z)>0$ and $v_2(z)<0$.  Put
\[
                    q_r=d+(1-d)\frac{z^r}{1+z^r}.
\]
The denominator is nonzero, and
\[
\begin{aligned}
 v_1(q_r/d-1)&=v_1((1-d)/d)+r v_1(z)\longrightarrow+\infty,\\
 v_2(q_r-1)&=v_2(d-1)-r v_2(z)\longrightarrow+\infty.
\end{aligned}
\]
For large $r$, the preceding square-root assertion gives
$\sqrt{q_r/d}$ and $\sqrt{q_r}$ in $M$.  Their quotient gives
$\sqrt d\in M$, a contradiction.  Thus no section fixes a
rank-one valuation in $L$.  Combined with Step 2, this proves
the proper-curve assertion.

\paragraph{Step 7: affine passage.}

For an affine hyperbolic curve $U$, the following elementary
cofinality fact supplies proper hyperbolic compactifications inside
its solvable tower.  In the free pro-$r$ geometric fundamental
group, each peripheral element has infinite order, for any fixed
prime $r$: with at least two punctures it has nonzero primitive
abelian class, and with one puncture it is the nontrivial product
of surface commutators.  Its orders in finite $r$-group quotients
are therefore unbounded.  Choose a quotient in which all the
finitely many peripheral orders exceed a common $M$, and take a
characteristic refinement.  If its degree is $d$ and the boundary
has $n$ points, the compactification has
\[
     2g'-2=d\left(2g-2+n-\sum_{j=1}^{n}\frac1{e_j}\right).
\]
Since $2g-2+n>0$, sufficiently large $e_j$ make $g'\ge2$.
The same construction over each finite solvable section-neighborhood,
followed by its solvable normal closure over $U$, gives cofinally
many such neighborhoods.  The section-induced arithmetic twists
are the covers defined by $V s(G_k)$; hence they are over $k$ and
the section lifts to them.

The section on each such affine cover induces one on the
geometrically prosolvable group of its proper compactification.
The proper-curve assertion gives a $k$-rational point there.
The compact spaces of these rational points have nonempty inverse
limit: any finite collection is dominated by a further section
neighborhood, whose rational point supplies compatible images.
The resulting point in the compactified solvable tower is
section-fixed.  Its stabilizer is the ordinary point decomposition
group if the base point lies in $U$, and the cusp decomposition
group otherwise.  This gives exactly point/cusp origin.

\end{proof}


\begin{statement}{D.55}{Proved}{The Legendre-relative prosolvable Fermat assertion and the three-power prosolvable tower}
For an odd integer \(n\ge3\), let
\[
 U_n=\{X^n+Y^n+Z^n=0,\ XYZ\ne0\},\qquad
 N=\pi_1(U_{n,\overline{\mathbb Q}}),\qquad
 \Pi_n=\pi_1(U_n).
\]
For an arithmetic-stable open normal subgroup \(R\subset N\), put
\(R^{(\infty)}=\bigcap_{m\ge1}R^{(m)}\), with closed derived groups.
Then the following statements hold.
\begin{enumerate}
\item If \(\ell\ge5\) is prime and divides \(n\), let
\(R_\ell=\ker(N\to\operatorname{SL}_2(\mathbb F_\ell))\)
for the Legendre family pulled back by \(\lambda=-(X/Z)^n\).
Every section of \(\Pi_n/R_\ell^{(\infty)}\to G_{\mathbb Q}\)
is cuspidal.  In particular the Legendre-relative prosolvable
clause of the Fermat hunch holds at every \(5^r\) level, using
\(\ell=5\).
\item If \(n=3^r\), \(r\ge1\), every section of
\(\Pi_n/R^{(\infty)}\to G_{\mathbb Q}\) is cuspidal for every
such \(R\).  This includes the plain geometrically prosolvable
quotient by taking \(R=N\).
\item Consequently the noncuspidal full metabelian families in
D.20 have no compatible lift to the plain prosolvable inverse
limit on the three-power tower, or to the indicated
Legendre-relative prosolvable inverse limit on the five-power
tower.  No first obstructed finite derived length is asserted.
\end{enumerate}
More generally, over every finite extension of \(\mathbb Q_p\),
D.54 implies the point-or-cusp conclusion for the relative
prosolvable quotient above any fixed arithmetic-stable finite
geometric Galois cover.  The proof below includes the arithmetic
twist which makes that assertion exact.
\end{statement}
\begin{proof}
\paragraph{Step 1: the exact arithmetic twist.}

Let $U/k$ be hyperbolic, $\Pi_U=\pi_1(U)$ and
$N=\pi_1(U_{\bar k})$.  Let $R\subset N$ be open and normal,
stable under $\Pi_U$-conjugation, and put
\[
 R^{(\infty)}=\bigcap_{m\ge1}R^{(m)},\qquad
 \Pi_{U,R}^{\mathrm{rel,sol}}
      =\Pi_U/R^{(\infty)}.
\]
All derived subgroups are closed.  The first intersection is the
kernel of the maximal prosolvable quotient of $R$: every finite
solvable image kills some $R^{(m)}$, and the quotient by the
intersection embeds in the product of the finite-derived quotients.

Fix an actual section $s:G_k\to\Pi_{U,R}^{\mathrm{rel,sol}}$.
The subgroup
\[
             H=R^{\mathrm{sol}}\,s(G_k)
                 \subset\Pi_{U,R}^{\mathrm{rel,sol}}
\]
is open, projects onto $G_k$, and has intersection
$R^{\mathrm{sol}}$ with the geometric subgroup.  Its preimage in
$\Pi_U$ defines a finite \'etale cover $Y_s\to U$ over $k$, with
constant field exactly $k$.  Geometrically it is the Galois cover
corresponding to $R$.  Since its full geometric fundamental group
is $R$, its geometrically prosolvable fundamental group is exactly
$H$.  In particular the section $s$ is an actual prosolvable section
of $Y_s/k$; no extension of $k$ and no descent of a geometric
conjugator is required.

For $k/\mathbb Q_p$ finite, D.54 makes this section
point or cusp induced on $Y_s$.  Its composite in
$\Pi_{U,R}^{\mathrm{rel,sol}}$ is therefore point or cusp induced
on $U$.  This proves the local relative assertion for \emph{every}
fixed finite geometric cover, including nonsolvable ones.

\paragraph{Step 2: finite support using only solvable covers.}

We use the following primary version to avoid any implicit real-place
assumption.  \cite[Theorem 4.1]{DE_Hoshi2012} applies to a locally
geometric pro-$\mathcal C$ section over a number field when
$\mathcal C$ contains all finite solvable groups.  It is geometric
if its underlying local compactification points lie in one finite
closed subscheme at all finite places outside a density-zero set.
\cite[Corollary 4.2]{DE_Hoshi2012} in particular proves geometricity when the curve
has a nonconstant map to an abelian variety $A$ with $A(k)$ and
$\operatorname{Sha}(A/k)$ finite.  Local geometricity in this statement is at
finite places.  The proof passes to a totally imaginary extension,
uses Stoll's Theorem 8.2 at every solvable section neighborhood,
and descends by Hoshi's Lemma 1.6.

Here is the cohomological content of the special case used below.
The induced section on $A$ is a class in
$H^1(k,T A)$, and its finite local restrictions are point Kummer
classes.  Therefore it lies in the inverse-limit finite Selmer
group.  The quotient of the finite-level Selmer group by
$A(k)/nA(k)$ is killed uniformly, using finiteness of $\operatorname{Sha}$ and
the fact that real local groups have exponent at most two.
Taking the inverse limit over all $n$ kills that quotient.  Hence
the class belongs to $A(k)$ if $A(k)$ is finite.  Full local Tate
Kummer injectivity makes each underlying local curve point lie in
the finite fibre above this one global point of $A$.  Stoll's
finite-support theorem \cite[Theorem 8.2]{DE_Stoll2007}, with the
finite-place correction in \cite{DE_StollErratum}, then supplies rational points on every
solvable section neighborhood.  Their compatible inverse limit
gives a point or cusp decomposition group.  Thus full finite
non-solvable torsors are not required in this argument.

For the modular target used below, Stoll's Corollary 8.8 and its
proof \cite{DE_Stoll2007} provide exactly the needed abelian variety: every
positive-genus $X_0(M)/\mathbb Q$ has a nontrivial quotient of its
Jacobian that is modular of analytic rank zero.  Its rational
points and Tate--Shafarevich group are finite.  The map of the
modular curve to that quotient is nonconstant, since the curve's
Abel--Jacobi image generates its Jacobian.  Thus any hyperbolic
curve with a nonconstant map to such an $X_0(M)$ satisfies the
preceding prosolvable finite-support principle.

\paragraph{Step 3: the Legendre-relative Fermat quotient.}

Let $\ell\ge5$ be prime, let $n$ be odd and divisible by $\ell$,
and let
\[
 U_n:\quad X^n+Y^n+Z^n=0,\qquad XYZ\ne0,
 \qquad\lambda=-(X/Z)^n.
\]
Pull the Legendre family $y^2=x(x-1)(x-\lambda)$ back to $U_n$.
Let $N=\pi_1(U_{n,\overline{\mathbb Q}})$, and set
\[
 R_\ell=\ker\bigl(N\longrightarrow
                         \operatorname{SL}_2(\mathbb F_\ell)\bigr).
\]
This geometric image is the whole $\operatorname{SL}_2$.
Indeed the tripod image is generated by its two standard
unipotent matrices and is $\operatorname{SL}_2(\mathbb F_\ell)$;
$N$ is normal in the tripod group with abelian quotient, whereas
$\operatorname{SL}_2(\mathbb F_\ell)$ is perfect for $\ell\ge5$.
The residual kernel is stable under the arithmetic group.

\medskip
\noindent\textbf{The required conclusion.}
Every section of
\[
 \Pi_{U_n}/R_\ell^{(\infty)}\longrightarrow G_{\mathbb Q}
                                                        \tag{D.55.1}
\]
is cuspidal.  In particular this settles the Legendre-relative
part of the Fermat hunch for the $5^r$ tower, using $\ell=5$.

To prove it, take a section $s$ of (D.55.1).  Step 1 makes every
finite local restriction point or cusp induced.  Consequently
every local residual calculation in D.23 applies to its actual
Legendre representation $\bar\rho_s$.  For clarity the divisibility
used there is $\ell\mid n$:
\[
  v_p(j)=8v_p(2)-2n|v_p(X/Z)|,
       \qquad [q(s)]_\ell=[256^{-1}]_\ell
\]
at a cusp, with the analogous formula using $Y/Z$ when necessary.
After one global quadratic twist, the representation is unramified
outside $\{2,\ell\}$, has conductor at most two away from $\ell$,
is finite flat of Serre weight two at $\ell$, and has determinant
$\bar\chi_\ell$.  It is odd because complex conjugation has
determinant $-1$.  An irreducible representation would therefore
come from a weight-two cusp form of level one or two, whose spaces
vanish.  The representation is reducible, exactly as in the
independently audited D.23 proof.  This argument uses the full
cuspidal constant $256^{-1}$; it does not presume that the cusp
representation splits.

Its invariant line gives a section-supported component $Z$ of
the finite cover of $U_n$ parametrizing a cyclic subgroup of order
$\ell$ of the Legendre curve.  Its geometric fundamental group
$R_B\subset N$ is the inverse image of a Borel subgroup
$B\subset\operatorname{SL}_2(\mathbb F_\ell)$, up to conjugacy.
Since $B$ is solvable,
\[
                       R_B^{(\infty)}=R_\ell^{(\infty)}.    \tag{D.55.2}
\]
For the two inclusions, any solvable quotient of $R_B$ restricts
to a solvable quotient of $R_\ell$, while
$R_B/R_\ell^{(\infty)}$ is an extension of the prosolvable group
$R_\ell^{\mathrm{sol}}$ by the finite solvable group $B$ and is
itself prosolvable.  Thus the subgroup of (D.55.1) belonging to $Z$
is \emph{exactly} $\Pi_Z^{\mathrm{sol}}$.  The lifted section is
a prosolvable section of this actual curve, not just a section of
an unspecified relative quotient.

As in D.23, $Z$ has a nonconstant map to $Y_0(4\ell)$: quotient
the Legendre curve by one rational two-torsion point and combine
the image of its order-$\ell$ subgroup with the cyclic order-four
subgroup supplied by the other rational two-torsion point.
The map extends to the compactifications, and
$g(X_0(4\ell))=(\ell-3)/2>0$.  D.54 makes the
section of $Z$ locally geometric.  Step 2, with the rank-zero
modular abelian quotient, makes it globally geometric.  Its image
in (D.55.1) is therefore a point or cusp section of $U_n$.
The prime-exponent argument of D.23, or the coordinate power map
$U_n\to U_\ell$ followed by that argument, gives $U_n(\mathbb Q)=
\varnothing$.  Thus the section is cuspidal, proving the stated
conclusion.

\paragraph{Step 4: the three-power tower and the exact limit limitation.}

For $n=3^r$, the compactification of $U_n$ has a nonconstant map
to the Fermat cubic.  The explicit cubic transformation in D.23
identifies that cubic with $X_0(27)$, whose rational points and
Tate--Shafarevich group are finite.  Therefore Step 2 and the
local candidate already show that every section of
$\Pi_{U_n}^{\mathrm{sol}}$ is cuspidal.  The same is true of
$\Pi_{U_n}/R^{(\infty)}$ for \emph{any} arithmetic-stable finite
index geometric subgroup $R$: its arithmetic twist $Y_s$ from
Step 1 still has a nonconstant map to $X_0(27)$, so its
prosolvable section is geometric.  The absence of ordinary
$U_n(\mathbb Q)$ points was checked in D.23.

Consequently the noncuspidal full metabelian classes of D.20 on
the three-power tower have no compatible lift to the full
prosolvable inverse limit.  For the five-power tower they have no
compatible lift to the relative prosolvable quotient (D.55.1).
Neither assertion identifies a first finite derived length.
In particular, over $G_{\mathbb Q}$ the sets of continuous sections
need not be compact when ramification is unrestricted; existence
of an independent lift at each finite derived length cannot be
silently replaced by existence of one compatible family.

The phrase ``only cuspidal'' here concerns the Fermat curves.
It does not assert this for the tripod itself, which has ordinary
rational point sections.  Also, the proof of Step 3 requires
$\ell\mid n$ and does not assert the same conclusion for an
unrelated Legendre prime without an additional argument.

\end{proof}


\begin{statement}{D.56}{Proved}{The local section theorem for full formations containing two primes}
Let \(k/\mathbb Q_p\) be finite and let \(\mathcal C\) be a nonempty
class of finite groups closed under subgroups, quotients, and
extensions.  Suppose \(C_p,C_\ell\in\mathcal C\) for some prime
\(\ell\ne p\).  Every continuous section of the geometrically
pro-\(\mathcal C\) fundamental group of a smooth geometrically
connected hyperbolic curve over \(k\) is induced by a rational
point or belongs to a rational cusp packet.  For a proper curve
only the point alternative occurs.  In particular this holds
for the maximal geometrically pro-\(\{p,\ell\}\) quotient.
There is no bound on derived length; no one-prime assertion or
fixed finite-derived assertion is included.
\end{statement}
\begin{proof}
We first take \(X/k\) proper and hyperbolic; the affine case is
proved in the final paragraph.  We check every change to D.54,
including the case \(\ell=2\), rather than assuming that the
solvable proof is automatically compatible with an arbitrary
formation.

\paragraph{1. The exact finite-group closure argument.}

Write $\Sigma(\mathcal C)=\{q:C_q\in\mathcal C\}$.  Closure under
extensions implies closure under finite direct products.  It also
implies that every finite solvable group whose prime divisors lie
in $\Sigma(\mathcal C)$ belongs to $\mathcal C$, by a composition
series with cyclic prime-order factors.  By Burnside's theorem,
every finite $\{p,\ell\}$-group therefore lies in $\mathcal C$.
Conversely, if $G\in\mathcal C$, Cauchy's theorem and subgroup
closure imply that every prime dividing $|G|$ lies in
$\Sigma(\mathcal C)$.

If $Y\to X$ is a geometric finite Galois cover with group in
$\mathcal C$, and a finite Galois cover of $Y$ has group in
$\mathcal C$, its normal closure over $X$ again has group in
$\mathcal C$.  Over $Y$ the normal closure has group a subgroup of
a finite product of the conjugate covering groups; its quotient
is the group of $Y/X$.  The two closure properties prove the
assertion.  The same proof holds for finite function-field
extensions.  Thus, for a characteristic open subgroup $V$ of
$\Delta=\pi_1(X_{\bar k})^{\mathcal C}$ and an actual section $s$,
the group $V s(G_k)$ is exactly the geometrically pro-$\mathcal C$
fundamental group of its arithmetic cover.  There is no loss of
covers in making this identification.

Fix the single set $\Sigma_0=\{p,\ell\}$ on every cover.  All the
pro-$\Sigma_0$ covers used in D.54, Step 2, belong to $\mathcal C$.
\cite[Theorem F(i), Proposition 2.4(iv)]{MT2024}
therefore gives the unique fixed pro-$\Sigma_0$ point on every
proper hyperbolic section neighborhood.  Pullback of the target
$\Sigma_0$ tower gives the same equivariant map as in D.54;
uniqueness in the target proves compatibility.  This argument
does not require the finite quotient $\Delta/V$ to be solvable.
A compatible point of type one gives a rational point by Tate's
fixed-field theorem exactly as before.  It remains to exclude a
nonclassical fixed valuation.

\paragraph{2. All collar and patching groups remain in the formation.}

The collar construction of D.54, Step 3, uses cyclic groups of
order $\ell$, normal closures of towers of such groups, and
geometrically pro-$\Sigma_0$ stable-model domination.  Each belongs
to $\mathcal C$ by Paragraph 1.  This also covers the separating-edge
construction: its repeated cyclic covers and their normal closure
are $\ell$-groups.  In the positive-genus component case an
unramified cyclic cover of order $\ell$ exists because the residue
characteristic is $p\ne\ell$.  The orders at the two marked points
are congruent to $+1$ and $-1$ modulo $\ell$; this gives the
opposite inertia characters required by the invariant annular
gluing, also when $\ell=2$.

In the equivariant patching step let the specified exterior
Galois group be $G\in\mathcal C$.  The construction produces an
actual $G$-torsor on the auxiliary curve; each connected component
has group a subgroup of $G$.  Its normal closure over the base
curve has group in $\mathcal C$ by Paragraph 1.  Thus replacing
``solvable'' with ``in $\mathcal C$'' preserves the entire patching
argument, including the group action and the algebraization.

For a general finite Galois extension $F/\bar k(X)$ with group
$G\in\mathcal C$, every horizontal ramification index divides
$|G|$.  Choose each auxiliary root order $m$ to be the least common
multiple of the relevant indices.  Every prime divisor of $m$
lies in $\Sigma(\mathcal C)$.  The auxiliary cover
\[
 Z:\quad y^\ell=T(T-1)/(T-\alpha),\qquad
 Z'=Z\times_{J,[m]}J
\]
has cyclic initial group $C_\ell$ and finite abelian division
group $J[m]$.  Over $\bar k$ the latter is a product of cyclic
groups with prime divisors among those of $m$, and hence belongs
to $\mathcal C$.  Its normal closure, all conjugates, all finite
composita, and the patched covers stay in $\mathcal C$.  No prime
outside $\Sigma(\mathcal C)$ is introduced by the division step.
Horizontal indices still divide $\ell$, for the characteristic-zero
local-field reason in D.54.

\paragraph{3. The case $\ell=2$ in the auxiliary curve.}

If $\ell$ is odd, the auxiliary computation is literally D.54,
Step 4.  If $\ell=2$, then $p$ is odd.  The curve
\[
                    y^2=T(T-1)/(T-\alpha)
\]
is a smooth genus-one double cover, with four distinct branch
points $0,\alpha,1,\infty$.  For $a=v_p(\alpha)>0$, the same
semistable model has two rational components and two parallel
annular edges of length $a$.  One component is obtained from the
outer coordinate $T$, the other from $T=\alpha u$; their reduced
equations are $y^2=T-1$ and $y^2=-u/(u-1)$.  At the common end
there are two distinct points, because $p\ne2$.  The skeleton is
a circle of total length $2a$.

This model is semistable rather than stable as an unmarked
genus-one curve.  Stability is not used in the auxiliary division
construction.  Its Jacobian is the same totally degenerate
elliptic curve and has a toric uniformization
\cite[Sections 4.1--4.2 and 4.5.2]{DE_BakerRabinoff2015}
$\mathbb G_m^{\mathrm{an}}/\Lambda$.  The cycle $e_1-e_2$ is a
primitive torus character; its Abel--Jacobi pullback on either
edge has slope one, with the choice of edge orientation.  The
same calculation as in D.54 gives
\[
 u_j=c_jT(1+h_j),\qquad
 |h_j|\le p^{-\min\{v_p(T),a-v_p(T)\}}.
\]
After division by $[m]$, the period adjustment on the clopen
$\Lambda/\Lambda^m$ pieces and the torsion-freeness of $\Lambda$
produce an analytic $m$-th root of $u_j$ on every entire inverse
edge.  On the middle collar with
$D>v_p(m)+1/(p-1)$ and $a>2D$, the log--exp root of $1+h_j$
then gives an $m$-th root of $T$.  The finite groups here are
still a cyclic group of order two followed by a subgroup of
$J[m]$.  Thus the case $\Sigma_0=\{p,2\}$ needs no third prime.

\paragraph{4. The literal completion inclusion.}

Let $B/\bar k(X)$ now be the maximal pro-$\mathcal C$ birational
extension, let $L$ be the universal geometrically pro-$\mathcal C$
\'etale extension, and let $H=\widehat L$ for the proposed
nonclassical valuation.  The proof of D.54, Step 5, applies
verbatim with the group checks of Paragraphs 1--3.  First every
finite $\mathcal C$-Galois extension with all horizontal indices
dividing $\ell$ embeds literally in $H$, using the collar lemma,
equivariant patching, and the complete-field-valued point
transfer.  General finite $\mathcal C$-Galois extensions are then
treated by the auxiliary $m$-division covers.  Their normal
closures have indices dividing $\ell$ and group in $\mathcal C$,
so the first step supplies their function fields in $H$.
Normality of each prescribed $F/\bar k(X)$ makes its resulting
embedding have the literal chosen image $F\subset H$.
Consequently $B\subset H$.

Exactly as in D.54, the map of decomposition groups
$D_w(B/k(X))\to D_w(L/k(X))$ is an isomorphism: prolongation
transitivity gives surjectivity, while an element in the kernel
fixes the completion $H$ pointwise and therefore fixes $B$.
The local section lifts to a geometrically pro-$\mathcal C$
birational section.  \cite[Proposition 1.7]{DE_Hoshi2012} applies because
$p\in\Sigma(\mathcal C)$, and places its whole image in the
horizontal decomposition group of a rational point.

\paragraph{5. The final contradiction uses any prime of the formation.}

This is the only alteration to D.54, Step 6, besides the group
names.  Choose any prime $q\in\Sigma(\mathcal C)$, for example
$q=p$.  The field $B$ is $q$-closed: if $b\in B$, put $b$ in a
finite $\mathcal C$-Galois extension of $\bar k(X)$ and adjoin
$q$-th roots of all conjugates of $b$.  Since $\bar k$ contains
$\mu_q$, the resulting normal extension has an elementary
abelian $q$-group as its relative Galois group; its full group
belongs to $\mathcal C$.  Thus all $q$-th roots of $b$ already
lie in $B$.

Put $M=B^{s(G_k)}$.  Both the vertical valuation $w_1$ and the
horizontal valuation $w_2$ at the rational point have unique
prolongation to $B/M$, because their decomposition group is all
of $\operatorname{Gal}(B/M)=s(G_k)$.  Hence
\[
                          B\cap M^{h_i}=M\quad(i=1,2).
\]
Indeed a henselization is the fixed field of the decomposition
group in a separable closure; restriction of that decomposition
group onto $B$ is all of $\operatorname{Gal}(B/M)$.
If $u\in M$ and
\[
                       w_i(u-1)>2w_i(q),
\]
Hensel's lemma for $Z^q-u$ at $Z=1$ gives a root in $M^{h_i}$.
Every root lies in $B$ by $q$-closedness, so the root lies in $M$.
This argument requires no assumption that $\mu_q\subset k$ or
$\mu_q\subset M$.

Choose a uniformizer $d$ of $k$.  It is not a $q$-th power in $k$,
and cannot become one in $M$ because $M\cap\bar k=k$.  Let $t$ be
a horizontal uniformizer at the rational point.  Choose $N$ with
$Nw_1(p)>w_1(t)$ and set $z=p^N/t$.  Then $w_1(z)>0$ and
$w_2(z)<0$.  For large positive integers $r$, set
\[
                  u_r=d+(1-d)\frac{z^r}{1+z^r}.
\]
We have $w_1(u_r/d-1)\to+\infty$ and
$w_2(u_r-1)\to+\infty$.  Applying the preceding root criterion
to $u_r/d$ for $w_1$ and to $u_r$ for $w_2$ makes both elements
$q$-th powers in $M$.  Their quotient is $d$, a contradiction.
The nonclassical fixed valuation is therefore impossible.

\paragraph{6. The affine passage and scope.}

The affine compactification argument of D.54, Step 7, only needs
cofinal characteristic finite covers with all peripheral orders
arbitrarily large.  The pro-$p$ covers used there belong to
$\mathcal C$ because $C_p\in\mathcal C$.  Riemann--Hurwitz gives
hyperbolic compactifications; the proper result, common refinements,
and compactness give the compatible point-or-cusp decomposition
in the original affine quotient.  This completes the proof.

The result uses a full formation, hence unrestricted derived
length.  It neither contradicts the nongeometric one-prime
sections constructed in D.12/D.52 nor settles a fixed finite
derived stage.  Its most restrictive specialization is the
geometrically pro-$\{p,\ell\}$ quotient.  It makes no assertion
about a pro-$\ell$ quotient when $p\notin\Sigma$, or about a
geometrically pro-$p$ quotient with no second prime.
\end{proof}


\begin{statement}{D.57}{Proved}{Every nonzero full metabelian Fermat parameter fails at a finite derived stage}
Let $\ell\in\{3,5\}$, $n=\ell^r$, and let
$s_a$ be the actual full metabelian Fermat family of D.20,
parametrized by $a\in\mathbb Z_\ell^{d_n}$ with
$d_n=(n^2-1)(n^2+3)/4$.  Then:
\begin{enumerate}
\item The restrictions of the $s_a$ to $G_{\mathbb Q_\ell}$
are pairwise nonconjugate in the full metabelian group.  For
$a\ne0$, the local section is neither point nor cusp induced.
\item No $a\ne0$ has a lift to the plain geometrically
prosolvable quotient, even locally at $\ell$.  This conclusion
holds for both the three-power and five-power towers.
\item For each $a\ne0$ there is a finite integer $m(a)\ge3$
such that its local section does not lift to $N/N^{(m(a))}$.
In particular its global section does not lift at every finite
derived length.  For each $h\ge1$ there is a uniform $m_h$ for
all $a\notin\ell^h\mathbb Z_\ell^{d_n}$.
\end{enumerate}
The assertion is existence of a finite obstruction, not a
calculation of the least $m(a)$ or of its obstruction cocycle.
It does not assert that one finite derived quotient has only
cuspidal sections, nor that the displayed family exhausts the
metabelian section set.

\end{statement}
\begin{proof}
\paragraph{1. The arithmetic source is a quotient of the local group.}

Put $K=\mathbb Q(\zeta_\ell)$ and let $\mathfrak l$ be its
unique prime over $\ell$.  Write
\[
 P_{\rm ar}=G_{K,S_\ell}^{(\ell)},\qquad
 \Gamma_\ell=P_{\rm ar}\rtimes C_{\ell-1}
\]
for the arithmetic quotient used in D.20.  Let
$D\subset P_{\rm ar}$ be the image of $G_{K_{\mathfrak l}}$.
It is a closed subgroup.  We claim $D=P_{\rm ar}$.

Every continuous homomorphism $P_{\rm ar}\to\mathbb F_\ell$
which vanishes on $D$ defines a cyclic degree-$\ell$ extension
of $K$ unramified outside $\mathfrak l$ and split at
$\mathfrak l$.  Thus it is unramified at every finite prime.
The field $K$ is totally imaginary and has class number one,
as proved by the explicit Minkowski bound in D.20.  Global
class field theory implies that the character is zero.
Consequently the image of $D$ in the Frattini quotient
$P_{\rm ar}/\Phi(P_{\rm ar})$ is the whole quotient: otherwise
a nonzero linear functional annihilating that image would
give such a character.  A closed subgroup of a finitely
generated pro-$\ell$ group that surjects on the Frattini quotient
is the whole group.  For completeness, a proper closed subgroup
has proper image in some finite $\ell$-group quotient and lies
in the preimage of a maximal subgroup of that quotient; that
preimage is an index-$\ell$ subgroup.  This contradicts the
Frattini surjectivity.

Since $K_{\mathfrak l}/\mathbb Q_\ell$ has degree $\ell-1$,
the image of $G_{\mathbb Q_\ell}$ in $\Gamma_\ell$ also
surjects onto the constant quotient $C_{\ell-1}$.  It contains
the preceding image $P_{\rm ar}$ and hence equals all of
$\Gamma_\ell$.  Thus
\[
                 G_{\mathbb Q_\ell}\twoheadrightarrow\Gamma_\ell.
                                                        \tag{D.57.1}
\]
No assertion that the global absolute Galois group itself is
free is involved.

\paragraph{2. Every parameter remains nonzero locally.}

Recall $P=N^{(\ell)}$, $V=P^{\rm ab}$,
$B=\bigwedge^2 V$, and the $\Gamma_\ell$-equivariant embedding
$j:P'/P''\hookrightarrow N'/N''$ from D.20.  The parameter
module was chosen as a free submodule
\[
 W\simeq\mathbb Z_\ell^{d_n}
          \subset H^1(\Gamma_\ell,B).
\]
The constructed cocycle $s_a s_O^{-1}$ projects to the class
$a\in W$.  Inflation along the surjection (D.57.1) is injective
on continuous first cohomology.  This can be checked without
an inverse-limit assertion: if the inflated cocycle is a
coboundary with a vector $b\in B$, evaluating on a lift of
each element of $\Gamma_\ell$ gives exactly the same
coboundary formula on $\Gamma_\ell$.  Therefore every nonzero
parameter, and every nonzero difference of parameters, stays
nonzero in $H^1(G_{\mathbb Q_\ell},B)$.

To control geometric conjugators, set
$A=N^{\rm ab}$ and $k=\mathbb Q_\ell$.  We have
\[
 H^0(G_k,A)=0,\qquad H^0(G_k,A/I_O)=0.                 \tag{D.57.2}
\]
Here $I_O=\widehat{\mathbb Z}(1)$ is the primitive boundary
line at the rational cusp $O$.  Indeed, for each coefficient
prime $q$, compactified curve homology is $T_qJ$.
The group $J(k)$ has a torsion-free open subgroup of finite
index, obtained from the formal group sufficiently near zero;
hence its torsion subgroup is finite.  It follows that
$H^0(G_k,T_qJ)=T_q(J(k)_{\rm tors})=0$.  The boundary part
is a finite permutation module tensored with $\mathbb Z_q(1)$,
modulo its diagonal.  After a finite extension splitting the
cusps, both it and its quotient by the primitive rational line
$I_O$ are sums of $\mathbb Z_q(1)$.  These have no invariants:
a finite local field has only finitely many roots of unity,
so its $q$-adic cyclotomic image is infinite.  The homology
exact sequences give (D.57.2), prime by prime and hence for the
full Tate module.

If two local $s_a$ are conjugate, the conjugator's image in
$A$ is invariant because both abelian cocycles are zero.
By (D.57.2) the conjugator lies in $N'/N''$.  Projecting it first
to $P'/P''$ and then to $B$ would make the difference class a
coboundary.  The preceding inflation injectivity forces the
two parameters to be equal.  Thus geometric conjugation does
not destroy local injectivity.

\paragraph{3. The only geometric local class with zero abelian cocycle.}

The full local Kummer map for an abelian variety is injective:
$J(k)$ is a profinite group, so
$J(k)\to\varprojlim_m J(k)/mJ(k)$ is injective, and the Kummer
exact sequences identify this limit with its image in
$H^1(k,TJ)$.  If a local point or cusp section of $U_n$ has
zero abelian cocycle relative to the unit tangent at $O$, its
compactified point $x\in X_n(k)$ satisfies
$[x-O]=0$ in $J(k)$.  The Abel--Jacobi map of the smooth proper
curve $X_n$, of genus at least one, is injective on points:
a principal divisor $x-O$ for $x\ne O$ would give a degree-one
map to $\mathbb P^1$, contrary to positive genus.  Hence
$x=O$.  Since $O\notin U_n$, the section must lie in the cusp
packet at $O$.

The cohomology exact sequence of
$0\to I_O\to A\to A/I_O\to0$ and (D.57.2) show that
\[
                H^1(k,I_O)\longrightarrow H^1(k,A)
\]
is injective.  Therefore the only member of that packet with
zero abelian class is the unit tangential class $s_O$.
The metabelian centralizer/normalizer computation of D.20
shows that this is still the exact packet inside the full
metabelian quotient.  By Paragraph 2, $s_a$ is not that class
when $a\ne0$.  This proves the local nongeometric assertion.

D.54 now rules out a local lift of any $s_a$, $a\ne0$, to the
plain geometrically prosolvable group.  A global lift would
restrict to such a local lift, so it is ruled out as well.
In particular the plain five-power conclusion for this
specific family does not require a finite-Shafarevich theorem
for the Fermat quintic Jacobian.

\paragraph{4. Why the finite-derived conclusion is legitimate locally.}

We first prove the compactness fact that is unavailable over
$G_{\mathbb Q}$ with unrestricted ramification.  Let $k$ be a
finite extension of $\mathbb Q_p$.  For every finite solvable
$G_k$-group $Q$, the set $Z^1_{\rm cts}(G_k,Q)$ is finite.
For finite abelian $Q$, finiteness of $H^1(k,Q)$ is the local
finite-module theorem in Milne, \emph{Arithmetic Duality
Theorems}, Chapter I, Corollary 2.3 (and the finiteness proof
of Theorem 2.1), printed pages 27--28 \cite{DE_MilneADT}; each coboundary orbit has at
most $|Q|$ elements, so $Z^1$ is finite as well.  For general
finite solvable $Q$, use its last nonzero characteristic
derived subgroup $C$.  The target $H^1(k,Q/C)$ is finite by
induction.  The fibre over a represented class is an orbit
set of $H^1(k,C_\xi)$ for the corresponding twisted finite
abelian module.  The latter is finite by the same local
abelian theorem.  Thus $H^1(k,Q)$, and hence $Z^1(k,Q)$, is
finite.  This is finiteness of actual continuous cocycles,
not only a finiteness assertion about discrete abstract maps.

Let $M$ be any finitely generated prosolvable geometric group
with continuous $G_k$ action.  Its characteristic open
subgroups are cofinal.  Therefore
\[
 Z^1_{\rm cts}(G_k,M)
       =\varprojlim_V Z^1_{\rm cts}(G_k,M/V)             \tag{D.57.3}
\]
is compact Hausdorff: every finite coordinate is finite, and
a compatible family defines a continuous cocycle because
each finite coordinate is continuous.  Surjectivity of the
transition maps is not being assumed.  Quotient by geometric
conjugation is compact Hausdorff as well, since the action of
the compact group $M$ has closed orbits.

Put $M_m=N/N^{(m)}$, using the convention $M_2=N/N''$.
Suppose that a fixed actual metabelian local cocycle $s_a$
lifts up to geometric conjugation to $M_m$ for every $m$.
Conjugating a lift by an element lifting the required
metabelian conjugator makes its projection equal to the
\emph{actual} chosen cocycle.  Thus its exact lift fibre
$\mathcal L_m$ is a nonempty closed subset of the compact
space in (D.57.3).  The maps $M_{m+1}\to M_m$ give continuous
maps $\mathcal L_{m+1}\to\mathcal L_m$.  The inverse limit
of these nonempty compact spaces is nonempty: any finite
set of compatibility conditions is satisfied by choosing
an element at one sufficiently late stage and projecting
back, and compactness gives the full intersection.  This
uses no surjectivity of the fibre transition maps.

A compatible family defines an actual continuous cocycle
in
\[
        \varprojlim_m N/N^{(m)}=N^{\rm sol}
\]
and hence a full local prosolvable lift of $s_a$.  Paragraph 3
excludes this for $a\ne0$.  Therefore some finite stage
$m(a)\ge3$ has empty local lift fibre.  This is a genuine
finite-derived obstruction, not only a failure of an
arbitrarily chosen incompatible branch.

\paragraph{5. Uniformity away from the zero parameter.}

Let $L_m\subset\mathbb Z_\ell^{d_n}$ consist of parameters
whose local class lifts to $M_m$.  The local class map is
continuous, by the jointly continuous construction in D.20.
The image of the compact section space for $M_m$ in the
Hausdorff metabelian section space is closed; thus $L_m$
is closed.  The sets decrease with $m$, and the previous
paragraphs show
\[
                           \bigcap_{m\ge2}L_m=\{0\}.
\]
For fixed $h\ge1$, the set
$\mathbb Z_\ell^{d_n}\setminus\ell^h\mathbb Z_\ell^{d_n}$
is compact.  If it met every $L_m$, nested compactness would
give a nonzero element in their intersection.  Hence some
$m_h$ satisfies
\[
                         L_{m_h}\subset\ell^h\mathbb Z_\ell^{d_n}.
                                                        \tag{D.57.4}
\]
The same bound applies to global liftable parameters.

\paragraph{Remaining exact blocker.}

The proof does not identify $m(a)$ or $m_h$, nor a finite
solvable quotient and a numerical nonzero obstruction class.
Those are still the explicit obligations in Task D.5.
It also does not decide whether arbitrarily small nonzero
parameters survive each fixed stage, or whether other
noncuspidal families survive there.  Consequently the hunch
that no solvable quotient of fixed length has only cuspidal
sections is not decided by (D.57.4).
\end{proof}


\begin{statement}{D.58}{Proved}{The full mixed fibre under a precise Selmer hypothesis}

Use the actual cubic twist and the preferred origin of D.8:
\[
 E_c:\ y^3=cx(1-x),\qquad O_c\text{ the cusp over }0,
 \qquad V_c=E_c\setminus\{\text{the three cusps}\}.
\]
All identifications below use that twist convention; replacing the
convention by its inverse merely relabels the cubic class.  Let
$\mathcal F_c\subset H^1(\mathbb Q,E_c[5])$ be the set of
classes in D.8's fibre which extend to a section of the full
geometric tripod group.  Put
\[
 \mathcal S_5(E_c)=\varprojlim_m\operatorname{Sel}_{5^m}(E_c/\mathbb Q),
 \qquad
 R_5(E_c)=\operatorname{im}\bigl(\mathcal S_5(E_c)
                  \longrightarrow H^1(\mathbb Q,E_c[5])\bigr),
\]
where the transition maps on torsion coefficients are multiplication
by five.  Then
\[
 \kappa_5\bigl(E_c(\mathbb Q)/5E_c(\mathbb Q)\bigr)
       \ \subset\ \mathcal F_c\ \subset\ R_5(E_c)
       \ \subset\ \operatorname{Sel}_5(E_c/\mathbb Q).       \tag{D.58.1}
\]
In particular the right hand filters are finite.  There is an exact
sequence of compact abelian groups
\[
 0\longrightarrow E_c(\mathbb Q)^{\wedge 5}
 \longrightarrow\mathcal S_5(E_c)
 \longrightarrow T_5\operatorname{Sha}(E_c/\mathbb Q)\longrightarrow0.  \tag{D.58.2}
\]
Consequently, if $T_5\operatorname{Sha}(E_c/\mathbb Q)=0$, then
\[
 \mathcal F_c=
 \kappa_5\bigl(E_c(\mathbb Q)/5E_c(\mathbb Q)\bigr).       \tag{D.58.3}
\]
The rational five-torsion group of every $E_c$ is zero.  Thus the
set in (D.58.3) has $5^{\operatorname{rank}E_c(\mathbb Q)}$
elements.  In particular (D.58.3) holds whenever $E_c$ has
analytic rank zero or one, with respectively one or five classes.
All classes on the left are realized by rational ordinary points
or rational cusp packets of the actual twist.  This classifies
the first-five image of the full section set under the displayed
hypothesis.  It does not assert a classification of the full
sections themselves, or a classification for the relative
pro-five quotient.

\end{statement}
\begin{proof}
\paragraph{1. The actual twisted cover and its proper Kummer image.}

Represent the fixed cubic cohomology class by the actual cocycle
used in D.8.  The inverse image of its graph in the arithmetic
tripod group is the arithmetic fundamental group of the
corresponding twist $V_c$.  A full section with that cubic
class can be conjugated so that its cubic cocycle equals this
representative; its image then lies in this covering subgroup.
Conversely every section of $V_c$ maps to a full tripod section
in the prescribed cubic fibre.  This is the finite etale
covering correspondence with the actual arithmetic twist,
not an inference from a maximal pro-five quotient.

Compactifying $V_c$ gives the pointed elliptic curve $E_c$.
Its geometric fundamental group is $TE_c$, so composing a
lifted section with the proper pro-five abelian quotient
gives a class
\[
                         z_s\in H^1(\mathbb Q,T_5E_c).
\]
At every finite place $v$, D.56 with the formation of all
finite groups makes the lifted full affine section point
or cusp induced.  Its proper image is therefore the Kummer
class of some $P_v\in E_c(\mathbb Q_v)$.  At the real place,
$H^1(\mathbb R,E_c[5^m])=0$ for every $m$, by averaging over
the group of order two; there is no missing five-primary
local condition.  Reducing $z_s$ modulo $5^m$ at each $m$
therefore gives a compatible class in
$\operatorname{Sel}_{5^m}(E_c/\mathbb Q)$.
Its first reduction is precisely the translation coordinate
of the section in D.8, by the origin and twist comparison
there.  This proves the middle containment in (D.58.1).

For the reverse containment on the left, take any
$P\in E_c(\mathbb Q)$.  If $P\in V_c$, its ordinary point
section gives the Kummer class $\kappa_5(P-O_c)$.  If $P$
is one of the three deleted points, a nonzero rational
tangent gives a cusp section with the same proper Kummer
class.  All three cusps are rational, and $O_c$ is the
origin, so this construction applies to every rational
point of $E_c$.  It proves the first containment in (D.58.1).

\paragraph{2. The inverse limit and the exact Shafarevich--Tate term.}

The finite Kummer sequences and the definition of the Selmer
local conditions give, for each $m\ge1$,
\[
 0\longrightarrow E_c(\mathbb Q)/5^mE_c(\mathbb Q)
 \longrightarrow\operatorname{Sel}_{5^m}(E_c/\mathbb Q)
 \longrightarrow\operatorname{Sha}(E_c/\mathbb Q)[5^m]
 \longrightarrow0.                                  \tag{D.58.4}
\]
The same sequence and finiteness of its middle term are
displayed in Kolyvagin's primary manuscript, Section 2,
equation (5), printed page 6 \cite{DE_Kolyvagin1990}.
The left hand transition maps are the natural quotient
maps and are surjective.  Thus this inverse system is
Mittag--Leffler, and its first derived inverse limit is
zero.  Taking inverse limits of (D.58.4) proves (D.58.2),
including surjectivity on the right.  All finite Selmer
groups carry their finite discrete topology; their inverse
limit is compact.  There is no assertion that the global
groups $H^1(\mathbb Q,E_c[5^m])$ are finite, and no
unjustified omission of a global derived-limit term.

If $T_5\operatorname{Sha}=0$, equation (D.58.2) identifies the inverse
Selmer group with the five-adic completion of the
Mordell--Weil group.  Its reduction modulo five has image
exactly $E_c(\mathbb Q)/5E_c(\mathbb Q)$, since a finitely
generated abelian group surjects onto this finite quotient.
The two containments already proved now give (D.58.3).
Notice that finite five-primary Shafarevich--Tate classes
may occur in $\operatorname{Sel}_5$; they do not survive
as a nonzero compatible Tate-module contribution.  This
is why the inverse Selmer group is the correct filter.

\paragraph{3. The torsion count and the unconditional analytic-rank cases.}

Choose a prime $q\equiv17\pmod{30}$ at which $c$ is a unit.
Such a prime exists by Dirichlet's theorem, and it is a
good reduction prime for the plane cubic $E_c$, since
$q\nmid3c$.  Because $q\equiv2\pmod3$, cubing is a
bijection on $\mathbb F_q$.  For each of the $q$ affine
values of $x$ there is exactly one $y$, and there is one
point at infinity.  Therefore
\[
                         |E_c(\mathbb F_q)|=q+1
                                      \equiv3\pmod5.
\]
Good reduction injects rational five-torsion into this
group, so $E_c(\mathbb Q)[5]=0$.  The structure theorem
for the finitely generated Mordell--Weil group gives the
stated cardinality in (D.58.3).

Every elliptic curve over $\mathbb Q$ is modular by
\cite[Theorem A]{DE_BCDT2001}.  Kolyvagin's primary
1990 manuscript, Theorem 1, printed page 4, states that
a modular elliptic curve of analytic rank at most one
has equal algebraic and analytic ranks and finite
Shafarevich--Tate group \cite{DE_Kolyvagin1990}.
The next page explains the Gross--Zagier and
quadratic-twist nonvanishing inputs that remove the
auxiliary Heegner-point hypothesis of the earlier theorem.
Thus $T_5\operatorname{Sha}(E_c/\mathbb Q)=0$ in these cases, with no
restriction that five be an ordinary or good reduction
prime.  Equations (D.58.2)--(D.58.3) finish the proof.
\end{proof}

\paragraph{Remaining scope of Task D.9.}

For arbitrary $c$ with a possible nonzero
$T_5\operatorname{Sha}(E_c/\mathbb Q)$, this argument supplies the
finite inverse-Selmer image $R_5(E_c)$ but does not decide
which of its non-Mordell--Weil classes lift nonabelianly.
It also leaves the relative pro-five question open even
when (D.58.3) settles the full-group first-five image.


\begin{statement}{D.59}{Proved}{Every second prime detects the Fermat family}

Let $\ell\in\{3,5\}$, $n=\ell^r$, $k=\mathbb Q_\ell$,
and use the Fermat curve, cusp and family of D.20/D.57:
\[
 X_n:X^n+Y^n+Z^n=0,\quad
 U_n=X_n\setminus\{XYZ=0\},\quad O=[0:-1:1],\quad
 s_a\quad(a\in\mathbb Z_\ell^{d_n}).
\]
Write $J_n$ for its Jacobian.  The single-prime proper Kummer map
\[
 X_n(k)\longrightarrow H^1(k,T_\ell J_n),\qquad
                   P\longmapsto\kappa_\ell([P-O])       \tag{D.59.1}
\]
is injective.  In particular its zero fibre is $\{O\}$.

For every prime $q\ne\ell$, put $\Sigma=\{\ell,q\}$ and
$N_\Sigma=\pi_1(U_{n,\overline{\mathbb Q}})^\Sigma$.
The images of the $s_a$ in $N_\Sigma/N_\Sigma''$ remain
pairwise nonconjugate over $k$.  Every nonzero parameter
has no local lift to $N_\Sigma$, and hence no global lift.
For each $a\ne0$ there is a finite $m(a,q)\ge3$ such that
it already has no local lift to
$N_\Sigma/N_\Sigma^{(m(a,q))}$.  For each $h\ge1$ the
stage can be chosen uniformly for
$a\notin\ell^h\mathbb Z_\ell^{d_n}$.
Thus $q=2$ suffices in both the three-power and five-power
towers.  The least obstructing derived length and an
explicit finite obstruction cocycle remain undetermined.

\end{statement}
\begin{proof}
\paragraph{1. All local points lie in three explicitly controlled disks.}

For a unit modulo $\ell^2$, its $n$th power equals its
$\ell$th power, because
\[
 n\equiv\ell\pmod{\ell(\ell-1)}.
\]
The possible unit powers are
\[
 \begin{array}{c|c|c}
 \ell&\ell^2&\{u^n\bmod\ell^2:u\in\mathbb Z_\ell^\times\}\\\hline
 3&9&\{1,8\}\\
 5&25&\{1,7,18,24\}.
 \end{array}                                          \tag{D.59.2}
\]
No three entries in either set, with repetition allowed,
sum to zero modulo $\ell^2$.  This finite assertion is
executed exhaustively in \path{code/de_fermat_two_prime.py}:
there are eight triples for $\ell=3$ and 64 for $\ell=5$.
The same code checks every primitive triple modulo
$\ell^2$ satisfying the Fermat equation; all 162 triples
modulo nine and all 1500 triples modulo 25 have exactly
one coordinate divisible by $\ell$.

Normalize a point of $X_n(k)$ to primitive integral
projective coordinates.  Equation (D.59.2) rules out
three unit coordinates.  Two coordinates divisible by
$\ell$ would leave the remaining unit term nonzero
modulo $\ell$, which is also impossible.  Therefore
exactly one coordinate is divisible by $\ell$.

Suppose that coordinate is $X$.  With $t=X/Z$ and
$u=Y/Z$ we have $t\in\ell\mathbb Z_\ell$, $u$ a unit,
and $u^n=-(1+t^n)$.  Reduction modulo $\ell$ gives
$u\equiv-1\pmod\ell$.  The analytic expression
\[
 h(t)=\exp\left(\frac1n\log(1+t^n)\right)              \tag{D.59.3}
\]
converges on the geometric open disk
\[
 v_\ell(t)>\rho_{\ell,r}:=
                    \frac{r+1/(\ell-1)}{\ell^r}<1.
\]
Indeed $v_\ell(t^n)>r+1/(\ell-1)$, so the logarithm
has that same valuation, and division by $n$ leaves
valuation greater than $1/(\ell-1)$, the convergence
threshold for the exponential.  The strict inequality
$\rho_{\ell,r}<1$ holds at $r=1$ and then follows by
induction.  Thus (D.59.3) is analytic on a disk strictly
larger than the closed disk $|t|\le\ell^{-1}$.

It satisfies $h(t)^n=1+t^n$ and $h(0)=1$.
The group $1+\ell\mathbb Z_\ell$ is torsion free for
$\ell=3,5$, as its logarithm injects it into
$\ell\mathbb Z_\ell$.  Hence $u=-h(t)$ is the unique
$k$-valued branch reducing to $-1$.  The map
\[
 t\longmapsto[t:-h(t):1]
\]
is an analytic embedding of that disk in $X_n^{\rm an}$:
projection to $t$ is its inverse, and the derivative
$nu^{n-1}$ in the defining equation is nonzero.  It
contains $O$ and every local point for which $X$ is
the nonunit coordinate.  Permuting coordinates gives
the other two disks, centred at
\[
                      O_1=[-1:0:1],\qquad
                      O_\infty=[-1:1:0].
\]
Their closed subdisks $|t|\le\ell^{-1}$ cover $X_n(k)$.
Only these embedded closed disks are needed below.

\paragraph{2. Prime-to-$\ell$ Kummer classes are constant on each disk.}

Fix a prime $b\ne\ell$, an integer $j\ge1$, and one
of the three disk centres $d$.  Let
$\iota_d:X_n\to J_n$ be $P\mapsto[P-d]$.  Pull back
$[b^j]:J_n\to J_n$ along $\iota_d$.  This is a finite
etale torsor under $J_n[b^j]$, including at the cusp $d$.
After base change to $\mathbb C_\ell$ it is a Galois
cover of degree a power of $b$.  Its restriction to
the embedded closed disk is trivial by Berkovich's
Theorem 6.3.2 and Corollary 6.3.3, printed pages
119--121 \cite{DE_Berkovich1993}.  The exact scope is
prime-to-residue-characteristic geometric etale
covers; no triviality assertion for $\ell$-power
covers of a disk is used.

The trivialization gives a canonical identification
of the fibres over $d$ and over any $P\in X_n(k)$
in the disk: pair points in the same connected
component of the covering disk.  This identification
commutes with the action of $G_k$, since the disk,
covering and its centre are defined over $k$.
The rational zero division point in the fibre over
$d$ belongs to a unique component.  That component
is $G_k$-stable, so its unique point over $P$ is
$G_k$-fixed.  This point belongs to the finite
algebraic fibre over $P$, hence is a $k$-rational
division point, not merely a point over a completed
extension.  It follows that
\[
                     [P-d]\in b^jJ_n(k)
                       \quad\text{for every }j.       \tag{D.59.4}
\]
This proves the vanishing of its full $b$-adic
Kummer class with the required rational descent.

The line $Y+Z=0$ cuts the Fermat curve in $nO$.
Similarly $X+Z=0$ and $X+Y=0$ cut it in $nO_1$
and $nO_\infty$.  Ratios of these lines therefore
show
\[
                   n[O_1-O]=n[O_\infty-O]=0.
\]
These differences are $\ell$-primary torsion and
are divisible by every power of $b\ne\ell$ inside
$J_n(k)$.  Combining this with (D.59.4) gives
\[
 \kappa_b([P-O])=0\quad
       \text{for every }P\in X_n(k),\ b\ne\ell.         \tag{D.59.5}
\]

If two points have the same class in (D.59.1),
their difference has zero Kummer class at every
prime, by (D.59.5).  The full local Kummer map of
$J_n(k)$ is injective by the profinite argument
in D.57.  Thus their degree-zero divisor classes
coincide.  The Abel--Jacobi map of $X_n$, whose
genus is at least one, is injective on points.
The two points are equal, proving (D.59.1).

\paragraph{3. Projection to any two-prime metabelian group retains the parameters.}

The geometric quotient map
$N_\Sigma\twoheadrightarrow N^{(\ell)}=P$
induces maps on abelianization, commutator
abelianization, and the class-two quotient
$B=\bigwedge^2P^{\rm ab}$ used in D.57.
Consequently the projection of the actual D.20
cocycle still has parameter $a$ in
$H^1(G_k,B)$.  This parameter is nonzero when
$a\ne0$, by the actual local surjection to
$\Gamma_\ell$ and inflation injectivity proved
in D.57.

Set $A_\Sigma=N_\Sigma^{\rm ab}$.  The proof
of (D.57.2), restricted to the two coefficient
primes, gives
\[
 H^0(G_k,A_\Sigma)=0,\qquad
 H^0(G_k,A_\Sigma/I_{O,\Sigma})=0.                    \tag{D.59.6}
\]
Here $I_{O,\Sigma}=\prod_{b\in\Sigma}\mathbb Z_b(1)$
is the primitive cusp line.  Thus a conjugator
between two projected family sections has zero
abelian image and lies in the commutator kernel.
Its projection to $B$ would make their difference
parameter a coboundary.  The two parameters must
therefore be equal.  This proves pairwise
nonconjugacy in the two-prime metabelian quotient.

Every family section has zero abelian cocycle.
If its local two-prime class were geometric, its
proper $\ell$-adic Kummer class would be zero.
By (D.59.1) the underlying compactification point
would be $O$.  Hence it would be cuspidal at
$O$.  The second vanishing in (D.59.6) makes
$H^1(k,I_{O,\Sigma})\to H^1(k,A_\Sigma)$
injective.  The zero abelian class forces the
unit cusp parameter, whose metabelian class is
$s_0$.  Pairwise nonconjugacy excludes this
for $a\ne0$.

\paragraph{4. A two-prime local lift has to be geometric.}

The class of all finite groups whose orders
have prime divisors in $\Sigma$ is a full
formation.  It contains the residue prime
$\ell$ and the distinct prime $q$.  D.56
therefore makes every local section of the
geometrically pro-$\Sigma$ arithmetic group
point or cusp induced.  Paragraph 3 excludes
a lift of every nonzero family member to
that group.  Restriction to $G_k$ excludes
a global lift as well.

Every finite $\Sigma$-group is solvable by
Burnside's two-prime theorem.  Therefore
\[
 N_\Sigma=\varprojlim_m
                 N_\Sigma/N_\Sigma^{(m)}.
\]
The actual local cocycle spaces for its
finite derived quotients are compact by
the finite-solvable-coefficient argument
of D.57.  If a fixed metabelian class
lifted at every derived length, conjugator
lifting makes all fibres exact, and
compactness gives one compatible lift to
$N_\Sigma$.  This is impossible for a
nonzero parameter.  The closed nested
liftable-parameter subsets and compactness
of the complement of
$\ell^h\mathbb Z_\ell^{d_n}$ give the
uniform finite stage exactly as in
(D.57.4).  No surjectivity of transition
maps of section fibres is assumed.
\end{proof}


% D.60; full independent root proof audit passed.
\subsection{An infinitely ramified Alexander module with zero first cohomology}

\begin{statement}{D.60}{Proved}{A completed cusp tangent with a unique metabelian lift}
Let \(M=F^{(5)}/(F^{(5)})''\) be the geometrically pro-five
metabelian quotient of D.35, with its actual tangential arithmetic
action. There is an infinitely supported completed Kummer class
\(a\in H^1(\mathbb Q,\mathbb Z_5(1))\), with valuation one at two,
such that for the actual twisted Alexander module of D.35,
\[
          H^1_{\mathrm{cts}}(G_{\mathbb Q},A_{(a,0)})=0.
                                                        \tag{D.60.1}
\]
Consequently the abelian section class \((a,0)\) has exactly one
metabelian section class above it: the class of its completed cusp
tangent at zero. It is the image of an actual full cusp section.

The construction gives a genuine rank-zero example in D.47, whereas
D.44 gives rank \(\aleph_0\) for every abelian pair supported only
at five. It does not determine the rank for every arbitrarily
ramified pair, or the least derived length of a Fermat obstruction.
\end{statement}

\begin{proof}
Write \(\Lambda=\mathbb Z_5[[u,v]]\),
\(\mathfrak m=(5,u,v)\), and \(t_1=1+u\).
Choose once and for all compatible roots defining continuous
integral Kummer cocycles \(\kappa(q)\) for each rational prime
\(q\). We shall choose pairwise disjoint finite sets of primes
\(Q_r\), for \(r\ge1\), disjoint from \(\{2,5\}\), and put
\[
 a=\kappa(2)+\sum_{r\ge1}5^r\sum_{q\in Q_r}\kappa(q).
                                                        \tag{D.60.2}
\]
Each inner sum is finite. The outer series converges uniformly
as a series of \(\mathbb Z_5(1)\)-valued continuous cocycles:
the entire \(r\)-th summand takes values in
\(5^r\mathbb Z_5(1)\), irrespective of the number of terms.
Thus (D.60.2) is an actual continuous cocycle, not only a formal
family of local valuations. Distinct prime supports give its
valuation one at two and \(5^r\) at each prime of \(Q_r\).

\paragraph{Finite construction at stage \(r\).}
Suppose \(Q_1,\ldots,Q_{r-1}\) have been chosen. Set
\[
 a_{<r}=\kappa(2)+
       \sum_{j<r}5^j\sum_{q\in Q_j}\kappa(q),\qquad
 S_r=\{2,5,\infty\}\cup\bigcup_{j<r}Q_j .
                                                        \tag{D.60.3}
\]
Give the finite additive module
\(W_r=\Lambda C/\mathfrak m^r\Lambda C\)
the actual action \(\rho_r\) of the abelian pair
\((a_{<r},0)\). This action factors through
\(G_{\mathbb Q,S_r}\). Indeed the untwisted tripod action is
unramified away from five, and the finite Kummer sum is
unramified outside the displayed set.

Any later additions in (D.60.2) leave this action unchanged.
By (D.47.4), for every \(c\in\mathbb Z_5\),
\[
           t_1^{5^r c}-1\in\mathfrak m^{r+1}.
                                                        \tag{D.60.4}
\]
Formula (D.35.1) therefore makes \(\rho_r\) exactly the
reduction of the final action, for every possible continuation
of our construction.

The set
\[
             E_r=Z^1(G_{\mathbb Q,S_r},W_r)
                                                        \tag{D.60.5}
\]
of actual continuous cocycles is finite. This is the finite
coefficient, finite ramification finiteness used in D.47:
\(H^1(G_{\mathbb Q,S_r},W_r)\) and the group of coboundaries
are both finite. Equivalently the corresponding finite affine
representations are bounded-degree extensions unramified
outside the fixed finite set.

For each \(z\in E_r\) whose restriction to
\(\ker\rho_r\) is nonzero, choose
\(\delta\in\ker\rho_r\) with \(z(\delta)\ne0\).
The finite affine representation
\[
 g\longmapsto\bigl(z(g),\rho_r(g)\bigr)
       \in W_r\rtimes\operatorname{Aut}(W_r)
                                                        \tag{D.60.6}
\]
contains the element \((z(\delta),1)\).
Chebotarev \cite[Chapter V, Theorem 3.23]{PBC_MilneCFT}
supplies a prime \(q\notin S_r\) whose Frobenius
belongs to its conjugacy class. Choose these primes distinct
for the finitely many cocycles under consideration. Then
\[
                 \rho_r(\operatorname{Frob}_q)=1,
                 \qquad z(\operatorname{Frob}_q)\ne0.
                                                        \tag{D.60.7}
\]
The second property is independent of the Frobenius
representative: conjugating \((w,1)\) in the affine group
sends \(w\) to an invertible linear transform of \(w\).
Let \(Q_r\) contain all the chosen witness primes and one
additional unused prime different from two and five. The
additional prime ensures infinite support, including when
there are no cocycles to eliminate. This completes the
induction and defines (D.60.2).

\paragraph{Every unrestricted integral cocycle is captured.}
Let \(z:G_{\mathbb Q}\to A_{(a,0)}\) be any continuous
integral cocycle. Its reduction \(z_r\) belongs to the
particular finite set (D.60.5), even if \(z\) has unrestricted
ramification. Here is the full uniform support argument.
At a prime \(q\ne5\) with zero abelian inertia, D.47 proves
that the inertia value of \(z\) is zero: wild inertia and
the prime-to-five tame factors have no homomorphism to the
pro-five coefficient group, and the first nonzero augmentation
grade of a putative tame value would satisfy
\((q^{j+2}-q)m=0\) in a torsion-free lattice.
At a prime \(q\in Q_j\), D.43 gives
\(H^1(G_{\mathbb Q_q},A_{(a,0)})=0\), so the local cocycle
is an actual integral coboundary. Its tame inertia value is
\[
                        (t_1^{5^j}-1)h_q,
                      \qquad h_q\in A_{(a,0)}.
                                                        \tag{D.60.8}
\]
If \(j\ge r\), this lies in
\(\mathfrak m^{j+1}A_{(a,0)}\subset\mathfrak m^rA_{(a,0)}\).
The action has trivial inertia modulo \(\mathfrak m^r\) there
as well. These are all the primes outside \(S_r\).
Hence both the action and the cocycle graph are unramified
outside \(S_r\), which is precisely the assertion that
\(z_r\in E_r\). This argument does not presume that a
general finite coefficient cocycle lifts integrally, or
choose a ramification set depending on \(z\).

If \(z_r\) were nonzero on \(\ker\rho_r\), stage \(r\)
would have selected its witness prime \(q\in Q_r\).
The final inertia vector there is \((5^r,0)\ne(0,0)\),
so D.43 again expresses the restriction of \(z\) as an
integral coboundary, including on a Frobenius lift.
But (D.60.7) says that this Frobenius acts trivially on
\(W_r\). Reduction of that coboundary therefore gives
\(z_r(\operatorname{Frob}_q)=0\), contradicting
(D.60.7). We have proved
\[
                         z_r|_{\ker\rho_r}=0
                         \quad\hbox{for every }r.
                                                        \tag{D.60.9}
\]

\paragraph{Descent to the actual operator image.}
Let \(I=\rho_{(a,0)}(G_{\mathbb Q})\), with its compact
inverse-limit topology. An element of the kernel of this
full action belongs to the kernel of every \(\rho_r\).
By (D.60.9) its cocycle value is zero modulo every
\(\mathfrak m^r\), and hence is zero. The cocycle \(z\)
is consequently constant on fibres of
\(G_{\mathbb Q}\twoheadrightarrow I\) and descends to a
continuous cocycle on \(I\). Continuity follows from the
compact quotient topology; no passage from finite-level
cohomology through an unexamined derived limit occurs.

We show directly that \(H^1(I,A_{(a,0)})=0\).
The image of five-primary inertia at two is the subgroup
\[
 N=\{m\longmapsto t_1^c m:c\in\mathbb Z_5\}\subset I.
                                                        \tag{D.60.10}
\]
It is the full displayed subgroup because the valuation of
\(a\) at two is one. The original tripod action and
cyclotomic character are unramified there. Every operator
in (D.35.1) is a multiplier followed by the substitution
\(t_i\mapsto t_i^{\chi_5(g)}\). It conjugates multiplication
by \(t_1^c\) to multiplication by
\(t_1^{c\chi_5(g)}\); therefore \(N\) is normal in \(I\).
Moreover
\[
                         A_{(a,0)}^N=0,
                                                        \tag{D.60.11}
\]
since \(t_1-1=u\) is a nonzero divisor in \(\Lambda\).

Given a continuous cocycle \(h:I\to A_{(a,0)}\), pull it
back to \(G_{\mathbb Q_2}\). D.43, applied to the primitive
inertia vector \((1,0)\), makes this pullback a coboundary
of one element \(b\in A_{(a,0)}\). Subtract the same
coboundary on all of \(I\). The resulting cocycle \(h'\)
vanishes on \(N\). For \(n\in N\) and \(g\in I\), its
cocycle identity and normality of \(N\) give
\[
 n h'(g)=h'(ng)
        =h'\bigl(g(g^{-1}ng)\bigr)=h'(g).
\]
Thus \(h'(g)\in A_{(a,0)}^N=0\) for every \(g\).
This proves the desired vanishing on \(I\), and the descent
already proved gives (D.60.1).

Finally the cusp section \(g\mapsto x^{a(g)}s_0(g)\)
exists in the actual geometrically pro-five fundamental
group. It is also the pro-five image of a full cusp section:
in the product decomposition of \(\widehat{\mathbb Z}(1)\),
use \(a\) in the five-adic component and zero in the others.
Its abelian class is \((a,0)\). By D.35 the nonempty
metabelian fibre is a torsor under (D.60.1), and hence is
a singleton. In particular this conclusion concerns actual
section conjugacy classes, with the abelian conjugator
stabilizer already removed in D.35.
\end{proof}


% D.61; independent root full proof audit passed.
\subsection{Density of the abelian twists with zero Alexander cohomology}

\begin{statement}{D.61}{Proved}{Zero-rank Alexander modules occur in every completed abelian neighborhood}
For every pair \(a_0,b_0\in H^1(\mathbb Q,\mathbb Z_5(1))\)
and every integer \(N\ge0\), there is
\[
 a\in a_0+5^N H^1(\mathbb Q,\mathbb Z_5(1))
 \quad\hbox{such that}\quad
 H^1_{\mathrm{cts}}(G_{\mathbb Q},A_{(a,b_0)})=0.
                                                        \tag{D.61.1}
\]
The class \(a\) can be chosen with infinitely many nonzero
prime valuations, while the second coordinate is unchanged.
The abelian pairs with zero first Alexander cohomology are
therefore dense for the displayed five-adic neighborhood
system. Existence of a metabelian section over a general
pair is not asserted; if the fibre is nonempty, it is a singleton.
For \(b_0=0\), the actual cusp tangent supplies that singleton
in the geometrically pro-five metabelian quotient \(M\) of
D.35. No uniqueness assertion about the full all-prime
metabelian quotient is made.
D.66 strengthens this density statement to a comeagre
locus in the same intrinsic five-adic topology, including
every fixed-second-coordinate space.
\end{statement}

\begin{proof}
Choose actual continuous integral representatives of \(a_0,b_0\),
and let \(a_{0,q},b_{0,q}\in\mathbb Z_5\) denote their tame valuations
at \(q\). The family of valuations tends to zero five-adically:
for each \(h\), only finitely many primes have
either \(a_{0,q}\) or \(b_{0,q}\) outside \(5^h\mathbb Z_5\).
Choose \(q_0\ne5\) with
\(a_{0,q_0}\in5^{N+1}\mathbb Z_5\), and begin with
\(a_0+5^N\kappa(q_0)\). Its valuation at \(q_0\) is
\[
                c=a_{0,q_0}+5^N
                   \in5^N\mathbb Z_5^\times,
                                                        \tag{D.61.2}
\]
where the notation means \(v_5(c)=N\).
We never select \(q_0\) again.

At stage \(r\ge1\), put
\[
 \begin{split}
 a_{<r}&=a_0+5^N\kappa(q_0)+
             \sum_{j<r}5^{N+j}\sum_{q\in Q_j}\kappa(q),\\
 S_r&=\{5,\infty,q_0\}\cup\bigcup_{j<r}Q_j
       \cup\{q\ne5:\min(v_5(a_{0,q}),v_5(b_{0,q}))<r-1\}.
 \end{split}                                           \tag{D.61.3}
\]
The set \(S_r\) is finite, although \(a_0\) itself may
have unrestricted ramification. The action \(\rho_r\)
on \(W_r=\Lambda C/\mathfrak m^r\Lambda C\) determined
by \((a_{<r},b_0)\) is unramified outside \(S_r\), by
the binomial estimate (D.47.4). All additions at stages
\(j\ge r\) multiply it by
\(t_1^{5^{N+j}d}\equiv1\pmod{\mathfrak m^r}\), for
every integral \(d\). Thus its equality with the final
reduced action is exact, not merely up to conjugacy.

Enumerate the finite set
\(Z^1(G_{\mathbb Q,S_r},W_r)\). For each member
nonzero on \(\ker\rho_r\), choose its Chebotarev
witness prime exactly as in (D.60.6)--(D.60.7), with
\(\rho_r(\operatorname{Frob}_q)=1\) and nonzero
cocycle value. In addition to \(S_r\) and all earlier
chosen primes, require the witness to avoid the finite set
\[
 \{q:a_{0,q}\notin5^{N+r+1}\mathbb Z_5
                  \ \hbox{or}\ b_{0,q}\notin5^{N+r+1}\mathbb Z_5\}.
                                                        \tag{D.61.4}
\]
Removing finitely many primes does not affect existence
in a nonempty Chebotarev class. Let \(Q_r\) consist of
these distinct witness primes and one additional unused
prime satisfying the same avoidance condition. Define
\[
 a=a_0+5^N\kappa(q_0)+
                  \sum_{r\ge1}5^{N+r}\sum_{q\in Q_r}\kappa(q).
                                                        \tag{D.61.5}
\]
Uniform convergence is the same integral argument as
in D.60, and the entire difference is \(5^N\) times
an actual continuous integral cocycle. At every new
prime \(q\in Q_r\), the final valuation is
\[
               a_q=a_{0,q}+5^{N+r},
                    \qquad v_5(a_q)=N+r.
                                                        \tag{D.61.6}
\]
It cannot cancel: the first summand has strictly greater
five-adic valuation by (D.61.4). Later additions have
different prime supports. The additional prime at every
stage consequently ensures infinitely many nonzero
valuations, even if some original valuations of \(a_0\)
were already nonzero.

Every integral cocycle for the final \(A_{(a,b_0)}\),
reduced modulo \(\mathfrak m^r\), lies in the finite
set enumerated at stage \(r\). Indeed, at any prime
outside \(S_r\), both original valuations are divisible
by \(5^{r-1}\); if the prime is selected later, the
new summand is divisible by \(5^{N+j}\) with \(j\ge r\).
Both final valuations are therefore still divisible by
\(5^{r-1}\). D.43 makes a nonzero-inertia local
cocycle an integral coboundary, whose inertia value is
in \(\mathfrak m^r\) by (D.47.4). A zero-inertia
prime contributes zero by D.47. This proves the full
unrestricted cocycle capture, including primes already
ramified in \(a_0\).

The new nonzero inertia at each witness prime now
contradicts its prescribed nonzero Frobenius value in
exactly the manner of D.60. Hence every integral cocycle
descends continuously to the actual compact operator
image \(I\). The image of inertia at \(q_0\) is
\[
 N=\{\text{multiplication by }
         (t_1^c t_2^{b_{0,q_0}})^d:d\in\mathbb Z_5\},
                                                        \tag{D.61.7}
\]
where \(c\ne0\) is (D.61.2). Every actual operator is
a multiplier followed by \(t_i\mapsto t_i^{\chi_5(g)}\),
so it conjugates the displayed multiplier to its
\(\chi_5(g)\)-th power. Hence \(N\) is normal in \(I\).
Primitivity is unnecessary: \(t_1^c t_2^{b_{0,q_0}}-1\)
is nonzero, as is already seen by setting \(t_2=1\),
and is therefore a nonzero divisor in \(\Lambda\).
Thus \(A_{(a,b_0)}^N=0\). D.43 at \(q_0\), followed by the
normal-subgroup cocycle identity of D.60, proves
\(H^1(I,A_{(a,b_0)})=0\). This gives (D.61.1).

The exact D.35 torsor description makes every nonempty
metabelian fibre over the constructed pair a singleton.
For \(b_0=0\), the full cusp section with five-adic
tangent component \(a\) supplies that metabelian lift.
The topology in the statement is explicitly the one
with neighborhoods \(a_0+5^N H^1\); no compactness
of this completed Kummer group is asserted.
\end{proof}


% D.62; independent root and FGH full proof audits passed.
\subsection{A nonzero Jannsen kernel for every Alexander twist}

\begin{statement}{D.62}{Proved}{Infinite ramification never removes the nonseparated second cohomology}
Let \(s=(a,b)\) be any completed abelian Kummer pair over
\(\mathbb Q\), and let \(A_s=\Lambda C\) have the actual
Alexander action of D.35. No finite ramification hypothesis
and no metabelian lifting hypothesis are imposed. Set
\(W_r=A_s/\mathfrak m^r A_s\), where
\(\mathfrak m=(5,u,v)\). Then
\[
 \varprojlim\nolimits_r^1 H^1(G_{\mathbb Q},W_r)\ne0.
                                                        \tag{D.62.1}
\]
Consequently the continuous-cohomology comparison fits into
the exact sequence with nonzero first term
\[
 0\longrightarrow
 \varprojlim\nolimits_r^1 H^1(G_{\mathbb Q},W_r)
 \longrightarrow H^2_{\mathrm{cts}}(G_{\mathbb Q},A_s)
 \longrightarrow\varprojlim_r H^2(G_{\mathbb Q},W_r)
 \longrightarrow0.
                                                        \tag{D.62.2}
\]
With the continuous-cochain quotient topology defined in
the proof, the first term is exactly the closure of zero.
In particular this \(H^2\) is not Hausdorff for any twist.
Every class in that first term restricts to zero in
\(H^2_{\mathrm{cts}}(G_{\mathbb Q_v},A_s)\) at every place
\(v\).

This proves nonvanishing of the unrestricted derived-limit
term even for infinitely ramified pairs, including the
rank-zero modules of D.60--D.61. It does not evaluate its
isomorphism type or the entire second cohomology for every
such pair.
\end{statement}

\begin{proof}
Write \(B_r=H^1(G_{\mathbb Q},W_r)\). Each \(B_r\)
is countable. Indeed there are countably many finite
Galois extensions of \(\mathbb Q\), and every continuous
cocycle with finite coefficients, together with its finite
action, factors through one of these finite quotients.

\paragraph{Every finite stage has infinitely many surviving
residues at the first coefficient level.}
Fix \(r\). Choose a finite set \(S\) containing five and
infinity such that the finite module \(W_r\) is unramified
outside \(S\). Such an \(S\) always exists, regardless of
the ramification of the full compact module \(A_s\).
Let \(X=\operatorname{Spec}\mathbb Z[S^{-1}]\), with
infinity omitted from the ring notation. Finite-coefficient
localization gives the exact segment
\[
 H^1(G_{\mathbb Q},W_r)
 \xrightarrow{\partial_r}
 \bigoplus_{q\notin S}H^0(\mathbb F_q,W_r(-1))
 \xrightarrow{\delta_r} H^2(G_{\mathbb Q,S},W_r).
                                                        \tag{D.62.3}
\]
For completeness, apply localization on \(X\) after
removing a finite set of closed points. At an unramified
prime, purity identifies the cohomology with support
with the residue-field cohomology shifted by two and
twisted by minus one. Taking the filtered union over
the removed finite sets gives (D.62.3), since continuous
finite-valued cochains factor through finite Galois
quotients. The comparison of the etale and Galois groups
is \cite[II, Proposition 2.9]{DE_MilneADT}. This is the
finite-coefficient calculation of (D.48.4); no inverse
limit is taken with the current \(S\) fixed across \(r\).
The last group in (D.62.3) is finite: its kernel of
localization is finite by
\cite[I, Theorem 4.10(a)]{DE_MilneADT}, and its image
lies in the finite product of the finite local second
cohomology groups at \(v\in S\)
\cite[I, Corollary 2.3]{DE_MilneADT}.

The exponent of \(W_r\) divides \(5^r\). Take a finite
Galois extension cutting out both its actual action
\(\rho_r\) and \(\mu_{5^r}\). Infinitely many primes
\(q\notin S\) split completely in this extension, by
Chebotarev \cite[Chapter V, Theorem 3.23]{PBC_MilneCFT}. At each such prime,
\[
 \rho_r(\operatorname{Frob}_q)=1,
 \qquad q\equiv1\pmod{5^r},
 \qquad H^0(\mathbb F_q,W_r(-1))=W_r(-1).
                                                        \tag{D.62.4}
\]
Choose at each such residue field a compatible basis
for the Tate twist and let \(w_q\) be the residue
corresponding to \(C\bmod\mathfrak m^r\). Its reduction
in \(W_1(-1)\) is nonzero.

Because the target of \(\delta_r\) is finite, an infinite
subset of these primes has one common value
\(\delta_r(w_q)\). Fix a prime \(q_*\) in that subset.
For every other \(q\) there, the element with residue
\(w_q\) at \(q\), residue \(-w_{q_*}\) at \(q_*\),
and zero at all other primes lies in \(\ker\delta_r\).
By (D.62.3) it is the residue of a class in \(B_r\).
Reduction to \(B_1\) leaves a nonzero residue at its
own prime \(q\) and none at the other chosen primes
except the common \(q_*\). These reduced classes are
linearly independent over \(\mathbb F_5\), as their
individual residues show. Consequently
\[
 \dim_{\mathbb F_5}\operatorname{im}(B_r\longrightarrow B_1)
                              =\aleph_0
                   \quad\hbox{for every }r\ge1.
                                                        \tag{D.62.5}
\]
The upper bound follows from the countability already
proved. This uses the actual finite action, not only
the cyclotomic character.

\paragraph{Integral cocycles have a uniformly finite first image.}
The groups \(H^0(G_{\mathbb Q},W_r)\) are finite, so
their inverse system is Mittag--Leffler. The continuous
cochain comparison therefore gives
\[
 H^1_{\mathrm{cts}}(G_{\mathbb Q},A_s)
                         \simeq\varprojlim_r B_r.
                                                        \tag{D.62.6}
\]
One may derive this directly from the cochain complexes:
\(C^\bullet(G_{\mathbb Q},A_s)=
\varprojlim_r C^\bullet(G_{\mathbb Q},W_r)\), and each
degreewise transition is surjective, since a continuous
finite-valued cochain can be lifted value by value.
This also proves (D.62.2), or one may apply
\cite[Proposition 1.6 and Theorem 2.2]{DE_Jannsen1988}.

For \(r=1\), the set \(S_r\) of (D.47.1) is just
\(\{5,\infty\}\): every Kummer valuation is integral,
so none has negative five-adic valuation. Thus every
integral \(A_s\)-cocycle, reduced to \(W_1\), factors
through \(G_{\mathbb Q,\{5,\infty\}}\). The module
\(W_1\simeq\mathbb F_5(2)\) is independent of the
abelian twist. Its finite-ramification first cohomology
is finite. In particular
\[
 \operatorname{im}\bigl(\varprojlim_r B_r\longrightarrow B_1\bigr)
                             \text{ is finite}.
                                                        \tag{D.62.7}
\]
The finiteness, rather than any vanishing calculation
for this group, is all that is needed.

If the countable tower \((B_r)\) were Mittag--Leffler,
let \(B'_r\) be the eventual image of \(B_j\) in
\(B_r\). Its transition maps would be surjective:
choose \(j\) large enough that both eventual images
have stabilized and lift from \(B_j\).
Recursive lifting then makes
\(\varprojlim_r B'_r\to B'_1\) surjective, and the
limits of \((B'_r)\) and \((B_r)\) agree. But by
(D.62.5), \(B'_1\) would be infinite-dimensional,
contradicting (D.62.7). Hence \((B_r)\) is not
Mittag--Leffler. For countable abelian groups the
converse to the usual derived-limit vanishing criterion
holds: \(\varprojlim^1 B_r=0\) would imply the
Mittag--Leffler condition. This is the Gray criterion
stated in \cite[Lemma 4.15, printed p.229]{DE_Jannsen1988}.
It proves (D.62.1).

\paragraph{The topology and the local invisibility.}
Give continuous \(A_s\)-valued cochains the uniform
\(\mathfrak m\)-adic topology: a basic neighborhood
of zero consists of cochains taking all their values
in \(\mathfrak m^rA_s\). Give the two-cocycle group
the subspace topology and \(H^2\) the quotient by
continuous coboundaries. Suppose a two-cocycle \(z\)
has zero class at every finite coefficient level.
For each \(r\), choose a continuous finite-valued
one-cochain \(b_r\) with
\(z\bmod\mathfrak m^r=db_r\), and lift its finitely
many values to \(A_s\). The resulting continuous
one-cochain \(\widetilde b_r\) satisfies
\[
                  z-d\widetilde b_r
                       \in C^2(G_{\mathbb Q},\mathfrak m^rA_s).
                                                        \tag{D.62.8}
\]
Thus \(z\) is in the closure of the coboundaries.
Conversely an approximating coboundary to this precision
makes the reduced class zero. The kernel in (D.62.2)
is therefore exactly the closure of zero for the
displayed quotient topology. Its nonzero size proves
non-Hausdorffness.

At every finite place \(v\), the groups
\(H^1(G_{\mathbb Q_v},W_r)\) are finite
\cite[I, Corollary 2.3]{DE_MilneADT}; at the real place
the analogous statement is immediate. Their derived
limit vanishes. The local analogue of (D.62.2) is
therefore injective into its inverse limit of finite
coefficient second cohomology. Restriction of a global
class in the first term of (D.62.2) is zero at every
such finite level, and consequently is zero locally.
This proves the final assertion without claiming a
global Hasse principle for the compact module.
\end{proof}


% D.63; independent root and PBC full proof audits passed.
\subsection{The separated global obstruction and its exact local convergence condition}

\begin{statement}{D.63}{Proved}{A noncompact separated obstruction for every Alexander twist}
Let \(s=(a,b)\), \(A_s=\Lambda C\), \(\mathfrak m\), and
\(W_r=A_s/\mathfrak m^rA_s\) be as in D.62. Arbitrary
infinite ramification is allowed. Give continuous cohomology
the uniform-cochain quotient topology of that statement and put
\[
 \mathcal H_s=H^2_{\mathrm{cts}}(G_{\mathbb Q},A_s)/
                                      \overline{\{0\}}.
                                                        \tag{D.63.1}
\]
For each rational prime \(q\ne5\), choose tame inertia and
an arithmetic Frobenius lift, and define the actual operators
\[
 \begin{split}
 \mu_q&=t_1^{a_q}t_2^{b_q},\qquad
 D_q=q^{-1}\rho_s(\sigma_q)-1,\\
 J_q&=(\mu_q-1)A_s+D_qA_s,\qquad L_q=A_s/J_q,\\
 F^rL_q&=(\mathfrak m^rA_s+J_q)/J_q.
 \end{split}                                             \tag{D.63.2}
\]
The sum defining \(J_q\) is a closed additive subgroup;
\(D_q\) is the actual arithmetic operator of D.35 and D.39,
not a weight approximation. Define
\[
 \mathcal L_s=\left\{(\lambda_q)_{q\ne5}\in\prod_{q\ne5}L_q:
  \text{for every }r,\quad
  \lambda_q\in F^rL_q\text{ for all but finitely many }q\right\}.
                                                        \tag{D.63.3}
\]
Its topology is the inverse-limit topology of the discrete
direct sums \(\bigoplus_{q\ne5}L_q/F^rL_q\).
There is a topologically exact sequence
\[
 0\longrightarrow\mathcal K_s\longrightarrow\mathcal H_s
       \xrightarrow{\mathrm{loc}_{\ne5}}\mathcal L_s
       \longrightarrow0,                               \tag{D.63.4}
\]
where \(\mathcal K_s\) is compact and the displayed map is
open. In particular every family satisfying precisely the
convergence condition (D.63.3) occurs as the localizations
away from five of a global continuous second-cohomology
class.

Here is the exact remaining compact term. Put
\[
 \begin{split}
 T_r&=H^0(G_{\mathbb Q},W_r^D)^*,\qquad
 W_r^D=\operatorname{Hom}(W_r,\mu_{5^\infty}),\\
 \mathcal T_s&=\varprojlim_r T_r,
 \qquad L_5=H^2_{\mathrm{cts}}(G_{\mathbb Q_5},A_s),\\
 \mathcal S_s&=\varprojlim_r
   \ker\left(H^2(G_{\mathbb Q},W_r)
                   \longrightarrow\bigoplus_vH^2(G_{\mathbb Q_v},W_r)\right).
 \end{split}                                             \tag{D.63.5}
\]
The star denotes the finite Pontryagin dual. All terms
in the inverse limit defining \(\mathcal S_s\) are finite.
There are exact sequences of compact groups
\[
 \begin{gathered}
 L_5\longrightarrow\mathcal T_s\longrightarrow0,\\
 0\longrightarrow\mathcal S_s\longrightarrow\mathcal K_s
     \longrightarrow\ker(L_5\longrightarrow\mathcal T_s)
     \longrightarrow0.
 \end{gathered}                                         \tag{D.63.6}
\]
These compact terms are retained, not asserted to vanish.

For every \(s\), \(\mathcal H_s\) is a complete,
countably based, separable Hausdorff group that is not
locally compact. It has a continuous quotient
\[
           \bigoplus_{q\equiv1\ (\mathrm{mod}\ 5)}\mathbb F_5.
                                                        \tag{D.63.7}
\]
Thus D.62's nonseparated kernel is accompanied by a large
separated obstruction, including when \(H^1(G_{\mathbb Q},A_s)=0\).
The inverse finite-localization kernel \(\mathcal S_s\), and
the isomorphism type of the nonseparated kernel for arbitrary
infinite ramification, remain separate unevaluated parts of
Task D.5. D.77 computes the other compact groups and the
actual map in (D.63.6) by the three boundary axes. D.64 proves
\(\mathcal K_s=0\) for every pair supported at five and
one prime of order twenty modulo twenty-five.
In contrast, D.75 exhibits split free factors of arbitrary
finite rank and a countable product in this compact kernel
for explicit cusp twists, while retaining its compact complement.
D.76 gives the joint three-axis free factor for every pair.
D.77 identifies its entire compact complement as precisely
\(\mathcal S_s\), giving a continuous product decomposition.
D.78 computes \(\mathcal S_s=\mathcal K_s=\mathbb Z/5\)
for the actual point pair \(\kappa(2),\kappa(3)\).
D.79 determines it for every invertible recombination of this
pair, and D.80 for arbitrary infinitely ramified pairs with
zero five-valuations, independent reductions modulo five,
and a primitive local unit-logarithm vector. The whole compact
kernel is respectively \(\mathbb Z/5\), \(\mathbb Z_5^d\),
or the countable compact product according to the exact local
boundary depth. D.81 removes the local-primitivity assumption and computes
the full compact kernel for every globally independent pair
with zero five-valuations, including the additional mixed
finite and free kernels. D.82 additionally evaluates every pair with two independent
residual inertia directions away from five, with arbitrary
five-valuations; the actual point \(t=10\) has zero compact
kernel. D.83 gives the stronger finite-level localization
isomorphism whenever the closed away-five valuation lattice
is generated by the local unit vector, including the zero
and nonprimitive cases. The last inverse limit remains open
for pairs outside these specified families.
\end{statement}

\begin{proof}
Write \(B_r=H^2(G_{\mathbb Q},W_r)\). The odd-primary
cohomological dimension of \(G_{\mathbb Q}\) is two;
equivalently apply \cite[I, Theorem 4.10(c)]{DE_MilneADT}
with all primes allowed, since the positive-degree
cohomology of a real decomposition group with five-primary
coefficients is zero. The coefficient long exact sequence
therefore makes \(B_{r+1}\to B_r\) surjective. Each
\(B_r\) is countable, by the finite-quotient argument in
D.62. That statement gives an algebraic identification
\[
                    \mathcal H_s\simeq\varprojlim_r B_r.
                                                        \tag{D.63.8}
\]
This is also topological. If a continuous two-cocycle has
zero class modulo \(\mathfrak m^r\), lift a finite-valued
one-cochain witnessing that fact and subtract its
coboundary. The resulting representative takes values
in \(\mathfrak m^rA_s\). Thus the kernel of reduction to
\(B_r\) is exactly the image of the uniform neighborhood
of zero at level \(r\). These kernels define both
topologies in (D.63.8). In particular \(\mathcal H_s\)
is a complete separable metrizable group.

\paragraph{The local quotients at every finite level.}
At \(q\ne5\), the original tripod action is unramified.
Its twisted action on inertia factors through five-primary
tame inertia, whose generator acts by \(\mu_q\).
Wild inertia and the other tame factors have no positive
five-primary cohomology here. The usual procyclic
resolution and the tame relation give
\[
 \begin{split}
 H^1(I_q,W_r)&=W_r/(\mu_q-1)W_r,\\
 H^2(G_{\mathbb Q_q},W_r)
   &=W_r/\bigl((\mu_q-1)W_r+D_qW_r\bigr)\
    =A_s/(J_q+\mathfrak m^rA_s).
 \end{split}                                             \tag{D.63.9}
\]
The Frobenius action on the first line is
\(q^{-1}\rho_s(\sigma_q)\), because the inertia generator
has cyclotomic weight one. Hochschild--Serre has no
other term in degree two: both the relevant inertia and
the residue-field group have five-cohomological dimension
one. This proves (D.63.9) even when \(\mu_q-1\) is a
zero divisor on the finite module, including \(\mu_q=1\).
Changing the Frobenius lift changes \(D_q\) only by an
operator whose image is contained in \((\mu_q-1)A_s\),
so the quotient is independent of that choice.

The two images defining \(J_q\) are compact, so their
sum is compact and closed. Finiteness of local first
cohomology eliminates its first derived inverse limit
\cite[I, Corollary 2.3]{DE_MilneADT}. Taking inverse
limits in (D.63.9) gives
\[
                   H^2_{\mathrm{cts}}(G_{\mathbb Q_q},A_s)=L_q.
                                                        \tag{D.63.10}
\]
In particular each finite local transition is onto,
also directly by its displayed quotient description.
Taking the inverse limit of the direct sums of (D.63.9)
gives exactly (D.63.3): its elements are compatible local
classes with finite support at every finite precision.
There is no common finite ramification set for \(A_s\)
in this argument.

\paragraph{Global reciprocity is absorbed at five.}
Each finite module \(W_r\) has finite ramification,
although the compact module \(A_s\) need not.
Apply finite-module global duality with the set of all
places. The last part of the Poitou--Tate sequence is
\[
 0\longrightarrow \operatorname{Sha}^2(\mathbb Q,W_r)
 \longrightarrow B_r
 \longrightarrow\bigoplus_v H^2(G_{\mathbb Q_v},W_r)
 \longrightarrow T_r\longrightarrow0.
                                                        \tag{D.63.11}
\]
Here \(\operatorname{Sha}^2\) denotes the kernel displayed
in (D.63.5), not an elliptic-curve Tate--Shafarevich group.
It is finite by \cite[I, Theorem 4.10(a)]{DE_MilneADT};
the exactness is part (b). The local summands are finite.
The last map is the sum of the local Tate pairings with
global invariant elements of \(W_r^D\).

The restriction
\(H^0(G_{\mathbb Q},W_r^D)\hookrightarrow
  H^0(G_{\mathbb Q_5},W_r^D)\)
is injective. Finite local duality therefore makes the
single component
\[
                  H^2(G_{\mathbb Q_5},W_r)\longrightarrow T_r
                                                        \tag{D.63.12}
\]
surjective. For any finite-support choice of local
classes away from five, add a class at five cancelling
its image in \(T_r\). Exactness of (D.63.11) then
realizes the resulting family globally. This is the
argument in \cite[I, Corollary 4.16 and its proof,
printed p.62]{DE_MilneADT}; it prescribes zero at all
other places as well, rather than merely prescribing
the originally chosen finite subset.

Let
\[
 K_r=\ker\left(B_r\longrightarrow
                      \bigoplus_{q\ne5}H^2(G_{\mathbb Q_q},W_r)\right).
\]
The real summand is zero. Equations (D.63.11)--(D.63.12)
give exact sequences
\[
 \begin{gathered}
 0\longrightarrow K_r\longrightarrow B_r
   \longrightarrow\bigoplus_{q\ne5}H^2(G_{\mathbb Q_q},W_r)
   \longrightarrow0,\\
 0\longrightarrow\operatorname{Sha}^2(\mathbb Q,W_r)
   \longrightarrow K_r
   \longrightarrow
       \ker\{H^2(G_{\mathbb Q_5},W_r)\longrightarrow T_r\}
   \longrightarrow0.
 \end{gathered}                                          \tag{D.63.13}
\]
In particular \(K_r\) is finite. All inverse systems
of the finite kernels in these rows are Mittag--Leffler.
Taking inverse limits proves (D.63.4) with
\(\mathcal K_s=\varprojlim_rK_r\), and proves (D.63.6).
No surjectivity of the individual transition maps on
the finite kernels is needed.

We also verify the asserted quotient topology. Fix
\(r\), and choose \(t\ge r\) so that the images
of \(K_j\to K_r\) have stabilized for \(j\ge t\).
The stable image equals the image of
\(\varprojlim K_j\to K_r\), by the finite inverse-limit
property. If \(\lambda\in\mathcal L_s\) is zero at
level \(t\), choose a global lift \(h\in\mathcal H_s\).
Then \(h_t\in K_t\), so \(h_r\) is the reduction of
some \(k\in\mathcal K_s\). The lift \(h-k\) is zero
at level \(r\). Thus the image of the \(r\)-th basic
open subgroup of \(\mathcal H_s\) contains the
\(t\)-th basic open subgroup of \(\mathcal L_s\).
This proves openness and topological exactness.

\paragraph{Noncompactness at every scale.}
At level one the action is always
\(W_1=\mathbb F_5(2)\), independent of \(s\).
Formula (D.63.9) gives
\[
 H^2(G_{\mathbb Q_q},W_1)=
 \begin{cases}\mathbb F_5,&q\equiv1\pmod5,\\
 0,&q\not\equiv1\pmod5
 \end{cases}\qquad(q\ne5).
                                                        \tag{D.63.14}
\]
The local transitions and therefore the transitions
between their direct sums are onto. Projection from
their inverse limit is onto, giving (D.63.7).

In fact no basic open subgroup of \(\mathcal L_s\)
is compact. Fix \(r\). Choose infinitely many primes
splitting the actual finite action on \(W_{r+1}\)
and \(\mu_{5^{r+1}}\); they exist by Chebotarev
\cite[Chapter V, Theorem 3.23]{PBC_MilneCFT}.
At each of these primes both terms in the local
differential vanish modulo \(\mathfrak m^{r+1}\),
so the local groups at levels \(r+1\) and \(r\)
are respectively \(W_{r+1}\) and \(W_r\).
Their nonzero kernel is
\(\mathfrak m^rA_s/\mathfrak m^{r+1}A_s\).
An arbitrary finite-support collection of these
kernel classes lifts to a compatible element of
\(\mathcal L_s\), using the surjectivity of its
finite-level transitions. The basic subgroup zero
at level \(r\) consequently has an infinite discrete
quotient at level \(r+1\). It cannot be compact.
Every neighborhood of zero contains one of these
closed open subgroups, so \(\mathcal L_s\) is not
locally compact. The open quotient map in (D.63.4)
then shows that \(\mathcal H_s\) is not locally
compact either. This completes the proof while
retaining exactly the compact ambiguity in (D.63.6).
\end{proof}


% D.64. Independent root and PBC full proof and primary-source audits passed.
\subsection{An evaluated family with one additional ramification prime}

\begin{statement}{D.64}{Proved}{The one-primitive-prime Alexander calculation and its exact section fibres}
Let \(q\ne5\) be a rational prime of order twenty in
\((\mathbb Z/25)^\times\), and put \(S=\{5,q,\infty\}\).
For arbitrary \(a_5,a_q,b_5,b_q\in\mathbb Z_5\), set
\[
 a=a_5\kappa(5)+a_q\kappa(q),\qquad
 b=b_5\kappa(5)+b_q\kappa(q),\qquad s=(a,b).
                                                        \tag{D.64.1}
\]
Use the actual Alexander module \(A_s=\Lambda C\) of D.35,
with \(\mathfrak m=(5,u,v)\). Its unrestricted first
cohomology and its finite-ramification second cohomology are
\[
 \begin{gathered}
 H^0(G_{\mathbb Q},A_s)=0,\\
 H^1(G_{\mathbb Q},A_s)=H^1(G_{\mathbb Q,S},A_s)
                   \simeq\prod_{\mathbb N}\mathbb Z_5,\\
 H^2(G_{\mathbb Q,S},A_s)
       \xrightarrow[\mathrm{res}_q]{\ \sim\ }
 H^2(G_{\mathbb Q_q},A_s)=L_q.
 \end{gathered}                                         \tag{D.64.2}
\]
The first isomorphism of compact modules is noncanonical.
The second is the actual restriction map, with
\[
 L_q=A_s/\bigl((t_1^{a_q}t_2^{b_q}-1)A_s+
                    (q^{-1}\rho_s(\sigma_q)-1)A_s\bigr),
                                                        \tag{D.64.3}
\]
where the sum is closed and the Frobenius operator is
the one in D.35 and D.39. In the primitive axis case
\(b_q=0\), \(a_q\in\mathbb Z_5^\times\), this compact
module is itself a countable product of \(\mathbb Z_5\)'s.
The same conclusion holds after permuting the punctures.

Let \({\cal C}_S\) and \(D_s\) be the actual residue
module and operator of D.48. There is an algebraic splitting
\[
 H^2(G_{\mathbb Q},A_s)
       \simeq L_q\oplus({\cal C}_S/D_s{\cal C}_S).
                                                        \tag{D.64.4}
\]
It is specified by the inverse of restriction in (D.64.2),
followed by inflation to \(G_{\mathbb Q}\); the retraction
is localization at \(q\). In particular no closure is
inserted in the second summand. At each finite coefficient
level \(W_r=A_s/\mathfrak m^rA_s\), localization is an
isomorphism
\[
 H^2(G_{\mathbb Q},W_r)
       \xrightarrow{\ \sim\ }
            \bigoplus_{v\ne5}H^2(G_{\mathbb Q_v},W_r).
                                                        \tag{D.64.5}
\]
Consequently D.63's compact kernel \(\mathcal K_s\) is zero:
the separated unrestricted \(H^2\) is topologically exactly
the local family \(\mathcal L_s\) of (D.63.3). The Jannsen
kernel is exactly
\[
 \frac{{\cal C}_S\cap
             \prod_{v\notin S}(v^{-1}\rho_s(\sigma_v)-1)A_s}
      {D_s{\cal C}_S}.
                                                        \tag{D.64.6}
\]
It is nonzero, as in D.48.

For the geometrically pro-five metabelian quotient
\(M=F_5/F_5''\), the entire unrestricted section fibre
above (D.64.1) is nonempty if and only if
\[
                       a_qb_q(a_q-b_q)=0.
                                                        \tag{D.64.7}
\]
Every section in a nonempty fibre is represented over
\(G_{\mathbb Q,S}\), and that fibre is a torsor under
\(\prod_{\mathbb N}\mathbb Z_5\). Its global cusp classes
are exactly the three explicit packet maps in D.35,
with parameters in
\(\mathbb Z_5\kappa(5)+\mathbb Z_5\kappa(q)\).
Thus this is an evaluated four-parameter family within
Task D.5; arbitrary finite ramification sets and arbitrary
infinite ramification remain outside this assertion.
\end{statement}

\begin{proof}
We first determine the actual arithmetic group needed
for the calculation. This step also supplies the integral
lifting control that would be missing from a comparison
of rational ranks of finite augmentation quotients.

\paragraph{The maximal pro-five group and its single relation.}
Put \(K=\mathbb Q(\zeta_5)\), let \(G=G_{K,\{5,q\}}\),
and let \(P\) be its maximal pro-five quotient. Since
\(q\) has order four modulo five, there is just one prime
of \(K\) above \(q\); its residue cardinality is
\(Q=q^4\). Moreover
\[
                         v_5(Q-1)=1.                  \tag{D.64.8}
\]
The class-number, units, and Brauer calculation in D.1
gives
\[
 \dim_{\mathbb F_5}H^1(G_{K,\{5\}},\mathbb F_5)=3,
 \qquad H^2(G_{K,\{5\}},\mathbb F_5)=0.
\]
Here \(\mu_5\subset K\), so the residual Tate twist is
trivial. Finite-coefficient localization at the one new
prime, whose residue field contains \(\mu_5\), now gives
\[
       \dim H^1(G,\mathbb F_5)=4,
       \qquad \dim H^2(G,\mathbb F_5)=1.               \tag{D.64.9}
\]
The residue map in degree one is onto, and the degree-two
residue map is an isomorphism. The degree-three term
vanishes by odd-primary cohomological dimension two
\cite[I, Theorem 4.10(c); II, Proposition 2.9]{DE_MilneADT}.

For clarity, the bound on the pro-five relation rank does
not presuppose a comparison theorem for all coefficients.
If \(N=\ker(G\to P)\), then
\(H^1(N,\mathbb F_5)^P=0\). Otherwise the kernel of a
nonzero invariant homomorphism would give a quotient of
\(G\) which is an extension of \(P\) by \(\mathbb F_5\),
hence a pro-five quotient with nontrivial image of \(N\),
contrary to maximality. Hochschild--Serre therefore gives
\[
 H^1(P,\mathbb F_5)=H^1(G,\mathbb F_5),\qquad
 H^2(P,\mathbb F_5)\lhook\joinrel\longrightarrow
                                      H^2(G,\mathbb F_5).
                                                        \tag{D.64.10}
\]
Choose the images \(\tau,\sigma\) of a five-primary tame
inertia generator and an arithmetic Frobenius of \(K_q\).
They satisfy the actual relation
\[
                        \sigma\tau\sigma^{-1}=\tau^Q.
                                                        \tag{D.64.11}
\]
Their Frattini classes in \(P/\Phi(P)\) are independent.
Indeed the surjective residue map just noted supplies
a character nonzero on \(\tau\). The cyclotomic degree-five
extension \(K(\mu_{25})/K\) supplies a character zero on
\(\tau\) and nonzero on \(\sigma\), since its Frobenius
is \(q^4\not\equiv1\pmod{25}\). Thus they extend to four
minimal generators. Present \(P\) as a quotient of the
free pro-five group \(F=F(\sigma,\tau,h_1,h_2)\), with
kernel \(R\subset\Phi(F)\).

The element
\(r_0=\sigma\tau\sigma^{-1}\tau^{-Q}\) belongs to \(R\).
Its image in \(F^{\rm ab}\simeq\mathbb Z_5^4\) is
\((1-Q)e_\tau\), which has valuation exactly one.
On the other hand the image of
\(R^5[R,F]\), with closures understood, is contained
in \(25\mathbb Z_5^4\), because
\(R\subset\Phi(F)\) has abelianization contained in
\(5\mathbb Z_5^4\). Hence
\[
       0\ne[r_0]\in R/\overline{R^5[R,F]}.
                                                        \tag{D.64.12}
\]
The dual of this relation module is
\(H^2(P,\mathbb F_5)\): apply Hochschild--Serre to the
free presentation, using the minimality of its four
generators. By (D.64.9)--(D.64.10) the module is
one-dimensional. Thus \(r_0\) normally generates \(R\).
One can verify the last implication directly: a nonzero
quotient of \(R\) by its closed normal closure would be
detected in a finite five-group quotient of \(F\).
The resulting nontrivial normal subgroup has a nonzero
coinvariant Frattini quotient under that finite five-group;
the augmentation ideal of its group algebra is nilpotent.
This gives an \(F\)-invariant homomorphism
\(R\to\mathbb F_5\) vanishing on \(r_0\), contradicting
(D.64.12) and dimension one.

It follows that, as actual pro-five groups with the
specified local subgroup,
\[
       P=D*F(h_1,h_2),\qquad
       D=\langle\sigma,\tau\mid
                  \sigma\tau\sigma^{-1}=\tau^{q^4}\rangle
                \simeq\mathbb Z_5\rtimes_{q^4}\mathbb Z_5.
                                                        \tag{D.64.13}
\]
The last presentation has exactly the indicated
semidirect-product meaning. Collecting words puts them
in the form \(\tau^a\sigma^b\); the compact image of
\(\mathbb Z_5^2\) is closed and contains the dense word
subgroup. The homomorphism to the displayed semidirect
product distinguishes these normal forms. In particular
\(D\) is the actual maximal pro-five local group of
\(K_q\), and is embedded as a free factor of \(P\).

\paragraph{The coefficient comparison and the local degree-two isomorphism.}
For a finite five-primary \(P\)-module \(W\), restriction
for the free product gives
\[
                 H^n(P,W)\simeq H^n(D,W)\quad(n\ge2).
                                                        \tag{D.64.14}
\]
Here the other factor is free and has cohomological
dimension one. This is the finite-index-set case of
\cite[Theorem 4.2]{DE_GaertnerWingberg}; take the
corestricting subgroups to be the factors themselves.
Alternatively the universal property of crossed
homomorphisms decomposes the completed augmentation
ideal of a free product as the direct sum of the two
induced augmentation ideals. Dimension shifting and
Shapiro give precisely (D.64.14). The group \(D\) is
an extension of two copies of \(\mathbb Z_5\), so its
cohomological dimension is two. Thus the same is true
of \(P\).

For \(W=\mathbb F_5\), inflation \(P\to G\) is now an
isomorphism in every degree: degrees zero and one were
already identified, degree two follows from (D.64.9),
(D.64.10), and (D.64.12), and higher degrees vanish.
Every finite five-primary \(P\)-module has a filtration
with trivial \(\mathbb F_5\) factors. Its action factors
through a finite five-group, whose simple modules in
characteristic five are trivial. Induction on such a
filtration and the two cohomology long exact sequences
therefore prove inflation isomorphisms in all degrees
for every such \(W\). This proves the needed full/pro-five
comparison; it was not assumed in the relation-rank
calculation.

Each actual module \(W_r=A_s/\mathfrak m^rA_s\), restricted
to \(G\), is a \(P\)-module. Its augmentation and
five-adic filtrations have trivial residual factors
over \(K\), since \(\chi_5\equiv1\pmod5\); its finite
action consequently has five-power image. The same
observation for the graph of a crossed homomorphism
applies to all its finite cocycles. We may therefore
use (D.64.14) and the comparison just proved for these
actual twisted coefficients.

Over \(\mathbb Q\) the relevant quotient is an extension
\(\mathcal P\) of \(C_4=\operatorname{Gal}(K/\mathbb Q)\)
by \(P\). The local decomposition subgroup at \(q\)
surjects onto \(C_4\) and intersects \(P\) in the
embedded factor \(D\). Choose a complement \(T\simeq C_4\)
inside that local subgroup, by the coprime profinite
Schur--Zassenhaus theorem. In particular \(T\) normalizes
\(D\). The isomorphism (D.64.14) is \(T\)-equivariant.
Taking its invariants is exact because four is a unit
on every coefficient module. Local inflation from
\(D\) over \(K_q\) is also a cohomology isomorphism:
the discarded wild and prime-to-five tame factors have
zero positive five-primary cohomology, and the
remaining residue action \(Q\) lies in \(1+5\mathbb Z_5\).
Equivalently its two-generator tame resolution gives
the same finite cohomology. We obtain
\[
 H^2(G_{\mathbb Q,S},W_r)
      \xrightarrow[\mathrm{res}_q]{\ \sim\ }
                         H^2(G_{\mathbb Q_q},W_r).
                                                        \tag{D.64.15}
\]
All finite-ramification and all local finite-coefficient
cohomology groups are finite. Their towers are
Mittag--Leffler. The continuous comparison therefore
passes (D.64.15) to \(A_s\) without a derived-limit
correction. This proves the second line of (D.64.2),
with (D.64.3) supplied by the explicit local resolution
in D.63. The primitive-axis evaluation is D.43, since
\(q-1\) is a five-adic unit.

\paragraph{The integral rank of first cohomology.}
The zero-inertia argument of D.47, or D.48, identifies
unrestricted \(H^1\) with \(H^1(G_{\mathbb Q,S},A_s)\).
It is compact, countably based, and torsion-free by
D.35 and D.47. If \(a_q=b_q=0\), the action is already
unramified outside five, and D.36 proves the asserted
countably infinite rank.

Suppose now \((a_q,b_q)\ne(0,0)\). Applied over the finite
local extension \(K_q/\mathbb Q_q\), D.43 gives the
integral vanishings
\[
                       H^0(D,A_s)=H^1(D,A_s)=0.
                                                        \tag{D.64.16}
\]
Consequently every \(P\)-cocycle has a unique cohomologous
representative whose restriction to \(D\) is zero: subtract
the unique coboundary that is its restriction. By the
actual free product (D.64.13), the two values on
\(h_1,h_2\) can then be arbitrary elements of \(A_s\).
Indeed the universal property applies to their graphs
inside the pro-five semidirect product \(A_s\rtimes P\).
It gives a continuous bijection of compact modules,
hence a topological isomorphism
\[
              H^1(P,A_s)\simeq A_s\oplus A_s.
                                                        \tag{D.64.17}
\]
No quotient by residual coboundaries remains, by the
zero \(H^0\) in (D.64.16). Since \(T\) normalizes \(D\),
the space of these normalized cocycles is \(T\)-stable.
Exact averaging identifies the desired global group
with \(H^1(P,A_s)^T\).

Here is an integral argument that these invariants
have infinite rank. Reduce (D.64.17) modulo five and
filter it by powers of \(I=(u,v)\). The filtration is
\(T\)-stable: its terms are exactly the normalized
cocycles with all values in \(I^jA_s\), modulo five.
Let
\[
 W=H^1\bigl(P/\overline{\langle\!\langle D\rangle\!\rangle},
                                      \mathbb F_5\bigr).
\]
This is a two-dimensional \(T\)-module. On the associated
graded piece in degree \(j\), the coefficient action is
\(\mathbb F_5(j+2)^{j+1}\), and the \(P\)-action is
trivial. Evaluating the cocycle on words consequently
uses only the Frattini abelianization of the free
quotient. Thus the associated graded \(T\)-module is
\[
                 \mathbb F_5(j+2)^{j+1}\otimes W.
                                                        \tag{D.64.18}
\]
The four characters of \(T\) are defined over
\(\mathbb F_5\), and every representation is semisimple.
Choose a character occurring in the nonzero module
\(W\). For infinitely many \(j\), its product with
\(\chi_5^{j+2}\) is trivial. At each such \(j\), the
invariant dimension in (D.64.18) is at least \(j+1\).
Invariants are exact, and every finite augmentation
quotient of the actual module in (D.64.17) is a
surjective quotient. Therefore
\(H^1(P,A_s)^T/5\) surjects onto finite modules of
unbounded dimension. Its compact torsion-free
\(\mathbb Z_5\)-rank is infinite. D.47's countability
then identifies it with \(\prod_{\mathbb N}\mathbb Z_5\).
This argument uses actual arbitrary integral values on
the free generators; it does not infer inverse-limit
surjectivity from finite-index images at successive
augmentation levels.

\paragraph{The unrestricted obstruction and its separated kernel.}
At a fixed coefficient level, localization outside
\(S\) identifies the residue at \(v\notin S\) with
\(H^2(G_{\mathbb Q_v},W_r)\). If a global class is zero
locally at every \(v\ne5\), its residues outside \(S\)
vanish, so it comes from \(H^2(G_{\mathbb Q,S},W_r)\).
Its localization at \(q\) is zero, and (D.64.15) then
makes that class zero. Thus the localization map in
(D.64.5) is injective. It is surjective by the exact
Poitou--Tate argument in D.63: the component at five
can cancel the reciprocity functional, leaving every
prescribed finite-support family away from five.
This proves (D.64.5).

In particular \(H^2(G_{\mathbb Q,S},W_r)\) injects
into the unrestricted group, already because its
restriction at \(q\) is an isomorphism. The compact
term \({\cal L}_s\) in (D.48.2) therefore vanishes,
as does every finite kernel whose inverse limit is
\(\mathcal K_s\) in D.63. D.48 gives (D.64.6), and
D.63 supplies the asserted exact topological
identification of the separated quotient. The
restriction map at \(q\), together with (D.64.2),
is a retraction of the first term in (D.48.1).
Splitting that short exact sequence proves (D.64.4)
as an assertion about the actual, possibly
non-Hausdorff continuous cohomology groups.

\paragraph{Existence and the complete metabelian fibre.}
The arithmetic action on \(M\) and the abelian pair
both factor through \(G_{\mathbb Q,S}\). The pullback
of \(M\rtimes G_{\mathbb Q,S}\) over the exact abelian
section has kernel \(A_s\). It has a continuous
set-theoretic section: use the normal-form lift
\(x^{a(g)}y^{b(g)}\) and the arithmetic element \(g\).
Its multiplication defect is consequently an actual
continuous two-cocycle in \(A_s\). The extension has
a continuous group section if and only if its class
in \(H^2(G_{\mathbb Q,S},A_s)\) is zero. By (D.64.2),
this happens if and only if its local class at \(q\)
is zero.

Since \(q\not\equiv1\pmod5\), local Kummer theory gives
\(H^1(G_{\mathbb Q_q},\mathbb Z_5(1))=\mathbb Z_5\),
via valuation, and every local unit has zero completed
Kummer class. The local abelian pair is therefore
exactly \((a_q\kappa(q),b_q\kappa(q))\). If it is
nonzero, D.42 forces any metabelian lift into a cusp
packet and gives the necessity of (D.64.7).
Conversely each of its three lines is the abelian
image of the corresponding actual local cusp packet,
with an arbitrary completed tangent parameter.
The zero pair also has a unit cusp lift. Hence
(D.64.7) is exactly the local, and therefore the
global, existence criterion.

At every prime outside \(S\) an unrestricted section
above (D.64.1) has zero abelian inertia. D.42 makes
its entire metabelian inertia zero. Thus every such
section factors through \(G_{\mathbb Q,S}\). The
conjugacy and stabilizer analysis in D.35 identifies
its nonempty fibre with a torsor under the first
group in (D.64.2), without an additional abelian
stabilizer quotient. Finally the three full packet
formulas of D.35 show that a cusp section has its
parameter supported on \(\{5,q\}\) exactly when its
abelian pair does. They give precisely the cusp
classes asserted in the statement.
\end{proof}


% D.65. Root and PBC independent full proof audits passed.
% The intrinsic five-adic topology is distinguished from cohomology topology.
\subsection{The abstract module carried by the nonseparated residue kernel}

\begin{statement}{D.65}{Proved}{The continuum-sized complete module in every finite-ramification Jannsen kernel}
Let \(s\) be any finitely ramified Alexander twist, and
retain \(S\), \({\cal C}_S\), \(D_{s,q}\), and the compact
group \({\cal L}_s\) of D.48. Put
\[
 \begin{split}
 {\cal R}_s&=
  \frac{{\cal C}_S\cap\prod_{q\notin S}D_{s,q}A_s}
       {D_s{\cal C}_S},\\
 {\cal J}_s&=\varprojlim\nolimits_r^1
                         H^1(G_{\mathbb Q},A_s/\mathfrak m^rA_s).
 \end{split}                                             \tag{D.65.1}
\]
For a set \(I\), define the complete free module
\[
 \widehat{\bigoplus}_{i\in I}\mathbb Z_5
   :=\varprojlim_n\bigoplus_{i\in I}\mathbb Z/5^n
   =\left\{(a_i)\in\prod_I\mathbb Z_5:
       \forall n,\ a_i\in5^n\mathbb Z_5
                         \text{ for all but finitely many }i\right\}.
                                                        \tag{D.65.2}
\]
There are noncanonical isomorphisms of abstract
\(\mathbb Z_5\)-modules
\[
 \begin{split}
 {\cal R}_s&\simeq
              \widehat{\bigoplus}_{2^{\aleph_0}}\mathbb Z_5,\\
 {\cal J}_s&\simeq {\cal L}_s\oplus
              \widehat{\bigoplus}_{2^{\aleph_0}}\mathbb Z_5.
 \end{split}                                             \tag{D.65.3}
\]
In particular the residue kernel is torsion-free, is
complete and separated for its intrinsic five-adic
topology, and has
\[
           \dim_{\mathbb F_5}({\cal R}_s/5{\cal R}_s)
                                      =2^{\aleph_0}.
                                                        \tag{D.65.4}
\]
The first isomorphism in (D.65.3) is also a homeomorphism
for these intrinsic five-adic topologies.
Equivalently its right-hand complete free module is
abstractly isomorphic to \(\prod_{\mathbb N}\mathbb Z_5\).
Under this equivalence the intrinsic topology on the
product is the uniform five-adic topology, not its
usual compact product topology.

These are not the subspace topologies from continuous
cohomology. In that topology \({\cal J}_s\) is the
closure of zero, as in D.62. The second splitting in
(D.65.3) is only an abstract module splitting; it does
not replace any cohomology group by a separated
quotient or assert a continuous retraction onto the
compact group with its original topology.

For the families of D.44 and D.64,
\({\cal L}_s=0\), so (D.65.3) gives the complete
abstract isomorphism type of their Jannsen kernel.
The same complete module is the kernel in each of
the three Tate weights of D.1. For a general finitely
ramified twist, the compact term \({\cal L}_s\)
is retained. This result makes no classification
claim for the whole kernel of an infinitely
ramified twist.
\end{statement}

\begin{proof}
We give the module and completeness arguments because
the two topologies here have different purposes.
Identify \(A_s\) additively with
\(\Lambda=\mathbb Z_5[[u,v]]\), and write
\[
 \begin{split}
 P&=\prod_{q\notin S}\Lambda,\qquad C={\cal C}_S,\\
 H&=\{(a_q)\in P:(D_{s,q}a_q)_q\in C\}.
 \end{split}                                             \tag{D.65.5}
\]
Each \(D_{s,q}\) is injective by D.48. Taking its unique
coordinatewise inverse on the image therefore identifies
\[
                             {\cal R}_s=H/C             \tag{D.65.6}
\]
as abstract modules. Also \(C\subset H\): every operator
\(D_{s,q}\) preserves \(\mathfrak m^n\), uniformly in
\(q\), because the actual arithmetic action does.

\paragraph{A stronger complete topology for the module calculation.}
For this paragraph only, give \(P\) the uniform
five-adic topology, whose neighborhoods are \(5^nP\).
The ring \(\Lambda\) is complete and separated for its
five-adic topology, so \(P\) is complete for this
uniform topology. The module \(C\) is closed in it.
Indeed if \(a\) can be approximated by elements of
\(C\) modulo every \(5^nP\), fix \(n\) and choose one
such approximation \(c\). Away from a finite set,
\(c_q\in\mathfrak m^n\), and
\(a_q-c_q\in5^n\Lambda\subset\mathfrak m^n\).
Thus \(a\in C\).

The coordinatewise operator \(D_s:P\to P\) is
continuous for the uniform five-adic topology:
each coordinate is \(\mathbb Z_5\)-linear and preserves
\(5^n\Lambda\), with the same bound for all coordinates.
Hence \(H=D_s^{-1}(C)\) is closed and complete as well.
Both \(C\) and \(H\) are five-saturated in \(P\).
For \(C\), use the exact order formula
\[
 \operatorname{ord}_{\mathfrak m}(5^n z)
          =n+\operatorname{ord}_{\mathfrak m}(z)
                        \quad(z\ne0).                  \tag{D.65.7}
\]
It follows either from coefficients or from the
polynomial integral-domain associated graded ring
of \(\mathbb Z_5[[u,v]]\) at \((5,u,v)\).
Thus \((5^na_q)_q\) tends to zero in \(\mathfrak m\)
if and only if \((a_q)_q\) does. For \(H\), apply the
same observation to \(D_s(5^na)=5^nD_s(a)\).
Consequently
\[
 H\cap5^nP=5^nH,\qquad C\cap5^nP=5^nC,
                  \qquad H/C\text{ is torsion-free}.
                                                        \tag{D.65.8}
\]
The topology on \(H\) induced from \(P\) is therefore
its intrinsic five-adic topology, and its quotient
topology on \(H/C\) has neighborhoods precisely
\(5^n(H/C)\).

This quotient is complete and separated. Separation
is the closedness of \(C\) just proved. For completeness,
take a compatible system of classes modulo
\(5^n(H/C)\). Choose representatives \(h_n\in H\)
successively. Compatibility lets one change the next
representative by an element of \(C\) so that
\(h_{n+1}-h_n\in5^nH\). Their limit exists in the
complete module \(H\), and its class gives the required
limit. This proves the asserted completeness directly;
it does not presume completeness of a quotient for
the original cohomology topology.

\paragraph{The exact number of five-adic generators.}
The Chebotarev construction in the last paragraph
of D.48 supplies distinct primes \(q_j\notin S\)
with
\[
                         D_{s,q_j}(1)\in\mathfrak m^j.
                                                        \tag{D.65.9}
\]
For any \(a=(a_j)\in\prod_j\mathbb Z_5\), put
\(h(a)_{q_j}=a_j\), with all other coordinates zero.
Then \(h(a)\in H\). Its class in \(H/C\) is the
copy of \(\prod_j\mathbb Z_5/c_0(\mathbb N,\mathbb Z_5)\)
already constructed there. It remains injective
after reduction modulo five. Indeed if
\(h(a)=5b+c\), with \(b\in H\), \(c\in C\), evaluation
at \(u=v=0\) gives
\[
                 a_j=5b_{q_j}(0,0)+c_{q_j}(0,0).
\]
The last constants tend to zero five-adically,
so \(a_j\bmod5\) is zero for all sufficiently large
\(j\). This is exactly the condition that \(a\)
be zero in the reduction modulo five of that
product-by-\(c_0\) quotient. Thus
\[
 \mathbb F_5^{\mathbb N}/\mathbb F_5^{(\mathbb N)}
             \lhook\joinrel\longrightarrow (H/C)/5(H/C).
                                                        \tag{D.65.10}
\]
The vector space on the left has dimension
\(2^{\aleph_0}\): over a finite field an infinite
dimensional vector space has cardinality equal
to its dimension, and quotienting by a countable
dimensional subspace does not change this dimension.
On the other hand \(P\), hence \(H/C\), has cardinality
at most \(2^{\aleph_0}\). This proves (D.65.4).

For completeness, a separated complete torsion-free
\(\mathbb Z_5\)-module \(M\) is the complete free module
on a basis of \(M/5M\). Choose arbitrary lifts of that
basis. A family of coefficients tending to zero
five-adically has a well-defined sum in \(M\).
Successively lifting the finitely supported expansion
of each residual class modulo five proves that these
sums exhaust \(M\). To check injectivity, take the
smallest valuation of a nonzero coefficient family,
divide by that power of five, and reduce modulo five;
the resulting nonempty finite basis combination is
nonzero. The map is an isomorphism and preserves
every subgroup \(5^nM\), hence is a homeomorphism
for the intrinsic topologies. Apply this to
\(M=H/C\), using (D.65.4), to obtain the first line
of (D.65.3).
The module \(\prod_{\mathbb N}\mathbb Z_5\) is also
intrinsically complete, separated, and torsion-free,
and its reduction is \(\mathbb F_5^{\mathbb N}\),
of the same continuum dimension. The same classification
proves the equivalent countable-product description.

\paragraph{Splitting off the retained compact term.}
D.48 gives the actual exact sequence of modules
\[
       0\longrightarrow{\cal L}_s
         \longrightarrow{\cal J}_s
         \longrightarrow{\cal R}_s\longrightarrow0.
                                                        \tag{D.65.11}
\]
Since the last module is torsion-free, the first
inclusion is pure: for every \(n\),
\({\cal L}_s\cap5^n{\cal J}_s=5^n{\cal L}_s\).
We spell out why compactness of \({\cal L}_s\), with
its original compact-module topology, gives an
abstract retraction. Introduce a variable \(x_j\)
with values in \({\cal L}_s\) for each
\(j\in{\cal J}_s\), and impose all equations
\[
 x_{j+j'}=x_j+x_{j'},\qquad
 x_{\lambda j}=\lambda x_j\ (\lambda\in\mathbb Z_5),
                 \qquad x_l=l\ (l\in{\cal L}_s).
                                                        \tag{D.65.12}
\]
Each finite subsystem, viewed as equations with
constants in \({\cal L}_s\), has the evident solution
\(x_j=j\) in \({\cal J}_s\). Smith normal form over
the discrete valuation ring \(\mathbb Z_5\) reduces
solvability of any such finite linear system to
divisibility of constants by powers of five and
zero equations. Purity makes those divisibilities
solvable already in \({\cal L}_s\). Thus every finite
subsystem has a solution in that compact module.
All equations are closed in the compact product
\({\cal L}_s^{{\cal J}_s}\). The finite-intersection
property gives a simultaneous solution, hence an
abstract \(\mathbb Z_5\)-linear retraction of the
inclusion in (D.65.11). This proves its splitting
and the second line of (D.65.3). No continuity of
this retraction for the nonseparated cohomology
topology has been asserted.

For D.44 the compact finite-support second cohomology
is zero, and for D.64 its maps to unrestricted finite
cohomology are injective. In both cases D.48's
\({\cal L}_s\) is therefore zero. Finally the three
Tate kernels in D.1 have the same proof with
\(\Lambda=\mathbb Z_5\), \(\mathfrak m=(5)\), and
\(D_q=q^{j-1}-1\). These operators are injective,
their arbitrarily divisible unit-input values are
given by the primes used in D.1, and there is no
compact base term. Thus their complete abstract
module type is the same one in (D.65.3).
\end{proof}


% D.66. Independent root, PBC, and FGH complete proof audits passed.
\subsection{Generic rigidity of completed cusp parameters}

\begin{statement}{D.66}{Proved}{A comeagre zero-cohomology locus and generic singleton cusp fibres}
Let \({\cal K}=H^1(G_{\mathbb Q},\mathbb Z_5(1))\), with
its intrinsic five-adic topology, and let \(A_s\) be the actual
Alexander module of D.35 for \(s=(a,b)\in{\cal K}^2\).
The set
\[
 {\cal Z}=\{(a,b):H^1_{\mathrm{cts}}(G_{\mathbb Q},A_{(a,b)})=0\}
                                                        \tag{D.66.1}
\]
contains a dense \(G_\delta\) subset of \({\cal K}^2\).
For every fixed \(b\in{\cal K}\), its intersection with
\({\cal K}\times\{b\}\) also contains a dense \(G_\delta\)
subset of that space. Thus the zero-cohomology locus is
comeagre both globally and in each such fibre. These dense
\(G_\delta\)'s may be chosen so that \(a\) is ramified at
infinitely many rational primes.

In particular, for a comeagre set of completed parameters
\(a\) in the cusp packet at zero, the abelian class \((a,0)\)
has exactly one geometrically pro-five metabelian section
class above it. It is supplied by the actual full cusp
section. The analogous statement holds in each of the other
two cusp packets, with its own completed parameter.

For comparison, the set of all abelian pairs which admit
any geometrically pro-five metabelian lift is contained
in a closed nowhere-dense subset of \({\cal K}^2\).
The assertion about generic singleton fibres is therefore
made in the cusp-parameter space, where existence is known;
no generic existence statement in \({\cal K}^2\) is asserted.
All these category statements use the intrinsic uniform
five-adic topology on completed Kummer classes, not the
product topology on their prime coordinates.
\end{statement}

\begin{proof}
\paragraph{The parameter space and its actual finite actions.}
Fix compatible roots of every positive rational prime and
write \(\kappa(q)\) for their continuous integral Kummer
cocycles. Kummer theory identifies the parameter space with
\[
 {\cal K}=c_0(\{q\text{ rational prime}\},\mathbb Z_5),
 \qquad a=\sum_q a_q\kappa(q).                         \tag{D.66.2}
\]
Indeed the finite Kummer groups are direct sums of
\(\mathbb Z/5^n\) indexed by the primes; the sign contributes
nothing at five. Their inverse limit is exactly the displayed
\(c_0\) module. The series in (D.66.2) converges uniformly as
a series of continuous cocycles, since \(a_q\to0\).
It supplies fixed representatives for all parameters,
linearly and continuously in their uniform five-adic topology.

The metric with balls \(a+5^n{\cal K}\) is complete.
A uniform Cauchy family has a coordinatewise limit, and its
uniform approximation by one \(c_0\) family shows that the
limit is still \(c_0\). Finite-support families with integer
coordinates form a countable dense subset. Thus \({\cal K}\),
\({\cal K}^2\), and each fixed-second-coordinate space are
complete separable metric spaces. We shall only use the
Baire conclusion that an intersection of countably many
dense open sets is dense. Here it follows directly by
choosing successively closed five-adic balls of radii tending
to zero inside the required open sets; completeness supplies
their common point.

Put \(\Lambda=\mathbb Z_5[[u,v]]\),
\(\mathfrak m=(5,u,v)\), and
\(W_r=A_s/\mathfrak m^rA_s\), for \(r\ge1\). For a
pair \(s=(a,b)\), set
\[
 S_r(s)=\{5,\infty\}\cup
       \{q:(a_q,b_q)\notin5^r\mathbb Z_5^2\}.          \tag{D.66.3}
\]
It is finite. The actual reduced action \(\rho_{s,r}\)
factors through \(G_{\mathbb Q,S_r(s)}\). Both this
action and the set \(S_r(s)\) are constant on every
coset of \(5^r{\cal K}^2\): the actual action formula
of D.35 and the estimate
\(t_i^{5^r c}-1\in\mathfrak m^{r+1}\) prove this
for the action, and (D.66.3) proves it for the support.
Define the finite set of actual cocycles
\[
 E_r(s)=Z^1(G_{\mathbb Q,S_r(s)},W_r).                  \tag{D.66.4}
\]
Its finiteness is the finite-coefficient finite-ramification
finiteness used in D.47, together with finiteness of the
group of coboundaries.

We need the exact uniform support conclusion of D.47:
every integral \(A_s\)-valued cocycle reduces to an
element of (D.66.4). This is a statement about its actual
finite cocycle, not merely its cohomology class. At a prime
outside (D.66.3), nonzero abelian inertia gives an integral
local coboundary by D.43. Its inertia values are multiples
of
\(t_1^{a_q}t_2^{b_q}-1\in\mathfrak m^{r+1}\).
When the abelian inertia is zero, the integral leading-weight
argument of D.47 makes the cocycle inertia zero as well.
Thus both the finite action and finite cocycle kill all
inertia outside the displayed finite set.

\paragraph{An open condition at each precision.}
Call a cocycle \(z\in E_r(s)\) exceptional if its
restriction to \(\ker\rho_{s,r}\) is nonzero. Let
\({\cal U}_r\subset{\cal K}^2\) be the set of pairs
such that, for every exceptional \(z\), there is a prime
\(q\notin S_r(s)\) with
\[
 \rho_{s,r}(\operatorname{Frob}_q)=1,
 \qquad z(\operatorname{Frob}_q)\ne0,
 \qquad (a_q,b_q)\ne(0,0).                            \tag{D.66.5}
\]
The first two conditions are well defined on a Frobenius
conjugacy class. In the finite affine representation
\(g\mapsto(z(g),\rho_{s,r}(g))\), conjugation preserves
the identity linear part and sends a nonzero translation
to a nonzero translation.

The set \({\cal U}_r\) is open. On a fixed
\(5^r{\cal K}^2\) coset the finite data, exceptional
cocycles, and the Frobenius tests are fixed. For each of
the finitely many exceptional cocycles, the last condition
in (D.66.5), at a chosen witnessing prime, is open in
the parameter coordinates. Existence of a choice of
witnesses is a union of these open conditions.

Moreover \({\cal U}_r\) is dense even when the second
coordinate is held fixed. Start from any pair
\(s_0=(a_0,b_0)\) and any neighborhood
\(a_0+5^N{\cal K}\), with \(N\ge0\). Choose
\(M\ge\max\{N,r\}\). For each exceptional cocycle
choose \(\delta\in\ker\rho_{s_0,r}\) with
\(z(\delta)\ne0\). Its finite affine image contains
\((z(\delta),1)\). Chebotarev, applied to this actual
finite Galois image, gives infinitely many primes
satisfying the first two tests in (D.66.5); see
\cite[Chapter V, Theorem 3.23]{PBC_MilneCFT}.
Choose distinct such primes \(q_z\), avoiding
\(S_r(s_0)\) and the finite set where
\(a_{0,q}\notin5^{M+1}\mathbb Z_5\). Put
\[
          a_1=a_0+5^M\sum_{z\text{ exceptional}}\kappa(q_z),
          \qquad b_1=b_0.                             \tag{D.66.6}
\]
This is a finite sum. It preserves all the finite data
at precision \(r\), stays in the prescribed neighborhood,
and makes \(a_{1,q_z}\) have valuation exactly \(M\).
No cancellation with the original parameter is possible.
It therefore gives all the required witnesses and lies in
\({\cal U}_r\). Openness also supplies a whole smaller
neighborhood with these same witnesses.

Let \({\cal U}\) be the open set of pairs with nonzero
abelian inertia at some prime \(q\ne5\). It is dense,
also in every fixed-second-coordinate space, by adding
an arbitrarily deep first-coordinate Kummer class at a
fresh prime. Hence
\[
            {\cal G}={\cal U}\cap\bigcap_{r\ge1}{\cal U}_r
                                                        \tag{D.66.7}
\]
is a dense \(G_\delta\), both in \({\cal K}^2\) and
in each fixed-second-coordinate space.

\paragraph{All integral cocycles vanish in cohomology on this set.}
Fix \(s\in{\cal G}\) and an actual continuous integral
cocycle \(c:G_{\mathbb Q}\to A_s\). Its reduction
\(c_r\) lies in \(E_r(s)\) by the uniform support
argument. It cannot be exceptional. Otherwise choose its
witness \(q\) in (D.66.5). Since \(s\) has nonzero
abelian inertia there, D.43 gives an integral element
\(v_q\in A_s\) with
\[
             c(g)=(\rho_s(g)-1)v_q
                         \quad(g\in G_{\mathbb Q_q}).
                                                        \tag{D.66.8}
\]
Evaluating at a Frobenius lift and reducing modulo
\(\mathfrak m^r\) gives zero, since its reduced action
is the identity. This contradicts the second test in
(D.66.5). Consequently \(c_r\) vanishes on
\(\ker\rho_{s,r}\), for every \(r\).

It follows that \(c\) is zero on \(\ker\rho_s\),
and descends to a continuous cocycle on the actual compact
operator image \(I_s=\rho_s(G_{\mathbb Q})\).
Continuity follows from the compact quotient map onto
\(I_s\); no abstract image is substituted for it.
Choose the nonzero inertia prime supplied by
\({\cal U}\), and put \(\mu=t_1^{a_q}t_2^{b_q}\).
Its operator image is the closed subgroup
\[
        N=\{\text{multiplication by }\mu^c:c\in\mathbb Z_5\}.
                                                        \tag{D.66.9}
\]
It is normal in \(I_s\): every actual operator is a
multiplier followed by \(t_i\mapsto t_i^{\chi_5(g)}\),
so conjugation sends \(\mu^c\) to
\(\mu^{c\chi_5(g)}\). Also \(A_s^N=0\), because
\(\mu\ne1\) and \(\Lambda\) is an integral domain.
The exponents need not be primitive.

By (D.66.8) at this prime, subtract one integral
coboundary so that the descended cocycle is zero on
\(N\). For \(g\in I_s\) and \(n\in N\), compare
its values at \(ng=g(g^{-1}ng)\). The two cocycle
identities give \(n c(g)=c(g)\). Thus every value
lies in \(A_s^N=0\). The cocycle is a coboundary,
which proves \({\cal G}\subset{\cal Z}\).

The finitely supported elements of \({\cal K}\) form
a countable union of closed nowhere-dense submodules:
fixing a finite support is closed, and a fresh prime
can be added in every neighborhood. Intersecting
\({\cal G}\) with the complement of this set in its
first coordinate proves the asserted infinite-ramification
refinement.

\paragraph{The section fibres and the local existence filter.}
For \(b=0\), each \((a,0)\) is the abelian image of
an actual full completed cusp section, as in D.60.
D.35 identifies its exact metabelian fibre with a torsor
under \(H^1(G_{\mathbb Q},A_{(a,0)})\), without an
additional abelian stabilizer quotient. It is a singleton
on the dense \(G_\delta\) just constructed. The
\(\mathbb Q\)-defined permutations of the three
punctures induce the corresponding continuous affine
homeomorphisms of completed packet parameters and
bijections of section fibres. Applying them gives the
same conclusion in each packet.

Finally D.42 implies that any abelian pair admitting a
metabelian lift satisfies, already at the single prime two,
\[
                        a_2b_2(a_2-b_2)=0.             \tag{D.66.10}
\]
The coordinate projection \({\cal K}^2\to\mathbb Z_5^2\)
is continuous and open. The zero set in (D.66.10) is
the union of three closed lines, with empty interior:
in any pair of five-adic balls one can choose two
nonzero unequal coordinates. Its inverse image is
therefore closed and nowhere dense. This proves the
last assertion and explains precisely where generic
singleton fibres have been established.
\end{proof}


% D.67. Root and PBC independent full proof and primary-source audits PASS.
\subsection{Generic cusp rigidity in the full free pro-five group}

\begin{statement}{D.67}{Proved}{A comeagre set of cusp abelian classes has a unique full pro-five section}
Let \(F_5=F^{(5)}\) be the actual geometric free pro-five group of
the tripod over \(\mathbb Q\), with the unit tangential arithmetic
action at zero. Give
\({\cal K}=H^1(G_{\mathbb Q},\mathbb Z_5(1))\) the intrinsic
uniform five-adic topology of (D.66.2). There is a dense
\(G_\delta\) subset \({\cal G}_{\rm full}\subset{\cal K}\)
such that, for every \(a\in{\cal G}_{\rm full}\), the fibre
\[
 H^1_{\mathrm{cts}}(G_{\mathbb Q},F_5)
       \longrightarrow H^1(G_{\mathbb Q},F_5^{\rm ab})
                       ={\cal K}^2
 \qquad\text{over }(a,0)                               \tag{D.67.1}
\]
is the singleton consisting of the completed cusp section
at zero. This section is supplied by the actual full
profinite cusp packet. The analogous assertion holds in
each of the other two cusp packets.

The set may be chosen so that every parameter is infinitely
ramified and also belongs to the generic singleton
metabelian-fibre locus of D.66. Thus both the full free
pro-five fibre and the metabelian fibre are then singletons.
Every compatible tower of nilpotent section classes
above such an abelian pair is the tower of this cusp class.
No assertion is made that individual lifts to arbitrarily
high finite nilpotent levels admit compatible choices,
or that a first obstructed finite level has been determined.
\end{statement}

\begin{proof}
Write \(P=F_5\), let \(H=P'\) be its closed commutator
subgroup, and let \(\sigma_0:G_{\mathbb Q}\to\operatorname{Aut}(P)\)
be the actual unit tangential action. In particular
\(\sigma_0(g)(x)=x^{\chi_5(g)}\). Use the fixed continuous
Kummer cocycles of D.66 for all parameters. The cusp
section with parameter \(a\) gives the genuine action
\[
 \sigma_a(g)=\operatorname{Int}(x^{a(g)})\sigma_0(g)
                  \quad\text{on }P\text{ and on }H.     \tag{D.67.2}
\]
The cocycle identity for \(a\) proves that (D.67.2) is
an action. A section class with abelian class \((a,0)\)
can be conjugated so that its abelian cocycle agrees
exactly with that of this cusp section. Its difference
is then an actual nonabelian cocycle
\[
                           c:G_{\mathbb Q}\longrightarrow H
                                                        \tag{D.67.3}
\]
for \(\sigma_a\). Conversely every such cocycle gives
a section in that fibre. A conjugator between two
such exact representatives has abelian image fixed
by \(G_{\mathbb Q}\), which is zero in
\(\mathbb Z_5(1)^2\). Hence the fibre is precisely
the nonabelian set \(H^1(G_{\mathbb Q},H;\sigma_a)\).
We shall show that this pointed set is trivial on a
dense \(G_\delta\) of parameters.

\paragraph{The exact local consequence of good reduction.}
Fix a rational prime \(q\ne5\). The tripod with its
three marked boundary points has good reduction over
\(\mathbb Z_q\), and the honest unit tangential action
on \(P\) is unramified. Let \(\Pi_q=P\rtimes G_{\mathbb Q_q}\),
and let \(\Pi_q^{\rm sp}\) be its geometrically pro-five
good-reduction specialization group over
\(G_{\mathbb F_q}\). The comparison diagram is
\[
\begin{array}{ccccccccc}
1&\longrightarrow&P&\longrightarrow&\Pi_q&\longrightarrow&G_{\mathbb Q_q}&\longrightarrow&1\\
 &&\Vert&&\downarrow&&\downarrow\\
1&\longrightarrow&P&\longrightarrow&\Pi_q^{\rm sp}&\longrightarrow&G_{\mathbb F_q}&\longrightarrow&1.
\end{array}                                               \tag{D.67.4}
\]
It is cartesian. The geometric isomorphism and this
description are exactly the prime-set \(\Sigma=\{5\}\)
case of \cite[Definition A.3 and Remark A.3.1, p.~38]{DE_Hoshi2012}.
The residue characteristic is \(q\), and is excluded
from \(\Sigma\), as that statement requires.

For every actual local section, its inertia image
in \(\Pi_q^{\rm sp}\) is either trivial or a
nontrivial subgroup of a conjugate cusp inertia
\(\mathbb Z_5(1)\). This is
\cite[Proposition A.6, pp.~41--42]{DE_Hoshi2012}.
We spell out the additional implication needed here,
rather than interpreting that proposition as an
unqualified local section conjecture.

In a free pro-five group the centralizer of a
nonzero power of the primitive generator \(x\)
is \(\langle x\rangle\). One can check this
directly by the Magnus embedding in
\(\mathbb Z_5\langle\!\langle X,Y\rangle\!\rangle\):
for nonzero \(b\in\mathbb Z_5\), the series
\((1+X)^b-1=bX+\cdots\) has a compositional inverse
over \(\mathbb Q_5\). Commuting with it therefore
implies commuting with \(X\). Comparing coefficients
of words containing \(Y\) shows that the centralizer
of \(X\) in the noncommutative power-series algebra
is \(\mathbb Q_5[[X]]\): repeatedly move initial
\(X\)'s to the end of a word containing \(Y\);
once it starts with \(Y\), the commutation equation
forces its coefficient to be zero. Evaluating at \(Y=0\)
then identifies a group element in this subalgebra
with the power of \(x\) determined by its first
abelian coordinate. The same argument applies to
the other two primitive cusp generators, after
a free-basis change. A geometric normalizer of a
nonzero subgroup of \(\langle x\rangle\) is its
centralizer: abelianization forces its exponent
automorphism to be the identity.

Consequently, if a local section has nonzero
specialized inertia in a cusp, its entire specialized
image lies in that cusp's decomposition group.
Indeed the image normalizes the inertia subgroup;
a unit tangential arithmetic lift normalizes the
whole cusp group, and the preceding normalizer
calculation forces the remaining geometric factor
to belong to it. Pullback through (D.67.4) puts the
actual local section in the actual pro-five cusp
packet. This conclusion is compatible with an
ordinary local point close to the cusp having the
same prime-to-\(q\) section image; no point/cusp
separation at the wrong coefficient prime is used.

Now apply this to a section whose exact abelian
cocycle is \((a,0)\). If \(a_q\ne0\), its
inertia image is nonzero and its abelian direction
forces the cusp to be zero: the three cusp
directions are the distinct axes and diagonal.
The parameter in that local packet has Kummer
class \(a|_{G_{\mathbb Q_q}}\), since the primitive
first-coordinate inclusion of its inertia induces
an injection on \(H^1\). Thus this local section
is conjugate to the reference cusp section. The
conjugator can be taken in \(H\). To see this
without ignoring a stabilizer, first combine the
geometric conjugator into the packet with the
cusp-inertia conjugator adjusting its Kummer
cocycle. The abelian image of the resulting
conjugator is fixed by \(G_{\mathbb Q_q}\), since
the two actual abelian cocycles agree. Frobenius
acts there by \(q\), and \(q-1\) is a nonzero
element of \(\mathbb Z_5\), so that fixed image
is zero. Therefore every cocycle (D.67.3) satisfies
\[
 a_q\ne0\quad\Longrightarrow\quad
 c|_{G_{\mathbb Q_q}}\text{ is an }H\text{-coboundary}.
                                                        \tag{D.67.5}
\]
If instead \(a_q=0\), Proposition A.6 forces
the actual section inertia to be zero in
\(\Pi_q^{\rm sp}\). Otherwise its nonzero
cuspidal inertia would have a nonzero abelian
image, contrary to the abelian inertia pair
\((0,0)\). The reference cusp section is also
unramified there, and hence
\[
                   a_q=0\quad\Longrightarrow\quad
                           c(I_q)=1.                   \tag{D.67.6}
\]

\paragraph{Uniform capture in actual finite cocycle sets.}
Use the descending iterated Frattini subgroups
\(P_0=P\), \(P_{r+1}=\Phi(P_r)\). They are
characteristic open subgroups with trivial
intersection. Indeed open subgroups of the
finitely generated free pro-five group are
finitely generated, and their Frattini quotients
are finite; every finite five-group has a
terminating Frattini series. Put
\[
       Q_r=P/P_r,\qquad H_r=\operatorname{im}(H\to Q_r).
                                                        \tag{D.67.7}
\]
The exponent of \(Q_r\) divides \(5^r\), since
\(g^{5^r}\in P_r\) for every \(g\in P\).
These quotients induce the actual topology of
\(H\), and are stable under (D.67.2).

Let \(\rho_{a,r}\) be the action on \(H_r\), and set
\[
 S_r(a)=\{5,\infty\}\cup\{q:a_q\notin5^r\mathbb Z_5\},
 \qquad E_r(a)=Z^1(G_{\mathbb Q,S_r(a)},H_r;\rho_{a,r}).
                                                        \tag{D.67.8}
\]
This is a finite set of actual nonabelian
cocycles. For example each cocycle defines a
homomorphism to the finite group
\(H_r\rtimes\operatorname{im}\rho_{a,r}\), with
fixed second coordinate; finite extensions of
bounded degree unramified outside a fixed
finite set form a finite set, by the same
Hermite--Minkowski argument used in D.47.

Both \(S_r(a)\) and the finite action are
constant on every coset of \(5^r{\cal K}\).
For the action this follows from (D.67.2) and
the exponent bound for \(Q_r\). Moreover every
integral cocycle (D.67.3) reduces to an element
of \(E_r(a)\). At a prime outside \(S_r(a)\)
with nonzero \(a_q\), (D.67.5) expresses its
inertia as a conjugation coboundary for
\(\operatorname{Int}(x^{a(\tau)})\). That
automorphism is the identity on \(Q_r\),
because \(a(\tau)\in5^r\mathbb Z_5\).
Thus the reduced cocycle inertia is trivial.
At a prime with \(a_q=0\), use (D.67.6).
The finite action itself has the same support
bound. This proves exact factorization of the
reduced affine representation through the
group in (D.67.8); it is stronger than an
unramified-cohomology-class assertion.

\paragraph{Dense open conditions eliminating nontrivial kernel restrictions.}
Call \(z\in E_r(a)\) exceptional if its
restriction to \(\ker\rho_{a,r}\) is nontrivial.
Define \({\cal V}_r\subset{\cal K}\) by the
following condition for each exceptional \(z\):
there exists a prime \(q\notin S_r(a)\) with
\[
 \rho_{a,r}(\operatorname{Frob}_q)=1,\qquad
 z(\operatorname{Frob}_q)\ne1,\qquad a_q\ne0.
                                                        \tag{D.67.9}
\]
This is open. On a \(5^r{\cal K}\) coset
all finite data are fixed, there are finitely
many exceptional cocycles, and nonzero
coordinates at a finite choice of witness
primes are open conditions.

It is dense by the actual finite affine-image
argument. Given \(a_0+5^N{\cal K}\), choose
\(M\ge\max\{N,r\}\). For every exceptional
\(z\), an element \(\delta\) in the action
kernel with \(z(\delta)\ne1\) gives an actual
element \((z(\delta),1)\) of
\(H_r\rtimes\operatorname{im}\rho_{a_0,r}\).
Chebotarev supplies infinitely many primes
in this conjugacy class. Both tests in
(D.67.9) are invariant under conjugacy, also
for this nonabelian translation group.
Choose distinct witness primes, outside
\(S_r(a_0)\) and outside the finite set where
\(a_{0,q}\notin5^{M+1}\mathbb Z_5\), and add
\(5^M\kappa(q)\) at each. The finite data
are unchanged, while each new witness
coordinate has valuation exactly \(M\).
Thus all tests hold in the prescribed
neighborhood. This uses precisely
\cite[Chapter V, Theorem 3.23]{PBC_MilneCFT}
on finite Galois quotients that have been
shown to contain the desired element.

Let \({\cal V}\) be the open dense set of
parameters with \(a_q\ne0\) at some
\(q\ne5\). Completeness of the parameter
space, proved in D.66, shows that
\[
                   {\cal G}_0={\cal V}\cap\bigcap_{r\ge1}{\cal V}_r
                                                        \tag{D.67.10}
\]
is a dense \(G_\delta\).

\paragraph{Descent to the actual compact action image.}
Fix \(a\in{\cal G}_0\) and a cocycle
\(c\) as in (D.67.3). Its reduction
\(c_r\in E_r(a)\) cannot be exceptional.
At its witness prime, (D.67.5) would give
an \(H_r\)-coboundary, whose value at a
Frobenius with identity action is one,
contradicting (D.67.9). Thus every \(c_r\)
is trivial on \(\ker\rho_{a,r}\).

It follows that \(c\) is trivial on the
kernel of the actual action on \(H\), and
descends continuously to the compact image
\[
          I_a=\operatorname{im}\left(
             G_{\mathbb Q}\longrightarrow
                         \prod_r\operatorname{Aut}(H_r)\right).
                                                        \tag{D.67.11}
\]
Here the compatible finite actions define
the actual action on \(H\); no enlargement
to an abstract automorphism group is made.
The quotient map from the compact Galois
group proves continuity of the descended
cocycle.

Choose \(q\ne5\) with \(a_q\ne0\).
The image of its inertia in \(I_a\) is
\[
 N=\{\operatorname{Int}(x^{a_q t})|_H:t\in\mathbb Z_5\}.
                                                        \tag{D.67.12}
\]
It is normal in \(I_a\), since every
\(\sigma_a(g)\) sends \(x\) to
\(x^{\chi_5(g)}\). Its fixed subgroup
in \(H\) is trivial:
\[
              H^N=H\cap C_P(x^{a_q})
                    =H\cap\langle x\rangle=1.          \tag{D.67.13}
\]
The first exponent need not be a unit.
By (D.67.5) the descended cocycle is an
\(H\)-coboundary on \(N\). Conjugate it
by that single element of \(H\) so that
it becomes identically one on \(N\).
For \(g\in I_a\), \(n\in N\), compare
the two cocycle formulas at
\(ng=g(g^{-1}ng)\). They imply
\(n(c(g))=c(g)\). Equation (D.67.13)
therefore makes every value one. This
proves that (D.67.3) is a global
\(H\)-coboundary, and hence proves
the singleton assertion in (D.67.1).

Intersect \({\cal G}_0\) with the dense
\(G_\delta\) of infinitely supported
parameters and with the fixed-\(b=0\)
dense \(G_\delta\) from D.66. This is
the asserted \({\cal G}_{\rm full}\).
The permutations of the punctures are
defined over \(\mathbb Q\), act by
homeomorphisms on completed packet
parameters, and transport the actual
section fibres. They give the assertion
for the other two packets.

Finally \(P=\varprojlim_n T_n\): every
finite five-group is nilpotent, so the
lower-central quotients detect every
finite quotient of \(P\). A compatible
sequence of section classes in \(T_n\)
has compatible actual representatives.
Inductively conjugate a representative
at stage \(n+1\) by a lift of the
conjugator identifying its projection
with the chosen representative at stage
\(n\). Their inverse limit is a
continuous section. Conversely two
sections conjugate at every stage are
conjugate in \(P\): their nonempty
closed conjugator sets in the compact
group \(P\) are nested. Thus compatible
nilpotent towers in the indicated fibre
are exactly the singleton full pro-five
fibre. This argument deliberately starts
from compatible classes and says nothing
about unrelated choices of finite-stage
lifts.
\end{proof}


% D.68. Root and PBC independent full proof and primary-source audits PASS.
\subsection{Compactness after fixing the actual abelian cocycle}

\begin{statement}{D.68}{Proved}{Fixed-abelian nilpotent fibres are compact and detect finite failure}
Retain the actual tangential arithmetic action on
\(T_n=F_5/\gamma_{n+1}F_5\). Fix an actual continuous
abelian cocycle
\(\alpha:G_{\mathbb Q}\to T_1=\mathbb Z_5(1)^2\),
with unrestricted ramification. Give continuous cochain
spaces their uniform topology. Then, for every \(n\),
\[
 \mathscr Z_n(\alpha)=
 \{z\in Z^1_{\rm cts}(G_{\mathbb Q},T_n):
                     z\bmod\gamma_2=\alpha\}
                                                        \tag{D.68.1}
\]
is compact. It is the image of a closed subset of a
finite power of \(\mathbb Z_5\) under a continuous
family of actual cocycles whose Hall coordinates
are rational polynomials in the parameters with
coefficients in fixed finite spaces of continuous
cochains. Consequently the section-class fibre
\[
 \mathscr S_n([\alpha])=
 \operatorname{Sec}(T_n)\mathbin{\times}_{\operatorname{Sec}(T_1)}
                                      \{[\alpha]\}
                                                        \tag{D.68.2}
\]
is compact and Hausdorff, as is the analogous full
free pro-five fibre. Every transition map between
these compact fibres is continuous and has closed
image.

If a fixed class \(u\in\mathscr S_m([\alpha])\)
lifts to \(T_n\) for every \(n\ge m\), then these
individual lifts can be chosen compatibly, and
\(u\) lifts to the actual full free pro-five group.
The same assertions hold for the actual nilpotent
quotients \(M_n=M/\gamma_{n+1}M\) of
\(M=F_5/F_5''\), with their limit \(M\).
In particular, individual nonemptiness at every
nilpotent level suffices for nonemptiness of the
corresponding actual limit fibre.

For \(a\) in the generic cusp set of D.67, every
noncuspidal class at a fixed finite nilpotent stage
above \((a,0)\) fails to lift at some higher finite
stage. More precisely, for every neighborhood \(V\)
of the cusp class in \(\mathscr S_m((a,0))\),
there is an \(N\ge m\) such that the image of
\(\mathscr S_N((a,0))\) in
\(\mathscr S_m((a,0))\) is contained in \(V\).
At class three the fibre is the compact torsor
\(\mathbb Z_5^2\), with cusp as origin; for every
\(h\ge1\), one finite stage excludes every
parameter outside \(5^h\mathbb Z_5^2\).
For the metabelian quotients this last assertion
already holds on the generic set of D.66.
These finite stages are not evaluated here; an
explicit first stage and its obstruction class
remain open.
\end{statement}

\begin{proof}
The issue is compactness inside the actual continuous
cochain complex. In particular, the global group
\(H^2\) may be nonseparated, as D.1 and D.62 prove.
We shall use a finite coefficient module inside
that complex and shall never replace the subgroup
of coboundaries by its closure.

\paragraph{Integral coordinates in every finite class.}
Let \(F_{\rm disc}\) be the discrete free group on
\(x,y\). The basic commutator theorem identifies
\(F_{\rm disc}/\gamma_{n+1}F_{\rm disc}\) with the
ordered integral powers of its basic commutators
of weights at most \(n\). Its successive
lower-central factors are free abelian on the
commutators of the corresponding weight; see
\cite[Theorem 5.6 and its corollary]{DE_HallEdmonton}.
In these coordinates multiplication and powers
are integer-valued rational polynomials
\cite[Theorem 6.5]{DE_HallEdmonton}. The same
source immediately after that theorem, case (iii),
identifies the coordinate construction over
\(\mathbb Z_5\) with the pro-five completion.

For clarity, this completion is exactly \(T_n\).
A homomorphism from the discrete free group to
a finite five-group of class at most \(n\)
factors through the displayed discrete quotient,
and extends uniquely from the free pro-five
group. Conversely every finite quotient of
\(T_n\) has class at most \(n\). Thus the two
pro-five completions have exactly the same
finite quotients with the specified generators.
It follows that the Hall coordinate map is a
homeomorphism from a finite power of
\(\mathbb Z_5\) onto \(T_n\), and that
\[
 L_n=\ker(T_n\longrightarrow T_{n-1})
       \simeq\mathbb Z_5(n)^{d_n},\qquad n\ge2,
                                                        \tag{D.68.3}
\]
where \(d_n\) is the number of basic commutators
of weight \(n\). The arithmetic action is scalar
\(\chi_5^n\) on this factor, since it is scalar
\(\chi_5\) on degree one and the graded Lie ring
is generated in degree one. Inner twisting by
any lower section does not change the action
on the central kernel \(L_n\).

Appending zero top Hall coordinates to the
lower coordinates gives a continuous polynomial set-section
\(j_n:T_{n-1}\to T_n\). No assertion of integral
Baker--Campbell--Hausdorff coordinates in class
at least five is needed. The rational denominators
of the Hall polynomials are bounded at each
fixed class, and their values on the actual
coordinate lattice remain integral.

The actual arithmetic action is polynomial in
the Hall coordinates of its input, with coefficients
that are fixed continuous functions of the
arithmetic element. Indeed an automorphism is
determined by its images of \(x,y\). Each basic
commutator maps to the same fixed group word
in these two images, and its arbitrary
\(\mathbb Z_5\)-power is evaluated by the power
polynomials just cited. Substituting these finitely
many coordinates in the multiplication polynomials
proves this assertion at every fixed class.
This describes the genuine action, including
all its higher arithmetic coefficients.

\paragraph{Finite generation of actual central cocycles.}
For every integer \(j\ge2\), every continuous
\(\mathbb Z_5(j)\)-valued one-cocycle kills
inertia at every \(q\ne5\). Its restriction to
inertia is a homomorphism because the coefficient
is unramified. Wild inertia and the prime-to-five
tame factors are killed. If \(b\) is its value
on a five-primary tame generator, the tame
Frobenius relation and the cocycle identity give
\[
                         q b=q^j b.
                                                        \tag{D.68.4}
\]
Since \(q^j-q\ne0\) in \(\mathbb Z_5\), we
have \(b=0\). The coefficient action itself is
unramified there. Thus the actual affine cocycle
factors through \(G_{\mathbb Q,\{5\}}\).

Finite-coefficient cohomology of this last
group is finite
\cite[I, Corollary 4.15]{DE_MilneADT}.
The finite systems of invariant groups are
Mittag--Leffler, so continuous comparison
\cite[Proposition 1.6 and Theorem 2.2]{DE_Jannsen1988}
identifies integral \(H^1\) with the inverse
limit of its finite-coefficient groups. Give it
this inverse-limit topology for the moment; it
is compact in that topology, before any assertion
about the uniform cochain quotient. The coefficient sequence
for multiplication by five injects
\(H^1/5H^1\) into the finite group
\(H^1(G_{\mathbb Q,\{5\}},\mathbb F_5(j))\).
Compact Nakayama then makes integral \(H^1\)
a finitely generated \(\mathbb Z_5\)-module.
This is the finiteness portion of D.1, now
spelled out for every \(j\ge2\); no finite
ramification is imposed on the original cocycle.

Choose finitely many actual continuous cocycles
representing generators of this \(H^1\), and
add the coboundary of a basis of the coefficient
lattice. They generate the module of all actual
one-cocycles. A finite integral linear combination
of these fixed continuous maps is continuous,
and the map from the finite parameter power of
\(\mathbb Z_5\) is continuous for the uniform
cochain topology. Hence
\[
 Z^1_{\rm cts}(G_{\mathbb Q},L_n)
       \text{ is a finitely generated compact }
                          \mathbb Z_5\text{-module}.
                                                        \tag{D.68.5}
\]
Torsion in \(H^1\) causes no problem: its
representative cocycles, together with the
coboundaries, still generate all actual cocycles.
The latter module is a submodule of a torsion-free
cochain module.

\paragraph{The finite-span coboundary lemma.}
Let \(L\) be a finite free \(\mathbb Z_5\)-module
with continuous arithmetic action. Fix finitely
many continuous two-cochains
\(\zeta_i\in C^2(G_{\mathbb Q},L\otimes\mathbb Q_5)\),
for \(1\le i\le r\), and define the actual
coefficient submodule
\[
 B=\left\{v\in\mathbb Z_5^r:
       \sum_i v_i\zeta_i\in
                 d C^1(G_{\mathbb Q},L)\right\}.
                                                        \tag{D.68.6}
\]
It is a \(\mathbb Z_5\)-submodule, even though
the coboundaries in the entire cochain space
need not be closed. Since \(\mathbb Z_5\) is
Noetherian, \(B\) is finite free. It is closed:
the map from \(\mathbb Z_5^e\) given by any
module basis of \(B\) has compact image in
\(\mathbb Z_5^r\).

Choose a basis \(v_1,\ldots,v_e\) of \(B\)
and integral continuous primitives \(b_k\) with
\(d b_k=\sum_i(v_k)_i\zeta_i\). The coordinate
functionals of this basis extend to linear
functionals over \(\mathbb Q_5\) on
\(\mathbb Q_5^r\): choose a left inverse over
\(\mathbb Q_5\) of the matrix with columns
\(v_k\). Consequently, if \(v=v(t)\) is a
polynomial coefficient vector taking values
in \(B\), the primitive
\[
                b(t)=\sum_k (R v(t))_k b_k
                                                        \tag{D.68.7}
\]
is integral and continuous, and is again a
rational-polynomial family in a fixed finite
space of continuous one-cochains. The coordinates
\((R v(t))_k\) are integral on the allowed
parameter set. Their denominators as polynomial
functions are fixed and bounded. This is an
actual choice of primitives, not a lifting
through a separated cohomology group.

\paragraph{Induction with all exact cocycles.}
We prove the stronger finite-parameter assertion
in the statement by induction on \(n\). A
polynomial family means a family on a closed
set \(X\subset\mathbb Z_5^r\), for a finite
\(r\), such that every Hall coordinate has
an expression
\[
                \sum_i P_i(t)\varphi_i(g),
 \quad P_i\in\mathbb Q_5[t_1,\ldots,t_r],\quad
       \varphi_i\in C(G_{\mathbb Q},\mathbb Q_5).
                                                        \tag{D.68.8}
\]
Only finitely many polynomials and fixed functions
occur. The coordinate values are integral on
\(X\). Since both parameter space and Galois
group are compact, joint continuity gives
continuity into the uniform cochain space.
At \(n=1\), the single family member is the
fixed actual abelian cocycle \(\alpha\).

Suppose that a polynomial family \(f_t\),
parametrized by a closed \(X\), covers all
of \(\mathscr Z_{n-1}(\alpha)\). Apply \(j_n\)
to obtain actual set-theoretic lifts \(\widehat f_t\).
Their multiplication defects are integral
continuous central two-cocycles in \(L_n\):
\[
 \eta_t(g,h)=\widehat f_t(g)\,
        \sigma_0(g)(\widehat f_t(h))\,
                   \widehat f_t(gh)^{-1}.
                                                        \tag{D.68.9}
\]
The coordinate facts proved above show that
this is a polynomial family in a fixed finite
space of two-cochains. To make the finite-span
claim explicit, multiplication and inversion
use finitely many Hall polynomials; arithmetic
substitution uses the finitely many continuous
coordinates of \(\sigma_0(g)(x)\) and
\(\sigma_0(g)(y)\); and substitution of
\(f_t(g),f_t(h),f_t(gh)\) uses finitely many
functions in (D.68.8). Expanding their finite
products gives finitely many monomials in
\(t\), whose coefficients are fixed continuous
functions of \((g,h)\). This proves the
finite-span assertion uniformly on all of \(X\),
including after every previous choice of
primitives.

Absorb bounded polynomial denominators into
the coefficient cochains. We may therefore
write
\[
                 \eta_t=\sum_i v_i(t)\zeta_i,
       \qquad v(t)\in\mathbb Z_5[t_1,\ldots,t_r]^k,
                                                        \tag{D.68.10}
\]
where the fixed \(\zeta_i\) may be rational
rather than integral cochains. Form \(B\) as
in (D.68.6), with the actual integral coefficient
lattice \(L_n\). Exactly those parameters in
\[
                         X'=X\cap v^{-1}(B)
                                                        \tag{D.68.11}
\]
admit an integral continuous cocycle lift.
This set is closed and compact. Choose the
primitive \(b(t)\) of (D.68.7); multiplying
\(\widehat f_t\) by the central cochain
\(-b(t)\) cancels its multiplication defect
and produces a genuine lifted cocycle.

Choose finite generators \(z_1,\ldots,z_e\)
of the actual central cocycle module (D.68.5).
Multiplying the chosen lift by
\(\sum_j\lambda_j z_j\), with
\(\lambda\in\mathbb Z_5^e\), gives a
polynomial family on the compact closed set
\(X'\times\mathbb Z_5^e\). It covers every
exact lift: the difference of any two lifts
with identical projection is an actual central
one-cocycle. By the inductive surjectivity onto
all lower exact cocycles, the new family covers
all of \(\mathscr Z_n(\alpha)\).
The chosen primitives are rational-polynomial
in the old parameters, and the new cocycles
are linear in the new parameters. Thus the
finite-span property (D.68.8) is preserved,
not merely continuity or a statement about
cohomology classes. This finishes the induction
and proves compactness in (D.68.1).

\paragraph{Conjugacy, limits, and closed images.}
The uniform cocycle space is Hausdorff, and
conjugation by the compact group \(T_n\) has
closed orbits and a Hausdorff quotient. Every
class with abelian class \([\alpha]\) admits
an exact representative with abelian cocycle
\(\alpha\): lift the required abelian
conjugator from \(T_1\) to \(T_n\).
A conjugator between two such exact
representatives has abelian image fixed by
\(G_{\mathbb Q}\), hence zero. Thus
\(\mathscr S_n([\alpha])\) is also the
quotient of (D.68.1) by the compact group
\(\gamma_2T_n\). It is compact Hausdorff.
Projection of cochains is continuous, so the
transition maps have closed image.

Now fix an actual representative \(z_m\) of
\(u\). At each \(n\ge m\), its exact lift
space is a closed subset of
\(\mathscr Z_n(\alpha)\), hence compact.
If \(u\) has a class lift at that level,
lifting a conjugator from \(T_m\) makes
this exact lift space nonempty. These compact
spaces form an inverse system. Every finite
compatibility condition can be met by a
member at the highest of the finitely many
levels. Compactness therefore gives a
compatible sequence of actual cocycles.
Since \(F_5=\varprojlim T_n\), their inverse
limit is an actual continuous \(F_5\)-cocycle.
The uniform cochain topology into \(F_5\) is
exactly the inverse-limit topology on these
exact fibres: every finite quotient of \(F_5\)
is nilpotent and therefore factors through
some \(T_n\), so every uniform neighborhood
is detected at a finite stage.
The same compact construction with \(m=1\)
proves compactness of the full exact fibre;
quotienting by the compact commutator subgroup
proves compactness of the full class fibre.
As in D.67, nested nonempty closed conjugator
sets give injectivity on the inverse limit
of section classes.

The metabelian assertion has exactly the same
hypotheses. To check them, the discrete free
metabelian Alexander module is
\(\mathbb Z[X^{\pm1},Y^{\pm1}]C\), and
its lower-central filtration is given by
powers of \((X-1,Y-1)\). Its nilpotent
quotients are therefore finitely generated
torsion-free nilpotent groups, with central
grade \(n\) free of rank \(n-1\).
Their pro-five completions are the actual
\(M_n\), by the same finite-quotient universal
property; equivalently use the completed
Alexander description of D.35. The central
grade has arithmetic action \(\chi_5^n\).
Hall's polynomial theorem applies to these
nilpotent groups too. The preceding proof
therefore supplies compact exact fibres and
compatible representatives. Here the limit is
indeed \(M\): the intersection of all augmentation
ideal powers in \(\mathbb Z_5[[u,v]]\) is zero,
and compactness then makes the map to the
inverse limit of the quotients surjective.
This also removes a possible compatibility
ambiguity in the individual-stage version
of the metabelian branch criterion of D.41.

\paragraph{The generic finite-stage obstruction.}
For the parameters of D.67, the intersection
of the descending compact images
\[
 \operatorname{im}\bigl(\mathscr S_N((a,0))
                      \longrightarrow\mathscr S_m((a,0))\bigr),
                         \qquad N\ge m,
                                                        \tag{D.68.12}
\]
is exactly the cusp class. Indeed a class
in their intersection lifts at every stage,
hence to \(F_5\), and D.67 identifies that
full fibre. If an open neighborhood \(V\)
of the cusp failed to contain every image
for large \(N\), the compact sets obtained
by intersecting (D.68.12) with the complement
of \(V\) would all be nonempty and nested.
Their intersection would contradict the
singleton intersection just proved.

For \(m=3\), the central kernel over \(T_2\)
is \(\mathbb Z_5(3)^2\). The exact fibre
above its cusp lift is a torsor under
\(H^1(G_{\mathbb Q},\mathbb Z_5(3))^2
=\mathbb Z_5^2\), by D.1 and the central
lifting calculation. There is no stabilizer
quotient: a fixed element of \(T_2\) for a
section action has a first nonzero degree
with positive cyclotomic weight, which is
impossible. Also the class-two fibre above
\((a,0)\) is a singleton, since
\(H^1(G_{\mathbb Q},\mathbb Z_5(2))=0\).
Explicitly, twist D.1's residual unit module
\(\mathbf1\oplus\omega\oplus\omega^2\) by
\(\omega\). The resulting characters
\(\omega,\omega^2,\omega^3\) have no invariant
part, so \(H^1(G_{\mathbb Q,\{5\}},\mathbb F_5(2))=0\).
Compact Nakayama gives the integral vanishing,
and (D.68.4) gives its unrestricted form.
The continuous torsor parametrization is a
bijection from a compact space to the
Hausdorff fibre, hence a homeomorphism.
Taking \(V=5^h\mathbb Z_5^2\) proves the
stated uniform finite failure. The same
argument with the singleton metabelian
fibre of D.66 proves its metabelian version.
No estimate for \(N\), or value of the
first obstruction, follows from this
compactness proof.
\end{proof}


% D.69. Root and PBC independent full proof and dependency audits PASS.
\subsection{Generic rigidity in the full profinite cusp packet}

\begin{statement}{D.69}{Proved}{Generic completed cusp classes have singleton full profinite fibres}
Let \(F=\pi_1((\mathbb P^1\setminus\{0,1,\infty\})_{\overline{\mathbb Q}})\)
be the actual free profinite geometric group, with the
unit tangential arithmetic action at zero. Put
\[
 \widehat{\mathcal K}=H^1(G_{\mathbb Q},\widehat{\mathbb Z}(1))
          =\prod_\ell\mathcal K_\ell,
 \qquad \mathcal K_\ell=H^1(G_{\mathbb Q},\mathbb Z_\ell(1)).
                                                        \tag{D.69.1}
\]
For odd \(\ell\), give \(\mathcal K_\ell\) its
uniform \(\ell\)-adic \(c_0\) topology; at two use
\[
 \mathcal K_2=\mathbb Z/2\ \oplus\
                   c_0(\{q\text{ rational prime}\},\mathbb Z_2),
                                                        \tag{D.69.2}
\]
with the sign factor discrete and the other factor
uniform two-adic. Give (D.69.1) the product topology.
There is a dense \(G_\delta\) subset
\(\widehat{\mathcal G}\subset\widehat{\mathcal K}\)
such that, for each \(a\in\widehat{\mathcal G}\),
the actual full profinite section fibre
\[
 H^1_{\rm cts}(G_{\mathbb Q},F)
       \longrightarrow H^1(G_{\mathbb Q},F^{\rm ab})
                      =\widehat{\mathcal K}^{2}
       \qquad\text{over }(a,0)
                                                        \tag{D.69.3}
\]
is the singleton consisting of its completed cusp
section at zero. The analogous assertion holds in
each of the other two rational cusp packets.

The parameters may be chosen infinitely ramified
in every prime component. At every prime \(\ell\),
the corresponding full geometric pro-\(\ell\)
section fibre is itself a singleton. At five the
set may simultaneously satisfy D.66--D.68, including
the singleton metabelian fibre and the non-effective
finite-stage obstruction away from its cusp.
No classification of all abelian pairs or of all
full profinite sections is asserted.
\end{statement}

\begin{proof}
We first check the coefficient topology and the
prime two; neither is inferred from the five-adic
statement by an implicit convention.

\paragraph{Power completion and the sign cocycle.}
The group \(\mathbb Q^\times\) is the product of
its sign group with the free abelian group on
the positive rational primes. Finite Kummer theory,
followed by the inverse limit of the finite
coefficient groups, gives
\[
 \mathcal K_\ell\simeq
 \begin{cases}
 c_0(\{q\},\mathbb Z_\ell),&\ell\ne2,\\
 \mathbb Z/2\oplus c_0(\{q\},\mathbb Z_2),&\ell=2.
 \end{cases}                                           \tag{D.69.4}
\]
Here the finite Kummer quotients at each exponent
are discrete direct sums over the rational primes;
their inverse limit is \(c_0\), not the unrestricted
product over those primes. The invariant groups
of the finite root-of-unity coefficients are
finite and Mittag--Leffler, so that continuous
\(H^1\) agrees with this inverse limit.
Also continuous cochains with values in
\(\widehat{\mathbb Z}(1)=\prod_\ell\mathbb Z_\ell(1)\)
are exactly the products of the component cochains.
Products are exact for these abelian complexes,
which proves (D.69.1). Equivalently it is the
power completion
\(\varprojlim_m\mathbb Q^\times/(\mathbb Q^\times)^m\)
with its inverse-limit topology, by the Chinese
remainder theorem. These Kummer groups need not
be compact: the finite-exponent quotients are
infinite discrete groups.

Fix compatible roots of every positive prime,
as in D.66, at each coefficient prime. At two
fix compatible primitive roots of unity. A
continuous representative for the sign class is
\[
            \kappa(-1)(g)=\frac{\chi_2(g)-1}{2}
                          \quad\text{in }\mathbb Z_2(1).
                                                        \tag{D.69.5}
\]
Indeed take \(\zeta_{2^{n+1}}\) as a compatible
\(2^n\)-th root of \(-1\). Its Kummer exponent
relative to \(\zeta_{2^n}\) is exactly (D.69.5)
modulo \(2^n\). Twice this cocycle is the
coboundary \(\chi_2-1\). Thus (D.69.5) represents
the nonzero class of order two; no linear section
of that order-two cohomology group into cocycles
is claimed. On each fixed-sign clopen component
we use the representative
\[
             a(g)=\varepsilon\kappa(-1)(g)
                         +\sum_q a_q\kappa(q)(g),
                    \qquad \varepsilon\in\{0,1\}.
                                                        \tag{D.69.6}
\]
The sum is uniformly convergent. Deep changes in
the prime coefficients leave the chosen sign
representative unchanged. For odd \(\ell\),
use the same formula without its sign term.

Every \(\mathcal K_\ell\) is a complete
separable metric group: the uniform \(c_0\)
argument of D.66 applies for every prime, and
at two there are just two clopen copies of it.
Their countable product is complete and
separable in a usual bounded product metric.
The Baire argument below therefore applies in
the precise topology specified in the statement.

\paragraph{The nonabelian construction at an arbitrary prime.}
Fix \(\ell\), let \(P_\ell=F^{(\ell)}\),
and put \(H_\ell=P_\ell'\). The completed
cusp parameter gives the honest action
\[
       \sigma_a(g)=\operatorname{Int}(x^{a(g)})\sigma_0(g),
       \qquad\sigma_0(g)(x)=x^{\chi_\ell(g)}.
                                                        \tag{D.69.7}
\]
As in D.67, the exact fibre above \((a,0)\)
is \(H^1(G_{\mathbb Q},H_\ell;\sigma_a)\).
The conjugator comparison uses
\(H^0(G_{\mathbb Q},\mathbb Z_\ell(1)^2)=0\),
which holds at two as well.

At every rational prime \(q\ne\ell\), the
tripod's marked smooth model and unit tangent
have good reduction. The specialization square
of \cite[Definition A.3 and Remark A.3.1]{DE_Hoshi2012}
is cartesian for \(\Sigma=\{\ell\}\), since the
residue characteristic is excluded from that
set. Proposition A.6 of the same primary source
then gives the specialized inertia alternative
used in D.67. Its hypotheses allow \(\ell=2\)
and every \(q\ne\ell\).

The Magnus proof of
\(C_{P_\ell}(x^b)=\langle x\rangle\) for
\(b\ne0\) works over \(\mathbb Q_\ell\) for
every prime. Abelianization turns the normalizer
of any nonzero subgroup of this cusp into that
centralizer. Thus the full local section is in
a cusp packet whenever its abelian inertia is
nonzero. For the pair \((a,0)\), that cusp is
zero. Equality of the actual abelian cocycles
puts the local conjugator in \(H_\ell\): the
local fixed abelian image is zero because
Frobenius acts by \(q\), with \(q-1\ne0\)
in the torsion-free lattice \(\mathbb Z_\ell^2\).
Consequently every integral nonabelian cocycle
\(c:G_{\mathbb Q}\to H_\ell\) satisfies
\[
 \begin{cases}
 c|_{G_{\mathbb Q_q}}\text{ is an integral }H_\ell
                         \text{-coboundary},& a_q\ne0,\\
 c(I_q)=1,&a_q=0.
 \end{cases}                                           \tag{D.69.8}
\]
The sign term of (D.69.6) is unramified at odd
\(q\), as are all local unit Kummer classes
at coefficient prime different from the residue
characteristic. Thus it changes no inertia
coordinate in (D.69.8).

Use the iterated Frattini quotients
\(Q_r=P_\ell/\Phi^r(P_\ell)\), of exponent
dividing \(\ell^r\), and the actual image
\(H_r\) of \(H_\ell\) therein. With
\[
 S_r(a)=\{\ell,\infty\}\cup
              \{q:a_q\notin\ell^r\mathbb Z_\ell\},
                                                        \tag{D.69.9}
\]
every reduced integral cocycle lies in the
finite actual set
\(Z^1(G_{\mathbb Q,S_r(a)},H_r;\rho_{a,r})\).
The proof is exactly the integral assertion
(D.69.8): at a later nonzero inertia prime,
its coboundary is for an inner power divisible
by \(\ell^r\), hence trivial on \(Q_r\).
At zero inertia it is already trivial.
The finite sets and actions are constant under
changes by \(\ell^r c_0\), on a fixed sign
component when \(\ell=2\).

For every reduced cocycle nontrivial on the
action kernel, apply Chebotarev to its actual
finite affine image and choose a prime with
identity action and nonidentity cocycle value.
Avoid the finite support of the low-valuation
coefficients of the original parameter, and
add \(\ell^M\kappa(q)\), with \(M\ge r\)
and arbitrarily large. This creates a nonzero
inertia coordinate of valuation exactly \(M\)
without changing any of the finite data.
The finite affine image really contains the
chosen kernel element, as in D.67; hence
\cite[Chapter V, Theorem 3.23]{PBC_MilneCFT}
applies without an assumption about the existence
of a Frobenius class. The finitely many tests
at each precision define an open dense set,
including in each fixed-sign component at two.

Intersect these sets over \(r\), and also
require nonzero inertia away from \(\ell\).
For every surviving parameter all cocycles
vanish on the finite action kernels, by
(D.69.8) at their witness primes. They descend
continuously to the actual compact action
image on \(H_\ell\). A nonzero inertia prime
gives a normal subgroup of inner powers of
\(x^{a_q}\), whose fixed subgroup in
\(H_\ell\) is trivial. Its restriction is
an integral coboundary by (D.69.8); conjugate
by its element of \(H_\ell\) to make that
restriction trivial. Comparing the cocycle at
\(ng=g(g^{-1}ng)\) makes every remaining
value fixed by that normal subgroup, hence
trivial. This proves a dense \(G_\delta\)
set \(\mathcal G_\ell\subset\mathcal K_\ell\)
with singleton full pro-\(\ell\) fibre.
Intersect it with the complement of the
countable union of finite-support submodules.
It remains dense \(G_\delta\), and every
parameter is infinitely ramified. At two
take the union of the constructions on the
two clopen sign components. At five we may
also intersect with the loci already proved
in D.66--D.68.

\paragraph{The full profinite comparison is a theorem about sections.}
Put
\[
                   \widehat{\mathcal G}
                           =\prod_\ell\mathcal G_\ell.
                                                        \tag{D.69.10}
\]
This is dense \(G_\delta\) in (D.69.1).
Indeed it is a countable intersection of
inverse images of dense open sets under the
open coordinate projections; every basic
product neighborhood restricts only finitely
many coefficient primes. Its prescribed cusp
section exists in the full profinite group:
the product of the component Kummer cocycles
is a continuous \(\widehat{\mathbb Z}(1)\)-cocycle,
which is inserted into the actual inertia
subgroup \(\langle x\rangle\) of the rational cusp.

Let \(s\) be any full profinite section with
abelian class \((a,0)\), for
\(a\in\widehat{\mathcal G}\). Its image in
every actual geometrically pro-\(\ell\)
quotient has the same class as this fixed
full cusp section, by the singleton result
just proved. D.24 applies to precisely these
two actual full sections: one is arbitrary,
and the other is a fixed member of the
packet of a rational cusp. It concludes
that their full section classes agree.
Its required local point-or-cusp theorem is
the full-formation case of D.56, as already
recorded in D.24. This is not an assertion
that the geometric pro-prime quotients
jointly detect every element of \(F\).
It uses the proved primewise-to-full theorem
for section classes and gives exactly the
singleton in (D.69.3).

Finally, the \(\mathbb Q\)-defined permutations
of the three punctures transport the actual
arithmetic section sets and cusp packets.
On packet parameters they give continuous
affine bijections, including the possible
sign constant at two. They therefore carry
the dense \(G_\delta\) just constructed to
the corresponding loci for the other two
cusps. No assertion outside these packet
abelian fibres is used or deduced.
\end{proof}


% D.70. Root and PBC independent full proof and primary-source audits PASS.
\subsection{Generic cusp rigidity over a number field}

\begin{statement}{D.70}{Proved}{Generic full pro-prime cusp fibres on affine curves and Fermat covers}
Let \(k\) be a number field, let \(U/k\) be a smooth geometrically
connected affine hyperbolic curve, and suppose that its smooth
compactification \(X\) has at least three geometric boundary
points and at least one \(k\)-rational boundary point.
Fix such a rational boundary point \(d\) and a nonzero rational
tangent there. No other boundary point is required to be rational. Use its
actual tangential section \(s_0\) as base point for section cocycles.
For each prime \(\ell\), put
\[
 P_\ell=\pi_1(U_{\bar k})^{(\ell)},\qquad
 V_\ell=P_\ell^{\rm ab},\qquad
 {\cal K}_{k,\ell}=H^1(G_k,\mathbb Z_\ell(1)).
                                                        \tag{D.70.1}
\]
The peripheral inclusion at \(d\) gives an injection
\[
 \iota_{\ell,*}:{\cal K}_{k,\ell}
                    \longrightarrow H^1(G_k,V_\ell).   \tag{D.70.2}
\]
Give \({\cal K}_{k,\ell}\) its intrinsic \(\ell\)-adic topology.
There is a dense \(G_\delta\) subset
\({\cal G}_{k,\ell,d}\subset{\cal K}_{k,\ell}\) such that, for
every \(a\) in this subset, the fibre
\[
 H^1_{\rm cts}(G_k,P_\ell)\longrightarrow H^1(G_k,V_\ell)
           \quad\hbox{over }\iota_{\ell,*}(a)             \tag{D.70.3}
\]
is the singleton consisting of the actual completed cusp section
with parameter \(a\). The subset can be chosen so that every
parameter is ramified at infinitely many finite places of \(k\).
This assertion uses no local pro-\(\ell\) section conjecture.
No corresponding assertion for one or two geometric cusps is made.

Give
\[
 \widehat{\cal K}_k=H^1(G_k,\widehat{\mathbb Z}(1))
                    =\prod_\ell{\cal K}_{k,\ell}         \tag{D.70.4}
\]
the product topology. There is a dense \(G_\delta\) subset of this
space for which the corresponding \emph{full profinite} abelian
section fibre is also the singleton cusp class. Every primewise
parameter can simultaneously be infinitely ramified. This last
assertion uses D.24 and its proved local input D.56.

In particular, this applies over \(\mathbb Q\) to every odd
Fermat curve
\[
 U_n=\{X^n+Y^n+Z^n=0,\ XYZ\ne0\},\qquad n\ge3,
                                                        \tag{D.70.5}
\]
at each of its three rational cusps, and at every coefficient
prime \(\ell\), including those not dividing \(n\). For the
\(3\)-power and \(5\)-power curves of D.20 it supplies generic
singleton full pro-\(\ell\) fibres at their actual cusp abelian
classes. It does not classify the metabelian fibres in D.20
or give a first finite obstruction for their nonzero members.
\end{statement}

\begin{proof}
\paragraph{The integral peripheral module.}
Write \(D=X-U\), let \(r=|D(\bar k)|\ge3\), and fix \(\ell\).
The group \(P=P_\ell\) is free pro-\(\ell\) of rank \(2g+r-1\).
Every cusp generator is primitive: in the usual surface
presentation eliminate a different peripheral generator.
Write \(c\) for the primitive generator at \(d\), \(H=P'\),
and \(V=P^{\rm ab}\). The boundary sequence is an exact
sequence of free integral \(\mathbb Z_\ell\)-modules
with continuous Galois action,
\[
 0\longrightarrow B
   \longrightarrow V\longrightarrow T_\ell J_X
   \longrightarrow0,\qquad
 B=\left(\bigoplus_{e\in D(\bar k)}\mathbb Z_\ell e\right)
                         \Big/\mathbb Z_\ell\sum_e e\ (1).
                                                        \tag{D.70.6}
\]
Here the final \((1)\) twists the whole permutation quotient.
Let \(I=\mathbb Z_\ell(1)c\subset B\). It is primitive, and
\(B/I\) is free of rank \(r-2\); hence \(V/I\) is also
torsion-free. Any two distinct peripheral lines have zero
intersection in \(B\). Indeed, if \(a e-b f\) is a multiple
of \(\sum_e e\), a third cusp coordinate forces that multiple
to be zero, and then \(a=b=0\).

Choose a finite set \(S\) of places containing the archimedean
places, the places above \(\ell\), all bad reduction places of
the marked curve \((X,D)\), and the places at which the chosen
tangent is not integral and nonzero on the special fibre.
At \(v\notin S\), the marked model is smooth with etale
boundary. A finite unramified extension of \(k_v\) splits
all its boundary points. On that extension \(B/I\) is a
direct sum of copies of \(\mathbb Z_\ell(1)\), and its
Frobenius acts by a positive power of \(N(v)\), which is
not one. Thus \((B/I)^{G_{k_v}}=0\), integrally.
Also
\[
                  (T_\ell J_X)^{G_{k_v}}=0.             \tag{D.70.7}
\]
For this, good reduction and \(\ell\ne\operatorname{char}
\kappa(v)\) inject \(J_X(k_v)[\ell^\infty]\) into the
finite group \(J_X(\kappa(v))\): the kernel of reduction
is pro-residue-characteristic. A finite \(\ell\)-primary
group has zero Tate module. Taking invariants in (D.70.6),
and then in its quotient by \(I\), gives
\[
 V^{G_{k_v}}=(V/I)^{G_{k_v}}=0.
                                                        \tag{D.70.8}
\]
The same vanishings hold for \(G_k\), by restriction to
one such decomposition group. The continuous cohomology
sequence of \(0\to I\to V\to V/I\to0\) therefore proves
both (D.70.2) and the local injections
\[
 H^1(k_v,I)\lhook\joinrel\longrightarrow H^1(k_v,V)
                    \qquad(v\notin S).                 \tag{D.70.9}
\]
No rationalization is used: the coefficient sequence
has a continuous integral linear splitting after
forgetting the Galois action.

\paragraph{Completed multiplicative groups and actual cocycles.}
The finite Kummer sequences and the Mittag--Leffler
property of the finite groups \(H^0(k,\mu_{\ell^n})\) give
\[
 {\cal K}_{k,\ell}=\varprojlim_n
                         k^\times/(k^\times)^{\ell^n}. \tag{D.70.10}
\]
Here each finite-exponent quotient is discrete and need
not be finite. We make the topology and representatives
explicit. The principal divisor subgroup
\(\operatorname{Prin}(k)\subset\bigoplus_{v\ {\rm finite}}
\mathbb Z v\) has finite index, with quotient the ideal
class group; its finiteness is
\cite[Theorem 4.4]{DE_MilneANT}. It is free abelian.
The exact sequence
\[
 1\longrightarrow\mathcal O_k^\times
 \longrightarrow k^\times\longrightarrow
             \operatorname{Prin}(k)\longrightarrow0
                                                        \tag{D.70.11}
\]
splits as abstract abelian groups because its quotient
is free. The unit theorem
\cite[Theorem 5.1]{DE_MilneANT} then gives a decomposition
\[
       k^\times=\mu(k)\times L,\qquad
       L\simeq\bigoplus_{i\in{\cal I}}\mathbb Z,
                                                        \tag{D.70.12}
\]
where \({\cal I}\) is countable; its finite unit-rank
part is included in this index set. Fix one such
decomposition. Consequently
\[
 {\cal K}_{k,\ell}\simeq
       \mu(k)(\ell)\times c_0({\cal I},\mathbb Z_\ell).
                                                        \tag{D.70.13}
\]
The finite root-of-unity factor is discrete, and the
second factor has the uniform \(\ell\)-adic topology.
This is the topology in (D.70.10), and also the intrinsic
topology with neighborhoods \(\ell^n{\cal K}_{k,\ell}\);
it is independent of the chosen splitting. It is a
complete separable metric topology.

Fix compatible roots for every generator of \(L\),
and thereby an integral continuous Kummer cocycle
\(\kappa_i:G_k\to\mathbb Z_\ell(1)\). Their uniformly
convergent sums
\(\sum_i a_i\kappa_i\), for \((a_i)\in c_0\), give a
continuous linear choice of actual cocycles on the
free factor. For each of the finitely many torsion
classes in \(\mu(k)(\ell)\), choose one actual integral
continuous Kummer representative. On each fixed
torsion component, add that representative to the
free-factor sum. This gives a continuous choice of
actual cocycle \(a(g)\) for every parameter. We do
not assert a linear section on the torsion factor,
at two or at any other prime.

For each finite place let \(a_v\in\mathbb Z_\ell\)
be its valuation coordinate. The valuation map is
continuous, and for each \(m\) only finitely many
\(a_v\) are nonzero modulo \(\ell^m\). To see this
directly, the class modulo \(\ell^m\) in (D.70.10)
has a representative in \(k^\times\), whose divisor
has finite support. At \(v\nmid\ell\), restriction
of the actual cocycle to inertia is precisely
\(a_v\) times the integral tame \(\ell\)-character.
Roots of unity and local units have zero such inertia.

Let \(h\) be the class number of \(k\). For every
finite \(v\), choose \(\alpha_v\in k^\times\) with
\[
                       \operatorname{div}(\alpha_v)=h v.
                                                        \tag{D.70.14}
\]
Multiply it by its inverse torsion component in
(D.70.12), so that \(\alpha_v\in L\) without changing
its divisor. Its cocycle is now the prescribed finite
integral linear combination of the \(\kappa_i\).
The perturbation \(\ell^M\kappa(\alpha_v)\) preserves
the torsion component, is uniformly divisible by
\(\ell^M\) as an actual cochain, and changes only
the valuation at \(v\), by \(h\ell^M\).

\paragraph{The actual local cusp comparison.}
The rational tangential action
\(\sigma_0:G_k\to\operatorname{Aut}(P)\) satisfies
\(\sigma_0(g)(c)=c^{\chi_\ell(g)}\). Its restriction
at \(v\notin S\) is unramified. This is the unit
tangential section of the marked good-reduction
specialization; a unit tangent has trivial
prime-to-residue-characteristic Kummer inertia.
The completed parameter gives the actual action
\[
 \sigma_a(g)=\operatorname{Int}(c^{a(g)})\sigma_0(g).
                                                        \tag{D.70.15}
\]
Every section class over \(\iota_{\ell,*}(a)\) can
be conjugated to have exactly that abelian cocycle.
Its difference from this cusp section is an actual
nonabelian cocycle \(z:G_k\to H\) for \(\sigma_a\).
Conversely every such cocycle supplies a section
in the fibre. By \(V^{G_k}=0\), every conjugator
between two exact representatives lies in \(H\).
Thus the fibre in (D.70.3) is exactly
\(H^1(G_k,H;\sigma_a)\).

Fix \(v\notin S\). Use the actual cartesian
specialization square of
\cite[Definition A.3 and Remark A.3.1, p.~38]{DE_Hoshi2012}
with \(\Sigma=\{\ell\}\). Proposition A.6 of that
source says that the specialized inertia of every
actual local section is zero or a nonzero subgroup
of one cusp inertia group. This is a statement for
ordinary geometrically pro-\(\ell\) sections of
every hyperbolic curve with good reduction, as the
setup on pp.~36--37 specifies.

For completeness the further group step is the
same integral one used in D.67. A primitive cusp
generator in a free pro-\(\ell\) group has
\[
 C_P(c^b)=\langle c\rangle
                    \quad(0\ne b\in\mathbb Z_\ell).
                                                        \tag{D.70.16}
\]
Choose a free basis containing \(c\). In its Magnus
embedding, commuting with \((1+C)^b\) implies
commuting with \(C\), because
\((1+C)^b-1\) has a compositional inverse over
\(\mathbb Q_\ell\). Comparing coefficients of
words involving the other basis variables gives
the centralizer \(\mathbb Q_\ell[[C]]\) in the
completed associative algebra. Retraction killing
those other variables then identifies a group
element in it as the power of \(c\) specified by
its integral first abelian coordinate. A geometric
normalizer of a nonzero subgroup of \(\langle c\rangle\)
centralizes it: its exponent automorphism is
forced to be one by abelianization.

If \(a_v\ne0\), the actual section has nonzero
specialized inertia, and its abelian image has
line \(\mathbb Z_\ell c\). The disjointness of
the geometric peripheral lines proved after
(D.70.6) forces the cusp in Hoshi's alternative
to be precisely \(d\). This comparison can also
be made after the finite unramified extension
splitting all boundary points; inertia is
unchanged by that extension. It distinguishes
every geometric cusp, including the extra
nonrational ones.

The whole specialized section normalizes its
inertia image. Comparing it with the rational
tangential lift at \(d\), (D.70.16) puts its
image in that cusp decomposition group.
The cartesian square pulls this conclusion
back to the actual local cusp packet. By
(D.70.9), its parameter equals \(a\) in
\(H^1(k_v,I)\). First conjugate into the
packet, and then conjugate by a cusp-inertia
element to adjust its Kummer representative
to the reference one. The resulting
conjugator has \(G_{k_v}\)-fixed abelian
image since the two actual abelian cocycles
agree. Equation (D.70.8) makes that image
zero, so the conjugator lies in \(H\).
If \(a_v=0\), Hoshi's nonzero alternative
would have nonzero abelian inertia, which
is impossible; both sections have zero
specialized inertia. We have proved the
two exact assertions
\[
 \begin{cases}
 z|_{G_{k_v}}\text{ is an }H\text{-coboundary},
                              &a_v\ne0,\\
 z(I_v)=1,                    &a_v=0,
 \end{cases}\qquad v\notin S.                         \tag{D.70.17}
\]
This argument uses the integral quotient \(V/I\),
not a rationalized abelian Kummer map.

\paragraph{Finite capture and fresh-prime witnesses.}
Let \(P_0=P\), \(P_{m+1}=\Phi(P_m)\), and put
\[
 Q_m=P/P_m,\quad H_m=\operatorname{im}(H\to Q_m),\quad
 S_m(a)=S\cup\{v:a_v\notin\ell^m\mathbb Z_\ell\}.
                                                        \tag{D.70.18}
\]
The \(P_m\) are characteristic open subgroups
with trivial intersection, and \(Q_m\) has
exponent dividing \(\ell^m\). These finite
quotients induce the topology of \(H\).
Write \(\rho_{a,m}\) for its actual action on
\(H_m\). Both \(S_m(a)\) and this action are
constant when the torsion component of \(a\)
is fixed and its free component is fixed
modulo \(\ell^m\).

Every actual cocycle \(z:G_k\to H\) reduces
to the finite set
\[
 E_m(a)=Z^1(G_{k,S_m(a)},H_m;\rho_{a,m}).
                                                        \tag{D.70.19}
\]
Indeed outside \(S_m(a)\) its inertia is zero
if \(a_v=0\). Otherwise (D.70.17) makes it
an integral conjugation coboundary; the
inertia action is
\(\operatorname{Int}(c^{a(\tau)})\), which is
trivial on \(Q_m\) because \(a(\tau)\) is
divisible by \(\ell^m\). The finite action
itself has this same support bound. Thus
the reduced \emph{actual} affine
representation factors through the stated
finite-ramification Galois group.
The set in (D.70.19) is finite by the
finite-extension finiteness theorem: its
cocycles are homomorphisms into the fixed
finite group \(H_m\rtimes\operatorname{im}\rho_{a,m}\),
with fixed second coordinate and ramification
in the fixed finite set \(S_m(a)\).

Call \(z\in E_m(a)\) exceptional if its
restriction to \(\ker\rho_{a,m}\) is nontrivial.
Let \({\cal V}_m\) be the set of parameters
for which each exceptional \(z\) has a
witness place \(v\notin S_m(a)\) with
\[
 \rho_{a,m}(\operatorname{Frob}_v)=1,\qquad
 z(\operatorname{Frob}_v)\ne1,\qquad a_v\ne0.
                                                        \tag{D.70.20}
\]
This is open: all finite data are locally
constant, there are only finitely many
exceptional cocycles, and nonvanishing of
finitely many specified valuation coordinates
is open.

It is dense on every fixed torsion component.
Given a neighborhood with free coordinates
fixed modulo \(\ell^N\), choose
\(M\ge\max\{m,N\}\) and write \(w=v_\ell(h)\).
For an exceptional \(z\), a group element
\(\delta\in\ker\rho_{a,m}\) with
\(z(\delta)\ne1\) gives the actual element
\((z(\delta),1)\) in its finite affine image.
Chebotarev, in precisely the finite-image
form of \cite[Chapter V, Theorem 3.23]{PBC_MilneCFT},
provides infinitely many places in its
conjugacy class. The two finite tests in
(D.70.20) are invariant under conjugacy.
Choose distinct such places for the
finitely many exceptional cocycles, avoiding
\(S_m(a)\) and the finite set where
\(a_v\notin\ell^{M+w+1}\mathbb Z_\ell\).
Add \(\ell^M\kappa(\alpha_v)\) at each
chosen place. By (D.70.14), the new
coordinate at that place has valuation
exactly \(M+w\); in particular it is
nonzero. All other chosen valuation
coordinates remain unchanged. The actual
finite action, support set, cocycle list,
torsion component, and prescribed
neighborhood are preserved. This proves
density. No principal-prime restriction
has been imposed on Chebotarev.

For every finite set \(T\) of finite places,
the locus \(a_v=0\) for all \(v\notin T\)
is closed and has empty interior: a fresh
perturbation (D.70.14), at sufficiently
large depth, leaves any neighborhood.
There are countably many such \(T\).
The Baire theorem applied to (D.70.13)
therefore gives a dense \(G_\delta\) set
\[
 {\cal G}_{k,\ell,d}\subset\bigcap_{m\ge1}{\cal V}_m
                                                        \tag{D.70.21}
\]
whose parameters all have infinite valuation
support. In particular each has \(a_v\ne0\)
at some \(v\notin S\).

\paragraph{Descent and the normal inner inertia subgroup.}
Fix \(a\) in (D.70.21) and an actual
cocycle \(z:G_k\to H\). None of its
reductions in (D.70.19) is exceptional.
Otherwise its witness in (D.70.20) and
the local coboundary assertion (D.70.17)
would give simultaneously a nontrivial
and a trivial value at a Frobenius with
identity action. Hence \(z\) is trivial
on the kernel of the actual action on
\(H\), and descends continuously to
\[
 {\cal I}_a=\operatorname{im}\left(
     G_k\longrightarrow\prod_m\operatorname{Aut}(H_m)\right).
                                                        \tag{D.70.22}
\]
This is the compact actual action image.
Continuity follows from the quotient map
of compact groups.

Choose \(v\notin S\) with \(a_v\ne0\).
Its inertia image in (D.70.22) is
\[
 N_v=\{\operatorname{Int}(c^{a_v t})|_H:
                                  t\in\mathbb Z_\ell\}.
                                                        \tag{D.70.23}
\]
It is normal in \({\cal I}_a\), since every
operator in (D.70.15) sends \(c\) to
\(c^{\chi_\ell(g)}\), and \(\chi_\ell(g)\)
is a unit. Equation (D.70.16) and the
primitive abelian image of \(c\) imply
\[
                    H^{N_v}=H\cap\langle c\rangle=1.
                                                        \tag{D.70.24}
\]
This includes nonprimitive \(a_v\); no
division by \(a_v\) is used in \(H\).
By (D.70.17) the descended cocycle on
\(N_v\) is an integral \(H\)-coboundary.
Conjugate it by that element of \(H\) to
make it trivial on \(N_v\). Comparing its
values at \(ng=g(g^{-1}ng)\) shows that
every value is fixed by \(N_v\), hence
is one by (D.70.24). Thus the whole
cocycle is a coboundary. This proves
the asserted singleton in (D.70.3).

\paragraph{The full profinite fibre and the actual Fermat maps.}
Continuous cochains into the product
\(\widehat{\mathbb Z}(1)=\prod_\ell\mathbb Z_\ell(1)\)
are exactly products of continuous
component cochains, and the same is true
of coboundaries. This proves (D.70.4).
It is also the power-completion topology
\(\varprojlim_n k^\times/(k^\times)^n\)
with discrete quotients. The countable
product of the complete separable metric
spaces in (D.70.13) is again such a space.
The product of the sets (D.70.21) is a
dense \(G_\delta\): write each as a
countable intersection of open dense
sets and pull those sets back by the
coordinate projections.

For every parameter in this product,
the product of the chosen component
Kummer cocycles gives an actual continuous
\(\widehat{\mathbb Z}(1)\)-valued cocycle,
hence a full cusp section. Any other
actual full section with its abelian
class has, by (D.70.3), the same class
as this one in every geometric pro-\(\ell\)
quotient. Apply precisely D.24, with its
D.56 local input, to these two actual
full sections. They have the same full
class. No assertion that primewise
geometric quotients detect arbitrary
geometric elements is used.

For (D.70.5), the actual finite etale
map to the tripod is
\[
 t=-\left(\frac XZ\right)^n,\qquad
 1-t=-\left(\frac YZ\right)^n.                          \tag{D.70.25}
\]
It has degree \(n^2\); its geometric
deck group is \(\mu_n^2\). The three
rational cusps are \([0:-1:1]\),
\([-1:0:1]\), and \([1:-1:0]\).
There are \(3n\) geometric cusps, so
the preceding local argument distinguishes
also all the additional nonrational ones.
At \(O=[0:-1:1]\), with coordinate
\(u=X/Z\), the local map is exactly
\(t=-u^n\). Consequently its actual
cusp abelian parameter maps to the
tripod pair
\[
                  \bigl(\kappa(-1)+na,\ 0\bigr).
                                                        \tag{D.70.26}
\]
The sign term is retained at two.
This identifies the cusp and its
abelian map with the ones used in
D.20, rather than introducing an
abstract free-group model. The proof
of (D.70.3) takes place in the actual
group of \(U_n\) itself, so it also
applies when \(\ell\nmid n\), where
one cannot infer injectivity merely
from the map to the tripod.
The zero-abelian metabelian families
of D.20 lie at \(a=0\), outside our
infinitely ramified generic set.
No claim about their metabelian
fibres is changed.
\end{proof}


% D.71. Root and PBC independent full proof and primary-source audits PASS.
\subsection{Generic class-two rigidity on the once-punctured elliptic curve}

\begin{statement}{D.71}{Proved}{The requested one-cusp curve has generic singleton class-two fibres}
Let \(E/k\) be an elliptic curve over a number field, let
\(U=E-\{O\}\), and fix a nonzero rational tangent at \(O\).
Use the resulting actual cusp section \(s_0\) as base point.
For every prime \(\ell\), put
\[
 P_\ell=\pi_1(U_{\bar k})^{(\ell)},\qquad
 T_{2,\ell}=P_\ell/\gamma_3P_\ell,\qquad
 V_\ell=T_\ell E,\qquad
 {\cal K}_{k,\ell}=H^1(G_k,\mathbb Z_\ell(1)).
                                                        \tag{D.71.1}
\]
The actual cusp packet gives a bijection
\[
 {\cal K}_{k,\ell}\ \xrightarrow{\ \sim\ }\
 \ker\!\left(H^1(G_k,T_{2,\ell})\longrightarrow
                                  H^1(G_k,V_\ell)\right).
                                                        \tag{D.71.2}
\]
For a dense \(G_\delta\) set of parameters in the intrinsic
\(\ell\)-adic topology on \({\cal K}_{k,\ell}\), the full
pro-\(\ell\) section fibre above its class in (D.71.2)
is a singleton, represented by that actual completed cusp.
The parameters may all be chosen infinitely ramified.

The same statement holds for the actual full profinite
class-two quotient and a dense \(G_\delta\) set in
\(\widehat{\cal K}_k=\prod_\ell{\cal K}_{k,\ell}\):
its full profinite section fibre is the singleton cusp
class. All primewise parameters may simultaneously
be infinitely ramified. This conclusion uses precisely
D.24 and D.56.

In particular these assertions hold for the specified
\[
 k=\mathbb Q(i),\qquad E:y^2=x^3-x,\qquad U=E-\{O\}.
                                                        \tag{D.71.3}
\]
Every cusp parameter has the same abelian section
class, but distinct parameters already have distinct
class-two classes. Thus the abelian fibre itself
contains an injected copy of the entire completed
cusp packet and cannot be a singleton. The new
generic uniqueness is at class two. No first finite
nilpotent or derived obstruction is asserted.
\end{statement}

\begin{proof}
\paragraph{The actual integral class-two group.}
Fix \(\ell\), write \(P=P_\ell\), and let \(c\)
generate the actual cusp inertia. The group \(P\)
is free pro-\(\ell\) on two generators. In a
surface basis the peripheral word is a commutator
or its inverse, according to orientation.
Consequently its class in
\(\gamma_2P/\gamma_3P\) is a primitive generator.
The commutator pairing and the Weil pairing give
the integral equivariant identifications
\[
 P^{\rm ab}=V=T_\ell E,\qquad
 I:=\gamma_2P/\gamma_3P
       =\bigwedge^2 V\simeq\mathbb Z_\ell(1).
                                                        \tag{D.71.4}
\]
In particular these statements hold at two; no
Baker--Campbell--Hausdorff denominator is used.
The actual central extension is
\[
 1\longrightarrow I\longrightarrow T_2
                     \longrightarrow V\longrightarrow1.
                                                        \tag{D.71.5}
\]
The map \(\langle c\rangle\to I\) is an isomorphism
of \(\mathbb Z_\ell(1)\)-modules. The tangential
action \(\sigma_0\) satisfies
\(\sigma_0(g)c=c^{\chi_\ell(g)}\).

Choose a finite set \(S\) containing all infinite
places, all places above \(\ell\), the bad
reduction places of \(E\), and the places where
the fixed tangent is not an integral nonzero
tangent of the smooth model. At \(v\notin S\),
\[
             V^{G_{k_v}}=0,\qquad I^{G_{k_v}}=0.
                                                        \tag{D.71.6}
\]
The first assertion follows from injection of
local prime-to-residue-characteristic torsion
into the finite reduction group, exactly as
proved in D.70. The second follows because
Frobenius acts on \(I\) by \(N(v)\ne1\).
Restriction to one such place gives the same
global vanishings. A fixed element of \(T_2\)
has zero image in \(V\) and then lies in
\(I^{G_{k_v}}\), or in \(I^{G_k}\), respectively.
Thus
\[
             T_2^{G_{k_v}}=T_2^{G_k}=1
                         \qquad(v\notin S).             \tag{D.71.7}
\]
These are vanishings for the actual integral
groups.

Any \(T_2\)-cocycle with zero abelian class
can be conjugated to have zero actual abelian
cocycle, by lifting a conjugator from \(V\).
It is then an \(I\)-valued cocycle. A conjugator
between two such representatives has
Galois-fixed image in \(V\), which is zero.
It therefore belongs to \(I\), and its
equivalence relation is exactly the usual
integral \(I\)-coboundary relation. This
proves (D.71.2), globally and also locally
at every \(v\notin S\). Every parameter is
supplied by the actual cusp cocycle
\(g\mapsto c^{a(g)}\).

\paragraph{The centralizer of the nonprimitive peripheral word.}
We prove the required free pro-\(\ell\) fact,
including its subgroup input. A finitely
generated closed subgroup \(L\) of a finite
rank free pro-\(\ell\) group \(P\) is a
free factor of some open subgroup. This
is \cite[Theorem 3.2, pp.~317--318]{DE_Lubotzky1982};
here is its finite-character proof in the
form needed here.

Choose a basis of the finite vector space
\(\operatorname{Hom}_{\rm cts}(L,\mathbb F_\ell)\).
For a common sufficiently small open normal
subgroup \(W\lhd P\), every chosen character
vanishes on \(L\cap W\). Each extends to
the open subgroup \(LW\) by
\(\phi(lw)=\phi(l)\). This is well-defined
and a homomorphism because \(W\) is normal.
Thus \(L/\Phi(L)\to LW/\Phi(LW)\) is injective.
The group \(LW\) is finite rank free pro-\(\ell\),
by the open-subgroup Schreier theorem.
A minimal generating set of \(L\) therefore
extends to a minimal generating set of \(LW\).
Such a set is a free basis: the resulting
surjective endomorphism of a finite rank free
pro-\(\ell\) group is an isomorphism. The
Hopf property follows directly because a
finitely generated profinite group has only
finitely many open subgroups of each index;
a surjective endomorphism permutes their
finite intersections and is injective on
the inverse limit of the resulting finite
quotients. The generators of \(L\) are
therefore part of a free basis, as claimed.

The peripheral element \(c\) is not an
\(\ell\)-th power in \(P\). If \(c=t^\ell\),
its zero abelian image forces \(t\in\gamma_2P\),
since \(V\) is torsion-free. Its primitive
class in \(\gamma_2P/\gamma_3P\) would then
be divisible by \(\ell\), a contradiction.
For \(0\ne b\in\mathbb Z_\ell\) and
\(g\in C_P(c^b)\), consider the closed
two-generated subgroup \(L=\langle g,c\rangle\).
The preceding lemma makes it free pro-\(\ell\).
If it had rank two, \(g,c\) would be a
free basis and \([g,c^b]\ne1\). One sees
the latter directly in the Magnus embedding:
the coefficient \(b\) of the degree-two
word \(GC-CG\) is nonzero in
\(\mathbb Z_\ell\). Hence \(L\) is procyclic.
As \(c\) has no \(\ell\)-th root in \(P\),
it generates \(L\). We obtain
\[
                   C_P(c^b)=\langle c\rangle
                         \quad(b\ne0).                 \tag{D.71.8}
\]
This proof does not treat \(c\) as a primitive
free generator.

The geometric normalizer of any nonzero
closed subgroup of \(\langle c\rangle\)
also equals \(\langle c\rangle\).
Indeed, if conjugation sends \(c^b\)
to \(c^{bu}\), its effect on
\(\gamma_2P/\gamma_3P\) is trivial.
Thus \(b[c]=bu[c]\); the primitive
torsion-free class forces \(u=1\),
and (D.71.8) applies.
Finally, for \(H=\gamma_3P\),
\[
                 H\cap\langle c\rangle=1.              \tag{D.71.9}
\]

\paragraph{The exact fibre and its local comparison.}
Use the actual Kummer representatives
on each fixed torsion component of
\({\cal K}_{k,\ell}\) constructed in
D.70. Define
\[
       \sigma_a(g)=\operatorname{Int}(c^{a(g)})\sigma_0(g)
                         \quad\hbox{on }P\hbox{ and }H.
                                                        \tag{D.71.10}
\]
On \(T_2\) this action equals \(\sigma_0\),
because the image of \(c\) is central.
A full pro-\(\ell\) section with the
same class-two class as this cusp can
be conjugated to have exactly its
class-two cocycle. Its difference is
then a continuous \(H\)-cocycle.
Conversely every such cocycle gives
a section in the fibre. A conjugator
between exact representatives projects
to a fixed element of \(T_2\); (D.71.7)
makes it trivial. Thus the exact fibre is
\[
                         H^1(G_k,H;\sigma_a).          \tag{D.71.11}
\]

At \(v\notin S\), apply the cartesian
good-reduction specialization square
and the inertia alternative of
\cite[Definition A.3, Remark A.3.1 and Proposition A.6]
{DE_Hoshi2012} with \(\Sigma=\{\ell\}\).
There is only one geometric cusp.
If \(a_v\ne0\), the prescribed
class-two inertia \(c^{a_v t}\)
is nonzero. The actual section
therefore has nonzero specialized
inertia in that cusp. Its whole
image normalizes this inertia;
(D.71.8) and its normalizer
consequence, compared with the
tangential arithmetic lift, place
the image in the cusp decomposition
group. The cartesian square puts
the actual local section in the
actual local pro-\(\ell\) packet.
The local version of (D.71.2)
shows that its Kummer parameter
is exactly \(a\) as a cohomology
class. Adjusting the packet
conjugator by a cusp-inertia
element makes the two class-two
cocycles agree exactly. Its
projection to \(T_2\) is then
fixed, hence zero by (D.71.7);
the conjugator lies in \(H\).

If \(a_v=0\), nonzero specialized
inertia in the sole cusp would
have a nonzero image in \(I\),
contrary to the prescribed zero
class-two inertia. Both actual
sections therefore have trivial
specialized inertia. For every
cocycle \(z:G_k\to H\) we have
the exact assertions
\[
 \begin{cases}
 z|_{G_{k_v}}\text{ is an }H\text{-coboundary},
                                  &a_v\ne0,\\
 z(I_v)=1,                        &a_v=0,
 \end{cases}\qquad v\notin S.                         \tag{D.71.12}
\]

\paragraph{The Baire construction in the genuine fibre.}
All parameter-space and arithmetic
constructions in D.70 now apply,
with the following precise
substitutions. Use the same complete
space
\(\mu(k)(\ell)\times c_0({\cal I},\mathbb Z_\ell)\),
with fixed actual torsion cocycles.
Use \(P_m=\Phi^m(P)\), \(Q_m=P/P_m\),
and now
\[
 H_m=\operatorname{im}(\gamma_3P\to Q_m),\qquad
 S_m(a)=S\cup\{v:a_v\notin\ell^m\mathbb Z_\ell\}.
                                                        \tag{D.71.13}
\]
The exponent of \(Q_m\) divides
\(\ell^m\). By (D.71.12), every
integral \(H\)-cocycle reduces to
the finite set
\[
                 Z^1(G_{k,S_m(a)},H_m;\rho_{a,m}).
                                                        \tag{D.71.14}
\]
In particular this finite capture
includes all unrestricted global
cocycles. At a later nonzero
inertia place, its integral local
conjugation coboundary is trivial
on \(Q_m\) because the exponent
of \(c\) is divisible by
\(\ell^m\); at a zero inertia
place use the second line of
(D.71.12). The finite action
and support are constant on
each fixed torsion component
and free-parameter coset
modulo \(\ell^m\).

Impose, for each finite cocycle
nontrivial on the action kernel,
a witness place with identity
Frobenius action, nontrivial
cocycle value, and nonzero
\(a_v\). These conditions are
open. They are dense by the
actual finite affine-image
Chebotarev argument of D.70:
an element in the action kernel
with nontrivial cocycle value
first gives an element of the
actual finite affine image.
Choose its Frobenius place
outside all current finite
supports, and perturb by
\(\ell^M\kappa(\alpha_v)\),
where \(\operatorname{div}(\alpha_v)=hv\).
Avoid valuations not divisible
by \(\ell^{M+v_\ell(h)+1}\).
The new nonzero valuation is
exactly \(M+v_\ell(h)\), while
the finite action and prescribed
neighborhood remain unchanged.
This works also on every torsion
component at two.

Intersect these countably many
open dense conditions and the
complements of all finite-support
loci. The Baire theorem gives
the asserted dense \(G_\delta\)
of infinitely ramified parameters.
For such a parameter, (D.71.12)
excludes every exceptional finite
cocycle at its witness place.
Every actual \(H\)-cocycle
therefore descends continuously
to the actual compact image
\[
 {\cal I}_a=\operatorname{im}\left(
       G_k\longrightarrow\prod_m\operatorname{Aut}(H_m)\right).
                                                        \tag{D.71.15}
\]
Choose \(v\notin S\) with \(a_v\ne0\).
Its inertia image is the normal
subgroup
\[
 \{\operatorname{Int}(c^{a_v t})|_H:t\in\mathbb Z_\ell\}
                    \ \lhd\ {\cal I}_a.                \tag{D.71.16}
\]
Normality follows from
\(\sigma_a(g)c=c^{\chi_\ell(g)}\).
Its fixed subgroup in \(H\)
is trivial by (D.71.8)--(D.71.9),
even for nonprimitive \(a_v\).
The local coboundary in (D.71.12)
allows one conjugation by an
element of \(H\) to kill the
cocycle on (D.71.16). Comparing
the two cocycle formulas at
\(ng=g(g^{-1}ng)\) then makes
every value fixed by (D.71.16);
every value is one. This proves
the singleton in (D.71.11).

\paragraph{Full profinite class two and the abelian obstruction.}
For the actual full geometric
group \(\pi=\pi_1(U_{\bar k})\),
\[
             \pi/\gamma_3\pi\simeq\prod_\ell T_{2,\ell}.
                                                        \tag{D.71.17}
\]
Indeed every finite class-two
quotient is nilpotent and is
the product of its Sylow groups.
The universal finite quotient
property and compactness give
the displayed isomorphism,
with the actual arithmetic
action on every factor.
Continuous cocycles and their
conjugacy relations into this
product are componentwise.
Thus (D.71.2) also holds with
full coefficients
\(\widehat{\mathbb Z}(1)\)
and parameter space
\(\widehat{\cal K}_k\).

Take the product of the
primewise dense \(G_\delta\)
sets just proved, with the
complete product topology
specified in D.70. The product
is dense \(G_\delta\), and each
parameter gives an actual
full cusp cocycle. Any actual
full section with that class-two
class has the same full
pro-\(\ell\) class as this
cusp at every \(\ell\).
D.24, with D.56 as its
proved local input, identifies
the two full section classes.
This proves the full assertion.

Finally \(c\in\gamma_2P\), and
the analogous full peripheral
element is in \(\gamma_2\pi\).
All cusp parameters therefore
have zero abelian cocycle.
Equation (D.71.2) and its
full version make distinct
parameters distinct already
at class two. The abelian
fibre contains the whole
injected packet, as claimed.
This applies without a change
of field or curve to (D.71.3).
\end{proof}


% D.72. Root, PBC and FGH full independent proof and primary-source audits PASS.
\subsection{Compact exact fibres after fixing class two}

\begin{statement}{D.72}{Proved}{Finite nilpotent failure on the requested curve and the Fermat curves}
Let \(U/k\) be a smooth geometrically connected affine hyperbolic
curve over a number field, with a chosen actual arithmetic section
\(s_0\). A rational point or rational tangential cusp supplies such
a choice in every application below. Fix a coefficient prime
\(\ell\), including \(\ell=2\), and use the action induced by \(s_0\).
Put
\[
 P=\pi_1(U_{\bar k})^{(\ell)},\qquad
 T_n=P/\gamma_{n+1}P.
                                                        \tag{D.72.1}
\]
For every fixed actual continuous cocycle
\(\alpha\in Z^1(G_k,T_2)\), with no ramification restriction,
the space
\[
 {\mathscr Z}_n(\alpha)=
   \{z\in Z^1(G_k,T_n):z\bmod\gamma_3=\alpha\},
                 \qquad n\ge2,                         \tag{D.72.2}
\]
is compact in the uniform topology on actual continuous cochains.
More precisely, it is the image of a closed subset of a finite
power of \(\mathbb Z_\ell\) under a continuous rational-polynomial
family of Hall coordinates with coefficients in fixed finite
spaces of continuous cochains.

The corresponding section-class fibres
\[
 {\mathscr S}_n([\alpha])=
   H^1(G_k,T_n)\mathbin{\times}_{H^1(G_k,T_2)}\{[\alpha]\}
                                                        \tag{D.72.3}
\]
are compact Hausdorff. So are the full pro-\(\ell\) fibre and
their exact versions. All transition maps have closed image.
If a fixed finite-stage class lifts separately to every higher
nilpotent stage, it lifts to an actual continuous \(P\)-section.
There is no assumption that the initially supplied lifts agree.
For a once-punctured elliptic curve these compactness and
compatibility assertions also hold for the actual metabelian
group \(P/P''\) and all its nilpotent quotients.

For every generic parameter furnished by D.71 on the requested
once-punctured elliptic curve, a finite-stage class different
from its cusp class fails to lift at some finite nilpotent stage.
The failure is uniform outside any fixed neighborhood of that
cusp in the compact fibre: for \(m\ge2\) and an open neighborhood
\(W\) of its cusp class, some \(N\ge m\) satisfies
\[
  \operatorname{im}\bigl({\mathscr S}_N([\alpha])
            \longrightarrow{\mathscr S}_m([\alpha])\bigr)
                      \subset W.                       \tag{D.72.4}
\]
The same assertion holds on every Fermat curve in D.70, for
its generic parameters and the fixed class-two image of the
chosen cusp. It holds at every coefficient prime. The all-prime
product statement concerns the actual pronilpotent tower; it
has the same compactness and finite-failure conclusions for
the product generic cusp parameters. It does not identify an
arbitrary full profinite section with a pronilpotent section.

No least stage, numerical bound, or first obstruction class is
computed. In particular this does not replace the open explicit
finite-derived obstruction in D.5.
\end{statement}

\begin{proof}
\paragraph{The central lattices and their weights.}
Write \(X\) for the smooth proper compactification, \(D=X-U\),
and \(V=P^{\rm ab}\). The geometric group \(P\) is free
pro-\(\ell\) of rank \(2g+|D(\bar k)|-1\). Its actual
lower-central factors are finite free lattices
\[
 L_j=\gamma_jP/\gamma_{j+1}P
       =\operatorname{Lie}_j(V),\qquad j\ge1,            \tag{D.72.5}
\]
where the right side is the degree-\(j\) part of the free
integral Lie algebra. Hall's basic-commutator theorem and
its pro-prime completion give this identification
\cite[Theorems 5.6 and 6.5, and completion case (iii)]{DE_HallEdmonton},
as checked in D.68. The integral Hall bases apply in every
class and at every prime. Inner automorphisms act trivially
on each \(L_j\); its arithmetic action therefore depends
only on the canonical action on \(V\), regardless of the
ramification of \(s_0\) or a twisting cocycle.

Choose one finite set \(S\) containing the archimedean places,
the places above \(\ell\), and all places of bad reduction
of the marked pair \((X,D)\). At \(v\notin S\), the
canonical action on \(V\) is unramified. Its boundary and
proper parts fit into the integral exact sequence
\[
 0\longrightarrow
 \left(\mathbb Z_\ell[D(\bar k)]/
                   \mathbb Z_\ell{\textstyle\sum_{d\in D(\bar k)}}d\right)(1)
 \longrightarrow V\longrightarrow T_\ell J_X
 \longrightarrow0.                                    \tag{D.72.6}
\]
The exact boundary description and good marked specialization
are the same as in D.70. Arithmetic Frobenius \(F_v\), with
\(q=N(v)\), has boundary eigenvalues \(q\zeta\) for roots
of unity \(\zeta\). Its proper eigenvalues are algebraic
numbers with every complex absolute value \(q^{1/2}\)
\cite[II, Theorem 1.1; IV, Section 3]{fghMilneAV2008}.
This convention agrees with \(F_v(\tau)F_v^{-1}=\tau^q\)
on tame inertia.

After tensoring with \(\mathbb Q_\ell\), \(L_j\) is a
subspace of \(V^{\otimes j}\), or equivalently a
subquotient suffices here. Its Frobenius eigenvalues are
therefore among products of \(j\) eigenvalues of \(V\).
Their complex absolute values belong to
\[
       \{q^{w/2}: j\le w\le2j\}.                       \tag{D.72.7}
\]
No semisimplicity of the extension (D.72.6) is needed:
characteristic polynomials in an invariant filtration
multiply, and tensor-product eigenvalues are the products
of the eigenvalues of its factors. Thus for every \(j\ge3\)
\[
     q\text{ is not an eigenvalue of }F_v\text{ on }
                                   L_j\otimes\mathbb Q_\ell.
                                                        \tag{D.72.8}
\]
The same reasoning excludes eigenvalue \(1\) in every
positive degree, including degrees one and two.

\paragraph{Every central one-cocycle has fixed finite support.}
Let \(j\ge3\) and let \(z:G_k\to L_j\) be any actual
continuous one-cocycle. No finite ramification has been
assumed. At \(v\notin S\), its restriction to inertia
is a homomorphism because the coefficient action there
is trivial. Wild inertia and the prime-to-\(\ell\) tame
factors map trivially to the additive pro-\(\ell\) lattice.
For an \(\ell\)-primary tame generator \(\tau\), put
\(b=z(\tau)\). The actual tame relation and the cocycle
identity give
\[
                       F_vb=q b.                     \tag{D.72.9}
\]
The terms involving \(z(F_v)\) cancel because inertia
acts trivially on the coefficient. Equation (D.72.8)
and torsion-freeness imply \(b=0\). Thus \(z\) kills
every inertia group outside \(S\). The coefficient
action itself is unramified outside \(S\), so the
affine representation \((z,\rho_{L_j})\) kills their
closed normal closure. Consequently
\[
 Z^1(G_k,L_j)=Z^1(G_{k,S},L_j),\qquad j\ge3,             \tag{D.72.10}
\]
as actual continuous cocycle modules.

For each \(r\), the group
\(H^1(G_{k,S},L_j/\ell^rL_j)\) is finite
\cite[I, Corollary 4.15]{DE_MilneADT}. This includes
\(\ell=2\) and real places; only degree-one finiteness
is used, not a cohomological-dimension assertion.
The finite invariant groups satisfy the
Mittag--Leffler condition. Continuous comparison
\cite[Proposition 1.6 and Theorem 2.2]{DE_Jannsen1988}
therefore identifies integral \(H^1\) with the compact
inverse limit of these finite first-cohomology groups.
The multiplication-by-\(\ell\) coefficient sequence
injects its quotient by \(\ell\) into the finite
group \(H^1(G_{k,S},L_j/\ell L_j)\). Compact Nakayama
makes this integral \(H^1\) finitely generated over
\(\mathbb Z_\ell\).

Choose actual continuous representatives of finitely
many module generators, and add the coboundaries of
a lattice basis of \(L_j\). These finitely many
actual maps generate every one-cocycle. The map
from the finite parameter power of \(\mathbb Z_\ell\)
to their linear combinations is continuous in the
uniform topology. Hence
\[
 Z^1(G_k,L_j)\text{ is a finitely generated compact }
              \mathbb Z_\ell\text{-module},\qquad j\ge3.
                                                        \tag{D.72.11}
\]
This argument retains actual coboundaries and works
even if \(H^1\) has finite torsion.

The starting level two is essential for this proof.
In degree two, the proper-proper tensor terms may
have Frobenius eigenvalue \(q\). On a once-punctured
elliptic curve that entire central factor is
\(\bigwedge^2T_\ell E=\mathbb Z_\ell(1)\), whose
unrestricted Kummer group is not finitely generated.
Fixing its actual class-two cocycle removes this
uncontrolled variation. For that curve (D.72.7)
is in fact pure weight \(j\).

\paragraph{Polynomial families and actual coboundaries.}
We give the finite-parameter step to specify fully
why the nonseparated global \(H^2\) creates no gap.
Hall coordinates identify every \(T_n\) homeomorphically
with a finite power of \(\mathbb Z_\ell\). Products,
inverses and powers are rational polynomials with
bounded denominators at fixed \(n\), integral on the
actual coordinate lattice. Appending zero top
coordinates gives a continuous polynomial
set-section \(j_n:T_{n-1}\to T_n\).
The actual action of \(G_k\) on \(T_n\) is polynomial
in the coordinates of its input, with coefficients
given by fixed continuous functions of the Galois
element: substitute the images of the finitely
many free generators in the Hall multiplication
and power polynomials. In particular this is the
actual arithmetic action, with all its higher
coefficients and any ramification of the chosen
reference section retained.

For a finite free lattice \(L\), fix finitely many
continuous rational two-cochains \(\zeta_i\) and put
\[
 B=\left\{u\in\mathbb Z_\ell^a:
          {\textstyle\sum_i}u_i\zeta_i
                       \in dC^1_{\rm cts}(G_k,L)\right\}.
                                                        \tag{D.72.12}
\]
This is the submodule of coefficients giving
\emph{actual integral continuous coboundaries}.
It is a submodule of the Noetherian module
\(\mathbb Z_\ell^a\), hence finite free and closed:
a finite module basis maps a compact finite
power of \(\mathbb Z_\ell\) onto it. Choose one
integral continuous primitive for each basis
vector of \(B\). The basis-coordinate inverse
extends to a \(\mathbb Q_\ell\)-linear map on
\(\mathbb Q_\ell^a\), by a left inverse of its
basis matrix. Applying that fixed linear map
to a polynomial vector taking values in \(B\)
therefore supplies an integral continuous
primitive, polynomial in the parameters and
lying in the finite span of the chosen
primitives. No closure of \(dC^1\) is taken
anywhere in this construction.

At level two, \(\mathscr Z_2(\alpha)=\{\alpha\}\)
is already a finite-parameter family. Suppose
all exact level-\((n-1)\) cocycles are covered
by a family \(f_t\) on a closed set
\(X\subset\mathbb Z_\ell^a\), every coordinate
being a finite sum
\[
             \sum_i P_i(t)\varphi_i(g),
      \quad P_i\in\mathbb Q_\ell[t],\quad
             \varphi_i\in C(G_k,\mathbb Q_\ell).
                                                        \tag{D.72.13}
\]
The set-lifts \(\widehat f_t=j_n(f_t)\) have
central multiplication defects
\[
 \eta_t(g,h)=\widehat f_t(g)\,
     \sigma_0(g)\bigl(\widehat f_t(h)\bigr)\,
                           \widehat f_t(gh)^{-1}
                       \quad\text{in }L_n.             \tag{D.72.14}
\]
Expanding the finitely many Hall polynomials
and the fixed arithmetic coordinates gives
finitely many parameter monomials, each with
a fixed continuous two-cochain as coefficient.
This finite span is uniform over \(X\).
Absorbing the bounded rational denominators
into those cochains, write
\(\eta_t=\sum_i v_i(t)\zeta_i\) with
\(v_i\in\mathbb Z_\ell[t]\).

Exactly the parameters in the closed set
\(X'=X\cap v^{-1}(B)\) lift, with \(B\)
defined by (D.72.12) for the actual lattice
\(L_n\). Its chosen primitive cancels the
defect. Adding arbitrary linear combinations
of the finite generators (D.72.11) gives
every exact lift, since two lifts with the
same projection differ by a central actual
one-cocycle. The new parameter set is
\(X'\times\mathbb Z_\ell^b\), closed and compact.
Its Hall coordinates again have the form
(D.72.13), including after the primitive
choice and every previous lifting step.
This proves the claimed finite-parameter
assertion by induction. Joint continuity
on the compact parameter set and \(G_k\)
gives continuity into the uniform cochain
space, and proves (D.72.2) compact.

\paragraph{Conjugators and compatible inverse limits.}
For any section action on \(T_2\), a fixed
element first projects to \(V^{G_k}\), which
is zero by the absence of eigenvalue one
at a good Frobenius. It then belongs to
\(L_2^{G_k}=0\), by the same argument.
These assertions are integral, since the
lattices are torsion-free. Twisting by
\(\alpha\) does not change the action on
either graded factor. Thus
\[
                       T_2^{G_k,\alpha}=1.             \tag{D.72.15}
\]
Every section class above \([\alpha]\)
has an exact representative above \(\alpha\):
lift its class-two conjugator to \(T_n\).
A conjugator between exact representatives
has fixed class-two image and hence lies
in \(K_n=\ker(T_n\to T_2)\).
Therefore (D.72.3) is the quotient of
the compact Hausdorff space (D.72.2)
by the continuous action of the compact
group \(K_n\). It is compact Hausdorff,
and every transition map has closed image.

Fix an exact representative \(z_m\) of
a class at a finite stage \(m\ge2\).
Its exact lift space at level \(n\ge m\)
is a closed subset of (D.72.2). A class
lift can be made exact by lifting the
required conjugator from \(T_m\), so
individual class liftability makes each
of these compact spaces nonempty.
Every finite compatibility condition
can be met at its highest level and
then projected down. Compactness gives
an inverse system of actual compatible
cocycles. Since every finite quotient
of \(P\) is a nilpotent \(\ell\)-group,
\[
                         P=\varprojlim_n T_n.          \tag{D.72.16}
\]
Their limit is an actual continuous
\(P\)-cocycle. The uniform topology
into \(P\) is exactly the inverse-limit
topology on these exact fibres: each
finite quotient detects a uniform
neighborhood and factors through some
\(T_n\). This proves compactness of
the full exact fibre and, on quotienting
by \(\gamma_3P\), of its class fibre.
Nested nonempty compact conjugator
sets show that two full classes with
identical classes at every level are
already conjugate. Thus no separate
conjugacy ambiguity remains in taking
the inverse limit of section classes.

For a once-punctured elliptic curve,
the metabelian group also has the
needed integral coordinate property.
In rank two the discrete Alexander
module is
\(\mathbb Z[X^{\pm1},Y^{\pm1}]C\),
and the lower-central filtration is
given by powers of \((X-1,Y-1)\),
as used in D.68. Its nilpotent
quotients are torsion-free finitely
generated nilpotent groups. Their
pro-\(\ell\) completions are the
actual metabelian nilpotent quotients,
by the finite-quotient universal
property. The central grade \(j\)
is finite free and, rationally,
a quotient of \(\operatorname{Lie}_j(V)\).
The eigenvalue exclusion (D.72.8)
therefore still applies for \(j\ge3\),
and every preceding polynomial,
cocycle and compactness argument
applies without change. The limit
is the actual metabelian group:
every finite quotient is nilpotent,
and compactness gives surjectivity
onto the inverse limit. This proves
the asserted metabelian compatibility;
it does not assert generic singleton
metabelian fibres on these elliptic
curves.

\paragraph{Finite failure in the named applications.}
For a generic parameter of D.71,
take \(\alpha\) to be its actual
class-two cusp cocycle. The full
pro-\(\ell\) fibre is a singleton.
The intersection of the descending
compact images
\[
 I_{N,m}=
 \operatorname{im}\bigl(\mathscr S_N([\alpha])
                \longrightarrow\mathscr S_m([\alpha])\bigr),
                              \qquad N\ge m,           \tag{D.72.17}
\]
is therefore exactly the cusp class:
a member in every image lifts to
the actual full group by the
compatibility result just proved.
For an open cusp neighborhood \(W\),
if all \(I_{N,m}\setminus W\) were
nonempty, they would be nested
nonempty closed subsets of a
compact space and would have
nonempty intersection. This
contradicts the singleton
intersection. Equation (D.72.4)
and individual finite failure
follow.

The same proof applies to D.70's
Fermat parameters. Their full
pro-\(\ell\) abelian fibre is
already a singleton, so its
subfibre over the specified
class-two cusp is a singleton.
Compactness is now taken after
fixing that class-two cocycle.
No assertion that all elliptic
or Fermat abelian fibres are
compact is used.

Finally the actual full geometric
class-\(n\) quotient is the product
of its pro-prime class-\(n\)
quotients, since finite nilpotent
groups split as products of their
Sylow groups. Cochains, conjugators
and section classes are componentwise.
The product of the compact fibres
just proved is compact, and its
limit is the actual pronilpotent
quotient \(\prod_\ell P_\ell\).
For the product generic cusp
parameters in D.70 or D.71 this
limit fibre is a singleton.
The same compact nested-image
argument gives finite nilpotent
failure and uniform exclusion
outside cusp neighborhoods in
that product tower. This step
concerns the pronilpotent quotient;
the full-profinite cusp conclusions
remain precisely those of D.70
and D.71, using D.24/D.56.
\end{proof}


% D.73. Root, FGH and PBC independent full proof/source audits PASS.
\subsection{Class-two cusp rigidity with any number of punctures}

\begin{statement}{D.73}{Proved}{Generic class-two cusp fibres on every affine hyperbolic curve}
Let \(U=X-D\) be a smooth geometrically connected affine
hyperbolic curve over a number field \(k\). Assume only that
\(D\) contains one \(k\)-rational point \(d\), and choose a
nonzero rational tangent there. No condition is imposed on
the number of other geometric punctures or their fields of
definition. Use this actual cusp section \(s_0\) as reference,
and put
\[
 P_\ell=\pi_1(U_{\bar k})^{(\ell)},\qquad
 T_{2,\ell}=P_\ell/\gamma_3P_\ell,\qquad
 {\cal K}_{k,\ell}=H^1(G_k,\mathbb Z_\ell(1)).
                                                        \tag{D.73.1}
\]
The chosen cusp packet injects into
\(H^1(G_k,T_{2,\ell})\), at every prime \(\ell\), including
two. In the intrinsic topology on \({\cal K}_{k,\ell}\)
specified in D.70 there is a dense \(G_\delta\) of
infinitely ramified parameters whose actual full pro-\(\ell\)
section fibre over this class-two image is the singleton
chosen cusp class.

The product generic locus in
\(\prod_\ell{\cal K}_{k,\ell}\) has the same singleton
property for the actual full profinite class-two fibre,
using precisely D.24 and D.56. Every coefficient-prime
component may simultaneously be infinitely ramified.
The compactness theorem D.72 then gives finite nilpotent
failure for every other class at a fixed finite stage,
uniformly outside any neighborhood of the cusp in the
fixed-class-two fibre. The all-prime version of this
last finite-failure assertion concerns the actual
pronilpotent tower, exactly as in D.72.

This includes curves with one or two geometric cusps.
For a one-cusp curve of genus greater than one it does
not identify the entire zero-abelian class-two fibre.
For a two-cusp curve it makes no singleton abelian
fibre claim. The exact zero-abelian parametrization
of D.71 remains specific to the once-punctured elliptic
curve. No first finite obstruction is computed.
\end{statement}

\begin{proof}
Fix \(\ell\), write \(r=|D(\bar k)|\), \(P=P_\ell\),
\(T_2=T_{2,\ell}\), \(V=P^{\rm ab}\), and
\(L_2=\gamma_2P/\gamma_3P=\bigwedge^2V\).
All these integral identifications use the actual
free pro-\(\ell\) group, with no logarithmic
coordinates. Let \(c\) generate the chosen
peripheral subgroup \(I_c\), so that
\[
                   \sigma_0(g)c=c^{\chi_\ell(g)}.       \tag{D.73.2}
\]
Choose a finite set \(S\) containing the places above
\(\ell\), the archimedean places, the bad places of
the marked pair, and the places where the chosen
tangent is not integral and nonzero.

\paragraph{Injection of the chosen packet.}
At a good \(v\notin S\), the Frobenius eigenvalues
on \(V\) have complex absolute values \(N(v)^{1/2}\)
or \(N(v)\), as proved in D.72. Those on \(L_2\)
have absolute values among \(N(v),N(v)^{3/2},N(v)^2\).
Thus neither lattice has a nonzero local fixed
element, and restriction to one such place gives
the corresponding global assertion.

If \(r\ge2\), the image of \(c\) in the boundary
part of \(V\) is primitive, and \(V/I_c\) is
torsion-free. Its boundary/proper sequence is
the quotient of (D.72.6) by this one boundary
line. Its Frobenius eigenvalues still have
positive weights, so
\((V/I_c)^{G_{k_v}}=(V/I_c)^{G_k}=0\).
The integral coefficient exact sequence gives
an injection
\[
 H^1(G,I_c)\lhook\joinrel\longrightarrow H^1(G,V),
                 \quad G=G_k\text{ or }G_{k_v}.        \tag{D.73.3}
\]
Consequently two parameters in the chosen
packet cannot have equal class-two classes
unless they are equal.

If \(r=1\), hyperbolicity gives \(g=g(X)\ge1\)
and \(V=T_\ell J_X\). In an actual surface
basis \(a_1,b_1,\ldots,a_g,b_g\), the boundary
word is the product of the commutators, up
to inversion. Its degree-two class is
\[
                  \Omega=\sum_{i=1}^g[a_i,b_i]
                        \quad\text{in }\bigwedge^2V.  \tag{D.73.4}
\]
It is primitive, including at two, and its
line \(I=\mathbb Z_\ell\Omega\) is the actual
image of \(I_c\simeq\mathbb Z_\ell(1)\).
The quotient \(L_2/I\) is torsion-free and
has only positive Frobenius weight, hence
zero local and global fixed subgroup.
Therefore \(H^1(G,I)\to H^1(G,L_2)\) is
injective. The map from this latter group
to the zero-abelian fibre of \(H^1(G,T_2)\)
is injective too: a conjugator between
two central cocycles has fixed image in
\(V\), hence belongs to \(L_2\), where
its relation is the usual coboundary
relation. This proves the chosen cusp
injection globally and locally for \(r=1\).
It does not say \(I=L_2\) unless \(g=1\).

More generally, for any section action
on \(T_2\), a fixed element first has
zero image in \(V\) and then lies in
the zero fixed subgroup of \(L_2\).
Inner twisting changes neither graded
action. We have proved
\[
               T_2^{G_k,\alpha}
                 =T_2^{G_{k_v},\alpha}=1
                              \quad(v\notin S)         \tag{D.73.5}
\]
for every actual section action \(\alpha\).

\paragraph{Peripheral centralizers and class-two separation.}
Every boundary generator has a primitive
first nonzero lower-central class in degree
at most two: degree one if \(r\ge2\), and
(D.73.4) if \(r=1\). It therefore has no
\(\ell\)-th root in \(P\), by torsion-freeness
of \(V\) and of the second graded lattice.
The closed two-generated subgroup argument
of D.71, proved there from
\cite[Theorem 3.2]{DE_Lubotzky1982}, applies
to every free finite-rank pro-\(\ell\) group.
For \(b\ne0\) it gives
\[
       C_P(c^b)=N_P(c^{b\mathbb Z_\ell})
                 =\langle c\rangle,\qquad
       \gamma_3P\cap\langle c\rangle=1.                \tag{D.73.6}
\]
For the normalizer equality, conjugation
acts trivially on that first nonzero
graded class, forcing its unit exponent
on \(c^b\) to be one. This also proves
the equalities for any other cusp.

For \(r\ge3\), two distinct geometric
peripheral lines have zero intersection
in \(V\), by the boundary permutation
module modulo its diagonal relation,
as checked in D.70. Nonzero powers of
different cusps therefore cannot be
conjugate in \(T_2\). For \(r=1\) there
is no other geometric cusp.

The new case is \(r=2\), when hyperbolicity
forces \(g\ge1\). Choose actual surface
generators with boundary words \(c_1,c_2\)
and relation
\[
          \left(\prod_{i=1}^g[a_i,b_i]\right)c_1c_2=1.
                                                        \tag{D.73.7}
\]
The images of the two boundary words
in \(V\) are opposite primitive vectors.
In \(T_2\), their exact relation is
\[
                    c_2=c_1^{-1}\Omega^{-1},
       \qquad\Omega=\sum_i[a_i,b_i]\in L_2.            \tag{D.73.8}
\]
Here additive notation is used for the
central lattice when writing \(\Omega\).
Let \(\overline V=T_\ell J_X\).
The canonical linear map
\[
              \bigwedge^2V\longrightarrow
                              \bigwedge^2\overline V  \tag{D.73.9}
\]
sends every commutator \([h,c_1]\) to
zero, and sends \(\Omega\) to the
nonzero primitive symplectic bivector
of \(\overline V\). This is the
exterior-square map of lattices; it
is not the class-two proper fundamental
group quotient, which would kill the
surface relation.

Suppose \(c_1^a=h c_2^b h^{-1}\) in
\(T_2\), with \(a,b\in\mathbb Z_\ell\).
Abelianization gives \(a=-b\).
Use (D.73.8) and cancel this common
power of \(c_1\). The remaining
central equality says
\[
                        b\Omega\in[c_1,V].            \tag{D.73.10}
\]
Applying (D.73.9) yields \(b\overline\Omega=0\).
Torsion-freeness and primitivity force
\(b=0\), and then \(a=0\).
Thus distinct geometric cusp inertia
groups are separated by conjugacy
of their nonzero powers already in
class two, also in the two-cusp case.
All comparisons may be made after a
finite unramified splitting extension
at a good place; nonzero inertia and
the geometric calculation are unchanged.

\paragraph{The exact local comparison.}
Use the actual continuous Kummer
representatives on each fixed torsion
component of \({\cal K}_{k,\ell}\)
constructed in D.70. For a parameter
\(a\), put
\[
 \sigma_a(g)=\operatorname{Int}(c^{a(g)})\sigma_0(g),
                      \qquad H=\gamma_3P.              \tag{D.73.11}
\]
A full pro-\(\ell\) section with its
prescribed class-two class can be
conjugated to have exactly the chosen
class-two cocycle. Its difference
from the cusp section is an actual
\(H\)-cocycle. By (D.73.5), a
conjugator between two exact
representatives lies in \(H\).
Consequently the fibre is exactly
\[
                         H^1(G_k,H;\sigma_a).          \tag{D.73.12}
\]

At \(v\notin S\), write \(a_v\) for
the valuation coordinate of \(a\).
If \(a_v\ne0\), the prescribed
class-two inertia is nonzero.
Hoshi's actual ordinary pro-\(\ell\)
specialization result
\cite[Appendix A, A.3, Remark A.3.1 and Proposition A.6]
{DE_Hoshi2012}
says that the section's nonzero
specialized inertia is peripheral.
The class-two separation just proved
forces it to be the chosen geometric
cusp, even if another cusp only
splits after an unramified extension.
The centralizer and normalizer
calculation (D.73.6), together with
the cartesian specialization square,
puts the entire local section in the
actual chosen cusp decomposition
group, up to geometric conjugacy.
The packet injection (D.73.3), or
the one-cusp argument following
(D.73.4), identifies its Kummer
parameter with \(a|_{G_{k_v}}\).
After this cusp adjustment the
conjugator has fixed image in
\(T_2\), so (D.73.5) puts it in
\(H\). Thus every global cocycle
in (D.73.12) is an integral local
\(H\)-coboundary at such a place.

If \(a_v=0\), a nonzero specialized
inertia group would be peripheral
and would retain a nonzero
class-two image, contradicting
the prescribed image. The same
cartesian square therefore makes
the actual inertia image zero.
The reference tangent is unramified
at this good place. Hence the
actual difference cocycle satisfies
\[
 \begin{cases}
 z|_{G_{k_v}}\text{ is an }H\text{-coboundary},
                                             &a_v\ne0,\\
 z|_{I_v}=1,                                &a_v=0.
 \end{cases}                                             \tag{D.73.13}
\]
The first assertion is over \(k_v\)
itself: the geometric comparison
over a splitting extension has
identified the chosen rational
cusp, whose full local arithmetic
decomposition group is defined
over \(k_v\). No descent of an
arbitrary geometric conjugator
from a field extension is used.

\paragraph{Finite support, witnesses, and exact descent.}
We spell out why the same Baire
argument applies to \(H=\gamma_3P\).
Let \(Q_m=P/\Phi^mP\); these are
finite characteristic quotients
of exponent dividing \(\ell^m\),
and their kernels form a
neighborhood basis. Let \(H_m\)
be the image of \(H\). Define
\[
         S_m(a)=S\cup\{v:a_v\notin\ell^m\mathbb Z_\ell\}.
                                                        \tag{D.73.14}
\]
This is finite. Outside this set,
the action \(\sigma_a\) on \(Q_m\)
is unramified. Equation (D.73.13)
makes every actual \(H\)-cocycle
unramified there modulo \(H_m\):
in the nonzero case the integral
local coboundary has trivial
inertia after reduction. Therefore
all actual reductions occur in
the finite set
\[
             Z^1(G_{k,S_m(a)},H_m;\sigma_a).            \tag{D.73.15}
\]
Finiteness is the finite-image
Hermite--Minkowski argument used
in D.70; it applies to the
associated finite semidirect
representation, not just to
abelian cohomology. Both the
finite action and this finite
support are locally constant
at sufficiently deep precision
on each fixed torsion component.

Here the action used for the
witness is precisely
\(\rho_m:G_k\to\operatorname{Aut}(H_m)\);
freezing the action on \(Q_m\)
also freezes this induced action.
If a cocycle in (D.73.15) is
nontrivial on \(\ker\rho_m\), its affine
image contains an actual element
\((z(\delta),1)\) with
\(z(\delta)\ne1\). Chebotarev
\cite[V, Theorem 3.23]{PBC_MilneCFT}
supplies fresh good primes whose
Frobenius has this action-trivial
but cocycle-nontrivial image.
The perturbations of D.70,
using elements \(\alpha_v\) with
\(\operatorname{div}(\alpha_v)=h v\)
for the class number \(h\),
give a nonzero new \(a_v\) while
preserving all previous finite
data. Choose depth beyond
\(m\) and beyond the prescribed
neighborhood, and avoid the
finite support where the old
valuation has depth at most
\(M+v_\ell(h)\). Then adding
\(\ell^M\kappa(\alpha_v)\)
has exactly nonzero valuation
of depth \(M+v_\ell(h)\)
at this fresh place, without
cancellation. The fixed torsion
component, including at two,
does not change.

The first case of (D.73.13)
now makes that Frobenius value
zero modulo \(H_m\), since its
finite action is the identity.
This contradicts the chosen
finite affine witness. As in
D.70, the finite lists produce
dense open conditions at each
precision eliminating all such
exceptional cocycles. Their
countable intersection is a
dense \(G_\delta\) in the
complete Kummer parameter space;
intersect also the dense open
conditions supplying infinitely
many nonzero valuation primes.

For a parameter in this locus,
every actual \(H\)-cocycle
factors through the actual
compact image
\(\mathcal I_a\) of \(G_k\)
in \(\operatorname{Aut}(H)\),
as detected by these finite
characteristic quotients.
Choose one nonzero valuation
place. Its inertia image is
\[
 N=\{\operatorname{Int}(c^{a_v t})|_H:
                             t\in\mathbb Z_\ell\}.
                                                        \tag{D.73.16}
\]
It is normal in \(\mathcal I_a\)
by (D.73.2). By (D.73.6),
\(H^N=1\), for primitive or
nonprimitive \(a_v\).
Its restriction is a coboundary
by (D.73.13); change the
global representative to kill
this restriction. For \(n\in N\)
and \(g\in\mathcal I_a\), the
two cocycle expansions of
\(ng=g(g^{-1}ng)\) show that
the value at \(g\) is fixed
by every \(n\). It is therefore
one. This proves the singleton
in (D.73.12).

Finally take the product of
the residual loci over all
coefficient primes. The
topology and actual full cusp
existence are precisely those
proved in D.70. If an actual
full section has the given
class-two class, its full
pro-\(\ell\) image is that
fixed cusp image for every
\(\ell\). D.24/D.56 then
compares the two actual full
sections and makes them equal.
Apply D.72 to this fixed
class-two cusp to obtain the
finite-stage and compact
nested-image assertions.
There is no claim that the
entire zero-abelian fibre has
been computed outside D.71's
elliptic case, or that
primewise geometric quotients
are jointly faithful.
\end{proof}



% D.74. Root, FGH, and PBC independent full proof and primary-source audits passed.
\subsection{Integral Alexander cocycles at every finite ramification set}

\begin{statement}{D.74}{Proved}{Every finitely ramified Alexander first cohomology is an infinite countable product}
Let \(s=(a,b)\) be any finitely ramified abelian pair for the
tripod over \(\mathbb Q\), and let \(A_s=\Lambda C\), with
\(\Lambda=\mathbb Z_5[[u,v]]\), carry the actual action of D.35.
Then, as compact \(\mathbb Z_5\)-modules,
\[
 H^1_{\rm cts}(G_{\mathbb Q},A_s)
                  \simeq\prod_{\mathbb N}\mathbb Z_5.
                                                        \tag{D.74.1}
\]
The isomorphism is noncanonical. Thus every nonempty
geometrically pro-five metabelian section fibre above a
finitely ramified pair is a torsor under this compact group.
Every finitely ramified cusp pair has uncountably many
metabelian classes which are neither cuspidal nor induced
by rational points.

The conclusion concerns first cohomology and already nonempty
fibres. It does not assert existence above an arbitrary pair,
evaluate the remaining compact second-cohomology terms, or
compute a first higher-derived obstruction. In particular the
zero-rank twists of D.60--D.66 must be infinitely ramified.
\end{statement}

\begin{proof}
The proof constructs integral continuous cocycles before any
rank comparison. We distinguish the operator ring from its
actual module throughout.

\paragraph{A finite arithmetic presentation of positive deficiency.}
Choose a finite rational set \(S\) containing \(5,\infty\)
and the support of \(a,b\). Put
\[
 K=\mathbb Q(\zeta_5),\quad G=G_{K,S},\quad P=G^{(5)}.
\]
Let \(d=\dim_{\mathbb F_5}H^1(G,\mathbb F_5)\) and
\(h=\dim_{\mathbb F_5}H^2(G,\mathbb F_5)\).
These are finite, and the global Euler characteristic is
\[
                         1-d+h=-2,\qquad d-h=3.
                                                        \tag{D.74.2}
\]
Here \(K\) is totally imaginary of degree four and \(S\)
contains every place over five. This is the finite-module
global Euler characteristic in Milne ADT, Chapter I,
Theorem 5.1 \cite{DE_MilneADT}, with trivial
\(\mathbb F_5\) coefficients.
There is no real-place correction.

Write \(N=\ker(G\to P)\). As in D.64, an invariant
nonzero character \(N\to\mathbb F_5\) would give a
pro-five quotient of \(G\) with nontrivial image of \(N\).
Consequently \(H^1(N,\mathbb F_5)^P=0\), and the
Hochschild--Serre sequence gives
\[
 H^1(P,\mathbb F_5)=H^1(G,\mathbb F_5),\qquad
 H^2(P,\mathbb F_5)\lhook\joinrel\longrightarrow
                                     H^2(G,\mathbb F_5).
                                                        \tag{D.74.3}
\]
Thus \(P\) has a finite pro-five presentation with \(d\)
generators and \(r\le h\) relations. One can obtain it
directly from a minimal free presentation: its relation
coinvariants modulo fifth powers are dual to
\(H^2(P,\mathbb F_5)\); lifts of a finite basis normally
generate the kernel by the finite-five-group Nakayama
argument reproduced in D.64. In particular
\[
                              d-r\ge3.                 \tag{D.74.4}
\]
No assertion that \(P\) has cohomological dimension two
is used here.

The actual Alexander action factors through
\(\mathcal G=P\rtimes C_4\), a quotient of
\(G_{\mathbb Q,S}\). Indeed its finite coefficient
actions over \(K\) are five-groups, by the augmentation
and five-adic filtrations. The same holds for the finite
affine graphs of actual cocycles. Hence all first
cohomology in question factors through \(\mathcal G\).
The extension by \(C_4=\operatorname{Gal}(K/\mathbb Q)\)
has a complement by the coprime complement theorem,
or by compatible complements in its finite quotients.
Fix a generator \(c\) of such a complement. Its actual
cyclotomic multiplier is a primitive fourth root
\(\omega\in\mathbb Z_5^\times\).

\paragraph{The actual completed operator ring.}
Put \(\epsilon=1+5\), and let
\(\gamma\) act on \(\Lambda\) by the substitution
\[
 \varphi(t_i)=t_i^\epsilon,\qquad t_1=1+u,\ t_2=1+v.
\]
Use the completed group ring
\[
 R=\mathbb Z_5[[\,\mathbb Z_5^2\rtimes_\epsilon\mathbb Z_5\,]],
 \qquad z=\gamma-1.
                                                        \tag{D.74.5}
\]
It has the unique ordered expansion
\[
 R=\left\{\sum_{k\ge0}a_k(u,v)z^k:a_k\in\Lambda\right\},
 \qquad
 z a=\varphi(a)z+\varphi(a)-a.
                                                        \tag{D.74.6}
\]
This is the actual semidirect-product Iwasawa ring, as in
Venjakob, \emph{A noncommutative Weierstrass preparation
theorem}, Example 2.3 and its preceding construction
\cite{DE_Venjakob2002Skew}.
The additive expansion also follows from the ordered
coordinates of the three copies of \(\mathbb Z_5\).
The topology is the \((5,u,v,z)\)-adic topology.

Every \(p\in P\) acts on \(A_s\) as
\[
 \rho_s(p)=b_p\gamma^{e(p)},\qquad
 b_p\in\Lambda^\times,\quad b_p(0,0)=\chi_5(p)^2,\quad
 \chi_5(p)=\epsilon^{e(p)}.
                                                        \tag{D.74.7}
\]
Formula D.35.1 gives \(b_p\), including its actual Ihara
factor. Thus \(p\mapsto b_p\gamma^{e(p)}\) is a continuous
homomorphism \(P\to R^\times\), with image in
\(1+(5,u,v,z)R\). This is an equality of actual operators,
not an approximation by a weight.

The ring \(R\) acts on \(\Lambda\) by multiplication and
the indicated substitution. This completed action is
well-defined: \(\varphi-1\) increases the
\((5,u,v)\)-adic order by at least one, so the series
in (D.74.6) converges uniformly on \(\Lambda\).
The action is continuous. We prove below both its
faithfulness and the stronger infinite-rank property
needed in the argument.

We will also use that \(R\) is a left and right
Noetherian domain with a division ring of fractions \(D\).
Here is a filtration proof in this particular ring.
Give \(5\) weight two and \(u,v,z\) weight one.
The resulting complete filtration is equivalent to the
maximal-ideal topology. The relation
\[
 \varphi(u)-u=5u+15u^2+20u^3+15u^4+6u^5+u^6
\]
and its \(v\)-analogue increase weight by at least two.
Consequently all commutators of the three weight-one
symbols have weight strictly greater than two, and the
ordered expansion gives
\[
             \operatorname{gr}R
                   =\mathbb F_5[P_0,U_0,V_0,Z_0],
 \quad \deg P_0=2,\quad
 \deg U_0=\deg V_0=\deg Z_0=1.                         \tag{D.74.8}
\]
This is the weighted filtration, not the generally
noncommutative maximal-ideal associated graded ring.
The graded domain proves that \(R\) is a domain.
For a right ideal, lift a finite homogeneous generating
set of its associated graded ideal. Successive
subtraction raises the weight; completeness sums the
coefficients into finitely many elements of \(R\).
This proves finite generation of the original right
ideal. The identical argument works on the left.

For clarity, the Ore property follows without an
unproved flatness claim about \(\Lambda\). In a right
Noetherian domain two nonzero principal right ideals
cannot have zero intersection: if \(aR\cap bR=0\),
then \(bR,abR,a^2bR,\ldots\) form a direct sum in
\(R\), giving an infinite strictly ascending chain
of right ideals. Thus common nonzero right multiples
exist; likewise on the left. These Ore fractions give
the division ring \(D\), and localization is exact on
right modules by the common-denominator construction.
The general Noetherian assertion also follows from
Venjakob's Iwasawa structure theorem, Theorem 3.26
\cite{DE_Venjakob2001Structure},
whose hypotheses hold for this torsion-free uniform
pro-five group.

\paragraph{Fox cocycles with operator coefficients.}
Give \(R\) the left \(P\)-action in (D.74.7).
Evaluation on the chosen generators identifies its
continuous cocycles and coboundaries with
\[
 Z_R=\ker\{L:R^d\to R^r\},\qquad
 B_R=\{\,\mathbf r\,t:t\in R\,\},\quad
 \mathbf r=(\rho_s(p_i)-1)_i.                           \tag{D.74.9}
\]
Here \(L\) is the actual completed Fox relation matrix.
These are right \(R\)-modules. To verify continuity,
arbitrary values on free pro-five generators give
compatible finite affine five-group homomorphisms
at every finite ring quotient. The relator conditions
are exactly \(L\mathbf v=0\), and then the cocycle
descends through the closed normal relation subgroup.
Thus no abstract discontinuous cocycles enter (D.74.9).

By D.35 there is \(\delta\in P\), with
\(\chi_5(\delta)=1\), such that
\[
                 h=\rho_s(\delta)-1\in\Lambda
                                      \quad\hbox{is nonzero}.
                                                        \tag{D.74.10}
\]
Its existence uses the primitive actual Soul\'e element
of D.35, which cannot be canceled by the linear Kummer
multiplier. Its image in the finite-ramification
quotient still has the same nonzero operator.
Add \(\delta\) as a redundant presentation generator
and its defining relation. This increases \(d,r\)
equally; write \(i\) for this generator's index.
The \(i\)-th entry of \(\mathbf r\) is the
nonzero multiplication operator \(h\).

After right localization, the vector space
\[
             M_D=(Z_R/B_R)\otimes_R D
\]
has dimension at least \(d-r-1\ge2\). Indeed the
relation matrix has at most \(r\) independent rows,
and the nonzero coboundary column has rank one.
This is a statement about the operator coefficient
ring. We do not tensor the actual module
\(\Lambda\) with \(D\).

\paragraph{Tame semilinear descent and an invariant denominator.}
Conjugation by the actual operator \(\rho_s(c)\)
defines an order-four automorphism \(\alpha\) of \(R\).
For \(a\in\Lambda\) it is \(a(t_i^\omega)\);
on \(\gamma\) it is conjugation by the unit
multiplier \(b_c\), since the two substitutions commute.
It preserves the preceding topology and weighted
filtration. It acts on the leading symbols \(u,v\)
by \(\omega\) modulo five. It extends to \(D\).
Indeed \(\varphi(b_c)-b_c\) has weight at least
three, so \(\alpha(z)-z\) does also; the leading
symbols of \(5,z\) are unchanged.
Every nonidentity power of \(\alpha\) is outer on
\(D\): the valuation from (D.74.8) extends to Ore
fractions, and its associated graded division ring
is commutative. Therefore an inner automorphism acts
identically on leading symbols, whereas
\(\alpha^j\), \(1\le j\le3\), multiplies the leading
symbol of \(u\) by \(\overline\omega^j\ne1\).

The action on actual \(R\)-valued cocycles is
\[
          (c\cdot z)(p)=\alpha\bigl(z(c^{-1}pc)\bigr).
                                                        \tag{D.74.11}
\]
It preserves \(Z_R,B_R\), and gives \(M_D\) a
semilinear order-four action.

We use the following elementary independence fact.
Distinct automorphisms of a division ring, no two
differing by an inner automorphism, are linearly
independent as functions with coefficients on the
left. The same holds when they are tested only on
a subring that generates the division ring.
For a proof take a shortest nonzero relation.
Evaluate it at \(yx\), subtract a suitable left
multiple of its value at \(x\), and cancel one
term. Minimality then forces any two remaining
automorphisms to differ by conjugation by a
coefficient, a contradiction. Testing \(x,y\) in
the generating subring suffices, since the resulting
identity extends to all fractions.

Choose a nonzero class \(m\in M_D\) represented by
a cycle in \(Z_R\). There is \(t\in R\) for which
\(\sum_{j=0}^3c^j(m t)\ne0\). Otherwise, in one
nonzero coordinate of the right \(D\)-vector space
\(M_D\), this would be a nontrivial left-coefficient
relation among \(\alpha^j(t)\), contrary to the
preceding independence. Averaging the representing
cycle gives an exactly \(C_4\)-invariant
\[
                 \mathbf w\in Z_R,\qquad
                    \mathbf w\notin\mathbf r D.       \tag{D.74.12}
\]
Division by four is allowed in \(\mathbb Z_5\).

We may clear a denominator without losing this exact
invariance. Let
\[
 J=\{t\in R:h^{-1}w_i t\in R\}.
 \]
It is a nonzero right ideal by the Ore condition.
The intersection \(J_*=\bigcap_{j=0}^3\alpha^j(J)\)
is again nonzero, by successive common right multiples,
and is stable under \(\alpha\). Choose \(0\ne j_*\in J_*\).
The same independence argument, now in \(D\), shows
that for some \(t\in R\) the element
\[
                    d_*=\sum_{j=0}^3\alpha^j(j_*t)
 \]
is nonzero. It belongs to \(J_*\) and is
\(\alpha\)-invariant. Replace \(\mathbf w\) by
\(\mathbf w d_*\). It is still an exactly invariant
cycle with nonzero class in \(M_D\), and now
\[
             w_i=h u_*\quad\hbox{for some }u_*\in R.
                                                        \tag{D.74.13}
\]
Since it is not a division-ring coboundary, for some
generator index \(j\) the actual operator
\[
                  T=r_j u_*-w_j\in R
                               \quad\hbox{is nonzero}. \tag{D.74.14}
\]
This supplies the denominator control needed below;
an arbitrary Reynolds average was not assumed nonzero.

\paragraph{Every nonzero operator has infinite rank on the actual tame fixed module.}
Write \(\rho_s(c)=b_c\varphi_\omega\) on \(\Lambda\).
Its constant term is \(\omega^2\), and it has
order four. The integral eigenprojector
\[
       e=\frac14\sum_{j=0}^3\omega^{-2j}\rho_s(c)^j(1)
                                                        \tag{D.74.15}
\]
has constant term one, hence is a unit, and
\(\rho_s(c)e=\omega^2e\).
Also put
\[
 U=\frac14\sum_{j=0}^3\omega^{-j}\varphi_\omega^j(u),
 \qquad
 V=\frac14\sum_{j=0}^3\omega^{-j}\varphi_\omega^j(v).
                                                        \tag{D.74.16}
\]
They are integral formal coordinates with leading
terms \(u,v\), and satisfy
\(\varphi_\omega(U)=\omega U\),
\(\varphi_\omega(V)=\omega V\).
Thus \(eU^n\in\Lambda^{\rho_s(c)}\) whenever
\(n\equiv2\pmod4\).

Let \(0\ne T\in R\). The operator \(Te\) is nonzero,
since \(R\) is a domain. Expand it uniquely in
\(U,V,z\), and let \(a\) be its least total
\((U,V)\)-degree. Its degree-\(a\) coefficient has
the form
\[
                  \sum_{i+j=a}U^iV^j f_{ij}(z),
          \qquad f_{ij}\in\mathbb Z_5[[z]],
 \]
with at least one nonzero \(f_{ij}\). The substitution
\(\varphi\) has linear part
\(U\mapsto\epsilon U,\ V\mapsto\epsilon V\).
Consequently the degree-\(n+a\) term of
\((Te)(U^n)\) is exactly
\[
       \sum_{i+j=a}f_{ij}(\epsilon^n-1)U^{n+i}V^j.
                                                        \tag{D.74.17}
\]
This identity includes the infinitely many completed
\(z\)-coefficients: they converge because
\(\epsilon^n-1\in5\mathbb Z_5\).

A nonzero series in \(\mathbb Z_5[[z]]\) has only
finitely many zeros in \(5\mathbb Z_5\), by applying
the elementary Strassmann theorem to \(f(5z)\).
Here this elementary form needs no hypothesis about
a rational formal series: \(f(5z)\) has coefficients
tending to zero. If \(N\) is the largest index at
which their maximal absolute value is attained,
division by \(z-a\), for a zero \(a\in\mathbb Z_5\),
is obtained by the convergent coefficient sums
\(h_i=\sum_{j>i}g_j a^{j-i-1}\). The resulting
strictly convergent quotient has largest maximal
index \(N-1\). Repeating proves that there are
at most \(N\) zeros, including multiplicity.
The values \(\epsilon^n-1\) are distinct. Thus for
all but finitely many positive \(n\equiv2\pmod4\),
(D.74.17) is nonzero. The corresponding actual
outputs \(T(eU^n)\) have different
\((U,V)\)-orders \(n+a\), and are linearly
independent over \(\mathbb Q_5\). In particular
\[
 \operatorname{rank}_{\mathbb Z_5}
          T\bigl(\Lambda^{\rho_s(c)}\bigr)=\infty.
                                                        \tag{D.74.18}
\]
The same argument without restricting \(n\) proves
faithfulness of the whole actual completed action.
There is no use of a purely formal operator field
acting faithfully by assumption.

\paragraph{Integral cocycles and their nonzero descended classes.}
For \(f\in\Lambda^{\rho_s(c)}\), evaluate every value
of the invariant \(R\)-valued cocycle \(\mathbf w\)
on \(f\). This gives an actual integral continuous
\(A_s\)-valued cocycle \(z_f\) on \(P\).
Equation (D.74.11) and \(c f=f\) make it exactly
\(C_4\)-equivariant. Extending it by zero on the
chosen complement gives a cocycle on
\(\mathcal G=P\rtimes C_4\), hence on
\(G_{\mathbb Q,S}\).

If its class is a coboundary, its values on \(P\)
have the form \(r_j m\) for one \(m\in\Lambda\).
The special component (D.74.13) gives
\[
                   h m=w_i f=h(u_*f).
 \]
Multiplication by the nonzero element \(h\) is
injective on the integral domain \(\Lambda\), so
\(m=u_*f\). Its \(j\)-th component now implies
\(T(f)=0\), where \(T\) is (D.74.14).
Therefore the kernel of the class map
\[
 \Lambda^{\rho_s(c)}\longrightarrow
                  H^1(G_{\mathbb Q,S},A_s),\quad
                       f\longmapsto[z_f]
 \]
is contained in the kernel of \(T\). By (D.74.18)
the image, and hence the cohomology group, has
infinite integral rank. Every step took place in
integral actual cochains.

Finally D.47 identifies this first cohomology with
the unrestricted group: outside \(S\) every integral
cocycle has zero inertia, by the tame positive-weight
argument already proved there. D.35 proves
torsion-freeness, and D.47 proves compactness and
that it is a product of at most countably many
copies of \(\mathbb Z_5\). Infinite rank therefore
gives precisely (D.74.1).

For a fixed pair each of the three cusp packets
contributes at most one metabelian class, by its
primitive abelian boundary coordinate and the
packet formulas of D.35. Rational points form a
countable set. A nonempty torsor under
\(\prod_{\mathbb N}\mathbb Z_5\) is uncountable,
so removing these classes still leaves uncountably
many. A finitely ramified cusp pair has an actual
cusp section, so its fibre is nonempty. No assertion
of existence is made for other pairs.
\end{proof}


% D.75. Root and FGH independent complete proof/source/certificate audits passed.
\subsection{Explicit free factors in the retained compact obstruction}

\begin{statement}{D.75}{Proved}{Local Kummer divisibility produces a compact global obstruction factor}
Use the actual Alexander action of D.35 and the compact kernel
\(\mathcal K_s\) in D.63. Let
\[
 0\ne a=\sum_p a_p\kappa(p)\in H^1(G_{\mathbb Q},\mathbb Z_5(1)),
 \qquad s=(a,0).
\]
The coefficients \(a_p\) tend to zero; finite support is not required.
Put
\[
 \begin{split}
 r&=\min_p v_5(a_p),\\
 N&=\sup\{n\ge0:a|_{G_{\mathbb Q_5}}\in
                   5^nH^1(G_{\mathbb Q_5},\mathbb Z_5(1))\},
 \qquad v_5(0)=+\infty .
 \end{split}                                               \tag{D.75.1}
\]
Thus \(r<\infty\), \(N\ge r\), and \(N=\infty\) means that the
local Kummer class is zero. There is a continuous surjection
\[
 \mathcal K_s\twoheadrightarrow
 \begin{cases}
       \mathbb Z_5^{\,N-r},&N<\infty,\\
       \displaystyle\prod_{\mathbb N}\mathbb Z_5,&N=\infty .
 \end{cases}                                               \tag{D.75.2}
\]
It splits continuously as a map of compact \(\mathbb Z_5\)-modules,
noncanonically. In particular these are actual direct factors of the
compact term, not merely classes at one finite coefficient level.

Both integers in (D.75.1) are explicit. With the usual convergent
five-adic logarithm set
\[
 \lambda_p=\frac{\log(p^4)}{4\log 6}\quad(p\ne5),\qquad
 b_5(a)=\sum_{p\ne5}a_p\lambda_p.
\]
Then
\[
                    N=\min\{v_5(a_5),v_5(b_5(a))\}.         \tag{D.75.3}
\]
The displayed sum converges, independently of order.
For example \(s=(\kappa(101),0)\) gives a split
\(\mathbb Z_5\) factor in \(\mathcal K_s\). More generally a prime
\(q\) with \(v_5(q-1)=N+1\), \(N\ge1\), gives a split
\(\mathbb Z_5^N\) factor.

There is also a finite-support example with an infinite free factor.
Put
\[
 c=\frac{\log 16}{\log 81}\in\mathbb Z_5^\times,\qquad
 a_*=\kappa(2)-c\kappa(3).
                                                               \tag{D.75.4}
\]
Here \(c\equiv3\pmod5\), \(r=0\), and \(N=\infty\).
Consequently \(\mathcal K_{(a_*,0)}\) has a continuous direct
factor \(\prod_{\mathbb N}\mathbb Z_5\). Both pairs in these
examples are cusp pairs and have actual sections.
The same assertions hold after permuting the three punctures.

This evaluates a free factor of the compact term in Task D.5.
The complementary compact factor and the general first
finite-derived obstruction are not determined here.
In particular the compact term cannot in general be discarded:
the zero calculation of D.64 does not extend to all finite
ramification sets.
\end{statement}

\begin{proof}
\paragraph{An actual one-variable quotient and its measure model.}
Write \(t=1+u\). The ideal \(vA_s\) is stable. In D.35's
normalization, the integral beta series is one when its second
argument is one. Thus on the actual quotient
\(M_a=A_s/vA_s\simeq\mathbb Z_5[[u]]\) the action is
\[
 g f(t)=t^{a(g)}\chi(g)S_{\chi(g)}(t)f(t^{\chi(g)}),
 \qquad S_\alpha(t)=\frac{t^\alpha-1}{t-1}.
                                                               \tag{D.75.5}
\]
This uses the exact Fox multiplier, including its factor
\(S_{\chi}(1)=\chi\), not a leading-weight approximation.
It follows directly from D.35 and D.39.
Multiplication by \(u\) is an isomorphism of compact modules
\[
 M_a(-1)\xrightarrow{\ \sim\ }I,\qquad
 I=\ker\{\mathbb Z_5[[\mathbb Z_5]]\longrightarrow\mathbb Z_5\},
 \qquad f(t)\longmapsto(t-1)f(t).                            \tag{D.75.6}
\]
On the right the action is the actual affine pushforward
\[
                  \nu\longmapsto t^{a(g)}\phi_{\chi(g)}(\nu),
 \qquad \phi_\alpha(t^z)=t^{\alpha z}.
                                                               \tag{D.75.7}
\]
Indeed \((t-1)S_\chi(t)=t^\chi-1\).
An integral change of Kummer cocycle by a coboundary corresponds
to conjugating this action by the translation \(t^d\), \(d\in\mathbb Z_5\).
Translation preserves \(I\).

Let \(D_a=M_a^D=\operatorname{Hom}_{\rm cts}(M_a,\mathbb Q_5/\mathbb Z_5)(1)\).
The dual of (D.75.6) identifies it with
\[
 D_a\simeq
 \frac{\operatorname{Cont}(\mathbb Z_5,\mathbb Q_5/\mathbb Z_5)}
                         {\mathbb Q_5/\mathbb Z_5}.
                                                               \tag{D.75.8}
\]
The denominator consists of constant functions. To verify this
identification integrally, a continuous linear functional on the
completed group ring is a locally constant function on \(\mathbb Z_5\).
The augmentation splits by the measure at zero, so every continuous
functional on \(I\) extends continuously to that ring, and two extensions
differ by a constant. Each functional has finite five-power image
and factors through a finite group ring. The action on (D.75.8)
is the dual of the affine action (D.75.7). Its invariant classes
are therefore exactly the functions whose change under each
affine transformation is constant.

\paragraph{The exact global and local affine images.}
We prove this step for either \(k=\mathbb Q\) or \(k=\mathbb Q_5\).
The cyclotomic character maps \(G_k\) onto
\(\Gamma=\mathbb Z_5^\times\). For a nonzero class
\(b\in H^1(G_k,\mathbb Z_5(1))\) let \(d\) be its largest
divisibility exponent, and write \(b=5^d b_0\) in cohomology
with \(b_0\) primitive. After an integral coboundary change,
the actual cocycle is \(5^d\) times a cocycle for \(b_0\).

If \(\Delta_k=\ker\chi\), then
\[
                       b_0(\Delta_k)=\mathbb Z_5.          \tag{D.75.9}
\]
Indeed the Kummer groups here are torsion-free, so primitiveness
means that the image in \(H^1(G_k,\mathbb F_5(1))\) is nonzero.
The kernel of restriction to \(\Delta_k\) is
\(H^1(\Gamma,\mathbb F_5(1))=0\). For completeness, the
order-four subgroup of \(\Gamma\) acts by a nontrivial scalar
on \(\mathbb F_5(1)\); averaging and the procyclic resolution
of \(1+5\mathbb Z_5\) kill this first cohomology. Equivalently,
the cocycle equation for the commuting element with
cyclotomic value two makes every cocycle principal.
Thus \(b_0|_{\Delta_k}\) is nonzero modulo five. Its image
is a compact additive subgroup of \(\mathbb Z_5\), so is all
of \(\mathbb Z_5\). This proves (D.75.9).

It follows that the closed affine image, after the indicated
integral translation, is exactly
\[
                  5^d\mathbb Z_5\rtimes\mathbb Z_5^\times.
                                                               \tag{D.75.10}
\]
It contains all the displayed translations by (D.75.9).
For every multiplier its translation coordinate is in
\(5^d\mathbb Z_5\), and can be removed using one of those
translations. If the class is zero, its image after an integral
translation is instead exactly \(\mathbb Z_5^\times\), with
zero translation part. This is the \(d=\infty\) convention.
Thus (D.75.10) applies globally with \(d=r\) and locally with
\(d=N\). No finite-level approximation to the affine image
has been substituted here.

\paragraph{All invariant functions modulo constants.}
For finite \(d\), suppose \([f]\) in (D.75.8) is fixed by
the group in (D.75.10). For \(h\in5^d\mathbb Z_5\), write
\[
                         f(z+h)-f(z)=c(h).
 \]
The function \(c\) is an additive continuous character.
Pure multiplication by two belongs to (D.75.10), and
\(f(2z)-f(z)\) is constant; evaluating at zero shows that
this constant is zero. Consequently
\[
                         c(2h)=c(h)=2c(h),\qquad c(h)=0.
 \]
Thus \(f\) is genuinely periodic modulo \(5^d\), not just
periodic modulo constants. Similarly each pure unit multiplier
fixes \(f\) exactly, again by evaluating at zero.
Its values are precisely arbitrary values on the \(d+1\)
unit orbits of \(\mathbb Z/5^d\): the valuation shells
\(0,\ldots,d-1\) and the zero orbit. Modulo constants,
\[
 D_a^{G_k}\simeq(\mathbb Q_5/\mathbb Z_5)^d.
                                                               \tag{D.75.11}
\]
A concrete basis consists of the classes of the characteristic
functions of \(5^j\mathbb Z_5\), \(1\le j\le d\), with arbitrary
coefficients in \(\mathbb Q_5/\mathbb Z_5\).

When \(d=\infty\), invariance means that \(f\) is constant
on every valuation shell and at zero. Continuity into the
discrete target says that the shell values eventually equal
the value at zero. Hence
\[
 D_a^{G_k}\simeq
                   \bigoplus_{j\ge1}\mathbb Q_5/\mathbb Z_5.
                                                               \tag{D.75.12}
\]
The same characteristic functions form a basis, now with
finite support. These calculations include every invariant
class in (D.75.8), not only the selected ball functions.

We also identify the inclusion of global in local invariants.
Normalize the global cocycle to be divisible by \(5^r\).
A further integral translation \(d_0\) normalizes the local
cocycle to be divisible by \(5^N\), or to be zero when
\(N=\infty\). This translation satisfies \(d_0\in5^r\mathbb Z_5\):
evaluate the coboundary equation on an element with
\(\chi=2\). Both cocycle values being compared lie in
\(5^r\mathbb Z_5\), and \(\chi-1=1\).
Translation by \(d_0\) therefore leaves the first \(r\)
ball functions unchanged. Thus the inclusion identifies
the first \(r\) basis vectors in (D.75.11)--(D.75.12), and
\[
 \frac{D_a^{G_{\mathbb Q_5}}}{D_a^{G_{\mathbb Q}}}
 \simeq
 \begin{cases}
       (\mathbb Q_5/\mathbb Z_5)^{N-r},&N<\infty,\\
       \displaystyle\bigoplus_{\mathbb N}\mathbb Q_5/\mathbb Z_5,
                                                    &N=\infty .
 \end{cases}                                               \tag{D.75.13}
\]
In particular a global-invariant class cannot conceal a
nonzero ball function in this quotient. The preceding
translation argument explains the possible additive slope,
and the tame multiplier two forces that slope to vanish.

\paragraph{Passage to the actual compact global term.}
Put \(D_s=A_s^D\), a discrete union of the finite Cartier
duals used in D.63. The stable quotient \(A_s\twoheadrightarrow M_a\)
gives a genuine inclusion \(D_a\hookrightarrow D_s\).
Moreover
\[
 D_a^{G_{\mathbb Q_5}}\cap D_s^{G_{\mathbb Q}}
                         =D_a^{G_{\mathbb Q}}.
 \]
Consequently (D.75.13) injects into
\[
                      D_s^{G_{\mathbb Q_5}}/D_s^{G_{\mathbb Q}}.
                                                               \tag{D.75.14}
\]
Local duality and finite-level inverse limits identify the
Pontryagin dual of (D.75.14) with
\[
 \ker\{H^2_{\rm cts}(G_{\mathbb Q_5},A_s)
                      \longrightarrow\mathcal T_s\},
 \]
where \(\mathcal T_s\) is D.63's global-invariant dual.
Here invariants commute with the increasing union of finite
discrete duals, and all local finite cohomology groups are
finite. These are exactly the local duality and inverse-limit
identifications in D.63, based on Milne ADT I.2 and I.4
\cite{DE_MilneADT}. In particular no finite-level class
is being assumed to survive an inverse limit.

D.63 gives a continuous surjection from \(\mathcal K_s\)
onto this kernel. Dualizing the genuine inclusion (D.75.13)
and composing proves (D.75.2).
The splitting assertion also holds topologically. The discrete
dual of its target is a direct sum of copies of
\(\mathbb Q_5/\mathbb Z_5\), hence is divisible and injective
as a \(\mathbb Z_5\)-module. For instance the extension
criterion for the ideals \(5^n\mathbb Z_5\) is exactly
divisibility by \(5^n\), and the usual maximal-extension
argument extends a map from any submodule. Thus its inclusion
in \(\mathcal K_s^D\) has a retraction. Pontryagin duality
turns that retraction into a continuous section of (D.75.2).
No splitting of the nonseparated cohomology quotient is asserted.

\paragraph{Exact arithmetic parameters and examples.}
The completed Kummer group of \(\mathbb Q_5\) has the basis
\(\kappa(5),\kappa(6)\). Indeed
\[
 \mathbb Q_5^\times=5^{\mathbb Z}\times\mu_4\times(1+5\mathbb Z_5),
 \]
the prime-to-five factor disappears in the pro-five completion,
and logarithm identifies \(1+5\mathbb Z_5\) with \(5\mathbb Z_5\);
\(\log6\) has valuation one. For \(p\ne5\), its local
Kummer coordinate is exactly \(\lambda_p\) above, since
raising to the fourth power removes its Teichmuller part.
All \(\lambda_p\) are integral and \(a_p\to0\), so the
sum defining \(b_5(a)\) converges independently of order.
This proves (D.75.3).

For \(q\equiv1\pmod5\), \(v_5(\log q)=v_5(q-1)\);
this follows immediately from the logarithm series, whose
first term has strictly smaller valuation than every later
term. Thus \(v_5(q-1)=N+1\) gives exactly the stated local
divisibility \(N\). In particular \(101-1=4\cdot25\).

Finally \(\log16\) and \(\log81\) both have valuation one,
since \(16-1=15\) and \(81-1=80\). Their ratio is a
five-adic unit, and its residue is \(3/16=3\) in \(\mathbb F_5\).
The local unit coordinate of (D.75.4) is exactly zero,
while its coefficient at the prime two is one. Thus
\(N=\infty\) and \(r=0\), as claimed.
These are completed tangential parameters in the actual
zero-cusp packet, so their section fibres are nonempty
by D.35. The transformations permuting the punctures
transport the same proof and parameters to the other packets.

The executed certificate \texttt{code/de\_compact\_axis.py}
checks the affine unit-orbit decomposition and the independence
of all displayed ball functions at depths zero through five.
For the latter it evaluates their translation-one defects at
\(z=5^i-1\); the matrix has entries \(1\) for \(j\le i\)
and zero otherwise, hence determinant one. At \(z=1\)
all defects are zero, excluding an unnoticed constant defect.
It also computes (D.75.4)'s ratio modulo \(5^{20}\) with a
proved logarithm-tail bound and checks sample primes of exact
depths one through five. The output and JSON are retained.
The continuous duality and infinite-factor assertion follow
from the proof above, not from extrapolating these finite checks.
\end{proof}


% D.76. Root and FGH independent full proof audits passed.
\subsection{The joint contribution of all three peripheral axes}

\begin{statement}{D.76}{Proved}{Three explicit compact factors for every abelian pair}
Let \(s=(a,b)\) be any abelian pair of the tripod over \(\mathbb Q\),
with arbitrary ramification, and let \(\mathcal K_s\) be the compact
term of D.63. For a completed Kummer class \(c\) define
\[
 \delta(c)=
 \begin{cases}
 0,&c=0,\\
 N(c)-r(c),&c\ne0,\quad N(c)<\infty,\\
 \infty,&c\ne0,\quad N(c)=\infty,
 \end{cases}                                               \tag{D.76.1}
\]
where \(r(c)\) and \(N(c)\) are the global and local-at-five
divisibility exponents of D.75. Put
\[
                     d(s)=\delta(a)+\delta(b)+\delta(a-b).
                                                               \tag{D.76.2}
\]
There is a continuous split surjection of compact modules
\[
 \mathcal K_s\twoheadrightarrow
 \begin{cases}
       \mathbb Z_5^{\,d(s)},&d(s)<\infty,\\
       \displaystyle\prod_{\mathbb N}\mathbb Z_5,&d(s)=\infty.
 \end{cases}                                               \tag{D.76.3}
\]
Thus the three contributions add: they are not merely three
possibly overlapping lower bounds.
The divisibility exponents are computed by the convergent
logarithm formula (D.75.3).

In particular \(s=(\kappa(101),0)\) has a split
\(\mathbb Z_5^2\) factor, and a cusp parameter of global
depth \(r\) and finite local depth \(N\) has a split
\(\mathbb Z_5^{2(N-r)}\) factor. The explicit logarithm
cancellation of D.75 still gives a countable product.
This theorem applies even when the metabelian section fibre
above \(s\) is empty. It asserts neither existence of such
a section nor a calculation of the complementary compact factor.
\end{statement}

\begin{proof}
\paragraph{All three are actual arithmetic quotients.}
Write \(\Lambda=\mathbb Z_5[[u,v]]\), \(t_1=1+u\),
\(t_2=1+v\), and
\[
 J_1=(v),\qquad J_2=(u),\qquad J_3=(t_1t_2-1).
 \]
Each is stable under the actual Alexander action.
The first quotient of \(A_s=\Lambda C\) is exactly
D.75's \(M_a\), for arbitrary \(b\): the factor \(t_2^{b(g)}\)
becomes one. The second quotient is \(M_b\).

For the third put \(t_1=t\), \(t_2=t^{-1}\).
The actual reflection identity of D.39 gives
\[
 \mathcal B_g(t,t^{-1})
             =\frac{\chi\,t^{(\chi-1)/2}}{S_\chi(t)},
 \qquad
 S_\chi(t^{-1})=t^{1-\chi}S_\chi(t).
 \]
Consequently the full Fox multiplier on that quotient is
\[
              \chi S_\chi(t)\,
                   t^{\,a(g)-b(g)+(1-\chi(g))/2}.       \tag{D.76.4}
\]
Since \(1/2\in\mathbb Z_5\), the extra exponent is an
integral coboundary and is removed by an integral translation.
Thus the third quotient is precisely \(M_{a-b}\), up to
the indicated actual gauge.

Let \(D_i\) be the Cartier duals of these three quotients,
viewed as submodules of \(D_s=A_s^D\).
The one-axis argument of D.75 therefore computes
\[
 Q_i:=D_i^{G_{\mathbb Q_5}}/D_i^{G_{\mathbb Q}}
 \simeq
 \begin{cases}
  (\mathbb Q_5/\mathbb Z_5)^{\delta(c_i)},&
                         \delta(c_i)<\infty,\\
  \displaystyle\bigoplus_{\mathbb N}\mathbb Q_5/\mathbb Z_5,&
                         \delta(c_i)=\infty,
 \end{cases}
 \quad(c_1,c_2,c_3)=(a,b,a-b).                           \tag{D.76.5}
\]
When \(c_i=0\), its global and local affine images are
the same unit group after one integral gauge, so \(Q_i=0\).
It remains to show that these quotients inject jointly.

\paragraph{The precise relation among the three axis duals.}
Consider the compact map
\[
 \Lambda\longrightarrow
       \Lambda/J_1\oplus\Lambda/J_2\oplus\Lambda/J_3.
                                                               \tag{D.76.6}
\]
Use \(u\) as coordinate on the first and third factors
and \(v\) on the second. A triple of power series
\((F(u),G(v),H(u))\) lies in its image exactly when
\[
 F(0)=G(0)=H(0),\qquad
                      H_1=F_1-G_1.                    \tag{D.76.7}
\]
Here the subscript denotes the linear coefficient.
Necessity follows from \(v=-u/(1+u)\) on the third axis.
For sufficiency, start with \(F(u)+G(v)-F(0)\).
Its discrepancy on the diagonal has zero constant and
linear coefficients. On that diagonal
\[
                         uv=-\frac{u^2}{1+u}.
 \]
Adding \(uv\,L(u)\) therefore corrects any such discrepancy
with an integral power series \(L(u)\).
This proves (D.76.7), and the cokernel of (D.76.6) is
exactly a free \(\mathbb Z_5\)-module of rank three.
All its defining functionals use only constant and linear
coefficients on each axis.

Duality now shows that the kernel of
\[
                       D_1\oplus D_2\oplus D_3
                             \longrightarrow D_s             \tag{D.76.8}
\]
has, in each component, a functional involving only the
constant and linear coefficients of the corresponding
one-variable module. Under the measure model of D.75,
such a functional is represented by a polynomial of degree
at most two in \(z\), modulo constant functions. Indeed
\[
       \frac{t^z-1}{t-1}
                =z+\binom z2(t-1)+\text{higher powers}.
                                                               \tag{D.76.9}
\]
Integral translation gauges replace a function by
\(f(z+d)-f(d)\), and so preserve this degree bound.
This remains true for the half-cyclotomic gauge in (D.76.4).
No independent-summands assertion has been assumed.

\paragraph{A polynomial defect cannot survive the tame unit.}
Suppose \(x_i\in D_i^{G_{\mathbb Q_5}}\) and that their
sum in \(D_s\) is globally invariant.
For each nonzero \(c_i\), make the compatible global and
local normalizations of D.75. Its global affine image is
still \(5^{r(c_i)}\mathbb Z_5\rtimes\mathbb Z_5^\times\),
and its local image has pure unit multipliers.
Represent \(x_i\) by a continuous function \(f_i\) modulo
constants, normalized by \(f_i(0)=0\).
Local invariance then gives
\[
                           f_i(2z)=f_i(z).             \tag{D.76.10}
\]
For \(\tau\in\ker\chi\), global invariance of the sum
puts the three component defects in the kernel of (D.76.8).
Writing \(h=c_i(\tau)\), and replacing \(\tau\) by its
inverse if necessary to match the dual-action convention,
we obtain
\[
 f_i(z+h)-f_i(z)=P_h(z)
                       =A_hz^2+B_hz+C_h,               \tag{D.76.11}
\]
with coefficients in \(\mathbb Q_5/\mathbb Z_5\).
A constant ambiguity is absorbed in \(C_h\).
The coefficients are unique: evaluation at \(0,1,2\)
has Vandermonde determinant two, a five-adic unit.

Equation (D.76.10) and the telescoping identity give
\[
              P_{2h}(2z)=P_h(z),\qquad
              P_{2h}(z)=P_h(z+h)+P_h(z).                \tag{D.76.12}
\]
The quadratic coefficient is therefore simultaneously
\(A_{2h}=2A_h\) and \(4A_{2h}=A_h\).
Thus \(7A_h=0\), whence \(A_h=0\).
The linear coefficients then give \(B_{2h}=2B_h\) and
\(2B_{2h}=B_h\), so \(3B_h=0\) and \(B_h=0\).
Finally the constants give \(C_{2h}=2C_h=C_h\),
so \(C_h=0\). Both seven and three are units on the
five-primary torsion target. Thus every defect (D.76.11)
is zero.

By the exact primitive-restriction argument of D.75,
the values \(h=c_i(\ker\chi)\) range over all of
\(5^{r(c_i)}\mathbb Z_5\). Hence \(f_i\) is fixed by
every global translation. It is also fixed by every pure
unit, from its local invariance and evaluation at zero.
The exact global affine image now makes \(x_i\) globally
invariant. If \(c_i=0\), this conclusion already holds
because the global and local affine images agree.
We have proved the genuine injection
\[
               Q_1\oplus Q_2\oplus Q_3
                  \lhook\joinrel\longrightarrow
                 D_s^{G_{\mathbb Q_5}}/D_s^{G_{\mathbb Q}}.
                                                               \tag{D.76.13}
\]

\paragraph{The compact direct factor.}
Compose the dual of (D.76.13) with D.63's continuous
surjection from \(\mathcal K_s\) onto the dual of its
right-hand side. Formula (D.76.5) gives exactly (D.76.3).
As in D.75, the direct sum on the left of (D.76.13) is
divisible, hence injective as a discrete \(\mathbb Z_5\)-module.
Its inclusion in \(\mathcal K_s^D\) has a retraction;
duality makes the compact surjection split continuously.
This determines the displayed factor only.
For a cusp pair \((a,0)\), the two equal nonzero Kummer
parameters \(a\) and \(a-0\) occupy different axes, and
the proved joint injection counts both of them.
This gives the final examples and explains why their ranks
double without an unsupported independence assumption.
The executed exact certificate \texttt{code/de\_three\_axis.py}
checks the integral jet interpolation, the translated degree-two
measure functional, and the successive coefficients seven, three
and one in (D.76.12). Its source, JSON and output are retained.
\end{proof}


% D.77. Root, FGH, and PBC full independent proof/source/certificate audits PASS.
\subsection{The entire contribution at five is peripheral}

\begin{statement}{D.77}{Proved}{Exact local and global invariant factors}
Let \(s=(a,b)\) be any completed abelian Kummer pair for the tripod over
\(\mathbb Q\), with arbitrary ramification, and use the actual Alexander
module \(A_s\) of D.35. Retain \(L_5,\mathcal T_s,\mathcal S_s,\mathcal K_s\)
from D.63. Set
\[
 (c_1,c_2,c_3)=(a,b,a-b).
\]
For \(c\ne0\), let \(r(c)\) and \(N(c)\) be the global and local-at-five
divisibility exponents of D.75. For \(c=0\), set \(r(c)=N(c)=\infty\).
Write
\[
 Z(d)=
 \begin{cases}
 \mathbb Z_5^d,&d<\infty,\\
 \displaystyle\prod_{j\ge1}\mathbb Z_5,&d=\infty .
 \end{cases}
 \qquad Z(0)=0.                                      \tag{D.77.1}
\]
There are isomorphisms of compact modules
\[
 L_5\simeq\prod_{i=1}^3 Z(N(c_i)),\qquad
 \mathcal T_s\simeq\prod_{i=1}^3 Z(r(c_i)).             \tag{D.77.2}
\]
They may be chosen so that \(L_5\to\mathcal T_s\) is, on each nonzero
axis, the projection to its first \(r(c_i)\) coordinates. On a zero
axis it is the identity between the two countable products.

Define \(\delta(c)\) as in D.76, with \(\delta(0)=0\), and put
\(d(s)=\delta(a)+\delta(b)+\delta(a-b)\). Then the \emph{whole} kernel, not just a split factor, is
\[
 \ker(L_5\longrightarrow\mathcal T_s)\simeq Z(d(s)),
 \qquad
 \mathcal K_s\simeq\mathcal S_s\times Z(d(s)).          \tag{D.77.3}
\]
The latter splitting is continuous and noncanonical.
Here \(\mathcal S_s\) is exactly the inverse limit of the finite
global localization kernels in (D.63.5). It is not asserted to vanish
or to be an elliptic-curve Tate--Shafarevich group.
D.79--D.80 subsequently evaluate \(\mathcal S_s\) on their
specified independent unit-pair families, including unrestricted
ramification in D.80. D.81 then removes local primitivity and
classifies every globally independent unit pair. D.82 allows
arbitrary five-valuations when two independent residual inertia
directions occur away from five. D.83 proves vanishing of the
whole compact kernel under its local-generator lattice
criterion. D.84 treats the primitive rank-one away lattice
with an independent valuation direction, and D.85 completes the
compact-kernel classification for every globally independent
residual pair. Globally dependent residual pairs remain open.
This makes no assertion that a section exists above \(s\),
and does not compute a first higher-derived obstruction.
\end{statement}

\begin{proof}
\paragraph{A primitive local Soul\'e character without a special-value input.}
Fix the embedding \(G_{\mathbb Q_5}\subset G_{\mathbb Q}\) and put
\[
 K=\mathbb Q_5(\zeta),\quad \zeta=\zeta_5,\quad
 \pi=\zeta-1,\quad
 \phi=-\zeta^2-\zeta^3=(1+\sqrt5)/2.
\]
The polynomial
\[
 \Phi_5(1+\pi)=\pi^4+5\pi^3+10\pi^2+10\pi+5
                                                               \tag{D.77.4}
\]
is Eisenstein. Normalize \(v_\pi(\pi)=1\), so \(v_\pi(5)=4\),
and let \(U_j=1+\pi^j\mathcal O_K\). The identity
\[
 \phi^4-1=3\phi+1=(5+3\sqrt5)/2
                                                               \tag{D.77.5}
\]
has valuation exactly two. Every fifth power in \(U_1\) lies in
\(U_5\): in \((1+x)^5-1\), the first four terms have valuations
at least \(4+jv_\pi(x)\ge5\), and the last has valuation at
least five. Also \(v_\pi(\zeta^e-1)=1\) for \(5\nmid e\).

These facts show that \(\phi,\zeta\) are independent in
\(K^\times/K^{\times5}\). Indeed, suppose
\(\phi^m\zeta^e=x^5\), where \(0\le m,e<5\).
The element \(x\) is a unit. Its fourth power is in \(U_1\),
since the residue field is \(\mathbb F_5\). Hence
\(\phi^{4m}\zeta^{4e}\in U_5\). If \(e\ne0\), its difference
from one has valuation one: the \(\phi^{4m}\) factor is in
\(U_2\), whereas \(\zeta^{4e}\) is not. If \(e=0\) and
\(m\ne0\), its difference from one has valuation two, by
(D.77.5) and \(5\nmid m\). Both are contradictions.

The raw Soul\'e normalization is (D.39.1), equivalently
\cite[(1.1)--(1.2)]{DE_NakamuraSakugawaWojtkowiak2017}.
D.35's exact first-level identity is
\[
 u_{3,1}=\prod_{j=1}^4(1-\zeta^j)^{j^2}
        =\phi^3\bigl((1-\zeta)^6\phi^2\bigr)^5.
                                                               \tag{D.77.6}
\]
The cyclotomic extension \(K(\mu_{5^\infty})/K\) has Galois group
\(\mathbb Z_5\), and its unique degree-five subextension is
\(K(\zeta_{25})\), with Kummer line generated by \(\zeta\).
The independence just proved implies
\(K(\phi^{1/5})\not\subset K(\mu_{5^\infty})\).
Consequently, if
\[
 \Delta_5=\ker\{\chi:G_{\mathbb Q_5}\to\mathbb Z_5^\times\},
\]
then \(\kappa_3\bmod5\) is nonzero on \(\Delta_5\).
In particular its actual integral image there contains a unit.
This proves the required local primitivity directly. No global
class-number inference or local \(L\)-value nonvanishing is used.

\paragraph{An actual local element annihilates the interior.}
On \(\Delta_5\), the three maps \(a,b,\kappa_3\) are continuous
additive homomorphisms. Their image
\[
 M=\{(a(\tau),b(\tau),\kappa_3(\tau)):\tau\in\Delta_5\}
       \subset\mathbb Z_5^3
\]
is a compact additive subgroup, and hence a closed
\(\mathbb Z_5\)-submodule. The Kummer and Soul\'e covariance gives,
under conjugation by an element with cyclotomic value \(\omega\),
the operator
\[
 T=\operatorname{diag}(\omega,\omega,\omega^3),
 \qquad \omega^2=-1,\quad\omega\in\mathbb Z_5^\times.
\]
Such an actual element of \(G_{\mathbb Q_5}\) exists, since its
cyclotomic character is onto. It is not required to have order four.
The integral operator
\[
 \frac{T-\omega}{\omega^3-\omega}
             =\operatorname{diag}(0,0,1)                \tag{D.77.7}
\]
preserves \(M\), because its denominator is a five-adic unit.
Apply it to an element whose last coordinate is a unit.
By the definition of the actual image \(M\), there is therefore
an actual \(\delta\in\Delta_5\) such that
\[
 a(\delta)=b(\delta)=0,\qquad \kappa_3(\delta)\in\mathbb Z_5^\times .
                                                               \tag{D.77.8}
\]
No formal assignment of Galois coefficients has occurred.

Put \(t_1=1+u,\ t_2=1+v,\ w=t_1t_2-1\), and use the actual beta
series \(\mathcal B_g\) of D.39. On \(\delta\), its exact boundary
values are
\[
 \mathcal B_\delta(t,1)=\mathcal B_\delta(1,t)
          =\mathcal B_\delta(t,t^{-1})=1.
\]
The first two follow from \(\mathcal B_g=1+uvB_g\), and the last
from the exact reflection formula of D.39 with \(\chi(\delta)=1\).
Thus \(B_\delta\) vanishes under \(v=-u/(1+u)\): the factor
\(uv=-u^2/(1+u)\) is a nonzero divisor in \(\mathbb Z_5[[u]]\).
The kernel of this substitution is precisely \((w)\), by division
by the linear polynomial
\(w=(1+u)(v+u/(1+u))\). Consequently
\[
 \mathcal B_\delta-1=uvw\,E_\delta,\qquad
                E_\delta\in\Lambda^\times.             \tag{D.77.9}
\]
To check the unit assertion, the degree-three term in D.35/D.39 is
\[
 \frac{\kappa_3(\delta)}{48}\,uv(u+v).
 \]
Since \(uvw\) has leading term \(uv(u+v)\), the constant term of
\(E_\delta\) is the unit \(\kappa_3(\delta)/48\).
This coefficient is the raw normalization in (D.77.6), not a
renormalized nonprimitive character.

Let \(D_s=A_s^D\) be the discrete Cartier dual, and let
\(H\) be either \(G_{\mathbb Q}\) or \(G_{\mathbb Q_5}\).
By (D.77.8), the action of \(\delta\) on \(A_s\) is exactly
multiplication by \(\mathcal B_\delta\). Its cyclotomic multiplier
is one, so Cartier twisting adds no scalar in this step.
Every \(H\)-invariant functional therefore annihilates
\((\mathcal B_\delta-1)A_s=uvwA_s\). The ideal \((uvw)\)
is stable, because its three constituent boundary ideals are stable.
With
\[
                         B=A_s/uvwA_s
\]
we have proved
\[
                         D_s^H=(B^D)^H.                \tag{D.77.10}
\]
The same actual local \(\delta\) proves this for the global group.

\paragraph{The finite relation among all three boundaries.}
Let \(M_i=A_s/J_iA_s\), where
\[
                  J_1=(v),\qquad J_2=(u),\qquad J_3=(w).
\]
D.76 proves that the actual modules \(M_i\) are respectively
\(M_a,M_b,M_{a-b}\), with the integral half-cyclotomic gauge on
the last axis. Explicitly, its multiplier before that gauge is
\[
 \chi S_\chi(t)\,t^{\,a(g)-b(g)+(1-\chi(g))/2}.
                                                               \tag{D.77.11}
\]
This follows from D.39's exact reflection identity, so includes
all arithmetic coefficients.

The integral restriction map fits in the exact sequence
\[
 0\longrightarrow B\longrightarrow
        M_1\oplus M_2\oplus M_3
          \longrightarrow T\longrightarrow0,
 \qquad T\simeq\mathbb Z_5^3.                            \tag{D.77.12}
\]
Indeed a triple \((F(u),G(v),H(u))\) is in its image exactly if
\[
 F_0=G_0=H_0,\qquad H_1=F_1-G_1.
\]
After matching the first two axes, any higher diagonal
discrepancy is corrected by \(uvL(u)\), since its diagonal
restriction is \(-u^2L(u)/(1+u)\). The intersection of the
three ideals is \((uvw)\): a multiple \(uvL\) vanishing on the
third axis has \(L\) vanishing there, hence is a multiple of \(uvw\).
This also proves exactness at the first term.

Here the action on the finite free cokernel is explicit. Use its
coordinates
\[
 q_1=G_0-F_0,\qquad q_2=H_0-F_0,\qquad q_3=H_1-F_1+G_1.
\]
For \(\alpha=\chi(g)\), before any individual translation gauges,
the matrix is
\[
 \begin{pmatrix}
 \alpha^2&0&0\\
 0&\alpha^2&0\\
 \alpha^2(b(g)+(\alpha-1)/2)&
 \alpha^2(a(g)-b(g))&\alpha^3
 \end{pmatrix}.                                        \tag{D.77.13}
\]
This comes by expanding the actual multipliers to first degree.
It also gives a filtration of \(T\) with characters
\(\chi^2,\chi^2,\chi^3\), without requiring that the extension
between these characters split.

Let \(D_i=M_i^D\). The compact sequence (D.77.12) splits as
\(\mathbb Z_5\)-modules by choosing lifts of the three basis vectors
of its finite free quotient; these lifts define a continuous map.
In particular dualizing gives the exact sequence
\[
 0\longrightarrow T^D\longrightarrow
       D_1\oplus D_2\oplus D_3
          \longrightarrow B^D\longrightarrow0.         \tag{D.77.14}
\]
The Cartier dual \(T^D\) has successive characters
\(\chi^{-1},\chi^{-1},\chi^{-2}\).
Choose an actual \(k\in H\) with \(\chi(k)=2\). Then \(k-1\) is
an automorphism of \(T^D\). For example the determinant of its
corresponding integral lattice matrix is
\[
             (2^{-1}-1)^2(2^{-2}-1)=-3/16,
                                                               \tag{D.77.15}
\]
a five-adic unit. This conclusion is unaffected by the off-diagonal
entries in (D.77.13).

\paragraph{Every invariant boundary sum has an invariant lift.}
Given \(x\in(B^D)^H\), choose any lift
\(z\in D_1\oplus D_2\oplus D_3\). Then \((k-1)z\in T^D\).
Subtracting its unique inverse image under \(k-1\) produces
a lift \(z'\) fixed by \(k\).

Use the measure model of D.75 on each axis. Its elements are
continuous functions on \(\mathbb Z_5\) with values in
\(\mathbb Q_5/\mathbb Z_5\), modulo constants. An integral
translation makes the affine action of \(k\) pure multiplication
by two, since \(\chi(k)-1=1\). This includes the integral gauge
in (D.77.11). Represent each component of \(z'\) by \(f_i\),
with \(f_i(0)=0\). Being fixed by \(k\) says
\[
                              f_i(2z)=f_i(z).           \tag{D.77.16}
\]
The possible constant defect is zero by evaluation at zero.

For \(\tau\in\ker(\chi|_H)\), the defect \(\tau z'-z'\)
belongs to \(T^D\), by invariance of \(x\). A component of \(T^D\)
involves only the constant and linear coefficients of its axis
power series. Under the augmentation-ideal measure model,
\[
 \frac{t^z-1}{t-1}=z+\binom z2(t-1)+\cdots,
\]
so that component is represented by a polynomial of degree at
most two in \(z\), modulo constants. Translation gauges preserve
this degree bound. Hence, with \(h=c_i(\tau)\), changing the sign
of \(h\) if required by the dual-action convention,
\[
 f_i(z+h)-f_i(z)=P_h(z)=A_hz^2+B_hz+C_h.                 \tag{D.77.17}
\]
The coefficients are unique in \(\mathbb Q_5/\mathbb Z_5\):
evaluation at \(0,1,2\) has determinant two. In particular
coefficient comparison is legitimate for this torsion target.

Equation (D.77.16) and telescoping give
\[
 P_{2h}(2z)=P_h(z),\qquad P_{2h}(z)=P_h(z+h)+P_h(z).
\]
The quadratic terms imply \(A_{2h}=2A_h\) and
\(4A_{2h}=A_h\), hence \(7A_h=0\) and \(A_h=0\).
The linear terms then imply \(B_{2h}=2B_h\) and
\(2B_{2h}=B_h\), hence \(3B_h=0\) and \(B_h=0\).
Finally \(C_{2h}=2C_h=C_h\), so \(C_h=0\).
Thus every component of \(z'\) is fixed by \(\ker(\chi|_H)\).

The element two topologically generates \(\mathbb Z_5^\times\):
its order modulo five is four, and \(v_5(2^4-1)=1\), so its
fourth power generates \(1+5\mathbb Z_5\).
Consequently the closed subgroup generated by \(k\) and
\(\ker(\chi|_H)\) is all of \(H\). The components of \(z'\)
are therefore \(H\)-invariant. Conversely the invariant kernel
in (D.77.14) is zero, since \(k-1\) is invertible on \(T^D\).
Together with (D.77.10), this proves the natural isomorphism
\[
 D_1^H\oplus D_2^H\oplus D_3^H
                  \xrightarrow{\ \sim\ }D_s^H
       \quad(H=G_{\mathbb Q},G_{\mathbb Q_5}).           \tag{D.77.18}
\]
The map is the actual sum of the three inclusions, so these two
isomorphisms commute with the inclusion of global in local
invariants. No vanishing of \(H^1(H,T^D)\) has been assumed.

\paragraph{All local and global invariant factors and the remaining term.}
D.75 computes each group in (D.77.18). For a depth \(d<\infty\),
it is \((\mathbb Q_5/\mathbb Z_5)^d\), with basis the functions
\(\mathbf1_{5^j\mathbb Z_5}\), \(1\le j\le d\), modulo constants.
For the zero local class, \(d=\infty\), it is the direct sum of
these groups over all \(j\ge1\): continuity forces finite support
in the ball basis. The compatible global and local gauges in
D.75 identify the first \(r(c_i)\) vectors under restriction
when \(c_i\ne0\). If \(c_i=0\), one global gauge makes both
actions pure unit actions and the inclusion is the identity.

Local duality and the finite-coefficient inverse limits of D.63
identify
\[
 L_5=(D_s^{G_{\mathbb Q_5}})^*,\qquad
 \mathcal T_s=(D_s^{G_{\mathbb Q}})^*.
                                                               \tag{D.77.19}
\]
The star here is ordinary Pontryagin duality of discrete torsion
groups. In particular it is not an additional Cartier twist.
The finite local cohomology groups are finite, so their inverse
systems are Mittag--Leffler; invariants commute with the increasing
union of finite duals. These are exactly the identifications
proved in D.63 from \cite[I.2.1, I.2.3, I.4.10]{DE_MilneADT}.
Thus this step does not assume that one finite-level class
survives to an infinite compatible system.

Dualizing (D.77.18), with the explicit ball bases, proves
(D.77.2) and its assertion about the actual map. Its entire
kernel is \(Z(d(s))\). D.63 now supplies the exact sequence
\[
 0\longrightarrow\mathcal S_s\longrightarrow\mathcal K_s
       \longrightarrow Z(d(s))\longrightarrow0.
                                                               \tag{D.77.20}
\]
The discrete dual of the last term is a direct sum of copies of
\(\mathbb Q_5/\mathbb Z_5\), and is divisible. It is therefore
injective as a \(\mathbb Z_5\)-module: the extension criterion
for ideals \(5^n\mathbb Z_5\) is exactly divisibility by \(5^n\).
Its inclusion in the discrete dual of \(\mathcal K_s\) has a
retraction. Pontryagin duality turns that retraction into a
continuous section of (D.77.20), proving (D.77.3).
Only the compact groups in D.63 are split here. No splitting
of the original nonseparated global cochain quotient is asserted.

The executed exact certificate
\texttt{code/de\_local\_boundary.py} verifies the 24
nontrivial Kummer combinations in the Eisenstein ring, the raw
Soul\'e unit factorization, the image projector, the actual jet
matrix and determinant \(-3/16\), and the order twenty of two
modulo twenty-five. Its JSON and output are retained. These finite
identities support the integral arguments above; they do not
replace the infinite invariant and duality proofs.
\end{proof}


% D.78; root and FGH complete independent proof/source/certificate audits PASS.
\subsection{A nonzero remaining compact kernel of exact order five}

\begin{statement}{D.78}{Proved}{The compact kernel at the point minus two has order five}
Take the actual Alexander twist
\[
 s=(\kappa(2),\kappa(3)),\qquad A_s=\mathbb Z_5[[u,v]]C
\]
of D.35. The remaining term of D.77 is
\[
 \mathcal S_s\simeq\mathbb Z/5,\qquad
 \mathcal K_s\simeq\mathbb Z/5,\qquad L_5=\mathcal T_s=0.
                                                               \tag{D.78.1}
\]
The first isomorphism is detected by an explicit class in a
\emph{single finite} Cartier-dual coefficient module, not just by
classes that become locally trivial at different unbounded levels.
Consequently the separated global second cohomology has the exact
topological description
\[
 0\longrightarrow\mathbb Z/5\longrightarrow\mathcal H_s
       \longrightarrow\mathcal L_s\longrightarrow0,     \tag{D.78.2}
\]
with the actual residue space and topology of D.63. No splitting
of (D.78.2) is asserted.

This is a nonempty metabelian section fibre: the rational point
\(t=-2\) has the specified abelian class. Its full metabelian
fibre is a torsor under \(\prod_{\mathbb N}\mathbb Z_5\), by D.74.
This does not determine higher-derived lifting of every
class in that fibre, or the remaining compact kernels for other
abelian pairs.
\end{statement}

\begin{proof}
Write \(\Lambda=\mathbb Z_5[[u,v]]\), \(t_1=1+u,t_2=1+v\),
and \(D=A_s^D\). Use ordinary Pontryagin duality, denoted by a
star, only after taking Cartier duals.
Let \(I_0\) be the actual compact operator image of the untwisted
action \(\rho_0\) in D.35; let \(I_s\) be the actual image of
\(\rho_s\). These groups include their cyclotomic multiplier.
The twist gives a continuous map
\[
 G_{\mathbb Q}\longrightarrow\mathbb Z_5^2\rtimes I_0,
       \qquad g\longmapsto(a(g),b(g),\rho_0(g)),
                                                               \tag{D.78.3}
\]
where \(I_0\) acts on the translation group by \(\chi\).

\paragraph{The global operator image contains every translation.}
The action of a triple in (D.78.3) is
\(t_1^a t_2^b\rho_0(g)\). This representation of its image is
faithful. If that operator is the identity, comparison of its
values on \(1\) and \(t_1\) first gives \(t_1^\chi=t_1\), hence
\(\chi=1\). Restriction to the first and second boundary axes
then gives \(t_1^a=t_2^b=1\), so \(a=b=0\). Its remaining
untwisted operator is the identity as well. Compactness therefore
identifies the actual image with the corresponding closed
subgroup of the semidirect product.

The untwisted action is unramified outside five, by D.39
(equivalently the actual unramified Jacobi-sum image in
\cite[(1.1)--(1.2) and the following paragraph]{DE_NakamuraSakugawaWojtkowiak2017}).
Inertia at two has Kummer vector \((1,0)\), up to a unit
choice of its five-primary tame generator; inertia at three
has vector \((0,1)\). Their closed images therefore contain
both full coordinate translation groups. Put \(H=\mathbb Z_5^2\).
For any element of the projected image \(I_0\), multiply a lift
by a suitable element of \(H\) to remove its two translation
coordinates. Consequently the actual equality is
\[
                         I_s=H\rtimes I_0.              \tag{D.78.4}
\]
This is an equality of compact operator groups; it is not an
assumption of independence of formal arithmetic coefficients.

\paragraph{The unramified image is entirely supplied at five.}
Put \(K=\mathbb Q(\zeta_5)\), let \(\mathfrak p\) be its unique
prime over five, and set \(P=G_{K,\{\mathfrak p\}}^{(5)}\).
Over \(K\), every finite quotient of the actual untwisted
operator image is a five-group, by its positive augmentation
filtration and \(\chi\equiv1\pmod5\). Thus this image factors
through \(P\).

D.1's class-number and unit calculation gives the following
basis of its Kummer Frattini dual:
\[
 H^1(P,\mathbb F_5)
  \simeq\mathcal O_K[1/5]^\times/
                    \mathcal O_K[1/5]^{\times5}
  =\langle\zeta_5,\phi,\pi\rangle_{\mathbb F_5},
 \quad \phi=-\zeta_5^2-\zeta_5^3,\quad \pi=\zeta_5-1.
                                                               \tag{D.78.5}
\]
Here the identification chooses a fifth root of unity, so the
trivial coefficient over \(K\) is identified with \(\mu_5\).
The valuation of \(\pi\) at \(K_{\mathfrak p}\) is one; D.77
proves independence of the two local unit classes
\(\zeta_5,\phi\). Hence localization is injective on (D.78.5).
The closed local decomposition image in the finitely generated
pro-five group \(P\) therefore maps onto its Frattini quotient
and is all of \(P\). The local cyclotomic image also supplies
the full order-four quotient over \(\mathbb Q\).
It follows that
\[
                G_{\mathbb Q_5}\longrightarrow I_0
                         \quad\hbox{is onto}.           \tag{D.78.6}
\]
This argument uses the exact local unit certificate, not a
surjectivity assertion based on rational dimensions.

\paragraph{The exact closed local logarithm line.}
The local completed Kummer group has basis
\(\kappa(5),\kappa(6)\), as in D.75. Set
\[
 \alpha=\frac{\log16}{4\log6},\qquad
 \beta=\frac{\log81}{4\log6}.
                                                               \tag{D.78.7}
\]
They satisfy
\[
 \alpha\equiv2,\qquad \beta\equiv4,\qquad
                    \alpha-\beta\equiv3\pmod5.
                                                               \tag{D.78.8}
\]
Thus all three are units. Locally the two Kummer classes are
\(\alpha\kappa(6),\beta\kappa(6)\). Apply an integral monomial
gauge to make these equalities hold for actual cocycles.
It preserves (D.78.4), the constant Cartier submodule, and the
finite augmentation quotients used below.

Let \(N\) be the kernel of \(G_{\mathbb Q_5}\to I_0\).
The image \(\kappa(6)(N)\) is all of \(\mathbb Z_5\).
If it were zero modulo five, the actual cocycle
\(\kappa(6)\bmod5\) would descend through the surjection
(D.78.6) to \(I_0\). Inflating to \(G_{\mathbb Q,\{5\}}\)
would give a class in
\[
 H^1(G_{\mathbb Q,\{5\}},\mathbb F_5(1))
                       =\mathbb F_5\kappa(5).
 \]
This is impossible locally: valuation distinguishes
\(\kappa(5)\) from the nonzero unit class \(\kappa(6)\).
The latter is nonzero because \(\log6\) has valuation one,
whereas a fifth power of a principal unit has logarithm
divisible by \(25\). Thus the closed additive subgroup
\(\kappa(6)(N)\) contains a unit and is \(\mathbb Z_5\).

It follows that the local operator image has precisely the
primitive translation line
\[
 H_5=\{(\alpha h,\beta h):h\in\mathbb Z_5\}.
 \]
All its other translation coordinates lie on this line as
well, so they can be removed without changing the projected
element of \(I_0\). Therefore, after the specified integral gauge,
\[
                  I_5=\rho_s(G_{\mathbb Q_5})
                            =H_5\rtimes I_0.            \tag{D.78.9}
\]
The common complement in (D.78.4) and (D.78.9) is the actual
untwisted image, with zero translation coordinates.

\paragraph{Continuous translation coinduction and the local restriction.}
Under \(H\), the Cartier dual \(D\) is the group of continuous
functions \(H\to\mathbb Q_5/\mathbb Z_5\), with translation
action. Indeed \(\Lambda\) is the completed group ring of
\(H\); on \(H\), the cyclotomic twist is one. This is a
continuous coinduced module, so
\[
 H^j(H,D)=0\quad(j>0),\qquad
 D^H=\lambda:=\mathbb Q_5/\mathbb Z_5(-1).               \tag{D.78.10}
\]
The character on the last group follows from the augmentation
character \(\chi^2\) on \(A_s/H\), with the single Cartier
twist. There is no ambiguity between this constant functional
and constant functions modulo constants in D.75's different,
augmentation-ideal model.

For clarity, the acyclicity assertion has an elementary
continuous proof. Homogeneous cochains with values in
continuous functions on the compact translation group identify
with continuous functions in one additional group variable.
Inserting the identity in that variable contracts the
positive-degree complex. All functions involved have finite
image locally, and compactness makes the required partitions
uniformly finite, so the contraction is continuous.
The same argument works for \(H_5\): a continuous linear
complement to its primitive line identifies \(H\) with
\(H_5\times\mathbb Z_5\), and the second factor is simply
part of the discrete function coefficient module.

Let \(D_{\rm lin}=D^{H_5}\). It is the Cartier dual of
\[
 B=\Lambda/(t_1^\alpha t_2^\beta-1)
                \simeq\mathbb Z_5[[z]],\qquad t=1+z,
 \quad t_1=t^\beta,\quad t_2=t^{-\alpha}.
                                                               \tag{D.78.11}
\]
These substitutions give an integral isomorphism, since
\(\alpha,\beta\) are units. Continuous Hochschild--Serre,
using the acyclicity just proved, now gives natural isomorphisms
\[
 \begin{split}
 H^1(I_s,D)&=H^1(I_0,\lambda),\\
 H^1(I_5,D)&=H^1(I_0,D_{\rm lin}).
 \end{split}                                               \tag{D.78.12}
\]
Under them restriction is induced by the actual inclusion
\(\lambda\subset D_{\rm lin}\). Neither group on the right
is being inferred from a formal Lie algebra.

\paragraph{Its entire kernel is of exact order five.}
Choose an actual \(\delta\in I_0\) with \(\chi(\delta)=1\)
and primitive \(\kappa_3(\delta)\), supplied locally by D.77.
Its multiplier on \(B\) is \(\mathcal B_\delta(t^\beta,t^{-\alpha})\).
Formula (D.77.9) and (D.78.8) give
\[
 \mathcal B_\delta(t^\beta,t^{-\alpha})-1
      =z^3\,e(z),\qquad e(0)\in\mathbb Z_5^\times .
                                                               \tag{D.78.13}
\]
The leading coefficient is
\[
 e(0)=\frac{\kappa_3(\delta)}{48}\,
                     \beta(-\alpha)(\beta-\alpha).
                                                               \tag{D.78.14}
\]
Thus \(D_{\rm lin}^{I_0}\) is contained in the dual of
\(B/z^3B\). Its three successive characters are
\(\chi^{-1},\chi^{-2},\chi^{-3}\). For an actual element
\(k\) with \(\chi(k)=2\), all three differences from one
are units modulo five. Hence
\[
                         D_{\rm lin}^{I_0}=0.           \tag{D.78.15}
\]

The quotient \(D_{\rm lin}/\lambda\) is the Cartier dual
of the stable ideal \(zB\). By (D.78.13), its invariant
part is contained in
\[
                       V=(zB/z^4B)^D.
 \]
On this module the cyclotomic kernel of \(I_0\) acts
trivially: its exact beta multiplier starts in degree
three, which acts trivially after multiplication by \(z\)
and reduction modulo \(z^4\). The three successive
characters of \(V\) are
\(\chi^{-2},\chi^{-3},\chi^{-4}\).
The three eigenvalues of \(k\) have pairwise unit
differences; its triangular integral matrix is therefore
diagonalized over \(\mathbb Z_5\), by successive elementary
changes of basis. Its invariant subgroup on the torsion
dual has order
\[
       \#\ker(2^{-4}-1:\mathbb Q_5/\mathbb Z_5
                         \to\mathbb Q_5/\mathbb Z_5)=5,
 \]
since \(v_5(2^4-1)=1\); the other two kernels are zero.
The element two topologically generates \(\mathbb Z_5^\times\).
We obtain
\[
                (D_{\rm lin}/\lambda)^{I_0}\simeq\mathbb F_5.
                                                               \tag{D.78.16}
\]
By the cohomology sequence of
\(0\to\lambda\to D_{\rm lin}\to D_{\rm lin}/\lambda\to0\)
and (D.78.15), (D.78.16) identifies the entire kernel of
the restriction map in (D.78.12), not merely a subgroup.

Its generator has a precise arithmetic description.
Let \(c(g)\in\mathbb F_5\) be the reduction of the normalized
integral weight-three cocycle \(48h(g)\) in D.31.
Its restriction to \(\ker\chi\) is the raw \(\kappa_3\),
and D.39 shows that it factors through \(I_0\).
Modulo five its character \(\chi^3\) equals \(\chi^{-1}\),
so it defines a cocycle with values in \(\lambda[5]\).
It is nonzero because its restriction to the local
cyclotomic kernel is primitive.

To identify the connecting homomorphism in (D.78.16),
lift its invariant vector to a functional on \(B/z^4B\)
of order five with nonzero top coefficient. Such a lift
can be made \(k\)-fixed by changing its constant coefficient,
since \(k-1\) is invertible on \(\lambda[5]\).
For every \(\tau\in\ker\chi\), (D.78.14) gives the defect
as a fixed nonzero scalar times
\(\kappa_3(\tau)\) times the constant functional.
Thus its connecting class is a nonzero scalar multiple
of \([c]\). Equality follows from restriction to
\(\ker\chi\), since
\(H^1(\mathbb Z_5^\times,\mathbb F_5(-1))=0\), by the
order-four scalar. All lifts and corrections can be
made in the order-five dual of \(B/z^4B\).
In particular the class \([c]\) has a local primitive at
five in a finite three-jet module.

\paragraph{An upper bound for the actual inverse finite-localization kernel.}
Finite Poitou--Tate duality in D.63,
using \cite[I, Theorem 4.10(a)]{DE_MilneADT}, identifies
\[
 \mathcal S_s^*=\varinjlim_r
       \operatorname{Sha}^1(G_{\mathbb Q},W_r^D),
       \qquad W_r=A_s/(5,u,v)^rA_s,                     \tag{D.78.17}
\]
where \(\operatorname{Sha}^1\) is the kernel of restriction to all places
at that \emph{one finite coefficient level}. The direct
limit injects into \(H^1(G_{\mathbb Q},D)\), by exactness
of filtered colimits of the finite-coefficient cohomology
groups and their subgroups. This does not commute a
colimit with an infinite product of local groups.

Every finite class in (D.78.17) descends to the actual
finite operator image. If a cocycle \(z_r\) were nonzero
on an element acting trivially on \(W_r^D\), the finite
affine image \((z_r,\rho_r)\) would contain
\((m,1)\) with \(m\ne0\). Chebotarev at a prime outside
its finite ramification set would realize a conjugate
Frobenius. A local coboundary there evaluates to zero on
that Frobenius, whereas this cocycle does not. This
contradicts its local triviality. Hence \(z_r\) factors
through that finite image, and therefore through \(I_s\).
Inflation from \(I_s\) to the global group is injective
in degree one. Likewise inflation from \(I_5\) to the
local group at five is injective. Thus (D.78.17) injects
into the restriction kernel computed in (D.78.16):
\[
                         \mathcal S_s^*
                                  \lhook\joinrel\longrightarrow\mathbb F_5.
                                                               \tag{D.78.18}
\]
The Chebotarev use is the exact finite-image theorem
\cite[V, Theorem 3.23]{PBC_MilneCFT}.

\paragraph{One bounded coefficient module works at every place.}
It remains to prove that \([c]\) survives in (D.78.17).
This is the uniformity step; pointwise division by \(q-1\)
would not prove it.
Let \(\epsilon_{ij}\) be the coefficient functional on
\(u^iv^j\), with values in \(\mathbb F_5\), and work in
the finite Cartier-dual module
\[
 M=\bigl(A_s/(5,(u,v)^4)A_s\bigr)^D
                  \subset W_4^D.
                                                               \tag{D.78.19}
\]
The constant functional is \(\epsilon_{00}\), and the
global cocycle is \(z(g)=c(g)\epsilon_{00}\).

For \(\bar\chi=\chi(g)\bmod5\ne1\), it has a primitive in
\(\mathbb F_5\epsilon_{00}\), since \(\bar\chi^{-1}-1\)
is a unit. Suppose now that \(\bar\chi=1\).
Through total degree three, D.31/D.39 and the actual
Fox formula give
\[
 \rho_s(g)f=
 t_1^{a(g)}t_2^{b(g)}
 \left(1+\frac{c(g)}{48}uv(u+v)\right)f
       \pmod{5,(u,v)^4}.                                \tag{D.78.20}
\]
Here substitution is the identity to this degree, and
both \(S_\chi\)-factors are one. Indeed their relevant
binomial denominators are units at five and
\(\chi-1\) is divisible by five. The weight-two beta
coefficient \((\chi^2-1)/24\) is zero modulo five;
conversion from logarithmic coordinates therefore
adds no cubic term. This verifies exactly the normalized
coefficient used in (D.78.20).

If \(\bar a\ne0\), then
\[
           (g-1)\epsilon_{10}=-\bar a\,\epsilon_{00},
 \]
and \(-c(g)\bar a^{-1}\epsilon_{10}\) is a primitive.
If \(\bar a=0,\bar b\ne0\), use
\(-c(g)\bar b^{-1}\epsilon_{01}\).
If both are zero, (D.78.20) gives
\[
           (g-1)\epsilon_{21}
                         =-\frac{c(g)}{48}\epsilon_{00},
 \]
so the single vector \(-48\epsilon_{21}\) is a primitive.
These equations use the inverse operator in the dual
action; their minus signs are therefore intentional.
All the primitives are in the one fixed finite module \(M\).

Outside \(\{2,3,5\}\), both the actual module and this
cocycle are unramified. The preceding primitive on
Frobenius consequently trivializes the whole local
cocycle. At two and three the constant primitive works
on the whole decomposition group: its Frobenius
multiplier is respectively two or three modulo five,
inertia fixes the constant functional, and \(c\) is
unramified there. At five the three-jet primitive
constructed after (D.78.16) belongs to \(M\), since
the quotient substitutions (D.78.11) and all integral
monomial gauges preserve the order-five three-jet
module. At the real place odd-primary cohomology is
zero (and the chosen normalized cocycle is zero on
complex conjugation). Thus
\[
                 [c]\in\operatorname{Sha}^1(G_{\mathbb Q},M)
                           \longrightarrow\operatorname{Sha}^1(G_{\mathbb Q},W_4^D).
                                                               \tag{D.78.21}
\]
This is a single global finite-level certificate.

The class remains nonzero in \(H^1(G_{\mathbb Q},D)\).
In (D.78.12) it comes from the nonzero class in
\(H^1(I_0,\lambda[5])\); this class remains nonzero in
\(H^1(I_0,\lambda)\), since \(\lambda^{I_0}=0\) and
the multiplication-by-five coefficient sequence makes
that inclusion injective in degree one. The first
isomorphism of (D.78.12) and global inflation preserve
it. Therefore (D.78.21) gives a nonzero element of
(D.78.17). Together with (D.78.18), it proves
\(\mathcal S_s^*\simeq\mathbb F_5\), hence
\(\mathcal S_s\simeq\mathbb Z/5\).

\paragraph{The exact compact term and actual section fibre.}
Each of \(a,b,a-b\) has global depth zero. By
(D.78.7)--(D.78.8), each also has local-at-five
depth zero. Thus D.77 gives \(L_5=\mathcal T_s=0\)
and \(\mathcal K_s=\mathcal S_s\), proving (D.78.1)
and (D.78.2).

For the actual rational point \(t=-2\), the two
Kummer coordinates are \(\kappa(-2),\kappa(3)\).
The first equals \(\kappa(2)\) in completed pro-five
cohomology, since \(-1\) has zero pro-five Kummer class.
Hence the metabelian fibre is nonempty, and D.35/D.74
give the asserted torsor. There is no assertion that
its nonzero second-cohomology classes obstruct this
particular point section; that section already exists.
The executed exact certificate \path{code/de_finite_sha.py}
derives the actual Fox-semilinear action through degree three
and checks all 500 residual primitive cases in (D.78.19).
It also computes the local line matrices and the determinant
\(-315/512\) of the quotient Cartier operator minus one,
whose five-adic valuation is one. The arithmetic image,
coinduction, and global duality assertions follow from the
proof above.
\end{proof}


% D.79. Root, FGH, and PBC complete independent proof/source/certificate audits PASS.
\subsection{The entire compact kernel on the two-prime invertible matrix family}

\begin{statement}{D.79}{Proved}{The complete two-prime matrix-family compact kernel}
Let \(M\in\operatorname{GL}_2(\mathbb Z_5)\), and put
\[
 \binom ab=M\binom{\kappa(2)}{\kappa(3)},\qquad
 \binom\alpha\beta
 =M\binom{\log16/(4\log6)}{\log81/(4\log6)}.
                                                               \tag{D.79.1}
\]
The vector \((\alpha,\beta)\) is primitive. With
\(v_5(0)=\infty\), set
\[
 d=\max\{v_5(\alpha),v_5(\beta),v_5(\alpha-\beta)\}.
                                                               \tag{D.79.2}
\]
At most one of these valuations is positive. For the actual
Alexander twist \(s=(a,b)\), we have
\[
 \mathcal S_s=
 \begin{cases}
  \mathbb Z/5,&d=0,\\
  0,&1\le d\le\infty,
 \end{cases}
 \qquad \mathcal T_s=0,
                                                               \tag{D.79.3}
\]
and, as compact groups,
\[
 \mathcal K_s\simeq
 \begin{cases}
  \mathbb Z/5,&d=0,\\
  \mathbb Z_5^d,&1\le d<\infty,\\
  \displaystyle\prod_{\mathbb N}\mathbb Z_5,&d=\infty.
 \end{cases}                                                   \tag{D.79.4}
\]
Here \(\mathcal S_s\) is exactly the inverse finite-global-localization
kernel of D.63, not a kernel formed only after passing to divisible
local coefficients. The separated second cohomology is the open
extension of the actual residue space by (D.79.4) in D.63.
No splitting of that latter extension is asserted.

This classifies the compact kernel for every matrix in this family,
including arbitrarily deep and exact boundary slopes. It does not
assert that every matrix pair has a metabelian section. If its fibre
is nonempty, D.74 makes it a torsor under the compact countable
product of \(\mathbb Z_5\). The identity matrix includes the actual
point \(t=-2\) of D.78.
\end{statement}

\begin{proof}
Use D.78's notation \(I_0,I_s,H,D,\lambda\). Inertia at two and
three gives the two columns of \(M\) in the translation group.
Because \(M\) is invertible, these columns generate the whole
\(H=\mathbb Z_5^2\). The faithful actual operator argument of
D.78 therefore gives
\[
 I_s=H\rtimes I_0.
                                                               \tag{D.79.5}
\]
The local surjection \(G_{\mathbb Q_5}\to I_0\) is independent
of the pair. After the same integral Kummer gauge as in D.78,
the local pair is \((\alpha\kappa(6),\beta\kappa(6))\).
The vector is primitive since the vector \((2,4)\bmod5\) is
nonzero and \(M\) is invertible. D.78's exact calculation
\(\kappa(6)(\ker(G_{\mathbb Q_5}\to I_0))=\mathbb Z_5\)
now gives the actual equality
\[
 I_5=H_5\rtimes I_0,
 \qquad H_5=\mathbb Z_5(\alpha,\beta).
                                                               \tag{D.79.6}
\]
Thus continuous translation coinduction gives, naturally,
\[
 H^1(I_s,D)=H^1(I_0,\lambda),\qquad
 H^1(I_5,D)=H^1(I_0,D_{\rm lin}),
 \quad\lambda=\mathbb Q_5/\mathbb Z_5(-1).
                                                               \tag{D.79.7}
\]
Every finite everywhere locally trivial class descends to the actual
operator image by D.78's finite affine Chebotarev argument.
Consequently
\[
 \mathcal S_s^*\ \subseteq\
 \ker\{H^1(I_0,\lambda)\longrightarrow H^1(I_0,D_{\rm lin})\}.
                                                               \tag{D.79.8}
\]
Here the star is Pontryagin dual; all coefficient modules inside
the cohomology groups are Cartier duals.

\paragraph{Boundary permutations are actual geometric comparisons.}
The rational automorphisms \(t\mapsto1-t\) and \(t\mapsto1/t\)
of the tripod permute its three peripheral directions. Their
maps on pro-five Kummer pairs are
\[
 (a,b)\longmapsto(b,a),\qquad
 (a,b)\longmapsto(-a,b-a),                              \tag{D.79.9}
\]
because the extra Kummer class of \(-1\) is zero at five.
These are comparisons of the actual arithmetic fundamental
extensions. Choose a geometric path between the transported unit
tangent and the original reference tangent. Its ambiguity is an
inner geometric automorphism, which gives a coefficient-module
gauge and hence the same continuous cohomology groups. The
transported unit cusp has zero abelian class, by the three actual
packet formulas in D.35; an integral abelian gauge therefore
identifies the reference abelian representatives. Thus (D.79.9)
transports the actual Alexander module and every local and global
cohomology map, including its finite-coefficient inverse limits.
It is not an inference from an abstract formal substitution.

The transformations are invertible over \(\mathbb Z_5\), keep
ramification inside \(\{2,3,5\}\), and permute the three
valuations in (D.79.2). They allow us, whenever \(d>0\), to
arrange that
\[
 v_5(\alpha)=d,\qquad \beta\in\mathbb Z_5^\times.
                                                               \tag{D.79.10}
\]
The exact diagonal comparison could equivalently be made using
D.39's half-cyclotomic monomial gauge, since two is a unit.

\paragraph{The case of three unit slopes.}
If \(d=0\), the proof of D.78 applies verbatim to the actual
images (D.79.5)--(D.79.6): the primitive local cubic coefficient
is \(\kappa_3\beta(-\alpha)(\beta-\alpha)/48\), a unit.
It computes the entire restriction kernel in (D.79.8) as
\(\mathbb F_5\). Its generator is the reduction of the integral
cocycle \(48h\) of D.31. The fixed ten-dimensional module of
D.78 still trivializes it at every prime. Outside two, three,
five everything is unramified; at two and three the constant
primitive works because the residual cyclotomic multipliers
are not one. The local unit-slope three-jet proof supplies
the primitive at five. Hence this is a single finite-level
everywhere locally trivial class, and
\(\mathcal S_s=\mathbb Z/5\).

\paragraph{An exact boundary gives an injective local restriction.}
If \(d=\infty\), make (D.79.10) with \(\alpha=0\).
Then
\[
 D_{\rm lin}=(\Lambda/(t_2-1))^D.
\]
Its action factors through \(\Gamma=\mathbb Z_5^\times\),
by the exact axis formula \(\chi S_\chi(t_1)\) in D.35.
Put \(N_0=\ker(I_0\to\Gamma)\). Both \(\lambda\) and
\(D_{\rm lin}\) have trivial \(N_0\)-action. The inclusion
\(\lambda\subset D_{\rm lin}\) is injective, so it induces
an injection on continuous homomorphisms from \(N_0\).
Moreover
\(H^1(\Gamma,\lambda)=0\), since its order-four subgroup
acts by a nontrivial scalar with unit difference from one.
Inflation--restriction therefore makes
\(H^1(I_0,\lambda)\to H^1(I_0,D_{\rm lin})\) injective.
Equation (D.79.8) gives \(\mathcal S_s=0\).

\paragraph{The exact finite module at a boundary depth.}
Now assume \(1\le d<\infty\) and (D.79.10). Write
\[
 N=5^d,\qquad \alpha=N a_0,
 \quad a_0,\beta\in\mathbb Z_5^\times,
 \qquad R=\mathbb Z_5[[z]],\quad t=1+z.
\]
The local line quotient is given integrally by
\(t_1=t^\beta,t_2=t^{-\alpha}\). Define the distinguished
polynomial
\[
 Q(z)=z^2(t^N-1),\qquad \deg Q=N+2.                   \tag{D.79.11}
\]
For every actual cyclotomic-kernel element \(\tau\in I_0\),
D.39's three-boundary factorization shows that its multiplier
minus one is divisible by \(Q\). For the actual primitive
local \(\delta\) of D.77 it is exactly
\[
                   \rho_0(\delta)-1=Q(z)E(z),
                          \qquad E\in R^\times.        \tag{D.79.12}
\]
Indeed \(t^{-\alpha}-1\) is \(t^N-1\) times a unit,
whereas \(t^\beta-1\) and \(t^{\beta-\alpha}-1\)
are \(z\) times units. The remaining beta factor is a unit
by the primitive cubic coefficient. All these statements are
integral, with no division by \(5^d\).

Consequently
\[
 D_{\rm lin}^{I_0}=(R/QR)^D{}^{\Gamma},\qquad
 (D_{\rm lin}/\lambda)^{I_0}
                    =(zR/zQR)^D{}^{\Gamma}.           \tag{D.79.13}
\]
The action on both finite-rank compact quotients factors
through \(\Gamma\), because every \(\tau\in N_0\)
acts by a multiplier congruent to one modulo \(Q\).
The second identity uses the dual of the stable ideal \(zR\).

Let \(V=(zR/zQR)^D\), and put \(J=(R/QR)^D\).
The actual inclusion of ideals gives the exact dual sequence
\[
 0\longrightarrow\lambda\longrightarrow J
       \longrightarrow V
       \longrightarrow (QR/zQR)^D\longrightarrow0.
                                                               \tag{D.79.14}
\]
On \(QR/zQR\), the arithmetic character is \(\chi^5\):
the multiplier has constant term \(\chi^2\), and the
ratio \(Q(t^\chi-1)/Q(t-1)\) has constant term \(\chi^3\).
The last Cartier dual is therefore
\(\mathbb Q_5/\mathbb Z_5(-4)\). Since
\(H^0(\Gamma,\lambda)=H^1(\Gamma,\lambda)=0\),
(D.79.14) yields an injection
\[
 \frac{V^\Gamma}{D_{\rm lin}^{I_0}}
       \lhook\joinrel\longrightarrow
 (\mathbb Q_5/\mathbb Z_5(-4))^\Gamma
                    =\mathbb F_5.                    \tag{D.79.15}
\]
The quotient on the left is exactly the full restriction
kernel in (D.79.8), by the coefficient long exact sequence.
Thus it remains to exclude this one possible order-five class.

\paragraph{A uniform actual cyclotomic operator kills that class.}
Let \(k\in I_0\) be actual with \(\chi(k)=2\). For a general
actual element with cyclotomic multiplier \(\gamma\), write its
operator on the line as \(U_\gamma(t)\varphi_\gamma\), where
\(\varphi_\gamma(t)=t^\gamma\). Its class modulo \(Q\) is
independent of the choice of that element, by (D.79.12).
The exact axis formula shows that
\(U_\gamma-\gamma S_\gamma(t^\beta)\) is divisible by
\(t^N-1\): each of its two differences contains
\(t^{-\alpha}-1\). Modulo \(Q=z^2(t^N-1)\) it therefore is
\[
 U_\gamma(t)\equiv
 \gamma S_\gamma(t^\beta)+(t^N-1)(A_\gamma+B_\gamma z)
                                                      \pmod Q,\tag{D.79.16}
\]
where the exact first two Taylor coefficients give
\[
 \begin{split}
 A_\gamma&=-\frac{a_0\gamma^2(\gamma-1)}2,\\
 B_\gamma&=a_0\gamma^2(\gamma-1)
       \left(\frac{\alpha(2\gamma-1)}{12}
                    +\frac N4-\frac{\beta(7\gamma-5)}{24}\right).
 \end{split}                                             \tag{D.79.17}
\]
To check the only beta input in this computation, its degree-two
coefficient is \((\gamma^2-1)uv/24\), exactly D.31/D.39.
All denominators in (D.79.17) are five-adic units. In particular
\(A_\gamma,B_\gamma\in5\mathbb Z_5\) whenever
\(\gamma\equiv1\pmod5\).

Take the actual power
\[
 g=k^{4\cdot5^{d-1}},\qquad
 \gamma=2^{4\cdot5^{d-1}}
                  \equiv1+3N\pmod{5N}.               \tag{D.79.18}
\]
The congruence follows by binomial expansion of
\((1+15)^{5^{d-1}}\). Modulo five, to degree \(N+2\),
\[
 \begin{split}
 \varphi_\gamma(z)&\equiv z+3z^N+3z^{N+1},\\
 U_\gamma(t)&\equiv S_\gamma(t^\beta)
                 \equiv1+3z^{N-1}+O(z^N)\pmod Q.
 \end{split}                                             \tag{D.79.19}
\]
For the second line, use \(\beta^{N-1}\equiv1\pmod5\).
The first line follows from
\((1+z)^\gamma\equiv(1+z)(1+3z^N)\) in this range;
the omitted next degree \(2N\) exceeds \(N+2\), since
\(N\ge5\).

In the compact module \(zR/zQR\), whose reduction modulo five
has basis \(z,\ldots,z^{N+2}\), these identities imply
\[
 (\gamma^{-1}\rho_0(g)-1)z^3
       \equiv (3\cdot3+3)z^{N+2}
       =2z^{N+2}\pmod{5,zQ}.                           \tag{D.79.20}
\]
The leading term of \(Q\) modulo five is exactly \(z^{N+2}\).
Every order-five \(\Gamma\)-invariant Cartier functional on
\(zR/zQR\) annihilates the left side of (D.79.20); hence it
annihilates \(Q\). Its image in the last term of (D.79.14)
is therefore zero. This argument uses an actual cyclotomic
power, not a formal graded eigenvalue.

Finally D.77 computes
\[
 D_{\rm lin}^{I_0}=D^{I_5}
                  \simeq(\mathbb Q_5/\mathbb Z_5)^d.   \tag{D.79.21}
\]
It is divisible. If the quotient in (D.79.15) were nonzero,
the inclusion of this divisible group into \(V^\Gamma\)
would split abstractly: a divisible abelian group is injective.
Its quotient, which has order five, would consequently have
an order-five lift in \(V^\Gamma\) mapping nontrivially to
\((QR/zQR)^D\). Equation (D.79.20) excludes that lift.
Thus the whole restriction kernel in (D.79.8) is zero.
This proves \(\mathcal S_s=0\) for every finite positive depth.

\paragraph{The compact groups and the remaining scope.}
Every one of \(a,b,a-b\) has global divisibility zero: the two
rows of an invertible matrix, and their difference, are
primitive. The local depths are respectively the valuations
in (D.79.2). D.77 therefore gives \(\mathcal T_s=0\),
\(L_5=\mathbb Z_5^d\) at finite depth, and the countable
compact product at infinite depth. Its continuous product
decomposition \(\mathcal K_s=\mathcal S_s\times L_5\)
proves (D.79.4). Finiteness of the ramification set is automatic
for the completed matrix combinations; D.74 applies to any
nonempty fibre. Section existence itself and the separated
extension class are not deduced from the cohomology calculation.
\end{proof}


% D.80. Root, FGH, and PBC complete independent proof/source/certificate audits PASS.
\subsection{An exact compact-kernel family with arbitrary infinite ramification}

\begin{statement}{D.80}{Proved}{Unrestricted independent unit pairs and actual point examples}
Let \(a,b\in H^1(G_{\mathbb Q},\mathbb Z_5(1))\) be arbitrary
completed Kummer classes. Assume that their valuations at five
are zero, their reductions modulo five are linearly independent,
and their local unit-logarithm vector is primitive:
\[
 a|_{G_{\mathbb Q_5}}=\alpha\kappa(6),\qquad
 b|_{G_{\mathbb Q_5}}=\beta\kappa(6),\qquad
                  (\alpha,\beta)\notin5\mathbb Z_5^2.
                                                               \tag{D.80.1}
\]
No finite ramification set for \(a,b\) is assumed. Define
\[
 d=\max\{v_5(\alpha),v_5(\beta),v_5(\alpha-\beta)\},
                    \qquad v_5(0)=\infty.
                                                               \tag{D.80.2}
\]
Then we have the full compact-kernel classification
\[
 \mathcal S_s=
 \begin{cases}\mathbb Z/5,&d=0,\\0,&d>0,\end{cases}
 \qquad\mathcal T_s=0,\qquad
 \mathcal K_s\simeq
 \begin{cases}
   \mathbb Z/5,&d=0,\\
   \mathbb Z_5^d,&1\le d<\infty,\\
   \displaystyle\prod_{\mathbb N}\mathbb Z_5,&d=\infty.
 \end{cases}                                                   \tag{D.80.3}
\]
These are compact topological groups in D.63's separated
second-cohomology sequence. The general nonseparated kernel,
arbitrary infinitely ramified first-cohomology ranks, and
section existence are not determined by this assertion.

There are actual rational-point fibres with compact kernel
\(\mathbb Z_5^d\) for every integer \(d\ge1\). Explicitly choose
the integer \(t_d\) by
\[
 \begin{split}
 t_d&=5^{d+1}-1+5^{d+2}k_d,\qquad 0\le k_d<36,\\
 k_d&\equiv(23-5^{d+1})(5^{d+2})^{-1}\pmod{36}.
 \end{split}                                                   \tag{D.80.4}
\]
The actual point \(t_d\) on the tripod has pair
\((\kappa(t_d),\kappa(1-t_d))\) satisfying the hypotheses
and depth \(d\). Its nonempty metabelian section fibre is
a torsor under \(\prod_{\mathbb N}\mathbb Z_5\), by D.74.
The point \(t=-2\) supplies the order-five case \(d=0\).
\end{statement}

\begin{proof}
\paragraph{The actual image from finitely many residual inertia vectors.}
The pro-five completed multiplicative group of \(\mathbb Q\)
has the valuation description
\[
 H^1(G_{\mathbb Q},\mathbb Z_5(1))
   =\{(x_q)_q\in\prod_q\mathbb Z_5:
                     x_q\longrightarrow0\text{ five-adically}\}.
                                                               \tag{D.80.5}
\]
Indeed \(\mathbb Q^\times=\{\pm1\}\oplus\bigoplus_q\mathbb Z\)
and the sign disappears in the pro-five completion. At each
finite coefficient level its valuation support is finite.
The local Kummer valuation of a class is its indicated
coordinate, so (D.80.5) applies to actual inertia images.

Modulo five, the vectors \((a_q,b_q)\) for \(q\ne5\)
span \(\mathbb F_5^2\). Otherwise a nonzero linear relation
on that span, together with the zero coordinate at five,
would contradict the assumed independence of \(a,b\).
Only finitely many such vectors are nonzero modulo five.
Choose two whose determinant is a unit in \(\mathbb Z_5\).
Their actual inertia images lie in the pure translation
kernel of the untwisted operator image \(I_0\), and
generate the whole translation group \(H=\mathbb Z_5^2\).
Thus D.78's faithful compact-image proof gives
\[
                         I_s=H\rtimes I_0.             \tag{D.80.6}
\]
This does not require bounded global ramification.

The local logarithm in (D.80.1) is continuous on the completed
group. Explicitly, for \(x_5=0\) its coefficient is
\[
      \sum_{q\ne5}x_q\frac{\log(q^4)}{4\log6}
                          \quad\hbox{in }\mathbb Z_5.
                                                               \tag{D.80.7}
\]
The summands tend uniformly to zero because the logarithm
ratios are integral and \(x_q\to0\). Thus the local unit
description does not involve a conditionally convergent
choice of ordering. After an integral cocycle gauge,
the exact local image is
\[
 I_5=\mathbb Z_5(\alpha,\beta)\rtimes I_0,              \tag{D.80.8}
\]
by the same local surjectivity and \(\kappa(6)\)-saturation
argument of D.78. The chosen local line is primitive.

The coinduction argument, the finite-image Chebotarev
injection of \(\mathcal S_s^*\), and the complete local
restriction-kernel calculation of D.79 now apply to
(D.80.6)--(D.80.8). None uses the global support after
these actual images have been established. In particular
\[
 \mathcal S_s^*=0\quad(d>0),\qquad
 \mathcal S_s^*\subseteq\mathbb F_5\quad(d=0).           \tag{D.80.9}
\]
For \(d=0\), the possible generator is the cocycle
\(c=48h\bmod5\) with constant Cartier coefficients,
as in D.78. It remains to prove that this class occurs
at one common finite coefficient level despite possible
infinite ramification.

\paragraph{Local vanishing on discrete Cartier coefficients.}
Fix \(q\ne5\) at which the abelian inertia vector is
nonzero. Put \(A=A_s\), \(W_r=A/\mathfrak m^r A\),
and \(D=A^D=\varinjlim W_r^D\), where
\(\mathfrak m=(5,u,v)\). D.43 proves
\(H^1(G_{\mathbb Q_q},A)=0\) with compact integral
coefficients, for every such inertia vector, whether
or not a metabelian section exists.

Here is the exact passage to the discrete dual; it is
not an unsupported formal duality assertion. Each
\(H^0(G_{\mathbb Q_q},W_r)\) and
\(H^1(G_{\mathbb Q_q},W_r)\) is finite. The compact
coefficient exact sequence and the Mittag--Leffler
property of the first tower therefore give
\[
 \varprojlim_r H^1(G_{\mathbb Q_q},W_r)
                    =H^1(G_{\mathbb Q_q},A)=0.          \tag{D.80.10}
\]
Finite local Tate duality
\cite[I, Corollary 2.3]{DE_MilneADT} identifies
\(H^1(G_{\mathbb Q_q},W_r^D)\) with the finite
Pontryagin dual of \(H^1(G_{\mathbb Q_q},W_r)\),
naturally for coefficient transitions. Pontryagin
duality between finite inverse and direct limits, and
continuity of cohomology on discrete filtered unions,
then give
\[
 \begin{split}
 H^1(G_{\mathbb Q_q},D)
  &=\varinjlim_r H^1(G_{\mathbb Q_q},W_r^D)\\
  &=\left(\varprojlim_r H^1(G_{\mathbb Q_q},W_r)\right)^*
   =0.
 \end{split}                                             \tag{D.80.11}
\]
In particular any class represented at one finite level
becomes a coboundary at some larger finite level.
No uniformity in \(q\) is concluded from this step alone.

\paragraph{Only finitely many places need those larger levels.}
Use the one fixed module of D.78,
\[
          M=(A_s/(5,(u,v)^4)A_s)^D\subset W_4^D.
                                                               \tag{D.80.12}
\]
Let
\[
 E=\{q\ne5:(a_q,b_q)\notin5\mathbb Z_5^2\}.
                                                               \tag{D.80.13}
\]
This set is finite by (D.80.5). For every \(q\notin E\cup\{5\}\),
the module \(M\) is unramified even if \(A_s\) is ramified:
its inertia multiplier satisfies
\[
 t_1^{a_q}t_2^{b_q}\equiv1\pmod{5,(u,v)^4},             \tag{D.80.14}
\]
because \((1+u)^{5x}=(1+u^5)^x\) modulo five, and similarly
for \(v\). The original untwisted action and \(c\) are
unramified at every \(q\ne5\). Hence D.78's universal
constant, linear and cubic Frobenius primitives trivialize
the local class of \(c\) in \(M\) at every prime outside
this finite exceptional set. The identities apply to the
actual Frobenius values of \(a,b,c,\chi\), without a
restriction on their ramification at higher precision.

At each of the finitely many primes in \(E\), the inertia
vector is nonzero. Equation (D.80.11) kills the local
class at some finite level \(W_{r_q}^D\). At five,
the \(d=0\) unit-slope primitive already lies in the
three-jet module \(M\), by D.78. Real cohomology is
zero at this odd prime. Choose
\[
                         R\ge\max(4,\max_{q\in E}r_q).
                                                               \tag{D.80.15}
\]
The primitives originally lying in \(M\) remain primitives
in \(W_R^D\), even where the larger module is ramified.
Thus the single global class \([c]\) is everywhere locally
trivial in this one finite module \(W_R^D\). Its nonzero
image in \(H^1(G_{\mathbb Q},D)\) follows from the actual
coinduction isomorphism and primitive raw Soul\'e restriction,
exactly as in D.78. Therefore \(\mathcal S_s^*=\mathbb F_5\)
when \(d=0\). This proves the first part of (D.80.3)
without commuting a direct limit with an infinite product.

Independence modulo five makes \(a,b,a-b\) all globally
primitive. D.77 therefore gives \(\mathcal T_s=0\), and
its exact local-depth calculation and continuous product
decomposition give the remaining parts of (D.80.3).

\paragraph{Actual ordinary fibres of every finite free rank.}
The inverse in (D.80.4) exists since five is prime to \(36\).
The resulting integer has
\[
 t_d\equiv22\pmod{36},\qquad
 t_d\equiv-1+5^{d+1}\pmod{5^{d+2}}.                    \tag{D.80.16}
\]
It is different from zero and one. The first congruence
gives \(v_2(t_d)=1\), \(v_3(1-t_d)=1\), while
\(v_2(1-t_d)=v_3(t_d)=0\). Thus the two global Kummer
classes are linearly independent modulo five, detected
by actual inertia at two and three. Both are units at five.

For \(d\ge1\), the second congruence gives
\[
 v_5(t_d^4-1)=d+1,\qquad 1-t_d\equiv2\pmod{25}.
\]
The five-adic logarithm preserves the positive valuation
of a principal unit minus one. Since \(v_5(\log6)=1\),
the local coefficient \(\alpha\) has valuation exactly
\(d\), and \(\beta\equiv2\pmod5\). Their difference
is a unit. The pair therefore has exactly the depth
claimed in (D.80.2), proving \(\mathcal K_s=\mathbb Z_5^d\).
The ordinary rational point supplies an actual full
section; its metabelian fibre is nonempty and the
finite-support statement D.74 applies. This proves
the final assertion, without asserting point-theoreticity
of every class in that much larger fibre.
\end{proof}


% D.81. Root, FGH, and PBC complete independent proof/source/certificate audits PASS.
\subsection{Every independent completed unit pair, including nonprimitive local lines}

\begin{statement}{D.81}{Proved}{The full compact kernel for every independent unit pair}
Let \(a,b\in H^1(G_{\mathbb Q},\mathbb Z_5(1))\) have zero
valuations at five and linearly independent reductions modulo five.
Their ramification may be infinite. Write their local unit classes as
\[
 a|_{G_{\mathbb Q_5}}=\alpha\kappa(6),\qquad
 b|_{G_{\mathbb Q_5}}=\beta\kappa(6),
 \qquad \alpha,\beta\in\mathbb Z_5.
                                                               \tag{D.81.1}
\]
There is no primitivity hypothesis on this local vector. Define
\[
 \begin{split}
 \epsilon(\alpha,\beta)&=
 \begin{cases}
  1,&(\bar\alpha,\bar\beta)=(0,0)
       \text{ or }\bar\alpha\bar\beta(\bar\alpha-\bar\beta)\ne0,\\
  0,&\text{otherwise},
 \end{cases}\\
 n(\alpha,\beta)&=
      v_5(\alpha)+v_5(\beta)+v_5(\alpha-\beta),
       \qquad v_5(0)=\infty.
 \end{split}                                                   \tag{D.81.2}
\]
Use \(Z(n)=\mathbb Z_5^n\) at finite \(n\), and
\(Z(\infty)=\prod_{\mathbb N}\mathbb Z_5\), as in D.77.
For the actual Alexander module \(A_s\), \(s=(a,b)\), we have
\[
 \mathcal S_s\simeq(\mathbb Z/5)^{\epsilon(\alpha,\beta)},
 \qquad \mathcal T_s=0,\qquad
 \mathcal K_s\simeq
       (\mathbb Z/5)^{\epsilon(\alpha,\beta)}
                         \times Z(n(\alpha,\beta)).
                                                               \tag{D.81.3}
\]
The last isomorphism is continuous and noncanonical. The finite term
is the inverse finite-global-localization kernel of D.63, not merely
a restriction kernel with divisible local coefficients.
There are thirteen residual vectors giving the order-five term:
the zero vector and the twelve vectors off the three boundary lines.
The twelve nonzero boundary vectors give no such term.

In particular, explicit finitely ramified actual Alexander modules
have compact kernels
\[
  \mathbb Z/5\times\mathbb Z_5^n\quad(n\ge3),
  \qquad
  \mathbb Z/5\times\prod_{\mathbb N}\mathbb Z_5.
                                                               \tag{D.81.4}
\]
Their completed Kummer parameters are given in (D.81.20) below.
These examples are assertions about actual Galois modules and their
cohomology. They do not assert that the corresponding abelian pairs
have metabelian sections. General first-cohomology ranks with
infinite ramification, the general nonseparated kernel, and the
first higher-derived lifting obstructions remain outside this
classification.
\end{statement}

\begin{proof}
\paragraph{The actual images do not require a primitive local line.}
Let \(\Lambda=\mathbb Z_5[[u,v]]\), \(t_1=1+u,t_2=1+v\), and let
\(I_0\) be the actual untwisted operator image in D.78.
The two global hypotheses give two residual inertia vectors away
from five with unit determinant, exactly as in D.80. Their pure
translation images generate \(H=\mathbb Z_5^2\). Hence
\[
                         I_s=H\rtimes I_0.             \tag{D.81.5}
\]
The convergence and exact Kummer interpretation of (D.81.1) are
also those proved in D.80. Apply its integral local monomial gauge
so that (D.81.1) holds for actual cocycles. D.78 proves both
surjectivity of \(G_{\mathbb Q_5}\to I_0\) and
\(\kappa(6)(\ker(G_{\mathbb Q_5}\to I_0))=\mathbb Z_5\).
Consequently, with its actual zero-translation complement,
\[
 I_5=H_5\rtimes I_0,\qquad
 H_5=\mathbb Z_5(\alpha,\beta).                         \tag{D.81.6}
\]
This holds when the line is nonprimitive and when \(H_5=0\).
Indeed every translation coordinate lies in \(H_5\), and the
kernel just described realizes every element of that subgroup,
so a lift of any \(I_0\)-element can be corrected to translation zero.
The gauge preserves (D.81.5), the constant dual submodule, and
all finite augmentation quotients.

Put \(D=A_s^D\), with the Cartier-dual convention of D.78.
Translation coinduction gives
\[
 H^j(H,D)=H^j(H_5,D)=0\quad(j>0).                       \tag{D.81.7}
\]
For the second assertion no saturated complement is presumed.
If \(H_5\ne0\), an integral linear change of coordinates identifies
it with \(5^e\mathbb Z_5\times0\). A continuous transversal is
\(\{0,\ldots,5^e-1\}\times\mathbb Z_5\); thus as a compact
translation space \(H\) is \(H_5\) times that transversal.
The contraction of continuous function-valued cochains in D.78
works with the transversal as an additional coefficient variable.
When \(H_5=0\), the positive-degree assertion is immediate.

Set
\[
 \begin{gathered}
 f=t_1^\alpha t_2^\beta-1,\qquad R=\Lambda/(f),\qquad
 J=(u,v)R,\\
 D_{\rm lin}=R^D=D^{H_5},\qquad
 \lambda=(R/J)^D=\mathbb Q_5/\mathbb Z_5(-1).
 \end{gathered}                                               \tag{D.81.8}
\]
If \(\alpha=\beta=0\), this means \(R=\Lambda\).
The augmentation character on \(R/J\) is \(\chi^2\), which
explains the single inverse cyclotomic character after Cartier duality.
Continuous Hochschild--Serre now identifies restriction with
\[
 H^1(I_0,\lambda)\longrightarrow H^1(I_0,D_{\rm lin}).
                                                               \tag{D.81.9}
\]
The finite affine-image Chebotarev argument of D.78/D.80 applies
without bounded ramification and gives the injection
\[
 \mathcal S_s^*\lhook\joinrel\longrightarrow
 \ker\{H^1(I_0,\lambda)\to H^1(I_0,D_{\rm lin})\}.
                                                               \tag{D.81.10}
\]
Here the star is ordinary Pontryagin duality. That argument is
at a finite coefficient level before taking its direct limit:
a locally trivial finite cocycle is zero on the kernel of its
actual finite operator image, because any nonzero affine value
there would occur at an unramified Frobenius. Thus no unrestricted
coefficient action or formal coefficient assignment is substituted
for the actual image in (D.81.10).

\paragraph{A connecting module valid even at zero or nonprimitive slope.}
Write \(w=t_1t_2-1\), \(P=uvw\), and
\(N_0=\ker(I_0\to\Gamma)\), where \(\Gamma=\mathbb Z_5^\times\).
For every \(\tau\in N_0\), the exact three boundary formulas in
D.39 give
\(\mathcal B_\tau-1\in P\Lambda\).
D.77 supplies an actual local element \(\delta\) for which
\[
              \mathcal B_\delta-1=P E,\qquad
                                  E\in\Lambda^\times. \tag{D.81.11}
\]
Its image in every \(R\) remains a unit. Since the \(N_0\)-action
is multiplication by the corresponding beta series, it follows
on the two actual dual modules that
\[
 \begin{split}
 D_{\rm lin}^{I_0}&=((R/PR)^D)^\Gamma,\\
 (D_{\rm lin}/\lambda)^{I_0}&=((J/PJ)^D)^\Gamma.
 \end{split}                                                 \tag{D.81.12}
\]
The second identity uses \(D_{\rm lin}/\lambda=J^D\).
On these quotient modules the action factors through \(\Gamma\);
no finite-rank hypothesis on either quotient is needed.

The actual ideal inclusion gives the compact exact sequence
\[
 0\longrightarrow C\longrightarrow J/PJ
       \longrightarrow R/PR\longrightarrow\mathbb Z_5(2)
       \longrightarrow0,\qquad
 C=PR/PJ=\frac{R}{J+\operatorname{Ann}_R(P)}.
                                                               \tag{D.81.13}
\]
We compute its first term explicitly, including every zero-divisor case.
The quotients \(\Lambda/(u),\Lambda/(v),\Lambda/(w)\) are domains.
For one of these boundary elements \(L\), if \(f\bmod L\ne0\),
then \(L\) is a nonzero divisor in \(R\): from \(Lr=fh\) in
\(\Lambda\), reduction modulo \(L\) forces \(h=Lh_1\);
cancellation gives \(r=fh_1\).

For a nonzero vector \((\alpha,\beta)\), the three reductions
of \(f\) vanish respectively exactly when
\(\beta=0\), \(\alpha=0\), and \(\alpha=\beta\).
At most one of these equalities occurs. If none occurs, all
three factors of \(P\) are nonzero divisors, so
\(\operatorname{Ann}_R(P)=0\). If one occurs, write
\(f=L S_\eta(L)\), where
\[
 S_\eta(L)=\frac{(1+L)^\eta-1}{L},\qquad S_\eta(0)=\eta\ne0.
 \]
Here \(\eta=\alpha\), \(\beta\), or their common value,
according to the boundary. The other two factors of \(P\)
are nonzero divisors in \(R\); cancellation in \(\Lambda\)
then gives
\(\operatorname{Ann}_R(P)=\operatorname{Ann}_R(L)=(S_\eta(L))\).
This argument also applies when \(\eta\) is a unit.
If \(\alpha=\beta=0\), the original ring \(\Lambda\) is a
domain and \(P\) is again a nonzero divisor. Thus
\[
 C\simeq
 \begin{cases}
  \mathbb Z_5(5),&
       \alpha\beta(\alpha-\beta)\ne0
                     \text{ or }(\alpha,\beta)=(0,0),\\
  (\mathbb Z_5/\eta\mathbb Z_5)(5),&
       \text{a nonzero exact boundary with exponent }\eta .
 \end{cases}                                                 \tag{D.81.14}
\]
To verify the indicated arithmetic character, the multiplier
of \(\rho_0(g)\) has constant term \(\chi(g)^2\), and
\(\varphi_\chi(P)/P\) has constant term \(\chi(g)^3\).
The action on the generator \(P\bmod PJ\) is therefore
\(\chi(g)^5\), including in the finite quotient case.
No division by \(\eta\), or by its five-adic valuation, occurs.

Cartier duality of (D.81.13) is exact and gives
\[
 0\longrightarrow\lambda\longrightarrow (R/PR)^D
       \longrightarrow (J/PJ)^D\longrightarrow C^D
       \longrightarrow0.                                \tag{D.81.15}
\]
The order-four subgroup of \(\Gamma\) acts nontrivially
on \(\lambda\), with unit difference from one. Averaging
and Hochschild--Serre give
\(H^0(\Gamma,\lambda)=H^1(\Gamma,\lambda)=0\).
Taking the two short exact sequences in (D.81.15), and
using (D.81.12), therefore injects the entire restriction
kernel in (D.81.10) into \((C^D)^\Gamma\).
In the infinite case this last group is
\((\mathbb Q_5/\mathbb Z_5(-4))^\Gamma=\mathbb F_5\).
In the finite case it is zero for \(v_5(\eta)=0\), and
has order five for \(v_5(\eta)\ge1\).
Indeed \(2\) topologically generates \(\mathbb Z_5^\times\),
and \(v_5(2^4-1)=1\). In particular
\[
                         |\mathcal S_s|\le5            \tag{D.81.16}
\]
for every pair in the statement. This bounds the full kernel,
not only its elements of order five.

\paragraph{One fixed primitive for every locally five-divisible pair.}
When \((\alpha,\beta)\notin5\mathbb Z_5^2\), D.80 gives
precisely (D.81.3). It remains to prove the nonzero lower bound
when \(\alpha,\beta\in5\mathbb Z_5\).

Let \(c=48h\bmod5\) be the actual integral weight-three cocycle
of D.31, as in D.78. Its coefficients may be viewed in
\(\lambda[5]\), because \(\chi^3=\chi^{-1}\) modulo five.
Its class in \(H^1(I_0,\lambda)\) is nonzero.
Indeed its restriction to \(N_0\) is primitive raw
\(\kappa_3\bmod5\), so its class with \(\lambda[5]\)-coefficients
is nonzero. The map from that cohomology group to
\(H^1(I_0,\lambda)\) is injective since \(\lambda^{I_0}=0\).
Actual coinduction and inflation therefore keep this class
nonzero in \(H^1(G_{\mathbb Q},D)\).

Use the one finite module
\[
 M=(A_s/(5,(u,v)^4)A_s)^D.
                                                               \tag{D.81.17}
\]
The exact local gauge in (D.81.1) makes every local translation
exponent divisible by five. On \(M\) these translations are
the identity, since \(t_i^{5z}=1\) modulo \((5,(u,v)^4)\).
The whole local action on \(M\) is therefore the untwisted one.

Let \(\varepsilon_{ij}\) extract the coefficient of \(u^iv^j\)
in the compact ten-dimensional quotient, with \(i+j\le3\).
The single functional
\[
 m=2\varepsilon_{01}+4\varepsilon_{10}
         +2\varepsilon_{11}+\varepsilon_{20}
         +2\varepsilon_{21}\quad\text{over }\mathbb F_5
                                                               \tag{D.81.18}
\]
satisfies \((g-1)m=c(g)\varepsilon_{00}\) for every actual
untwisted action through this quotient. Here is a finite
coefficient verification, retaining the actual arithmetic
normalization. Its compact operator is
\[
 \begin{split}
 \rho_0(g)F&=S_\chi(t_1)S_\chi(t_2)\mathcal B_g
                         F(t_1^\chi-1,t_2^\chi-1),\\
 \mathcal B_g&\equiv
  1+\frac{\chi^2-1}{24}uv
   +\left(\frac{c(g)}{48}-\frac{\chi^2-1}{48}\right)
                                  (u^2v+uv^2)
                 \pmod{5,(u,v)^4}.
 \end{split}                                                 \tag{D.81.19}
\]
This is exactly the D.31/D.39 beta series; all denominators
are units at five. The coefficients of \(S_\chi\) and of
the substitutions through degree three are binomial
polynomials with unit denominators. On the ten monomials,
expansion gives
\(\chi m-\rho_0(g)^t m=c(g)\chi^2\varepsilon_{00}\)
over \(\mathbb F_5[\chi,c]/(\chi^4-1)\).
More explicitly, before imposing \(\chi^4=1\), the ten
residuals over \(\mathbb F_5[\chi,c]\), in the monomial
order \(1,v,u,v^2,uv,u^2,v^3,uv^2,u^2v,u^3\), are
\[
 \chi(\chi^4-1)
             (0,-2,1,0,-2,-1,0,0,-2,0).
\]
Cartier action is \(\chi\rho_0(g)^{-t}\), so this identity
is precisely the asserted coboundary equation.
The accompanying certificate verifies the ten symbolic
remainders and all twenty residual pairs \((\chi,c)\);
it does not assign any unknown higher beta coefficient.
Thus a single primitive in \(M\) works for the entire
local decomposition group at five.

The finite-level uniformity at all other primes is now
exactly the argument proved in D.80. The set
\(E=\{q\ne5:(a_q,b_q)\notin5\mathbb Z_5^2\}\) is finite.
At every prime outside \(E\cup\{5\}\), \(M\) is unramified
even if the full module is ramified, and D.78's universal
Frobenius primitives apply. At each prime in \(E\),
D.43 and finite local Tate duality, with the finite
Mittag--Leffler comparison in D.80, give
\(H^1(G_{\mathbb Q_q},D)=0\). Its one local class therefore
dies at some larger finite coefficient level.
Taking the maximum of these finitely many levels and
the level containing \(M\) gives one common finite
global coefficient module. The primitive (D.81.18)
handles five in that same module; real cohomology is zero.
The nonzero global class \(c\) consequently belongs to
the direct limit of finite everywhere locally trivial
classes. Equation (D.81.16) gives
\(\mathcal S_s=\mathbb Z/5\).
No direct limit has been interchanged with an infinite
product of local cohomology groups.

Global independence modulo five makes \(a,b,a-b\)
all primitive. D.77 therefore gives
\(\mathcal T_s=0\) and the continuous decomposition in
(D.81.3), with exactly the three local valuations in
(D.81.2). Counting the residual vectors gives twelve
off-boundary vectors, one zero vector, and twelve
nonzero boundary vectors, as stated.

\paragraph{Explicit mixed finite and free compact kernels.}
For \(q\ne5\), put
\(\lambda_q=\log(q^4)/(4\log6)\in\mathbb Z_5\).
We have \(\lambda_{31}\equiv1\pmod5\):
\(31^4\equiv21\pmod{25}\), and \(4\log6\equiv20\pmod{25}\).
For any prescribed \(A,B\in\mathbb Z_5\), define actual
completed Kummer classes
\[
 \begin{split}
 a&=\kappa(7)+
       \frac{A-\lambda_7}{\lambda_{31}}\kappa(31),\\
 b&=\kappa(43)+
       \frac{B-\lambda_{43}}{\lambda_{31}}\kappa(31).
 \end{split}                                                 \tag{D.81.20}
\]
All coefficients are integral. These classes are supported
on \(\{7,31,43\}\), have zero five-valuations, and their
inertia vectors at seven and forty-three are the standard
basis. Their local logarithms are exactly \(A,B\).
For each integer \(n\ge3\), take
\[
                 A=5,\qquad B=5(1+5^{n-3}).
                                                               \tag{D.81.21}
\]
Their two valuations are one, and
\(v_5(A-B)=n-2\). Their residual vector is zero.
Equation (D.81.3) gives the first family in (D.81.4).
Taking \(A=B=0\) gives its countable-product example.
These are actual finitely ramified Galois modules, so
D.74 also computes their unrestricted first cohomology
as \(\prod_{\mathbb N}\mathbb Z_5\).
The existence of a section above an arbitrary one of
these completed abelian pairs is not inferred from
either cohomology computation.
\end{proof}


% D.82. Root, FGH, and PBC full independent proof/source/certificate audits PASS.
\subsection{The valuation direction and the full local Kummer matrix}

\begin{statement}{D.82}{Proved}{Every pair with two independent inertia directions away from five}
Let \(a,b\) be arbitrary completed Kummer classes for \(\mathbb Q\).
Assume that the vectors
\[
             \{(a_q,b_q)\bmod5:q\ne5\}
                         \quad\hbox{span }\mathbb F_5^2.
                                                               \tag{D.82.1}
\]
Their ramification may be infinite, and their five-valuations
\(A=a_5,B=b_5\) need not vanish. Write the exact local classes as
\[
 a=A\kappa(5)+\alpha\kappa(6),\qquad
 b=B\kappa(5)+\beta\kappa(6)
                         \quad\hbox{on }G_{\mathbb Q_5}.
                                                               \tag{D.82.2}
\]
Set
\[
 \begin{gathered}
 W=\operatorname{span}_{\mathbb F_5}
       \{(\bar A,\bar B),(\bar\alpha,\bar\beta)\},\\
 \epsilon(W)=
 \begin{cases}
  1,&W=0\text{ or }W\text{ is a line other than }
       \langle(1,0)\rangle,\langle(0,1)\rangle,\langle(1,1)\rangle,\\
  0,&\text{otherwise},
 \end{cases}\\
 n=\min\{v_5(A),v_5(\alpha)\}
   +\min\{v_5(B),v_5(\beta)\}
   +\min\{v_5(A-B),v_5(\alpha-\beta)\},
 \qquad v_5(0)=\infty .
 \end{gathered}                                               \tag{D.82.3}
\]
For the actual Alexander twist \(s=(a,b)\), the invariants are
\[
 \mathcal S_s\simeq(\mathbb Z/5)^{\epsilon(W)},\qquad
 \mathcal T_s=0,\qquad
 \mathcal K_s\simeq(\mathbb Z/5)^{\epsilon(W)}\times Z(n),
                                                               \tag{D.82.4}
\]
with the compact convention \(Z(n)\) of D.77.
In particular a two-dimensional \(W\) gives
\(\mathcal K_s=0\); a nonboundary line gives
\(\mathcal K_s=\mathbb Z/5\); a boundary line gives the free
compact group \(Z(n)\), \(n\ge1\); and \(W=0\) gives
\(\mathbb Z/5\times Z(n)\), \(n\ge3\).

The actual rational point \(t=10\) has \(\mathcal K_s=0\).
Its nonempty metabelian fibre is a torsor under
\(\prod_{\mathbb N}\mathbb Z_5\).
The general assertion is a compact-kernel calculation,
not a section-existence criterion. D.84--D.85 subsequently remove the away-five spanning
restriction when the complete residual pair is globally independent.
Globally dependent residual pairs, arbitrary infinitely ramified
first-cohomology ranks, and the first higher-derived obstruction
remain outside that completed classification.
\end{statement}

\begin{proof}
\paragraph{The Kummer cocycle of five factors through the actual image.}
Let \(I_0\) be D.78's actual untwisted operator image, and
write \(\kappa_5=\kappa(5)\), to distinguish it from the raw
weight-five Soul\'e character. We first prove that \(\kappa_5\)
is an actual continuous cocycle on \(I_0\).

If \(g\) is in the kernel of the operator representation
\(\rho_0\), comparison on \(1,t_1\) gives \(\chi(g)=1\);
then its beta series is one. D.39's exact raw-character
formula consequently gives
\(\kappa_m(g)=0\) for every odd Soul\'e weight \(m\ge3\).
Use precisely the Kummer definition
\cite[(1.1)--(1.2)]{DE_NakamuraSakugawaWojtkowiak2017}.
For each \(r\ge1\) choose
\[
                 m_r=1+4\cdot5^{r-1}.
                                                               \tag{D.82.5}
\]
For every \(a\) prime to five, \(a^{m_r-1}\equiv1\pmod{5^r}\).
Thus the defining element at this level satisfies
\[
 \prod_{\substack{1\le a<5^r\\5\nmid a}}
      (1-\zeta_{5^r}^{a})^{a^{m_r-1}}
       \equiv
 \prod_{\substack{1\le a<5^r\\5\nmid a}}
      (1-\zeta_{5^r}^{a})
       =\Phi_{5^r}(1)=5
       \pmod{\mathbb Q(\zeta_{5^r})^{\times 5^r}}.
                                                               \tag{D.82.6}
\]
An element of \(\ker\chi\) fixes both that cyclotomic field
and the roots of unity, so the Kummer root choices do not
affect the resulting equality of characters.
It follows that
\(\kappa_5(g)\equiv\kappa_{m_r}(g)=0\pmod{5^r}\) for every \(r\).
Hence \(\kappa_5\) vanishes on \(\ker\rho_0\).
The compact quotient property makes its descended cocycle
on \(I_0\) continuous. This argument takes no limit of
unverified global integral extensions of high-weight cochains:
all the raw characters used are on \(\ker\chi\).

Its image on the local cyclotomic kernel is primitive.
Over \(\mathbb Q_5(\zeta_5)\), normalize
\(v_{\zeta_5-1}(\zeta_5-1)=1\). Then \(v_{\zeta_5-1}(5)=4\).
The Kummer class of five is therefore nonzero modulo fifth
powers and is independent of the class of \(\zeta_5\), whose
valuation is zero. The unique degree-five subextension of
the cyclotomic \(\mathbb Z_5\)-extension has Kummer line
\(\langle\zeta_5\rangle\). Thus a fifth root of five is not
in that cyclotomic extension, and \(\kappa_5\bmod5\) is nonzero
on \(G_{\mathbb Q_5}\cap\ker\chi\). Its compact additive image
there is consequently all of \(\mathbb Z_5\).
D.78's local surjection onto \(I_0\) now proves the actual
quotient statement
\[
 Q=I_0/\{\chi=1,\ \kappa_5=0\}
       \simeq\mathbb Z_5(1)\rtimes\Gamma,\qquad
                  \Gamma=\mathbb Z_5^\times .
                                                               \tag{D.82.7}
\]
Indeed the cyclotomic map is onto and its kernel supplies
every translation, so every projected element can be
corrected to zero translation.

\paragraph{The exact local image and restriction diagram.}
Put \(\mathbf r=(A,B)\), \(\mathbf z=(\alpha,\beta)\), and
\(\Lambda=\mathbb Z_5[[u,v]]\), \(t_1=1+u,t_2=1+v\).
Condition (D.82.1) gives two pure inertia translations with
unit determinant away from five. The faithful actual-image
argument of D.78 therefore gives
\[
                          I_s=\mathbb Z_5^2\rtimes I_0.
                                                               \tag{D.82.8}
\]
After an integral local monomial gauge, (D.82.2) holds for
actual cocycles. D.78 proves that
\(\kappa(6)(\ker(G_{\mathbb Q_5}\to I_0))=\mathbb Z_5\).
The valuation cocycle is zero on that kernel by the first
paragraph. Hence the actual local image consists exactly of
\[
 I_5=\{(h\mathbf z+\kappa_5(g)\mathbf r,g):
                              h\in\mathbb Z_5,\ g\in I_0\}.
                                                               \tag{D.82.9}
\]
It is the semidirect product of \(H_5=\mathbb Z_5\mathbf z\)
with the graph of the descended cocycle
\(\kappa_5\mathbf r\). This graph is an actual subgroup.
It is not replaced by the zero-translation complement.

Let \(D=A_s^D\), and set
\[
 \begin{gathered}
 f_{\mathbf z}=t_1^\alpha t_2^\beta-1,\qquad
 R=\Lambda/(f_{\mathbf z}),\qquad J=(u,v)R,\\
 D_{\rm lin}=R^D,\qquad
 \lambda=(R/J)^D=\mathbb Q_5/\mathbb Z_5(-1).
 \end{gathered}                                               \tag{D.82.10}
\]
The \(I_0\)-action on \(R\) is now the actual twisted action
\[
 \rho_{\mathbf r}(g)
        =t_1^{A\kappa_5(g)}t_2^{B\kappa_5(g)}\rho_0(g).
                                                               \tag{D.82.11}
\]
Continuous coinduction for \(H_5\), including nonprimitive
and zero lines, is proved in D.81. Together with global
translation coinduction it identifies the local restriction
map with the inclusion-induced map
\[
 H^1(I_0,\lambda)\longrightarrow
             H^1(I_0,D_{\rm lin}) .
                                                               \tag{D.82.12}
\]
The coefficient action in the second group is (D.82.11).
The finite affine Chebotarev argument of D.78/D.80 thus gives
\[
 \mathcal S_s^*\lhook\joinrel\longrightarrow
 \mathcal R_s:=
 \ker\{H^1(I_0,\lambda)\to H^1(I_0,D_{\rm lin})\}.
                                                               \tag{D.82.13}
\]

\paragraph{The whole restriction kernel still has order at most five.}
Put \(w=t_1t_2-1\), \(P=uvw\), and
\(N_1=\{\chi=1,\kappa_5=0\}\subset I_0\).
For every element of \(N_1\), (D.82.11) is multiplication
by a beta series congruent to one modulo \(P\).
There is an actual local element \(\delta\in N_1\) with
\(\mathcal B_\delta-1=P E\), \(E\in\Lambda^\times\).
For completeness, on the local cyclotomic kernel the pair
\((\kappa_5,\kappa_3)\) has compact additive image. Conjugation
by an actual element with cyclotomic value
\(\omega\), \(\omega^2=-1\), acts by
\(\operatorname{diag}(\omega,\omega^3)\).
The integral projector
\((T-\omega)/(\omega^3-\omega)\) kills the first coordinate
and preserves a primitive second coordinate. D.77 proves
that the latter occurs locally. The resulting actual element
has the asserted unit multiplier by the same exact three
boundary restrictions and cubic coefficient as D.77.

Consequently
\[
 D_{\rm lin}^{I_0}=((R/PR)^D)^Q,\qquad
 (D_{\rm lin}/\lambda)^{I_0}=((J/PJ)^D)^Q.
                                                               \tag{D.82.14}
\]
The quotient is \(Q\), not just \(\Gamma\).
Its translation element \(h\in\mathbb Z_5(1)\) acts on
the compact quotients by multiplication by
\((t_1^A t_2^B)^h\). Its beta factor disappears modulo
\(P\), or modulo \(PJ\), respectively.

The compact connecting module \(C=PR/PJ\) and its exact
four-term sequence are computed in D.81 without a
primitivity assumption. It is \(\mathbb Z_5(5)\), or
\((\mathbb Z_5/\eta)(5)\) at an exact nonzero boundary of the unit vector.
The additional translation action is trivial on \(C\),
because its augmentation is one.
Also
\[
                      H^0(Q,\lambda)=H^1(Q,\lambda)=0.
                                                               \tag{D.82.15}
\]
To verify the first-cohomology assertion, use
\(1\to\mathbb Z_5(1)\to Q\to\Gamma\to1\).
The subgroup acts trivially on \(\lambda\).
Its \(H^1\) is the continuous homomorphism module
\(\mathbb Q_5/\mathbb Z_5(-2)\), since conjugation on the
source has weight one and the target has weight minus one.
Its \(\Gamma\)-invariants vanish, as do
\(H^0(\Gamma,\lambda)\) and \(H^1(\Gamma,\lambda)\);
the order-four subgroup has a unit difference from one
on each relevant character. The low-degree
Hochschild--Serre sequence proves (D.82.15).

Dualizing D.81's four-term sequence and taking \(Q\)-invariants
therefore injects \(\mathcal R_s\) into
\((C^D)^Q\), a group of order at most five.
This is a bound on the entire kernel, including possible
higher-order classes.

\paragraph{An exact order-five witness if that kernel is nonzero.}
Suppose \(\mathcal R_s\ne0\). In the coefficient long exact
sequence, the subgroup \(D_{\rm lin}^{I_0}=D^{I_5}\) maps
injectively into \((D_{\rm lin}/\lambda)^{I_0}\), since
\(\lambda^{I_0}=0\). D.77 computes \(D^{I_5}\) as a direct
sum of copies of \(\mathbb Q_5/\mathbb Z_5\), so it is
divisible. Its quotient \(\mathcal R_s\) has order five.
Injectivity of divisible abelian groups, or subtraction
of a divisible lift of five times a representative, gives
an order-five representative of its nonzero quotient class.

By (D.82.14), that representative is a continuous
\(\mathbb F_5\)-valued Cartier functional \(\ell\) on \(J/PJ\).
The injection into \(C^D\) shows that \(\ell(P)\ne0\).
Pull it back to the augmentation ideal
\(\mathfrak a=(u,v)\Lambda\). It has the following exact properties:
\[
 \begin{gathered}
 \ell(f_{\mathbf z}\Lambda)=0,\qquad
 \ell(P\mathfrak a)=0,\qquad \ell(P)\ne0,\\
 \ell(f_{\mathbf r}\mathfrak a)=0,\qquad
 f_{\mathbf r}=t_1^A t_2^B-1,\\
 \ell(\rho_0(k)j)=2\ell(j)\quad(j\in\mathfrak a),
       \qquad \chi(k)=2,\quad\kappa_5(k)=0.
 \end{gathered}                                               \tag{D.82.16}
\]
Such an actual \(k\) exists by (D.82.7).
For the fourth identity, use the actual quotient translation
with \(\chi=1,\kappa_5=1\). Its beta factor minus one lies
in \(P\Lambda\), so its effect on the ideal \(J\) is zero
after applying \(\ell\). This gives precisely
\(\ell(f_{\mathbf r}J)=0\).
All evaluations take place modulo five; no division by
a local depth is involved.

\paragraph{The primitive boundary test from D.79.}
We recall its stronger finite-functional form. If
\(\mathbf v\in\mathbb Z_5^2\) is primitive and its reduction
lies on a boundary line, then every order-five
Cartier functional on
\[
 J_{\mathbf v}/P J_{\mathbf v},\qquad
 J_{\mathbf v}=\ker\{\Lambda/(t^{\mathbf v}-1)\to\mathbb Z_5\},
 \]
which is fixed by an actual untwisted \(k\) with \(\chi(k)=2\),
annihilates \(P\).
Here is the exact step of that proof, so the assertion does
not rely on a rational rank computation. The actual boundary
permutations and integral gauges of D.79 preserve the
augmentation ideal and the product boundary ideal.
They reduce to a primitive line with one exponent of
valuation \(d\ge1\) and another a unit. If the line is an
exact boundary, \(P=0\) in its quotient. Otherwise that
quotient is \(\mathbb Z_5[[z]]\) and
\[
 P=Q E,\quad E\in\mathbb Z_5[[z]]^\times,\qquad
 Q=z^2((1+z)^{5^d}-1).
 \]
For the actual power \(g=k^{4\cdot5^{d-1}}\), D.79.20 proves
\[
 (\chi(g)^{-1}\rho_0(g)-1)z^3
          \equiv2z^{5^d+2}\pmod{5,zQ}.
                                                               \tag{D.82.17}
\]
The Cartier-invariant functional annihilates its left side,
and \(Q\equiv z^{5^d+2}\pmod5\). It thus annihilates \(Q\),
and then \(P\), because it annihilates \(Qz\).
This holds for any actual choice of \(k\); the exact boundary
formula and two Taylor coefficients used in D.79 determine
the required operator independently of unknown higher coefficients.

\paragraph{Exclusion for rank two and for a boundary line.}
If \(W=\mathbb F_5^2\), the unit vector \(\mathbf z\)
is primitive. Modulo five its quotient ring is
\(\mathbb F_5[[z]]\), with augmentation ideal \((z)\).
The independent valuation vector makes \(f_{\mathbf r}\)
have nonzero linear coefficient in this ring. Hence
\(f_{\mathbf r}J=(z^2)\), whereas \(P\in J^3\).
The fourth identity of (D.82.16) forces \(\ell(P)=0\),
a contradiction.

Now suppose that \(W\) is a nonzero boundary line.
If \(\mathbf z\) is primitive, the preceding primitive
boundary test applies directly to \(\ell\), again giving
a contradiction. Otherwise \(\mathbf r\) is primitive.
The fourth identity of (D.82.16) makes the restriction of
\(\ell\) to \(f_{\mathbf r}\Lambda\) factor through
\(f_{\mathbf r}\Lambda/f_{\mathbf r}\mathfrak a\).
On this compact cyclic module, \(\rho_0(k)\) has character
\(\chi(k)^3\): its multiplier has augmentation \(\chi^2\),
and \(f_{\mathbf r}(t^\chi)/f_{\mathbf r}(t)\) has
augmentation \(\chi\). Cartier invariance thus gives
\[
                  (2^3-2)\ell(f_{\mathbf r})=0.
 \]
The coefficient six is a five-adic unit. Consequently
\(\ell\) vanishes on \(f_{\mathbf r}\Lambda\), not only
on \(f_{\mathbf r}\mathfrak a\). It descends to the
primitive \(\mathbf r\)-line quotient, where (D.82.17)
forces \(\ell(P)=0\). This contradiction proves
\[
           \mathcal S_s=0
       \quad\text{if \(W\) has rank two or is a boundary line.}
                                                               \tag{D.82.18}
\]

\paragraph{The two nonzero cases have one common finite-level class.}
Let \(c=48h\bmod5\) be D.81's nonzero class with
\(\lambda[5]\)-coefficients. The global actual image
(D.82.8) makes its image in \(H^1(G_{\mathbb Q},D)\)
nonzero, by the same coinduction and inflation argument.
Use \(M=(A_s/(5,(u,v)^4))^D\).
If \(W=0\), every local translation vanishes on \(M\),
and the single functional in D.81.18 is a primitive
for the entire local class at five.

If \(W\) is a nonboundary line, choose an integral lift
\(\mathbf e=(e_1,e_2)\) with
\(e_1e_2(e_1-e_2)\in\mathbb Z_5^\times\).
All local translation vectors modulo five lie in
\(\mathbb F_5\bar{\mathbf e}\).
They therefore disappear on the actual stable quotient
\[
 \Lambda/(5,(u,v)^4,t_1^{e_1}t_2^{e_2}-1)
       \simeq\mathbb F_5[z]/(z^4),\quad
             t_1=(1+z)^{e_2},\quad t_2=(1+z)^{-e_1}.
                                                               \tag{D.82.19}
\]
Its Cartier dual is a submodule of \(M\).
The image of \(P\) has nonzero cubic coefficient
\(e_2(-e_1)(e_2-e_1)\). The unit-line cubic calculation
of D.78 consequently supplies a primitive for \(c\)
in this four-dimensional dual, for the entire local action.
One can see the existence directly: prescribe its top
functional to detect \(P\) with coefficient \(-48\);
the three lower Cartier eigenvalues of an actual
residual \((\chi,c)=(2,0)\) are \(2^{-1},2^{-2},2^{-3}\),
all with unit differences from one. There is a unique
correction in those lower degrees making it fixed by
that element. The residual \((1,1)\) element has exactly
the required constant defect. These two actual residual
operators generate the group with parameters
\((\chi,c)\in\mathbb F_5^\times\ltimes\mathbb F_5(3)\),
so the identity holds for every local element.
The accompanying certificate verifies all three
nonboundary projective lines and all their actual
translation, cyclotomic and cubic residual values.

The common finite-level global argument of D.80 now
applies without change. Away from five only finitely
many inertia vectors are nonzero modulo five.
The universal cubic-module Frobenius primitives handle
all remaining places uniformly; at each exceptional
place D.43 and finite local Tate duality kill the one
class at some larger finite level. Their finite maximum
works globally and preserves the just-constructed
primitive at five. Therefore \(c\) is a nonzero
finite-level everywhere locally trivial class. The
whole-kernel bound proves
\(\mathcal S_s=\mathbb Z/5\) in both remaining cases.

Condition (D.82.1) makes \(a,b,a-b\) all globally
primitive. The local completed Kummer basis is
\(\kappa(5),\kappa(6)\), so their local divisibilities
are exactly the three minima in (D.82.3).
D.77 proves the remaining compact statements and their
continuous product splitting. The four descriptions
by \(W\) follow directly from these minima.

\paragraph{An actual rational point and arbitrary local matrices.}
For \(t=10\), the actual point section has pair
\[
        a=\kappa(2)+\kappa(5),\qquad b=2\kappa(3).
                                                               \tag{D.82.20}
\]
The sign in \(1-t=-9\) vanishes in pro-five Kummer cohomology.
Inertia at two and three has determinant two, a unit.
Its valuation vector is \((1,0)\), while D.78's exact
unit logarithms give
\((\alpha,\beta)\equiv(2,3)\pmod5\).
Thus \(W\) has rank two and \(\mathcal K_s=0\).
The full point section proves nonemptiness, and D.74
computes its finite-support metabelian fibre as the
stated countable-product torsor.

More generally every integral pair of local vectors
\((A,B),(\alpha,\beta)\) occurs in the hypotheses:
add \(A\kappa(5),B\kappa(5)\), respectively, to
D.81.20 with its prescribed local unit values
\(\alpha,\beta\). The away-five inertia vectors at
seven and forty-three remain the standard basis.
This observation constructs actual finite-support
Galois modules in every local matrix stratum; it makes
no further section-existence assertion.
\end{proof}


% D.83. Root, FGH, and PBC full independent proof/source audits PASS.
\subsection{A local image criterion with zero compact kernel}

\begin{statement}{D.83}{Proved}{One local generator for the entire away-five lattice}
Let \(s=(a,b)\) be any completed Kummer pair over \(\mathbb Q\),
with arbitrary infinite ramification, and write its prime coefficients as
\[
 \mathbf v_q=(a_q,b_q)\in\mathbb Z_5^2,
 \qquad \mathbf r=\mathbf v_5.
\]
Define the closed away-five valuation lattice and its local unit vector by
\[
 \begin{split}
 L&=\overline{\sum_{q\ne5}\mathbb Z_5\mathbf v_q}
                                      \ \subset\mathbb Z_5^2,\\
 \lambda_q&=\frac{\log(q^4)}{4\log6}\in\mathbb Z_5,
 \qquad
 \mathbf z=\sum_{q\ne5}\lambda_q\mathbf v_q\in L.
 \end{split}                                                   \tag{D.83.1}
\]
The sums converge in the usual five-adic topology. Suppose
\[
                           L=\mathbb Z_5\mathbf z.              \tag{D.83.2}
\]
The zero lattice is allowed; a nonzero lattice need not be
saturated in \(\mathbb Z_5^2\). For the actual Alexander module
\(A_s\) of D.35 and every finite coefficient quotient
\(W_j=A_s/(5,u,v)^jA_s\), the actual restriction map is an
isomorphism
\[
 H^2(G_{\mathbb Q},W_j)
       \xrightarrow{\ \sim\ }
       \bigoplus_{q\ne5}H^2(G_{\mathbb Q_q},W_j).                 \tag{D.83.3}
\]
Consequently
\[
 \mathcal S_s=\mathcal K_s=0,
 \qquad
 H^2_{\rm cts}(G_{\mathbb Q},A_s)/\overline{\{0\}}
       \xrightarrow{\ \sim\ }\mathcal L_s                       \tag{D.83.4}
\]
topologically, with exactly D.63's local convergence condition.
The actual continuous second cohomology still has D.62's
nonzero nonseparated kernel.

This includes every pair supported only at five and, more generally,
\[
 (a,b)=\mathbf u\,c+\mathbf r\,\kappa(5),                         \tag{D.83.5}
\]
where \(\mathbf u,\mathbf r\in\mathbb Z_5^2\) are arbitrary,
\(c\) has zero five-valuation, its away-five coefficients generate
\(\mathbb Z_5\), and its local unit logarithm coefficient is a unit.
No independence condition on \(\mathbf u,\mathbf r\) is required.
The actual rational point \(t=5\) has this property:
\[
                  (a,b)=(\kappa(5),2\kappa(2)).                  \tag{D.83.6}
\]
Its metabelian fibre is nonempty and is a torsor under
\(\prod_{\mathbb N}\mathbb Z_5\), by D.74.
There is no assertion of section existence for all pairs in
(D.83.2), no general infinitely ramified first-cohomology rank,
and no explicit first higher-derived obstruction.
\end{statement}

\begin{proof}
\paragraph{The infinite sums and their exact cochain representatives.}
The pro-five completion of \(\mathbb Q^\times\) has no sign factor.
In the prime coordinates used in D.35 it is
\(c_0(\{q:q\text{ prime}\},\mathbb Z_5)\). In particular, for
each \(n\), only finitely many vectors \(\mathbf v_q\) fail to
lie in \(5^n\mathbb Z_5^2\). Choose once and for all compatible
roots for each rational prime, giving integral continuous
Kummer cocycles \(\kappa(q)\). They all take values in
\(\mathbb Z_5(1)\), so the sums
\[
 \mathbf s(g)=\mathbf r\kappa(5)(g)
           +\sum_{q\ne5}\mathbf v_q\kappa(q)(g)                 \tag{D.83.7}
\]
converge uniformly in \(g\). They are therefore actual continuous
cocycles, representing the prescribed classes. This argument
does not replace the infinitely ramified class by a finite sum.

Every \(q\ne5\) satisfies \(q^4\in1+5\mathbb Z_5\), and
\(v_5(\log6)=1\). Hence all \(\lambda_q\) in (D.83.1)
are integral and uniformly bounded. The sum defining
\(\mathbf z\) converges and belongs to the closed module \(L\).
The latter is a finitely generated \(\mathbb Z_5\)-module,
although infinitely many primes may contribute to it.

Locally at five, the completed unit class of \(q\) is
\(\lambda_q\kappa(6)\). Thus there is an integral
\(h_q\in\mathbb Z_5\) for which the difference of these two
actual local cocycles is the coboundary of \(h_q\).
The elements \(h_q\) are automatically uniformly bounded.
Therefore
\[
                       \mathbf h=\sum_{q\ne5}\mathbf v_qh_q
                                  \ \in L                       \tag{D.83.8}
\]
converges. A single integral monomial gauge by \(\mathbf h\)
makes the local equality
\[
              \mathbf s(g)=\mathbf z\kappa(6)(g)
                                +\mathbf r\kappa(5)(g)          \tag{D.83.9}
\]
hold for actual cocycles. The gauge lies in \(L\), a point
needed to preserve the global operator subgroup below.

\paragraph{The actual global image and its graph complement.}
Let \(I_0\) be the compact actual untwisted operator image from
D.78. D.82 proves that the actual Kummer cocycle \(\kappa(5)\)
descends continuously to \(I_0\); denote the descended cocycle
by \(k_5\). D.78 also proves that the map
\(G_{\mathbb Q_5}\to I_0\) is surjective and that
\(\kappa(6)\) maps its kernel onto \(\mathbb Z_5\).
These are assertions about the actual arithmetic operators.

The representation of the image of
\[
 g\longmapsto(\mathbf s(g),\rho_0(g))
           \quad\hbox{in }\mathbb Z_5^2\rtimes I_0
\]
by \(t_1^{a(g)}t_2^{b(g)}\rho_0(g)\) is faithful, by D.78:
comparison on \(1,t_1\) determines \(\chi\), the two boundary
axes determine the two translation coordinates, and the
remaining operator determines the untwisted component.
Thus its compact image is the actual operator group \(I_s\).

The untwisted action is unramified at every \(q\ne5\).
Five-primary tame inertia at such a prime has translation
image \(\mathbb Z_5\mathbf v_q\); the chosen generator may
multiply the displayed vector by a unit, which does not
change that image. The closed global image consequently
contains the entire pure translation subgroup \(L\).
Conversely (D.83.7) shows that every translation coordinate,
minus \(\mathbf r k_5\) of its untwisted component, belongs
to \(L\). Surjectivity onto \(I_0\), followed by removal of
this coordinate using the contained subgroup \(L\), proves
the exact compact subgroup identity
\[
 I_s=\{(\mathbf l+\mathbf r k_5(g),g):
                       \mathbf l\in L,\ g\in I_0\}.             \tag{D.83.10}
\]
The set \(\{(\mathbf r k_5(g),g)\}\) is an actual subgroup
because \(k_5\) is a cyclotomic cocycle. It is the graph
complement; it has not been replaced by zero translations.
Conjugation on \(L\) is multiplication by \(\chi\).
A monomial gauge in \(L\), such as (D.83.8), preserves
the subgroup in (D.83.10).

Equation (D.83.9), the local surjection onto \(I_0\),
and the exact image of \(\kappa(6)\) on its kernel now
give the equally exact local identity
\[
 I_5=\{(h\mathbf z+\mathbf r k_5(g),g):
                            h\in\mathbb Z_5,\ g\in I_0\}.      \tag{D.83.11}
\]
Hypothesis (D.83.2) makes (D.83.10) and (D.83.11) equal.
This argument includes nonprimitive lattices and the zero
lattice; no division by a generator of \(L\) occurs.

\paragraph{One decomposition image detects every finite global class.}
We record the finite-level consequence explicitly. Let \(M\)
be any finite continuous module whose global action factors
through \(I_s\), and let \(J\) be its actual finite action
image. A class in
\[
 \operatorname{Sha}^1(\mathbb Q,M)
    =\ker\{H^1(G_{\mathbb Q},M)
                         \longrightarrow\prod_vH^1(G_{\mathbb Q_v},M)\}
\]
has an actual cocycle representative \(z\). This cocycle
vanishes on the kernel of \(G_{\mathbb Q}\to J\).
Indeed, if \(z(g)=m\ne0\) for an element acting trivially
on \(M\), its finite affine image contains \((m,1)\).
The exact finite-image Chebotarev theorem supplies an
unramified prime whose Frobenius is conjugate to that element
\cite[V, Theorem 3.23]{PBC_MilneCFT}. Its action on \(M\)
is the identity, and its cocycle value is nonzero: conjugation
replaces \(m\) by its image under the invertible linear
action of an element of \(J\).
A local coboundary on that Frobenius is zero, a contradiction.
Thus \(z\) descends to a cocycle on \(J\).

Equality \(I_5=I_s\) says that \(G_{\mathbb Q_5}\to J\)
is surjective. Inflation
\(H^1(J,M)\to H^1(G_{\mathbb Q_5},M)\) is injective:
an inflated cocycle which is the coboundary of \(m\)
is already that coboundary on every element of \(J\).
Local triviality at five therefore makes the descended
class zero. We have proved at its original finite level
\[
                       \operatorname{Sha}^1(\mathbb Q,M)=0.
                                                                  \tag{D.83.12}
\]
Also \(M^{G_{\mathbb Q}}=M^{G_{\mathbb Q_5}}\), directly
from equality of their action images.

Apply these facts to the finite Cartier dual \(M=W_j^D\).
Its action factors through \(I_s\), which retains \(\chi\).
Finite global duality identifies
\(\operatorname{Sha}^2(\mathbb Q,W_j)\) with the dual of
\(\operatorname{Sha}^1(\mathbb Q,W_j^D)\), hence it is zero
\cite[I, Theorem 4.10(a)]{DE_MilneADT}.
Finite local duality and equality of the invariant groups
make the component at five
\[
 H^2(G_{\mathbb Q_5},W_j)
            \longrightarrow H^0(G_{\mathbb Q},W_j^D)^*
\]
an isomorphism. The two exact rows of D.63.13 consequently
have zero kernel \(K_j\), proving (D.83.3). In particular
this proves a finite-level isomorphism before taking
any inverse limit. D.63's topological inverse-limit
comparison then gives (D.83.4).

\paragraph{The same conclusion in boundary-depth coordinates.}
For completeness, D.77's explicit three-axis map gives
the same compact vanishing. Let \(e:\mathbb Z_5^2\to\mathbb Z_5\)
be any of the three forms \((x,y)\mapsto x,y,x-y\).
The ideal of global prime coefficients of \(e(s)\) is
generated by \(e(\mathbf r)\) and \(e(L)\). Its local
ideal at five is generated by \(e(\mathbf r)\) and
\(e(\mathbf z)\). They are equal under (D.83.2).
Thus the global and local divisibility exponents agree
on each nonzero axis; on a zero axis the two maps are
the same identity prescribed in D.77. No free compact
factor has been lost in (D.83.4).

\paragraph{Explicit finite and infinite ramification.}
For (D.83.5), write \(c=\sum_{q\ne5}c_q\kappa(q)\).
The condition on its coefficients gives
\(L=\mathbb Z_5\mathbf u\). If its local coefficient
\(\lambda(c)=\sum\lambda_qc_q\) is a unit, then
\(\mathbf z=\lambda(c)\mathbf u\) generates this lattice.
This includes \(\mathbf u=0\) and arbitrary nonprimitive
\(\mathbf u\). It also gives infinitely ramified examples:
for any distinct primes \(q_i\notin\{2,5\}\), take
\[
                       c=\kappa(2)+\sum_{i\ge1}5^i\kappa(q_i).
                                                                  \tag{D.83.13}
\]
The series converges in the completed Kummer space and
\(\lambda(c)\equiv\lambda_2\equiv2\pmod5\), by D.78's
exact logarithm calculation. Its away-five coefficients
generate \(\mathbb Z_5\), as required.

For \(t=5\), the Kummer coordinates are
\(\kappa(5),\kappa(-4)=2\kappa(2)\); the sign has zero
pro-five class. Here \(L=\mathbb Z_5(0,2)\) and
\(\mathbf z=(0,2\lambda_2)\) generates it. The full point
section supplies actual existence, and D.74 supplies
the entire nonempty metabelian fibre. Neither equality
of the operator images nor the separated vanishing is
used to assert arbitrary-pair lifting: the continuous
nonseparated obstruction is still present in general.
\end{proof}


% D.84. Root, FGH, and PBC full proof/source/certificate audits PASS; actual-generator clarification included.
\subsection{The local index of a primitive rank-one global translation lattice}

\begin{statement}{D.84}{Proved}{The whole compact kernel for a primitive rank-one away-five lattice}
Let \(\mathbf u,\mathbf r\) be a basis of \(\mathbb Z_5^2\).
Let \(c=\sum_{q\ne5}c_q\kappa(q)\) be a completed Kummer
class with zero five-valuation whose coefficients generate
\(\mathbb Z_5\). Arbitrary infinite \(c_0\) support is allowed.
Set
\[
 (a,b)=\mathbf u c+\mathbf r\kappa(5),\qquad
 \lambda=\sum_{q\ne5}c_q\frac{\log(q^4)}{4\log6},\qquad
 e=v_5(\lambda),
                                                                  \tag{D.84.1}
\]
where \(v_5(0)=\infty\). For
\(\mathbf r=(r_1,r_2)\), put
\[
 d=\max\{v_5(r_1),v_5(r_2),v_5(r_1-r_2)\}.
                                                                  \tag{D.84.2}
\]
Since \(\mathbf r\) is primitive, either \(d=0\), or exactly
one of these three valuations is positive. Retain the actual
Alexander groups of D.63 and the notation \(Z(n)\) of D.77.
Then
\[
 \mathcal T_s=0,\qquad
 \mathcal S_s\simeq
 \begin{cases}
  \mathbb Z/5,&e>0\text{ and }d=0,\\
  0,&\text{otherwise},
 \end{cases}
 \qquad
 \mathcal K_s\simeq
 \begin{cases}
  0,&e=0,\\
  \mathbb Z/5,&e>0,\ d=0,\\
  Z(\min(e,d)),&e>0,\ d>0.
 \end{cases}                                                     \tag{D.84.3}
\]
The last free factor has countably infinite compact rank when
\(e=d=\infty\). Thus this evaluates a family not covered by
D.82's two independent away-five directions.

The actual point \(t=-24\) has
\[
 (a,b)=(3\kappa(2)+\kappa(3),2\kappa(5)),\qquad
 e=1,\quad d=\infty,\quad \mathcal K_s\simeq\mathbb Z_5.
                                                                  \tag{D.84.4}
\]
Its metabelian fibre is nonempty and is a torsor under
\(\prod_{\mathbb N}\mathbb Z_5\). The theorem makes no
section-existence claim for all the completed parameters,
no general infinitely ramified first-cohomology computation,
and no first higher-derived obstruction computation.
\end{statement}

\begin{proof}
Use \(I_0,k_5,N_1\) from D.82, so
\(N_1=\ker\chi\cap\ker k_5\), the quotient
\(I_0/N_1=\mathbb Z_5(1)\rtimes\Gamma\) is the actual full
affine group, and \(\Gamma=\mathbb Z_5^\times\).
There is an actual element \(k\in I_0\) with
\(\chi(k)=2\) and \(k_5(k)=0\).
Let
\[
 f_{\mathbf w}=t_1^{w_1}t_2^{w_2}-1,\qquad
 P=(t_1-1)(t_2-1)(t_1t_2-1).
\]
The unit statement about the cubic multiplier in D.82 supplies
an actual \(\delta\in N_1\) for which
\(\mathcal B_\delta-1=P E\), \(E\in\Lambda^\times\).
Every element of \(N_1\) has beta multiplier congruent to one
modulo \(P\).

\paragraph{The exact images and the coefficient inclusion.}
D.83's uniformly convergent Kummer cochains and integral local
gauge apply here. The closed away-five lattice is
\(H=\mathbb Z_5\mathbf u\), and the actual two images are
\[
 \begin{split}
 I_s&=\{(h\mathbf u+\mathbf r k_5(g),g):h\in\mathbb Z_5,g\in I_0\},\\
 I_5&=\{(\lambda h\mathbf u+\mathbf r k_5(g),g):
                                      h\in\mathbb Z_5,g\in I_0\}.
 \end{split}                                                     \tag{D.84.5}
\]
Both formulas retain the same actual graph complement. If
\(e=0\), these groups coincide and D.83 proves (D.84.3).
Assume henceforth \(e>0\), allowing infinity.

Put \(D=A_s^D\), and define compact modules and their Cartier
duals by
\[
 B=\Lambda/(f_{\mathbf u}),\qquad
 R=\Lambda/(f_{\lambda\mathbf u}),\qquad
 J=f_{\mathbf u}R,\qquad D_B=B^D,\quad D_R=R^D.
                                                                  \tag{D.84.6}
\]
When \(\lambda=0\), the relation defining \(R\) is absent.
Translation coinduction, with D.81's continuous transversal
when \(\lambda\mathbf u\) is nonprimitive, gives natural
identifications
\[
 H^1(I_s,D)=H^1(I_0,D_B),\qquad
 H^1(I_5,D)=H^1(I_0,D_R).
                                                                  \tag{D.84.7}
\]
If \(\lambda=0\), the second identification simply uses
the graph isomorphism \(I_5\simeq I_0\).
The coefficient action on both compact modules in (D.84.6)
is the actual action
\[
             \rho_{\mathbf r}(g)=
                 t_1^{r_1k_5(g)}t_2^{r_2k_5(g)}\rho_0(g).
                                                                  \tag{D.84.8}
\]
The restriction in (D.84.7) is induced by \(D_B\subset D_R\).
Its entire kernel is therefore
\[
 \mathcal R=
       (J^D)^{I_0}/\operatorname{im}(D_R^{I_0}),
                                                                  \tag{D.84.9}
\]
by the exact sequence \(0\to D_B\to D_R\to J^D\to0\).
The finite affine Chebotarev descent of D.78 gives an injection
\(\mathcal S_s^*\hookrightarrow\mathcal R\). It applies to
each finite globally locally trivial coefficient class before
taking its filtered union, exactly as in D.80 and D.82.

Because \(\mathbf u,\mathbf r\) form a basis, each of the
three boundary linear forms has global coefficient ideal
\(\mathbb Z_5\). Thus D.77 gives
\[
 D_B^{I_0}=D^{I_s}=0,
 \qquad D_R^{I_0}=D^{I_5}\simeq Z(\min(e,d))^*
                       \quad(d>0),                              \tag{D.84.10}
\]
and \(D_R^{I_0}=0\) if \(d=0\).
Here and below a star is ordinary Pontryagin duality after
the Cartier dual has already been taken. In particular
the map from \(D_R^{I_0}\) in (D.84.9) is injective.

\paragraph{The actual compact connecting module.}
On \(N_1\), (D.84.8) is multiplication by beta alone.
The unit cubic element \(\delta\) and the divisibility by
\(P\) for every other element show that
\[
                      (J^D)^{N_1}=(J/PJ)^D.
\]
On this quotient the translation subgroup
\(\mathbb Z_5(1)\subset I_0/N_1\) acts by pure multiplication
\(t_1^{r_1h}t_2^{r_2h}\). Indeed its beta factor is one
on \(J/PJ\). Taking its invariants consequently gives
\[
 (J^D)^{I_0}=(T(-1)_\Gamma)^*,
 \qquad T=J/(PJ+f_{\mathbf r}J).                                 \tag{D.84.11}
\]
The symbol \((-1)\) in this compact formula accounts for
the one Cartier twist: its operator is \(\chi^{-1}\rho\).
Coinvariants are the compact quotient by the closed span of
the operator differences. The \(\Gamma\)-action is the one
induced by actual lifts with \(k_5=0\); differences of lifts
belong to \(N_1\) and have already been killed. Thus (D.84.11)
uses neither a formally assigned beta series nor an assumed
normal coordinate subgroup.

We compute this compact coinvariant module. The actual
permutations of the three punctures, with their transported
frames and the integral gauge in D.79, permit the following
normalization when \(d>0\):
\[
 \mathbf r=(\alpha,\beta),\qquad
 \beta\in\mathbb Z_5^\times,
 \quad v_5(\alpha)=d,
 \quad t_1=t^\beta,\quad t_2=t^{-\alpha},\quad t=1+z.
                                                                  \tag{D.84.12}
\]
The substitution is the exact quotient by \(f_{\mathbf r}\).
The exponent of \(t\) in the image of \(t^{\mathbf u}\)
is \(\beta u_1-\alpha u_2\), a unit. Hence the image of
\(f_{\mathbf u}\) generates \((z)\). Multiplication by its
unit factor merely changes the coordinate of this same
actual ideal. For finite \(n\ge0\), write
\[
 q_n=t^{5^n}-1,\qquad S_n=q_n/z,
 \qquad S_0=1.
                                                                  \tag{D.84.13}
\]
For finite \(e\), the relation obtained after cancelling
the generator \(f_{\mathbf u}\) is \(S_e\), up to a unit.
This remains true for a unit factor in \(\lambda/5^e\):
\((X^\epsilon-1)/(X-1)\) has unit constant coefficient
when \(\epsilon\in\mathbb Z_5^\times\).
For \(e=\infty\) that relation is absent.

In the quotient (D.84.12), \(P\) is a unit times
\(z^2q_d\) when \(0<d<\infty\), it is a unit times
\(z^3\) when \(d=0\), and it is zero when \(d=\infty\).
Thus, as a compact module with its actual action on the
ideal generated by \(z\), (D.84.11) has the scalar description
\[
 T\simeq
 \begin{cases}
  \mathbb Z_5[[z]]/(S_e,z^3),&d=0,\\
  \mathbb Z_5[[z]]/(S_e,z^2q_d),&0<d<\infty,\\
  \mathbb Z_5[[z]]/(S_e),&d=\infty.
 \end{cases}                                                     \tag{D.84.14}
\]
An occurrence of \(S_e\) is omitted for \(e=\infty\).
The scalar description does not discard the generator's
extra linear weight.

\paragraph{The exact finite boundary module.}
For \(n\ge1\), let
\(C_n=\mathbb Z/5^n\), let \(R_n=\mathbb Z_5[C_n]\),
and let \(I_n\) be its augmentation ideal. The multiplication
action \(t\mapsto t^\chi\) gives
\[
                        (I_n^2)_\Gamma\simeq\mathbb Z_5^n.
                                                                  \tag{D.84.15}
\]
Here is an integral verification. The quotient
\(I_n/I_n^2\simeq\mathbb Z/5^n(1)\) has nontrivial
order-four character. Exact order-four coinvariants therefore
identify the coinvariants of \(I_n^2\) and \(I_n\).
The augmentation sequence for \(I_n\subset R_n\) splits
\(\Gamma\)-equivariantly by the group-ring element at zero.
The \(\Gamma\)-orbits in \(C_n\) are zero and the \(n\)
nonzero valuation shells. The coinvariants of the permutation
module \(R_n\) are consequently free on \(n+1\) orbit classes,
and their augmentation kernel is free of rank \(n\).
This proves (D.84.15), with no rational rank inference.

Whenever a quotient in (D.84.14) is
\(Q_n=\mathbb Z_5[[z]]/(S_n)\), its actual Cartier-untwisted
action identifies it with \(I_n^2\). To see the precise map,
regard its scalar element as an element of the ideal generated
by \(f_{\mathbf u}\) and multiply by \(t_1-1\).
Both factors are unit multiples of \(z\). In
\(R_n=\mathbb Z_5[[z]]/(q_n)\), multiplication by their
product identifies \(Q_n\) with \(I_n^2\).
Injectivity follows without dividing by \(5\): if
\(z^2f\in zS_n\mathbb Z_5[[z]]\), then
\(zf=S_nh\); augmentation gives \(5^nh(0)=0\), so
\(h=zh_1\), and cancellation gives \(f=S_nh_1\).
On this quotient \(t_2=1\), so the exact Fox multiplier
is \(\chi S_\chi(t_1)\). Multiplication by \(t_1-1\)
changes it to \(\chi\) times substitution; the Cartier
untwist removes \(\chi\). This proves the asserted
\emph{actual} equivariance with the permutation action.

If \(e\le d\), with finite \(e>0\), then \(q_d\) is
divisible by \(q_e\), or \(P=0\) when \(d=\infty\).
The extra relation in (D.84.14) is redundant. The preceding
comparison and (D.84.15) give
\[
                         T(-1)_\Gamma\simeq\mathbb Z_5^e.
                                                                  \tag{D.84.16}
\]

\paragraph{The three-jet extension when \(0<d<e\).}
Assume \(d\) finite and positive, and \(e>d\), including
\(e=\infty\). There is an exact sequence of compact actual
\(\Gamma\)-modules
\[
                 0\longrightarrow K\longrightarrow T
                       \longrightarrow Q_d\longrightarrow0.
                                                                  \tag{D.84.17}
\]
Before the Cartier untwist, its kernel is
\[
 K\simeq
 \begin{cases}
  (\mathbb Z/5^{e-d})[z]/(z^3),&e<\infty,\\
  \mathbb Z_5[z]/(z^3),&e=\infty,
 \end{cases}                                                     \tag{D.84.18}
\]
with actual generator \(q_d\) in the ideal under discussion.
Indeed \(S_e=S_d g\), where
\(g=\sum_{j=0}^{5^{e-d}-1}t^{j5^d}\), and
\(z^2q_d=z^3S_d\). Cancelling the nonzero element \(S_d\)
in the power-series domain identifies the kernel with
\(\mathbb Z_5[[z]]/(g,z^3)\). Through degree two the
coefficients of \(g\) are all divisible by \(5^{e-d}\),
and its constant quotient is one. This follows directly
by summing \(1,j,\binom j2\), whose denominators two
and six are units at five. Thus \(g\) is
\(5^{e-d}\) times a unit modulo \(z^3\), proving (D.84.18).
The infinite case follows by omitting \(S_e\).

On \(q_d,q_dz,q_dz^2\), the actual compact action has
diagonal characters \(\chi^3,\chi^4,\chi^5\).
The factor \(\chi^2\) is the augmentation of the actual
Alexander multiplier, the factor \(\chi\) comes from
\(q_d(t^\chi)/q_d(t)\), and the remaining powers come
from the displayed monomials. Hence \(K(-1)\) has
diagonal characters \(\chi^2,\chi^3,\chi^4\).
For the actual element with \(\chi=2\), their differences
are units pairwise, so the triangular operator is
diagonalized integrally. The first two eigenvalues minus
one are units; the third, \(2^4-1\), has valuation one.
Since two topologically generates \(\Gamma\),
\[
                         K(-1)_\Gamma\simeq\mathbb F_5,
                                                                  \tag{D.84.19}
\]
generated by the class of \(q_dz^2\).

This possible torsion is killed in the coinvariants of
\(T(-1)\). Set \(g=k^{4\cdot5^{d-1}}\). Its multiplier
is one modulo \(5^d\), so it acts \emph{exactly} as the
identity on \(Q_d(-1)\simeq I_d^2\), by the permutation
description above. Therefore
\((\chi(g)^{-1}\rho_0(g)-1)z^3\) belongs to \(K(-1)\).
Here \(z^3\) is an element of the actual ideal \(J\);
its scalar representative in (D.84.14) is \(z^2\).
Thus the scalar relation \(z^2q_d\) corresponds to
the actual-ideal relation \(z^3q_d\) used below.
D.79's exact actual-power identity gives
\[
 (\chi(g)^{-1}\rho_0(g)-1)z^3
             \equiv 2z^{5^d+2}
                  \pmod{5,\ z^3q_d}.                           \tag{D.84.20}
\]
Changing the primitive coordinate or boundary frame multiplies
the indicated top coefficient by a unit, which suffices here.
Since \(q_dz^2\equiv z^{5^d+2}\pmod5\), this is a nonzero
class of the one-dimensional module (D.84.19). It is an actual
operator difference in \(T(-1)\), so the map from (D.84.19)
to \(T(-1)_\Gamma\) is zero. Right exactness of compact
coinvariants in (D.84.17) now gives
\[
                         T(-1)_\Gamma\simeq\mathbb Z_5^d.
                                                                  \tag{D.84.21}
\]
This step uses an exact identity on the finite boundary
quotient, not merely a congruence of its eigenvalues.

\paragraph{Off the boundary and at the simultaneous infinite limit.}
If \(d=0\), (D.84.14) is just three jets, reduced modulo
\(5^e\) when \(e\) is finite. Its generator has the same
three characters as (D.84.18), so
\[
                         T(-1)_\Gamma\simeq\mathbb F_5.
                                                                  \tag{D.84.22}
\]
There is no finite boundary quotient and no relation of
the kind used in (D.84.20).

If \(e=d=\infty\), the natural map in (D.84.9) must be
checked, rather than comparing two infinite coranks.
Use the actual boundary permutation to take
\(\mathbf r=(0,\beta)\), \(\beta\) a unit.
The translation relation is \(t_2=1\).
The map \(D_R^{I_0}\to(J^D)^{I_0}\), after the natural
axis identifications of D.77 and multiplication by
\(t_1-1\), is dual to the inclusion
\[
 I^2\lhook\joinrel\longrightarrow I,
 \qquad I=(t-1)\mathbb Z_5[[t-1]],
                                                                  \tag{D.84.23}
\]
with pure substitution \(t\mapsto t^\chi\) on both compact
modules. Their quotient \(I/I^2\simeq\mathbb Z_5(1)\)
has nontrivial order-four character. Exact order-four
coinvariants make (D.84.23) an isomorphism on coinvariants,
and hence on full \(\Gamma\)-coinvariants. Its invariant
dual map is consequently an isomorphism. This proves
\(\mathcal R=0\) also in this infinite case.

For finite \(\min(e,d)>0\), (D.84.16)--(D.84.21) show
that \((J^D)^{I_0}\) is a divisible group of that finite
corank. The injective map in (D.84.10) has source a
divisible group of the same corank, so it is an isomorphism.
Indeed its Pontryagin dual is a surjection between free
\(\mathbb Z_5\)-modules of the same finite rank, hence an
isomorphism. Together with (D.84.22) and
\(D_R^{I_0}=0\) when \(d=0\), this computes the whole
restriction kernel:
\[
 \mathcal R\simeq\mathbb F_5\quad(e>0,d=0),
 \qquad \mathcal R=0\quad\text{otherwise}.
                                                                  \tag{D.84.24}
\]

\paragraph{The nonzero class and a common finite global level.}
Let \(c_3\) denote D.78's normalized raw cubic cocycle
modulo five, viewed in the constant Cartier submodule
\(\mathbb F_5(-1)\subset D_B\).
Its global class is nonzero in the present, nonconstant
coinduced coefficient. To prove this directly, choose
an actual \(\tau\in I_0\) with
\(\chi(\tau)=1,k_5(\tau)=1,c_3(\tau)=0\).
Such an element exists: first prescribe the two affine
coordinates, then multiply by a suitable element of
\(N_1\), on which the cubic character is primitive.
On \(B\simeq\mathbb Z_5[[z]]\), its compact operator is
multiplication by \(t^{\mathbf r}\mathcal B_\tau\).
The difference from one is \(z\) times a unit, because
\(\mathbf u,\mathbf r\) are a basis and beta starts in
degree three on the cyclotomic kernel.
Thus the \(\tau\)-fixed part of \(D_B\) is precisely the
constant functional module. If \(c_3\) were the coboundary
of an element of \(D_B\), that element would be fixed
by \(\tau\), hence constant. An element of \(N_1\) with
unit cubic character fixes constants and contradicts
the coboundary identity. This proves nonvanishing in
\(H^1(I_0,D_B)\), hence by coinduction and inflation in
\(H^1(G_{\mathbb Q},D)\).

Suppose \(e>0,d=0\). Modulo five the local unit vector
\(\lambda\mathbf u\) is zero, and the valuation vector
\(\mathbf r\) is a nonboundary direction. D.82's actual
four-dimensional line quotient of the fixed cubic module
therefore gives a local primitive for \(c_3\) at five.
This local construction uses only the residual vectors
and the actual beta coefficients; it does not require
D.82's global away-five rank-two hypothesis.

At every prime away from five with nonzero abelian
inertia, D.80's finite local Tate-duality argument and
D.43 give \(H^1(G_{\mathbb Q_q},D)=0\).
Only finitely many such primes have nonzero residual
inertia modulo five. At those primes a primitive occurs
in some finite coefficient module, and one takes a
common larger level for this finite set. At every
other prime, the original ten-dimensional cubic module
is unramified and D.78's explicit Frobenius primitive
works at the same fixed level. The real place contributes
no odd-primary cohomology. This is precisely D.80's
finite-exception argument; it gives a single finite
global coefficient module in which \(c_3\) is locally
trivial everywhere. Its nonzero image just proved gives
\(\mathcal S_s^*\ne0\). The upper bound (D.84.24)
then proves the finite assertion in (D.84.3).

\paragraph{The whole compact kernel and actual examples.}
The global three boundary depths are zero. If \(d=0\),
the local valuation vector alone makes every local
depth zero. If \(d>0\), there is one boundary form
whose value on \(\mathbf r\) has valuation \(d\);
its value on \(\mathbf u\) is a unit because these
vectors are a basis. Its local depth is therefore
\(\min(e,d)\), and the other two are zero. D.77 gives
all the remaining conclusions of (D.84.3), continuously.

For \(t=-24\), the global Kummer pair is exactly
(D.84.4), since signs vanish in pro-five completion.
Its away-five coefficients \(3,1\) generate \(\mathbb Z_5\),
and the basis vectors are \((1,0),(0,2)\).
The exact identity
\[
                      24^4-1=25\cdot13271
\]
shows that its unit logarithm coefficient has valuation
one. The actual full point section supplies existence;
D.74 then computes its nonempty metabelian fibre.

For explicit completed modules at every index, let
\(\lambda_2,\lambda_3\) be the units in D.78 and set
\[
 c_m=\kappa(2)+\frac{5^m-\lambda_2}{\lambda_3}\kappa(3)
       \quad(m\ge1),\qquad
 c_\infty=\kappa(2)-\frac{\lambda_2}{\lambda_3}\kappa(3).
                                                                  \tag{D.84.25}
\]
These classes are globally primitive and have the stated
local coefficient \(5^m\), or zero.
Choosing \(\mathbf u=(1,0),\mathbf r=(0,1)\) realizes
\(Z(m)\), including the infinite product, while
\(\mathbf r=(1,2)\) realizes \(\mathbb Z/5\).
They are actual finite-support Galois modules; no further
section existence is inferred.

The audited certificate is reproduced by
\path{code/de_sha_rank_one_index.py}; its full output is retained
and agrees with the separately preserved candidate execution.
\end{proof}


% D.85. Root, FGH, and PBC full independent proof/source/certificate audits PASS.
\subsection{Every globally primitive rank-two Kummer pair}

\begin{statement}{D.85}{Proved}{The whole compact kernel for every globally independent residual pair}
Let \(s=(a,b)\) be arbitrary completed Kummer classes over
\(\mathbb Q\), with arbitrary infinite ramification, and assume
\[
 \operatorname{span}_{\mathbb F_5}
       \{(a_q,b_q)\bmod5:q\text{ prime}\}=\mathbb F_5^2.
                                                               \tag{D.85.1}
\]
The prime five is included in this hypothesis. Write the actual
local classes at five as
\[
 a=A\kappa(5)+\alpha\kappa(6),\qquad
 b=B\kappa(5)+\beta\kappa(6),
 \qquad \mathbf r=(A,B),\quad\mathbf z=(\alpha,\beta).
                                                               \tag{D.85.2}
\]
Put
\[
 \begin{gathered}
 W=\operatorname{span}_{\mathbb F_5}\{\overline{\mathbf r},
                                      \overline{\mathbf z}\},\\
 \epsilon(W)=
 \begin{cases}
 1,&W=0\text{ or }W\text{ is a line other than }
 \langle(1,0)\rangle,\langle(0,1)\rangle,\langle(1,1)\rangle,\\
 0,&\text{otherwise},
 \end{cases}\\
 n=\min\{v_5(A),v_5(\alpha)\}
  +\min\{v_5(B),v_5(\beta)\}
  +\min\{v_5(A-B),v_5(\alpha-\beta)\},\qquad v_5(0)=\infty.
 \end{gathered}                                                \tag{D.85.3}
\]
For D.63's actual Alexander invariants and D.77's compact
convention \(Z(n)\), the complete answer is
\[
 \mathcal S_s\simeq(\mathbb Z/5)^{\epsilon(W)},\qquad
 \mathcal T_s=0,\qquad
 \mathcal K_s\simeq(\mathbb Z/5)^{\epsilon(W)}\times Z(n).
                                                               \tag{D.85.4}
\]
Thus D.82's classification holds under global residual
independence, without requiring two independent directions
away from five. The splitting in (D.85.4) is continuous and
noncanonical.

The new finite-index case includes the actual point
\[
 t=-799,\qquad
 (a,b)=(\kappa(17)+\kappa(47),\,5\kappa(2)+2\kappa(5)),
 \qquad \mathcal K_s\simeq\mathbb Z_5.
                                                               \tag{D.85.5}
\]
Its metabelian fibre exists and is a torsor under
\(\prod_{\mathbb N}\mathbb Z_5\), by D.74. No existence is
asserted for every pair in (D.85.1). Globally dependent
residual pairs, arbitrary infinitely ramified first-cohomology
ranks, and the first higher-derived obstruction are not
computed by this statement.
\end{statement}

\begin{proof}
Throughout, \(\Lambda=\mathbb Z_5[[t_1-1,t_2-1]]\), and
\[
 f_{\mathbf v}=t_1^{v_1}t_2^{v_2}-1,
 \qquad P=(t_1-1)(t_2-1)(t_1t_2-1).
\]
Use the actual untwisted image \(I_0\), its descended Kummer
cocycle \(k_5\), and \(N_1=\ker\chi\cap\ker k_5\) from
D.82. In particular \(I_0/N_1=\mathbb Z_5(1)\rtimes\Gamma\),
where \(\Gamma=\mathbb Z_5^\times\), is the actual quotient.
The compact action has augmentation \(\chi^2\).
There is an actual element of \(N_1\) whose beta multiplier
minus one is \(P\) times a unit. Every element of \(\ker\chi\)
has beta multiplier congruent to one modulo \(P\).

\paragraph{The only remaining shape of the away-five lattice.}
Let
\[
 L=\overline{\sum_{q\ne5}\mathbb Z_5(a_q,b_q)},
 \qquad
 \mathbf z=\sum_{q\ne5}\lambda_q(a_q,b_q),
 \qquad \lambda_q=\frac{\log(q^4)}{4\log6}.
                                                               \tag{D.85.6}
\]
These are D.83's closed lattice and convergent local unit
vector. If \(L\bmod5\) has rank two, D.82 applies. It
cannot have rank zero under (D.85.1), because the only
remaining vector would be \(\mathbf r\). Suppose therefore
that it has rank one. Then \(\overline{\mathbf r}\) is
independent of that line. Choose \(\mathbf u\in L\) whose
reduction generates the line. The vectors
\(\mathbf u,\mathbf r\) form a basis of \(\mathbb Z_5^2\).

Projection onto the \(\mathbf r\)-coordinate in this basis
sends \(L\) to an ideal contained in \(5\mathbb Z_5\).
If it is zero, then \(L=\mathbb Z_5\mathbf u\), and D.84
applies. Otherwise it is \(5^h\mathbb Z_5\), \(h\ge1\).
Subtracting the \(\mathbf u\)-coordinate of an element
with primitive projected image, and multiplying by a unit,
gives the exact decomposition
\[
 L=\mathbb Z_5\mathbf u\oplus5^h\mathbb Z_5\mathbf r,
 \qquad h\ge1.
                                                               \tag{D.85.7}
\]
Thus it remains to treat this finite-index, nonsaturated
away-five lattice.

Write each prime vector uniquely as
\((a_q,b_q)=c_q\mathbf u+5^h d_q\mathbf r\).
The sequences \((c_q),(d_q)\) are \(c_0\): division by
the fixed scalar \(5^h\) preserves convergence to zero.
Consequently the actual Kummer sums
\[
 c=\sum_{q\ne5}c_q\kappa(q),\quad
 d_0=\sum_{q\ne5}d_q\kappa(q),\quad
 (a,b)=\mathbf u c+\mathbf r(\kappa(5)+5^h d_0)
                                                               \tag{D.85.8}
\]
converge uniformly. Put
\[
 \begin{gathered}
 \lambda=\sum_{q\ne5}\lambda_qc_q,\qquad
 \nu=\sum_{q\ne5}\lambda_qd_q,\qquad
 \mathbf z=\lambda\mathbf u+5^h\nu\mathbf r,\\
 e=v_5(\lambda),\qquad
 d=\max\{v_5(r_1),v_5(r_2),v_5(r_1-r_2)\}.
 \end{gathered}
                                                               \tag{D.85.9}
\]
Here \(d_0\) denotes the Kummer class, while \(d\) is
the boundary depth. The latter is zero unless the primitive
vector \(\mathbf r\) has exactly one boundary form divisible
by five; that form then has depth \(d\), possibly infinite.

\paragraph{The actual images and the natural restriction map.}
D.83's integral gauge is a uniformly convergent sum in
\(L\). Its image computation therefore applies verbatim,
including infinite ramification, and gives
\[
 \begin{split}
 I_s&=\{(\mathbf l+\mathbf r k_5(g),g):
                               \mathbf l\in L,\ g\in I_0\},\\
 I_5&=\{(j\mathbf z+\mathbf r k_5(g),g):
                               j\in\mathbb Z_5,\ g\in I_0\}.
 \end{split}                                                   \tag{D.85.10}
\]
Both retain the same actual graph complement. They are
the faithful compact operator images, not abstract groups
with an assigned action.

Let \(D=A_s^D\) be the discrete Cartier dual and define
\[
 I_L=(f_{\mathbf u},f_{5^h\mathbf r}),\qquad
 B=\Lambda/I_L,\qquad R=\Lambda/(f_{\mathbf z}),\qquad
 J=I_LR.
                                                               \tag{D.85.11}
\]
The defining relation for \(R\) belongs to \(I_L\).
The actual \(I_0\)-action on both coefficient quotients is
\[
             \rho_{\mathbf r}(g)=t^{\mathbf r k_5(g)}\rho_0(g).
                                                               \tag{D.85.12}
\]
Translation coinduction gives natural identifications
\[
 H^1(I_s,D)=H^1(I_0,B^D),\qquad
 H^1(I_5,D)=H^1(I_0,R^D).
                                                               \tag{D.85.13}
\]
For the first identity, one may choose the finite coset
transversal of \(L\) in \(\mathbb Z_5^2\); the continuous
function model for \(D\) is then coinduced for \(L\).
For the second, the line transversal, including nonprimitive
or zero \(\mathbf z\), is the one constructed in D.81.
In both cases the translation group's higher cohomology
vanishes, so the Hochschild--Serre edge map gives (D.85.13).
These function identifications show that restriction is
exactly induced by \(B^D\hookrightarrow R^D\); no unrelated
choice of the two coinduction isomorphisms is made.

The exact sequence dual to \(0\to J\to R\to B\to0\)
therefore identifies the whole restriction kernel with
\[
 \mathcal R_s=(J^D)^{I_0}/\operatorname{im}((R^D)^{I_0}).
                                                               \tag{D.85.14}
\]
Finite affine Chebotarev descent, at each finite coefficient
level as in D.78 and D.83, injects
\(\mathcal S_s^*\) into (D.85.14).
Every global boundary depth is zero by (D.85.1), so
\(D^{I_s}=(B^D)^{I_0}=0\) by D.77. In particular the
map in the denominator of (D.85.14) is injective.

\paragraph{An exact two-generator module calculation.}
The \(\mathbb Z_5\)-basis \(\mathbf u,\mathbf r\) gives
genuine completed group-ring coordinates
\[
 x=f_{\mathbf u},\qquad y=f_{\mathbf r},\qquad
 \Lambda=\mathbb Z_5[[x,y]],\qquad
 g_h=(1+y)^{5^h}-1,
 \quad I_L=(x,g_h).
                                                               \tag{D.85.15}
\]
Let \(V=I_L/yI_L\). There is an exact topological
\(\mathbb Z_5[[x]]\)-module identification
\[
 V=x\mathbb Z_5[[x]]\oplus\mathbb Z_5\,g_h,
 \qquad xg_h=0.
                                                               \tag{D.85.16}
\]
Indeed \(g_h=yS_h(y)\), where \(S_h(0)=5^h\ne0\).
In \(V\), both \(yg_h\) and \(xg_h\) vanish; the latter
belongs to \(yx\Lambda\). Every element thus has the
form in (D.85.16). To prove uniqueness, suppose
\(xf(x)+a g_h\in y(x,g_h)\). Evaluation at \(y=0\)
gives \(f=0\). Cancel \(y\) and set \(x=0\); then
\(aS_h(y)=yS_h(y)b(y)\). Cancellation in the power-series
domain gives \(a=yb(y)\), hence \(a=0\). The map and its
inverse are continuous, since they are continuous maps
between compact Hausdorff modules.

The exact local relation has the form
\[
 f_{\mathbf z}=(1+x)^\lambda(1+g_h)^\nu-1.
                                                               \tag{D.85.17}
\]
Modulo \(yI_L\), the products \(g_h^2\) and \(xg_h\)
vanish, so its image is
\[
                 f_\lambda(x)+\nu g_h,
 \qquad f_\lambda(x)=(1+x)^\lambda-1.
                                                               \tag{D.85.18}
\]
This identity holds for every \(\lambda,\nu\in\mathbb Z_5\)
by the convergent integral binomial series, not only for
ordinary nonnegative exponents.

\paragraph{The extra generator has a nontrivial tame character.}
The actual \(N_1\) multiplier and the valuation-translation
subgroup of \(I_0/N_1\) give exactly as in D.84
\[
 (J^D)^{I_0}=(T(-1)_\Gamma)^*,
 \qquad T=J/(PJ+yJ).
                                                               \tag{D.85.19}
\]
The Cartier untwist means that its compact operator is
\(\chi^{-1}\rho\). Coinvariants use the closed span of
actual operator differences.

Let \(G\) be the image of \(\mathbb Z_5g_h\) in \(T\).
This is a cyclic compact quotient, possibly zero or finite.
It is stable under the actual \(\Gamma\)-action and has
\emph{exact} character \(\chi^3\). To verify this, use
lifts with \(k_5=0\). Their compact multiplier has
augmentation \(\chi^2\); every positive augmentation term
annihilates \(g_h\) in (D.85.16). Substitution sends
\[
 g_h\longmapsto (1+g_h)^\chi-1\equiv\chi g_h
                                      \pmod{yI_L}.
\]
Thus the multiplier and substitution together give
\(\chi^3g_h\), independently of all higher actual beta
coefficients. Consequently \(G(-1)\) has character
\(\chi^2\). Its coinvariants vanish, since an actual lift
with \(\chi=2\) acts minus the identity by the unit three.

Evaluation at \(y=0\), followed by quotienting out this
generator, gives a natural exact sequence
\[
 0\longrightarrow G\longrightarrow T\longrightarrow T_0
                       \longrightarrow0,
 \qquad
 T_0=\frac{x\mathbb Z_5[[x]]}
 {f_\lambda(x)\mathbb Z_5[[x]]+
                        P(x,0)x\mathbb Z_5[[x]]}.
                                                               \tag{D.85.20}
\]
Here the product by \(P\) kills the cyclic summand of
(D.85.16), because \(P\) has zero augmentation. Equations
(D.85.16)--(D.85.18) prove exactness, including when
\(\lambda=0\), \(\nu=0\), or \(\nu\) is nonprimitive.
The map is equivariant for the actual action restricted to
\(f_{\mathbf r}=0\). It is precisely D.84's rank-one
connecting ideal with the same basis \(\mathbf u,\mathbf r\)
and local coefficient \(\lambda\).

Compact coinvariants are right exact, and the kernel just
computed has zero coinvariants. Hence the natural map gives
\[
                    T(-1)_\Gamma\simeq T_0(-1)_\Gamma.
                                                               \tag{D.85.21}
\]
This removes the extra finite-index direction integrally.
It does not discard the whole \(g_h\) term before taking
coinvariants or replace a group action by its rational rank.

\paragraph{Local depths and the actual invariant comparison.}
The three local boundary depths are determined by
\(\mathbf r\) and \(\mathbf z\). If \(d=0\), they are
all zero. If \(d>0\), exactly one boundary form, say
\(\ell\), has \(v_5(\ell(\mathbf r))=d\); its value on
\(\mathbf u\) is a unit. The identity
\[
 \ell(\mathbf z)=\lambda\ell(\mathbf u)
                       +5^h\nu\ell(\mathbf r)
\]
therefore gives, including cancellation and infinite values,
\[
 \min\{v_5(\ell(\mathbf r)),v_5(\ell(\mathbf z))\}
                         =\min(d,e).
                                                               \tag{D.85.22}
\]
The second summand cannot change this minimum; when \(d\)
is finite it has valuation at least \(h+d>d\).
All other boundary depths are zero. D.77 consequently gives
\((R^D)^{I_0}\simeq Z(\min(e,d))^*\) for \(d>0\),
and zero for \(d=0\).

D.84's full computation of the actual module \(T_0\),
including its three-jet extension and actual cyclotomic
power relation, now applies to (D.85.21). If \(e=0\),
the module \(T_0\) is zero. If \(e>0,d=0\), its
coinvariants are \(\mathbb F_5\), whereas the local
invariants are zero. If \(0<\min(e,d)<\infty\), both
invariant groups in (D.85.14) are divisible of that same
finite corank. The already proved injection between them
is an isomorphism: its dual is a surjection between free
\(\mathbb Z_5\)-modules of the same finite rank.

The case \(e=d=\infty\) requires the natural map, not
this finite-corank argument. Use the actual boundary
permutation and frame transport of D.79 to arrange
\(\mathbf r=(0,\beta)\), \(\beta\) a unit. Then
\(\lambda=0\) and \(\mathbf z=5^h\nu\mathbf r\).
The unique nonzero local invariant axis of D.77 is the
quotient \(t_2=1\), whose map is induced by the actual
surjection \(R\to\Lambda/(t_2-1)\). On this quotient,
\(g_h=0\) and \(f_{\mathbf u}\) is a unit multiple of
\(t_1-1\). The map \(J\to R\to\Lambda/(t_2-1)\)
therefore factors through exactly the evaluation map
in (D.85.20), with image the augmentation ideal.

After multiplication by \(t_1-1\) on both modules and
the Cartier untwist, this natural comparison is
\[
 I^2\hookrightarrow I,\qquad
 I=(t_1-1)\mathbb Z_5[[t_1-1]],
                                                               \tag{D.85.23}
\]
with pure substitution action. The possible extra generator
\(G\) maps to zero and has zero coinvariants by (D.85.21).
The quotient \(I/I^2=\mathbb Z_5(1)\) has nontrivial
order-four character. Exact order-four coinvariants make
(D.85.23) an isomorphism on full \(\Gamma\)-coinvariants.
Thus the natural map in (D.85.14) is an isomorphism also
at this simultaneous infinite endpoint.

We have proved the complete restriction-kernel bound
\[
 \mathcal R_s\simeq
 \begin{cases}
 \mathbb F_5,&e>0,\ d=0,\\
 0,&\text{otherwise}.
 \end{cases}                                                    \tag{D.85.24}
\]

\paragraph{The nonzero cubic class survives the finite-index global quotient.}
For the remaining case \(e>0,d=0\), the normalized cubic
cocycle \(c_3\) of D.78 has values in the constant
\(\mathbb F_5(-1)\subset B^D\). It remains nonzero in
\(H^1(I_0,B^D)\). Indeed
\[
                    B\simeq\mathbb Z_5[[y]]/(g_h).
                                                               \tag{D.85.25}
\]
Choose an actual \(\tau\in I_0\) with
\(\chi(\tau)=1,k_5(\tau)=1,c_3(\tau)=0\), as in D.84,
by correcting the cubic value with a primitive element of
\(N_1\). Its compact multiplier on (D.85.25) is
\((1+y)\mathcal B_\tau\). Since beta starts in degree
three on \(\ker\chi\), its difference from one is
\(y\) times a unit, before and after quotienting by \(g_h\).
Its fixed Cartier functionals are therefore exactly the
constant module \((B/yB)^D\).

If \(c_3\) were a coboundary, its primitive would be fixed
by \(\tau\), and hence constant. An actual element of
\(N_1\) with nonzero cubic character fixes constants and
contradicts that coboundary identity. This proves the claimed
nonvanishing. Natural coinduction and inflation in (D.85.13)
preserve it in \(H^1(G_{\mathbb Q},D)\).

Modulo five, \(\mathbf z=0\) when \(e>0\), whereas
\(\mathbf r\) is a nonboundary direction when \(d=0\).
The actual local cubic primitive of D.82 therefore applies.
Its local construction requires only these local residual
vectors, not D.82's separate global hypothesis.
At every prime away from five with nonzero abelian inertia,
D.43 and the finite-local-duality passage of D.80 give
\(H^1(G_{\mathbb Q_q},D)=0\). Only finitely many primes
have nonzero inertia in the fixed cubic coefficient module.
Enlarge the finite coefficient level to contain their local
primitives. At every other prime D.78's unramified cubic
Frobenius primitive works at its original common level.
This is D.80's finite-exception argument, and produces one
finite global coefficient module with local triviality at
all places. The global nonzero class above then shows
\(\mathcal S_s^*\ne0\). Equation (D.85.24) proves
\(\mathcal S_s=\mathbb Z/5\) in precisely this case.

\paragraph{Assembly and the actual new point.}
In the finite-index case, \(\overline{\mathbf r}\ne0\)
and \(\overline{\mathbf z}=\bar\lambda\overline{\mathbf u}\).
If \(e=0\), the local residual span has dimension two.
If \(e>0\), it is the valuation line, which is nonboundary
exactly when \(d=0\). Thus (D.85.24) and its attained
nonzero case say exactly \(\mathcal S_s=(\mathbb Z/5)^{\epsilon(W)}\).
Together with the already proved cases D.82 and D.84, this
covers every possibility under (D.85.1). The global boundary
depths are zero, and D.77's continuous splitting supplies
the remaining assertions of (D.85.4) with the exact local
sum \(n\).

For (D.85.5), \(799=17\cdot47\) and
\(800=2^5\cdot5^2\). Hence
\(L=\mathbb Z_5(1,0)+5\mathbb Z_5(0,1)\),
\(\mathbf u=(1,0)\), \(\mathbf r=(0,2)\), and \(h=1\).
The unit coefficient of \(c=\kappa(17)+\kappa(47)\)
has depth
\[
 e=v_5(799^4-1)-1=1,
 \qquad d=\infty.
                                                               \tag{D.85.26}
\]
The full rational point supplies actual section existence;
D.74 then gives the asserted entire metabelian fibre.
The supporting exact certificate
\path{code/de_sha_global_independence.py} checks the new integral
generator/action identities, the exhaustive residual case
partition, depth identities including cancellation, and
this point. The full cohomology and infinite-endpoint
claims are proved above, not inferred from those finite checks.
\end{proof}


% D.86 integrated after complete root and independent proof audits; historical candidate retained.
\subsection{Infinite compact five-torsion in an actual cusp fibre}

\begin{statement}{D.86}{Proved}{An infinite torsion compact kernel for a cusp parameter}
Let
\[
 a=\sum_{q\ne5}a_q\kappa(q),\qquad a_q\longrightarrow0,
 \qquad \overline{\langle a_q:q\ne5\rangle}=\mathbb Z_5,
 \qquad
 \lambda=\sum_{q\ne5}a_q\frac{\log(q^4)}{4\log6},
 \qquad v_5(\lambda)=1.                                  \tag{D.86.1}
\]
Thus the five-valuation of \(a\) is zero, its global away-five
valuation image is primitive, and its local unit class is
\(\lambda\kappa(6)\). Arbitrary infinite \(c_0\) support is allowed.
For the actual Alexander module \(A_s\) of the pair \(s=(a,0)\),
the compact groups of D.63 satisfy
\[
 \begin{split}
 \mathcal S_s&\simeq\prod_{n\ge0}\mathbb F_5,\\
 \mathcal T_s&\simeq\prod_{n\ge0}\mathbb Z_5,\qquad
 L_5\simeq\mathcal T_s\times\mathbb Z_5^2,\\
 \mathcal K_s&\simeq
       \left(\prod_{n\ge0}\mathbb F_5\right)\times\mathbb Z_5^2.
 \end{split}                                             \tag{D.86.2}
\]
All isomorphisms are of compact topological groups. The map
\(L_5\to\mathcal T_s\) is the identity on the zero-parameter
boundary factor and kills the other two factors. The last
splitting is continuous and noncanonical.

Every pair \((a,0)\) here has an actual cusp section. In particular
\(a=\kappa(26)\) satisfies (D.86.1); its nonempty metabelian
section fibre is a torsor under
\(\prod_{\mathbb N}\mathbb Z_5\), by D.74. The infinite torsion
in (D.86.2) concerns the compact second-cohomology kernel, not
torsion in that first-cohomology translation group. Each finite
collection of the constructed nonzero dual classes occurs at
one finite coefficient level and is locally trivial there at
\emph{every} place. No single finite level is asserted to contain
all infinitely many independent classes.

This settles the stated globally dependent family only. It
makes no classification of all dependent pairs, no arbitrary
infinitely ramified first-cohomology rank computation, and no
first higher-derived obstruction computation.
\end{statement}

\begin{proof}
Write
\[
 \Lambda=\mathbb Z_5[[u,v]],\quad t_1=1+u,\quad t_2=1+v,
 \quad w=t_1t_2-1,\quad P=uvw,\quad D=A_s^D.
\]
A superscript \(D\) means Cartier dual; a star means ordinary
Pontryagin dual after the Cartier twist has been accounted for.
Let \(I_0\) be the actual compact untwisted operator image of
D.78, including its cyclotomic multiplier. Put
\(N_0=\ker(\chi:I_0\to\Gamma)\), with
\(\Gamma=\mathbb Z_5^\times\). The actual local surjectivity
\(G_{\mathbb Q_5}\to I_0\), the saturation of \(\kappa(6)\)
on its kernel, and the primitive cubic element of D.77--D.78
are used with their already proved integral meanings.

\paragraph{Actual global and local images.}
Let \(H=\mathbb Z_5\) act by multiplication by \(t_1^h\).
The untwisted action is unramified outside five. Actual inertia
at the primes in (D.86.1) therefore generates the whole pure
translation subgroup \(H\). The uniformly convergent cocycle
and local gauge construction of D.80/D.83 gives actual equalities
\[
              I_s=H\rtimes I_0,
       \qquad I_5=5H\rtimes I_0.                       \tag{D.86.3}
\]
The second equality uses \(\lambda\mathbb Z_5=5\mathbb Z_5\),
and both use the same zero-translation complement. The local
gauge lies in \(H\), so preserves the first equality. These
are faithful compact operator images, not formal semidirect
products assigned independently of Galois.

Put
\[
 \begin{gathered}
 q(u)=(1+u)^5-1,\qquad S(u)=q(u)/u,\\
 S(u)=5+10u+10u^2+5u^3+u^4,\\
 B=\Lambda/(u),\quad R=\Lambda/(q),\quad J=uR.
 \end{gathered}
                                                               \tag{D.86.4}
\]
The continuous translation coinduction of D.81/D.84 gives
natural identifications
\[
 H^1(I_s,D)=H^1(I_0,B^D),\qquad
 H^1(I_5,D)=H^1(I_0,R^D).                               \tag{D.86.5}
\]
Their restriction map is induced by the actual inclusion
\(B^D\subset R^D\). In particular its \emph{whole} kernel is
\[
 \mathcal R=(J^D)^{I_0}/\operatorname{im}((R^D)^{I_0}). \tag{D.86.6}
\]
The finite affine Chebotarev descent of D.78/D.83, at every
finite coefficient level before taking a union, gives
\[
                         \mathcal S_s^*\lhook\joinrel\longrightarrow
                         \mathcal R.                  \tag{D.86.7}
\]
This gives an upper bound; we will separately construct
infinitely many classes satisfying the common finite-level
local-triviality condition defining \(\mathcal S_s^*\).

\paragraph{The exact compact connecting module.}
Every \(g\in N_0\) has \(\mathcal B_g-1\in P\Lambda\).
There is an actual \(\delta\in N_0\) with
\(\mathcal B_\delta-1=P E\), \(E\in\Lambda^\times\),
by the local raw cubic primitivity of D.77. Consequently
\[
 (J^D)^{I_0}=(T(-1)_\Gamma)^*,\qquad T=J/PJ.            \tag{D.86.8}
\]
Here \((-1)\) means that the compact operator is
\(\chi^{-1}\rho_0\), exactly one Cartier untwist. The
\(\Gamma\)-action is induced by actual lifts in \(I_0\);
differences between lifts have already been killed by \(PJ\).
Compact coinvariants mean the quotient by the closed span of
operator differences.

Multiplication by the actual generator \(u\) identifies the
scalar module underlying \(J\) with
\[
             A=\Lambda/(S(u))=\mathcal O[[v]],
             \qquad \mathcal O=\mathbb Z_5[\zeta_5],
             \quad u=\zeta_5-1.                       \tag{D.86.9}
\]
The last equality follows from the Eisenstein polynomial in
D.77. Thus \(A\) is a domain, \(u\) is a uniformizer in
\(\mathcal O\), and \(A/(u)=\mathbb F_5[[v]]\).
In scalar coordinates \(T=A/(uvw)\), but its elements still
have the extra actual factor \(u\). We retain that factor in
all action calculations.

There is an exact sequence of actual compact modules
\[
 0\longrightarrow E\longrightarrow T\longrightarrow Q
       \longrightarrow0,
 \quad E=vwA/(uvwA)\simeq\mathbb F_5[[v]],
 \quad Q=A/(vw).                                      \tag{D.86.10}
\]
The generator of \(E\) in the \emph{actual} ideal is \(P=uvw\),
not the scalar \(vw\). The isomorphism for \(E\) uses
cancellation of the nonzero divisor \(vw\) in \(A\).
The ideals \((v)\) and \((w)\) have intersection \((vw)\):
modulo \(v\), \(w=u\) is a nonzero divisor in \(\mathcal O\),
and the converse divisibility follows by cancellation.
Therefore
\[
 0\longrightarrow Q\longrightarrow A/(v)\oplus A/(w)
       \longrightarrow A/(v,w)=\mathbb F_5\longrightarrow0.
                                                               \tag{D.86.11}
\]
Each term here carries the actual factor \(u\). In particular
\(Q\) is finite free of rank eight over \(\mathbb Z_5\).

\paragraph{The infinite torsion layer and the finite boundary.}
Let \(\gamma\in\Gamma\) and write
\(S_\gamma(t)=(t^\gamma-1)/(t-1)\). The actual multiplier
of \(\rho_0\) is
\(S_\gamma(t_1)S_\gamma(t_2)\mathcal B_g\).
On the generator \(P\), substitution contributes the additional
factor
\(S_\gamma(t_1)S_\gamma(t_2)S_\gamma(t_1t_2)\).
On \(E\), set \(u=0\) and reduce modulo five. Its compact
Cartier-untwisted action is consequently
\[
        f(v)\longmapsto
        \overline\gamma S_\gamma(1+v)^3
                   f((1+v)^\gamma-1).                \tag{D.86.12}
\]
There is no remaining beta term on this axis. Multiplication
by \(v^3\) identifies (D.86.12) with the ideal
\(v^3\mathbb F_5[[v]]\) inside the completed group ring
\(C=\mathbb F_5[[\mathbb Z_5]]\), with action
\[
                    t^x\longmapsto
                       \overline\gamma t^{\gamma x}.  \tag{D.86.13}
\]
The quotient \(C/v^3C\) has successive tame characters
\(\overline\chi,\overline\chi^2,\overline\chi^3\).
All are nontrivial on the order-four subgroup. Exact averaging
on that subgroup therefore identifies the coinvariants of
\(v^3C\) and \(C\).

The discrete ordinary dual of \(C\) is the continuous
\(\mathbb F_5\)-valued function space on \(\mathbb Z_5\).
Invariant functionals for (D.86.13) are functions satisfying
\[
 f(\gamma x)=\overline\gamma^{-1}f(x).
                                                               \tag{D.86.14}
\]
They vanish at zero. On each shell \(5^n\mathbb Z_5^\times\)
they are determined freely by their value at \(5^n\).
Continuity at zero forces those values to vanish for all
sufficiently large \(n\); conversely every finitely supported
sequence gives a continuous function satisfying (D.86.14).
It follows, by duality, that
\[
                         E(-1)_\Gamma
                               \simeq\prod_{n\ge0}\mathbb F_5.
                                                               \tag{D.86.15}
\]
This is the product compact topology, not a direct-sum topology.

For the two boundary modules in (D.86.11), after the Cartier
untwist the exact multiplier is respectively
\(S_\gamma(t)^2\) and
\(S_\gamma(t)^2t^{(1-\gamma)/2}\), on a scalar element
with actual generator \(u\). The second formula is the exact
reflection formula of D.39. Multiplication by \(u^2\), and
on the second boundary also by \(t^{-1/2}\), identifies both
modules with
\[
 I_5^2\subset\mathbb Z_5[\mathbb Z/5],
 \qquad t\longmapsto t^\gamma.                         \tag{D.86.16}
\]
The integral injectivity and equivariance are the same
cancellation argument as D.84.15, with the stated half-integral
power legitimate because two is a unit in \(\mathbb Z_5\).
The cokernel in (D.86.11), retaining the actual generator \(u\),
has character \(\overline\chi^2\) after untwisting. Its
order-four invariants and coinvariants vanish. Thus
\[
                      Q(-1)_\Gamma\simeq\mathbb Z_5^2.
                                                               \tag{D.86.17}
\]
Moreover the action on \(Q(-1)\) factors through the finite
order-four quotient: this is true on both finite group-ring
modules in (D.86.16), and (D.86.11) is an actual equivariant
injection.

Choose an actual lift \(k\) of \(2\in\Gamma\). Two
topologically generates \(\Gamma\), so all the compact
coinvariants above are cokernels of \(k-1\). The snake lemma
for (D.86.10) shows that the kernel of
\(E(-1)_\Gamma\to T(-1)_\Gamma\) is an image of
\(Q(-1)^{k=1}\). This last module is finite free over
\(\mathbb Z_5\), so its image in the exponent-five module
(D.86.15) is finite. A quotient of
\(\prod_{\mathbb N}\mathbb F_5\) by a finite subgroup is
again such a product: dualizing removes only a finite
codimension from a countably infinite-dimensional vector space.
Consequently (D.86.10), (D.86.15), and (D.86.17) give
\[
 T(-1)_\Gamma\simeq
             \left(\prod_{\mathbb N}\mathbb F_5\right)
                           \times\mathbb Z_5^2.        \tag{D.86.18}
\]
The split is continuous: choose lifts of the two basis vectors
of the free compact quotient, then use continuous
\(\mathbb Z_5\)-multiplication. This argument keeps the
possible finite snake-map image; it does not assume that
\(E\)-coinvariants inject.

\paragraph{The entire restriction kernel has exponent five.}
D.77 computes the actual invariant groups and their natural
inclusions. Globally only the \(b=0\) axis contributes; locally
the \(a\) and \(a-b\) axes each contribute one additional
copy of \(\mathbb Q_5/\mathbb Z_5\). Thus
\[
 (R^D)^{I_0}/(B^D)^{I_0}
                  \simeq(\mathbb Q_5/\mathbb Z_5)^2.   \tag{D.86.19}
\]
Its map into \((J^D)^{I_0}\) is injective by exactness.
By (D.86.8) and (D.86.18), the latter group is
\[
             \bigoplus_{\mathbb N}\mathbb F_5
                         \ \oplus\
                     (\mathbb Q_5/\mathbb Z_5)^2.
\]
Every divisible subgroup lies in its maximal divisible
subgroup. The image in (D.86.19) fills that subgroup:
dualizing its injection gives a surjection between free
\(\mathbb Z_5\)-modules of rank two, necessarily an
isomorphism. Therefore the actual quotient (D.86.6) satisfies
\[
                    \mathcal R\simeq
                         \bigoplus_{\mathbb N}\mathbb F_5.
                                                               \tag{D.86.20}
\]
In particular (D.86.7) implies that \(\mathcal S_s^*\)
has exponent five and at most countable dimension. This
upper calculation alone does not prove that its dimension
is infinite; the next steps address that separately.

\paragraph{Infinitely many actual global cocycles.}
The quotient \(J/uJ\), with its actual generator \(u\), is
\(\mathbb F_5[[v]]\). Its untwisted compact action is
\[
 f(v)\longmapsto
       \overline\gamma S_\gamma(1+v)
                      f((1+v)^\gamma-1).              \tag{D.86.21}
\]
All elements of \(N_0\) act trivially on this quotient.
Multiplication by \(v\) puts (D.86.21) inside the same
module (D.86.13); its missing constant quotient has
nontrivial character \(\overline\chi\). The shell argument
therefore gives an infinite-dimensional countable vector space
\[
 W:=((J/uJ)^D)^{I_0}\simeq\bigoplus_{\mathbb N}\mathbb F_5.
                                                               \tag{D.86.22}
\]
View \(W\) as a subspace of \((J^D)^{I_0}\).
Its intersection with the image of (D.86.19) has dimension
at most two, because it consists of elements of order five
in a divisible group of corank two. Hence its image in
\(\mathcal R\) is infinite-dimensional.

For \(\varphi\in W\), choose a concrete order-five extension
\(F\in R^D\) as follows. In
\[
                  R/5R=\mathbb F_5[[u,v]]/(u^5),
\]
let \(F\) apply the functional defining \(\varphi\) to the
coefficient of \(u^1\). Thus \(F|_J=\varphi\); it depends
on finitely many coefficients in \(v\), and on no powers
of \(u\) other than the first. Define
\[
 z_\varphi(g_0)=g_0F-F\in B^D\quad(g_0\in I_0),
 \qquad z_\varphi(h,g_0)=z_\varphi(g_0)\quad((h,g_0)\in I_s).
                                                               \tag{D.86.23}
\]
The first inclusion holds by invariance of \(\varphi\).
The second formula is an actual cocycle because \(H\)
acts trivially on \(B^D\). Its class is the connecting
image of \(\varphi\), under (D.86.5). Inflation from
\(I_s\) to \(G_{\mathbb Q}\) is injective in degree one
for these coefficients. Thus the infinite-dimensional image
just found consists of genuine nonzero global cohomology
classes. All cochains in (D.86.23) are continuous and have
values in a finite coefficient module, because \(F\) has
finite coefficient support and the actual action preserves
the augmentation filtration.

\paragraph{A bound for every unramified Frobenius primitive.}
Put \(\mathfrak a=(u,v)\), and let
\[
              M_n=(A_s/(5,\mathfrak a^n)A_s)^D.
\]
Choose \(n\ge2\) such that \(F\) and all values of
\(z_\varphi\) lie in \(M_n\). For any actual element
\(g=(h,g_0)\in I_s\), a primitive of
\(z_\varphi(g)\) for the cyclic group generated by \(g\)
exists in the single module \(M_{n+1}\). Here is a uniform
proof with all three cases distinguished.

If \(h\equiv0\pmod5\), the first-\(u\) functional \(F\)
is fixed by multiplication by \(t_1^h\): modulo five this
power has no term of degree one through four in \(u\).
Since \(H\) also fixes \(B^D\), (D.86.23) gives
\(gF-F=z_\varphi(g)\).

If \(\chi(g)\not\equiv1\pmod5\), integral conjugation
by an element of \(H\) removes the coordinate \(h\),
since \(\chi(g)-1\) is a unit. The cocycle is zero on
\(H\), and \(H\) fixes its values. Transporting the
primitive \(F\) from \((0,g_0)\) back by this conjugation
gives a primitive for \(g\) in \(M_n\). The conjugating
translation preserves this finite module. This is a
cocycle conjugation calculation, not an assumption that
the raw Frobenius translation was zero.

It remains that \(h\not\equiv0\pmod5\) and
\(\chi(g)\equiv1\pmod5\). On the compact coefficient
space modulo five set
\(U_g=\chi(g)^{-1}\rho_s(g)\). The exact low-degree
Fox/Ihara formulas of D.31/D.35 imply
\[
 U_g-1\text{ raises augmentation degree by at least one,}
 \qquad \operatorname{gr}(U_g-1)=\overline h\,u.
                                                               \tag{D.86.24}
\]
Indeed \(t_1^h=1+\overline h u+O(\mathfrak a^2)\);
\(S_\chi(t_i)\equiv1\pmod{\mathfrak a^2,5}\) when
\(\chi\equiv1\pmod5\); beta has no linear term;
and substitution \(t_i\mapsto t_i^\chi\) differs from
identity only in degree at least five on each variable.
The Cartier scalar is one modulo five. Multiplication by
\(\overline h u\) is injective on every homogeneous
polynomial piece of \(\mathbb F_5[u,v]\). Comparing the
least nonzero degree therefore proves that
\[
 U_g-1:\mathbb F_5[[u,v]]/\mathfrak a^n
          \longrightarrow
          \mathbb F_5[[u,v]]/\mathfrak a^{n+1}
                    \quad\hbox{is injective}.          \tag{D.86.25}
\]
Duality makes the actual \(g-1\) map
\(M_{n+1}\to M_n\) surjective. More explicitly, its
matrix differs from the transpose dual of \(U_g-1\)
by the invertible factor \(-U_g^{-T}\). Hence it has
a preimage of \(z_\varphi(g)\), with the same bound
\(n+1\) for every such \(g\). This proves the uniform
cyclic primitive assertion.

\paragraph{One finite global level for every place.}
Choose \(b\) with \(5^b\ge n+1\). By the \(c_0\)
condition, the set
\[
          E_n=\{q\ne5:a_q\notin5^b\mathbb Z_5\}
                                                               \tag{D.86.26}
\]
is finite. At every \(q\notin E_n\cup\{5\}\), the
module \(M_{n+1}\) is unramified: its inertia multiplier
is \(t_1^{a_q}\equiv1\pmod{5,\mathfrak a^{n+1}}\).
The cocycle (D.86.23) is also unramified there, because
it is inflated from the untwisted image \(I_0\).
Apply the preceding primitive construction to the actual
Frobenius element. It gives local triviality in the same
finite module \(M_{n+1}\) at all these infinitely many
primes.

At five, every element of the translation subgroup \(5H\)
fixes \(F\). Equation (D.86.23) is consequently the
coboundary of \(F\) on the actual whole group
\(I_5=5H\rtimes I_0\), already in \(M_n\).
At every \(q\in E_n\), the abelian inertia is nonzero.
D.43 and D.80's finite local Tate-duality/Mittag--Leffler
argument give
\(H^1(G_{\mathbb Q_q},D)=0\). Thus the local class
becomes a coboundary in some finite coefficient module
\((A_s/\mathfrak m^{r_q}A_s)^D\), where
\(\mathfrak m=(5,u,v)\). This is used for the finite
exceptional set only. Enlarge to one level containing
\(M_{n+1}\) and all those finitely many primitives.
The earlier primitives remain valid after inclusion, even
when this larger module has additional ramification.
Odd-primary cohomology at the real place vanishes.

We have proved that each class from (D.86.23) has a
representative in one finite global coefficient module
that is locally trivial at every place. Taking maxima
gives the same assertion for any finite collection.
No interchange of a filtered union with an infinite
product of local cohomology groups has occurred.
Their infinite-dimensional image therefore lies in
\(\mathcal S_s^*\). Combined with the exponent-five,
countable-dimensional upper bound (D.86.20), this gives
\[
                     \mathcal S_s^*\simeq
                            \bigoplus_{\mathbb N}\mathbb F_5.
\]
Pontryagin duality proves the first assertion of (D.86.2)
with its stated compact product topology.

\paragraph{The remaining factors and an actual cusp example.}
The nonzero classes \(a\) and \(a-b=a\) have global
depth zero and local depth one; the zero class \(b\)
has infinite depth both globally and locally. D.77's
natural boundary description therefore gives exactly
\(L_5\to\mathcal T_s\) and its kernel \(\mathbb Z_5^2\)
as stated. Its continuously split compact sequence gives
\(\mathcal K_s\simeq\mathcal S_s\times\mathbb Z_5^2\).
This concerns only D.63's compact separated kernel and
asserts no splitting of the unrestricted nonseparated
cochain quotient.

The cusp at zero has inertia parameter \(a\) and exact
abelian pair \((a,0)\); its inertia cocycle supplies an
actual geometrically pro-five section. One can also take
zero parameters at all other coefficient primes to get
an actual full profinite cusp section.
For \(a=\kappa(26)=\kappa(2)+\kappa(13)\), the away
coefficients are primitive and the five-valuation is zero.
Since \(26=1+25\),
\[
                 v_5(\log26)=2,\qquad v_5(\log6)=1,
                 \qquad v_5(\lambda)=1.
\]
Thus this explicit cusp section has the compact kernel
in (D.86.2). Its twist is finitely ramified, so the
already evaluated D.74 theorem computes its nonempty
metabelian fibre as stated, without identifying that
whole fibre with the cusp packet.
\end{proof}


% D.87 integrated after complete root and independent proof audits; historical candidate retained.
\subsection{A uniform finite-level comparison with the actual operator image}

\begin{statement}{D.87}{Proved}{The exact finite-localization kernel of the actual operator image}
For every completed pair \(s=(a,b)\) of five-adic Kummer classes
over \(\mathbb Q\), let \(I_s\) and \(I_5\) be its actual
compact global and local-at-five Alexander operator images,
including the cyclotomic multiplier, and let \(D=A_s^D\).
There is a natural identification
\[
 \mathcal S_s^*\simeq
 \ker\{H^1(I_s,D)\longrightarrow H^1(I_5,D)\}.       \tag{D.87.1}
\]
It is an identification with the filtered union of the finite
localization kernels defining \(\mathcal S_s^*\), not merely
with the unrestricted kernel of restriction to infinitely many
local groups. Every given class on the right has a representative
at one finite coefficient level which is locally trivial there
at every place.

Here is its explicit coefficient form. Let
\[
 \mathbf r=(a_5,b_5),\qquad
 L=\overline{\sum_{q\ne5}\mathbb Z_5(a_q,b_q)},\qquad
 \mathbf z=\sum_{q\ne5}\frac{\log(q^4)}{4\log6}(a_q,b_q)
                                                     \quad\in L.
                                                               \tag{D.87.2}
\]
The sums have their actual uniform \(c_0\) meaning. Write
\(f_{\mathbf v}=t_1^{v_1}t_2^{v_2}-1\) and let \(I_L\)
be the closed ideal generated by these elements for
\(\mathbf v\in L\). With
\[
 B=\Lambda/I_L,\qquad R=\Lambda/(f_{\mathbf z}),
                 \qquad J=I_LR,
                                                               \tag{D.87.3}
\]
where the relation is omitted for \(\mathbf z=0\), (D.87.1)
becomes the natural exact formula
\[
 \mathcal S_s^*\simeq
       (J^D)^{I_0}/\operatorname{im}((R^D)^{I_0}).       \tag{D.87.4}
\]
The actual \(I_0\)-action here is
\(t^{\mathbf r\kappa_5(g)}\rho_0(g)\). The formula
includes zero and nonprimitive lattices of either rank.
It does not claim an isomorphism type for every resulting
module quotient, a general section-existence theorem, or a
first higher-derived obstruction.

\end{statement}

\begin{proof}
\paragraph{Exact images and the two inclusions.}
D.83's actual closed-image and integral local-gauge argument
applies with precisely (D.87.2). It gives
\[
 I_s=L\rtimes I_0,\qquad
 I_5=\mathbb Z_5\mathbf z\rtimes I_0,                 \tag{D.87.5}
\]
where the common embedded complement is the actual graph
\(g\mapsto(\mathbf r\kappa_5(g),g)\). Its conjugation on
\(L\) is multiplication by \(\chi(g)\).
In particular the common complement is contained in \(I_5\).
These are faithful actual operator groups; no formal series
or abstract image is substituted for them.

A finite everywhere-locally-trivial class descends through
\(I_s\), by the finite affine Chebotarev argument already
proved in D.78/D.83. For completeness, if its cocycle were
nonzero on the kernel of the actual operator representation,
its finite affine image would contain \((v,1)\), \(v\ne0\).
Chebotarev supplies an unramified Frobenius in that conjugacy
class. Its linear action is the identity while its cocycle
value is a nonzero invertible linear image of \(v\), so its
local unramified class cannot vanish. This is a contradiction.
The cocycle therefore descends continuously through the compact
quotient \(I_s\). Its local-at-five class is zero on \(I_5\),
since \(G_{\mathbb Q_5}\to I_5\) is onto.
Inflation in degree one is injective. Passing to the filtered
union of finite coefficients gives the natural injection from
left to right in (D.87.1). If an element becomes zero globally
in the union, it is already zero at a later finite global
level, hence zero in that finite localization kernel; this
also checks injectivity of the union into \(H^1(G_{\mathbb Q},D)\).

For the reverse inclusion, fix a class in the right-hand kernel
and a continuous actual cocycle \(z:I_s\to D\) representing
it. Its restriction to the common graph complement \(I_0\)
is a coboundary, since \(I_0\subset I_5\). We first prove
that for every \(g\in I_s\), the element \(z(g)\) is a
\(g-1\) coboundary in the \emph{full} module \(D\).

\paragraph{The full torsion-free coefficient is essential.}
We use here \(A_s=\Lambda C\), not any nonprimitive quotient
\(R\), \(B\), or \(J\). It is a torsion-free formal power-series
module over \(\mathbb Z_5\). Let
\[
                     U_g=\chi(g)^{-1}\rho_s(g)
\]
be its compact Cartier-untwisted operator. If
\(\chi(g)\in1+5\mathbb Z_5\) is nontrivial, then
\[
                         \ker(U_g-1)=0.               \tag{D.87.6}
\]
Indeed the augmentation graded piece of degree \(n\ge0\)
is free over \(\mathbb Z_5\), and the operator on it is
multiplication by \(\chi(g)^{n+1}\): the actual Alexander
multiplier has augmentation \(\chi^2\), substitution has
weight \(\chi^n\), and the Cartier untwist removes one
factor \(\chi\). A nonzero formal series has a least
nonzero augmentation degree. Its fixedness would imply
\((\chi(g)^{n+1}-1)F_n=0\) in that torsion-free piece.
The scalar is nonzero because \(1+5\mathbb Z_5\) is
torsion-free. This is impossible. No argument on the possibly
torsion graded pieces of a cyclotomic quotient is used.

If \(\chi(g)=1\), the full operator is multiplication by
an actual unit \(\mu_g\in\Lambda^\times\). If
\(\mu_g\ne1\), then \(U_g-1\) is injective because
\(\Lambda\) is a domain. If \(\mu_g=1\), faithfulness
of the actual image, including its cyclotomic coordinate,
means that \(g\) is the identity in \(I_s\); then
\(z(g)=0\).

For every injective case just discussed, ordinary Pontryagin
duality gives surjectivity of \(g-1:D\to D\). In matrices
or operator notation this is the dual of
\(U_g^{-1}-1\), whose injectivity is equivalent to that of
\(U_g-1\). Continuous duality on compact abelian groups is
exact, so no density or rationalization is substituted for
actual surjectivity on the discrete dual.

The remaining case is \(\chi(g)\not\equiv1\pmod5\).
Write \(g=(h,g_0)\) in (D.87.5). Since
\(\chi(g)-1\in\mathbb Z_5^\times\), the translation
\(h/(\chi(g)-1)\) belongs to the actual lattice \(L\).
Conjugation by it sends \(g\) into the common complement
\(I_0\). The cocycle class is zero on that complement,
so its restriction to the cyclic subgroup generated by
\(g\) is a coboundary as well. This is invariance of
restriction under actual group conjugation; it does not
assume that the cocycle is unchanged pointwise by conjugation.
Explicitly, for a conjugating element \(b\) and
\(g'=bgb^{-1}\),
\[
 z(g')=b\,z(g)+(1-g')z(b).
\]
If \(z(g')=(g'-1)F\), then
\(z(g)=(g-1)b^{-1}(F+z(b))\).
Together these cases prove cyclic triviality for every actual
\(g\in I_s\), including all torsion cyclotomic values.

\paragraph{Compactness supplies one finite bound for all elements.}
Put
\[
       W_r=A_s/\mathfrak m^r A_s,\qquad
       D_r=W_r^D,\qquad \mathfrak m=(5,u,v).
\]
These are finite \(I_s\)-stable modules, and their increasing
union is \(D\). Continuity of \(z\) on the compact group
\(I_s\) gives finite image, so all its values lie in some
\(D_{r_0}\). For \(r\ge r_0\), define
\[
 U_r=\{g\in I_s:z(g)\in(g-1)D_r\}.                    \tag{D.87.7}
\]
Each \(U_r\) is clopen: it is the finite union, over
\(d\in D_r\), of the sets where the two continuous
finite-valued functions \(z(g)\) and \((g-1)d\) agree.
The action really factors through a finite group on this
\(I_s\)-stable coefficient module. The preceding paragraphs
show that the increasing sets \(U_r\) cover \(I_s\).
Compactness therefore gives one integer \(r\) with
\[
                              U_r=I_s.                 \tag{D.87.8}
\]
This asserts a common finite coefficient bound, not merely
existence of unrelated unbounded primitives. Continuous
selection of a primitive is not needed.

\paragraph{From the uniform cyclic bound to all local groups.}
Continuous translation coinduction in (D.87.5) permits choosing
the representative \(z\) as an inflation from
\(H^1(I_0,B^D)\), so \(z\) is zero on \(L\).
This remains a continuous cocycle in finite coefficients,
and the argument (D.87.6)--(D.87.8) applies to this representative.
The coinduction proof is valid for arbitrary closed lattices:
a Smith decomposition gives a continuous transversal for
\(\mathbb Z_5^2/L\), including finite-index and rank-zero
cases, and insertion of the identity contracts the positive
translation cochain degrees, as in D.81/D.84.

By the \(c_0\) condition, outside a finite set of primes
\(E\) away from five, the pair of inertia coordinates is
so divisible that its multiplier is trivial on \(W_r\).
Thus \(D_r\) is unramified there. The cocycle is unramified
there as well: untwisted inertia is trivial and it is zero
on the pure translation group \(L\). Apply (D.87.8) to an
actual Frobenius element. Its primitive in \(D_r\) gives
local triviality on the entire unramified local quotient.

Every prime in \(E\) where enlargement is needed has nonzero
abelian inertia. D.43 and D.80's finite local Tate-duality
argument give \(H^1(G_{\mathbb Q_q},D)=0\). Each such local
class therefore has a primitive in some finite \(D_{r_q}\).
There are only finitely many such primes. At five the assumed
vanishing on \(I_5\) gives one primitive in \(D\), hence in
some finite \(D_{r_5}\), for the whole local group. Enlarge
\(r\) to dominate these finitely many levels. All the old
primitives remain valid after inclusion, including at primes
where the enlarged coefficient has extra ramification.
The real place has no positive odd-primary cohomology.

The global cocycle is consequently locally trivial everywhere
at this one finite coefficient level. It belongs to
\(\mathcal S_s^*\), proving the reverse inclusion in (D.87.1).
No filtered union was interchanged with an infinite product
of local groups.

Finally the natural coinduction identifications in (D.87.5)
give \(H^1(I_s,D)=H^1(I_0,B^D)\) and
\(H^1(I_5,D)=H^1(I_0,R^D)\), with restriction induced by
\(B^D\subset R^D\). The exact coefficient sequence
\(0\to B^D\to R^D\to J^D\to0\) proves (D.87.4).
All nonprimitive cyclotomic components are retained in that
exact quotient. The only invariant-vanishing argument used
the full torsion-free \(A_s\), as emphasized above.
\end{proof}


% D.88 integrated after complete root, FGH and PBC all-depth proof/certificate audits; candidates retained.
\subsection{Every local depth in a zero-five-valuation cusp packet}

\begin{statement}{D.88}{Proved}{The compact kernel at every cusp index with zero five-valuation}
Let \(a\ne0\) be an arbitrary completed Kummer class over
\(\mathbb Q\) with zero valuation at five. Write
\[
 \begin{aligned}
 r&=\min_{q\ne5}v_5(a_q),&
 \lambda&=\sum_{q\ne5}a_q\frac{\log(q^4)}{4\log6},\\
 N&=v_5(\lambda)\in\mathbb Z_{\ge0}\cup\{\infty\},&
 e&=N-r\in\mathbb Z_{\ge0}\cup\{\infty\}.
 \end{aligned}                                           \tag{D.88.1}
\]
The valuation support may be infinite \(c_0\) support. For the
actual cusp pair \(s=(a,0)\), the compact groups of D.63 are
\[
 \begin{aligned}
 \mathcal S_s&\simeq
 \begin{cases}
  0,&e=0,\\
  \displaystyle\prod_{j\ge0}\mathbb Z/5^e,&0<e<\infty,\\
  \displaystyle\prod_{j\ge0}\mathbb Z_5,&e=\infty,
 \end{cases}\\
 \mathcal K_s&\simeq
 \begin{cases}
  \mathcal S_s\times\mathbb Z_5^{2e},&e<\infty,\\
  \displaystyle\prod_{j\ge0}\mathbb Z_5,&e=\infty.
 \end{cases}
 \end{aligned}                                           \tag{D.88.2}
\]
These are compact topological isomorphisms; the splitting is
continuous and noncanonical. More precisely D.77 identifies
\(L_5\to\mathcal T_s\) with the identity on the zero-parameter
countable product and, on the other two boundaries, the two
projections \(\mathbb Z_5^N\to\mathbb Z_5^r\),
where \(\mathbb Z_5^\infty=\prod_{j\ge0}\mathbb Z_5\)
has the compact product topology.

For every pair of integers \(r\ge0,e\ge1\), the actual rational tangential
cusp parameter
\[
                    a=5^r\kappa(1+5^{e+1})             \tag{D.88.3}
\]
has precisely these indices. Its nonempty metabelian fibre is
a torsor under \(\prod_{\mathbb N}\mathbb Z_5\), by D.74.
Thus every bounded five-primary exponent in (D.88.2) occurs
in a genuine cusp section with finite ramification.

The statement requires \(a_5=0\).
It does not handle the additional actual valuation-graph
relation when \(a_5\ne0\). The infinite endpoint below
is proved from the actual kernel map, not equal infinite coranks. It makes no claim about the first
higher-derived obstruction or all noncuspidal abelian pairs.
\end{statement}

\begin{proof}
The new input D.87 identifies the finite-global-localization
kernel with the \emph{entire} actual operator-image restriction
kernel. We will compute that kernel, retaining its finite
torsion and all actual coefficient actions.

\paragraph{Exact images and the relative cyclotomic module.}
Use the same notation \(\Lambda,t_1,t_2,u,v,w,P,I_0,\Gamma\)
as D.86. For \(n\ge0\), put
\[
                     q_n=(1+u)^{5^n}-1,
                     \qquad q_0=u.
\]
The actual global and local images, after the integral local
gauge of D.83, are
\[
 I_s=5^r\mathbb Z_5\rtimes I_0,
       \qquad I_5=5^N\mathbb Z_5\rtimes I_0.           \tag{D.88.4}
\]
Here \(5^\infty\mathbb Z_5=0\).
The common complement has the untwisted action \(\rho_0\),
because \(a_5=0\). If \(e=0\), the two actual images agree
and D.83, or D.87, gives \(\mathcal S_s=0\). D.77 then
also gives \(\mathcal K_s=0\). First assume \(0<e<\infty\);
the infinite endpoint is treated separately below.

Set
\[
 B=\Lambda/(q_r),\quad R=\Lambda/(q_N),\quad J=q_rR,
 \qquad F_{r,N}=q_N/q_r,
 \qquad A=\Lambda/(F_{r,N}).                           \tag{D.88.5}
\]
Multiplication by the \emph{actual} generator \(q_r\)
identifies the scalar module underlying \(J\) with \(A\).
The polynomial \(F_{r,N}\) is monic, has constant term
\(5^e\), and satisfies
\(F_{r,N}\bmod5=u^{5^N-5^r}\). It is therefore
distinguished and is independent of \(v\). Thus \(A\) is
finite free over \(\mathbb Z_5[[v]]\), and
\[
                     A/(u)=(\mathbb Z/5^e)[[v]].       \tag{D.88.6}
\]
Multiplication by \(u\) is injective in \(A\). Indeed,
if \(uf=F_{r,N}g\) in the full power-series ring, evaluation
at \(u=0\) gives \(5^e g(0,v)=0\), so \(g=ug_1\);
cancellation in the full domain gives \(f=F_{r,N}g_1\).
This is valid even though \(A\) has several cyclotomic
components and its quotient (D.88.6) has five-torsion.

Continuous coinduction and D.87 give the natural formula
\[
 \mathcal S_s^*\simeq
       (J^D)^{I_0}/\operatorname{im}((R^D)^{I_0}).       \tag{D.88.7}
\]
On the cyclotomic kernel, the images of the actual operator
differences generate exactly \(PJ\): the primitive cubic
element of D.77 has difference \(P\) times a unit, and all
other beta differences are multiples of \(P\). Thus the
invariant discrete dual is precisely the annihilator of
\(PJ\). Consequently
\[
 (J^D)^{I_0}=(T(-1)_\Gamma)^*,\qquad
             T=J/PJ\simeq A/(uvw)                    \tag{D.88.8}
\]
in scalar coordinates. The retained actual generator is
\(q_r\). In particular its substitution factor is
\(S_\chi(t_1^{5^r})\), with constant coefficient \(\chi\),
not a scalar obtained by dividing by \(5^r\).

\paragraph{The exact torsion layer; the snake map is zero.}
As \(v\) and \(w\) form a regular sequence in \(A\),
there is an exact sequence of compact actual modules
\[
 0\longrightarrow E\longrightarrow T\longrightarrow Q
      \longrightarrow0,
 \quad E=vwA/(uvwA)\simeq(\mathbb Z/5^e)[[v]],
 \quad Q=A/(vw).                                      \tag{D.88.9}
\]
For the regular-sequence assertion, \(v\) is a nonzero
 divisor in the formal variable and \(w\bmod v=u\) is
injective by the previous paragraph. The actual generator of
\(E\) is \(q_rvw\). At \(u=0\), the Alexander multiplier
contributes \(\chi S_\chi(t_2)\), substitution on \(q_r\)
contributes \(\chi\), and substitution on each of \(v,w\)
contributes \(S_\chi(t_2)\). After the one Cartier untwist,
the action on (D.88.6) is
\[
 f(v)\longmapsto\chi S_\chi(1+v)^3
                         f((1+v)^\chi-1).             \tag{D.88.10}
\]
Multiplication by \(v^3\) identifies this with
\(v^3C\), where \(C=(\mathbb Z/5^e)[[\mathbb Z_5]]\)
has action \(t^x\mapsto\chi t^{\chi x}\).
The missing first three degrees have nontrivial tame
characters \(\chi,\chi^2,\chi^3\). Exact order-four
averaging therefore gives
\[
                     E(-1)_\Gamma\simeq
                           \prod_{j\ge0}\mathbb Z/5^e.
                                                               \tag{D.88.11}
\]
Indeed its invariant discrete dual is the space of continuous
functions \(f:\mathbb Z_5\to\mathbb Z/5^e\) satisfying
\(f(\gamma x)=\gamma^{-1}f(x)\). Such a function vanishes
at zero, has one freely chosen value on each nonzero valuation
shell, and has only finitely many nonzero shell values by
continuity at zero. This proves the asserted compact topology.

There is no hidden finite loss in (D.88.11). On
\[
             Q/uQ=(\mathbb Z/5^e)[[v]]/(v^2)
\]
the actual Cartier-untwisted action has successive characters
\(\chi^2,\chi^3\): the retained generator is \(q_r\).
Their order-four invariant parts are zero. Hence
\(Q(-1)^\Gamma\subset uQ(-1)\). The natural map
\[
                              uT\longrightarrow uQ
                                                               \tag{D.88.12}
\]
is an equivariant isomorphism. Surjectivity is immediate;
its kernel is \(uT\cap E\), which is zero because modulo
\(u\) the product \(vw\) is \(v^2\), a nonzero divisor
even over \(\mathbb Z/5^e\). Equivalently,
\((u)\cap(vw)=(uvw)\) inside \(A\).
Thus every invariant element of the Cartier-untwisted
\(Q(-1)\) has an invariant lift in \(uT(-1)\). All
invariants in this snake argument use that untwisted action.

Choose an actual lift of \(2\in\Gamma\), which topologically
generates \(\Gamma\). Apply the snake lemma to its operator
minus one in (D.88.9). The preceding invariant lifting makes
the connecting map zero, so there is an exact sequence
\[
 0\longrightarrow E(-1)_\Gamma\longrightarrow T(-1)_\Gamma
       \longrightarrow Q(-1)_\Gamma\longrightarrow0. \tag{D.88.13}
\]
This proves injectivity of the infinite torsion layer; it does
not discard a finite snake-map image by comparing ranks.

\paragraph{The actual finite boundary and its torsion.}
Let \(R_n=\mathbb Z_5[\mathbb Z/5^n]\), including
\(R_0=\mathbb Z_5\), and let
\[
       J_{r,N}=\ker(R_N\longrightarrow R_r)=q_rR_N.
\]
The action is the genuine permutation action
\(t\mapsto t^\chi\). Restriction to the two boundaries
\(v=0,w=0\) gives an injection of \(Q\) into their two
scalar \(A\)-modules. Its cokernel is
\(A/(v,w)=\mathbb Z/5^e\), with actual untwisted
character \(\chi^2\), so its tame invariant and coinvariant
parts vanish.

On the first boundary, multiplication by \(u q_r\)
identifies the actual Cartier-untwisted scalar module with
\(uJ_{r,N}\subset R_N\). On the diagonal boundary one
also multiplies by \(t^{-1/2}\), using D.39's exact
reflection formula. These identifications are integral:
if \(u q_r f\in q_N\Lambda\), cancellation of \(q_r\),
evaluation at \(u=0\), and cancellation of \(u\) show
\(f\in F_{r,N}\Lambda\). They preserve the exact action
because substitution supplies the factors
\(S_\chi(t)S_\chi(t^{5^r})\).

The quotient
\(J_{r,N}/uJ_{r,N}\simeq\mathbb Z/5^e(1)\)
has nontrivial tame character. Exact tame averaging gives
\[
       (uJ_{r,N})_\Gamma\simeq(J_{r,N})_\Gamma.        \tag{D.88.14}
\]
We compute the latter integrally. The \(\Gamma\)-orbits in
\(\mathbb Z/5^n\) are zero and its \(n\) nonzero valuation
shells. Its group-ring invariants are the orbit sums, and its
coinvariants are free on the orbit classes. On invariants,
the map \(R_N\to R_r\) sends the zero basis vector to
itself and sends the shell of valuation \(j<r\) to
\(5^e\) times the corresponding shell sum. Shells with
\(j\ge r\) map to a multiple of the zero basis vector.
Its invariant cokernel is therefore
\( (\mathbb Z/5^e)^r\).
On coinvariants the same map is onto: it maps the first
\(r\) shells to the corresponding shells and all remaining
ones to zero's orbit class. Its kernel is free of rank
\(N-r=e\). The exact sequence for the operator \(2-1\)
thus gives
\[
 0\longrightarrow(\mathbb Z/5^e)^r
   \longrightarrow(J_{r,N})_\Gamma
   \longrightarrow\mathbb Z_5^e\longrightarrow0.
\]
The free quotient splits continuously, yielding
\[
 Q(-1)_\Gamma\simeq
                (\mathbb Z/5^e)^{2r}\times\mathbb Z_5^{2e}.
                                                               \tag{D.88.15}
\]
The finite summand is essential when \(r>0\); it has not
been removed by a rational dimension calculation.

\paragraph{The whole restriction kernel and its exponent.}
Equations (D.88.11), (D.88.13), and (D.88.15) imply
\[
 T(-1)_\Gamma\simeq T_{\rm tor}\times\mathbb Z_5^{2e},
\]
where \(T_{\rm tor}\) is a compact bounded-exponent module
and fits naturally in an exact sequence
\[
 0\longrightarrow\prod_{j\ge0}\mathbb Z/5^e
    \longrightarrow T_{\rm tor}
    \longrightarrow(\mathbb Z/5^e)^{2r}\longrightarrow0.
                                                               \tag{D.88.16}
\]
For example boundedness by \(5^{2e}\) follows immediately
from this sequence. The free quotient is split by choosing
lifts of its finite basis.

By D.77's natural map, the quotient
\((R^D)^{I_0}/(B^D)^{I_0}\) is the divisible group
\((\mathbb Q_5/\mathbb Z_5)^{2e}\). Its injective image
in \((J^D)^{I_0}\) is the whole maximal divisible subgroup:
that subgroup has the same finite corank \(2e\), and
ordinary duality turns the injection into a surjection between
free \(\mathbb Z_5\)-modules of the same rank. Thus
(D.88.7) identifies \(\mathcal S_s\) with \(T_{\rm tor}\).

In fact this whole module is killed by \(5^e\), not merely
by \(5^{2e}\). The actual subgroup \(I_5\subset I_s\)
is open of index \(5^e\). Restriction followed by
corestriction on degree-one continuous cohomology is
multiplication by that index.\cite[II, Proposition 1.30]{PBC_MilneCFT}
Here is the cocycle verification
with no normality assumption. For left coset representatives
\(t\), write \(gt=t_g h(g,t)\). If \(z\) is a global
cocycle, the corestriction of its restriction has value
\[
 \sum_t t_g z(h(g,t))
   =5^e z(g)+(g-1)\sum_t z(t),
                                                               \tag{D.88.17}
\]
using the cocycle identity and the permutation \(t\mapsto t_g\).
The same finite-sum formula is a continuous cocycle and sends
coboundaries to coboundaries. Therefore every class in the
restriction kernel is killed by \(5^e\). By D.87 this is
the entire \(\mathcal S_s^*\), hence also
\(5^e\mathcal S_s=0\).

The exact sequence (D.88.16) is now a sequence of modules over
\(\mathbb Z/5^e\), whose quotient is finite free over that
ring. Choosing lifts of its basis splits it continuously.
Its finite free summand is absorbed by reindexing the countable
product. We obtain
\[
                       \mathcal S_s\simeq
                            \prod_{j\ge0}\mathbb Z/5^e.
\]
This is the asserted exact compact group, including its exponent.
The argument does not silently identify a quotient by an
arbitrary finite subgroup with the same bounded product.

\paragraph{The free local factor and actual examples.}
The two nonzero boundary classes \(a,a-b\) have global
depth \(r\) and local depth \(N\); the zero class \(b\)
contributes the identity on its countable product. D.77 therefore
gives the stated natural map, its kernel \(\mathbb Z_5^{2e}\),
and the continuous splitting in (D.88.2).

For (D.88.3), the integer \(1+5^{e+1}\) is congruent to two
modulo four, so its valuation at two is one. It is a unit at
five and its principal-unit logarithm has valuation \(e+1\).
Since \(v_5(\log6)=1\), its Kummer class is globally primitive
and has local unit coefficient of valuation \(e\).
Multiplication by \(5^r\) gives exactly (D.88.1).
The cusp inertia construction supplies an actual full section
with this parameter, and finite support permits application
of D.74 to its nonempty metabelian fibre.

\paragraph{The infinite local index uses the actual kernel map.}
Now let \(\lambda=0\), so \(N=e=\infty\), while \(r\)
is still finite because \(a\ne0\). The actual local translation
subgroup is zero. In (D.88.5) replace \(R\) by \(\Lambda\)
and put \(J=q_r\Lambda\). Multiplication by its actual
generator identifies its scalar module with \(A=\Lambda\).
The same primitive cubic operator and the same actual coefficient
sequence give
\[
 \begin{aligned}
 \mathcal S_s&=\ker\left\{
      T(-1)_\Gamma\longrightarrow(\Lambda/(P))(-1)_\Gamma
                         \right\},\\
 T&=J/PJ\simeq\Lambda/(uvw).
 \end{aligned}                                         \tag{D.88.18}
\]
Here the map is induced by the actual inclusion
\(q_r\Lambda\hookrightarrow\Lambda\). In scalar coordinates
it is multiplication by \(q_r\), and all actions retain the
corresponding generator factor. Equation (D.88.18) is the
ordinary compact dual of D.87's exact quotient; hence its kernel
has the actual closed-subgroup topology.

Equations (D.88.9)--(D.88.13) remain valid with
\[
 \begin{aligned}
 E&\simeq\mathbb Z_5[[v]],\qquad Q=\Lambda/(vw),\\
 E(-1)_\Gamma&\simeq\prod_{j\ge0}\mathbb Z_5.
 \end{aligned}                                         \tag{D.88.19}
\]
For the last isomorphism, the continuous weighted shell functions
now take values in \(\mathbb Q_5/\mathbb Z_5\). Each shell
has one arbitrary torsion value, and continuity at zero forces
all but finitely many shell values to vanish. Their direct sum
has the compact dual displayed in (D.88.19). The zero-snake
argument still holds: \(Q/uQ=\mathbb Z_5[[v]]/(v^2)\)
has the two nontrivial tame characters, and multiplication by
\(v^2\) is injective modulo \(u\).

It remains to prove torsion-freeness of the boundary quotient;
no integral torsion is discarded in passing to infinity.
Let
\[
 R_\infty=\mathbb Z_5[[\mathbb Z_5]],\qquad
 J_{r,\infty}=\ker(R_\infty\longrightarrow R_r)=q_rR_\infty.
\]
The two actual boundary modules identify, after multiplication
by \(u q_r\) and the same diagonal half gauge, with
\(uJ_{r,\infty}\). The cokernel of the two-axis injection
has tame character \(\chi^2\); the quotient
\(J_{r,\infty}/uJ_{r,\infty}\simeq\mathbb Z_5(1)\)
has nontrivial tame character as well. Therefore
\[
                 Q(-1)_\Gamma\simeq
                       (J_{r,\infty})_\Gamma^{\,2}.   \tag{D.88.20}
\]

The invariant measures in \(R_\infty\) for the pure multiplier
action are exactly \(\mathbb Z_5\delta_0\). To see this
integrally, project to \(R_m\). An invariant coefficient on
a nonzero valuation shell must, under every further refinement
to \(R_{m+k}\), be divisible by \(5^k\), because its orbit
splits into fibres of that size. Thus that coefficient is zero.
The zero point supplies the stated invariant measure. By
contrast the invariants of \(R_r\) are its \(r+1\) orbit
sums. The actual map on invariants has image exactly
\(\mathbb Z_5\delta_0\), so its cokernel is \(\mathbb Z_5^r\).

The compact coinvariants of \(R_\infty\) are torsion-free.
Indeed their discrete dual consists of continuous functions
constant on each multiplier orbit. Such a function has a value
at zero and only finitely many shell values different from it.
This is a countable direct sum of \(\mathbb Q_5/\mathbb Z_5\),
whose compact dual is a countable product of \(\mathbb Z_5\).
Apply the exact invariant/coinvariant sequence for the actual
generator \(2\) to \(0\to J_{r,\infty}\to R_\infty\to R_r\to0\).
It gives an exact sequence
\[
 0\longrightarrow\mathbb Z_5^r\longrightarrow
    (J_{r,\infty})_\Gamma\longrightarrow
    \ker\{(R_\infty)_\Gamma\to(R_r)_\Gamma\}
                         \longrightarrow0.            \tag{D.88.21}
\]
The right-hand kernel is a closed subgroup of a torsion-free
compact module. Thus \((J_{r,\infty})_\Gamma\), then
\(Q(-1)_\Gamma\) and \(T(-1)_\Gamma\), are torsion-free.
All these modules are countably based compact modules; the
operator images are closed because they are continuous images
of compact modules.

Finally the entire injected factor \(E(-1)_\Gamma\) lies in
the actual kernel (D.88.18). Its actual generator is
\[
                     q_rvw=P\,(q_r/u),                 \tag{D.88.22}
\]
which maps to zero in \(\Lambda/(P)\). Hence
\(\mathcal S_s\) is a closed torsion-free countably based
compact \(\mathbb Z_5\)-module containing
\(\prod_{j\ge0}\mathbb Z_5\). The dual of any torsion-free
compact \(\mathbb Z_5\)-module is a divisible five-primary
group: this follows from exact compact duality applied to
multiplication by five. Countable baseness makes that discrete
dual countable, so the elementary classification of divisible
five-primary groups makes it a finite or countable direct sum
of \(\mathbb Q_5/\mathbb Z_5\). The displayed closed subgroup
rules out a finite number of summands. Therefore
\[
                    \mathcal S_s\simeq\prod_{j\ge0}\mathbb Z_5.
                                                               \tag{D.88.23}
\]
This conclusion uses the actual zero map (D.88.22), not a
comparison of infinite coranks. D.77's two boundary projections
now have countably infinite free compact kernel. Its exact
compact extension by (D.88.23) is again torsion-free, countably
based, and infinite, so \(\mathcal K_s\simeq\prod\mathbb Z_5\)
as asserted. The continuous splitting may also be obtained by
dualizing the injectivity of the divisible dual of the compact
quotient in this extension.

This endpoint occurs in actual completed cusp parameters. Put
\(\lambda_q=\log(q^4)/(4\log6)\). The nonzero finite-support
class
\[
          a=5^r\bigl(\lambda_3\kappa(2)
                         -\lambda_2\kappa(3)\bigr)    \tag{D.88.24}
\]
has global depth \(r\) because \(\lambda_2\equiv2\) and
\(\lambda_3\equiv4\pmod5\), and has local coefficient zero.
It is an actual completed tangential parameter; no assertion
that its two five-adic exponents are ordinary integers is needed.
The full cusp construction supplies its section, and D.74
again determines the torsion-free translation group of its
nonempty metabelian fibre.
The exact distinguished-polynomial, weighted-shell refinement,
invariant-orbit transition, actual-kernel-generator and rational-tangent
checks are retained in
\path{code/de_sha_cusp_finite_depth.py}. Independent
literal integral boundary matrices, including the augmentation
multiple and diagonal half gauge, and their Smith forms are retained in
\path{code/de_cusp_depth_boundary_independent.py}. These finite checks
support the actual formulas; the all-depth and infinite assertions
follow from the written proofs above.
\end{proof}


% D.89 integrated after complete root, FGH and PBC structural proof audits; original candidate retained.
\subsection{Cusp parameters with primitive valuation at five}

\begin{statement}{D.89}{Proved}{The entire compact kernel for a primitive five-valuation cusp}
Let
\[
 a=A\kappa(5)+\sum_{q\ne5}a_q\kappa(q),\qquad
 A\in\mathbb Z_5^\times,\qquad a_q\longrightarrow0,
\]
and consider the actual cusp pair \(s=(a,0)\).
If all \(a_q\) vanish, then \(\mathcal S_s=\mathcal K_s=0\).
Otherwise put
\[
 \begin{aligned}
 r&=\min_{q\ne5}v_5(a_q),&
 \lambda&=\sum_{q\ne5}a_q\frac{\log(q^4)}{4\log6},\\
 N&=v_5(\lambda),&e&=N-r\in\mathbb Z_{\ge0}\cup\{\infty\}.
 \end{aligned}                                             \tag{D.89.1}
\]
Then the compact groups in D.63 satisfy
\[
 \mathcal S_s\simeq\mathcal K_s\simeq
 \begin{cases}
  0,&e=0,\\
  \displaystyle\prod_{j\ge0}\mathbb Z/5^e,&0<e<\infty,\\
  \displaystyle\prod_{j\ge0}\mathbb Z_5,&e=\infty.
 \end{cases}                                               \tag{D.89.2}
\]
These are isomorphisms with the actual compact topologies.
In particular the finite-depth compact kernel has no additional
free factor. The pair always has an actual full cusp section.
For every pair of integers \(r\ge0,e\ge1\), the rational
tangent
\[
                  5(1+5^{e+1})^{5^r}                     \tag{D.89.3}
\]
realizes these indices and the exact bounded exponent \(5^e\).
Its nonempty metabelian fibre has translation group
\(\prod_{\mathbb N}\mathbb Z_5\), by D.74.

The proof retains the actual \(\kappa(5)\)-graph action.
It does not extend this formula to a nonzero nonunit \(A\),
which imposes a different cyclotomic relation. It makes no
general section-existence or first higher-derived-obstruction claim.
\end{statement}

\begin{proof}
Use \(\Lambda=\mathbb Z_5[[u,v]]\), \(t_1=1+u\),
\(t_2=1+v\), \(w=t_1t_2-1\), and \(P=uvw\).
Let \(I_0\) be the actual untwisted Alexander operator image,
including \(\chi\). As in D.82, write \(\kappa_5=\kappa(5)\)
for its actual continuous weight-one cocycle; this is not the
weight-five Soul\'e character. A superscript \(D\) means
Cartier dual, and a star means ordinary Pontryagin dual after
the Cartier twist is accounted for.

\paragraph{The actual groups retain their common valuation graph.}
If the away-five class is zero, D.83 applies with its zero
lattice and proves the statement. Otherwise its closed
translation lattice is \(5^r\mathbb Z_5\) on the first axis.
The exact image and integral gauge construction in D.83/D.87
gives
\[
 I_s=5^r\mathbb Z_5\rtimes I_0,
 \qquad I_5=5^N\mathbb Z_5\rtimes I_0,                  \tag{D.89.4}
\]
where \(5^\infty\mathbb Z_5=0\). In these descriptions
the common embedded complement is the actual graph
\[
             g\longmapsto(A\kappa_5(g),g).
\]
Its action on every coefficient is
\[
                \rho_A(g)=t_1^{A\kappa_5(g)}\rho_0(g).
                                                               \tag{D.89.5}
\]
The integral local gauge lies in the actual global translation
lattice. Uniform \(c_0\) convergence gives the displayed
\(\lambda\) and \(N\ge r\), including infinite support.
No independent formal action has replaced the Galois image.

If \(e=0\), the actual images in (D.89.4) agree, so D.87
gives \(\mathcal S_s=0\). We will check the compact free
term at the end. Suppose henceforth \(e>0\), finite or infinite.
Put \(q_n=(1+u)^{5^n}-1\), including \(q_0=u\), and set
\[
 B=\Lambda/(q_r),\qquad
 R=\begin{cases}\Lambda/(q_N),&N<\infty,\\
                 \Lambda,&N=\infty,\end{cases}
 \qquad J=q_rR.                                         \tag{D.89.6}
\]
The actual coinduced inclusion and D.87 give
\[
 \mathcal S_s^*\simeq
       (J^D)^{I_0}/\operatorname{im}((R^D)^{I_0}),        \tag{D.89.7}
\]
with precisely the graph action (D.89.5).

\paragraph{The operator-difference ideal is exactly \((u)\).}
Let
\[
 N_0=\ker(\chi:I_0\to\mathbb Z_5^\times),\qquad
 N_1=\{g\in N_0:\kappa_5(g)=0\}.
\]
The already proved actual statements in D.77/D.82 supply
an element \(\delta\in N_1\) with
\(\beta_\delta-1=P\varepsilon\), where
\(\varepsilon\in\Lambda^\times\), and an element
\(\tau\in N_0\) with \(\kappa_5(\tau)=1\).
For clarity, the first assertion is the integral projection
of the primitive cubic character to the weight-three part of
the compact actual image of \((\kappa_5,\kappa_3)\):
the tame weight-one part is removed, while the cubic coordinate
remains a unit. The second uses the actual saturated image of
\(\kappa(5)\) on the cyclotomic kernel proved in D.82.
Both elements are actual operator-image elements, not assigned
values of unrelated characters.

On \(N_0\), the operator (D.89.5) is multiplication by
\(t_1^{A\kappa_5(g)}\beta_g\). Every \(\beta_g-1\)
is divisible by \(P\), by D.39. The element \(\delta\)
therefore puts \(P\) in the ideal generated by the operator
differences. The element \(\tau\), modulo that ideal,
puts \(t_1^A-1\) in it, since
\[
 t_1^A\beta_\tau-1=(t_1^A-1)+t_1^A(\beta_\tau-1).
\]
Conversely, every difference belongs to
\((P,t_1^A-1)\): for each \(c\in\mathbb Z_5\),
\(t_1^{Ac}-1\) is divisible by \(t_1^A-1\) by the
integral binomial identity. As \(A\) is a unit,
\((t_1^A-1)/u\) has unit constant coefficient \(A\).
Thus the exact closed ideal is
\[
       (\rho_A(g)-1:g\in N_0)=(P,t_1^A-1)=(u).         \tag{D.89.8}
\]
The equality is integral. In fact the two specified differences
already generate it as a \(\Lambda\)-ideal, so no closure
of a merely dense image is substituted for this equality.

Consequently, for \(X=R\) and for the actual submodule
\(X=J\), invariance of the discrete dual under \(N_0\)
means exactly annihilation of \(uX\). In particular
\[
                (R^D)^{N_0}=(R/uR)^D,
                \qquad (J^D)^{N_0}=(J/uJ)^D.            \tag{D.89.9}
\]

\paragraph{The natural invariant map vanishes.}
The map in (D.89.7) is induced by the actual inclusion
\(J=q_rR\hookrightarrow R\). Since \(q_r\) is
divisible by \(u\), its image lies in \(uR\).
Every \(N_0\)-invariant functional on \(R\) annihilates
this image by (D.89.9). Therefore
\[
                 (R^D)^{I_0}\longrightarrow(J^D)^{I_0}
                           \quad\text{is zero}.        \tag{D.89.10}
\]
This is an equality of the actual maps, not a dimension or
corank comparison. Equations (D.89.7)--(D.89.10) identify
the whole compact finite-localization term with
\[
                   \mathcal S_s\simeq(J/uJ)(-1)_\Gamma,
                   \qquad\Gamma=\mathbb Z_5^\times.    \tag{D.89.11}
\]
D.87 supplies the one finite global coefficient level for
each dual class, including every local place.

\paragraph{Finite local index and the retained generator.}
When \(N<\infty\), multiplication by the actual generator
\(q_r\) identifies the scalar module underlying \(J\)
with \(\Lambda/(q_N/q_r)\). Its quotient by \(u\) is
\[
                         J/uJ\simeq(\mathbb Z/5^e)[[v]].
                                                               \tag{D.89.12}
\]
Indeed \((q_N/q_r)(0)=5^{N-r}\); the distinguished
polynomial and cancellation details are proved in D.88.
The substitution on the retained generator is
\(q_r\mapsto q_rS_\chi(t_1^{5^r})\), which contributes
\(\chi\) at \(u=0\). The actual graph multiplier in
(D.89.5) is one at \(u=0\); the untwisted Alexander
multiplier there is \(\chi S_\chi(t_2)\).
After the Cartier untwist, (D.89.12) therefore has the exact
action
\[
       f(v)\longmapsto\chi S_\chi(1+v)
                       f((1+v)^\chi-1).                \tag{D.89.13}
\]
No division by \(5^r\) or \(5^e\) occurs.

Multiplication by \(v\) identifies this module with the
augmentation ideal in
\((\mathbb Z/5^e)[[\mathbb Z_5]]\), acted on by
\(t^x\mapsto\chi t^{\chi x}\). The missing constant
quotient has tame character \(\chi\), so tame averaging
identifies their compact coinvariants. As in D.88, the
invariant discrete dual consists of continuous functions
\(f(\gamma x)=\gamma^{-1}f(x)\), one arbitrary value
in \(\mathbb Z/5^e\) on each nonzero valuation shell,
zero at the origin, and only finitely many nonzero shell
values. Its compact dual is precisely
\(\prod_{j\ge0}\mathbb Z/5^e\). This proves the middle
case of (D.89.2), including its exact exponent and topology.

\paragraph{The infinite local index is computed directly.}
If \(N=\infty\), then \(R=\Lambda\) and
\(J=q_r\Lambda\). Since \(\Lambda\) is a domain,
division by its displayed generator is an actual scalar
identification, and
\[
                            J/uJ\simeq\mathbb Z_5[[v]].
                                                               \tag{D.89.14}
\]
The same retained-generator calculation gives (D.89.13).
Now the weighted continuous shell functions take values in
\(\mathbb Q_5/\mathbb Z_5\); continuity at the origin
again gives a direct sum, and their compact dual is
\(\prod_{j\ge0}\mathbb Z_5\). Together with the exact
zero map (D.89.10), this proves the infinite case. It is not
an extrapolation from finite exponents or infinite coranks.

\paragraph{No free local factor, and actual tangents.}
The global depth and local depth at five of the nonzero
boundary class \(a\) are both zero, because its valuation
coefficient \(A\) is a unit. The same holds for \(a-b=a\).
The class \(b=0\) contributes the identical countable
product on both sides of D.77's local/global boundary map.
Thus that map is an isomorphism, and the compact exact
sequence in D.63 identifies \(\mathcal K_s=\mathcal S_s\)
topologically. This also completes the zero-index and pure
valuation cases.

For (D.89.3), take \(A=1\). Its away-five coefficient
at two is \(5^r\), since \(1+5^{e+1}\equiv2\pmod4\).
Its local unit logarithm coefficient has valuation \(r+e\),
as proved and exactly checked in D.88. The zero local unit
case is realized by adding \(\kappa(5)\) to the completed
parameter in (D.88.24). These are actual cusp parameters,
so the full cusp construction supplies sections. Each has
finite valuation support in the stated examples, permitting
the D.74 fibre computation. The infinite torsion in the
finite-depth compact kernel is a second-cohomology feature,
not torsion in the first-cohomology translation group.
\end{proof}


% D.90. Root, FGH and PBC full proof/source/whole-certificate audits PASS.
\subsection{A nonunit five-valuation dominates the away lattice}

\begin{statement}{D.90}{Proved}{All compact elementary divisors when the five-valuation dominates}
Let
\[
 a=A\kappa(5)+\sum_{q\ne5}a_q\kappa(q),\qquad
 0\ne A\in\mathbb Z_5,\qquad a_q\longrightarrow0,
\]
and suppose the away-five class is nonzero. Put
\[
 \begin{aligned}
 m&=v_5(A),&r&=\min_{q\ne5}v_5(a_q),&0&\le m\le r,\\
 \lambda&=\sum_{q\ne5}a_q\frac{\log(q^4)}{4\log6},&
 N&=v_5(\lambda),&e&=N-r\in\mathbb Z_{\ge0}\cup\{\infty\}.
 \end{aligned}                                               \tag{D.90.1}
\]
For the actual cusp pair \(s=(a,0)\), the entire compact
kernel in D.63 is \(\mathcal K_s=\mathcal S_s\). Its compact
topological isomorphism type is
\[
 \mathcal S_s\simeq
 \begin{cases}
 0,&e=0,\\
 \displaystyle
 \left(\prod_{n\ge0}\mathbb Z/5^e\right)
 \times\prod_{j=1}^{\min(m,e-1)}(\mathbb Z/5^j)^2,
       &0<e<\infty,\\
 \displaystyle
 \left(\prod_{n\ge0}\mathbb Z_5\right)
 \times\prod_{j=1}^{m}(\mathbb Z/5^j)^2,
       &e=\infty.
 \end{cases}                                               \tag{D.90.2}
\]
Empty finite products are zero. The finite smaller-exponent
summands are retained. At \(m=0\) this recovers D.89.
The actual rational tangent \(630^5\) has
\((m,r,N,e)=(1,1,3,2)\), hence compact kernel
\[
             \left(\prod_{n\ge0}\mathbb Z/25\right)
                         \times(\mathbb Z/5)^2.        \tag{D.90.3}
\]
Every integer \(m\le r\) and positive finite integer \(e\)
is realized by the actual rational tangent
\(5^{5^m}(1+5^{e+1})^{5^r}\). The infinite endpoint is
realized by a finite-support completed cusp parameter.
Their actual metabelian fibres are nonempty and have
translation group \(\prod_{\mathbb N}\mathbb Z_5\), by D.74.

This is restricted to \(m\le r\). The intervening range
\(r<m<N\), arbitrary other dependent pairs, and the first
higher-derived obstruction are not settled here.
\end{statement}

\begin{proof}
Keep the actual notation \(\Lambda,u,v,t_1,t_2,w,P,I_0,N_0,N_1\)
of D.89. Set \(q_j=t_1^{5^j}-1\), including \(q_0=u\).
The actual common complement still acts as
\(\rho_A(g)=t_1^{A\kappa_5(g)}\rho_0(g)\).
Its two actual cyclotomic-kernel operators, one with
\(\kappa_5=0\) and primitive cubic coefficient and one
with \(\kappa_5=1\), give exactly
\[
       (\rho_A(g)-1:g\in N_0)=(P,t_1^A-1)=(P,q_m).
                                                               \tag{D.90.4}
\]
The proof is D.89's integral operator-ideal calculation before
its final use of \(A\) being a unit: if \(A=5^m\eta\),
then \((t_1^A-1)/q_m\) has unit constant coefficient
\(\eta\). Both generating elements are actual image elements.

\paragraph{The exact scalar module and its natural map.}
The actual global and local images are
\(5^r\mathbb Z_5\rtimes I_0\) and
\(5^N\mathbb Z_5\rtimes I_0\), with that same valuation
graph and \(5^\infty\mathbb Z_5=0\). If \(e=0\), they
agree and D.87 gives \(\mathcal S_s=0\).
Suppose \(e>0\). Put
\[
 R=\begin{cases}\Lambda/(q_N),&N<\infty,\\
                 \Lambda,&N=\infty,\end{cases}
                  \qquad J=q_rR.
\]
Every invariant functional on \(R\) annihilates \(q_mR\)
by (D.90.4). Since \(m\le r\), the actual inclusion
\(J\hookrightarrow R\) has image inside \(q_mR\).
Thus the natural invariant map is zero, exactly as a map,
and D.87 identifies
\[
 \mathcal S_s\simeq
              \bigl(J/(PJ+q_mJ)\bigr)(-1)_\Gamma,
                  \qquad \Gamma=\mathbb Z_5^\times.    \tag{D.90.5}
\]
In particular the whole finite-localization kernel is used;
no unrestricted everywhere-local kernel is substituted for it.

For finite \(N\), scalar division by the retained actual
generator \(q_r\) gives the relations
\(q_N/q_r,q_m,P\). Modulo \(q_m\), one has
\(t_1^{5^r}=1\), and hence the exact identity
\[
               q_N/q_r=\sum_{i=0}^{5^e-1}t_1^{i5^r}
                              \equiv5^e\pmod{q_m}.     \tag{D.90.6}
\]
At \(N=\infty\) there is no relative norm relation.
Consequently define
\[
 C=\begin{cases}\mathbb Z/5^e,&e<\infty,\\
                \mathbb Z_5,&e=\infty,\end{cases}
 \quad H_m=C[\mathbb Z/5^m],\quad
 B_C=H_m[[v]],\quad T_C=B_C/(uvw).                      \tag{D.90.7}
\]
Here \(H_0=C\). The compact module in (D.90.5) is
exactly \(T_C\) with the following actual action.
Substitution on \(q_r\) contributes
\(S_\chi(t_1^{5^r})=\chi\), which cancels the Cartier
untwist. The valuation graph multiplier is one modulo
\(q_m\). Therefore the action on the scalar module is
precisely the original \(\rho_0\), namely
\[
     S_\chi(t_1)S_\chi(t_2)\beta_g(t_1,t_2)
                            f(t_1^\chi,t_2^\chi).      \tag{D.90.8}
\]
It factors through \(\Gamma\) modulo \(P\).
This computation retains \(q_r\); treating it as a
weight-zero basis element would change the answer.

\paragraph{The exact torsion layer and zero snake map.}
Both \(v\) and \(w\) are nonzero divisors in \(B_C\).
For \(w\), the continuous change of variables
\(v\mapsto u+v+uv\), fixing \(u\), is an automorphism
of the completed ring and preserves the relation \(q_m\).
Thus it carries the injective multiplication by \(v\)
to multiplication by \(w\). In particular \(vw\)
is a nonzero divisor. We have an exact sequence
\[
 0\longrightarrow E_C\longrightarrow T_C\longrightarrow Q_C
                 \longrightarrow0,
 \quad E_C=vwB_C/(uvwB_C)\simeq C[[v]],
 \quad Q_C=B_C/(vw).                                    \tag{D.90.9}
\]
The action on \(E_C\), using its actual scalar generator
\(vw\), is \(\chi S_\chi(1+v)^3\) followed by
substitution. D.88's exact weighted-shell argument gives
\[
                         (E_C)_\Gamma\simeq\prod_{n\ge0}C.
                                                               \tag{D.90.10}
\]
For \(C=\mathbb Z_5\) this uses continuous functions valued
in \(\mathbb Q_5/\mathbb Z_5\); for finite \(C\) it uses
the finite coefficient dual. Both have the stated compact topology.

The quotient \(Q_C/uQ_C=C[[v]]/(v^2)\) has successive
tame characters \(\chi^2,\chi^3\). Thus every
\(\Gamma\)-invariant element of \(Q_C\) lies in \(uQ_C\).
The natural map \(uT_C\to uQ_C\) is an equivariant
isomorphism: modulo \(u\), the element \(vw\) is
\(v^2\), a nonzero divisor in \(C[[v]]\), so
\((u)\cap(vw)=(uvw)\). Every invariant lifts invariantly.
The snake lemma for an actual lift of the topological
generator \(2\in\Gamma\) therefore gives
\[
 0\longrightarrow(E_C)_\Gamma\longrightarrow(T_C)_\Gamma
                  \longrightarrow(Q_C)_\Gamma\longrightarrow0.
                                                               \tag{D.90.11}
\]
This exact injection remains valid over every \(C\) in (D.90.7).

\paragraph{The additional axis-intersection kernel is retained.}
The two-axis map for \(Q_C\) is not injective before tame
averaging, because \(u\) is a zero divisor in \(H_m\).
Write
\[
                    \mathcal N_m=\sum_{i=0}^{5^m-1}t_1^i.
\]
The literal cyclic coefficient equations give
\(\operatorname{Ann}_{H_m}(u)=C\mathcal N_m\).
Moreover \(w\mathcal N_m=v\mathcal N_m\).
If \(vf=wg\), then modulo \(v\),
\(u g(u,0)=0\), so \(g=c\mathcal N_m+vh\).
It follows that the kernel of the two-axis map is
generated by \(v\mathcal N_m\). This is an actual copy
of \(C\): if \(c v\mathcal N_m\) were in \((vw)\),
cancel \(v\) and evaluate on \(w=0\) to get
\(c\mathcal N_m=0\), hence \(c=0\).
There is consequently an exact sequence
\[
 0\longrightarrow C\,v\mathcal N_m\longrightarrow Q_C
 \longrightarrow H_m\oplus H_m\longrightarrow C\longrightarrow0.
                                                               \tag{D.90.12}
\]
The last map is the difference of augmentations.

The first term has exact tame character \(\chi^3\).
Indeed \(\sigma_\chi\mathcal N_m=\mathcal N_m\),
\(u\mathcal N_m=0\), and
\(v^2\mathcal N_m=vw\mathcal N_m=0\) in \(Q_C\).
Using (D.90.8) therefore gives
\(v\mathcal N_m\mapsto\chi^3 v\mathcal N_m\).
The final term has character \(\chi^2\). Both their
invariant and coinvariant parts under the order-four tame
subgroup vanish. Exact averaging removes these two terms
before taking the remaining procyclic coinvariants.

On the first boundary the action is
\(\chi S_\chi(t_1)\sigma_\chi\); on the diagonal it is
\(\chi S_\chi(t_1)t_1^{(1-\chi)/2}\sigma_\chi\).
The integral gauge \(t_1^{-1/2}\) identifies them, by D.39.
Multiplication by \(u\) transports the first action to
\(\chi\sigma_\chi\) on \(uH_m\).
Its kernel \(C\mathcal N_m\) has character \(\chi^2\);
the quotient \(H_m/uH_m\) for the target weighted action
has character \(\chi\). Exact tame averaging therefore
identifies each boundary's coinvariants with those of the
literal weighted permutation module
\[
                H_m,\qquad t^x\longmapsto\chi t^{\chi x}.
                                                               \tag{D.90.13}
\]

The zero orbit contributes nothing, since \(2-1\) is a unit.
The orbit of valuation \(i<m\) has length
\(4\cdot5^{m-i-1}\). Eliminating the basis vectors along
its actual weighted permutation cycle leaves the relation
\(2^{4\cdot5^{m-i-1}}-1\), whose five-valuation is
\(m-i\). Thus, naturally with respect to coefficient
reduction, (D.90.12)--(D.90.13) give
\[
                  (Q_C)_\Gamma\simeq
                            \bigoplus_{j=1}^{m}(C/5^jC)^2.
                                                               \tag{D.90.14}
\]
This integral calculation explains the smaller finite cyclic
factors; a rational rank computation would lose all of them.

\paragraph{Finite coefficients: split every finite-order lift.}
Let \(0<e<\infty\), so \(C=\mathbb Z/5^e\).
The entire module \((T_C)_\Gamma\) is killed by \(5^e\).
Choose a cyclic generator of order \(5^d\) in (D.90.14),
where \(d=\min(e,j)\), and any lift \(x\) to
\((T_C)_\Gamma\). Then \(5^d x\in(E_C)_\Gamma\)
and is killed by \(5^{e-d}\). In the actual product
\((E_C)_\Gamma\simeq\prod C\),
\[
      \operatorname{Ann}(5^{e-d})=5^d\prod C.
\]
Choose \(y\) in that product with \(5^dy=5^dx\).
Now \(x-y\) is a lift of the original cyclic generator
with order dividing \(5^d\). Doing this for the finitely
many independent cyclic generators produces a section of
(D.90.11). Its finite domain makes the section continuous;
the resulting compact-to-Hausdorff bijection is a topological
isomorphism. Hence
\[
 (T_C)_\Gamma\simeq
       \left(\prod_{n\ge0}\mathbb Z/5^e\right)
             \times\prod_{j=1}^{m}(\mathbb Z/5^{\min(e,j)})^2.
                                                               \tag{D.90.15}
\]
Only the factors with exponent exactly \(5^e\) can be
absorbed into the countable product. Retaining the others
gives the finite case of (D.90.2).

\paragraph{Infinite coefficients: the exact reduction maps split the extension.}
Now \(C=\mathbb Z_5\). Put
\(M=(T_{\mathbb Z_5})_\Gamma\),
\(E=(E_{\mathbb Z_5})_\Gamma\), and
\(Q=(Q_{\mathbb Z_5})_\Gamma\). Equation (D.90.11) is
\[
 0\longrightarrow E\simeq\prod\mathbb Z_5
       \longrightarrow M\longrightarrow
       Q\simeq\prod_{j=1}^{m}(\mathbb Z/5^j)^2
                                  \longrightarrow0.    \tag{D.90.16}
\]
It is not permissible to infer its splitting just from these
abstract end groups. We use the actual coefficient reductions.
For \(h\ge m\) (take \(h\ge1\) if \(m=0\)), reduction
of the displayed rings and operators gives
\[
 M/5^hM=(T_{\mathbb Z/5^h})_\Gamma,
 \qquad E/5^hE=(E_{\mathbb Z/5^h})_\Gamma.              \tag{D.90.17}
\]
These are the natural quotient maps: taking coinvariants is
the quotient by the image of the same actual operator minus
one and therefore commutes with coefficient reduction.
That image is compact, hence closed. The explicit cycle
calculation (D.90.14) identifies the corresponding map on
\(Q\) with \(Q\to Q/5^hQ=Q\).

The zero-snake injection at the finite coefficient level in
(D.90.11) implies
\[
                         E\cap5^hM=5^hE.
\]
Since \(5^hQ=0\), one also has \(5^hM\subset E\).
Consequently
\[
                              5^hM=5^hE.               \tag{D.90.18}
\]
The finite-coefficient splitting already proved gives an
order-\(5^j\) lift of each cyclic generator of \(Q\)
in \(M/5^hM\). Lift it to \(x\in M\). Then
\(5^j x\in5^hM=5^hE\subset5^jE\), because \(j\le m\le h\).
An adjustment by an element of \(E\) makes \(5^j x=0\)
without changing its image in \(Q\). Choosing the finitely
many such lifts gives an actual continuous section of
(D.90.16). This proves its splitting and the infinite case
of (D.90.2). All maps here are the actual coefficient
reductions; neither matching cardinalities nor an unproved
interchange of limits is used.

\paragraph{The complete compact kernel and actual examples.}
Both the global depth and the full local depth of \(a\)
are \(m\): its away depth is \(r\ge m\), its unit local
depth is \(N\ge r\), and its five-valuation depth is
\(m\). The same holds for \(a-b=a\). The zero boundary
\(b\) contributes an identical countable product on both
sides. D.77's natural boundary map is therefore an
isomorphism, and D.63 gives \(\mathcal K_s=\mathcal S_s\).
This also completes the zero-index case.

For the stated finite examples take \(A=5^m\) and the
away class \(5^r\kappa(1+5^{e+1})\). Its valuation at
two is exactly \(5^r\), and its local unit depth is
\(r+e\), by the proved D.88 logarithm calculation.
For \(m=r=1,e=2\), the tangent is
\(5^5\cdot126^5=630^5\), giving (D.90.3).
For infinite index take
\[
 a=5^m\kappa(5)+
    5^r\bigl(\lambda_3\kappa(2)-\lambda_2\kappa(3)\bigr),
 \qquad \lambda_q=\log(q^4)/(4\log6).
\]
The away class is nonzero of depth \(r\) and has zero
local coefficient. These are actual full cusp parameters;
their finite valuation support permits the D.74 fibre
calculation. No assertion about all other abelian fibres follows.

The literal integral boundary matrices, exact norm-line and
half-gauge identities, and Smith forms through cyclic depth
three are retained in \path{code/de_sha_dominant_valuation_cusp.py}
and its executed outputs. These finite checks support the
actual formulas; the extension and infinite-endpoint proofs
above supply the all-depth conclusion.
\end{proof}


% D.91. Root, FGH and PBC complete proof/source/certificate audits PASS.
\subsection{The compact kernel for every completed cusp parameter}

\begin{statement}{D.91}{Proved}{The complete compact-kernel classification on the three cusp packets}
Let
\[
 a=A\kappa(5)+\sum_{q\ne5}a_q\kappa(q),\qquad
 A,a_q\in\mathbb Z_5,\qquad a_q\longrightarrow0,
                                                               \tag{D.91.1}
\]
and consider the actual cusp pair \(s=(a,0)\). The assertion
includes \(a=0\). If all away-five coefficients vanish, then
\(\mathcal S_s=\mathcal K_s=0\), for every \(A\), including
\(A=0\). Thus the trivial Kummer parameter, represented by
the rational tangents \(1\) and \(-1\), is included separately.
Suppose henceforth that the away-five class is nonzero, and put
\[
 \begin{gathered}
 r=\min_{q\ne5}v_5(a_q),\qquad m=v_5(A)\in
                  \mathbb Z_{\ge0}\cup\{\infty\},\\
 \lambda=\sum_{q\ne5}a_q\frac{\log(q^4)}{4\log6},\qquad
 N=v_5(\lambda)\ge r,\qquad e=N-r.
 \end{gathered}                                                \tag{D.91.2}
\]
If \(e=0\), again \(\mathcal S_s=\mathcal K_s=0\).
If \(0<e<\infty\), define
\[
 H=\min\{e,\max(m-r,0)\},\qquad t=\min(m,r),
 \quad H=e\text{ when }m=\infty.
                                                               \tag{D.91.3}
\]
Then, as compact topological modules,
\[
 \begin{aligned}
 \mathcal S_s&\simeq
 \left(\prod_{n\ge0}\mathbb Z/5^e\right)
 \times\prod_{j=1}^{\min(t,e-H-1)}
                      (\mathbb Z/5^{H+j})^2,\\
 \mathcal K_s&\simeq\mathcal S_s\times\mathbb Z_5^{\,2H}.
 \end{aligned}                                                \tag{D.91.4}
\]
A finite product with upper limit less than one is zero.
If \(N=\infty\) and \(m<\infty\), put
\(h=\max(m-r,0)\) and \(t=\min(m,r)\). Then
\[
 \begin{aligned}
 \mathcal S_s&\simeq
 \left(\prod_{n\ge0}\mathbb Z_5\right)
          \times\prod_{j=1}^{t}(\mathbb Z/5^{h+j})^2,\\
 \mathcal K_s&\simeq\mathcal S_s\times\mathbb Z_5^{\,2h}.
 \end{aligned}                                                \tag{D.91.5}
\]
If \(N=m=\infty\), both \(\mathcal S_s\) and
\(\mathcal K_s\) are isomorphic to \(\prod_{n\ge0}\mathbb Z_5\).

The subgroup \(\mathcal S_s\subset\mathcal K_s\) and
the map \(L_5\to\mathcal T_s\) are the actual natural
maps of D.63. The displayed splittings and identifications
with products are continuous and noncanonical. In (D.91.5)
one may absorb the finite free factor into the countable
free product as an abstract compact-module isomorphism;
this does not identify that factor with the natural kernel
\(\mathcal S_s\).

The same classification applies to all three peripheral
markings: the actual tripod automorphisms transport
\((a,0)\) to \((0,a)\) and to \((-a,-a)\), together
with the tangential frame. They induce isomorphisms of the
entire global/local cohomology diagrams. In particular no
zero-parameter packet at one of the other two cusps is
omitted or identified with another section packet.

Every parameter above has an actual cusp section. For finite
indices the rational tangent
\[
                5^{5^m}(1+5^{e+1})^{5^r}
                                                               \tag{D.91.6}
\]
realizes every \(r,m\ge0\) and \(e\ge1\); omit the first
factor for \(m=\infty\). For example, \(5^{25}626^5\)
has \((r,m,N,e)=(1,2,4,3)\), and hence
\[
 \mathcal S_s\simeq
 \left(\prod_{n\ge0}\mathbb Z/125\right)\times(\mathbb Z/25)^2,
 \qquad \mathcal K_s\simeq\mathcal S_s\times\mathbb Z_5^2.
                                                               \tag{D.91.7}
\]
Their finitely ramified metabelian fibres have translation
group \(\prod_{\mathbb N}\mathbb Z_5\), by D.74.
The theorem computes the compact second-cohomology kernel
on the cusp families, not every abelian pair, every
infinitely ramified first-cohomology rank, or a first
higher-derived obstruction.
\end{statement}

\begin{proof}
Write \(\Lambda=\mathbb Z_5[[u,v]]\), \(t=1+u\),
\(w=t(1+v)-1\), \(P=uvw\), and \(q_j=t^{5^j}-1\),
including \(q_0=u\). Every quotient below is compact,
with its quotient topology. All actions are the actual
operator actions of D.35, D.39 and D.87.

\paragraph{Degeneracies and the ranges already computed.}
If the away-five class is zero, the actual global and local
images are the same graph of \(A\kappa(5)\) over \(I_0\).
D.83 and D.87 give \(\mathcal S_s=0\), and D.77 gives
\(\mathcal K_s=0\). This includes \(a=0\); it is not
obtained by subtracting two infinite valuations.
For nonzero away class the actual images are
\[
 I_s=5^r\mathbb Z_5\rtimes I_0,\qquad
 I_5=5^N\mathbb Z_5\rtimes I_0,
                                                               \tag{D.91.8}
\]
with the same graph action
\(\rho_A(g)=t^{A\kappa_5(g)}\rho_0(g)\).
Here \(5^\infty\mathbb Z_5=0\). D.83's uniformly
convergent integral gauge belongs to the actual global
translation lattice, so (D.91.8) holds also for infinite
\(c_0\) support. If \(N=r\), the images coincide.

For \(A=0\), D.88 proves the asserted answer, including
\(N=\infty\). For finite \(m\le r\), D.90 proves it.
If \(r\le N\le m<\infty\), the graph multiplier is
already one on \(R=\Lambda/(q_N)\) and on \(J=q_rR\).
The entire natural invariant map in D.87 is therefore
literally the D.88 map, not just a map between isomorphic
abstract modules. D.88 computes \(\mathcal S_s\), and
D.77 adds its free factor \(\mathbb Z_5^{2(N-r)}\).
It remains to prove the genuinely new range
\[
                 0\le r<m<N\le\infty.
                                                               \tag{D.91.9}
\]
Put \(h=m-r>0\). For finite \(N\), write \(b=N-m>0\),
so \(e=h+b\).

\paragraph{The actual scalar module and the map that must be kept.}
The two actual operators of D.89/D.90 give
\[
       (\rho_A(g)-1:g\in\ker\chi)=(P,q_m).
                                                               \tag{D.91.10}
\]
Let \(R=\Lambda/(q_N)\) for finite \(N\), and \(R=\Lambda\)
otherwise; put \(J=q_rR\). D.87 identifies \(\mathcal S_s\)
with the kernel of the natural compact map
\[
 M:=\bigl(J/(PJ+q_mJ)\bigr)(-1)_\Gamma
       \longrightarrow
 U:=\bigl(R/(PR+q_mR)\bigr)(-1)_\Gamma,
 \qquad \Gamma=\mathbb Z_5^\times.
                                                               \tag{D.91.11}
\]
The map is induced by \(J\hookrightarrow R\).
In particular it is not discarded when \(r<m\).
The target is the ordinary dual of the actual local
Cartier-invariant group, so D.77 identifies it with
\(L_5\); it is torsion-free.

Retain \(q_r\) as the actual generator of \(J\).
For finite \(N\), set \(F=q_m/q_r\). The exact norm
identity modulo \(q_m\) gives
\[
 q_N/q_r\equiv5^b F\pmod{q_m},\qquad
 B=\mathbb Z_5[[u,v]]/(q_m,5^bF),\qquad T=B/(P).
                                                               \tag{D.91.12}
\]
For \(N=\infty\), omit the relation \(5^bF\) and put
\(B=\mathbb Z_5[\mathbb Z/5^m][[v]]\), again \(T=B/(P)\).
After the Cartier untwist, the action on the scalar module
\(T\) is
\[
 \chi(g)^{-1}S_{\chi(g)}(t^{5^r})\rho_0(g).
                                                               \tag{D.91.13}
\]
The graph multiplier is one modulo \(q_m\). Thus
\(M=T_\Gamma\) with precisely (D.91.13). The target
has scalar coefficient \(\Lambda/(q_m,P)\), action
\(\chi^{-1}\rho_0\), and the map in (D.91.11) is
\emph{multiplication by \(q_r\)}.

\paragraph{The infinite layer injects before any boundary calculation.}
Let \(C=\mathbb Z/5^e\) for finite \(N\), and
\(C=\mathbb Z_5\) for \(N=\infty\). Since \(F(0)=5^h\),
we have \(B/uB=C[[v]]\). Both \(v\) and \(w\) are
nonzero divisors in \(B\): multiplication by \(v\) is
injective, and the continuous change of variables
\(v\mapsto u+(1+u)v\) in the complete \((5,u,v)\)-adic
ring carries it to multiplication by \(w\). Consequently
\[
 0\longrightarrow E\longrightarrow T\longrightarrow Q\longrightarrow0,
 \qquad E=vwB/(uvwB)\simeq C[[v]],\quad Q=B/(vw).
                                                               \tag{D.91.14}
\]
The retained generator \(q_rvw\) gives on \(E\) the
action \(\chi S_\chi(1+v)^3\) followed by substitution.
D.88's exact weighted-shell calculation gives
\(E_\Gamma\simeq\prod_{n\ge0}C\).

The two successive tame characters of \(Q/uQ\) are
\(\chi^2\) and \(\chi^3\). Therefore every invariant
element of \(Q\) lies in \(uQ\). Moreover
\(uT\to uQ\) is an equivariant isomorphism: modulo
\(u\), the element \(vw\) becomes the nonzero divisor
\(v^2\) in \(C[[v]]\), so \((u)\cap(vw)=(uvw)\).
Every invariant lifts invariantly. The snake lemma for
the actual lift of \(2\in\Gamma\) yields the exact sequence
\[
       0\longrightarrow E_\Gamma\longrightarrow M
                 \longrightarrow Q_\Gamma\longrightarrow0.
                                                               \tag{D.91.15}
\]
The same proof applies to coefficient reductions modulo
\(5^H\) in the infinite case. This is an actual injection,
not a deduction from a comparison of ranks.

\paragraph{Axis intersections and the literal boundary modules.}
Write \(\mathcal N_m=q_m/u\). In the finite case
\[
                 \operatorname{Ann}_B(u)
                   =(\mathbb Z/5^b)[[v]]\mathcal N_m.
                                                               \tag{D.91.16}
\]
Indeed, if \(uf=F(q_rg+5^b k)\) in the original
torsion-free power-series ring, evaluation at \(u=0\)
forces \(k(0,v)=0\). Cancel \(u\) to express \(f\),
modulo the relations, as \(\mathcal N_m g\). Since
\(u\mathcal N_m=0\), only \(g(0,v)\) remains. Its
coefficient has exact modulus \(5^b\): cancel \(F\)
and reduce modulo \(q_r\), where the norm
\(q_r/u\) has every cyclic coefficient equal to one.
For \(N=\infty\), the same literal cyclic coefficient
argument gives \(\operatorname{Ann}_B(u)=\mathbb Z_5[[v]]\mathcal N_m\).

The two-axis map for \(Q\) consequently has kernel
\(Cv\mathcal N_m\), with \(C\) here replaced by
\(\mathbb Z/5^b\) in the finite case, and cokernel
\(B/(u,v)\). The first has exact tame weight \(\chi^3\),
the second weight \(\chi^2\). To see the kernel directly,
write \(vf=wg\), reduce modulo \(v\), use (D.91.16),
and note \(w\mathcal N_m=v\mathcal N_m\) and
\(v^2\mathcal N_m=0\) in \(Q\). Both kernel and
cokernel have zero invariants and coinvariants under the
order-four tame subgroup. They therefore contribute no
\(\Gamma\)-homology.

Let \(R_j=\mathbb Z_5[\mathbb Z/5^j]\) and
\(J_{r,N}=\ker(R_N\to R_r)\). On the first axis,
(D.91.13) is
\(S_\chi(t)S_\chi(t^{5^r})\sigma\), with
\(\sigma(t)=t^\chi\). Multiplication by \(u q_r\)
transports it exactly to the ordinary permutation action on
\(uJ_{r,N}/q_m uJ_{r,N}\). Multiplication by \(u\)
is injective on \(J_{r,N}\): its only annihilator in
\(R_N\) is the full norm, which has nonzero image in
the torsion-free ring \(R_r\). The comparison
\[
       uJ_{r,N}/q_m uJ_{r,N}
                      \longrightarrow J_{r,N}/q_mJ_{r,N}
                                                               \tag{D.91.17}
\]
has kernel \(q_mJ_{r,N}/u q_mJ_{r,N}\), of tame
weight \(\chi^2\), and cokernel \(J_{r,N}/uJ_{r,N}\),
of tame weight \(\chi\). They vanish in all
\(\Gamma\)-homology. The diagonal axis has the same
comparison after multiplication by \(t^{-1/2}\), from
the exact reflection identity of D.39. Thus
\[
 Q_\Gamma\simeq
          \left((J_{r,N}/q_mJ_{r,N})_\Gamma\right)^2.
                                                               \tag{D.91.18}
\]
At \(N=\infty\), use \(J_{r,\infty}=q_r\mathbb Z_5[[t-1]]\);
the same argument and tame weights apply.

\paragraph{The middle boundary is a nonsplit cyclic extension.}
For finite \(N\), let \(D=\mathbb Z/5^b\) act on
\(R_N\) by translation by \(5^m\), and let
\(\mathcal Q=D\rtimes\Gamma\). Its action on \(R_r\)
has trivial \(D\)-part. Compact group homology of
\(0\to J_{r,N}\to R_N\to R_r\to0\) gives
\[
 H_1(\mathcal Q,R_N)\longrightarrow H_1(\mathcal Q,R_r)
 \longrightarrow (J_{r,N})_{\mathcal Q}
 \longrightarrow (R_N)_{\mathcal Q}\longrightarrow(R_r)_{\mathcal Q}
 \longrightarrow0.                                           \tag{D.91.19}
\]
These are permutation modules on finite sets. For complete
control of continuity and the map, dualize to the discrete
modules of functions valued in \(\mathbb Q_5/\mathbb Z_5\).
Continuous Shapiro at each open point stabilizer identifies
first cohomology with its continuous characters. Dualizing
back gives the pro-five abelianization of that stabilizer.
The finite-index cochain proof is the usual Shapiro map
\cite[II, Proposition 1.11 and Example 1.27(a)]{PBC_MilneCFT};
it applies to continuous cochains and discrete coefficients.
Pullback of functions along an equivariant point-set map
becomes restriction of characters. Its compact dual is
therefore the stabilizer inclusion, not a transfer. This
also constructs the exact sequence (D.91.19) directly by
dualizing the discrete cohomology sequence.

There are \(m+1\) upper orbits, the valuation shells
\(0,\ldots,m-1\) and the ball divisible by \(5^m\),
and \(r+1\) lower orbits. The kernel on free orbit
coinvariants is \(\mathbb Z_5^h\). For a lower nonzero
orbit represented by \(5^j\), put \(k=r-j\), so
\(1\le k\le r\). Its stabilizer is
\(D\rtimes(1+5^k\mathbb Z_5)\), whose pro-five
abelianization is
\[
             \mathbb Z/5^{\min(b,k)}\oplus\mathbb Z_5.
                                                               \tag{D.91.20}
\]
The sole upper orbit over it has stabilizer the graph
\(d=(1-\chi)/5^{k+h}\), with
\(\chi\in1+5^{k+h}\mathbb Z_5\). Choose its generator
\(\chi=(1+5^k)^{5^h}\). The actual stabilizer inclusion
maps it to \((d,5^h)\), and \(d\) is a five-adic unit.
Its cokernel in (D.91.20) is therefore a \emph{single}
cyclic group of order \(5^{h+\min(b,k)}\).
For the lower zero orbit, tame conjugation kills the
translation contribution to the abelianization, and the
upper zero-orbit stabilizer \(\Gamma\) surjects on the
remaining first homology. There is no contribution there.
It follows that
\[
 (J_{r,N}/q_mJ_{r,N})_\Gamma
       \simeq\mathbb Z_5^h\oplus
                       \bigoplus_{k=1}^r\mathbb Z/5^{h+\min(b,k)}.
                                                               \tag{D.91.21}
\]
The finite free quotient splits continuously. In particular
\((r,m,N)=(1,2,3)\) gives \(\mathbb Z_5\oplus\mathbb Z/25\),
not \(\mathbb Z_5\oplus(\mathbb Z/5)^2\).

\paragraph{The actual boundary map has exactly that torsion kernel.}
On the target first axis the untwisted action is
\(S_\chi(t)\sigma\). Multiplication by \(u\) transports
it to ordinary permutation on \(uR_m\), with a kernel
of tame weight \(\chi\). The map \(uR_m\to R_m\)
is injective on coinvariants: the fixed element \(\delta_0\)
makes \(H_1(\Gamma,R_m)\to H_1(\Gamma,\mathbb Z_5)\)
surjective. Its image is the kernel of augmentation in
the free orbit module \((R_m)_\Gamma\).

The source and target transports commute with the actual
map \(f\mapsto q_rf\): both composites are multiplication
by \(u q_r\). The diagonal half-gauges commute with
\(q_r\) as well. Thus the induced boundary map is
exactly the literal reduction and inclusion
\[
 (J_{r,N}/q_mJ_{r,N})_\Gamma\longrightarrow(R_m)_\Gamma.
                                                               \tag{D.91.22}
\]
In (D.91.19), it kills the finite homology cokernel and
identifies the free orbit-kernel quotient with
\(\ker((R_m)_\Gamma\to(R_r)_\Gamma)\). Its kernel
is precisely the finite torsion in (D.91.21), with no
kernel on the free quotient. The target axis-intersection
kernel and cokernel have tame weights \(\chi^2\) and
\(\chi\), so the same assertion holds for both actual
boundary axes. No comparison by abstract coranks is used.

\paragraph{Finite depth: the entire extension is killed before it is split.}
Let \(F_b\) denote two copies of the finite torsion in
(D.91.21). Equations (D.91.15) and (D.91.22) show that
the kernel in (D.91.11) is contained in the preimage
of \(F_b\). That preimage is exactly the torsion subgroup
of \(M\), because \(E_\Gamma\) has exponent \(5^e\)
and the remaining quotient is finite free. Conversely all
this torsion maps to zero in the torsion-free actual target
\(U\). Therefore
\[
  \mathcal S_s=M_{\rm tor},\qquad
 0\longrightarrow\prod_{n\ge0}\mathbb Z/5^e
   \longrightarrow\mathcal S_s\longrightarrow F_b\longrightarrow0.
                                                               \tag{D.91.23}
\]
The finite-index actual images in (D.91.8) have index
\(5^e\). Restriction followed by corestriction on their
actual continuous first cohomology, with its actual
coefficient action, is multiplication by \(5^e\).
D.88 gives the explicit finite-coset proof, including
nonnormal subgroups; see also \cite[II, Proposition 1.30]{PBC_MilneCFT}.
Consequently the \emph{whole} group \(\mathcal S_s\)
in (D.91.23) is killed by \(5^e\).

For a cyclic quotient generator of order \(5^c\), choose
a lift \(x\). Then \(5^cx\) lies in the product layer
and is annihilated by \(5^{e-c}\). In that product,
\(\operatorname{Ann}(5^{e-c})=5^c E_\Gamma\). Subtract
an element of the layer to make the lift have order
\(5^c\). Lifting the finitely many cyclic generators
this way splits (D.91.23) continuously. The factors with
\(h+\min(b,k)=e\) are absorbed by the countable product
of \(\mathbb Z/5^e\); the smaller factors are retained.
This proves the strict-middle instance of (D.91.4).

\paragraph{The infinite boundary and its actual reduction maps.}
Keep \(m,r\) fixed and let the finite integer \(b\)
increase. The scalar boundary modules are
\(R_m/(5^bF)\) with the fixed retained boundary action
\(\varphi_r=S_\chi(t)S_\chi(t^{5^r})\sigma\). Their
inverse limit is \(R_m\): their defining ideals tend
uniformly to zero. Compactness also gives
\[
  (R_m,\varphi_r)_\Gamma
       \simeq\varprojlim_b(R_m/(5^bF),\varphi_r)_\Gamma.
                                                               \tag{D.91.24}
\]
Explicitly, compact intersections give surjectivity onto
compatible quotient cosets, and
\(\bigcap_b((2-1)R_m+5^bFR_m)=(2-1)R_m\), where
\(2-1\) denotes the actual operator difference. Its
image is compact and closed; the ideals tend to zero.

The stabilizer maps in (D.91.20) reduce the translation
coordinate and leave the free orbit-kernel map unchanged.
For \(b\ge r\), all finite cyclic factors have stabilized,
and the transition maps are isomorphisms: they are onto,
are the identity on the free orbit quotient, and have
equal finite torsion orders. Hence the infinite boundary is
\[
 (J_{r,\infty}/q_mJ_{r,\infty})_\Gamma
       \simeq\mathbb Z_5^h\oplus
                              \bigoplus_{k=1}^r\mathbb Z/5^{h+k}.
                                                               \tag{D.91.25}
\]
Its natural map to \((R_m)_\Gamma\) still kills exactly
the finite torsion and embeds the free orbit part, by
the same actual maps, or by their stabilized limit.

\paragraph{Integral torsion lifts at infinity, and the actual kernel.}
Now \(E_\Gamma=\prod\mathbb Z_5\). Put
\(F_\infty=\bigoplus_{k=1}^r(\mathbb Z/5^{h+k})^2\).
Split the finite free orbit quotient of \(M\), and let
\(M_0\) be the preimage of \(F_\infty\). Then
\[
 0\longrightarrow E_\Gamma\longrightarrow M_0
              \longrightarrow F_\infty\longrightarrow0,
 \qquad M\simeq M_0\times\mathbb Z_5^{2h}.
                                                               \tag{D.91.26}
\]
For every \(H\ge1\), reduction of the actual scalar
module modulo \(5^H\) commutes with coinvariants. The
zero-snake injection proved after (D.91.15) remains valid
over \(\mathbb Z/5^H\). It gives the exact purity identity
\[
             E_\Gamma\cap5^H M=5^H E_\Gamma.
                                                               \tag{D.91.27}
\]
For \(H\ge m\), the quotient \(F_\infty\) is killed
by \(5^H\), so \(5^HM_0\subset E_\Gamma\). Equation
(D.91.27) therefore gives \(5^HM_0=5^HE_\Gamma\).
The finite-coefficient extension
\(M_0/5^HM_0\) is killed by \(5^H\), has the product
layer \(E_\Gamma/5^HE_\Gamma\), and quotient
\(F_\infty\). The same finite-order adjustment as above
provides a lift of each quotient generator of its exact
order modulo \(5^HM_0\). Lift it integrally. If its
order is \(5^c\), its \(5^c\)-multiple lies in
\(5^HM_0=5^HE_\Gamma\subset5^cE_\Gamma\); one more
integral adjustment gives an actual order-\(5^c\) lift.
Thus (D.91.26) splits continuously, retaining its finite
torsion rather than inferring a splitting from abstract ranks.

The actual map in (D.91.11) kills the whole product layer:
its actual generator maps as
\[
                      q_rvw=P\,(q_r/u)=0
                          \quad\text{in }\Lambda/(P,q_m).
                                                               \tag{D.91.28}
\]
It also kills the finite torsion, because its target is
torsion-free. Conversely (D.91.22)/(D.91.25) embeds the
free orbit quotient, so nothing outside \(M_0\) can
lie in its kernel. Therefore
\(\mathcal S_s=M_0\simeq E_\Gamma\times F_\infty\).
This proves the strict-middle case of (D.91.5) by actual
maps, including the simultaneous absence of any local
unit translation.

\paragraph{The remaining boundary factor, all cusp markings, and examples.}
D.77 identifies the zero-parameter boundary map with the
identity on its countable product. The other two boundary
depths are globally \(\min(m,r)\) and locally
\(\min(m,N)\). For finite \(N\) their difference is
exactly \(H\) in (D.91.3); for infinite \(N\) and
finite \(m\) it is \(h\). This gives the natural compact
quotient \(\mathbb Z_5^{2H}\), or \(\mathbb Z_5^{2h}\),
and the continuous noncanonical split sequence of D.63.
When \(m=N=\infty\), D.88 gives the stated product
identifications with their precise natural-map interpretation.

The actual automorphism \(t\mapsto1-t\) exchanges the
first two cusp packets. The actual automorphism
\(t\mapsto1/t\) transforms the Kummer pair as
\((a,b)\mapsto(-a,b-a)\), because \(\kappa(-1)=0\)
in pro-five Kummer cohomology. It therefore sends
\((a,0)\) to \((-a,-a)\). Transport the tangential
frame along the automorphism; changing back to the standing
frame is an inner gauge. These are actual isomorphisms of
the twisted Alexander modules and of every finite global
and local cohomology map. Replacing \(a\) by \(-a\)
does not change \(m,r,N\). This proves the asserted
three-packet scope, without identifying their sections.

For (D.91.6), \(1+5^{e+1}\equiv2\pmod4\), so
its exponent at two is exactly one and its Kummer class
is globally primitive. Its logarithmic unit coefficient
has valuation \(e\). Raising it to \(5^r\) gives
the required \(r,N=r+e\), and the independent five-power
factor supplies \(m\). For \(N=\infty\), use
\(5^r(\lambda_3\kappa(2)-\lambda_2\kappa(3))\),
where \(\lambda_2\equiv2\) and \(\lambda_3\equiv4\pmod5\),
and add \(5^m\kappa(5)\) if desired. All are actual
completed cusp parameters, with the full cusp construction
supplying section existence. D.74 supplies the asserted
finitely ramified fibre statement.

The literal scalar matrices, norm relations, actual
source-to-target maps and diagonal gauge are checked in
\path{code/de_sha_all_cusp_parameters.py}.
The independent stabilizer calculation is retained in
\path{notes/pbc_middle_boundary_derivation.md} and its
exact certificate. Those finite computations support the
formulas; the all-depth compact and arithmetic arguments
are the proof above.
\end{proof}


% D.92. Root and PBC complete independent proof/dependency audits PASS.
\subsection{The exact metabelian abelian image on the global rank-one locus}

\begin{statement}{D.92}{Proved}{A complete abelian-image criterion on every global valuation line}
Let \(M=F^{(5)}/(F^{(5)})''\) carry the actual unit-tangent
arithmetic action. Write the unique completed Kummer coordinates
of an abelian pair as
\[
 (a,b)=\sum_{q\text{ prime}}\mathbf v_q\kappa(q),
 \qquad \mathbf v_q=(a_q,b_q)\in\mathbb Z_5^2,
 \qquad \mathbf v_q\longrightarrow0,
                                                               \tag{D.92.1}
\]
and put
\[
 d(a,b)=\dim_{\mathbb Q_5}
       \operatorname{span}_{\mathbb Q_5}\{\mathbf v_q:q\text{ prime}\}.
                                                               \tag{D.92.2}
\]
On the locus \(d(a,b)\le1\), the exact image of
\(\operatorname{Sec}(M)\) in abelian Kummer cohomology is
\[
 \begin{aligned}
 &\{(\eta,0),(0,\eta),(\eta,\eta):
          \eta\in H^1(G_{\mathbb Q},\mathbb Z_5(1))\}\\
 &\qquad{}\cup
 \{(A\kappa(5),B\kappa(5)):(A,B)\in\mathbb Z_5^2\}.
 \end{aligned}                                                \tag{D.92.3}
\]
In particular, a nonperipheral global valuation line can
support an actual metabelian section only when every
away-five coefficient is zero. Every pair in (D.92.3)
does have an actual section.

The entire separated compact second-cohomology kernel on
every nonempty fibre in this rank-one locus is therefore
computed: D.91 applies on the three peripheral families,
and it is zero on the pure-five family by D.83. These
descriptions agree on their intersections. This statement
does not determine the first-cohomology rank at every
infinitely ramified peripheral parameter, or the first
obstruction to a higher-derived lift.

Among ordinary rational points of the tripod, exactly
\[
                         t=-1,\quad 2,\quad \tfrac12
                                                               \tag{D.92.4}
\]
have \(d(\kappa(t),\kappa(1-t))\le1\). Their abelian
images are respectively
\[
       (0,\kappa(2)),\qquad(\kappa(2),0),\qquad
                         (-\kappa(2),-\kappa(2)).
                                                               \tag{D.92.5}
\]
No assertion identifies any of these point-section classes
with a cusp-section class having the same abelian image.
Every section above a pure-five pair off the three
peripheral lines is neither cuspidal nor the section of
an ordinary rational point, as already established in D.36.
\end{statement}

\begin{proof}
The Kummer coordinate space is
\(c_0(\{q:q\text{ prime}\},\mathbb Z_5)\):
\(\mathbb Q^\times=\{\pm1\}\times\bigoplus_q\mathbb Z q\),
and the factor \(\{\pm1\}\) has trivial pro-five
completion. Taking the inverse limit modulo \(5^n\)
gives precisely the families that are finitely supported
modulo every \(5^n\), equivalently the displayed \(c_0\)
condition. This proves uniqueness in (D.92.1).

If \(d(a,b)=1\), choose a primitive integral vector
\(\mathbf w=(w_1,w_2)\) on its one-dimensional
\(\mathbb Q_5\)-span. Then
\[
                  \mathbf v_q=c_q\mathbf w,
                   \qquad c_q\in\mathbb Z_5,
                   \qquad c_q\longrightarrow0.                \tag{D.92.6}
\]
Indeed one coordinate of \(\mathbf w\) is a unit,
so the corresponding coordinate of \(\mathbf v_q\)
determines an integral \(c_q\), with the same convergence
property. Thus \((a,b)=(w_1\eta,w_2\eta)\), where
\(\eta=\sum_qc_q\kappa(q)\) is an actual completed
Kummer class. No choice of a rational generator of the
line, nor a finite ramification hypothesis, is used.

Suppose that the pair is the abelianization of an actual
continuous metabelian section and that
\(\mathbf v_q\ne0\) for some \(q\ne5\).
Restrict that section to a decomposition group at \(q\).
On a chosen five-primary tame inertia generator, its
abelian coordinate is \(\mathbf v_q\), up to the common
unit used to identify \(\mathbb Z_5(1)\) with \(\mathbb Z_5\).
This follows directly from the valuation description of
local Kummer theory: \(\kappa(q)\) has primitive inertia
coordinate, whereas the Kummer classes of the other
rational primes have zero inertia coordinate. The
arithmetic unit-tangent action and the cyclotomic character
are unramified at \(q\ne5\); conjugating a section cannot
alter this abelian inertia vector.

The full local theorem D.42 consequently applies to this
\emph{actual} restricted section. Its nonzero inertia vector
lies on one of
\[
          \mathbb Q_5(1,0),\qquad
          \mathbb Q_5(0,1),\qquad
          \mathbb Q_5(1,1).                                  \tag{D.92.7}
\]
Since the nonzero \(\mathbf v_q\) spans the same line as
every \(\mathbf v_{q'}\), including \(\mathbf v_5\),
the whole pair belongs to the corresponding family
\((\eta,0)\), \((0,\eta)\), or \((\eta,\eta)\).
This is a global necessity deduced from one nonzero
away-five inertia coordinate. It does not extrapolate
local section existence into global section existence.

If there is no such prime, then all away-five coordinates
vanish and the pair is exactly
\((A\kappa(5),B\kappa(5))\). The rank-zero pair is
\((0,0)\) and is included in both descriptions.
Conversely, the three actual complete cusp packets of
D.35 supply sections at every peripheral parameter,
including the three zero-parameter packets. D.36 supplies
an actual metabelian section above every pure-five pair.
This proves both directions of (D.92.3).

The compact-kernel assertion is now a case distinction
using the already proved natural maps. D.91 covers every
peripheral parameter, with all finite and infinite
divisibilities. D.83 gives \(\mathcal S_s=\mathcal K_s=0\)
when the away-five lattice is zero. The pure-five
peripheral intersection is the zero-away case in D.91,
so the two results agree. The argument supplies no
additional first-cohomology or higher-derived conclusion.

Finally let \(t\in\mathbb Q\setminus\{0,1\}\).
Its actual point section has abelian coordinates
\((\kappa(t),\kappa(1-t))\). If these are supported
only at five, then
\(t,1-t\in\{\pm5^n:n\in\mathbb Z\}\); reducing
in \(\mathbb Z[1/5]\) modulo two contradicts
\(t+(1-t)=1\). Thus a point in the rank-one locus
must lie on a peripheral line by (D.92.3).
For rational \(x\ne0\), the equality \(\kappa(x)=0\)
is equivalent to \(x\in\{1,-1\}\), because every
prime valuation is zero. The first coordinate line
therefore gives \(1-t=-1\), hence \(t=2\); the
second gives \(t=-1\). On the diagonal,
\(\kappa(t/(1-t))=0\). The sign \(+1\) gives
\(t=1/2\), whereas the sign \(-1\) would give
\(0=-1\). The excluded alternatives \(t=0,1\)
do not give tripod points. Conversely, direct substitution
gives (D.92.5), proving (D.92.4). This calculation is
solely a calculation of abelian images. D.36 gives the
last assertion about the nonperipheral pure-five fibres.
\end{proof}


% D.93 integrated after full root, FGH and PBC proof/code audits of the original theorem and all-depth Smith addendum.
\subsection{Unbounded finite compact kernels on actual ordinary fibres}

\begin{statement}{D.93}{Proved}{All finite compact kernels in an ordinary-point family}
For every integer \(r\ge1\), let
\[
 t_r=2^{5^r},\qquad
 s_r=\bigl(\kappa(t_r),\kappa(1-t_r)\bigr),
 \qquad n_r=5^r.                                      \tag{D.93.1}
\]
This is the abelian image of the actual rational-point section
of the tripod at \(t_r\). Retain the compact groups
\(\mathcal S_s\) and \(\mathcal K_s\) of D.63.
On the free \(\mathbb Z_5\)-module with basis
\(e_j\), \(j\in\mathbb Z/n_r\), define the two literal
permutation matrices
\[
 T_re_j=e_{j+1},\qquad
 \Pi_re_j=e_{2j},
 \qquad
 M_r=2(I+T_r)^3\Pi_r.                                 \tag{D.93.2}
\]
Then the entire separated compact kernel has the exact finite
presentation
\[
 \begin{gathered}
 \mathcal K_{s_r}\simeq\mathcal S_{s_r}
 \simeq\operatorname{coker}_{\mathbb Z_5}(M_r-I)
 \simeq
 \left(\bigoplus_{j=1}^{r-1}\mathbb Z/5^j\right)
       \oplus\mathbb Z/5^{r+1},\\
 \#\mathcal S_{s_r}=5^{\,1+r(r+1)/2}.
 \end{gathered}                                      \tag{D.93.3}
\]
In particular these finite orders are unbounded on actual
ordinary rational-point fibres. The first two integral Smith
calculations give
\[
 \mathcal K_{s_1}\simeq\mathbb Z/25,\qquad
 \mathcal K_{s_2}\simeq
        \mathbb Z/5\ \oplus\ \mathbb Z/125.             \tag{D.93.4}
\]
Equation (D.93.2) is a fully explicit finite integral matrix
for every \(r\). The all-depth elementary divisors follow from
the integral norm-extension proof below; the two small Smith
calculations are independent checks.

Every indicated metabelian fibre is nonempty and has translation
group \(\prod_{\mathbb N}\mathbb Z_5\), with its compact product
topology. Exactly one ordinary rational point has its abelian
image, namely \(t_r\); no cusp has that abelian image. Thus every
other class in this metabelian fibre is neither a rational-point
class nor a cusp class. This is a statement about actual
geometrically pro-five metabelian sections. It does not classify
their higher-derived lifts or general abelian-pair existence.
\end{statement}

\begin{proof}
\paragraph{The global lattice and the exact local depths.}
Write \(B_r=2^{5^r}-1>0\). Since the sign has zero pro-five
Kummer class, the pair in (D.93.1) is
\[
        s_r=\bigl(5^r\kappa(2),\kappa(B_r)\bigr).
                                                               \tag{D.93.5}
\]
The integer \(B_r\) is odd and prime to five. It is not a
fifth power: put \(h=2^{5^{r-1}}\), so that
\((h-1)^5<h^5-1<h^5\). Consequently at least one prime
valuation of \(B_r\) is a unit in \(\mathbb Z_5\).
The first coordinate is supported only at two, and two does
not divide \(B_r\). Hence the actual closed away-five lattice is
\[
        L=5^r\mathbb Z_5(1,0)\oplus\mathbb Z_5(0,1),
             \qquad (a_5,b_5)=(0,0).                  \tag{D.93.6}
\]
This uses the complete valuation lattice, not only its reduction
modulo five. That reduction has rank one, so the earlier
residual-rank-two results do not give (D.93.3).

For a rational five-adic unit \(x\), use the integral coefficient
\[
                   \lambda_x=\frac{\log(x^4)}{4\log6}.
\]
It is the coefficient of \(\kappa(6)\) in the local unit class.
The denominator has valuation one. Direct reduction modulo
twenty-five gives
\[
 \begin{aligned}
 \lambda_2&\equiv2\pmod5,&
 2^{5^r}&\equiv7\pmod{25},\\
 1-2^{5^r}&\equiv-6\pmod{25},&
 \lambda_{1-2^{5^r}}&\equiv1\pmod5.
 \end{aligned}                                      \tag{D.93.7}
\]
Indeed \(5^r\equiv5\pmod{20}\), while two has order twenty
modulo twenty-five. Also
\(16-1\equiv15\pmod{25}\) and \(6^4-1\equiv20\pmod{25}\),
which give the two logarithmic congruences.
Put
\[
 \alpha=\lambda_2,\qquad
 \beta=\lambda_{1-2^{5^r}},\qquad
 \mathbf z=(5^r\alpha,\beta),\qquad
 \gamma=-\alpha/\beta\in\mathbb Z_5^\times .
                                                               \tag{D.93.8}
\]
These are the actual local coefficients in D.87.
The three global depths of \(a,b,a-b\) are respectively
\((r,0,0)\), and their three local depths are also \((r,0,0)\).
The first assertion follows from (D.93.6); the second follows
from (D.93.7), including the unit \(\beta-5^r\alpha\).
D.77 therefore gives
\[
                       \mathcal K_{s_r}\simeq
                       \mathcal S_{s_r}.              \tag{D.93.9}
\]
There is no discarded local free factor in this equality.

\paragraph{The actual image comparison and its natural map.}
Use \(\Lambda=\mathbb Z_5[[u,v]]\), \(t_1=1+u\),
\(t_2=1+v\), \(w=t_1t_2-1\), and \(P=uvw\).
Let \(I_0\) be the actual untwisted Alexander operator image,
including its cyclotomic coordinate \(\chi\), and let
\(N_0=\ker(\chi:I_0\to\mathbb Z_5^\times)\).
There is no valuation graph here, since (D.93.6) has zero
five-coordinate. D.87 supplies the actual common-complement
images
\[
              I_s=L\rtimes I_0,\qquad
              I_5=\mathbb Z_5\mathbf z\rtimes I_0 .
                                                               \tag{D.93.10}
\]
It also identifies the finite-localization kernel with its
natural coefficient map. In particular one finite coefficient
level suffices for all the local trivializations used in this
identification; no unrestricted locally trivial group is being
substituted for \(\mathcal S_s\).

Set \(q=(1+u)^{5^r}-1\). The local coefficient quotient and
translation ideal are
\[
 R=\Lambda/(t_1^{5^r\alpha}t_2^\beta-1),
                    \qquad J=(q,v)R.                 \tag{D.93.11}
\]
Since \(\beta\) is a unit, the first relation eliminates \(t_2\)
integrally. Thus, as actual topological rings,
\[
 R=\mathbb Z_5[[u]],\qquad
 t_2=t_1^{5^r\gamma},\qquad
 v=(1+q)^\gamma-1,\qquad J=qR.                         \tag{D.93.12}
\]
The series \(v/q\) has unit constant coefficient \(\gamma\).
Likewise
\[
 w=t_1^{1+5^r\gamma}-1
\]
is \(u\) times a unit, since \(1+5^r\gamma\) is a unit.
Consequently the ideal generated by the actual \(P\) in \(R\) is
\[
                              PR=u^2qR.               \tag{D.93.13}
\]
No rationalization or splitting into cyclotomic characters
has been used in these integral identities.

For every \(g\in N_0\), the exact beta boundary identities
of D.39 give \(\mathcal B_g-1\in(P)\). The primitive actual
local cubic element in D.77 has
\(\mathcal B_\delta-1=P E_\delta\), with \(E_\delta\) a unit.
Hence the ideal of all these actual operator differences is
exactly \((P)\), on \(R\) and on the invariant ideal \(J\).
Ordinary duality applied to D.87.4 therefore gives
\[
 \mathcal S_s=
 \ker\left\{
       (J/PJ)(-1)_\Gamma\longrightarrow
       (R/PR)(-1)_\Gamma\right\},
                    \qquad \Gamma=\mathbb Z_5^\times .
                                                               \tag{D.93.14}
\]
The map is still induced by the actual inclusion \(J\subset R\).
The twist \((-1)\) records precisely the Cartier-to-ordinary
duality convention. These are compact coinvariants, with closed
images, and (D.93.14) is a map of compact modules.

\paragraph{Retain the generator before taking the two low jets.}
Scalar division by the retained generator \(q\) identifies
\(J/PJ\) with
\[
                         T=R/(u^2q).
\]
On this scalar module the actual Cartier-untwisted action is
\[
        \chi(g)^{-1}S_{\chi(g)}(t_1^{5^r})\rho_0(g),
          \qquad S_c(t)=\frac{t^c-1}{t-1},             \tag{D.93.15}
\]
where \(\rho_0\) includes its substitution \(t_i\mapsto
t_i^{\chi(g)}\). The natural map in (D.93.14) is multiplication
by \(q\), not the identity on two isomorphic scalar quotients.

The quotient \(T/u^2T=R/(u^2)\) is free of rank two over
\(\mathbb Z_5\). Its successive actual characters are
\(\chi^2,\chi^3\): the Alexander multiplier has augmentation
\(\chi^2\), the retained \(q\) contributes \(\chi\), the
Cartier untwist removes one factor \(\chi\), and the linear
substitution supplies the next factor. Both characters are
nontrivial on the order-four tame subgroup. Exact averaging
by that subgroup shows that every continuous
\(\Gamma\)-homology group of this two-jet quotient is zero.
The compact exact homology sequence consequently gives an
isomorphism
\[
                 (u^2T)_\Gamma\xrightarrow{\ \sim\ }T_\Gamma.
                                                               \tag{D.93.16}
\]
This proves both surjectivity and injectivity; no claim that a
coinvariant functor is generally left exact is needed.

Since \(R\) is a domain, scalar division by \(u^2\) gives
\[
                u^2T\simeq R/(q)
                    \simeq\mathbb Z_5[C_{5^r}].        \tag{D.93.17}
\]
On this quotient \(t_1^{5^r}=t_2=1\).
The exact untwisted Alexander action is
\[
 \rho_0(g)F=
 S_\chi(t_1)S_\chi(t_2)\mathcal B_g(t_1,t_2)
                 F(t_1^\chi,t_2^\chi).
\]
The boundary identity \(\mathcal B_g(t,1)=1\) is exact for
every multiplier. Combining it with (D.93.15) and the
substitution on the retained \(u^2\) gives on (D.93.17)
\[
                       F(t)\longmapsto
                  \chi\,S_\chi(t)^3 F(t^\chi).         \tag{D.93.18}
\]
Thus neither an unspecified Frobenius constant nor an
approximate weight action remains.

Multiplication by \(q\) sends the actual \(u^2\)-layer to
zero in \(R/PR\), since \(u^2q=0\) there. By (D.93.16)
the entire natural map in (D.93.14) is zero.
It follows that
\[
             \mathcal S_{s_r}\simeq
       \bigl(\mathbb Z_5[C_{5^r}],
                  \chi S_\chi(t)^3\sigma_\chi\bigr)_\Gamma .
                                                               \tag{D.93.19}
\]
The element two topologically generates \(\mathbb Z_5^\times\):
its order modulo five is four and
\(v_5(2^4-1)=1\). Its actual action in (D.93.18) is precisely
the matrix \(M_r\) of (D.93.2). The image of \(M_r-I\)
on this finite free compact module is closed. Density of the
cyclic subgroup generated by two therefore identifies its
cokernel with the full compact \(\Gamma\)-coinvariants.
This proves the finite presentation in (D.93.3).

\paragraph{The all-depth determinant.}
Compute the determinant after adjoining the \(5^r\)-th roots
of unity in characteristic zero; this does not change it.
The evaluation at \(t=1\) has scalar \(16\).
For each \(1\le j\le r\), multiplication of the exponents by
two is one orbit on the primitive \(5^j\)-th roots, of length
\(\ell_j=4\cdot5^{j-1}\). On that orbit the product of the
weights in (D.93.18) is
\[
 2^{\ell_j}\prod_{k=0}^{\ell_j-1}
       (1+\zeta^{2^k})^3
 =2^{\ell_j},
 \qquad
 1+\zeta^{2^k}
     =\frac{\zeta^{2^{k+1}}-1}{\zeta^{2^k}-1}.          \tag{D.93.20}
\]
The product telescopes because the orbit closes and none
of these primitive roots is one. For a weighted permutation
cycle the determinant of identity minus the cycle is one
minus its weight product. Consequently
\[
 \det(I-M_r)
      =(1-16)\prod_{j=1}^r
                 (1-2^{\,4\cdot5^{j-1}}).             \tag{D.93.21}
\]
All factors are nonzero. The elementary lifting-the-exponent
calculation gives
\[
 v_5(16-1)=1,\qquad
 v_5(2^{\,4\cdot5^{j-1}}-1)=j.
 \]
For example this follows inductively by expanding
\((1+5c)^5-1\), which increases the valuation by exactly
one when \(c\) is a unit. Thus the cokernel is finite of the
order in (D.93.3), with its finite compact topology.
The exact literal integral Smith forms for \(r=1,2\) are
executed in the accompanying certificate and give (D.93.4).
Those two executions are not used to prove (D.93.21).

\paragraph{An exact extension by one order-five factor.}
Put \(n=5^r\), \(R=\mathbb Z_5[t]/(t^n-1)\), \(u=t-1\),
and \(I=uR\). Let \(\sigma_\chi(t)=t^\chi\).
On the same underlying free module \(R\), distinguish
\[
 \mathscr A:\ \chi\sigma_\chi,\qquad
 \mathscr B:\ \chi S_\chi(t)^3\sigma_\chi,
       \qquad S_\chi(t)=\frac{t^\chi-1}{t-1}.
\]
The module \(\mathscr B_\Gamma\), with
\(\Gamma=\mathbb Z_5^\times\), is exactly the compact kernel
in D.93. Write
\[
                  \mathcal N=1+t+\cdots+t^{n-1}.
\]
Multiplication by \(u^3\) gives the exact equivariant sequence
\[
 0\longrightarrow \mathbb Z_5\mathcal N(\chi^4)
   \longrightarrow \mathscr B
   \xrightarrow{u^3} I^3\mathscr A
   \longrightarrow0.                                \tag{D.93.22}
\]
Indeed the equivariance is
\(u^3 S_\chi(t)^3=\sigma_\chi(u^3)\).
The kernel of multiplication by \(u^3\) in \(R\) is precisely
\(\mathbb Z_5\mathcal N\): over a characteristic-zero splitting
field it is supported only at \(t=1\), hence is killed by \(u\);
the coefficient equations for \(u f=0\) then say that all
coefficients of \(f\) are equal. This also proves the integral
kernel assertion. The action on this generator is \(\chi^4\).

The quotient \(\mathscr A/I^3\mathscr A\) has the three
successive tame characters \(\chi,\chi^2,\chi^3\), on its
augmentation filtration. These assertions include any
five-torsion in its two positive-degree pieces. All three
characters are nontrivial on the order-four subgroup of
\(\Gamma\). Exact tame averaging therefore kills every
\(\Gamma\)-homology group of the quotient. Consequently
\[
 (I^3\mathscr A)_\Gamma\simeq\mathscr A_\Gamma,
 \qquad H_1(\Gamma,I^3\mathscr A)=H_1(\Gamma,\mathscr A)=0.
                                                               \tag{D.93.23}
\]
For the last equality, the completed procyclic resolution
after tame averaging identifies \(H_1\) with invariants.
The action of two on \(\mathscr A\) is twice a finite
permutation, so its fixed module is zero. All modules here
are finite free compact \(\mathbb Z_5\)-modules before the
displayed finite quotients; the compact exact resolution is
being used with its actual topology.

Decompose the basis \(t^a\) of \(\mathscr A\) into its orbits
under \(a\mapsto2a\) in \(\mathbb Z/n\).
The zero orbit has weight two and contributes nothing.
For each \(1\le j\le r\), the orbit of
\(t^{5^{r-j}}\) has length \(\ell_j=4\cdot5^{j-1}\);
its weighted-cycle cokernel is
\(\mathbb Z_5/(2^{\ell_j}-1)\simeq\mathbb Z/5^j\).
The exact homology sequence of (D.93.22) now gives
\[
 0\longrightarrow\mathbb Z/5\,[\mathcal N]
 \longrightarrow\mathscr B_\Gamma
 \longrightarrow\bigoplus_{j=1}^r\mathbb Z/5^j
 \longrightarrow0.                                  \tag{D.93.24}
\]
In particular the norm class has exact order five. It is
not discarded by a rational calculation.

\paragraph{An explicit lift of the largest orbit generator.}
Let \(\sigma=\sigma_2\), and write
\[
 A=2\sigma,\qquad B=2(1+t)^3\sigma
\]
for the actual operators on \(\mathscr A,\mathscr B\).
The integral-at-five polynomials
\[
 G=-1-\frac{u}{3}+\frac{2u^2}{21},
       \qquad f=\frac4{21}(t+3)
\]
satisfy the literal polynomial identity
\[
                 u^3 f=t+(A-1)G.                     \tag{D.93.25}
\]
Thus \([f]\) maps in (D.93.24) to the largest orbit
generator \([t]\), up to the stated natural identification.

Put
\[
 \ell=4\cdot5^{r-1},\qquad
 c=2^\ell-1,\qquad
 Q=\sum_{i=0}^{\ell-1}2^i t^{2^i},\qquad H=Q+cG .
                                                               \tag{D.93.26}
\]
Here \(v_5(c)=r\) and \(c/n\equiv3\pmod5\).
The weighted orbit sum gives
\[
                       (A-1)H=c\,u^3f
                          \quad\hbox{in }R.           \tag{D.93.27}
\]
Indeed \(A^\ell t=2^\ell t^{2^\ell}=(1+c)t\) in
\(R\), since \(n\mid c\). This argument is in the actual
finite group ring; it does not set \(t^{2^\ell}=t\)
in the unrestricted polynomial ring.

We give an integral lift of \(H\) through multiplication
by \(u^3\). Write \(h_i\) for its coefficient of \(u^i\)
in the unrestricted polynomial ring. Its constant
coefficient is zero. For
\[
        S_k=\sum_{i=0}^{\ell-1}2^{ki}
               =\frac{(1+c)^k-1}{2^k-1},
\]
its first three coefficients are
\[
 \begin{aligned}
 h_1&=\frac{c+c^2}{3},\\
 h_2&=-\frac c{42}+\frac{c^2}{21}+\frac{c^3}{14},\\
 h_3&=\frac{S_4-3S_3+2S_2}{6}.
 \end{aligned}                                       \tag{D.93.28}
\]
The first two lie in \(n\mathbb Z_5\). Therefore
\[
 W_0=h_1/n,\qquad
 W_1=\bigl(h_2-\tbinom n2W_0\bigr)/n
                       \quad\hbox{belong to }\mathbb Z_5,
\]
and there is an actual polynomial \(Z\in\mathbb Z_5[u]\)
with
\[
 H-(t^n-1)(W_0+W_1u)=u^3Z.                            \tag{D.93.29}
\]
Since \(v_5(\binom n2)=v_5(\binom n3)=r\), its
constant coefficient \(z_0=Z|_{u=0}\) satisfies
\[
                         z_0\equiv h_3\pmod{5^r}.      \tag{D.93.30}
\]

\paragraph{The norm extension is nonzero on the largest factor.}
Equations (D.93.27)--(D.93.29), and the integral kernel
calculation in (D.93.22), give
\[
             c f-(B-1)Z=k\mathcal N
                   \quad\hbox{in }R,\quad k\in\mathbb Z_5.
                                                               \tag{D.93.31}
\]
Evaluation at \(t=1\) yields the exact equality
\[
                 nk=c\,\frac{16}{21}-15z_0.            \tag{D.93.32}
\]
We compute \(k\) modulo five. By (D.93.30), replacing
\(z_0\) by \(h_3\) changes \(15z_0/n\) by a multiple
of five. The terms \(S_2,S_3\) in (D.93.28) have
valuation at least \(r\), so they disappear after this
operation modulo five. The \(S_4\) term gives exactly
\[
 \frac{15h_3}{n}
   \equiv\frac{(1+c)^4-1}{6n}
   \equiv\frac{2c}{3n}\pmod5.
\]
All expressions compared here are integral at five.
Consequently
\[
       k\equiv\frac cn\left(\frac{16}{21}-\frac23\right)
        =\frac{2c}{21n}\equiv1\pmod5.                 \tag{D.93.33}
\]
Thus \(c[f]\) is a nonzero norm class in (D.93.24).
Since \(v_5(c)=r\), the lift \([f]\) has exact order
\(5^{r+1}\). This proves the required nonsplitting over
the largest cyclic quotient.

\paragraph{Splitting every smaller factor.}
Rescale this lift by a five-adic unit to obtain \(x_r\)
with \(5^r x_r=[\mathcal N]\). Choose arbitrary lifts
\(x_j\) of the smaller orbit generators, \(1\le j<r\).
For some \(a_j\in\mathbb Z/5\), they satisfy
\(5^jx_j=a_j[\mathcal N]\).
Choose integral representatives of \(a_j\) and replace
\[
                     x_j\ \longmapsto\
                     x_j-a_j5^{r-j}x_r .
 \]
The new lift has order exactly \(5^j\). On the quotient
in (D.93.24), these replacements are invertible triangular
basis changes: the \(r\)-th generator is kept, and each
smaller one is changed only by a permissible multiple of
that generator. The modified lifts and \(x_r\) generate
the whole module. Their product of orders equals the
order supplied by (D.93.24). They therefore give precisely
the direct-sum decomposition (D.93.3).

This closes the all-depth elementary-divisor calculation.
The literal symbolic identities and moment checks are retained
in \path{code/de_ordinary_kernel_smith.py}.
The small integral matrices are retained in
\path{code/de_ordinary_finite_kernel.py}.
The extension argument above supplies the proof for every \(r\).

\paragraph{The actual section fibres and their ordinary points.}
The rational point \(t_r\) supplies a full section, hence an
actual metabelian class above \(s_r\). The pair has finite
prime support. D.74 therefore identifies the translation
group of its metabelian fibre with
\(\prod_{\mathbb N}\mathbb Z_5\).
The two Kummer coordinates are linearly independent over
\(\mathbb Q_5\), by their disjoint nonempty prime supports.
They lie on none of the three completed cusp lines in D.35.

For a nonzero rational number \(x\), its pro-five Kummer
class vanishes exactly when \(x=\pm1\); this follows directly
from its prime valuations. If an ordinary rational point
\(z\) has abelian pair \(s_r\), it follows that
\[
                  z=\pm t_r,\qquad
                  1-z=\pm(1-t_r).
 \]
The choice \(z=-t_r\) makes the second equality impossible
in either sign. Hence \(z=t_r\), as claimed. These statements
classify the ordinary point images in this fibre, not the
full section classes of the geometric fundamental group.
\end{proof}

